\documentclass[11pt,a4paper]{article}
\usepackage[T1]{fontenc}
\usepackage[utf8]{inputenc}
\usepackage{lmodern}
\ifdefined\pdfmapfile
  \pdfmapfile{=lm.map}
  \pdfmapfile{+cm.map}
  \pdfmapfile{+symbols.map}
  \pdfmapfile{+euler.map}
  \pdfmapfile{+rsfs.map}
\fi
\usepackage[a4paper,margin=27mm]{geometry}
\usepackage{amsmath,amssymb,amsthm,mathtools,mathrsfs}
\usepackage{booktabs,tabularx,longtable,array,enumitem}
\usepackage{microtype,xcolor,etoolbox}
\usepackage{hyperref}
\hypersetup{colorlinks=true,linkcolor=blue!45!black,citecolor=blue!45!black,
 urlcolor=blue!45!black,pdfauthor={Research record prepared with ChatGPT},pdftitle={Relative Bedford--Taylor theory on Berkovich spaces: third consolidated research record}}
\usepackage[capitalise,noabbrev]{cleveref}
\numberwithin{equation}{section}
\newtheorem{theorem}{Theorem}[section]
\newtheorem{lemma}[theorem]{Lemma}
\newtheorem{proposition}[theorem]{Proposition}
\newtheorem{corollary}[theorem]{Corollary}
\theoremstyle{definition}
\newtheorem{definition}[theorem]{Definition}
\newtheorem{example}[theorem]{Example}
\theoremstyle{remark}
\newtheorem{remark}[theorem]{Remark}
\newcommand{\Cb}{\mathcal C_b}
\newcommand{\R}{\mathbf R}
\newcommand{\Q}{\mathbf Q}
\newcommand{\Z}{\mathbf Z}
\newcommand{\PSH}{\operatorname{PSH}}
\newcommand{\BTcap}{\operatorname{Cap}}
\newcommand{\supp}{\operatorname{supp}}
\newcommand{\dd}{d'd''}
\newcommand{\an}{\mathrm{an}}
\setlist{topsep=4pt,itemsep=2pt,parsep=0pt}
\setlength{\parindent}{1.3em}
\setlength{\parskip}{2pt}
\setlength{\emergencystretch}{2em}
\makeatletter
\patchcmd{\l@section}{1.5em}{2.1em}{}{}
\renewcommand*\l@subsection{\@dottedtocline{2}{1.5em}{3.6em}}
\makeatother

\title{Relative Bedford--Taylor theory on Berkovich spaces\\[5pt]
\large Third research record: energy, polar extension, relative capacity,\\
Dirichlet existence and stability, and obstructions on intersecting supports}
\author{Research record prepared with ChatGPT\\
\small Based on the authorized manuscripts R, B, and A2}
\date{10 November 2026}
\begin{document}
\maketitle
\begin{abstract}
The bounded calculus for every fixed original CLD-positive separately
closed background is retained from the preceding records. This continuation
proves an all-degree mixed-energy estimate controlling relative capacity,
strong convergence tested by normalized classical psh functions, and
simultaneous tail-convex-block compactness in the relative completion.
It constructs a polar decomposition and proves finite-mass zero extension
across closed locally complete pole sets. Analytic separation of carriers
gives exact capacity additivity; a uniform CLN tail estimate extends
regularity to order limits of backgrounds. On specified extendible
exhaustion domains, compact bounded-subsolution data and Dini-capacity data
have unique ambient zero-boundary solutions, with continuity for Dini data.
A radial construction gives a matching universal obstruction for every
non-Dini modulus. Arbitrary-relative classical Dirichlet existence is
false even with continuous matching subsolutions and power-capacity bounds:
intersecting analytic supports force incompatible actual branch values.
The relative completion supplies explicit solutions and a finite-branch
gluing theorem. Every result retains the original finite $p$-adic field,
its absolute value, and the full original CLD smooth test category.
General relative raw-capacity regularity, the unrestricted potential
quotient, and general boundary and relative completed Dirichlet problems
remain open.
\end{abstract}

\section*{Scope and reading guide}
\addcontentsline{toc}{section}{Scope and reading guide}
The current theorem correspondence and the comparison with R2 are in
Section~\ref{sec:r3-status}. Part I keeps the mathematical text
from the frozen R1/R2 record verbatim and preserves its numbering.
Those historical status tables are superseded only by the particular
qualified result cited in Part II. This record does not declare the
whole Bedford--Taylor program complete.

The internal manuscripts R and B are accepted inputs in their actual
categories; their proofs have not been comprehensively re-audited.
A2's same-model geometric adelic nefness theorem is already complete.
The provenance labels \texttt{INTERNAL\_INPUT},
\texttt{EXTERNAL\_KNOWN}, \texttt{DIRECT\_COROLLARY}, and
\texttt{NEW\_PROVED} are distinct from the logical statuses
\texttt{PROVED}, \texttt{CONDITIONAL}, \texttt{OPEN}, and
\texttt{FALSE}. A failed route is marked \texttt{FAILED\_ATTEMPT}.
Here \texttt{NEW\_PROVED} records a complete proof in this research
project, not an assertion of worldwide priority or external refereeing.

An arbitrary background is allowed in the fixed-background calculus,
energy, and stated extension theorems. Ambient Dirichlet theorems have
their own domain and measure hypotheses. Ordinary-current compatibility
is not equality on all $\delta$-tests. Raw capacity is not identified
with its open outer regularization without a theorem establishing it.

The accompanying archive records exact source versions, internal checks,
the recovery of files after an execution-environment reset, and the UTC
working intervals. Parallel branches are not added to elapsed time.
Pure proof time was not separately instrumented. Reconstructed source
versions are identified and checked as such rather than assigned their
earlier file hashes.
\tableofcontents
\clearpage
\part*{Part I. Frozen R1 and R2 mathematical dependencies}
\addcontentsline{toc}{part}{Part I. Frozen R1 and R2 mathematical dependencies}

% BEGIN INCLUDED SOURCE: base_r2/base_r1/sections/inputs
\section{Scope, inputs, and conventions}
\label{sec:inputs}

Fix a finite extension $K/\mathbf Q_p$ with its given absolute value.
Let $X$ be a good, Hausdorff, boundaryless, quasi-smooth $K$-analytic
space of pure dimension $n\geq1$, and let $U\subset X$ be an ordinary open
set. We use only the original Chambert--Loir--Ducros sheaf
$A^{a,b}$, its compactly supported smooth tests, and its original
calibration. We assume $0\leq p\leq n$. A current of bidegree $(p,p)$ acts on
$A_c^{n-p,n-p}$. We write $D=d'd''$.

The class $\Cb(U)$ consists of the actual real-valued, locally bounded,
upper-semicontinuous classical psh functions of \textup{B}, Definition
2.2: pullback to every analytic curve after every complete isometric
extension is subharmonic. No function is changed on a null or polar
set. A norm in this report always means the actual pointwise bound
\[
 \|u\|_{V,\infty}=\sup_{x\in V}|u(x)|,
\]
even where the traditional notation $L^\infty(V)$ is used. There is
no unspecified reference volume measure behind this notation.

\subsection{The accepted internal results}

The user has authorized the following theorems as established inputs.
They have been checked for scope and dependencies, but their proofs
have not been re-audited in this project.

\begin{description}
\item[\textup{R}.] \emph{Continuous CLD-positive potentials and adelic
mixed measures}, 10 November 2026, 19 pages.
Theorem 2.9 gives local uniform smooth positive approximation of finite
continuous original-current-positive potentials on smooth
quasi-projective algebraic analytifications over a $p$-adic completion.
Lemma 3.1 gives finite coefficient Radon control for the original
positive cone and stability under a fixed continuous multiplier.
Lemma 3.2 smooths pure Fubini--Study metrics on their original line.
Proposition A.1 gives regularized maxima with exact inactive branches.
Theorem 2.10 identifies positivity of the prescribed Yuan--Zhang
integrable adelic mixed measure under its stated hypotheses.

\item[\textup{B}.] \emph{Locally bounded classical plurisubharmonic
functions and mixed currents}, 10 November 2026, 32 pages.
Theorems 2.5--2.7 identify $\Cb$ with the local decreasing PL class,
and construct strongly positive mixed currents and weighted joint
decreasing continuity with background $[1]$.
Definition 2.4 and Lemma 7.2 specify actual-space coefficient Radon
measures and Borel multiplication.
Theorem 4.4, Lemma D.2, and Proposition 7.1 provide the common
projective line, frame, ordinary neighborhood, and decreasing pure
Fubini--Study approximants for a finite family of decreasing sequences.
Lemma D.4 supplies ordinary algebraic neighborhoods without changing
$K$. These common-domain inputs are needed in
Lemma~\ref{lem:common-majorants}; separate local regularizations with
shrinking domains would not suffice.

\item[\textup{A2}.] \emph{A non-Archimedean criterion for geometric
adelic nefness}, revised 12 October 2026, 31 pages.
Theorem 1.1 is the completed same-model geometric nefness theorem under
its number-field, integrable adelic line, and full Cartier-boundary
hypotheses. It preserves the open line and its prescribed
identifications. This is an accepted global application, not an open
problem in the present project. It is not used to prove the local
relative BT theorems below.
\end{description}

Any use of the broader analytic setting passes through the local
algebraic neighborhoods already supplied by \textup{B}. A finite
extension of $\mathbf Q_p$ can be realized as a number-field completion;
the use of \textup{R} does not change the prescribed absolute value.
No argument below makes a new base change or renormalizes integration.

\subsection{Three positivity cones}

On a finite original moment presentation, strongly positive forms are
generated by decomposable squares. Positive forms are represented by
positive semidefinite symmetric coefficient tensors. Weakly positive
forms are nonnegative on decomposable vectors. These cones need not
coincide in middle degree. A \emph{positive current} in this report is
nonnegative on positive complementary original forms, as in
\textup{R}, Convention 2.2. A strongly positive current is nonnegative
on weakly positive complementary forms, as in \textup{B}, Convention
2.3. Thus the arbitrary background considered here is not assumed
strongly positive.

We use the finite-matrix domination from \textup{R}, not a decomposition
of every symmetric test into decomposable squares. The latter would
lose a middle-degree annihilator unless one had strong positivity and
used the separate argument in \textup{B}, Lemma 5.1. Wedge with a
positive $(1,1)$-form preserves the original positive cone. More
generally, wedge with a strongly positive form preserves it; these are
the only cone-preserving wedge operations used in the proofs.

\subsection{Meaning of the result labels}

\begin{center}
\begin{tabularx}{\textwidth}{@{}lX@{}}
\toprule
Label & Meaning in this record\\
\midrule
\texttt{EXTERNAL\_KNOWN} & A verified statement of a cited external source.\\
\texttt{INTERNAL\_INPUT} & A theorem accepted from the supplied manuscripts.\\
\texttt{DIRECT\_COROLLARY} & A consequence obtained directly from named accepted inputs or preceding proved results.\\
\texttt{NEW\_PROVED} & A statement with a complete proof in this record, relative to the adopted inputs; this does not claim bibliographic priority or external peer review.\\
\texttt{CONDITIONAL} & A stated implication whose extra hypothesis is not established here.\\
\texttt{OPEN} & An unresolved assertion in the declared scope.\\
\texttt{FALSE} & An assertion ruled out by a specified counterexample.\\
\texttt{FAILED\_ATTEMPT} & A proof route with an identified invalid or missing step.\\
\bottomrule
\end{tabularx}
\end{center}

The separate analogue-status table uses \texttt{PROVED},
\texttt{CONDITIONAL}, \texttt{OPEN}, and \texttt{FALSE}.
Here \texttt{PROVED} includes the user's accepted internal inputs and
proofs derived from them. It does not mean that the input manuscripts
have been independently validated by this research run.

\subsection{What is not being constructed}

The product has an arbitrary \emph{fixed} positive closed background
and bounded classical psh potentials. It is not a wedge product of two
entirely arbitrary positive currents. It is not the full GK
$\delta$-form algebra, and its current identities do not quantify over
all $\delta$-test forms. Later applications to Green forms or
Beilinson--Bloch heights do not enter the proof of the present calculus.

% END INCLUDED SOURCE: base_r2/base_r1/sections/inputs


% BEGIN INCLUDED SOURCE: base_r2/base_r1/sections/relative_cln_agent
% Independent contribution: /root/relative_cln, 10 November 2026.
% Fragment; the master file supplies theorem environments and bibliography.

\section{Ordered products and relative finite-test estimates}
\label{sec:relative-cln-agent}

Throughout this section, $K/\mathbf Q_p$ is finite with its prescribed
absolute value. The space $X$ is good, Hausdorff, boundaryless and
quasi-smooth, of pure dimension $n$. All open sets are ordinary open sets,
and all forms are original CLD forms. We write $D=d'd''$.
The word \emph{positive} means dual positivity for the original positive
form cone used in R, Convention 2.2. It does not mean strong positivity.
The background $S$ is a symmetric positive $(p,p)$-current satisfying
$d'S=d''S=0$.
All $L^\infty$ bounds below mean actual pointwise suprema on the stated
sets; no ambient measure is used to discard exceptional values.

\subsection{The precise order-zero input}

\begin{lemma}[Positive controls on a finite test stage]
\label{lem:finite-positive-controls}
Let $W$ be an ordinary open set, let $L\subset W$ be compact, and fix a
finite original moment presentation of a smooth $(q,q)$-form $\eta$ near
$L$. After replacing $\eta$ by a compactly supported extension agreeing
with it near $L$, there are positive forms
$\Theta_1,\ldots,\Theta_N\in A_c^{q,q}(W)$ and a constant $C_\eta$ such that
every positive $(n-q,n-q)$-current $T$ satisfies
\begin{equation}\label{eq:finite-positive-controls}
 |\mu_{T,\eta}|(L)
 \leq C_\eta\sum_{a=1}^N\langle T,\Theta_a\rangle.
\end{equation}
The same controls work for all $T$. For forms on one fixed finite
presentation, with supports in one fixed compact set, $C_\eta$ is bounded
by a constant times the supremum norm of the coefficient matrices.
\end{lemma}

\begin{proof}
This is the local domination argument of R, Lemma 3.1. We recall the
details that matter for subsequent estimates. A finite smooth partition
of unity reduces the assertion to one moment presentation with real
coordinates $t_1,\ldots,t_r$. Put
\[
 \beta=\sum_{a=1}^r d't_a\wedge d''t_a,
 \qquad \Omega_q=\frac{\beta^q}{q!},
\]
with $\Omega_0=1$. Choose a nonnegative compact smooth cutoff $\rho$,
equal to one near the relevant support. The symmetric coefficient
matrix of $\Omega_q$ is positive definite. Hence, with a constant
$C_\eta$ controlled by the coefficient supremum norm,
\[
 C_\eta\|h\|_\infty\rho\Omega_q\pm h\eta_{\rm sym}\geq0
\]
for every compact smooth scalar $h$. Antisymmetric forms and their
negatives belong to the null part of this positivity convention; a
positive current annihilates them. Therefore
\[
 |\langle T,h\eta\rangle|
 \leq C_\eta\|h\|_\infty\langle T,\rho\Omega_q\rangle.
\]
If $h$ uses further moment coordinates, pull the displayed cone
inequalities to a common refinement. Positivity is preserved, so the
constant does not depend on those additional scalar coordinates.

Original smooth scalars are uniformly dense on compact sets. Riesz
representation thus gives the coefficient Radon measure and the stated
total-variation bound. Summing over the finite partition gives
\eqref{eq:finite-positive-controls}. The local algebraic neighborhoods in
B, Proposition 7.1 and Lemma D.4, let us apply R's argument on the
analytic spaces considered here.
\end{proof}

\paragraph{Status and dependence.}
Lemma~\ref{lem:finite-positive-controls} is
\texttt{DIRECT\_COROLLARY}, from R, Lemma 3.1 and B's local algebraic
transport. In particular it must not be presented as a new positivity
theorem. Strong positivity, and B's decomposable-square annihilator in
middle degree, are not used in this argument.

The scalar coefficient relations hold as identities of Radon measures:
test each relation after multiplication by every original smooth scalar
and use uniqueness of the representing measure. Consequently, for every
locally bounded Borel function $a$,
\[
 \langle aT,\eta\rangle=\int a\,d\mu_{T,\eta}
\]
defines an intrinsic coefficient-Radon current. It commutes with
restriction and with common refinements. The proof is the measure
uniqueness argument of B, Lemma 7.2, applied to the original-positive
coefficient measures supplied by R. It does not replace $a$ by an
almost-everywhere representative chosen using another measure.

\subsection{One multiplication and ordered iteration}

\begin{proposition}[One bounded potential over an arbitrary background]
\label{prop:one-potential-arbitrary-background}
For every $u\in\mathcal C_b(W)$ the current $D(uS)$ is symmetric,
positive, coefficient Radon, and separately $d'$- and $d''$-closed.
If $a_j$ are locally uniformly bounded Borel functions converging
pointwise to $a$, then $a_jS\to aS$ on every original compact test.
In particular,
\[
 D(u_jS)\longrightarrow D(uS)
\]
whenever $u_j\to u$ pointwise and the $u_j$ have a common local bound.
In the last assertion the $u_j$ need not be monotone.
\end{proposition}

\begin{proof}
The convergence assertions are dominated convergence for the fixed
finite coefficient measures of $S$, followed by continuity of
distributional differentiation. For a continuous classical psh
function $v$, the accepted local approximation results of R and B give
smooth psh functions $v_m\to v$ locally uniformly on a fixed local
neighborhood. The identity
\[
 D(v_mS)=Dv_m\wedge S
\]
and preservation of the original positive cone under wedge with a
positive $(1,1)$-form show positivity. Passing to the limit proves
positivity of $D(vS)$. For general $u$, use B, Theorem 2.5 to obtain
continuous semipositive PL functions $u_j\downarrow u$ on a fixed local
neighborhood. They are uniformly bounded on each compact subdomain by
their first term and the locally bounded limit. Dominated convergence
for the fixed coefficient measures of $S$ gives
$D(u_jS)\to D(uS)$. The cone of positive currents is weakly closed.

Closedness follows from $d'D=d''D=0$, and symmetry follows from the
CLD involution. Lemma~\ref{lem:finite-positive-controls} supplies the
coefficient-Radon property. A finite localization proves the assertions
on $W$.
\end{proof}

\begin{corollary}[Canonical ordered products]
\label{cor:ordered-arbitrary-background}
For $0\leq k\leq n-p$, the recursion
\[
 T_{S,0}=S,\qquad
 T_{S,k}(u_1,\ldots,u_k)
 =D\bigl(u_1T_{S,k-1}(u_2,\ldots,u_k)\bigr)
\]
defines a positive, symmetric-as-a-current, coefficient-Radon, separately
closed current for every ordered tuple $u_i\in\mathcal C_b(W)$.
The recursion is intrinsic, multilinear on the cone, positively
homogeneous, unchanged by adding a constant in a differentiated slot,
and compatible with ordinary restriction and gluing.
\end{corollary}

\begin{proof}
Apply Proposition~\ref{prop:one-potential-arbitrary-background}
inductively. Borel multiplication and distributional differentiation
are intrinsic and local. Their linearity proves additivity and positive
homogeneity by induction. A constant in the outer slot contributes
$D(cT)=0$ because the inner current is separately closed; the assertion
for other slots follows inductively. On overlaps, equal coefficient
currents have equal Borel multiples, so the recursion glues.
\end{proof}

Both preceding statements are \texttt{DIRECT\_COROLLARY}. At this
point ``symmetric-as-a-current'' refers only to the CLD involution.
It does \emph{not} assert invariance under permutations of two
discontinuous differentiated potentials. Similarly, intrinsic recursion
does not yet assert convergence along every simultaneous regularization
of all slots. These are separate analytical assertions.

\begin{lemma}[Continuous slots]
\label{lem:continuous-slot-exchange}
Let $u\in\mathcal C_b(W)$ and let $v$ be continuous and classical psh.
Then
\[
 D\bigl(uD(vS)\bigr)=D\bigl(vD(uS)\bigr).
\]
\end{lemma}

\begin{proof}
First take both potentials continuous. Smooth local regularization,
R, Lemma 3.1, and the coefficient mass bounds show that joint locally
uniform convergence of continuous potentials passes through the
recursion. Explicitly, if $Q_j\to Q$ are positive currents and
$a_j\to a$ locally uniformly, with $a$ continuous, then
\[
 a_jQ_j-aQ=(a_j-a)Q_j+a(Q_j-Q)\longrightarrow0.
\]
The first term is controlled by the locally uniform coefficient mass
bound, and the second by R, Lemma 3.1. Smooth commutativity therefore
passes to continuous potentials over the fixed arbitrary background.

Choose continuous psh $u_j\downarrow u$ on one fixed local
neighborhood. The left side of
\[
 D\bigl(u_jD(vS)\bigr)=D\bigl(vD(u_jS)\bigr)
\]
converges by dominated convergence against the \emph{fixed} coefficient
measures of $D(vS)$. On the right, $D(u_jS)\to D(uS)$ by
Proposition~\ref{prop:one-potential-arbitrary-background}; multiplication
by the fixed continuous function $v$ passes to the limit by R,
Lemma 3.1. Differentiate and localize.
\end{proof}

This is again \texttt{DIRECT\_COROLLARY}. No simultaneously varying
discontinuous multiplier is used in its proof.

\subsection{Relative Chern--Levine--Nirenberg bounds}

\begin{theorem}[Finite-test relative CLN estimate]
\label{thm:finite-test-cln}
Fix $K_0\Subset V\Subset W$, an integer $k\leq n-p$, and a finite
original moment presentation of a $(q,q)$-test $\eta$ near $K_0$, where
$q=n-p-k$. There are positive compact tests
\[
 \theta_1,\ldots,\theta_N\in A_c^{n-p,n-p}(V)
\]
and a constant $C$, determined by these fixed data, such that for every
positive separately closed background $S$ and every tuple
$u_1,\ldots,u_k\in\mathcal C_b(W)$,
\begin{equation}\label{eq:relative-cln-product}
 \bigl|\mu_{T_{S,k}(u_1,\ldots,u_k),\eta}\bigr|(K_0)
 \leq C\prod_{i=1}^k\|u_i\|_{L^\infty(V)}
       \sum_{\ell=1}^N\langle S,\theta_\ell\rangle.
\end{equation}
One may replace each norm by $\operatorname{osc}_V(u_i)$.
The empty product is one when $k=0$.

For a locally bounded Borel weight $a$, the same right side multiplied
by $\|a\|_{L^\infty(V)}$ bounds
$|\mu_{aT_{S,k},\eta}|(K_0)$.
The controls and constants are independent of the potentials and of
every regularization parameter. On a fixed finite test stage, the
constant is bounded by a fixed constant times the coefficient
supremum norm of $\eta$.
\end{theorem}

\begin{proof}
Choose nested compact neighborhoods between $K_0$ and $V$, with enough
room for $k+1$ cutoff choices. Lemma~\ref{lem:finite-positive-controls}
first bounds the left side by finitely many evaluations
\[
 C_0\sum_a\langle T_{S,k},\Theta_a^{(0)}\rangle,
 \qquad \Theta_a^{(0)}\in A_c^{q,q}(V),\quad
 \Theta_a^{(0)}\geq0.
\]
All these controls are chosen before the potentials. For every one of
them, the defining recursion gives
\[
 \langle T_{S,k},\Theta_a^{(0)}\rangle
 =\int u_1\,d\mu_{T_{S,k-1}(u_2,\ldots,u_k),D\Theta_a^{(0)}}.
\]
Apply Lemma~\ref{lem:finite-positive-controls} to the fixed smooth test
$D\Theta_a^{(0)}$, using the next compact neighborhood. We obtain
positive tests $\Theta_{ab}^{(1)}$ of degree $(q+1,q+1)$ and constants
independent of the potentials such that
\[
 \langle T_{S,k},\Theta_a^{(0)}\rangle
 \leq C_a\|u_1\|_{L^\infty(V)}
       \sum_b\langle T_{S,k-1},\Theta_{ab}^{(1)}\rangle.
\]
Repeat this operation $k$ times. At each stage only a finite collection
of fixed smooth compact tests is differentiated and controlled.
The final forms have degree $(q+k,q+k)=(n-p,n-p)$, are positive,
and are supported in $V$. Absorb the finitely many constants into $C$
and enumerate the final controls as $\theta_\ell$. This proves
\eqref{eq:relative-cln-product}.

Adding a constant to any differentiated potential leaves the ordered
product unchanged. Subtracting $\inf_Vu_i$ therefore replaces the norm
by its oscillation. Local boundedness and $\overline V\subset W$
make these extrema finite. The weighted estimate follows from
$|a\mu|\leq\|a\|_\infty|\mu|$.

For the last assertion, choose the initial cutoff controls using the
fixed stage and use the matrix norm only in its first domination
constant. All later differentiated controls are then fixed. No
derivative of $\eta$ enters a varying constant. This proves the stated
uniformity.
\end{proof}

\paragraph{Status.}
The estimate is a relative extension needed in this project, with a
complete proof above. Under the user's stipulated inputs its logical
status is \texttt{DIRECT\_COROLLARY}: it follows by finite iteration of
R's fixed-test domination and the ordered recursion. It should not be
advertised as a new independent analytical mechanism.

\paragraph{Charts and refinements.}
The auxiliary forms need not come from a single ambient Euclidean ball,
and no strictly psh exhaustion is assumed. Each domination is an
inequality of intrinsic original CLD forms. Pulling it to a common
moment refinement preserves positivity; R's estimate already permits
scalar coefficients using additional coordinates. Different choices
of finite charts can give different numerical constants and controls,
but they bound the same intrinsic coefficient measures. The proof
therefore establishes compatibility of the estimate without asserting
a chart-independent numerical norm.

\paragraph{What the estimate does not prove by itself.}
A uniform total-variation bound does not allow a discontinuous
multiplier to pass through weakly convergent coefficient measures.
Consequently \eqref{eq:relative-cln-product}, by itself, proves neither
exchange of two discontinuous slots nor weighted joint decreasing
continuity. A further comparison, energy, or capacity argument is
required. In particular, one cannot replace the fixed continuous
multiplier in R, Lemma 3.1 by an arbitrary bounded psh function merely
by invoking this bound.


% END INCLUDED SOURCE: base_r2/base_r1/sections/relative_cln_agent


% BEGIN INCLUDED SOURCE: base_r2/base_r1/sections/relative_monotone
\section{A common collar and weighted monotone convergence}
\label{sec:monotone}

The results of this section use the internal inputs listed in
Section~\ref{sec:inputs}. In particular, they do not assert an independent
proof of the regularization theorems in \textup{R} and \textup{B}.
All currents and test forms belong to the original CLD calculus.
Positivity means positivity in the cone used in \textup{R}, not strong
positivity. Set $D=d'd''$.

We first record the exact local consequence of the internal approximation
results that is needed below. A smooth psh function in this section has a
positive original $(1,1)$ curvature \emph{form}.

\begin{lemma}[A common domain for smooth majorants]
\label{lem:common-majorants}
Let $x\in U$, and let $u_{i,j}\downarrow u_i$ be finitely many decreasing
families in $\Cb(U)$ with locally bounded limits. After shrinking around
$x$, there is one ordinary neighborhood $W$, independent of $i,j$, with
the following property. Each of the functions $a=u_{i,j}$ and $a=u_i$
has smooth psh majorants $a_\nu$ on $W$ such that
\[
 a_\nu\geq a,\qquad a_\nu(x')\longrightarrow a(x')\quad(x'\in W).
\]
The majorants are bounded uniformly in $\nu$ on each fixed compact subset
after a finite prefix is discarded. If $a\leq M$ on a compact neighborhood
$C\Subset W$, its majorants may be chosen with $a_\nu\leq M+1$ on $C$.
The index discarded may depend on the function $a$.
\end{lemma}

\begin{proof}
Use \textup{B}, Proposition 7.1, to work on one ordinary algebraic
neighborhood, with the original field and calibration. Apply \textup{B},
Lemma D.2, to the entire finite family of sequences, before choosing
approximants. It gives one projective variety, one ample line, one frame
on $W$, and bounded classical metrics whose weights on $W$ are the given
functions plus one fixed constant. By \textup{B}, Theorem 4.4, every such
metric is a decreasing limit of pure Fubini--Study metrics on this same
line. Their weights $p_\nu$ are continuous, satisfy $p_\nu\geq a$, and
decrease pointwise to $a$ after the fixed constant is removed.

Smooth each pure Fubini--Study metric uniformly, using \textup{R},
Lemma 3.2 and Proposition A.1. If the uniform error is at most
$\epsilon_\nu\downarrow0$, add $\epsilon_\nu$ to the smoothed weight.
The resulting smooth psh function $a_\nu$ satisfies
\[
 p_\nu\leq a_\nu\leq p_\nu+2\epsilon_\nu.
\]
This proves pointwise convergence and the majorant property.
On $C$, the increasing open sets $\{p_\nu<M+1/2\}$ cover $C$ whenever
$a\leq M$. Compactness gives a common index on $C$. Discard this prefix
and require $2\epsilon_\nu<1/2$. The same argument with any bound for
$a$ gives local uniform boundedness of the row. No monotonicity of the
\emph{smoothed} row is asserted or needed.
\end{proof}

\begin{lemma}[A smooth psh collar]
\label{lem:smooth-collar}
Under the assumptions of Lemma~\ref{lem:common-majorants}, there are
compact polynomial norm boxes
\[
 x\in\operatorname{Int}C_1,\qquad
 C_1\subset\operatorname{Int}C_{\mathrm{mid}},\qquad
 C_{\mathrm{mid}}\subset\operatorname{Int}C_2\Subset W,
\]
a smooth psh function $b$ on a neighborhood of $C_2$, and a nonnegative
original smooth cutoff $\chi$ compactly supported in
$\Omega=\operatorname{Int}C_2$, with the following properties.
After adding one constant $c$ to all the original functions, put
\[
 v_{i,j}=\max\{u_{i,j}+c,b\},\qquad
 v_i=\max\{u_i+c,b\}.
\]
Then:
\begin{enumerate}
\item $v_{i,j},v_i\in\Cb(\Omega)$ and $v_{i,j}\downarrow v_i$;
\item $v_{i,j}=u_{i,j}+c$ and $v_i=u_i+c$ on a neighborhood of $C_1$;
\item $\chi=1$ on a neighborhood of $C_{\mathrm{mid}}$;
\item all $v_{i,j},v_i$ equal $b$ on one fixed neighborhood of
      $\operatorname{supp}d'\chi\cup\operatorname{supp}d''\chi$;
\item each patched function $v$ has smooth psh majorants $v_\nu\geq v$
      converging pointwise to $v$, uniformly bounded on $\operatorname{supp}\chi$,
      and equal to $b$ on that same collar.
\end{enumerate}
The majorant rows in (5) need not be decreasing.
\end{lemma}

\begin{proof}
Choose the boxes by the construction in \textup{B}, Lemma D.2. In
particular, lower-bound polynomials have no zero near $C_2$. Choose
thresholds strictly between the defining thresholds for $C_1$ and
$C_{\mathrm{mid}}$. The finite maximum
\[
 b_0=\max\{0,\log|P_a^{N_a}/c_a|,
                     \log|d_l/Q_l^{M_l}|\}
\]
is zero on a neighborhood of $C_1$ and at least some $\delta>0$ on a
neighborhood in $C_2$ of $C_2\setminus\operatorname{Int}C_{\mathrm{mid}}$.
All its branches other than the zero branch are logarithms of analytic
units where they are active. A branch with a zero is strictly inactive
near that zero.

Apply the regularized maximum of \textup{R}, Proposition A.1, with
smoothing radius smaller than all the fixed gaps. The result $\widetilde b$
is an original smooth psh function: zero branches at polynomial zeros
are deleted by the exact inactivity property. On a neighborhood of
$C_1$, only the constant zero branch is active, so $\widetilde b=0$.
On the outer collar, $\widetilde b\geq\delta$.

The endpoint bounds give a constant $c$ and a finite $M$ such that
\[
 1\leq u_i+c\leq u_{i,j}+c\leq M
 \quad\hbox{on a neighborhood of }C_2.
\]
Multiply $\widetilde b$ by a positive constant so that the resulting $b$
is greater than $M+3$ on the outer collar. Choose $\chi$, with
$0\leq\chi\leq1$, equal to one near $C_{\mathrm{mid}}$ and compactly
supported in $\Omega$. Smooth cutoff existence and a finite localization
are part of the original CLD test calculus. Its derivative support lies
where the fixed outer-collar inequality holds.

Classical psh functions are stable under maxima and addition of constants.
This proves (1)--(4). To prove (5), apply Lemma~\ref{lem:common-majorants}
to the \emph{unpatched} function $a=u_{i,j}+c$ or $u_i+c$. Discard a
prefix so that its smooth psh majorants $a_\nu$ are at most $M+1$ on
$C_2$. Define
\[
 v_\nu=\operatorname{rmax}_{\epsilon_\nu}(a_\nu,b),
 \qquad \epsilon_\nu\downarrow0,
\]
with all radii less than the fixed collar gap. The regularized maximum
majorizes the ordinary maximum and has error at most $\epsilon_\nu$.
It is exactly $b$ on the collar because the other branch is inactive
there. These facts give all the assertions in (5).
\end{proof}

\begin{remark}
Neither lemma supplies a strictly psh exhaustion of an arbitrary analytic
domain. The collar is constructed on a chosen small norm box. The closed
positive form used to detect a test below is chosen separately and may
be degenerate in directions not occurring in that test.
\end{remark}

\begin{lemma}[Green exchange on a common collar]
\label{lem:collar-green}
Let $Q$ be a symmetric positive $d'$- and $d''$-closed current of
bidegree $(n-1,n-1)$ on $\Omega$. Let $\chi\in A_c^{0,0}(\Omega)$.
Suppose that $a$ is smooth, $g$ is locally bounded Borel, and both equal
one smooth function $b$ on a neighborhood of
$\operatorname{supp}d'\chi\cup\operatorname{supp}d''\chi$.
Whenever $D(gQ)$ has coefficient Radon type, one has
\[
 \langle aD(gQ),\chi\rangle
 =\langle gD(aQ),\chi\rangle.
\]
In this statement positivity of $a$ and $g$ is not required.
\end{lemma}

\begin{proof}
Since $Q$ is closed and $a$ is smooth, $D(aQ)=Da\wedge Q$. The
definition of the distributional derivative gives
\begin{align*}
 \langle aD(gQ)-gD(aQ),\chi\rangle
  &=\langle gQ,D(\chi a)-\chi Da\rangle\\
  &=\langle Q,g\{aD\chi+d'\chi\wedge d''a
                         +d'a\wedge d''\chi\}\rangle.
\end{align*}
Every term in braces is supported in the collar. Substituting $a=g=b$
there makes the last expression
\[
 \langle bQ,D(\chi b)-\chi Db\rangle=0.
\]
The final equality follows by applying the same distributional identity
to the smooth pair $(b,b)$. The substitutions are valid identities of
coefficient measures. No derivative of the Borel function $g$ is used.
\end{proof}

\begin{lemma}[Monotonicity of the ordered trace integral]
\label{lem:trace-monotonicity}
Let $S$ have bidegree $(p,p)$, let $0\leq k\leq n-p$, and let
$\alpha$ be a smooth, positive, $d'$- and $d''$-closed form of bidegree
$(n-p-k,n-p-k)$ such that wedging it with a positive current preserves
positivity. Let $\chi\geq0$ be compactly supported in $\Omega$.

For $0\leq i\leq k$, suppose $f_i,g_i\in\Cb(\Omega)$,
$f_i\leq g_i$, and all these functions equal the same smooth psh
function $b$ near the derivative support of $\chi$. Assume that each
$g_i$ has the smooth majorants with this exact collar supplied by
Lemma~\ref{lem:smooth-collar}. Define, using the ordered products,
\[
 J(f_0;f_1,\ldots,f_k)
   =\langle f_0T_{S,k}(f_1,\ldots,f_k),\chi\alpha\rangle.
\]
Then
\[
 J(f_0;f_1,\ldots,f_k)
 \leq J(g_0;g_1,\ldots,g_k).
\]
Only existence and positivity of the ordered products are used here.
\end{lemma}

\begin{proof}
The case $k=0$ is measure monotonicity. For $k>0$, fix smooth psh
majorants $G_i\geq g_i$ with the common collar. First replace the weight
$f_0$ by $G_0$, using positivity of
$T_{S,k}(f_1,\ldots,f_k)\wedge\alpha$. Apply
Lemma~\ref{lem:collar-green} with
\[
 Q=T_{S,k-1}(f_2,\ldots,f_k)\wedge\alpha
\]
to move $G_0$ into a differentiated position and $f_1$ into the weight.
The resulting current $DG_0\wedge Q$ is positive and closed. Replace
the weight $f_1$ by $G_1$, and repeat. At the end one obtains
\[
 J(f_0;f_1,\ldots,f_k)
 \leq\langle G_k\,DG_0\wedge\cdots\wedge DG_{k-1}\wedge S,
                     \chi\alpha\rangle.
\]
A final application of the same Green identity, together with
commutation of smooth even-degree factors, makes the right-hand side
$J(G_0;G_1,\ldots,G_k)$.

Now let the index of $G_k$ tend to infinity, keeping every other index
fixed. Bounded convergence against the fixed coefficient measures of
$S$ implies $D(G_kS)\to D(g_kS)$. The outer factors and the weight are
smooth and fixed, so this convergence can be tested directly. Next let
the index of $G_{k-1}$ tend to infinity. The background
$D(g_kS)$ is now fixed, so the same argument applies. Continue in the
order $k,k-1,\ldots,1$. Finally, with the ordered product fixed, let the
weight $G_0$ tend to $g_0$ by bounded convergence. The result is exactly
$J(g_0;g_1,\ldots,g_k)$ in its original order. At no step are two
discontinuous limits interchanged.
\end{proof}

\begin{lemma}[An upper bound for decreasing weights]
\label{lem:upper-weight}
Let $P_j,P$ be positive coefficient Radon currents of one bidegree,
with $P_j\to P$ on all original compact smooth tests. Let $a_j\downarrow a$
be locally uniformly bounded functions. Suppose that, near the support
of a fixed positive compact test $\Theta$, every fixed $a_l$ has bounded
continuous majorants converging pointwise to $a_l$.
Then
\[
 \limsup_{j\to\infty}\langle a_jP_j,\Theta\rangle
       \leq\langle aP,\Theta\rangle.
\]
\end{lemma}

\begin{proof}
For $j\geq l$ and a continuous majorant $a_{l,\nu}\geq a_l$,
positivity gives
\[
 \langle a_jP_j,\Theta\rangle
       \leq\langle a_{l,\nu}P_j,\Theta\rangle.
\]
Keep $l,\nu$ fixed and pass to the limit in $j$, using \textup{R},
Lemma 3.1. Bounded convergence against the fixed positive Radon measure
$\mu_{P,\Theta}$ first allows $\nu\to\infty$ and then $l\to\infty$.
\end{proof}

\begin{theorem}[Weighted decreasing continuity for a fixed arbitrary background]
\label{thm:relative-monotone}
Assume the internal inputs \textup{R} and \textup{B}. Let $S$ be any
symmetric positive original CLD current of bidegree $(p,p)$ on $U$,
separately $d'$- and $d''$-closed. Let $0\leq k\leq n-p$.
If $u_{i,j}\downarrow u_i$ in $\Cb(U)$ for $0\leq i\leq k$, with
locally bounded limits, then
\[
 u_{0,j}T_{S,k}(u_{1,j},\ldots,u_{k,j})
 \longrightarrow
 u_0T_{S,k}(u_1,\ldots,u_k)
\]
on every original compact smooth complementary test form. The ordered
unweighted products have the corresponding joint decreasing continuity.
The background $S$ is fixed throughout.
\end{theorem}

\begin{proof}
Write $W_k$ for the weighted assertion with $k$ differentiated slots.
Bounded convergence for the fixed coefficient measures of $S$ proves
$W_0$. Distributional differentiation shows that $W_{k-1}$ implies
joint unweighted convergence in degree $k$. We prove $W_k$ assuming
this unweighted convergence.

The statement is local on compact test supports. Use a finite original
smooth partition to reduce to one test $\eta$ supported in a small
norm box $C_1$ as in Lemma~\ref{lem:smooth-collar}. Include in the chosen
ordinary neighborhood all the unit coordinates needed to present
$\eta$. Let these coordinates be $t_1,\ldots,t_N$, and set
\[
 \beta=\sum_{a=1}^N d't_a\wedge d''t_a,
 \qquad r=n-p-k,\qquad \alpha=\frac{\beta^r}{r!}.
\]
For $r=0$ set $\alpha=1$. The form $\alpha$ is closed and strongly
positive. In the fixed finite coefficient space it is a positive
definite trace form. After replacing $\eta$ by its symmetric part,
there is $C>0$ such that
\[
 \Theta_\pm=C\chi\alpha\pm\eta
\]
are positive original compact tests. This is precisely the matrix
domination used in \textup{R}, Lemma 3.1; additional scalar coordinates
are handled by pullback to a common refinement. No definiteness in
all possible moment directions is required.

Patch all the potentials to $v_{i,j},v_i$ by
Lemma~\ref{lem:smooth-collar}, and put
\[
 P_j=T_{S,k}(v_{1,j},\ldots,v_{k,j}),\qquad
 P=T_{S,k}(v_1,\ldots,v_k).
\]
The induction hypothesis gives $P_j\to P$. Applying
Lemma~\ref{lem:trace-monotonicity} to the ordered tuples
$(v_0,\ldots,v_k)\leq(v_{0,j},\ldots,v_{k,j})$ yields
\[
 \langle v_0P,\chi\alpha\rangle
       \leq\langle v_{0,j}P_j,\chi\alpha\rangle.
\]
Lemma~\ref{lem:upper-weight} gives the reverse limsup inequality.
Consequently
\begin{equation}
 \label{eq:trace-limit}
 \langle v_{0,j}P_j,\chi\alpha\rangle
       \longrightarrow\langle v_0P,\chi\alpha\rangle.
\end{equation}
Apply Lemma~\ref{lem:upper-weight} to $\Theta_+$ and $\Theta_-$.
Together with \eqref{eq:trace-limit}, the two resulting inequalities
give respectively the upper and lower bounds for
$\langle v_{0,j}P_j,\eta\rangle$. Thus this pairing converges to
$\langle v_0P,\eta\rangle$.

On a neighborhood of $\operatorname{supp}\eta$, the patched potentials
are $u_{i,j}+c$ and $u_i+c$. Ordinary locality of the recursion identifies
the differentiated products there with the original products, since
constants in differentiated slots vanish. Remove the constant $c$ from
the weight by the unweighted degree-$k$ convergence already established.
Finite summation over the test partition proves $W_k$.

This induction proves all $W_k$ and all allowed unweighted assertions.
It uses neither capacity nor energy convergence, nor any interchange
of two discontinuous approximation parameters.
\end{proof}

\begin{corollary}[Symmetry and the first-stage calculus]
\label{cor:full-calculus}
Under the hypotheses of Theorem~\ref{thm:relative-monotone}, the currents
$T_{S,k}$ are symmetric in the differentiated arguments,
positive, symmetric as currents, and separately $d'$- and $d''$-closed.
They are additive and positively homogeneous in each potential, commute
with ordinary restriction and gluing, and are independent of the
regularization used to compute them. In particular,
\[
 D\bigl(uD(vS)\bigr)=D\bigl(vD(uS)\bigr)
 \qquad(u,v\in\Cb(U),\ p\leq n-2).
\]
\end{corollary}

\begin{proof}
All assertions except permutation symmetry were proved for the ordered
recursion. Choose, locally on one fixed domain, continuous semipositive
PL approximants decreasing to each potential, as in \textup{B}.
Continuous arbitrary-background products are symmetric, by uniform
smooth approximation and fixed-continuous-multiplier convergence.
Theorem~\ref{thm:relative-monotone}, with a constant weight, passes this
identity to the given potentials. The recursion identifies the result
independently of the approximants and on overlaps.
\end{proof}

\begin{corollary}[Top-degree convergence as measures]
\label{cor:top-measures}
If $p+k=n$, the weighted convergence in
Theorem~\ref{thm:relative-monotone} holds against every compactly
supported continuous scalar function, locally as signed Radon measures.
\end{corollary}

\begin{proof}
Unweighted convergence against one nonnegative smooth cutoff equal to
one near a prescribed compact set bounds the masses of the positive
measures there. The weights have a common actual pointwise bound on a
slightly larger compact neighborhood. Thus the weighted variations are
uniformly bounded. Uniform density of original smooth scalar functions,
with a fixed cutoff, extends convergence to continuous scalar tests.
\end{proof}

\begin{remark}[What independence means here]
The construction is intrinsic to the \emph{given potentials} and to $S$.
Independence of regularization does not assert that $D(uS)$ depends only
on the ambient current $D[u]$ for every possible pair of potentials.
That stronger quotient-by-potentials statement is distinct from the
claims proved here and from the special Bedford--Taylor-current
condition in CLD 2025, Definition 8.6.2. Equality on all original tests
also does not assert equality on all $\delta$-tests.
\end{remark}

% END INCLUDED SOURCE: base_r2/base_r1/sections/relative_monotone


% BEGIN INCLUDED SOURCE: base_r2/base_r1/sections/functional_weights
\section{Continuous functions of several monotone weights}
\label{sec:functional-weights}

This section records a \texttt{DIRECT\_COROLLARY} of
Theorem~\ref{thm:relative-monotone}. It is useful when a localization
factor is a continuous function of several psh potentials, although that
factor need not itself be psh.

\begin{lemma}[Increasing convex compositions]
\label{lem:convex-compositions}
Let $h\in\Cb(V)$, and suppose that its values lie in a compact interval
$I$. If $\varphi:I\to\mathbf R$ is continuous, convex, and
nondecreasing, then $\varphi\circ h\in\Cb(V)$. In particular, if
$h\geq0$, every positive integral power of $h$ is classical psh.
If $h_j\downarrow h$ have values in $I$, then
$\varphi(h_j)\downarrow\varphi(h)$.
\end{lemma}

\begin{proof}
Interpolate $\varphi$ linearly on successively refined dyadic partitions
of $I$. The resulting functions $\varphi_m$ are convex and
nondecreasing, are finite maxima of affine functions with nonnegative
slopes, and decrease uniformly to $\varphi$. Each
$\varphi_m(h)$ is therefore a finite maximum of nonnegative multiples
of $h$ plus constants. The defining classical class is stable under
these operations. Its decreasing-limit property, with the local
boundedness of the limit, gives $\varphi(h)\in\Cb(V)$.
The last assertion follows from continuity and monotonicity of
$\varphi$.
\end{proof}

\begin{proposition}[Functional weighted convergence]
\label{prop:functional-weights}
Let $S$ be a fixed background as in
Theorem~\ref{thm:relative-monotone}, and let $0\leq k\leq n-p$.
Suppose that
\[
 u_{i,j}\downarrow u_i\quad(1\leq i\leq k),
 \qquad h_{a,j}\downarrow h_a\quad(1\leq a\leq s)
\]
are finite families in $\Cb(U)$ with locally bounded limits. Set
\[
 P_j=T_{S,k}(u_{1,j},\ldots,u_{k,j}),\qquad
 P=T_{S,k}(u_1,\ldots,u_k).
\]
For every continuous function $F:\mathbf R^s\to\mathbf R$,
\[
 F(h_{1,j},\ldots,h_{s,j})P_j
 \longrightarrow F(h_1,\ldots,h_s)P
\]
on every original compact smooth complementary test. The factors
$F(h_{1,j},\ldots,h_{s,j})$ are the actual Borel functions, and no
psh assertion about those factors is required.
\end{proposition}

\begin{proof}
Fix a compact test $\eta$ and a relatively compact neighborhood $V$ of
its support. The endpoint bounds provide one compact interval for all
values of each $h_{a,j}$ and $h_a$ on $V$. Add fixed constants so that
the shifted functions $x_{a,j}=h_{a,j}+c_a$ and $x_a=h_a+c_a$ are
nonnegative there. Absorb these shifts into $F$.

We first prove convergence for a polynomial in the $x_a$. Every
monomial of total degree $d\geq1$ can be regarded as a product
$z_1\cdots z_d$ in which variables may repeat. The polarization
identity
\[
 z_1\cdots z_d
 =\frac1{d!}\sum_{A\subset\{1,\ldots,d\}}
       (-1)^{d-|A|}\left(\sum_{a\in A}z_a\right)^d
\]
expresses it as a linear combination of powers of sums with
nonnegative coefficients. Each nonzero sum of shifted psh functions
is nonnegative and classical psh. By
Lemma~\ref{lem:convex-compositions}, its $d$th power is also classical
psh and decreases along the given sequences. The weighted theorem
therefore gives convergence for every term in the displayed identity.
Constants are covered by its constant-weight case. Finite linearity
proves the assertion for every polynomial.

The unweighted convergence $P_j\to P$, together with \textup{R},
Lemma 3.1, supplies a finite bound
\[
 \sup_j|\mu_{P_j,\eta}|(\operatorname{supp}\eta)
       +|\mu_{P,\eta}|(\operatorname{supp}\eta)<\infty.
\]
Approximate $F$ uniformly by polynomials on the common compact
rectangle of shifted values. An approximation error of at most
$\epsilon$ changes each pairing by at most $\epsilon$ times this
fixed variation bound. First pass to the limit in $j$ for the
polynomial, and then let $\epsilon\downarrow0$. This proves the
assertion on $\eta$.
\end{proof}

\begin{remark}
The proposition permits, for example, bounded continuous cutoffs of
$u_j-v_j$ and products of finitely many bounded psh weights.
It does not permit replacing $F$ by a discontinuous function without
an additional argument. In applications to strict maximum sets, one
first uses continuous scalar cutoffs and only afterwards takes a
bounded pointwise limit against the already fixed coefficient measures.
\end{remark}


% END INCLUDED SOURCE: base_r2/base_r1/sections/functional_weights


% BEGIN INCLUDED SOURCE: base_r2/base_r1/sections/scope_stability
\section{A sufficient condition for varying backgrounds}
\label{sec:dominated-backgrounds}

The main theorem fixes the background. The following separate result
gives a sufficient additional hypothesis under which a background may
vary. It is an elementary \texttt{DIRECT\_COROLLARY} of the
coefficient-Radon domination argument and the ordered construction;
it does not assert unrestricted varying-background stability.

\begin{lemma}[Order domination and Borel multipliers]
\label{lem:order-dominated-borel}
Let $P_j,P,R$ be original positive currents of one bidegree on $U$.
Suppose that $P_j\to P$ on all original compact smooth tests and that
\[
 0\leq P_j\leq C R
\]
for one constant $C<\infty$. Let $a_j$ be Borel functions converging
pointwise to $a$, with one common bound on each compact subset.
Then
\[
 a_jP_j\longrightarrow aP
\]
on every original compact smooth test. Closedness of these currents
is not required in this lemma.
\end{lemma}

\begin{proof}
Weak closure of the positive cone gives $0\leq P\leq CR$.
Fix a compact test $\eta$. The finite positive control construction
in R, Lemma 3.1 supplies a finite positive Radon measure $\nu_\eta$,
depending only on $R$, $C$ and the test, such that
\begin{equation}\label{eq:fixed-measure-domination}
 |\mu_{P_j,\eta}|\leq\nu_\eta,
 \qquad |\mu_{P,\eta}|\leq\nu_\eta.
\end{equation}
Here these are inequalities of measures, not only bounds on their
total masses. Indeed, on each moment piece choose a positive compact
control $\Theta$ with $A\Theta\pm\eta\geq0$. For every nonnegative
smooth scalar $h$,
\[
 |\langle P_j,h\eta\rangle|
 \leq A\langle P_j,h\Theta\rangle
 \leq AC\langle R,h\Theta\rangle.
\]
Radon uniqueness and scalar density give domination by
$AC\mu_{R,\Theta}$. Sum over the finite pieces to obtain
\eqref{eq:fixed-measure-domination}.

For a fixed bounded Borel function $b$, approximate $b$ in
$L^1(\nu_\eta)$ by continuous functions on a compact neighborhood
of the test support. Such approximation is available for a finite
Radon measure; cutoffs keep the functions on that neighborhood.
R, Lemma 3.1 gives convergence with each fixed continuous multiplier,
and \eqref{eq:fixed-measure-domination} bounds the approximation error
uniformly in $j$. It follows that
\[
 \int b\,d\mu_{P_j,\eta}\longrightarrow
 \int b\,d\mu_{P,\eta}.
\]
Finally,
\[
 \left|\int(a_j-a)\,d\mu_{P_j,\eta}\right|
 \leq\int|a_j-a|\,d\nu_\eta\longrightarrow0
\]
by bounded convergence for the fixed measure $\nu_\eta$. Apply the
fixed-multiplier assertion to $b=a$ and combine the two limits.
\end{proof}

\begin{proposition}[Order-dominated varying backgrounds]
\label{prop:dominated-background-stability}
Let $R_0,S_j,S$ be symmetric positive separately closed original
$(p,p)$-currents on $U$. Assume
\[
 S_j\longrightarrow S,
 \qquad 0\leq S_j\leq C R_0
\]
on $U$, with one finite constant $C$. Fix
$u_1,\ldots,u_k\in\Cb(U)$, where $0\leq k\leq n-p$.
Then
\[
 T_{S_j,k}(u_1,\ldots,u_k)
 \longrightarrow T_{S,k}(u_1,\ldots,u_k).
\]
Moreover, for every locally uniformly bounded sequence of Borel
functions $a_j\to a$ pointwise,
\[
 a_jT_{S_j,k}(u_1,\ldots,u_k)
 \longrightarrow aT_{S,k}(u_1,\ldots,u_k).
\]
All convergence is on the full original compact smooth test sheaf.
The differentiated potential tuple is fixed.
\end{proposition}

\begin{proof}
For a fixed potential tuple the ordered construction is linear in
the background. Positivity therefore gives
\begin{equation}\label{eq:product-order-domination}
 0\leq T_{S_j,r}(u_{k-r+1},\ldots,u_k)
 \leq C T_{R_0,r}(u_{k-r+1},\ldots,u_k)
\end{equation}
for every $0\leq r\leq k$: the difference is the ordered product
with the positive closed background $CR_0-S_j$.

Induct on $r$. The case $r=0$ is the assumed weak convergence.
Assume convergence at degree $r-1$. Its currents are dominated by
the one fixed positive current on the right of
\eqref{eq:product-order-domination}. Apply
Lemma~\ref{lem:order-dominated-borel} with the fixed multiplier
$u_{k-r+1}$, and then apply $D$. This gives convergence at degree
$r$ by the recursion. At $r=k$, apply the same lemma once more,
now with the varying Borel weights $a_j$, to obtain the final
assertion.
\end{proof}

\begin{remark}
The hypothesis $S_j\leq CR_0$ supplies domination by fixed coefficient
measures. It is stronger than a uniform bound on local masses.
The conclusion permits discontinuous Borel weights precisely because
this additional domination is available. No varying differentiated
tuple, and no assertion for arbitrary weakly convergent backgrounds,
is included.
\end{remark}

\subsection{Continuous changes of potential}

\begin{lemma}[Continuous harmonic multipliers]
\label{lem:continuous-harmonic-multipliers}
Assume the local regularization theorem of \textup{R}. Let $S$ be an
allowed positive separately closed original current of bidegree
$(p,p)$, with $p\leq n-1$. If $h$ is finite continuous and
$D[h]=0$ as an ambient original current, then
\[
 D(hS)=0.
\]
Consequently, if $u,v$ are finite continuous and $D[u]=D[v]$, then
$D(uS)=D(vS)$. This conclusion concerns continuous potentials.
\end{lemma}

\begin{proof}
Work on the ordinary local algebraic neighborhoods supplied by
\textup{B}, Proposition 7.1. Both $h$ and $-h$ have positive ambient
curvature currents, since those currents vanish. By \textup{R},
Theorem 2.9, each is locally uniformly approximable by smooth
functions with positive curvature forms. For a smooth approximant
$h_m$, the current
\[
 D(h_mS)=Dh_m\wedge S
\]
is positive. Locally uniform convergence passes through multiplication
by the fixed coefficient-Radon current $S$ and then through $D$.
Hence $D(hS)$ is positive. Apply the same argument to $-h$: the
current $-D(hS)$ is positive as well. The original positive cone of
currents is pointed, by the finite-test domination in R, Lemma 3.1,
so $D(hS)=0$. Take $h=u-v$ for the final assertion and localize.
No classical-psh assertion about $h$ is needed.
\end{proof}

This is a \texttt{DIRECT\_COROLLARY} of \textup{R} and the fixed-
background one-potential construction. Applying it to the positive
closed current formed by the remaining slots also proves potential
independence for a continuous slot in a mixed product. It does not
establish the same statement for an arbitrary discontinuous
difference $u-v$ with $D[u-v]=0$.


% END INCLUDED SOURCE: base_r2/base_r1/sections/scope_stability


% BEGIN INCLUDED SOURCE: base_r2/base_r1/sections/potential_independence_scope
\section{A limited quotient-by-potentials statement}
\label{sec:potential-independence-scope}

The main theorem constructs products of specified, actual pointwise
potentials. The assertion that $D(uS)$ depends only on the ambient
current $D[u]$ is stronger. The following cases do follow from the
available inputs. They are \texttt{DIRECT\_COROLLARY} statements, not
an assertion of this stronger property for every background.

\begin{lemma}[Ambient mixed backgrounds]
\label{lem:potential-independence-ambient-mixed}
Let $u,v\in\Cb(U)$ satisfy $D[u]=D[v]$ on all original compact smooth
tests. Let $0\leq p\leq n-1$, and suppose locally that
\[
 S=\sum_{\ell=1}^{N}c_\ell
      T_{[1],p}(a_{\ell,1},\ldots,a_{\ell,p}),
 \qquad a_{\ell,i}\in\Cb(U),\quad c_\ell\in\mathbf R.
\]
Then
\[
 D(uS)=D(vS).
\]
The conclusion also holds for a locally finite sum of this kind.
In particular, it applies when such a current is an allowed positive
closed background.
\end{lemma}

\begin{proof}
It suffices to consider one summand. By the symmetry already present
in \textup{B}, Theorem 2.6,
\begin{align*}
 D\bigl(uT_{[1],p}(a_1,\ldots,a_p)\bigr)
  &=T_{[1],p+1}(a_1,\ldots,a_p,u)\\
  &=D\bigl(a_1D(a_2\cdots D(a_pD[u])\cdots)\bigr).
\end{align*}
Replace the innermost current $D[u]$ by $D[v]$. At every subsequent
step, equality of coefficient-Radon currents implies equality after
multiplication by the same actual Borel function, by \textup{B},
Lemma 7.2. Distributional differentiation preserves that equality.
Reversing the preceding symmetry gives $D(vS)$. Finite linearity and
localization prove the stated extensions. For $p=0$, the assertion is
the original equality $D[u]=D[v]$.
\end{proof}

\begin{lemma}[Smooth closed backgrounds]
\label{lem:potential-independence-smooth}
Under the same hypothesis on $u,v$, let $S=[\alpha]$, where $\alpha$
is an original smooth $(p,p)$-form satisfying
$d'\alpha=d''\alpha=0$. Then $D(uS)=D(vS)$.
\end{lemma}

\begin{proof}
Set $h=u-v$. The given local boundedness makes all products below
well defined. Borel multiplication of the integration current agrees
with smooth-form multiplication:
\[
 h[\alpha]=\alpha\wedge[h].
\]
The distributional Leibniz rule and the two closedness assumptions give
\[
 D(hS)=D(\alpha\wedge[h])=\alpha\wedge D[h]=0.
\]
This is an equality on the entire original compact smooth test sheaf.
\end{proof}

\begin{proposition}[Passage through a dominated weak closure]
\label{prop:potential-independence-dominated-closure}
Fix $u,v\in\Cb(U)$ with $D[u]=D[v]$. Let $S_j,S,R_0$ be allowed
positive separately closed backgrounds of degree $(p,p)$, with
\[
 S_j\longrightarrow S,\qquad 0\leq S_j\leq C R_0
\]
for one finite constant $C$. Suppose that $D(uS_j)=D(vS_j)$ for every
$j$, for example by either preceding lemma. Then
\[
 D(uS)=D(vS).
\]
\end{proposition}

\begin{proof}
The function $h=u-v$ is an actual locally bounded Borel function.
Lemma~\ref{lem:order-dominated-borel}, with the fixed multiplier $h$,
gives $hS_j\to hS$ on all original compact smooth tests.
Apply $D$ to this convergence. Every $D(hS_j)$ is zero by hypothesis,
so $D(hS)=0$.
\end{proof}

\begin{remark}[The remaining scope question]
These arguments give neither a representation of an arbitrary
original-positive closed current by ambient bounded-potential products,
nor a dominated approximation of every such current by the backgrounds
above. They also do not show that a locally bounded difference
$h=u-v$ satisfying $D[h]=0$ is a continuous function at its actual
point values. This original-current kernel regularity question remains
open in this investigation. The previously proved continuous-harmonic
multiplier lemma applies if that additional continuity is known.
\end{remark}


% END INCLUDED SOURCE: base_r2/base_r1/sections/potential_independence_scope


% BEGIN INCLUDED SOURCE: base_r2/base_r1/sections/relative_energy_agent
\section{Relative energy and convergence in capacity}
\label{sec:relative-energy-capacity}

Throughout this section, the fixed-background calculus of
Theorem~\ref{thm:relative-monotone} is available. Its proof uses no energy
or capacity assertion from this section. All measures below are measures
on the actual Berkovich space, and all multipliers use their given point
values. We continue to use the original positive cone of \textup{R}.

\subsection{A relative energy form}

\begin{lemma}[Regularization on a prescribed smaller norm box]
\label{lem:compact-box-regularization}
Work in one ordinary algebraic neighborhood supplied by \textup{B}.
Let $C_{\mathrm{in}}\Subset\operatorname{Int}C_{\mathrm{out}}$ be
compact polynomial norm boxes, and let finitely many bounded classical
psh functions be defined on a neighborhood of $C_{\mathrm{out}}$.
They admit smooth decreasing psh approximants on one neighborhood of
the entire $C_{\mathrm{in}}$. Thus the smaller domain need not be
chosen separately at each point of a prescribed compact test support.
\end{lemma}

\begin{proof}
This is the compact-box form of the construction in \textup{B},
Lemma D.2, followed by its Theorem 4.4. We spell out the distinction
from applying local regularization separately at every point.
Insert a middle norm box and construct the finite polynomial maximum
$b$ used in that lemma: it is zero on $C_{\mathrm{in}}$ and is at
least a fixed $\delta>0$ on the outer collar. Homogenizing its branches
and adding inactive global generating sections realizes $b$ as the
weight of one pure-Fubini--Study metric $q$, in a frame nonzero near
$C_{\mathrm{out}}$. After adding one common constant, all given
functions lie in $[1,M]$ there. Choose $m\delta>M+1$ and give each
extended metric the local weight $\max\{a+c,mb\}$, setting it equal
to $mq$ outside the middle box. These prescriptions agree on their
overlap. They define bounded classical metrics on one projective
variety and one ample line, with weights exactly $a+c$ on a
neighborhood of $C_{\mathrm{in}}$.

Theorem 4.4 of \textup{B} gives decreasing global pure-Fubini--Study
approximants for each metric. Their smooth global approximations from
\textup{R}, Lemma 3.2, restrict to the same frame near the whole
$C_{\mathrm{in}}$. If the uniform smoothing error is at most
$\epsilon_m=4^{-m}$, adding $2\epsilon_m$ gives a decreasing smooth
row: its values lie between $p_m+\epsilon_m$ and
$p_m+3\epsilon_m$, while $p_{m+1}\leq p_m$.
These particular smooth approximants are classical psh. Indeed, after
any complete field extension and pullback to an analytic curve, their
active branches are logarithms of norms of analytic units, hence
harmonic functions. The regularized maximum is convex and
nondecreasing in each variable, so its composition with these branches
is subharmonic, by the curve chain rule in \textup{B}, Lemma B.1.
A section vanishing identically or at a point gives
an inactive branch nearby, which is deleted. Positive scaling and addition of
constants preserve this property. Thus the rows used below belong to
$\Cb$; no general identification of original-form positivity with
classical plurisubharmonicity is being assumed. The same convex,
coordinatewise nondecreasing composition argument applies to the
additional regularized maxima with the smooth collar below.
Remove the common constant $c$ to finish. No extension of a background
current is used in this argument.

\end{proof}

Let $Q$ be a symmetric positive, separately $d'$- and $d''$-closed
current of bidegree $(n-1,n-1)$ on an ordinary open $\Omega$. Let
$\mathcal D_b(\Omega)$ be the real vector space generated locally by
$\Cb(\Omega)$. If $f,g\in\mathcal D_b(\Omega)$, define
\begin{equation}
\label{eq:relative-gamma}
 \Gamma_Q(f,g)
   =\frac12\bigl\{D(fgQ)-fD(gQ)-gD(fQ)\bigr\},
 \qquad \Gamma_Q(f)=\Gamma_Q(f,f).
\end{equation}
Here $D(fQ)$ is a signed coefficient Radon current, by linearity of
the ordered one-potential construction. The first term in
\eqref{eq:relative-gamma} is initially a distribution. The following
lemma proves that the full expression is a measure.

\begin{lemma}[Relative energy measures]
\label{lem:relative-energy-measures}
For $f,g\in\mathcal D_b(\Omega)$, $\Gamma_Q(f,g)$ is a signed Radon
measure and $\Gamma_Q(f)$ is nonnegative. For every nonnegative compact
smooth $\chi$, one has
\begin{equation}
\label{eq:relative-energy-CS}
 \left|\int\chi\,d\Gamma_Q(f,g)\right|^2
 \leq\left(\int\chi\,d\Gamma_Q(f)\right)
       \left(\int\chi\,d\Gamma_Q(g)\right).
\end{equation}
Suppose, for the following integration formula, that $f=a-b$, where
$a,b\in\Cb$ are functions on one common norm-box domain and agree
outside a compact $L$ contained in a strictly smaller norm box.
Suppose also that the psh summands of $g$ are defined on that larger
domain. This includes the common-collar differences used below.
Then
\begin{equation}
\label{eq:relative-green-energy}
 \int f\,D(gQ)=\int g\,D(fQ)
       =-\int\Gamma_Q(f,g),
 \qquad
 \int\Gamma_Q(f)=-\int f\,D(fQ).
\end{equation}
The integrals in \eqref{eq:relative-green-energy} are supported on a
fixed compact neighborhood of $\operatorname{supp}f$.
\end{lemma}

\begin{proof}
First suppose that all functions are smooth. Then
\[
 \Gamma_Q(f,g)
  =\frac12(d'f\wedge d''g+d'g\wedge d''f)\wedge Q.
\]
Symmetry of $Q$ makes the pairing symmetric. Since
$d'f\wedge d''f$ is a strongly positive $(1,1)$-form, the measure
$\Gamma_Q(f)$ is positive. Positivity of
$\Gamma_Q(f+t g)$ for every real $t$ proves
\eqref{eq:relative-energy-CS} by the discriminant inequality.
Stokes proves \eqref{eq:relative-green-energy}.

We explain carefully the approximation used to pass these identities
to bounded potentials. On a common smaller domain, let $p_m\downarrow a$
be the continuous pure-Fubini--Study weights supplied by \textup{B} for
a given $a\in\Cb$. Smooth $p_m$ uniformly with error at most
$\epsilon_m=4^{-m}$, and call the result $s_m$. Then
\[
 a_m=s_m+2\epsilon_m,
 \qquad
 p_m+\epsilon_m\leq a_m\leq p_m+3\epsilon_m,
 \qquad a_{m+1}<a_m.
\]
Thus $a_m$ are smooth psh functions decreasing pointwise to $a$.
The same construction is made for finitely many psh summands.

If an exact common collar is required, apply a regularized maximum
with the fixed smooth collar function from
Lemma~\ref{lem:smooth-collar}. Its smoothing radii may decrease to zero.
For completeness, the regularized maximum is nondecreasing in the
smoothing radius: for fixed input, its expression as the integral of
a maximum of affine functions of the radius is convex, and symmetry
of the mollifier makes it even. Hence it is nondecreasing for
nonnegative radius. The patched smooth rows remain decreasing and
equal the prescribed collar function exactly.

Write $f=\sum c_i a_i$ and $g=\sum d_l b_l$, with classical psh
$a_i,b_l$, and use such simultaneous smooth rows. Products of bounded
functions converge pointwise and are uniformly bounded on each compact
set. Therefore $D(f_mg_mQ)\to D(fgQ)$ by bounded convergence against
the fixed coefficient measures of $Q$. Each term in $f_mD(g_mQ)$
is a constant multiple of $a_{i,m}D(b_{l,m}Q)$, so it converges by
Theorem~\ref{thm:relative-monotone} with one differentiated slot.
The same holds with $f$ and $g$ exchanged. Consequently
\[
 \Gamma_Q(f_m,g_m)\longrightarrow\Gamma_Q(f,g)
\]
as distributions. A positive distribution in top degree is a positive
Radon measure. Apply this to $f+t g$, and polarize, to obtain the
measure assertion and \eqref{eq:relative-energy-CS}.

It remains to justify compact-support integration by parts without
assuming differentiability of $f$. Write $f=a-b$ as in the statement.
Choose several nested norm boxes between $L$ and the boundary of the
given domain. First apply Lemma~\ref{lem:compact-box-regularization}
with a smaller box still containing all the boxes to be used below.
This gives smooth decreasing rows for $a,b$ and all psh summands of
$g$ on a neighborhood of that entire box. This step is essential:
applying a pointwise regularization theorem separately would not
produce a domain containing the full support of $f$.

Inside this common regularization domain, choose a second smooth psh
collar which is zero near $L$ and dominates the bounded shifted values
of $a$ and $b$ on an outer collar. Patch their smooth rows by a
regularized maximum with this collar, with radii decreasing to zero.
The limiting patched difference is still exactly $f$: on $L$ the
collar is inactive, and outside $L$ the two limiting functions agree.
The smooth differences $f_m$ have support in one fixed compact set
because their two rows have the same exact outer collar. The smooth
rows of the psh summands of $g$ are already defined near this compact
set. No approximation or global extension of the fixed $Q$ is involved.

For these smooth rows, insert one compact smooth cutoff equal to one
on the common support of $f_m$. Stokes gives the two Green equalities.
Every term is a finite sum of weighted one-potential products, so
Theorem~\ref{thm:relative-monotone} passes it to the limit. The
energy terms converge by the already proved distributional formula
and positivity, with the same cutoff. All resulting multipliers are the
Borel multipliers in \eqref{eq:relative-gamma}; no exceptional-set
representative is chosen. This local integration formula is the only
compact-support version used below.
\end{proof}

\begin{remark}[A noncircular estimate before joint continuity]
\label{rem:energy-before-continuity}
There is a useful estimate that precedes the preceding lemma. Suppose
$a_j\downarrow a$ are \emph{smooth} psh functions, uniformly bounded
near a compact set, and all equal one smooth function outside a fixed
compact $L\Subset\Omega$. Choose $\chi\geq0$, compactly supported,
with $\chi=1$ near $L$, and put
\[
 E_\chi(a)=\int\chi\,d'a\wedge d''a\wedge Q.
\]
For $j<l$, $g=a_j-a_l\geq0$ has compact support in $L$, and smooth
integration by parts gives
\begin{equation}
\label{eq:smooth-energy-cauchy}
 E_\chi(a_l)-E_\chi(a_j)
   =E_\chi(g)+2\int g\,D(a_jQ)\geq E_\chi(g).
\end{equation}
The identity
\[
 E_\chi(a_j)=\frac12\langle Q,a_j^2D\chi\rangle
                  -\langle D(a_jQ),\chi a_j\rangle
\]
and the relative CLN estimate bound these energies uniformly. They
therefore increase to a finite limit, and the gradients are Cauchy
in the local Hilbert seminorm induced by $Q$. In fact the smooth
energy measures are Cauchy in total variation:
\[
 \|\chi\Gamma_Q(a_j)-\chi\Gamma_Q(a_l)\|_{\mathrm{TV}}
 \leq E_\chi(a_j-a_l)^{1/2}
       \bigl(E_\chi(a_j)^{1/2}+E_\chi(a_l)^{1/2}\bigr).
\]
To see this, polarize the difference and use the measure-valued
Cauchy--Schwarz inequality. That inequality follows by writing the
positive $2\times2$ matrix of energy measures relative to the sum of
its two diagonal measures, applying the pointwise discriminant
inequality to its densities, and then ordinary Cauchy--Schwarz.
This statement uses only the one-potential construction and relative
CLN. We do not use
it to prove Theorem~\ref{thm:relative-monotone}; that theorem has the
independent trace-squeeze proof.
\end{remark}

\subsection{A localized relative capacity}

Fix a positive closed background $S$ of bidegree $(p,p)$, and put
$q=n-p$. Work on a relatively compact ordinary moment domain $\Omega$
on which a smooth closed strongly positive $(1,1)$-form $\beta$ is
defined. A typical choice is the trace form in the proof of
Theorem~\ref{thm:relative-monotone}. No strict positivity of $\beta$
in all analytic directions is assumed. For $1\leq k\leq q$ and
a Borel subset $E\subset\Omega$, define
\begin{equation}
\label{eq:relative-capacity}
 \operatorname{Cap}_{S,\beta,k}(E;\Omega)
  =\sup_{\substack{v\in\Cb(\Omega)\\0\leq v\leq1}}
    \int_ET_{S,k}(v,\ldots,v)\wedge\beta^{q-k}.
\end{equation}
For $k=q$, $\beta$ drops out, and we write
$\operatorname{Cap}_S(E;\Omega)$. In particular, $S=[1]$ gives the
usual candidate non-Archimedean Bedford--Taylor capacity in top degree.
We make no nondegeneracy assertion for an arbitrary $S$ or $\Omega$.

\begin{proposition}[Elementary properties and mixed domination]
\label{prop:capacity-elementary}
The capacity in \eqref{eq:relative-capacity} is monotone and countably
subadditive on Borel sets, finite on compact subsets of $\Omega$, and
continuous from below on increasing Borel sets. Restriction to a
smaller domain can only increase it. If
$a_1,\ldots,a_k\in\Cb(\Omega)$ satisfy $0\leq a_i\leq M_i$, then
\begin{equation}
\label{eq:mixed-capacity-domination}
 \int_ET_{S,k}(a_1,\ldots,a_k)\wedge\beta^{q-k}
 \leq\frac{k^k}{k!}\left(\prod_{i=1}^kM_i\right)
       \operatorname{Cap}_{S,\beta,k}(E;\Omega).
\end{equation}
If some $M_i=0$, the left side vanishes. Constants may first be
subtracted from the $a_i$, so bounds on their oscillations suffice.
\end{proposition}

\begin{proof}
The first assertions follow from the same properties of positive
measures and the supremum in \eqref{eq:relative-capacity}; for an
increasing sequence, the two suprema over the sequence index and the
admissible functions commute. Compact finiteness is the uniform
relative CLN bound. Domain monotonicity follows by restricting every
admissible function on the larger domain.

If all $M_i>0$, set $v_i=a_i/M_i$ and
$v=k^{-1}\sum_i v_i$. Then $0\leq v\leq1$. Symmetry, multilinearity,
and positivity show that the expansion of
$T_{S,k}(v,\ldots,v)$ contains
$k!k^{-k}T_{S,k}(v_1,\ldots,v_k)$, with all other terms
positive. This gives \eqref{eq:mixed-capacity-domination}.
\end{proof}

\begin{theorem}[A quantitative relative capacity estimate]
\label{thm:relative-capacity-estimate}
For this theorem, take $\Omega$ to be the ordinary interior of one
polynomial norm box in an algebraic neighborhood, and assume that
the compact set $L$ below is contained in a strictly smaller nested
norm box. Let $u_j,u\in\Cb(\Omega)$ satisfy $u_j\downarrow u$.
Suppose that there are such a compact $L\Subset\Omega$ and one smooth psh function $h$
such that $u_j=u=h$ on $\Omega\setminus L$, and that the functions
have a common bound on a neighborhood of $L$. Set $g_j=u_j-u\geq0$.
For $1\leq k\leq q$, there is a finite constant $C$, independent of
$j$, $v$, and $\delta>0$, such that
\begin{equation}
\label{eq:capacity-rate}
 \operatorname{Cap}_{S,\beta,k}(\{g_j>\delta\};\Omega)
 \leq\frac{C}{\delta}
   \left(\int g_j\,
      T_{S,k}(u,\ldots,u)\wedge\beta^{q-k}\right)^{2^{-k}}.
\end{equation}
The constant depends on the compact collar, $\beta$, a bound for $u$
on a slightly larger compact set, and finitely many positive cutoff
evaluations of $S$. Consequently $u_j\to u$ in this relative capacity.
\end{theorem}

\begin{proof}
Choose a nonnegative compact smooth $\chi$ equal to one near $L$.
Fix an admissible $v$, and abbreviate products using the symmetric
calculus already proved. For $0\leq r\leq k$, set
\[
 I_r=\int g_j\,(Dv)^r(Du)^{k-r}\wedge S\wedge\beta^{q-k}.
\]
These numbers are finite and nonnegative. For $r\geq1$, put
\[
 Q_r=(Dv)^{r-1}(Du)^{k-r}\wedge S\wedge\beta^{q-k}.
\]
This is an arbitrary allowed positive closed current of bidegree
$(n-1,n-1)$; the energy lemma therefore applies to it. Since $g_j$
has compact support in $L$, Green's identity and Cauchy--Schwarz give
\begin{align}
 I_r
  &= -\int\Gamma_{Q_r}(g_j,v),\notag\\
 I_r^2
  &\leq\left(\int\Gamma_{Q_r}(g_j)\right)
          \left(\int\chi\,\Gamma_{Q_r}(v)\right). 
 \label{eq:capacity-energy-step}
\end{align}
The first factor satisfies the one-sided estimate
\begin{align}
 0\leq\int\Gamma_{Q_r}(g_j)
  &= -\int g_jD(g_jQ_r)\notag\\
  &=\int g_jD(uQ_r)-\int g_jD(u_jQ_r)
    \leq I_{r-1}.
 \label{eq:capacity-one-sided}
\end{align}
The final inequality uses $g_j\geq0$ and positivity of
$D(u_jQ_r)$. Thus there is no need to estimate a signed curvature
by exchanging two limits.

For the second factor, the energy identity gives
\[
 \int\chi\,\Gamma_{Q_r}(v)
   =\frac12\langle Q_r,v^2D\chi\rangle
                         -\langle D(vQ_r),\chi v\rangle.
\]
It is nonnegative. Since $0\leq v\leq1$, it is bounded above by
the absolute coefficient estimate for the first term. Apply the
relative CLN bound to $Q_r$, whose differentiated arguments consist
of $r-1$ copies of $v$ and $k-r$ copies of the fixed bounded $u$.
The smooth $\beta$ factors are fixed. Hence
\[
 0\leq\int\chi\,\Gamma_{Q_r}(v)\leq A_r,
\]
where $A_r$ is finite and independent of $v$ and $j$. Combining this
with \eqref{eq:capacity-energy-step}--\eqref{eq:capacity-one-sided}
and iterating yields
\[
 I_k\leq
 A_k^{1/2}A_{k-1}^{1/4}\cdots A_1^{1/2^k}I_0^{1/2^k}.
\]
On $\{g_j>\delta\}$, the inequality
$1\leq g_j/\delta$ gives \eqref{eq:capacity-rate} after taking the
supremum over $v$.

Finally, $I_0$ is the integral of $g_j\downarrow0$ against the
\emph{fixed} locally finite positive measure
$T_{S,k}(u,\ldots,u)\wedge\beta^{q-k}$. The functions
$g_j$ have a common bound and fixed compact support. Bounded
convergence gives $I_0\to0$. This last step does not use a claim about
varying backgrounds.
\end{proof}

\begin{corollary}[Local convergence in capacity and quasicontinuity]
\label{cor:relative-quasicontinuity}
On a sufficiently small ordinary norm-box neighborhood, let
$V_0\Subset\Omega$ be a smaller core and let $u_j\downarrow u$ in
$\Cb$ on a neighborhood of $\overline\Omega$. Then, for every
$\delta>0$,
\[
 \operatorname{Cap}_{S,\beta,k}
   (V_0\cap\{u_j-u>\delta\};\Omega)\longrightarrow0.
\]
Every $u\in\Cb$ is locally quasicontinuous for this capacity: for
every $\epsilon>0$, there is an ordinary open $G\subset\Omega$
with $\operatorname{Cap}_{S,\beta,k}(G;\Omega)<\epsilon$ such that
$u|_{V_0\setminus G}$ is continuous in the subspace topology.
For $k=q$, this statement uses the intrinsic relative capacity
$\operatorname{Cap}_S$ and needs no trace form.
\end{corollary}

\begin{proof}
Patch the entire decreasing family by the common collar of
Lemma~\ref{lem:smooth-collar}. The patched functions agree with the
original ones plus a fixed constant on a neighborhood of $V_0$,
and their differences have support in one fixed compact set. Apply
Theorem~\ref{thm:relative-capacity-estimate} to them.

For quasicontinuity, take continuous psh approximants decreasing to
the patched $u$, with the same collar. Choose indices $j_l$ such that
\[
 \operatorname{Cap}_{S,\beta,k}
       (\{u_{j_l}-u>2^{-l}\};\Omega)
       <\epsilon 2^{-l-1}.
\]
Each exceptional set is ordinary open: $u_{j_l}$ is continuous and
$u$ is upper semicontinuous, so $u_{j_l}-u$ is lower semicontinuous.
Their union $G$ is open and has capacity less than $\epsilon$ by
countable subadditivity. On $\Omega\setminus G$, the chosen
subsequence converges uniformly to the patched $u$. Its restriction
to $V_0$ is the original $u$ up to the fixed constant.
\end{proof}

\begin{remark}[Scope of the capacity conclusions]
The preceding results concern the Borel capacity defined by the
supremum of the constructed measures. They do not establish outer
regularity, continuity from above on compact sets, Choquet
capacitability, relative extremal functions, or an identification of
capacity-zero sets with ambient pluripolar sets. Degeneracy is allowed.
In particular, an arbitrary $S$ can be supported on an ambient
pluripolar set, and the relative measures can charge that set.

The relative energy lemma is a direct consequence of the first-stage
calculus. The quantitative estimate \eqref{eq:capacity-rate} and its
local quasicontinuity consequence provide an additional general
relative theorem. Their proof depends on the already proved
fixed-background symmetry, so they are not used retrospectively in
the proof of that symmetry.
\end{remark}

% END INCLUDED SOURCE: base_r2/base_r1/sections/relative_energy_agent


% BEGIN INCLUDED SOURCE: base_r2/base_r1/sections/capacity_stability
\section{Stability under local convergence in capacity}
\label{sec:capacity-stability}

This section is a \texttt{DIRECT\_COROLLARY} of the relative calculus,
mixed domination, and local quasicontinuity proved above. In particular,
it is not used in establishing any of those assertions. The background
$S$ remains fixed. Put $q=n-p$ and assume $q\geq1$.

\begin{lemma}[Top capacity controls each finite coefficient stage]
\label{lem:top-capacity-coefficients}
Let $\eta$ be an original compact smooth test of bidegree
$(q-r,q-r)$, $0\leq r\leq q$, presented by unit moment coordinates
$t_1,\ldots,t_N$ on an ordinary domain $\Omega$. Suppose these
coordinates are bounded on $\Omega$. Let
$a_1,\ldots,a_r\in\Cb(\Omega)$ have actual bounds $|a_i|\leq M$.
There is a constant $C_\eta$, independent of these potentials, such
that, for every Borel set $E\subset\operatorname{supp}\eta$,
\[
 |\mu_{T_{S,r}(a_1,\ldots,a_r),\eta}|(E)
       \leq C_\eta\operatorname{Cap}_S(E;\Omega).
\]
The constant may depend on $M$, the test, and the bounded moment
coordinates. For $r=0$ the tuple is empty.
\end{lemma}

\begin{proof}
Set $\rho=\frac12\sum_{a=1}^Nt_a^2$ and
$\beta=D\rho=\sum_a d't_a\wedge d''t_a$. The function $\rho$ is
bounded, smooth and classical psh: after every complete field extension
and curve pullback, each unit logarithm is harmonic, and this convex
function of their values is subharmonic. Finite coefficient-matrix
domination, as in Lemma~\ref{lem:finite-positive-controls}, gives an
inequality of measures on the test support,
\[
 |\mu_{T_{S,r}(a_1,\ldots,a_r),\eta}|
 \leq C\,T_{S,r}(a_1,\ldots,a_r)\wedge\beta^{q-r}.
\]
This is a measure inequality, not merely a total-mass bound: multiply
the positive test inequalities by every nonnegative smooth scalar and
use Radon uniqueness. Smooth factors commute with the recursion, so
the measure on the right is
\[
 T_{S,q}(a_1,\ldots,a_r,
             \underbrace{\rho,\ldots,\rho}_{q-r}).
\]
Shift the $a_i$ by $M$, and apply the mixed domination inequality
\eqref{eq:mixed-capacity-domination} in degree $q$. The bounds $2M$
and $\sup_\Omega\rho$ give the required constant. A zero bound
contributes a constant differentiated slot, whose product vanishes.
No positivity of $D\rho$ outside the test's moment directions is
needed.
\end{proof}

\begin{lemma}[A capacity-controlled multiplier limit]
\label{lem:capacity-multiplier}
Let $P_j\to P$ be positive currents of one fixed bidegree on $\Omega$,
with convergence on all original compact smooth complementary tests.
Fix such a compact test $\eta$ supported in $L\Subset\Omega$.
Assume that the coefficient measures for $P_j$ and $P$ satisfy
\[
 |\mu_{P_j,\eta}|(E),\ |\mu_{P,\eta}|(E)
       \leq C_\eta\operatorname{Cap}_S(E;\Omega)
       \qquad(E\subset L\text{ Borel}).
\]
Let $a_j,a$ be real-valued Borel functions, uniformly bounded near $L$, with
\[
 \operatorname{Cap}_S(L\cap\{|a_j-a|>\delta\};\Omega)
       \longrightarrow0\qquad(\delta>0).
\]
Assume that $a$ is continuous outside ordinary open sets of arbitrarily
small $\operatorname{Cap}_S(\,\cdot\,;\Omega)$ on a compact
neighborhood of $L$. Then
\[
 \langle a_jP_j,\eta\rangle\longrightarrow\langle aP,\eta\rangle.
\]
\end{lemma}

\begin{proof}
Weak convergence of the positive currents bounds their coefficient
variations on $L$ by one number $A$. If $|a_j|,|a|\leq M$, then
\[
 \left|\int(a_j-a)\,d\mu_{P_j,\eta}\right|
 \leq \delta A+2MC_\eta
       \operatorname{Cap}_S(L\cap\{|a_j-a|>\delta\};\Omega).
\]
First let $j\to\infty$, and then let $\delta\downarrow0$.

To handle the fixed multiplier $a$, choose an ordinary open exceptional
set $G$ of capacity less than $\epsilon$. On a compact neighborhood
$H$ of $L$, the restriction of $a$ to $H\setminus G$ is continuous.
The space $H$ is compact Hausdorff, hence normal. Tietze extension and
a cutoff give a continuous function $a_\epsilon$ near $L$, bounded
by $M$, which equals $a$ on $L\setminus G$. Thus
\[
 \left|\int(a-a_\epsilon)\,d\mu_{P_j,\eta}\right|,
 \left|\int(a-a_\epsilon)\,d\mu_{P,\eta}\right|
 \leq 2MC_\eta\epsilon.
\]
For fixed $a_\epsilon$, R, Lemma 3.1 gives convergence. Let
$\epsilon\downarrow0$ and combine the two estimates. Neither outer
regularity of the capacity nor a choice of a different representative
for $a$ is used.
\end{proof}

\begin{definition}[Local convergence in the relative capacity]
\label{def:local-capacity-convergence}
We say that real-valued Borel functions $a_j$ on $U$ converge to $a$ locally in
$S$-capacity if, for every relatively compact polynomial norm-box
domain $\Omega$ whose closure is contained in an ordinary algebraic
neighborhood $W\subset U$, every compact $L\Subset\Omega$, and
every $\delta>0$,
\[
 \operatorname{Cap}_S(L\cap\{|a_j-a|>\delta\};\Omega)
       \longrightarrow0.
\]
Within any prescribed finite unit-moment presentation one may first
choose such a neighborhood $W$, and then choose boxes with closure
contained in $W$. The quantifier covers all these boxes, including
their further shrinks. We do not identify this condition with
convergence for a single larger-domain capacity: such an identification
would require an additional comparison between capacities of domains.
\end{definition}

\begin{theorem}[Weighted stability in local relative capacity]
\label{thm:capacity-stability}
Let $0\leq k\leq q$. For $0\leq i\leq k$, let
$u_{i,j},u_i\in\Cb(U)$ have one actual uniform bound on every
compact subset of $U$, independent of $j$. Suppose
$u_{i,j}\to u_i$ locally in $S$-capacity in the sense of
Definition~\ref{def:local-capacity-convergence}. Then
\[
 u_{0,j}T_{S,k}(u_{1,j},\ldots,u_{k,j})
       \longrightarrow u_0T_{S,k}(u_1,\ldots,u_k)
\]
on every original compact smooth complementary test. No monotonicity
is required. Unweighted products have the same stability.
\end{theorem}

\begin{proof}
Localize a test to finitely many relatively compact norm boxes on
which its unit moment coordinates are defined and bounded. The assumed
local uniform bounds apply on their closures. On each piece,
Lemma~\ref{lem:top-capacity-coefficients} controls the coefficient
measures of every degree-$r$ product, $0\leq r\leq k$, uniformly
in $j$, and also controls its proposed limit.

Each $u_i$ is locally quasicontinuous by
Corollary~\ref{cor:relative-quasicontinuity}. If that corollary is
applied on still smaller boxes, domain monotonicity gives
$\operatorname{Cap}_S(G;\Omega)\leq
\operatorname{Cap}_S(G;\Omega')$ for $G\subset\Omega'\subset\Omega$.
A finite cover and countable subadditivity therefore give the ordinary
open exceptional set required near the test support, with arbitrarily
small capacity measured in $\Omega$. Choose a compact neighborhood of
the test support inside the union of the smaller cores; continuity
glues on its complement of that exceptional set. This supplies the
compact-neighborhood hypothesis of Lemma~\ref{lem:capacity-multiplier}.

Now induct on the number of differentiated slots. In degree zero the
background is fixed, and Lemma~\ref{lem:capacity-multiplier} applies
with $P_j=P=S$. At each later degree, weighted convergence at the
preceding degree passes through $D$ and gives unweighted convergence
of the new products. Apply Lemma~\ref{lem:capacity-multiplier} with
weight $u_{0,j}$ to obtain the weighted assertion at this degree.
The capacity controls and quasicontinuity just verified supply all its
hypotheses. Sum the finitely many localized tests.
\end{proof}

\begin{remark}
Theorem~\ref{thm:capacity-stability} strengthens the mode of convergence
available once the monotone theory has been established. It neither
gives compactness of arbitrary bounded psh families in this capacity,
nor proves that every weak limit of their Monge--Amp\`ere measures
comes from a bounded psh function. Those are different assertions.
The result does not change the fixed-background requirement.
\end{remark}

% END INCLUDED SOURCE: base_r2/base_r1/sections/capacity_stability


% BEGIN INCLUDED SOURCE: base_r2/base_r1/sections/outer_regularity_attempt
\section{The remaining compact-regularity obstruction}
\label{sec:outer-regularity-attempt}

This section records a precise reduction, not a proof of outer
regularity for the full Bedford--Taylor capacity. The initial capacity estimate concerns decreasing potential sequences.
The later nonmonotone stability theorem assumes convergence in capacity;
it does not extract a convergent subsequence from an arbitrary normalized
family. Compact continuity of the capacity instead permits a different
maximizing potential at every index. This distinction is essential.

\subsection{A measure-family reduction}

\begin{proposition}[Weak closure and compact outer regularity]
\label{prop:capacity-weak-closure}
Let $L$ be a compact metrizable space and let $\mathcal M$ be a
nonempty family of positive Radon measures on $L$ satisfying
\[
 \sup_{\mu\in\mathcal M}\mu(L)<\infty.
\]
For a Borel set $E\subset L$, put
\[
 c(E)=\sup_{\mu\in\mathcal M}\mu(E),
 \qquad
 c^*(K)=\inf_{\substack{G\supset K\\G\text{ open in }L}}c(G)
 \quad(K\subset L\text{ compact}).
\]
Let $\overline{\mathcal M}^{\,w}$ be the weak closure in the space
of finite measures. Then:
\begin{enumerate}
\item $c$ is inner regular on Borel sets:
\[
 c(E)=\sup_{\substack{K\subset E\\K\text{ compact}}}c(K).
\]
\item For every compact $K\subset L$,
\begin{equation}
\label{eq:compact-capacity-closure}
 c^*(K)=\sup_{\mu\in\overline{\mathcal M}^{\,w}}\mu(K).
\end{equation}
\item The following conditions are equivalent:
\begin{enumerate}
\item $c(K_j)\downarrow c(K)$ whenever $K_j\downarrow K$ are compact;
\item $c^*(K)=c(K)$ for every compact $K$;
\item every $\mu\in\overline{\mathcal M}^{\,w}$ satisfies
      $\mu(K)\leq c(K)$ for every compact $K$.
\end{enumerate}
\end{enumerate}
\end{proposition}

\begin{proof}
For the first assertion, Radon inner regularity gives
\[
 \sup_{\mu\in\mathcal M}\mu(E)
  =\sup_{\mu\in\mathcal M}
       \sup_{K\Subset E}\mu(K)
  =\sup_{K\Subset E}\sup_{\mu\in\mathcal M}\mu(K).
\]
This interchanges two suprema, which is legitimate. It does not
interchange a supremum with an infimum over open neighborhoods.

The mass bound makes $\overline{\mathcal M}^{\,w}$ compact. Suppose
$\mu_j\in\mathcal M$ converges weakly to $\mu$. For an open
$G\supset K$, the open-set inequality for weak convergence yields
\[
 \mu(K)\leq\mu(G)\leq\liminf_j\mu_j(G)\leq c(G).
\]
Taking the infimum over $G$ proves that the right side of
\eqref{eq:compact-capacity-closure} is at most the left side.

Choose decreasing open neighborhoods $G_j$ of $K$, with compact
closures $F_j$, such that $F_j\downarrow K$ and the $G_j$ form a
neighborhood basis of $K$. Choose $\mu_j\in\mathcal M$ with
\[
 \mu_j(G_j)\geq c(G_j)-1/j.
\]
After passing to a subsequence, $\mu_j$ converges weakly to some
$\mu\in\overline{\mathcal M}^{\,w}$. For each fixed $l$, all
sufficiently large $j$ satisfy $G_j\subset F_l$, so the closed-set
inequality gives
\[
 \lim_j c(G_j)
 \leq\limsup_j\mu_j(F_l)\leq\mu(F_l).
\]
Let $l\to\infty$. The measure is finite, hence
$\mu(F_l)\downarrow\mu(K)$, while
$\lim_jc(G_j)=c^*(K)$. This proves
\eqref{eq:compact-capacity-closure}.

Condition (a) implies (b) by applying it to compact neighborhoods
$F_j\downarrow K$ and using $c(K)\leq c^*(K)\leq c(F_j)$.
Conversely, assume (b) and let $K_j\downarrow K$. Given $\epsilon>0$,
choose an open $G\supset K$ with $c(G)\leq c(K)+\epsilon$.
The nested compact sets $K_j$ eventually lie in $G$; otherwise their
intersections with the closed complement of $G$ would have a common
point. Thus $\lim_jc(K_j)\leq c(K)+\epsilon$. The reverse inequality
is monotonicity. Finally, equivalence with (c) follows from
\eqref{eq:compact-capacity-closure}.
\end{proof}

\noindent\textbf{Status.} All assertions of the proposition are
\texttt{DIRECT\_COROLLARY} elementary reductions from Radon regularity
and weak compactness of bounded positive measure families. They locate
the missing analytic assertion and are not counted as new BT theorems.

\paragraph{Application to the constructed relative capacity.}
Fix its domain $\Omega$, its background $S$, and its trace form
$\beta$, and write
\[
 \mu_v=T_{S,k}(v,\ldots,v)\wedge\beta^{n-p-k},
 \qquad v\in\Cb(\Omega),\quad0\leq v\leq1.
\]
Choose an ordinary open core $W\Subset\Omega$ lying in one smaller
compact norm box, and an original smooth cutoff $\chi$, with
$0\leq\chi\leq1$, supported inside that box, and
$\chi=1$ on a neighborhood of $\overline W$. The measures
\[
 \mathcal M_\chi=\{\chi\mu_v:0\leq v\leq1\}
\]
have uniformly bounded total mass by relative CLN and can be viewed
on one compact metrizable norm-box neighborhood of
$\operatorname{supp}\chi$. Their capacity on subsets of $W$ is
exactly $\operatorname{Cap}_{S,\beta,k}(\,\cdot\,;\Omega)$.
Proposition~\ref{prop:capacity-weak-closure} therefore identifies the
precise remaining local assertion.

\begin{quote}
\textbf{Weak-closure domination (open).}
If $0\leq v_j\leq1$ are classical psh functions on the
\emph{fixed} domain $\Omega$, and
$\chi\mu_{v_j}\to\nu$ weakly, then, for every compact
$K\subset W$,
\[
 \nu(K)\leq
 \operatorname{Cap}_{S,\beta,k}(K;\Omega).
\]
\end{quote}

This condition is weaker than representing $\nu|_W$ as the
Monge--Amp\`ere measure of one admissible potential. It also retains
the original capacity domain: representing the limit by a bounded
psh function only on a smaller domain would bound it by that
\emph{larger} domain-restricted capacity, which is insufficient.
If weak-closure domination is proved, compact continuity from above
and outer regularity on compact sets follow immediately. Outer
regularity on all Borel sets and a capacitability theorem still need
their own argument; neither is assumed here.

\subsection{Cases where the reduction closes}

\begin{proposition}[Finite families and uniformly compact potential families]
\label{prop:restricted-family-regularity}
\begin{enumerate}
\item If $\mathcal M$ is finite, its capacity $c$ is outer regular
      on every Borel set.
\item If $\mathcal M$ is weakly compact, its capacity is continuous
      from above on compact sets and outer regular on compact sets.
\item In the relative Bedford--Taylor setting, let $\mathcal F$ be
      a compact family, for local uniform convergence, of continuous
      classical psh functions on $\Omega$ with values in $[0,1]$.
      The capacity obtained by restricting the supremum to
      $v\in\mathcal F$ has the two compact-set regularity properties
      on every relatively compact core.
\end{enumerate}
\end{proposition}

\begin{proof}
For (1), write $\mathcal M=\{\mu_1,\ldots,\mu_N\}$. Given a
Borel $E$ and $\epsilon>0$, choose open $G_i\supset E$ with
$\mu_i(G_i)\leq\mu_i(E)+\epsilon$. The finite intersection
$G=\bigcap_iG_i$ is open and satisfies
$c(G)\leq c(E)+\epsilon$. Assertion (2) follows directly from
Proposition~\ref{prop:capacity-weak-closure}.

For (3), uniform stability of the continuous fixed-background
products makes $v\mapsto\chi\mu_v$ continuous from local uniform
convergence to weak convergence of finite measures on the compact
cutoff support. Its image is compact, so (2) applies. These are
restricted-family results. No local-uniform compactness of the
entire normalized classical psh class is asserted.
\end{proof}

\subsection{Why the direct extensions fail}

\paragraph{Radon regularity is not uniform family regularity.}
On $L=[0,1]$, take
\[
 A=\mathbf Q\cap(0,1],\qquad
 \mathcal M=\{\delta_x:x\in A\}.
\]
The set $A$ is dense in $L$, but $0\notin A$. Every measure is Radon
and has mass one. Nevertheless,
\[
 c(\{0\})=0,\qquad c([0,1/m])=1,\qquad c^*(\{0\})=1.
\]
The weak closure contains $\delta_0$, for example as the limit of
$\delta_{1/j}$, exactly accounting for the
defect in \eqref{eq:compact-capacity-closure}. Thus positivity,
Radon regularity, uniform mass bounds, countable subadditivity, and
continuity from below do not by themselves imply compact continuity
from above. This is an abstract obstruction to a proposed proof,
not a counterexample to the full non-Archimedean BT capacity.

\paragraph{Finite maxima do not supply an admissible upper envelope.}
On the ordinary unit disc in $\mathbf A_K^{1,\mathrm{an}}$, the
bounded continuous classical psh functions
\[
 a_j=\max\{-1,j^{-1}\log|T|\}
\]
increase pointwise. Their limit is zero at every point with $T\ne0$
and is $-1$ at the rigid zero. It is not upper semicontinuous.
Thus closure under finite maxima and bounded decreasing limits
does not justify taking the pointwise supremum of a bounded family.
An upper-semicontinuous envelope theorem, with its all-extension
curve condition, remains an additional assertion.

\paragraph{Potential order does not order curvature measures.}
Even if a tail envelope is available, $v_j\leq w_j$ does not imply
$\mu_{v_j}\leq\mu_{w_j}$. On an annulus crossing the unit circle,
a suitable positive rescaling of $\max\{0,\log|T|\}$ lies in
$[0,1]$ and has a nonzero positive Laplacian, whereas the larger
constant function $1$ has zero Laplacian. An envelope route therefore
also needs a contact or extremal-measure comparison; pointwise
domination alone is insufficient.

\paragraph{Status.}
The measure-family reduction and the restricted-family propositions
have complete proofs above. They isolate the weak-closure domination
lemma as the next analytic target. For the full normalized
$\Cb(\Omega)$ family, that lemma, compact continuity from above,
and outer regularity remain \texttt{OPEN}. The failed arguments
recorded here are not used elsewhere in the report.

% END INCLUDED SOURCE: base_r2/base_r1/sections/outer_regularity_attempt


% BEGIN INCLUDED SOURCE: base_r2/base_r1/sections/maximum_locality
\section{Maximum locality and relative comparison}
\label{sec:maximum-locality}

This section uses the fixed-background mixed-current theorem proved above. Thus
\(S\) is a fixed symmetric positive original CLD current of bidegree
\((p,p)\), with \(d'S=d''S=0\), and all potentials belong to
\(\mathcal C_b\). Positivity means the cone used in R, not strong positivity.
No varying-background convergence is used.

For a coefficient-Radon current \(T\) and a Borel set \(E\), write
\(\mathbf 1_E T\) for its Borel multiple. In particular,
\[
 \mu_{\mathbf 1_E T,\eta}=\mathbf 1_E\mu_{T,\eta}
\]
for every original compact smooth complementary test \(\eta\). This notation
does not assert that \(E\) is ordinarily open, nor that \(\mathbf 1_E T\) is
closed.

\begin{theorem}[Maximum locality for a fixed background]
\label{thm:maximum-locality-relative}
Let \(1\leq k\leq n-p\), let \(u,v,w_2,\ldots,w_k\in\mathcal C_b(V)\), and
put \(M=\max(u,v)\) and \(E=\{u>v\}\). Then
\[
 \mathbf 1_E T_{S,k}(M,w_2,\ldots,w_k)
 =\mathbf 1_E T_{S,k}(u,w_2,\ldots,w_k).
\tag{ML}
\]
The equality holds as coefficient-Radon currents on the original CLD test
sheaf. The corresponding identity on \(\{v>u\}\) holds with \(v\) in the
first differentiated slot.
\end{theorem}

\noindent\textbf{Status.} \texttt{NEW\_PROVED} consequence of the
fixed-background theorem. Its \([1]\)-background case was already
a \texttt{DIRECT\_COROLLARY} of B; see below.

\begin{proof}
The assertion is local on ordinary open sets. By B, Theorem 2.5, choose on one
ordinary neighborhood simultaneous decreasing continuous semipositive PL
approximants
\[
 u_j\downarrow u,\qquad v_j\downarrow v,\qquad w_{i,j}\downarrow w_i.
\]
Finite intersections of the fixed neighborhoods suffice. Set
\[
 M_j=\max(u_j,v_j),\qquad h_j=M_j-v_j,
 \qquad h=M-v=(u-v)_+.
\]
Both \(M_j\) and \(v_j\) decrease to locally bounded classical psh functions.
The functions \(h_j\) themselves need not be monotone; no monotonicity of
\(h_j\) is asserted or used.

Write
\[
 A_j=T_{S,k}(M_j,w_{2,j},\ldots,w_{k,j}),\qquad
 B_j=T_{S,k}(u_j,w_{2,j},\ldots,w_{k,j}).
\]
The set \(E_j=\{u_j>v_j\}\) is ordinarily open. Ordinary locality gives
\(A_j=B_j\) there, while \(h_j=0\) on its complement. Hence, as
coefficient-Radon currents,
\[
 h_j A_j=h_j B_j.
\tag{1}
\]
Apply weighted joint decreasing continuity to each of
\(M_j A_j,v_j A_j,M_j B_j,v_j B_j\). Subtracting the resulting limits in
(1) gives
\[
 h(A-B)=0,
 \quad A=T_{S,k}(M,w_2,\ldots,w_k),\quad
 B=T_{S,k}(u,w_2,\ldots,w_k).
\tag{2}
\]
This passage uses two admissible psh weights, not a convergence theorem for
an arbitrary changing Borel weight.

For an original compact test \(\eta\), multiply (2) by arbitrary original
smooth scalar functions. Uniqueness of its Radon coefficient measure gives
\(h\mu_{A-B,\eta}=0\). On \(E_m=\{h\geq 1/m\}\), the Borel function
\(1/h\) is bounded. Therefore
\(\mathbf 1_{E_m}\mu_{A-B,\eta}=0\). Since
\(E=\bigcup_{m\geq1}E_m\), it follows that
\(\mathbf 1_E\mu_{A-B,\eta}=0\). This proves (ML) on every original
test, and the local identities glue. Interchanging \(u\) and \(v\) proves
the final assertion.
\end{proof}

\begin{theorem}[Relative comparison with a compact boundary gap]
\label{thm:relative-comparison}
Let \(m=n-p\geq1\). For \(a\in\mathcal C_b(V)\), put
\[
 \mu_a=T_{S,m}(a,\ldots,a).
\]
Let \(u,v\in\mathcal C_b(V)\). Assume that for every \(\varepsilon>0\),
\[
 \overline{\{u<v-\varepsilon\}}\ \text{is compact and contained in }V.
\tag{BG}
\]
The closure in (BG) is the ordinary topological closure of a Borel set. Then
\[
 \int_{\{u<v\}}d\mu_v\ \leq\ \int_{\{u<v\}}d\mu_u.
\tag{CP}
\]
The integrals are allowed to be infinite. If \(V\Subset U\) and
\(\liminf_{x\to\partial V}(u(x)-v(x))\geq0\), condition (BG) holds.
\end{theorem}

\noindent\textbf{Status.} \texttt{NEW\_PROVED} consequence of the
fixed-background theorem. It requires no strictly psh exhaustion and no
capacity theorem.

\begin{proof}
Fix \(\varepsilon>0\) and let \(w=\max(u,v-\varepsilon)\). Form the
positive closed current
\[
 Q=\sum_{j=0}^{m-1}
 T_{S,m-1}(\underbrace{w,\ldots,w}_{j},
                  \underbrace{u,\ldots,u}_{m-1-j}).
\]
For \(m=1\), this means \(Q=S\). Multilinearity, symmetry, and the
recursive definition give the telescoping identity
\[
 \mu_w-\mu_u=D\bigl((w-u)Q\bigr).
\tag{3}
\]
Here \(w-u\) is a bounded Borel function on its compact support; the identity
is obtained by subtracting two well-defined Borel multiples. It does not
require \(w-u\) to be psh.

By (BG), the support of \((w-u)Q\) is contained in a compact subset of
\(V\). Choose an original smooth compact cutoff \(\chi\) equal to one
near that support. Equation (3) and distributional differentiation give
\[
 \langle\mu_w-\mu_u,\chi\rangle
 =\langle(w-u)Q,D\chi\rangle=0.
\tag{4}
\]
The signed measure \(\mu_w-\mu_u\) has compact support, since the two
potentials agree on the ordinary complement of
\(\overline{\{u<v-\varepsilon\}}\). Thus (4) is its total signed mass.
No subtraction of two possibly infinite total masses on \(V\) is made.

Set
\[
 A_\varepsilon=\{u<v-\varepsilon\},\qquad
 B_\varepsilon=\{u=v-\varepsilon\},\qquad
 C_\varepsilon=\{u>v-\varepsilon\}.
\]
The set \(A_\varepsilon\cup B_\varepsilon\) is contained in the compact
set \(\overline{\{u<v-\varepsilon/2\}}\), so every measure below has
finite mass on it. Maximum locality and symmetry give
\[
 \mathbf 1_{A_\varepsilon}\mu_w
   =\mathbf 1_{A_\varepsilon}\mu_v,
 \qquad
 \mathbf 1_{C_\varepsilon}\mu_w
   =\mathbf 1_{C_\varepsilon}\mu_u.
\]
For the first equality, constants in differentiated slots have zero
curvature. Applying (4) and partitioning into these three Borel sets yields
\[
 \mu_v(A_\varepsilon)+\mu_w(B_\varepsilon)
   =\mu_u(A_\varepsilon\cup B_\varepsilon).
\]
Positivity therefore implies
\[
 \mu_v(\{u<v-\varepsilon\})
 \leq\mu_u(\{u\leq v-\varepsilon\}).
\tag{5}
\]
Along any positive sequence \(\varepsilon\downarrow0\), both sets in
(5) increase to \(\{u<v\}\). Continuity from below of positive measures
proves (CP).

For the final assertion, the boundary hypothesis means that every boundary
point has a neighborhood in which \(u-v>-\varepsilon\). Compactness of
\(\partial V\) gives a finite collar of this kind. Its complement in
\(\overline V\) is compact and contained in \(V\), proving (BG).
\end{proof}

\paragraph{What comparison does not prove.}
The theorem does not by itself imply pointwise domination or uniqueness
for an arbitrary background. Section~\ref{sec:capacity-domination}
proves precise conclusions modulo capacity zero. For instance, \(S=0\) makes both Monge--Amp\`ere measures zero
for all potentials. More generally, relative products can only detect
the support and the curvature directions of \(S\). Domination requires
an additional mechanism that detects a strict inequality of potentials;
the complex proof uses strictly psh perturbations for this purpose.

\subsection{The part already implied by B}
\label{subsec:baseline-maximum-locality}

\begin{proposition}[Maximum locality for the previously constructed products]
On every analytic space permitted by B, identity (ML) for B's products
\(T_k=T_{[1],k}\) already follows from B, Theorems 2.5--2.7, without the
fixed-background extension proved in this report. The same is true when
the background is itself one of B's mixed currents.
\end{proposition}

\noindent\textbf{Status.} \texttt{DIRECT\_COROLLARY}.

\begin{proof}
Repeat the proof of Theorem~\ref{thm:maximum-locality-relative} with
\(S=[1]\). Every input of that proof is already supplied by B:
Theorem 2.5 supplies decreasing continuous regularizations on a fixed
ordinary neighborhood; Theorem 2.6 supplies ordinary locality and
coefficient-Radon products; Theorem 2.7 supplies weighted joint decreasing
continuity on all original complementary tests. The last Borel division
argument uses the Borel multiplication from B, Definition 2.4 and
Lemma 5.1. Thus the proof does not use any theorem first established in
this report. If the background is a B mixed current, absorb its fixed
potentials into the differentiated slots and use B's symmetry and weighted
continuity in those additional slots.
\end{proof}

% END INCLUDED SOURCE: base_r2/base_r1/sections/maximum_locality


% BEGIN INCLUDED SOURCE: base_r2/base_r1/sections/capacity_domination
\section{Domination and uniqueness modulo capacity zero}
\label{sec:capacity-domination}

There is a precise consequence of comparison that does not require a
nondegeneracy theorem. It separates the proved measure-theoretic
conclusion from the additional condition needed for pointwise order.
The results in this section are \texttt{DIRECT\_COROLLARY} of
Theorem~\ref{thm:relative-comparison} and the positive mixed calculus.

Let $S$ have bidegree $(p,p)$ and put $q=n-p\geq1$. On an ordinary
open $V$, use the top-degree capacity
\[
 \operatorname{Cap}_S(E;V)
   =\sup_{\substack{h\in\Cb(V)\\0\leq h\leq1}}
       \int_E T_{S,q}(h,\ldots,h),
       \qquad E\subset V\text{ Borel}.
\]
This definition makes sense without a moment frame. Relative CLN gives
compact finiteness, and countable subadditivity is unchanged. Write
$\mu_a=T_{S,q}(a,\ldots,a)$. For two potentials $a,b$, the compact-gap
condition means
\[
 \overline{\{a<b-\delta\}}\Subset V
             \qquad\text{for every }\delta>0,
\]
where the closure is the ordinary topological closure.

\begin{proposition}[Domination modulo capacity zero]
\label{prop:capacity-domination}
Let $u,v\in\Cb(V)$ satisfy the compact-gap condition in the order
$(u,v)$. If $\mu_u(\{u<v\})=0$, then
\[
 \operatorname{Cap}_S(\{u<v\};V)=0.
\]
This does not assert that the strict inequality set is empty.
\end{proposition}

\begin{proof}
Fix $\delta>0$, $0<\epsilon<\delta$, and any admissible
$h\in\Cb(V)$ with $0\leq h\leq1$. Put
\[
 v_{\delta,\epsilon}=v+\epsilon(h-1)-\delta,
 \qquad E=\{u<v_{\delta,\epsilon}\}.
\]
Then
\[
 \{u<v-2\delta\}\subset E\subset\{u<v-\delta\}.
\]
The pair $(u,v_{\delta,\epsilon})$ satisfies the compact-gap condition.
Comparison and positivity of the mixed expansion give
\[
 0\leq\epsilon^q\int_{\{u<v-2\delta\}}\mu_h
 \leq\int_E\mu_{v_{\delta,\epsilon}}
 \leq\int_E\mu_u=0.
\]
Constants disappear in differentiated slots, which explains the
$\epsilon^q\mu_h$ term. Since $h$ was arbitrary, the capacity of
$\{u<v-2\delta\}$ is zero. Take $\delta$ through a sequence
decreasing to zero and use countable subadditivity. All sets are Borel;
none is assumed ordinary open.
\end{proof}

\begin{proposition}[Uniqueness modulo capacity zero]
\label{prop:capacity-uniqueness}
Let $u,v\in\Cb(V)$ satisfy the compact-gap condition in both orders.
If $\mu_u=\mu_v$ as Radon measures, then
\[
 \operatorname{Cap}_S(\{u\ne v\};V)=0.
\]
\end{proposition}

\begin{proof}
Fix $\delta>0$, $0<\epsilon<\delta$, and an admissible $h$ as above.
Set $u_{\delta,\epsilon}=u+\epsilon(h-1)-\delta$ and
$E=\{v<u_{\delta,\epsilon}\}$. Then
\[
 \{v<u-2\delta\}\subset E\subset\{v<u-\delta\}\Subset V,
\]
where the last compact containment refers to its closure. Comparison
for $(v,u_{\delta,\epsilon})$ and positivity yield
\[
 \int_E\mu_u+\epsilon^q\int_E\mu_h
 \leq\int_E\mu_{u_{\delta,\epsilon}}
 \leq\int_E\mu_v=\int_E\mu_u.
\]
The term being cancelled is finite, because $E$ has compact closure
in $V$. Hence $\mu_h(\{v<u-2\delta\})=0$. Taking the supremum
over $h$ and then a countable union in $\delta$ gives capacity zero
for $\{v<u\}$. Interchange $u$ and $v$ and take a union.
\end{proof}

\begin{corollary}[A conditional pointwise conclusion]
\label{cor:conditional-pointwise-domination}
Assume, in addition, the following property for $(S,V)$: every
nonempty strict gap $\{a<b-\delta\}$, for $a,b\in\Cb(V)$ satisfying
the compact-gap condition and $\delta>0$, has positive
$\operatorname{Cap}_S(\,\cdot\,;V)$. Then
Proposition~\ref{prop:capacity-domination} gives $u\geq v$
pointwise, and Proposition~\ref{prop:capacity-uniqueness} gives
$u=v$ pointwise.
\end{corollary}

\begin{proof}
If a prohibited pointwise inequality held at a point, it would hold
there with some positive margin $\delta$. Its strict gap is a subset
of the capacity-zero set already proved, contradicting the additional
property.
\end{proof}

The last corollary has status \texttt{CONDITIONAL}. The condition on
strict gaps is a new research obligation for any proposed class of
domains and backgrounds; positivity on nonempty ordinary opens alone
does not identify it. It fails for $S=0$, and may fail away from the
support of a nonzero background. The propositions retain their literal
meaning even when the capacity degenerates.

% END INCLUDED SOURCE: base_r2/base_r1/sections/capacity_domination


% BEGIN INCLUDED SOURCE: base_r2/base_r1/sections/counterexamples
\section{Counterexamples to unrestricted extensions}
\label{sec:counterexamples}

All three examples use the original field \(K\), with its prescribed
absolute value. They concern statements stronger than the fixed-background
theorem.

\subsection{A bounded discontinuous classical psh function over \texorpdfstring{$K$}{K}}

\begin{proposition}[A sequence of branch wells]
\label{prop:branch-wells}
There exist an ordinary open disc \(V\subset\mathbf A_K^{1,\mathrm{an}}\), a
function \(w\in\mathcal C_b(V)\), a Gauss point \(\xi\in V\), and rigid
points \(x_i\in V\) such that
\[
 x_i\longrightarrow\xi,\qquad w(\xi)=0,\qquad w(x_i)=-1.
\]
The function \(w\) is the pointwise decreasing limit of continuous classical
psh functions on one fixed disc, uniformly bounded there.
\end{proposition}

\begin{proof}
Let \(\widetilde K\) be the finite residue field. Choose pairwise distinct
monic irreducible polynomials
\(\widetilde P_i\in\widetilde K[T]\), and lift them to monic polynomials
\(P_i\in K^\circ[T]\). Write \(d_i=\deg P_i\) and \(a_i=2^{-i}\), with
indices starting at \(i=1\). Thus \(\sum_i a_i=1\). Choose \(R>1\), put
\(C=\log R\), and work on
\[
 V=\{x:|T(x)|<R\}.
\]
For example one may take \(R=|\pi|^{-1}\); no normalization of the given
absolute value is made. Set
\[
 f_i(x)=\frac{a_i}{d_i}
       \max\!\left\{\log|P_i(x)|,-\frac{d_i}{a_i}\right\}.
\tag{7}
\]
At a zero of \(P_i\), the first entry is interpreted as \(-\infty\).
Each \(f_i\) is continuous and classical psh: logarithms of analytic
functions are classical psh, and this class is closed under maxima with
constants and positive scalar multiples (Gubler--Rabinoff,
Proposition 4.2).

On \(\{|T|\leq1\}\), one has \(-1\leq f_i\leq0\). Moreover, at any
point of this closed disc at most one \(f_i\) is nonzero. Indeed, for
\(i\ne j\), coprimality of the reductions gives integral polynomials
\(A,B,H\) with
\[
 1=A P_i+B P_j+\pi H.
\]
If both \(|P_i(x)|<1\) and \(|P_j(x)|<1\), evaluation at a point with
\(|T(x)|\leq1\) contradicts the non-Archimedean triangle inequality.
Consequently every finite sum of the \(f_i\) is bounded below by \(-1\)
on the closed unit disc.

If \(1<|T(x)|<R\), monicity and integrality give
\[
 |P_i(x)|=|T(x)|^{d_i},\qquad f_i(x)=a_i\log|T(x)|.
\]
In particular, \(f_i\leq a_i C\) everywhere on \(V\). Define
\[
 w_N=C+\sum_{i=1}^N(f_i-a_i C)
     =\sum_{i=1}^N f_i+C\sum_{i>N}a_i.
\tag{8}
\]
The functions \(w_N\) are continuous classical psh on the same open disc,
and
\[
 w_{N+1}-w_N=f_{N+1}-a_{N+1}C\leq0,
 \qquad -1\leq w_N\leq C.
\]
Thus \(w=\lim_N w_N\) is locally bounded classical psh, by the decreasing
limit property in Gubler--Rabinoff, Proposition 4.2(4).

Let \(\xi\) be the Gauss point of the closed unit disc. Since
\(|P_i(\xi)|=1\), equations (7)--(8) give \(w(\xi)=0\). Let \(x_i\)
be the rigid Berkovich point defined by \(P_i\). The irreducibility of
\(\widetilde P_i\) implies irreducibility of \(P_i\) over \(K\).
The point \(x_i\) is a point of the original \(K\)-analytic line; using
its finite residue extension to describe its seminorm is not a base
change of the analytic space. At this point, \(f_i(x_i)=-1\), while
\(f_j(x_i)=0\) for \(j\ne i\). Therefore \(w(x_i)=-1\).

It remains to verify the convergence of the points. For a nonzero
\(F\in K[T]\), multiply by a scalar so that its Gauss norm equals one.
Its reduction is then a nonzero polynomial over \(\widetilde K\).
Only finitely many distinct irreducibles \(\widetilde P_i\) divide that
reduction. For every other \(i\), evaluation at \(x_i\) is a unit,
so \(|F(x_i)|=\|F\|_{\mathrm{Gauss}}\). Hence
\(|F(x_i)|\to|F(\xi)|\) for every polynomial \(F\). This is precisely
convergence to \(\xi\) in the topology of the Berkovich affine line.
It follows in particular that \(w\) is discontinuous at \(\xi\).
\end{proof}

\begin{corollary}[Failure for varying backgrounds]
\label{cor:varying-background-false}
Weak convergence of positive closed backgrounds does not, by itself,
imply weighted convergence for a fixed locally bounded classical psh
weight. This fails already with no differentiated potential.
\end{corollary}

\noindent\textbf{Status of the unrestricted assertion.} \texttt{FALSE}.

\begin{proof}
On the curve in Proposition~\ref{prop:branch-wells}, take
\(S_i=\delta_{x_i}\) and \(S=\delta_\xi\). These are symmetric positive
top-degree original currents, and are separately \(d'\)- and
\(d''\)-closed for degree reasons. The convergence \(x_i\to\xi\)
gives \(S_i\to S\) on every compact original smooth scalar test.
Nevertheless,
\[
 wS_i=-\delta_{x_i}\longrightarrow-\delta_\xi,
 \qquad wS=0.
\]
Thus the failure occurs for \(p=n=1\) and \(k=0\), with a constant
sequence of admissible weights. It does not involve changing a function
on a null set, changing the base field, or changing the test sheaf.
\end{proof}

\begin{remark}[A strict inequality set need not be ordinary open]
The function $w$ above satisfies $w(\xi)=0$, $w(x_i)=-1$, and
$x_i\to\xi$. Thus $\{w>-1/2\}$ contains $\xi$ but contains no
$x_i$, and is not an ordinary open neighborhood of $\xi$. Both $w$
and the constant $-1/2$ are bounded classical psh functions. This
checks explicitly the distinction used in maximum locality.
\end{remark}

\subsection{A nonzero relative product carried by an ambient pluripolar set}

\begin{proposition}[Concentration on a smooth divisor]
\label{prop:polar-relative-mass}
For every \(m\geq1\), there are a smooth \((m+1)\)-dimensional ordinary
analytic open \(V\), a nonzero positive closed \((1,1)\)-current \(S\),
and bounded continuous classical psh functions \(u_1,\ldots,u_m\) such
that \(T_{S,m}(u_1,\ldots,u_m)\) is a nonzero positive measure supported
on an ambient pluripolar set.
\end{proposition}

\noindent\textbf{Status of universal ambient non-concentration.}
\texttt{FALSE}, even for continuous potentials and an integration-current
background.

\begin{proof}
Choose \(0<r<1<R\) and \(R_0>0\). With coordinates
\(z_0,z_1,\ldots,z_m\), put
\[
 V=\{|z_0|<R_0\}\times
       \prod_{i=1}^m\{r<|z_i|<R\},
 \qquad Y=\{z_0=0\}\cap V.
\]
These are ordinary opens in the relevant smooth algebraic
analytifications. Let \(\iota:Y\hookrightarrow V\) be the closed
embedding and take \(S=\iota_*[1]_Y\). This is a symmetric positive
closed original CLD current. The set \(Y\) is ambient pluripolar,
since the nontrivial classical psh function \(\log|z_0|\) equals
\(-\infty\) precisely there.

Set
\[
 u_i=\max\{0,\log|z_i|\},\qquad 1\leq i\leq m.
\]
Each \(u_i\) is bounded and continuous on \(V\). Pushforward commutes
with distributional \(D\), and continuous multiplication satisfies the
projection formula. Recursion therefore gives
\[
 T_{S,m}(u_1,\ldots,u_m)
   =\iota_* T^Y_{[1],m}(u_1|_Y,\ldots,u_m|_Y).
\tag{9}
\]
The right side can be calculated on
\((\mathbf P^1_K)^m\). For the standard model
\((\mathbf P^1_{K^\circ})^m\), take the pullback of
\(\mathcal O(1)\) from each of its \(m\) factors, with its standard
model metric. In the affine frames, their weights are the displayed
\(u_i\). The special fiber is irreducible and reduced, and the mixed
intersection degree of these \(m\) line bundles is one. The model
measure formula, together with R, Lemma 3.3, identifies their original
CLD mixed measure with the Dirac mass at the product Gauss point
\(\xi_Y\). This point lies in the displayed polyannulus. Restricting
to \(Y\) and using (9) gives
\[
 T_{S,m}(u_1,\ldots,u_m)=\delta_{\iota(\xi_Y)}.
\]
In particular it has positive mass on \(Y\).
\end{proof}

\paragraph{Implication for capacity.}
An ambient pluripolar set need not be negligible for a relative
capacity. Indeed, in this example a positive rescaling of
\(u_1+\cdots+u_m\) lies between zero and one on \(V\), and its
\(m\)-fold curvature relative to \(S\) has a positive contribution
on \(Y\). Any non-concentration theorem must impose a condition on
the background, or formulate polarity intrinsically on its support.

\subsection{The unqualified capacity can vanish identically}

\begin{proposition}[Compact-domain degeneration]
\label{prop:compact-capacity-zero}
Let \(X\) be a connected compact boundaryless \(K\)-analytic space on
which the classical maximum principle holds, for example
\(X=\mathbf P_K^{n,\mathrm{an}}\) with \(n\geq1\). Every finite
locally bounded classical psh function on \(X\) is constant. Hence
\[
 \sup_{\substack{u\in\mathcal C_b(X)\\0\leq u\leq1}}
       \int_E (Du)^n=0
\]
for every Borel set \(E\subset X\).
\end{proposition}

\noindent\textbf{Status of nondegeneracy on every ordinary domain.}
\texttt{FALSE}. The calculation itself is a
\texttt{DIRECT\_COROLLARY} of the classical maximum principle.

\begin{proof}
An upper-semicontinuous real function on a compact space attains its
maximum \(M\). Its maximum set \(F=\{u=M\}\) is nonempty and closed.
Every point of \(F\) is interior because \(X\) is boundaryless.
Gubler--Rabinoff, Theorem 6.1, makes \(u\) locally constant at such a
point. Thus \(F\) is also open, and connectedness gives \(F=X\).
The curvature of a constant vanishes, proving the capacity assertion.
\end{proof}

The entire space \(X\) is not pluripolar when the definition excludes
identically \(-\infty\) local witnesses. Thus an unqualified equivalence
between zero capacity and pluripolarity also fails for this capacity.
Two different constants have the same zero Monge--Amp\`ere measure;
with empty boundary, no Dirichlet uniqueness follows without a
normalization or additional geometric data.

% END INCLUDED SOURCE: base_r2/base_r1/sections/counterexamples


% BEGIN INCLUDED SOURCE: base_r2/base_r1/sections/comparison
\section{Theorem-by-theorem comparison with the classical theory}
\label{sec:comparison}

\subsection{Scope and source conventions}

The comparison below separates the statement of a theorem from the category
in which it is proved. All complex statements use their authors' normalization
of \(dd^c\); replacing it by a positive constant multiple changes numerical
constants, not the conclusions listed here. All statements proved in this
report use \(D=d'd''\), the prescribed absolute value on the finite extension
\(K/\mathbf Q_p\), and the original compactly supported CLD smooth test
forms. No renormalization of the valuation is implicit.

The following source versions have been read and pinned: BT76 and BT82 are
the original journal articles; Dem12 is the June 21, 2012 version; CLD12 is
arXiv:1204.6277v1; CLD25 is arXiv:1204.6277v2; GR25 is
arXiv:2506.13548v2; Wu26 is arXiv:2609.27610v1; GK17 is the
arXiv:1406.7637v3 text of the subsequently published article; BGJK21 is
arXiv:2102.07392v1. The theorem numbers below refer to these versions.
The source register records URLs, hashes, and verification notes. Page
numbers are printed page numbers, which agree with PDF page numbers in the
quoted classical and CLD files.

\paragraph{The 2025 test-category issue.}
CLD25 develops Chapter 8 on its tropical spaces and subsequently applies
that formalism to analytic spaces; it also develops harmonic tropicalizations
and additional differential structures. Its theorems are recorded as
\texttt{EXTERNAL\_KNOWN} in their own stated categories. An arbitrary
functional on the original CLD12 test sheaf is not automatically a current
on an enlarged test sheaf. Accordingly, no CLD25 assertion is used here to
extend an arbitrary original current \(S\) without a proved comparison of
the relevant current categories. The proofs in this report instead use the
original-category inputs R and B. Products arising from common smooth or
model data can be compared by their actual definitions; this does not prove
an extension theorem for arbitrary currents.

\paragraph{Internal provenance.}
R, Theorem 2.9, provides continuous-potential regularization in its stated
algebraic setting, with R, Lemma 3.1, supplying the original positive-cone
coefficient-Radon and continuous-multiplier control. B, Theorems 2.5--2.7,
provide classical-psh regularization and the entire bounded-potential
calculus with background \([1]\), including weighted joint decreasing
continuity. These are \texttt{INTERNAL\_INPUT}, not new results of this
round. A2, Theorem 1.1, is already an established input under its number-field,
integrable-adelic, fixed-boundary hypotheses. It is not an open direction in
the present comparison.

\subsection{Status matrix}

Each status below refers to the precise assertion in its cell. In particular,
an \texttt{OPEN} entry is not a claim that every special case is unknown in
the literature, and a \texttt{FALSE} entry concerns the explicitly unrestricted
version. Local capacity and energy refinements proved elsewhere in this
report have their own hypotheses; they are not silently identified with a
nondegenerate global Bedford--Taylor capacity.

\begin{longtable}{@{}>{\raggedright\arraybackslash}p{.05\textwidth}>{\raggedright\arraybackslash}p{.24\textwidth}>{\raggedright\arraybackslash}p{.31\textwidth}>{\raggedright\arraybackslash}p{.31\textwidth}@{}}
\hline
 & Classical target & Established non-Archimedean baseline & Assessment in this project\\
\hline
\endfirsthead
\hline
 & Classical target & Established non-Archimedean baseline & Assessment in this project\\
\hline
\endhead
1 & Locally bounded psh functions and decreasing regularization &
GR25 gives the classical curve-test class; B 2.5 identifies it with its
decreasing semipositive PL regularization class. &
\texttt{PROVED / INTERNAL\_INPUT}. Continuous original-current positivity
is covered by R in R's scope; unrestricted distributional positivity is not
silently equated with classical psh.\\[3pt]
2 & Positive closed currents and mixed products &
CLD12 treats approachable continuous potentials; CLD25 treats continuous
potentials on arbitrary backgrounds in its category; B 2.6 treats \([1]\). &
Arbitrary fixed original positive closed \(S\), bounded classical psh
potentials: \texttt{PROVED / NEW\_PROVED} in the fixed-background theorem.
Ordered existence alone is a \texttt{DIRECT\_COROLLARY} of R/B.\\[3pt]
3 & Chern--Levine--Nirenberg estimates &
CLD12 5.6.2 gives the smooth estimate; CLD25 8.4.1 gives a relative smooth
estimate in its category. &
Original coefficient-Radon relative mass control for arbitrary fixed
\(S\): \texttt{PROVED / DIRECT\_COROLLARY}. No global strictly psh form is assumed.\\[3pt]
4 & Symmetry and weighted monotone convergence &
B 2.6--2.7 gives the full \([1]\)-background result; CLD25 8.7.2 concerns
local uniform convergence of continuous potentials. &
Fixed arbitrary original \(S\): \texttt{PROVED / NEW\_PROVED}.
Unrestricted varying-background weighted stability:
\texttt{FALSE}, Corollary~\ref{cor:varying-background-false}.\\[3pt]
5 & Locality and maximum principles &
Ordinary locality is in B; strict-set maximum locality is already a
\texttt{DIRECT\_COROLLARY} of B 2.5--2.7. GR25 6.1 gives the classical
psh maximum principle. &
The extension of Borel maximum locality to arbitrary fixed \(S\):
\texttt{PROVED / NEW\_PROVED},
Theorem~\ref{thm:maximum-locality-relative}. The set \(\{u>v\}\) is Borel;
ordinary openness is not asserted.\\[3pt]
6 & Bedford--Taylor capacities &
The mixed measures supply candidate set functions; neither a formal
supremum nor CLN alone proves nondegeneracy or full Choquet continuity. &
The stated Borel axioms and mixed domination are \texttt{PROVED}; the
norm-box estimate in Theorem~\ref{thm:relative-capacity-estimate} is
\texttt{PROVED / NEW\_PROVED}. Outer regularity, continuity from above
on compact sets, and Choquet capacitability remain \texttt{OPEN}.
Nondegeneracy on every analytic domain: \texttt{FALSE},
Proposition~\ref{prop:compact-capacity-zero}. A complete intrinsic domain
classification remains \texttt{OPEN}.\\[3pt]
7 & Quasicontinuity and convergence in capacity &
No unrestricted such theorem follows just from weak monotone convergence.
The complex energy proof supplies an exact template, with a dependency
warning explained below. &
Local bounded convergence in capacity and quasicontinuity on smaller
cores of sufficiently small norm boxes:
\texttt{PROVED / NEW\_PROVED}, Corollary~\ref{cor:relative-quasicontinuity}.
They use the stated Borel capacity and ordinary open exceptional sets;
the decreasing family is defined near the closure of the box.
With actual local uniform bounds, convergence in top relative capacity
on all prescribed local norm boxes gives nonmonotone weighted stability:
\texttt{PROVED / DIRECT\_COROLLARY},
Theorem~\ref{thm:capacity-stability}.
Full global/unbounded analogues and outer regularity remain \texttt{OPEN}.\\[3pt]
8 & Pluripolar and negligible sets &
CLD25 16.9.2--4 gives integrability and non-concentration for continuous
approachable potentials against \([1]\)/smooth tests, with hypotheses on
the polar witness. &
Universal ambient non-concentration for arbitrary \(S\): \texttt{FALSE},
Proposition~\ref{prop:polar-relative-mass}. The full zero-capacity/polar/
negligible equivalence is \texttt{OPEN} in the intended domain class.\\[3pt]
9 & Comparison and domination &
The compact-gap comparison for \([1]\) and B mixed backgrounds is a
\texttt{DIRECT\_COROLLARY} of B. Wu26 4.15--4.16 gives separate global
projective results with its stated hypotheses. &
Extension of compact-gap comparison to arbitrary fixed \(S\):
\texttt{PROVED / NEW\_PROVED}, Theorem~\ref{thm:relative-comparison}.
Domination and uniqueness modulo capacity zero are
\texttt{PROVED / DIRECT\_COROLLARY},
Propositions~\ref{prop:capacity-domination}--\ref{prop:capacity-uniqueness}.
Pointwise conclusions require strict-gap capacity detection and are
\texttt{CONDITIONAL}; unrestricted arbitrary-background pointwise
domination remains \texttt{FALSE}.\\[3pt]
10 & Relative extremal functions &
Global envelope theorems in Wu26 have residue-characteristic and global
geometric hypotheses; they do not give the required local theorem over
finite extensions of \(\mathbf Q_p\). &
Local envelope regularity, equilibrium support, and capacity/mass
identities in a suitable domain class: \texttt{OPEN}. Proposed statements
under additional barrier/envelope assumptions are \texttt{CONDITIONAL}.\\[3pt]
11 & Monge--Amp\`ere and Dirichlet problems &
CLD supplies the continuous operator and model comparisons; BGJK21 solves
specified invariant toric/abelian equations. &
The bounded \([1]\) operator is \texttt{PROVED / INTERNAL\_INPUT} (B);
its fixed arbitrary-background extension is also proved here. Appropriate general local
Dirichlet existence and uniqueness are \texttt{OPEN}. Unqualified
existence/uniqueness on every domain is \texttt{FALSE}.\\
\hline
\end{longtable}

\subsection{1. Locally bounded plurisubharmonic functions}

\paragraph{Classical statement and location.}
Dem12, Chapter I, Definition 5.1 and Theorems 5.4--5.8, pp.~39--40,
define psh functions by upper semicontinuity and subharmonicity on complex
lines, prove decreasing-limit stability, decreasing smooth convolution
regularization on smaller opens, convex increasing composition, and the
distributional curvature characterization. In Theorem 5.5 the function is
not identically \(-\infty\) on any component, and the regularization
\(u*\rho_\varepsilon\) is defined on \(\Omega_\varepsilon\), not on an
arbitrary unchanged domain. The upper-envelope theorem is I.5.7; its
proof uses Choquet's lemma I.4.23 and the submean inequality.
\cite[Ch.~I, pp.~39--40]{Dem12}

\paragraph{Non-Archimedean comparison.}
GR25, Definition 4.1 and Proposition 4.2, pp.~21--22, use pullbacks to
analytic curves after every complete valued extension. Their basic
properties include maxima, positive linear combinations, logarithms of
analytic functions, and decreasing nets. Theorem 4.3, pp.~22--23, relates
classical psh PL functions to semipositivity under its good/strict/nontrivial
valuation hypotheses. B, Theorem 2.5, is the stronger accepted input
needed here: actual pointwise locally bounded classical psh functions admit
decreasing continuous semipositive PL approximations on a fixed ordinary
neighborhood over the prescribed field.\cite[\S4]{GR25}

CLD12, Definition 5.5.1, p.~84, uses a distributional definition for
continuous functions. Definition 5.6.3 and Remark 5.6.4, p.~87, introduce
local uniform psh approachability and explicitly identify the then-missing
implication from continuous current positivity to approachability. R,
Theorem 2.9, resolves the stated continuous problem in R's setting.
CLD25, Definition 8.1.1, p.~195, instead defines a sheaf by closure under
decreasing sequences; Definition 8.1.7 and Proposition 8.1.11,
pp.~196--198, specify its continuous class. Theorem names or similar
notation do not identify these classes with classical psh without a bridge.
\cite[\S\S5.5--5.6]{CLD12}\cite[\S8.1]{CLD25}

\paragraph{Required change.}
The present theory fixes B's actual pointwise class \(\mathcal C_b\).
There is no replacement of a function by an almost-everywhere representative:
an arbitrary background can charge sets invisible to ambient volume.

\subsection{2. Positive currents and the recursive product}

\paragraph{Classical statement and proof ingredients.}
Dem12, III.1.14, p.~133, gives measure coefficients and trace control for
positive currents; III.1.23, p.~136, gives local weak compactness under mass
bounds. III.3.2, p.~145, states that for locally bounded psh \(u\) and
arbitrary positive closed \(T\), \(dd^c(uT)\) is positive closed. Its proof
uses decreasing smooth regularization, bounded convergence for the fixed
coefficient measures of \(T\), and continuity of distributional
differentiation. Iteration gives an ordered product before symmetry is
known. BT82, equation (2.1), p.~4, records the recursive definition for the
\([1]\)-background products.\cite[III.1.14, III.1.23, III.3.2]{Dem12}
\cite[\S2, p.~4]{BT82}

\paragraph{Non-Archimedean comparison.}
CLD12, Definition 5.4.1, pp.~81--82, distinguishes the three positivity
cones. Proposition 5.4.6 and Corollary 5.4.7, pp.~83--84, treat top positive
currents as Radon measures and continuous multiplication. Corollaries
5.6.5--5.6.6, p.~88, construct approachable continuous mixed products.
Paragraph 5.6.8, p.~89, explicitly discusses \(D(uT)\) for bounded Borel
\(u\) and explains that positivity needs appropriate decreasing psh
regularization; it is not a completed locally bounded theory.
\cite[\S\S5.4, 5.6]{CLD12}

CLD25, Lemma 8.5.5, p.~207, and \S8.7.1, p.~209, give continuous
recursive products on arbitrary positive closed backgrounds in that
formalism. Its Definition 8.6.2, p.~208, calls a current ``Bedford--Taylor''
if the resulting operation depends only on the curvature of its potential.
That extra property is different from the fixed-potential recursion asked
for here and is not assumed for our arbitrary \(S\).
\cite[\S\S8.5--8.7]{CLD25}

R's coefficient-Radon theorem concerns the original positive cone, whereas
B's \([1]\)-products are strongly positive. These assertions are not
interchangeable in middle degree. The arbitrary-background ordered
construction uses R's exact cone, B's regularization, bounded convergence
against fixed coefficient measures, and continuous-stage positivity. Thus
ordered existence is correctly labeled \texttt{DIRECT\_COROLLARY}; its
later joint continuity and symmetry require further analysis.

\subsection{3. Relative Chern--Levine--Nirenberg estimates}

\paragraph{Classical statement.}
For compact \(L\subset K^\circ\) on a complex manifold and fixed positive
closed \(T\), Dem12, III.3.3, p.~146, proves
\[
 \|dd^cu_1\wedge\cdots\wedge dd^cu_q\wedge T\|_L
 \leq C_{K,L}\prod_{i=1}^q\|u_i\|_{L^\infty(K)}\,\|T\|_K.
\]
The mass seminorm is the sum of absolute variations of local coefficient
measures; its equivalence to a positive trace is explained on p.~145.
The proof iterates a cutoff integration by parts and domination of the
cutoff's Hessian by a local positive trace. BT82, Theorem 2.10, p.~10,
also records mass, bounded-weight, and energy estimates.
\cite[III.3.3]{Dem12}\cite[Theorem 2.10]{BT82}

\paragraph{Non-Archimedean comparison and direct consequence.}
CLD12, Proposition 5.6.2, pp.~86--87, proves its smooth estimate.
CLD25, Proposition 8.4.1, pp.~202--203, gives a smooth relative version;
Proposition 8.4.3, pp.~203--205, gives an additional weighted estimate
under a smooth psh defining-function hypothesis. These are not statements
that every analytic open has a strictly psh exhaustion.
\cite[Proposition 5.6.2]{CLD12}\cite[\S8.4]{CLD25}

For the original arbitrary \(S\), Theorem~\ref{thm:finite-test-cln}
provides a finite family of positive original tests controlling each
requested coefficient test, uniform constants after differentiation and
cutoff, and compatibility with tropical refinements. Its status is
\texttt{DIRECT\_COROLLARY}: finite iteration of R's positive coefficient
control and the ordered recursion supplies the estimate. It is not counted
as a new analytic breakthrough. The result controls arbitrary original
complementary tests, not just positive scalar tests against top-degree
measures.

\subsection{4. Symmetry and weighted joint decreasing continuity}

\paragraph{Classical statement.}
BT82, Theorem 2.1 and Corollary 2.2, p.~4, give joint decreasing
continuity and symmetry for the \([1]\)-products. Theorem 2.4,
pp.~7--9, includes a simultaneously decreasing locally bounded psh weight.
For arbitrary fixed positive closed \(T\), Dem12, III.3.7,
pp.~147--149, proves
\[
 u_{0,j}\,dd^cu_{1,j}\wedge\cdots\wedge dd^cu_{q,j}\wedge T
 \longrightarrow
 u_0\,dd^cu_1\wedge\cdots\wedge dd^cu_q\wedge T,
\]
when all \(u_{i,j}\downarrow u_i\) and the limits are locally bounded.
III.3.10, p.~149, then deduces symmetry.\cite[Theorems 2.1, 2.4]{BT82}
\cite[III.3.7--3.10]{Dem12}

\paragraph{The decisive classical mechanism.}
Dem12 first max-patches every potential to one smooth psh function in a
common boundary collar. CLN yields compactness. Lemma III.3.9, p.~148,
gives a one-sided inequality for a weak limit of decreasing
upper-semicontinuous multipliers times varying positive measures.
Integration by parts moves one smooth regularization at a time; ordered
limits then give the reverse inequality for the total positive trace.
Trace positivity eliminates the defect. This is more than formal
integration by parts, and it is not valid to exchange two discontinuous
limits without the trace argument. The present proof supplies a
test-adapted original-CLD replacement for the required collar and trace.

\paragraph{Baseline and forbidden strengthening.}
B, Theorem 2.7, already proves the whole weighted assertion for background
\([1]\). CLD25, Proposition 8.7.2, pp.~209--210, instead proves local
\emph{uniform} convergence of continuous potentials with weakly varying
backgrounds; Proposition 8.7.4 gives continuous symmetry. The change from
uniform to decreasing convergence is substantive. The branch-well
counterexample in Corollary~\ref{cor:varying-background-false} shows why
one must keep the background fixed in the bounded discontinuous theorem.
\cite[Propositions 8.7.2, 8.7.4]{CLD25}

\subsection{5. Ordinary locality, maximum locality, and maximum principles}

\paragraph{Exact locality assertion.}
For \(u,v,w_2,\ldots,w_k\) locally bounded psh, the desired maximum
identity is
\[
 \mathbf1_{\{u>v\}}\,
 dd^c\max(u,v)\wedge dd^cw_2\wedge\cdots\wedge T
 =\mathbf1_{\{u>v\}}\,
 dd^cu\wedge dd^cw_2\wedge\cdots\wedge T.
\]
It follows directly from the verified weighted theorem BT82, Theorem 2.4,
by the continuous-approximation and weighted-difference argument proved
in Theorem~\ref{thm:maximum-locality-relative}. We do not assign it an
invented numbered location in BT82. Ordinary-open locality follows
immediately from recursion; Borel strict-set locality is a separate
statement.

Wu26, Theorem 4.11, p.~27, supplies local maximum locality for its
top-degree products. Part (1) in the pinned v1 has \(\mathrm{MA}(v)\)
where the right-hand side on \(\{u>v\}\) must be
\(\mathrm{MA}(u)\); part (2) has the correct orientation. This textual
issue is recorded, not used as a premise. B 2.5--2.7 alone already imply
the \([1]\) maximum identity by the proof above, including all
complementary tests.\cite[Theorem 4.11]{Wu26}

The classical local maximum principle for psh functions has a distinct
non-Archimedean counterpart: GR25, Theorem 6.1, pp.~30--31, states local
constancy at an interior local maximum on the connected situation of that
theorem; Theorem 6.4, pp.~33--34, gives a Shilov-boundary maximum
principle on affinoids. Neither assertion identifies ordinary opens with
sets defined by inequalities of discontinuous psh functions.
\cite[\S6]{GR25}

\subsection{6. Capacities and their elementary axioms}

\paragraph{Classical definition and exact axioms.}
BT82, Definition 3.1, p.~11, defines for compact \(F\subset\Omega\)
\[
 \begin{aligned}
 C(F,\Omega)&=\sup_{0<u<1}\int_F(dd^cu)^n,\\
 C(E,\Omega)&=\sup_{\substack{F\subset E\\F\text{ compact}}}C(F,\Omega).
 \end{aligned}
\]
For Borel \(E\), the latter agrees with the supremum of the integrals
over \(E\). Proposition 3.2 gives monotonicity, contravariance in the
domain, countable subadditivity, and continuity from below for Borel
sets. Section 6, p.~22, defines the outer capacity
\(C^*(E)=\inf_{\substack{G\supset E\\G\text{ open}}}C(G)\);
Proposition 6.1 gives its elementary properties.
\cite[Definition 3.1, Proposition 3.2, \S6]{BT82}

Theorem 8.2, p.~29, proves the full Choquet axioms on strongly
pseudoconvex domains, including continuity along decreasing compact sets
and along arbitrary increasing sets. The Choquet theorem stated there
then gives capacitability of analytic sets. These last assertions are
not formal consequences of defining a supremum over measures.
\cite[Theorem 8.2 and the following Choquet theorem]{BT82}

\paragraph{Established non-Archimedean properties and remaining requirements.}
Proposition~\ref{prop:capacity-elementary} proves monotonicity, countable
subadditivity on Borel sets, compact finiteness, continuity from below on
increasing Borel sets, domain contravariance, and mixed-measure domination
for \(\operatorname{Cap}_{S,\beta,k}\). It uses the constructed positive
measures, symmetry, and relative CLN. For this Borel capacity, equality
with its outer envelope remains \texttt{OPEN}. Defining a separate outer
capacity does not establish that equality or transfer the proved estimates
to it. Continuity from above on compact sets and Choquet capacitability
also remain \texttt{OPEN}. Nondegeneracy requires enough admissible psh curvature on
the chosen domain; it fails identically on a connected proper space for
the untwisted capacity, by Proposition~\ref{prop:compact-capacity-zero}.
No theorem in the checked CLD12 \S\S5.4--5.7 or CLD25 Chapter 8 supplies
the full bounded-discontinuous relative capacity package on all ordinary
domains. A useful local norm-box capacity is a more limited statement and
must retain its localization data.

\subsection{7. Quasicontinuity and convergence in capacity}

\paragraph{Classical theorem and dependency.}
BT82, Proposition 3.3, p.~12, gives the energy Cauchy--Schwarz inequality.
Theorem 3.4, pp.~12--13, proves capacity convergence for decreasing
potentials in the common-collar class \(P^*(\Omega,F,\rho,A,B)\),
defined in (2.2), pp.~4--5, by agreement with \(A\rho+B\) off a
compact set. The proof repeatedly applies Cauchy--Schwarz and ends by
using the \emph{already proved} weighted Theorem 2.4; see p.~13.
Consequently it cannot be imported to prove that weighted theorem in a
circle. Theorem 3.5, pp.~13--14, gives quasicontinuity. Its unbounded
case uses the sublevel-capacity estimate Theorem 6.10, p.~26.
\cite[Proposition 3.3; Theorems 3.4, 3.5, 6.10]{BT82}

Theorem 7.2, pp.~26--27, gives quasi-uniform convergence for monotone
psh sequences converging almost everywhere to a psh limit. Theorem 7.4,
p.~27, gives a stronger mixed continuity theorem with all but one
sequence monotone. These require the intervening negligible-set and
capacity theory; they are not corollaries of fixed-background weak
decreasing continuity alone.\cite[Theorems 7.2, 7.4]{BT82}

\paragraph{Transfer achieved in this round and its scope.}
After weighted continuity and symmetry have been established independently,
Theorem~\ref{thm:relative-capacity-estimate} proves the uniform estimate
over the relative capacity tests on a polynomial norm box. Its hypotheses
include a common smooth psh collar and a compact support for the differences
inside a strictly smaller nested box. Corollary~\ref{cor:relative-quasicontinuity}
then proves convergence in capacity on a smaller core for a decreasing
\(\mathcal C_b\) family defined near the closure of the box, and local
bounded quasicontinuity with ordinary open exceptional sets. These are
\texttt{PROVED / NEW\_PROVED} relative results. They concern the defined
Borel capacity; no outer-regularity or nondegeneracy assertion is implicit.
For \(k=n-p\), the trace form \(\beta\) drops out. A general sequence of
Borel multipliers cannot replace the admissible psh differences. Unbounded
truncation estimates, global capacity nondegeneracy criteria, and the full
monotone-a.e. theorem remain \texttt{OPEN}.

\paragraph{Nonmonotone stability now available.}
Theorem~\ref{thm:capacity-stability} is a
\texttt{DIRECT\_COROLLARY} of this established theory. Its hypotheses
are actual bounds uniform in the sequence index on every compact set,
and convergence of every potential and the psh weight in
\(\operatorname{Cap}_S(\,\cdot\,;\Omega)\) on every relatively compact
polynomial norm-box domain \(\Omega\) whose closure lies in an ordinary
algebraic neighborhood in \(U\). The quantified family includes boxes
inside any prescribed finite moment presentation. It proves weighted
joint convergence on all original compact smooth complementary tests
without monotonicity. It does not assert equivalence with convergence
for one larger-domain capacity. On each finite presentation, bounded
\(\rho=\tfrac12\sum t_a^2\) fills the missing differentiated slots
and controls lower-degree coefficient variations by the top relative
capacity. Quasicontinuity then supplies the fixed-multiplier limit.
These arguments use neither outer regularity nor a capacity compactness
theorem, and are not used in proving the earlier monotone or capacity
estimates.

\subsection{8. Pluripolar sets and negligible sets}

\paragraph{Classical statements.}
BT82, Proposition 5.1, p.~18, proves that a negligible set
\(\{h<h^*\}\), where \(h\) is a locally upper bounded psh envelope,
has capacity zero. Theorem 6.9, pp.~25--26, identifies outer-capacity
zero sets with pluripolar sets on strongly pseudoconvex domains;
Corollary 6.6 and Proposition 6.3 supply the equivalent extremal-function
condition. Theorem 7.1, p.~26, concludes that negligible sets are
pluripolar. Theorem 8.3, p.~29, characterizes pluripolarity of analytic
sets by their compact subsets.\cite[\S\S5--8]{BT82}

\paragraph{What CLD already proves.}
CLD25, Theorem 8.12.1 and Corollary 8.12.3, pp.~223--224, give
integrability and polar non-concentration for continuous locally
approachable potentials, assuming suitable smooth psh compact
sublevels and an approachable locally integrable polar witness. In the
analytic setting, Proposition 16.9.1 and Theorem 16.9.2,
p.~466, and Corollaries 16.9.3--16.9.4, p.~467, give the corresponding
conclusions for \(Du_1\cdots Du_m\wedge\alpha\), where \(\alpha\)
is a smooth complementary form. In particular they prove the indicated
measures do not charge lower-dimensional Zariski-closed sets.
\cite[\S8.12, \S16.9]{CLD25}

These theorems are not an assertion for an arbitrary positive closed
background. Proposition~\ref{prop:polar-relative-mass} gives a positive
integration-current background whose nonzero differentiated relative
product is entirely carried by an ambient polar divisor. Thus the
correct relative theory must impose a condition on \(S\), use a
background-relative notion of negligible set, or make a statement on
the support of \(S\). Even for \([1]\), extending the full
zero-capacity/polar/negligible equivalence to the intended bounded theory
requires envelope and unbounded-witness analysis not supplied by the
first-stage calculus.

\subsection{9. Comparison, domination, and uniqueness}

\paragraph{Classical theorem.}
BT82, Theorem 4.1, pp.~14--16, proves for bounded psh \(u,v\) on a
bounded open \(\Omega\subset\mathbf C^n\), with
\(\liminf_{z\to\partial\Omega}(u-v)\geq0\), that
\[
 \int_{\{u<v\}}(dd^cv)^n\leq
 \int_{\{u<v\}}(dd^cu)^n.
\]
Corollaries 4.4--4.5, p.~17, give uniqueness with common boundary data
and the domination principle under the stated zero-mass hypothesis.
The printed proof of uniqueness uses a strictly psh perturbation and
the fact that a nonempty strict inequality set of psh functions has
positive ambient volume. This detection step is additional to comparison.
\cite[Theorem 4.1, Corollaries 4.4--4.5]{BT82}

\paragraph{Non-Archimedean statement now available.}
Theorem~\ref{thm:relative-comparison} proves the relative inequality for
arbitrary fixed original \(S\) when every
\(\overline{\{u<v-\varepsilon\}}\) is compactly contained in the
domain. Its proof uses Borel maximum locality and the compact-support
identity
\[
 T_{S,m}(w^m)-T_{S,m}(u^m)
  =D\!\left((w-u)\sum_{j=0}^{m-1}
       T_{S,m-1}(w^j,u^{m-1-j})\right),
 \qquad w=\max(u,v-\varepsilon).
\]
There is no subtraction of two infinite total masses. However,
\(S=0\), or a background missing directions or regions, prevents
comparison from implying domination. Wu26, Theorems 4.15--4.16,
pp.~29--30, gives separate global projective comparison/domination
results with its polarization and envelope-property hypotheses; those
are not the arbitrary-background local assertion.\cite[Theorems 4.15--4.16]{Wu26}

\paragraph{Comparison already implied by B.}
For \(S=[1]\), the same compact-gap comparison proof uses only B's
symmetry, recursion, coefficient-Radon multiplication, and the maximum
locality already derived from B 2.5--2.7. Thus this case is
\texttt{DIRECT\_COROLLARY}, as is the case where \(S\) is already a B
mixed current, by absorbing its fixed potentials into the differentiated
slots. The \texttt{NEW\_PROVED} assertion is the extension to arbitrary
fixed original positive closed backgrounds.

\paragraph{The measure-theoretic domination conclusion.}
Proposition~\ref{prop:capacity-domination} now proves, under the same
compact-gap condition in the order \((u,v)\), that
\[
 \mu_u(\{u<v\})=0
 \quad\Longrightarrow\quad
 \operatorname{Cap}_S(\{u<v\};V)=0.
\]
If \(\mu_u=\mu_v\) and the compact-gap condition holds in both
orders, Proposition~\ref{prop:capacity-uniqueness} proves
\(\operatorname{Cap}_S(\{u\ne v\};V)=0\). These are
\texttt{DIRECT\_COROLLARY} statements for the defined Borel capacity;
they do not identify a capacity-zero set with the empty set or an ambient
pluripolar set. The additional property that every nonempty strict gap
of such potentials has positive relative capacity gives the pointwise
conclusions in Corollary~\ref{cor:conditional-pointwise-domination},
which is explicitly \texttt{CONDITIONAL}. Establishing that detection
property for a useful domain/background class remains \texttt{OPEN};
positivity of capacity on ordinary open sets alone is insufficient.

\subsection{10. Relative extremal functions and equilibrium measures}

\paragraph{Classical statement.}
For \(E\subset\Omega\), set
\[
 h_E=\sup\{v\in\operatorname{PSH}(\Omega):v\leq0,\ v|_E\leq-1\}.
\]
BT82, Proposition 5.3, p.~19, proves for compact \(E\) in a
strictly pseudoconvex domain that \((dd^ch_E^*)^n\) is supported on
\(E\) and that its total mass equals \(C(E,\Omega)\). The envelope
is psh after upper regularization; the proof also needs increasing
convergence (Proposition 5.2), negligible sets, comparison, and local
homogeneous Dirichlet replacements. The simultaneous induction is
explicit on pp.~19--22. Merely writing the envelope does not prove its
equilibrium identity.\cite[\S5, especially Proposition 5.3]{BT82}

\paragraph{Actual non-Archimedean frontier.}
One must first specify ordinary domains, admissible boundary barriers,
and a class for which upper regularized envelopes stay classical psh.
Then one must prove maximality off the obstacle and identify its
Monge--Amp\`ere measure. Wu26's principal envelope theorem,
Theorem 7.1, p.~51, concerns geometrically unibranch projective
varieties (in particular normal varieties) over discretely or trivially
valued fields of equicharacteristic zero; the residue
characteristic excludes the finite extensions of \(\mathbf Q_p\)
fixed here. B's accepted local regularization remains available but
does not by itself prove the required obstacle-envelope theorem.
\cite[Theorem 7.1]{Wu26}

\subsection{11. Monge--Amp\`ere operators and Dirichlet problems}

\paragraph{Original classical theorem.}
BT76, Theorem D, p.~4, states: on a bounded strictly pseudoconvex
\(\Omega\subset\mathbf C^n\), for continuous boundary data
\(\varphi\in C(\partial\Omega)\) and nonnegative
\(f\in C(\overline\Omega)\), there is a unique
\(u\in\operatorname{PSH}(\Omega)\cap C(\overline\Omega)\)
with the prescribed boundary data and \((dd^cu)^n=f\), using the
paper's volume-form convention. Its existence proof is in \S8,
pp.~39--43, and uses the regularity result Theorem C on balls plus
spherical modifications. The theorem does not assert solvability for
every positive Radon measure.\cite[Theorems C--D and \S8]{BT76}

\paragraph{Non-Archimedean consequence and remaining problem.}
B already supplies the bounded \([1]\)-background operator. The present
first-stage theorem extends its calculus to arbitrary fixed original
backgrounds, and the comparison theorem proves the relative compact-gap
inequality; that is necessary but not
sufficient for existence or uniqueness. One still needs a domain
geometry, barriers, envelope regularity, and a measure class compatible
with the actual capacity. BGJK21 proves invariant toric and abelian
Monge--Amp\`ere existence results, with specified global geometry and
data; it is a model for selected mechanisms, not an arbitrary local
Dirichlet theorem. A proper connected untwisted space already shows
why empty-boundary uniqueness fails for this formulation.
\cite[\S\S6--8]{BGJK21}

\subsection{Compatibility with tropical, \texorpdfstring{$\delta$}{delta}-form, and adelic theories}

BGJK21, Theorem 3.1.4, pp.~13--15, shows that locally bounded tropical psh
functions on its toric partial compactifications are finite and
continuous. Theorem 4.1.2, pp.~20--21, gives arbitrary positive closed
Lagerberg-background products, with local uniform convergence of the
potentials and weak convergence of the background. Theorem 6.1.1,
p.~33, compares the corresponding tropical and non-Archimedean
continuous products. Its bounded class is therefore not a source of
the missing discontinuous-potential argument on general Berkovich spaces.
\cite[Theorems 3.1.4, 4.1.2, 6.1.1]{BGJK21}

GK17, Proposition 4.15, constructs the differential algebra of
\(\delta\)-forms; its Poincar\'e--Lelong theorem is Theorem 7.2, and
Theorem 10.5 identifies formal-metric mixed measures with
Chambert--Loir measures. These statements include nonpositive,
nonclosed, and nondiagonal objects beyond the present BT calculus.
Agreement of two ordinary currents on the original smooth tests does
not imply agreement on every \(\delta\)-test. On common semipositive
PL/model data, ordinary-current agreement is supplied by the
continuous-stage identification, R's model-measure comparison, and
recursion; a broader \(\delta\)-test comparison remains a separate
theorem.\cite[Proposition 4.15, Theorems 7.2, 10.5]{GK17}

R, Theorems 2.9--2.10 and Lemma 3.3, already supply the stated continuous
regularization, positivity, and model/adelic mixed-measure comparisons.
A2, Theorem 1.1, already obtains same-model geometric adelic nefness
from one-place original-current positivity for the specified
integrable arithmetic adelic rational line. Its fixed compactification,
full Cartier boundary, prescribed open-line isomorphisms, and number-field
hypotheses remain part of that result. None is removed by the local BT
extension, and none needs to be proved again in this project.

\paragraph{Dependency discipline.}
The useful classical order is
\[
 \begin{aligned}
 &\text{regularization and coefficient measures}\\
 &\quad\Longrightarrow\ \text{ordered products and CLN}\\
 &\quad\Longrightarrow\ \text{weighted monotone continuity and symmetry}\\
 &\quad\Longrightarrow\ \text{energy/capacity estimates and maximum locality}.
 \end{aligned}
\]
Extremal functions, pluripolar equivalences, and Dirichlet existence
require further inputs. The comparison proof given here follows directly
from maximum locality and compact-support mass conservation and does
not assume those later inputs. This separation prevents capacity from
being used circularly to justify the weighted theorem on which its
classical proof already depends.

% END INCLUDED SOURCE: base_r2/base_r1/sections/comparison


% BEGIN INCLUDED SOURCE: base_r2/base_r1/sections/frontiers
\section{Minimal remaining problems and the next research step}
\label{sec:frontiers}

The fixed-background bounded-potential calculus is proved relative to
the accepted inputs. The following questions are not hidden hypotheses
of that proof. They concern genuinely later parts of the proposed
theory, or a stronger notion of independence than the one requested
for the recursive product.

\subsection{Compact regularity, then envelopes}

The first local analytic target is the weak-closure domination statement
in Section~\ref{sec:outer-regularity-attempt}. Its quantifiers are
important: the normalized potentials live on one fixed domain,
their measures need only converge weakly, and the resulting bound must
use the capacity of that same domain. It suffices for compact outer
regularity and compact continuity from above. It is weaker than a
compactness-and-representation theorem for the potentials, and is
therefore the more economical initial target.

The next envelope problem is to specify a class of ordinary domains
and prove that, for each compact obstacle $E$, the function
\[
 h_E^*=\left(\sup\{v\in\Cb(V):v\leq0,\ v|_E\leq-1\}\right)^*
\]
belongs to the intended classical psh class, with all complete
extensions in its curve test. The upper regularization can change
values on $E$. One must therefore prove, rather than assume, that
the obstacle condition is retained in the required pointwise or
capacity sense. Maximality away from $E$, support of its equilibrium
measure, and a formula relating that measure to capacity are separate
steps. The finite-max counterexample in the preceding section explains
why a formal supremum is insufficient.

Neither compact regularity alone nor local bounded quasicontinuity
has been used to assert Borel outer regularity, negligible-set
theorems, or Choquet capacitability. Each of these remains
\texttt{OPEN} in the full stated relative setting.

\subsection{An exact sufficient condition for polar non-concentration}

\begin{proposition}[A conditional truncation criterion]
\label{prop:polar-truncation-criterion}
Let $q=n-p\geq1$, let $K_0\Subset V$, and let
$\varphi:V\to[-\infty,0]$ be a Borel function such that
$\varphi_N=\max\{\varphi,-N\}\in\Cb(V)$ for every integer $N\geq1$.
Suppose that one number $A<\infty$ satisfies
\begin{equation}
\label{eq:uniform-polar-weight}
 \sup_{N\geq1}\ sup_{\substack{v\in\Cb(V)\\0\leq v\leq1}}
       \int_{K_0}(-\varphi_N)T_{S,q}(v,\ldots,v)\leq A.
\end{equation}
Then, for every $t>0$,
\[
 \operatorname{Cap}_S(K_0\cap\{\varphi<-t\};V)\leq A/t,
 \qquad
 \operatorname{Cap}_S(K_0\cap\{\varphi=-\infty\};V)=0.
\]
Every mixed relative measure of bounded potentials whose oscillations
are bounded on $V$ gives zero mass to the latter set.
\end{proposition}

\begin{proof}
Choose an integer $N\geq t$. On $\{\varphi<-t\}$ one has
$-\varphi_N\geq t$. For each admissible $v$, Markov's inequality
and \eqref{eq:uniform-polar-weight} give
\[
 t\int_{K_0\cap\{\varphi<-t\}}T_{S,q}(v,\ldots,v)\leq A.
\]
Take the supremum over $v$. The pole set is contained in every
sublevel set, so its capacity is at most $A/t$ for every $t>0$.
Let $t\to\infty$. The mixed-measure assertion follows from
\eqref{eq:mixed-capacity-domination}, after shifting and normalizing
the bounded potentials.
\end{proof}

\noindent\textbf{Status.} \texttt{CONDITIONAL}. The implication is
proved, but the uniform weighted bound is not. The bounded CLN estimate
alone gives a bound growing with $N$ and cannot supply
\eqref{eq:uniform-polar-weight}. A plausible next analytic problem is
to derive such a bound from integrability of $\varphi$ against a
specified finite collection of coefficient controls of $S$, together
with an appropriate local barrier. The exact sufficient assumptions
must be proved. In the integration-current counterexample supported
on a polar divisor, $-\varphi_N=N$ on the relative mass; the left
side grows linearly. Thus an unrestricted assertion for all $S$ is
false, not merely unproved.

\subsection{From capacity-zero conclusions to pointwise order}

Section~\ref{sec:capacity-domination} proves domination and uniqueness
modulo zero capacity. The remaining pointwise problem is to determine
for which pairs $(S,V)$ the capacity detects every nonempty strict
psh gap in Corollary~\ref{cor:conditional-pointwise-domination}.
Positivity on nonempty ordinary open sets is not automatically the
same assertion: differences of discontinuous psh functions need not
have ordinary open strict-sign sets. An appropriate relative notion
may be needed if $S$ is carried by a proper analytic subset.

The local stability theorem in Section~\ref{sec:capacity-stability}
now handles nonmonotone sequences that converge in the specified local
capacity. It does not prove precompactness of arbitrary bounded psh
families, nor replace the domain quantifier by a single larger-domain
capacity. Those would require additional theorems.

\subsection{Dirichlet problems and compatibility}

A local Dirichlet theorem must begin by choosing the domain class and
the boundary on which data are prescribed. Ordinary topological,
Shilov, and projective model boundaries serve different purposes.
For a proposed domain class, the next concrete tasks are to construct
barriers, prove the required envelope property, identify measures
controlled by the relative capacity, and establish pointwise
uniqueness under a verified detection condition. General existence
for arbitrary positive Radon data is not assumed.

There is also an elementary obstruction to a domain-free exact
Monge--Amp\`ere existence statement. On a proper boundaryless space,
the positive top product has
\[
 \int_X T_{S,q}(u,\ldots,u)
   =\left\langle D\bigl(uT_{S,q-1}(u,\ldots,u)\bigr),1\right\rangle
   =0
\]
for $q\geq1$, since $1$ is a compactly supported smooth test and
$D1=0$. Thus no nonzero positive measure is such an exact product
there. A curvature class or a suitable local boundary geometry is
substantive additional data.

Ordinary-current compatibility on the already identified continuous
PL/model overlap is retained. Extending identities to all
$\delta$-tests requires its own test-space and multiplication theorem.
The full GK algebra includes objects outside the positive closed
diagonal setting. A2's same-model geometric adelic nefness theorem is
already established and is not included among the open problems.

\subsection{Failed routes retained for continuation}

\begin{enumerate}
\item \texttt{FAILED\_ATTEMPT}: replace fixed-continuous-multiplier
convergence by bounded-measurable-multiplier convergence using only
uniform masses. The branch-well example disproves the resulting
varying-background assertion.
\item \texttt{FAILED\_ATTEMPT}: exchange two discontinuous limits in
the smooth Green formula. The proved replacement is the ordered
one-parameter trace argument; the order of limits is explicit.
\item \texttt{FAILED\_ATTEMPT}: prove weighted monotone continuity
using a capacity convergence theorem whose proof already invokes
weighted monotone continuity. The proof in Section~\ref{sec:monotone}
does not use capacity or the later energy measures.
\item \texttt{FAILED\_ATTEMPT}: use a globally strictly positive
Euclidean trace on every analytic domain. The replacement is a
polynomial norm-box collar and a trace chosen for each finite test
presentation.
\item \texttt{FAILED\_ATTEMPT}: represent every symmetric middle-degree
test by decomposable positive squares. The proof instead uses R's
positive coefficient-matrix cone and respects its distinction from
strong positivity.
\item \texttt{FAILED\_ATTEMPT}: pass compact-support energy identities
through pointwise local approximations or multiply a difference of
psh functions by an arbitrary smooth partition. The latter need not
remain a difference of psh functions. The corrected proof first
regularizes on a prescribed compact box and then uses a second
collar that leaves the original difference unchanged.
\item \texttt{FAILED\_ATTEMPT}: infer outer regularity of a supremum
of measures from the individual Radon properties. The Dirac-family
example shows the exact failure; weak-closure domination is the
remaining analytic statement.
\item \texttt{FAILED\_ATTEMPT}: treat $\{u>v\}$ as an ordinary open
set or infer ambient pluripolar non-concentration for every $S$.
The explicit examples in Section~\ref{sec:counterexamples} rule out
these assertions.
\end{enumerate}

\subsection{Recommended order for the next research session}

First attack fixed-domain weak-closure domination, keeping track of
where an envelope/contact statement is actually needed. In parallel,
test the exact uniform truncation estimate
\eqref{eq:uniform-polar-weight} under an explicit $S$-integrability
hypothesis. Then establish a domain and background criterion that
detects strict psh gaps, so that the proved comparison and capacity-zero
uniqueness become pointwise statements. Only after those steps should
an obstacle-envelope or Dirichlet existence theorem be attempted.
The archived main proof and inputs should be reused; no general
re-audit of R, B, or A2 is part of this plan.

% END INCLUDED SOURCE: base_r2/base_r1/sections/frontiers


% BEGIN INCLUDED SOURCE: base_r2/sections/upper_envelopes
\section{Upper envelopes supplied by the internal projective theory}
\label{sec:r2-upper-envelopes}

\noindent\textbf{Status: \texttt{DIRECT\_COROLLARY}.}
The following assertions use \textup{B}, Theorems C.6 and 4.4, as
accepted inputs. They do not require another proof of envelope
continuity or of the internal contact theorem. The first record did
not use the full force of these inputs in discussing local envelopes.
This section corrects that omission; it is not counted as an
independent new regularization theorem.

\begin{lemma}[A bounded classical competitor lies below the continuous envelope]
\label{lem:r2-envelope-competitor}
Let $F$ be one of the fields allowed in \textup{B}, let $Y/F$ be smooth,
projective and geometrically integral, and let $L$ be ample. If $r$ is
a bounded classical metric on $L$, $f$ is a continuous metric on $L$,
and $r\leq f$, then $r\leq P_L f$ at every point.
\end{lemma}

\begin{proof}
By \textup{B}, Theorem 4.4(2), there are continuous model-semipositive
metrics $r_j\downarrow r$ on the same identified line. Fix
$\epsilon>0$. The open sets
\[
 \{r_j<f+\epsilon\}
\]
increase and cover the compact space $Y^{\an}$. Consequently one of
them is all of $Y^{\an}$. The metric $r_j-\epsilon$ is a continuous
model-semipositive minorant of $f$ and is therefore at most $P_Lf$.
Thus $r\leq P_Lf+\epsilon$. Letting $\epsilon\downarrow0$ proves the
assertion with the actual pointwise representatives. Real translations
preserve continuous model-semipositivity; equivalently they can be
obtained as uniform limits of the permitted rational valuation
translations. No comparison of Monge--Amp\`ere measures is used.
\end{proof}

\begin{proposition}[Envelopes of bounded upper-semicontinuous obstacles]
\label{prop:r2-usc-obstacle}
In the preceding setting, assume $Y^{\an}$ is metrizable, as it is
for the finite $p$-adic fields and the finite unit-Gauss towers used
here. Let $f$ be an upper-semicontinuous metric whose difference from
one continuous metric is bounded. There is a greatest bounded
classical metric $P_L^{\mathrm{b}}f\leq f$. If $f_j\downarrow f$ are
continuous metrics, with a common bounded range relative to a fixed
continuous reference metric, then
\[
 P_L^{\mathrm{b}}f=\lim_j P_L f_j.
\]
In particular, this limit is independent of the continuous majorants.
\end{proposition}

\begin{proof}
Choose a continuous model-semipositive reference metric $a$. A bounded
upper-semicontinuous function on a compact metric space is the
decreasing limit of uniformly bounded continuous functions: for a
compatible metric $d$, the functions
\[
 x\longmapsto\sup_y\bigl((f-a)(y)-j d(x,y)\bigr)
\]
give one such choice. Add $a$ to obtain $f_j$.
By \textup{B}, Theorem C.6, $p_j=P_L f_j$ is continuous
model-semipositive. The sequence decreases. For a sufficiently large
constant $M$, the semipositive metric $a-M$ lies below every $f_j$;
thus $a-M\leq p_j\leq f_j$, and $p=\lim_jp_j$ is bounded and
classical by decreasing closure. Moreover $p\leq f$.

If $r\leq f$ is any bounded classical metric, then $r\leq f_j$ for
every $j$. Lemma~\ref{lem:r2-envelope-competitor} gives $r\leq p_j$;
hence $r\leq p$. Thus $p$ is the greatest permitted minorant, which
also proves independence. The metrizability used here follows, for
these fields, from a countable dense coefficient subfield and
polynomial seminorm inequalities with rational radii on a finite
affine cover. It does not posit metrizability for arbitrary complete
non-Archimedean fields.
\end{proof}

\begin{corollary}[Upper envelopes of bounded projective families]
\label{cor:r2-projective-upper-envelope}
Let $(r_\alpha)_{\alpha\in A}$ be a nonempty family of bounded
classical metrics on $L$, uniformly bounded relative to one continuous
reference metric. Write $h=\sup_\alpha r_\alpha$, with its upper
semicontinuous regularization $h^*$. Then $h^*$ is a bounded classical
metric on $L$.
\end{corollary}

\begin{proof}
The regularization is taken in local frames; changes of frame add a
continuous function and therefore commute with this operation. Put
$p=P_L^{\mathrm b}(h^*)$. Every $r_\alpha$ is a bounded classical
minorant of $h^*$, so $h\leq p$. Since $p$ is upper semicontinuous,
$h^*\leq p$. The defining obstacle inequality gives $p\leq h^*$.
Thus $p=h^*$.
\end{proof}

\begin{theorem}[Local bounded upper-envelope theorem]
\label{thm:r2-local-upper-envelope}
Let $U$ have the fixed analytic hypotheses of the first record. Let
$(u_\alpha)_{\alpha\in A}$ be a nonempty family in $\Cb(U)$ that is
locally uniformly bounded above. Then
\[
 \left(\sup_{\alpha\in A}u_\alpha\right)^*\in\Cb(U).
\]
The regularization is the ordinary upper-semicontinuous regularization
on the actual Berkovich space. The conclusion includes the complete
field-extension and curve tests in the definition of classical psh.
\end{theorem}

\begin{proof}
Work near an arbitrary point. By \textup{B}, Lemmas D.4--D.5 and the
compactification used in Proposition 4.5, work on an ordinary open
of a smooth projective geometrically integral variety over an allowed
finite $p$-adic constant field. Choose the nested compact polynomial
norm boxes $C_1\subset\operatorname{Int}C_{\rm mid}
\subset C_{\rm mid}\subset\operatorname{Int}C_2\subset U$ used
in \textup{B}, Lemma D.2.

Fix one member $u_{\alpha_0}$. It has a finite lower bound on $C_2$.
Choose $M$ with $u_{\alpha_0}\geq-M$ there and replace each
$u_\alpha$ on a neighborhood of $C_2$ by
$v_\alpha=\max\{u_\alpha,-M\}$. These functions have one common
two-sided bound on $C_2$, and
\[
 \sup_\alpha v_\alpha=\sup_\alpha u_\alpha\quad\hbox{on }C_2.
\]
The collar construction in the proof of \textup{B}, Lemma D.2
depends only on these bounds and the chosen boxes. Explicitly, its
pure-FS metric $q$ has a weight $b$ that is zero on $C_1$ and bounded
below by $\delta>0$ on the outer collar. Choose a single translation
$c$ and one integer $m$ with
$0<v_\alpha+c\leq M'$ and $m\delta>M'+1$ for every $\alpha$.
On the same ample line and in the same frame, put
\[
 R_\alpha=\max\{v_\alpha+c,mb\}
       \quad\hbox{on }\operatorname{Int}C_2,
 \qquad R_\alpha=mq\quad\hbox{off }C_{\rm mid}.
\]
The prescriptions agree on the ordinary overlap. Maximum closure and
locality give bounded classical metrics; their differences from $mq$
all lie in $[0,M']$. This argument applies to an arbitrary index set:
the finiteness assumption in the statement of the cited decreasing
family lemma is replaced here by the displayed common bound, and no
simultaneous regularization of this family is being presumed.

Corollary~\ref{cor:r2-projective-upper-envelope} applies to these
metrics. On $\operatorname{Int}C_1$ their weights are
$v_\alpha+c$. Upper regularization is local on ordinary open sets;
subtracting $c$ therefore gives the asserted function near the point.
All metrics in this argument are classical in the original sense,
and restrictions preserve all its field-extension and curve tests.
Classicality is local, so the conclusion holds on $U$.
\end{proof}

\begin{remark}[What this result does not assert]
The theorem does not say that the unregularized supremum is upper
semicontinuous, or that its exceptional set has zero capacity for an
arbitrary background. The latter assertion is false, even when the
family is an increasing sequence of continuous bounded potentials;
see Section~\ref{sec:r2-support-envelopes}.
\end{remark}

% END INCLUDED SOURCE: base_r2/sections/upper_envelopes


% BEGIN INCLUDED SOURCE: base_r2/sections/weak_closure
\section{Quasi-open sets and the weak closure of bounded BT families}
\label{sec:r2-weak-closure}

The fixed field, original CLD test sheaf, positivity cone, and actual
point-value convention are those of the frozen first-round report.
We use its fixed-background calculus, relative coefficient CLN bounds,
local capacity convergence, and local quasicontinuity as inputs. We do
not assume outer regularity on arbitrary Borel sets or a theorem about
upper envelopes. The capacity domain is fixed throughout each argument.

Fix an ordinary norm-box domain $\Omega$, a symmetric positive closed
background $S$ of bidegree $(p,p)$, and put $q=n-p\geq1$. Write
\[
 C(E)=\operatorname{Cap}_S(E;\Omega)
   =\sup_{\substack{v\in\mathcal C_b(\Omega)\\0\leq v\leq1}}
      \int_E T_{S,q}(v,\ldots,v)
\]
for Borel $E\subset\Omega$. The same arguments apply to the localized
$\operatorname{Cap}_{S,\beta,k}$ capacities of the first round.
Define its outer regularization by
\[
 C^*(E)=\inf_{\substack{G\supset E\\G\text{ ordinary open in }\Omega}}
          C(G).
\]
This is a definition; equality with $C(E)$ is not presumed.

\subsection{Outer regularity on the proved class of quasi-open sets}

\begin{definition}
A Borel set $E\subset\Omega$ is $C$-quasi-open if for every
$\epsilon>0$ there is an ordinary open $G\subset\Omega$ such that
$C(G)<\epsilon$ and $E\cup G$ is ordinary open.
\end{definition}

\begin{proposition}[Quasi-open sets are outer regular]
\label{prop:r2-quasiopen-outer}
Every $C$-quasi-open Borel set satisfies $C^*(E)=C(E)$.
Countable unions and finite intersections of $C$-quasi-open sets
are $C$-quasi-open.

Let $V\Subset\Omega$ be an ordinary open core on which real-valued
functions $a,b$ are $C$-quasicontinuous in the sense proved in the
first round. Then
\[
 V\cap\{a>b\},\qquad V\cap\{|a-b|>\delta\}\quad(\delta>0)
\]
are $C$-quasi-open, and hence outer regular for $C$.
In particular this applies to locally bounded classical psh
functions on the permitted local cores.
\end{proposition}

\begin{proof}
Monotonicity gives $C(E)\leq C^*(E)$. If $E\cup G$ is open and
$C(G)<\epsilon$, countable subadditivity gives
\[
 C^*(E)\leq C(E\cup G)\leq C(E)+\epsilon.
\]
Letting $\epsilon\downarrow0$ proves the first assertion. For a
countable family $E_i$, choose exceptional opens $G_i$ with
$C(G_i)<\epsilon2^{-i-1}$ and take $G=\bigcup_iG_i$. Then
$\bigcup_iE_i\cup G$ is open. For a finite intersection, choose
such a common $G$ and use
\[
 \left(\bigcap_{i=1}^N E_i\right)\cup G
     =\bigcap_{i=1}^N(E_i\cup G).
\]

For the final assertion, choose an ordinary open $G$ of capacity
less than $\epsilon$ such that both $a$ and $b$ are continuous on
$V\setminus G$. Their strict inequality set there is relatively
open in $V\setminus G$. Thus there is an ordinary open $O\subset V$
such that
\[
 (V\cap\{a>b\})\setminus G=O\setminus G.
\]
Consequently $(V\cap\{a>b\})\cup G=O\cup G$ is ordinary open.
The proof for $|a-b|>\delta$ is identical. Notice that the argument
does not declare the original strict inequality sets ordinarily open.
\end{proof}

\begin{corollary}[Outer-small exceptional sets]
\label{cor:r2-outer-zero}
If a local strict inequality set between two bounded classical psh
functions has zero capacity, it has ordinary open neighborhoods
with arbitrarily small capacity. Thus the exceptional sets in the
first-round domination and uniqueness statements modulo capacity
zero may also be taken to have outer capacity zero.
\end{corollary}

\begin{proof}
Apply Proposition~\ref{prop:r2-quasiopen-outer} to each strict
inequality set. Equality of two functions fails on the union of
the two strict inequality sets; use countable subadditivity.
\end{proof}

\subsection{What weak limits already preserve}

\begin{lemma}[Open-set and quasi-open-set bounds survive weak limits]
\label{lem:r2-open-limit}
Let $\nu_j$ be positive Radon measures on a relatively compact
ordinary region $W\subset\Omega$, converging locally weakly to
$\nu$. Suppose that for every ordinary open $G\subset W$,
\[
 \nu_j(G)\leq A C(G)
\]
with one finite constant $A$. Then
\[
 \nu(G)\leq A C(G)
\]
for such opens, and
\[
 \nu(E)\leq A C(E)
\]
for every $C$-quasi-open $E\subset W$ whose outer neighborhoods
are taken inside $W$.

The analogous statement holds for total variations of signed
Radon measures converging locally weakly and satisfying the same
uniform open-set bounds.
\end{lemma}

\begin{proof}
For positive measures, the weak-convergence inequality on opens
gives $\nu(G)\leq\liminf_j\nu_j(G)$. For signed measures, use
\[
 |\nu|(G)
  =\sup\left\{\left|\int f\,d\nu\right|:
       f\in C_c(G),\ |f|\leq1\right\}
  \leq\liminf_j|\nu_j|(G).
\]
The local form follows by exhaustion with compactly supported
continuous tests. For a Borel $E\subset W$, Radon outer regularity
of the individual limiting measure gives
\[
 |\nu|(E)\leq A\inf_{G\supset E} C(G)=A C^*(E),
\]
with the neighborhood infimum restricted to $W$.
Proposition~\ref{prop:r2-quasiopen-outer} identifies this with
$A C(E)$ for the stated quasi-open sets.
\end{proof}

\begin{corollary}[Partial weak-closure domination]
\label{cor:r2-quasiopen-domination}
Let $0\leq v_j\leq1$ be classical psh functions on the fixed
domain $\Omega$, with no monotonicity assumption. If
\[
 T_{S,q}(v_j,\ldots,v_j)\longrightarrow\nu
\]
locally weakly, then $\nu(E)\leq C(E)$ on every local
$C$-quasi-open set. In particular this holds on the strict
inequality sets of Proposition~\ref{prop:r2-quasiopen-outer}.
The assertion applies to all members of the local weak closure
of the normalized BT measure family, including limits for which
no potential representation is known.
\end{corollary}

\begin{proof}
Every measure in the original family satisfies the capacity bound
on all Borel sets by the definition of $C$. Apply
Lemma~\ref{lem:r2-open-limit} locally. Positive smooth cutoffs equal
to one on a prescribed core give finite uniformly bounded measures,
so the same argument applies to the weak closure there.
\end{proof}

\subsection{A multiplier theorem for a capacity-controlled weak closure}

\begin{theorem}[Quasicontinuous multipliers through weak closure]
\label{thm:r2-weak-closure-multipliers}
Let $P_j\to P$ be positive coefficient-Radon original CLD currents
of one bidegree, with convergence on all original compact smooth
tests. Fix a compact test $\eta$. Suppose that on a compact
neighborhood $H$ of its support there is one finite constant $A_\eta$
such that
\begin{equation}
\label{eq:r2-open-coefficient-bound}
 |\mu_{P_j,\eta}|(G)\leq A_\eta C(G)
\end{equation}
for every ordinary open $G$ in a fixed slightly larger region.
Let $a$ be bounded and $C$-quasicontinuous on that region. Then
\[
 \langle aP_j,\eta\rangle\longrightarrow\langle aP,\eta\rangle.
\]
No Borel capacity bound for $P$ is required as an extra hypothesis.

More generally, suppose $a_j,a$ are uniformly bounded,
$C$-quasicontinuous there, and $a_j\to a$ locally in $C$-capacity
on every compact subset of the region. Then
\[
 \langle a_jP_j,\eta\rangle\longrightarrow\langle aP,\eta\rangle.
\]
All multipliers are their actual Borel functions.
\end{theorem}

\begin{proof}
The positive-current coefficient theorem implies weak convergence
of the signed coefficient measures against fixed continuous compact
scalar tests. Their total variations are locally uniformly bounded.
Lemma~\ref{lem:r2-open-limit} passes
\eqref{eq:r2-open-coefficient-bound} to $\mu_{P,\eta}$ on opens.

Let $M$ bound $|a|$ on $H$. For $\epsilon>0$, choose an ordinary
open $G$ with $C(G)<\epsilon$ such that $a$ is continuous on
$H\setminus G$. By Tietze extension on the compact Hausdorff
neighborhood $H$, followed by a cutoff equal to one near
$\operatorname{supp}\eta$, there is a continuous compact scalar
$a_\epsilon$, with $|a_\epsilon|\leq M$, equal to $a$ on
$\operatorname{supp}\eta\setminus G$. Therefore
\[
 \left|\int(a-a_\epsilon)\,d\mu_{P_j,\eta}\right|
 +\left|\int(a-a_\epsilon)\,d\mu_{P,\eta}\right|
 \leq4M A_\eta\epsilon.
\]
The integrals with $a_\epsilon$ converge by the fixed-continuous
multiplier theorem. Let $\epsilon\downarrow0$.

For the varying-weight assertion, choose an ordinary open core
$V$ with $\operatorname{supp}\eta\subset V\Subset H$. The Borel
set
\[
 E_{j,\delta}=V\cap\{|a_j-a|>\delta\}
\]
is $C$-quasi-open by Proposition~\ref{prop:r2-quasiopen-outer}.
Individual Radon outer regularity and the open-set bound imply
\[
 |\mu_{P_j,\eta}|(E_{j,\delta})
       \leq A_\eta C(E_{j,\delta}).
\]
The latter capacity tends to zero, because it is bounded by the
capacity of the corresponding set on the larger compact $H$.
If $N$ bounds the coefficient variations on
$\operatorname{supp}\eta$, then
\[
 \left|\int(a_j-a)\,d\mu_{P_j,\eta}\right|
 \leq N\delta+2M A_\eta C(E_{j,\delta}).
\]
First let $j\to\infty$ and then $\delta\downarrow0$. Combine
this estimate with the fixed-weight assertion just proved.
\end{proof}

\begin{corollary}[Fixed additional BT slots on a weak closure]
\label{cor:r2-fixed-slots-closure}
On a fixed local norm-box domain, let
\[
 P_j=T_{S,r}(v_{1,j},\ldots,v_{r,j}),\qquad0\leq v_{i,j}\leq1,
\]
and suppose $P_j\to P$ on original compact smooth tests. The
generating potentials need not converge and need not be monotone.
For every fixed bounded classical psh function $u$,
\[
 uP_j\longrightarrow uP,
 \qquad D(uP_j)\longrightarrow D(uP)
\]
on all original compact tests, whenever the indicated degrees
are allowed. More generally, any fixed tuple of additional bounded
classical psh differentiated slots, and a fixed bounded classical
psh weight, passes through this weak convergence.
\end{corollary}

\begin{proof}
The coefficient-capacity domination proved in the first round
gives \eqref{eq:r2-open-coefficient-bound}, uniformly in $j$, on
each compact original test stage. More precisely, first divide a
given compact test into finitely many moment presentations and
choose the local norm boxes required by that domination statement.
The capacities used for these coefficient estimates have those
fixed local domains. A domain may depend on the test presentation,
but never on $j$; this localization makes no comparison assertion
between capacities on different domains. The restrictions of all
$v_{i,j}$ remain in $[0,1]$ there. The given $u$ is locally
quasicontinuous by the first-round theorem. Apply
Theorem~\ref{thm:r2-weak-closure-multipliers}, and then use
continuity of distributional $D$.

For iteration, after one additional slot the approximating currents
remain bounded mixed BT currents with the same original background
$S$. The additional fixed potential has a finite bound on the
slightly larger local domain. Positive rescaling and addition of
constants put it into a normalized bounded range; the first-round
mixed coefficient-capacity estimate therefore remains uniform.
The limit is the corresponding positive closed ordered product
with background $P$. Repeat one slot at a time and finally apply
the multiplier theorem to the fixed weight.

This is a capacity-controlled varying-background theorem, not an
assertion for arbitrary weakly convergent positive backgrounds.
The branch-well/Dirac counterexample in the first round lacks the
uniform capacity control required here.
\end{proof}

\subsection{The unresolved compact-set step}

The preceding arguments do not establish
\[
 \nu(K)\leq C(K)
 \qquad(K\text{ compact})
\]
for an arbitrary weak limit of normalized BT measures. A compact
set is ordinarily closed; it need not be quasi-open. The theorem
above does not turn its characteristic function into a
quasicontinuous multiplier. The remaining lemma is precisely
weak-closure domination on compact sets from the first-round
reduction, with the original capacity domain unchanged.

Two proposed stronger shortcuts fail. First, individual
quasicontinuity gives a small exceptional open for each potential;
even a common exceptional open for a countable family supplies
continuity of each member, not equicontinuity or compactness of the
normalized family. Second, replacing a sequence by finite or tail
maxima does not order its curvature measures. Moreover the tail
supremum need not be upper semicontinuous. No upper-envelope or
contact theorem is assumed in this section.

The new proved assertions concern quasi-open sets and weak-closure
calculus under capacity bounds. Full compact continuity, full
Borel outer regularity, and general weak-closure domination for an
arbitrary background remain \texttt{OPEN}. For the ambient background
on the later exhaustion domains, the extremal arguments in
Sections~\ref{sec:r2-ambient-extremals} and
\ref{sec:r2-polar-converse-choquet} prove the corresponding compact
and relatively compact Borel assertions. They are additional
arguments, not consequences of weak convergence alone.

% END INCLUDED SOURCE: base_r2/sections/weak_closure


% BEGIN INCLUDED SOURCE: base_r2/sections/capacitary_closure
\section{A bounded capacitary closure and its BT calculus}
\label{sec:r2-capacitary-closure}

This section introduces an auxiliary class for a specific purpose:
an increasing bounded sequence of classical psh functions has an
actual Borel limit which need not be upper semicontinuous, and its
ambient upper regularization can change values on a set detected by
$S$. We construct the calculus on limits in outer relative capacity
without making either of those replacements. The original class
$\Cb$ and all of its pointwise theorems remain as in the first record.
The new class is relative to the fixed current $S$.

Put $q=n-p\geq1$. On an ordinary norm-box domain $\Omega$, write
$C_\Omega(E)=\operatorname{Cap}_{S}(E;\Omega)$ and let
$C_\Omega^*$ denote its outer regularization, defined in
Section~\ref{sec:r2-weak-closure}. All sets occurring as function
level sets below are Borel. The outer capacity is also defined on
arbitrary subsets.

\begin{lemma}[Elementary outer-capacity facts]
\label{lem:r2-outer-facts}
The function $C_\Omega^*$ is monotone and countably subadditive,
and $C_\Omega^*(G)=C_\Omega(G)$ on ordinary open sets. If
$\Omega'\subset\Omega$ is ordinary open and $E\subset\Omega'$, then
\[
 C_\Omega^*(E)\leq C_{\Omega'}^*(E).
\]
For decreasing compact sets $K_j\downarrow K$ contained in
$\Omega$, one has
\[
 C_\Omega^*(K_j)\downarrow C_\Omega^*(K).
\]
The last assertion is about the outer regularization; it does not
establish $C_\Omega^*(K)=C_\Omega(K)$.
\end{lemma}

\begin{proof}
For countable subadditivity choose open covers whose capacities
approximate the respective infima, with errors summing to
$\epsilon$, and take their union. Monotonicity and the assertion
on opens follow directly from the definition. For domain
monotonicity, an open $G\subset\Omega'$ is also open in $\Omega$,
and every normalized potential on $\Omega$ restricts to an
admissible potential on $\Omega'$. Thus
$C_\Omega(G)\leq C_{\Omega'}(G)$, and taking infima proves the
claim. Finally, every open neighborhood $G$ of $K$ contains all
sufficiently large $K_j$, by compactness of $K_1$. Hence
$\lim_j C_\Omega^*(K_j)\leq C_\Omega(G)$. Take the infimum
over $G$ and use monotonicity for the reverse inequality.
\end{proof}

\begin{definition}[Local outer-capacity convergence]
\label{def:r2-outer-convergence}
Real Borel functions $a_j$ converge to $a$ locally in outer
$S$-capacity if, for every norm-box domain $\Omega$ whose compact
closure lies in an allowed ordinary algebraic neighborhood, every
compact $H\Subset\Omega$, and every $\delta>0$,
\[
 C_\Omega^*(H\cap\{|a_j-a|>\delta\})\longrightarrow0.
\]
The Cauchy condition is the corresponding limit with $a_j-a_l$,
as both indices tend to infinity. We quantify over all these
domains, not just one larger domain. Two functions are locally
equal quasi-everywhere if the outer capacity of their disagreement
set on each such compact $H$ is zero.
\end{definition}

\begin{lemma}[Quasicontinuity of capacity limits]
\label{lem:r2-qc-limit}
Let $a_j$ be real, locally bounded, locally $S$-quasicontinuous
functions, and let $a$ be a locally bounded Borel function. If
$a_j\to a$ locally in outer $S$-capacity, then $a$ is locally
$S$-quasicontinuous. The conclusion requires neither compact
outer regularity of $C_\Omega$ nor nondegeneracy of the capacity.
\end{lemma}

\begin{proof}
Fix $H\Subset\Omega$ and a slightly larger compact neighborhood
$H'$ in $\Omega$. On $H'$ choose a subsequence $a_{j_l}$ such that
\[
 C_\Omega^*(H'\cap\{|a_{j_l}-a|>2^{-l}\})
       <\epsilon2^{-l-3}.
\]
Choose ordinary open covers $G_l$ of these sets with
$C_\Omega(G_l)<\epsilon2^{-l-2}$. Local quasicontinuity and a
finite cover of $H'$ give ordinary open sets $F_l$, with the same
capacity bound, such that $a_{j_l}$ is continuous on
$H'\setminus F_l$. Domain monotonicity allows all these estimates
to be measured in the fixed domain $\Omega$.
The union $G=\bigcup_l(G_l\cup F_l)$ has capacity less than
$\epsilon$. On $H'\setminus G$ the functions $a_{j_l}$ are
continuous and converge uniformly to $a$. Their uniform limit is
continuous there. Restriction to $H$ proves the assertion.
\end{proof}

\begin{definition}[Bounded capacitary closure]
\label{def:r2-capacitary-class}
For an ordinary algebraic norm-box domain $W$ and $M<\infty$,
let $\mathcal A^M_S(W)$ consist of the Borel functions $a$ with
$|a|\leq M$ which are limits, locally in outer $S$-capacity on
$W$, of sequences
\[
 a_j\in\Cb(W),\qquad |a_j|\leq M\quad\hbox{on }W.
\]
We retain actual Borel representatives but may also take the
quotient by local quasi-everywhere equality. Let
$\mathcal A_{S,b}(U)$ be the local class: near each point a member
restricts to some $\mathcal A^M_S(W)$ as above. The particular
bound and approximating domain may depend on the neighborhood.
\end{definition}

\begin{proposition}[Existence of capacitary limits]
\label{prop:r2-completion}
On a norm-box domain $W$ over the fixed finite $p$-adic field,
every sequence $a_j\in\Cb(W)$ with $|a_j|\leq M$ which is
locally Cauchy in $S$-capacity has a Borel limit
$a\in\mathcal A^M_S(W)$ locally in outer $S$-capacity. It is
unique up to local quasi-everywhere equality.

The quotient of $\mathcal A^M_S(W)$ is complete for this local
capacity convergence. Addition, nonnegative scalar multiplication,
and finite maxima preserve the local class, with the corresponding
bounds. Every member is locally quasicontinuous.
\end{proposition}

\begin{proof}
The space $W$ has a countable basis of norm-box neighborhoods:
use a countable dense coefficient subfield on an affine algebraic
chart and polynomial seminorm inequalities with rational
thresholds. Choose countably many pairs
\[
 V_l\Subset\Omega_l\Subset W
\]
of ordinary opens, with $\Omega_l$ permitted norm boxes, cofinal
in the following sense: whenever $H\Subset\Omega\subset W$ is
a permitted compact/domain pair, finitely many $V_l$ cover $H$
and their corresponding $\Omega_l$ lie in $\Omega$.
The $V_l$ can be chosen to have compact closure in $\Omega_l$.

Each gap $V_l\cap\{|a_j-a_k|>\delta\}$ is quasi-open for
$C_{\Omega_l}$ by the proved quasicontinuity of $\Cb$ and
Proposition~\ref{prop:r2-quasiopen-outer}. Its outer capacity is
therefore its capacity. The assumed Cauchy condition on
$\overline V_l$ lets us choose a subsequence $a_{j_m}$ such that,
for all $l\leq m$,
\[
 C_{\Omega_l}^*(V_l\cap
       \{|a_{j_{m+1}}-a_{j_m}|>2^{-m}\})<2^{-m}.
\]
For each pair $(l,m)$ with $l\leq m$, choose an ordinary open
cover $G_{l,m}\subset\Omega_l$ of that set with capacity less
than $2^{1-m}$. For fixed $l$ the tails
$G_l^{(N)}=\bigcup_{m\geq N}G_{l,m}$, $N\geq l$, satisfy
\[
 C_{\Omega_l}(G_l^{(N)})\leq 2^{2-N}.
\]
Outside their intersection the subsequence is eventually pointwise
Cauchy on $V_l$, since $\sum_m2^{-m}<\infty$.
Define the Borel function $a=\limsup_m a_{j_m}$ at every point;
it lies in $[-M,M]$. Outside $G_l^{(N)}$,
\[
 |a-a_{j_N}|\leq 2^{1-N}\quad\hbox{on }V_l.
\]
This proves convergence of the subsequence in outer capacity on
each $V_l$. The Cauchy condition and the triangle inequality for
level sets extend convergence to the entire sequence. Cofinality,
finite subadditivity and domain monotonicity then give the stated
convergence on every permitted compact/domain pair. Indeed, a
cover inside each $\Omega_l\subset\Omega$ is an ordinary open
cover in $\Omega$ with no greater capacity.

Two limits have disagreement sets covered by the countable union
of their positive level gaps; the level-set triangle inequality
gives zero outer capacity for each such gap. This proves
uniqueness. Lemma~\ref{lem:r2-qc-limit} gives quasicontinuity.

For completeness, start with a Cauchy sequence in
$\mathcal A^M_S(W)$. For its $j$th member choose a defining
classical approximant within $2^{-j}$ in capacity, at threshold
$2^{-j}$, on the first $j$ compact/domain pairs of a cofinal
enumeration. Such a choice is possible by its definition, and
every chosen approximant has the same bound $M$. The resulting
classical sequence is Cauchy on every pair. The preceding
construction gives its limit, which is also the limit of the
original sequence by the level-set triangle inequality.
Finally, sums, nonnegative multiples and maxima of defining
classical sequences converge to the corresponding pointwise
operations: use the inequalities
\[
 |(a+b)-(c+d)|\leq|a-c|+|b-d|,
 \qquad |\max(a,b)-\max(c,d)|\leq\max(|a-c|,|b-d|).
\]
The classical class has these closure properties. Localizing
proves the assertions for $\mathcal A_{S,b}$.
\end{proof}

\begin{theorem}[BT calculus on the bounded capacitary closure]
\label{thm:r2-completed-calculus}
Let $S$ be the arbitrary fixed original-positive separately closed
background. For $0\leq k\leq q$ and
$a_1,\ldots,a_k\in\mathcal A_{S,b}(U)$, the recursion
\[
 \mathcal T_{S,0}=S,\qquad
 \mathcal T_{S,k}(a_1,\ldots,a_k)
  =D\bigl(a_1\mathcal T_{S,k-1}(a_2,\ldots,a_k)\bigr)
\]
defines positive, symmetric, separately closed coefficient-Radon
currents. It agrees with the original BT product on $\Cb$,
is independent of defining approximants and of local
quasi-everywhere changes of representatives, and is symmetric in
all differentiated slots. It is additive and positively homogeneous
in those slots and respects restriction and gluing.

If $a_{i,j}\in\Cb(U)$, $0\leq i\leq k$, have actual uniform
bounds on compact subsets and converge locally in outer
$S$-capacity to locally bounded Borel functions $a_i$, then
\begin{equation}
\label{eq:r2-completed-weighted-limit}
 a_{0,j}T_{S,k}(a_{1,j},\ldots,a_{k,j})
 \longrightarrow
 a_0\mathcal T_{S,k}(a_1,\ldots,a_k)
\end{equation}
on all original compact smooth complementary tests. The same
stability holds for sequences in the capacitary closure if, locally,
their defining approximants admit one common bound.
\end{theorem}

\begin{proof}
Work on a compact test stage, localizing into finitely many smaller
norm boxes $\Omega$ where the test's unit moment coordinates are
defined and bounded. R1 finite coefficient domination gives, for
every bounded classical approximating product $P_j$ and each fixed
test $\eta$ on this stage,
\begin{equation}
\label{eq:r2-approx-coefficient-cap}
 |\mu_{P_j,\eta}|(E)\leq C_\eta C_\Omega(E)
       \quad(E\text{ Borel}),
\end{equation}
with a constant depending only on the common bounds and the test
stage. In degree zero this holds for $P_j=S$ itself, by filling
the complementary degree with the bounded quadratic moment
potential. If positive currents $P_j$ converge to $P$, their
coefficient measures converge weakly against compact continuous
scalars by \textup{R}. Lower semicontinuity of total variation on
opens and individual Radon outer regularity imply
\begin{equation}
\label{eq:r2-limit-coefficient-outer}
 |\mu_{P,\eta}|(E)\leq C_\eta C_\Omega^*(E)
       \quad(E\text{ Borel}).
\end{equation}
This is precisely the weak-closure argument of
Lemma~\ref{lem:r2-open-limit}; it does not replace $C_\Omega^*$
by $C_\Omega$ on compact sets.

We induct on $k$. The case $k=0$ follows from Borel multiplication
of $S$, the coefficient bound, and the multiplier theorem
\ref{thm:r2-weak-closure-multipliers}. Limits $a_i$ are
quasicontinuous by Lemma~\ref{lem:r2-qc-limit}. Suppose the
assertions through degree $k-1$ are known. Weighted convergence
there, followed by distributional $D$, gives convergence of the
degree-$k$ products to the displayed recursive expression. A weak
limit of positive symmetric closed currents retains these
properties. The coefficient-Radon theorem of \textup{R} applies
to it, and \eqref{eq:r2-limit-coefficient-outer} supplies its
outer-capacity control. Apply the multiplier theorem once more
to get \eqref{eq:r2-completed-weighted-limit} in degree $k$.
Its varying-weight proof also follows directly from
\[
 \left|\int(a_{0,j}-a_0)\,d\mu_{P_j,\eta}\right|
 \leq\delta A+2M C_\eta
      C_\Omega^*(H\cap\{|a_{0,j}-a_0|>\delta\}),
\]
and continuous extension of $a_0$ off an ordinary exceptional
open of arbitrarily small capacity. Here $A$ bounds coefficient
variation on the compact test support and $M$ bounds the weights.

The recursion uses the actual Borel values of $a_i$ and therefore
does not contain any approximation choice. On overlaps the same
recursion gives the same current, proving gluing. Formula
\eqref{eq:r2-limit-coefficient-outer} shows that changing a slot
on a locally outer-capacity-zero set leaves every relevant Borel
product unchanged; induction proves representative independence.
Symmetry, addition and positive homogeneity follow by taking
limits of the corresponding identities for classical defining
sequences, with simultaneous finite-slot approximation.

For sequences already in the closure, their coefficient measures
obey \eqref{eq:r2-limit-coefficient-outer}, with uniform constants
when the defining bounds are uniform. The same multiplier
argument and induction apply. All localizations are over finitely
many actual original test stages, so the conclusion holds for
every original compact complementary test, including tests not
coming from one initially chosen moment map.
\end{proof}

\paragraph{Status and role.}
The analytic extension theorem is recorded as \texttt{NEW\_PROVED},
relative to the R1 fixed-background theory and the R2 weak-closure
multiplier lemma. The outer-capacity and completion lemmas are
\texttt{DIRECT\_COROLLARY} measure-capacity arguments with proofs
included. The construction is useful even when $C_\Omega$ is
degenerate. It does not assert that every member of
$\mathcal A_{S,b}$ has a classical psh representative on the
ambient space, or that outer capacity equals the original Borel
capacity. In particular it does not solve the unresolved compact
weak-closure domination problem by definition. For the ambient
background alone, Corollary~\ref{cor:r2-ambient-classical-representative}
later identifies each equivalence class with a unique classical
representative; the arbitrary-background distinction remains.

% END INCLUDED SOURCE: base_r2/sections/capacitary_closure


% BEGIN INCLUDED SOURCE: base_r2/sections/envelope_capacity_attempt
\section{Increasing sequences and relative Aubin--Mabuchi energy}
\label{sec:r2-increasing-energy}

This section uses the frozen first-round calculus and its common-collar
Green identity. Let $S$ be any allowed positive closed original current
of bidegree $(p,p)$, put $q=n-p\geq1$, and fix one ordinary polynomial
norm-box domain $\Omega$. The capacity is always
$C(E)=\operatorname{Cap}_S(E;\Omega)$, with this same domain.

Choose a bounded smooth classical psh function $b$ on $\Omega$ and a compact
$L\Subset\Omega$, contained in a strictly smaller norm box. Let
$\mathcal P_b$ denote the locally bounded classical psh functions $u$
which equal $b$ on $\Omega\setminus L$. All families below have a common
bound on a fixed compact neighborhood of $L$, hence on all of $\Omega$
because of the common collar. Such families arise from
the common-collar construction; its maximum operation also preserves
increasing sequences.

\subsection{A compactly supported relative energy}

\begin{definition}
For $u\in\mathcal P_b$, define
\[
 \mathcal E_b(u)=\frac1{q+1}\sum_{i=0}^q
       \int_\Omega (u-b)(Du)^i\wedge(Db)^{q-i}\wedge S.
\]
Every integral is a Borel integral supported in $L$. Products use the
proved fixed-background calculus; in particular this definition does
not require a smooth background or a strict psh exhaustion.
\end{definition}

\begin{theorem}[Relative energy difference formula]
\label{thm:r2-energy-difference}
For $u,v\in\mathcal P_b$,
\begin{equation}
\label{eq:r2-energy-difference}
 \mathcal E_b(v)-\mathcal E_b(u)
   =\frac1{q+1}\sum_{i=0}^q
      \int_\Omega(v-u)(Dv)^i\wedge(Du)^{q-i}\wedge S.
\end{equation}
Consequently $u\leq v$ implies $\mathcal E_b(u)\leq\mathcal E_b(v)$,
and in that case
\begin{equation}
\label{eq:r2-lower-energy-control}
 \int_\Omega(v-u)(Du)^q\wedge S
     \leq(q+1)\bigl(\mathcal E_b(v)-\mathcal E_b(u)\bigr).
\end{equation}
On a uniformly bounded subset of $\mathcal P_b$, the energies are
uniformly bounded in absolute value.
\end{theorem}

\begin{proof}
Put $g=v-u$ and $u_t=(1-t)u+tv$, for $0\leq t\leq1$. All these
functions have the same collar. Multilinearity makes
$F(t)=\mathcal E_b(u_t)$ a finite polynomial in $t$; differentiating it
is therefore an algebraic operation on finitely many finite integrals.
Write $\alpha_t=Du_t$ and $\beta=Db$. The derivative is
\begin{align*}
 (q+1)F'(t)
 &=\sum_{i=0}^q\int g\,\alpha_t^i\beta^{q-i}S\\
 &\quad+\sum_{i=1}^q i\int(u_t-b)Dg\,
                                  \alpha_t^{i-1}\beta^{q-i}S.
\end{align*}
In each term of the second sum, the background
$\alpha_t^{i-1}\beta^{q-i}S$ is positive closed of bidegree
$(n-1,n-1)$, and both differences have compact support in the common
norm box. The first-round Green identity gives
\[
 \int(u_t-b)Dg\,\alpha_t^{i-1}\beta^{q-i}S
  =\int g(\alpha_t-\beta)\alpha_t^{i-1}\beta^{q-i}S.
\]
The resulting sums telescope, leaving
\[
 F'(t)=\int g\,\alpha_t^qS.
\]
Expand $\alpha_t^q=((1-t)Du+tDv)^q$ and integrate from zero to one.
The identity
\[
 \binom qi\int_0^1t^i(1-t)^{q-i}\,dt=\frac1{q+1}
\]
proves \eqref{eq:r2-energy-difference}.

If $g\geq0$, every summand is nonnegative. Retaining the term $i=0$
gives \eqref{eq:r2-lower-energy-control}. Finally, the relative CLN
estimate bounds the masses of all mixed products on $L$ uniformly
when the potentials have one common bound. Multiplying by the bounded
function $u-b$ gives the asserted energy bound.
\end{proof}

\begin{lemma}[The pairwise capacity bound]
\label{lem:r2-pair-capacity}
On a uniformly bounded subset of $\mathcal P_b$, there is a constant
$A$ such that $u\leq v$ implies, for every $\delta>0$,
\begin{equation}
\label{eq:r2-pair-capacity}
 C(\{v-u>\delta\})
 \leq\frac A\delta
   \left((q+1)\bigl(\mathcal E_b(v)-\mathcal E_b(u)\bigr)
   \right)^{2^{-q}}.
\end{equation}
The same estimate holds with $C^*$ on the left.
\end{lemma}

\begin{proof}
The proof of the first-round quantitative capacity estimate applies
to any single pair $u\leq v$ with compactly supported difference.
Indeed its iteration uses only $g=v-u\geq0$, Green's identity,
Cauchy--Schwarz, positivity of $D(vQ)$, and uniform CLN bounds.
Before its final limit step it gives
\[
 C(\{v-u>\delta\})
 \leq\frac A\delta\left(\int(v-u)(Du)^qS\right)^{2^{-q}}.
\]
Apply \eqref{eq:r2-lower-energy-control}. The strict inequality set is
quasi-open by the preceding section, so its outer capacity equals its
Borel capacity. This proves the last assertion without assuming compact
outer regularity.
\end{proof}

\subsection{Increasing limits converge to their actual pointwise limit}

\begin{theorem}[Increasing sequences are Cauchy in outer capacity]
\label{thm:r2-increasing-cauchy}
Let $u_j\in\mathcal P_b$ be increasing and uniformly bounded on the
fixed compact neighborhood. Put
\[
 u_\infty(x)=\sup_j u_j(x),\qquad
 e_\infty=\sup_j\mathcal E_b(u_j).
\]
Then $e_\infty$ is finite. For every $\delta>0$,
\begin{equation}
\label{eq:r2-increasing-actual-rate}
 C^*(\{u_\infty-u_j>\delta\})
 \leq\frac A\delta
     \left((q+1)(e_\infty-\mathcal E_b(u_j))\right)^{2^{-q}}.
\end{equation}
The actual pointwise limit $u_\infty$ is locally
$C$-quasicontinuous on each ordinary core. In particular
$u_j\to u_\infty$ in outer capacity on the common difference support,
with the original domain $\Omega$ unchanged.

The theorem does not assert that $u_\infty$ is upper semicontinuous,
classically psh, or equal outside a capacity-zero set to its upper
semicontinuous regularization.
\end{theorem}

\begin{proof}
Theorem~\ref{thm:r2-energy-difference} makes the energies increasing
and uniformly bounded. For $l\geq j$, apply
Lemma~\ref{lem:r2-pair-capacity} to $(u_j,u_l)$. Its right side is
bounded by the right side of \eqref{eq:r2-increasing-actual-rate},
independently of $l$. The strict gap sets increase with $l$ and their
union is $\{u_\infty-u_j>\delta\}$. Continuity from below of the
Borel capacity, already proved in the first round, gives the stated
bound with $C$ on the left. Each strict gap set is quasi-open, and
countable unions of quasi-open sets are quasi-open. Therefore the
union has outer capacity equal to its Borel capacity. This proves
\eqref{eq:r2-increasing-actual-rate} directly, before any regularity
of $u_\infty$ is known.

Here is a direct quasicontinuity argument, with no assumption about
upper envelopes. Fix an ordinary core $V$ and $\epsilon>0$. Choose
a subsequence $u_{j_l}$ so late that the pairwise estimate guarantees
\[
 C^*(\{u_{j_{l+1}}-u_{j_l}>2^{-l}\})
       <\epsilon2^{-l-3}.
\]
This choice can be made successively from the bound involving
$e_\infty-\mathcal E_b(u_{j_l})$. Choose ordinary open covers $G_l$
of these gap sets with $C(G_l)<\epsilon2^{-l-2}$. Also use
first-round quasicontinuity to choose ordinary opens $H_l$ with
$C(H_l)<\epsilon2^{-l-2}$ such that $u_{j_l}$ is continuous on
$V\setminus H_l$. The union
\[
 G=\bigcup_l(G_l\cup H_l)
\]
has capacity less than $\epsilon$. On $V\setminus G$, the continuous
functions $u_{j_l}$ have consecutive increments bounded by $2^{-l}$.
They therefore converge uniformly there. Their pointwise limit is
the specified function $u_\infty$, so this function is continuous
on $V\setminus G$.

The convergence assertion now follows from
\eqref{eq:r2-increasing-actual-rate}.
\end{proof}

\begin{corollary}[Local increasing convergence without a prescribed collar]
\label{cor:r2-local-increasing}
Let $u_j\in\Cb(U)$ increase pointwise and be uniformly bounded on
every compact subset of $U$. Set $u_\infty=\sup_j u_j$, with its
actual Borel values. For every permitted norm-box domain
$\Omega\Subset U$, every compact $H\Subset\Omega$, and every
$\delta>0$,
\[
 \operatorname{Cap}_S^*(H\cap\{u_\infty-u_j>\delta\};\Omega)
       \longrightarrow0.
\]
Thus the convergence is local outer-capacity convergence in the
all-domain sense of Definition~\ref{def:r2-outer-convergence}.
The limit is locally quasicontinuous and belongs to
$\mathcal A_{S,b}(U)$.
\end{corollary}

\begin{proof}
Fix $(H,\Omega)$. Cover $H$ by finitely many ordinary smaller cores
$V_a$ contained in norm-box domains $\Omega_a\Subset\Omega$.
On each such domain use the common-collar construction on the whole
sequence. Its construction requires a common two-sided bound on a
slightly larger box; it does not use the direction of monotonicity.
After one fixed shift $c_a$, it gives a bounded smooth classical psh
barrier $b_a$ and
\[
 v_{j,a}=\max\{u_j+c_a,b_a\},
\]
equal to $u_j+c_a$ on $V_a$, and equal to $b_a$ outside a fixed
compact subset of $\Omega_a$. The maximum preserves increasing
order. Its pointwise supremum equals $u_\infty+c_a$ on $V_a$.
Theorem~\ref{thm:r2-increasing-cauchy} therefore gives
\[
 \operatorname{Cap}_S^*\bigl(
 V_a\cap\{u_\infty-u_j>\delta\};\Omega_a\bigr)
       \longrightarrow0.
\]
Domain monotonicity of outer capacity, with
$\Omega_a\subset\Omega$, and finite subadditivity imply the
asserted estimate on $H$. This repeats the collar construction
inside each requested domain; convergence for one larger capacity
domain has not been substituted for convergence on its smaller
subdomains.

The actual limit is a locally bounded Borel function. Applying the
quasicontinuity lemma for outer-capacity limits gives local
quasicontinuity. On each relatively compact neighborhood the
original sequence itself is uniformly bounded and gives the
defining approximation for $\mathcal A_{S,b}$.
\end{proof}

\begin{corollary}[Weighted increasing limits with actual representatives]
\label{cor:r2-increasing-weighted}
Let $0\leq k\leq q$, and suppose that
$u_{i,j}\in\Cb(U)$, $0\leq i\leq k$, increase with $j$ and
have a common bound on every compact subset for each fixed slot.
Write $u_i=\sup_j u_{i,j}$ with actual Borel values. Then
\[
 u_{0,j}T_{S,k}(u_{1,j},\ldots,u_{k,j})
   \longrightarrow
 u_0\mathcal T_{S,k}(u_1,\ldots,u_k)
\]
on every original compact smooth complementary test form, where
$\mathcal T$ is the calculus on the bounded capacitary closure.
No ambient upper regularization of any $u_i$ is made.
\end{corollary}

\begin{proof}
Apply Corollary~\ref{cor:r2-local-increasing} in each slot, followed
by Theorem~\ref{thm:r2-completed-calculus}. These are successive
uses: the energy estimate and the local increasing convergence were
proved using only the frozen classical-potential calculus, before
applying its extension to the capacitary closure.
\end{proof}

\subsection{The precise upper-envelope question}

Suppose an upper-envelope theorem supplies a classical psh function
\[
 U=(\sup_j u_j)^*
\]
on the same local domain. Because all $u_j$ equal the same smooth
$b$ outside $L$, so does $U$ there. The following formulation separates
the available energy argument from the remaining analytic step.

\begin{proposition}[The energy no-jump criterion]
\label{prop:r2-envelope-energy-criterion}
Under these hypotheses, the following are equivalent:
\begin{enumerate}
\item $\mathcal E_b(U)=\sup_j\mathcal E_b(u_j)$;
\item $u_j\to U$ locally in the relative capacity, on a core
      containing the common compact difference support;
\item $U=u_\infty$ outside a set of relative capacity zero on that core.
\end{enumerate}
The zero set in (3) has outer capacity zero, since both functions
are quasicontinuous on the core.
\end{proposition}

\begin{proof}
For (1)$\Rightarrow$(2), apply
Lemma~\ref{lem:r2-pair-capacity} to $(u_j,U)$; the energy difference
tends to zero. For (2)$\Rightarrow$(1), use
\eqref{eq:r2-energy-difference} with $(u_j,U)$. Every measure on its
right is a mixed top-degree measure of uniformly bounded psh
potentials on the \emph{same} $\Omega$. After translating and
rescaling the potentials, the first-round mixed-capacity bound gives
one uniform bound by a constant times $C(E)$ on every Borel set $E$.
Their masses on $L$ are uniformly bounded by CLN. Split each integral
of $U-u_j$ into the sets where this gap is at most $\delta$ and
greater than $\delta$. The first part is bounded by a uniform mass
times $\delta$, and the second by a uniform bound times
$C(\{U-u_j>\delta\}\cap L)$. First let $j\to\infty$, then
$\delta\downarrow0$. This proves energy convergence without changing
the capacity domain or invoking a stronger convergence on smaller
domains.

Theorem~\ref{thm:r2-increasing-cauchy} gives capacity convergence to
the actual limit $u_\infty$. If (2) holds, subadditivity applied to
the two error-level sets gives $C(\{|U-u_\infty|>\delta\})=0$
for every $\delta>0$ on a slightly smaller core. Taking their
countable union proves (3). Conversely, (3), subadditivity, and the
already proved convergence to $u_\infty$ give (2). Quasicontinuity
of the two functions makes the strict error-level sets quasi-open,
so their outer capacities also vanish.
\end{proof}

For arbitrary backgrounds $S$, the upper-envelope identification
cannot be assumed: an ambient envelope can raise the limit on a
subvariety which carries positive relative capacity. For $S=[1]$,
the criterion isolates a step which is proved later by the ordinary
divisorial barrier in Theorem~\ref{thm:r2-increasing-no-jump}.
Continuity and orthogonality for a continuous obstacle in
\textup{B}, Theorem C.6, and the upper contact statement in
\textup{B}, Theorem 4.4, do not by their statements identify that
energy. The later barrier proof supplies the additional argument;
it was not inserted into the criterion's proof.

% END INCLUDED SOURCE: base_r2/sections/envelope_capacity_attempt


% BEGIN INCLUDED SOURCE: base_r2/sections/relative_capacitary_envelopes
\section{Relative capacitary upper envelopes}
\label{sec:r2-capacitary-envelopes}

The upper envelope constructed here belongs to the bounded
capacitary closure. Its order is the order modulo relative
outer-capacity-zero sets. The construction does not replace it by
an ambient upper semicontinuous regularization. In particular it
continues to apply when the background is carried by an ambient
pluripolar subvariety.

\subsection{One common collar and one capacity domain}

Use the notation of Section~\ref{sec:r2-increasing-energy}:
$q=n-p\geq1$, a fixed norm box $\Omega$, a bounded smooth
classical psh collar $b$, and the class $\mathcal P_b$ with one
compact difference support $L\Subset\Omega$. Write
$C_\Omega^*=\operatorname{Cap}_S^*(\,\cdot\,;\Omega)$.
For Borel functions define $a\preceq_\Omega d$ to mean
$C_\Omega^*(\{a>d\})=0$.

\begin{theorem}[A capacitary envelope from a maximizing sequence]
\label{thm:r2-one-domain-envelope}
Let $\mathcal F\subset\mathcal P_b$ be nonempty and uniformly
bounded. Let $\mathcal F^\vee$ be its closure under nonempty finite
maxima, and set
\[
 e_* = \sup_{v\in\mathcal F^\vee}\mathcal E_b(v).
\]
There is an increasing sequence $u_j\in\mathcal F^\vee$ such that
$\mathcal E_b(u_j)\uparrow e_*$. Its actual pointwise limit
\[
 H=\sup_j u_j
\]
belongs to $\mathcal A_{S,b}(\Omega)$ and has the following properties:
\begin{enumerate}
\item For each $v\in\mathcal F$, one has $v\preceq_\Omega H$.
\item If a Borel function $h$ satisfies $v\preceq_\Omega h$ for
      every $v\in\mathcal F$, then $H\preceq_\Omega h$.
\item The equivalence class of $H$ modulo $C_\Omega^*$-zero sets
      is independent of the maximizing sequence.
\end{enumerate}
Moreover $H=b$ outside $L$, and the extended BT calculus gives
$\mathcal E_b(H)=e_*$, with the same finite-sum definition of the
energy.
\end{theorem}

\begin{proof}
The energy is uniformly bounded on $\mathcal F^\vee$ by the
common bounds and CLN. Choose $v_j\in\mathcal F^\vee$ with
$\mathcal E_b(v_j)>e_*-2^{-j}$, and put
$u_j=\max(v_1,\ldots,v_j)$. Monotonicity of the energy gives
\[
 e_*-2^{-j}<\mathcal E_b(v_j)\leq\mathcal E_b(u_j)\leq e_*.
\]
The preceding increasing convergence theorem, including its
all-domain local version, applies to $H$ and gives membership in
the capacitary closure.

Fix $v\in\mathcal F$ and put $w_j=\max(u_j,v)$. These functions
are in $\mathcal F^\vee$, and
\[
 0\leq\mathcal E_b(w_j)-\mathcal E_b(u_j)
       \leq e_*-\mathcal E_b(u_j)\longrightarrow0.
\]
For every $\delta>0$, the pairwise energy estimate implies
\[
 C_\Omega^*(\{v>H+\delta\})
 \leq C_\Omega^*(\{w_j-u_j>\delta\})\longrightarrow0.
\]
Taking the countable union over $\delta=1/m$ proves (1).

Every $u_j$ is a finite maximum of members of $\mathcal F$.
If $h$ is as in (2), then $u_j\leq h$ off a finite union of
outer-capacity-zero sets. Taking the countable union of these
exceptional sets over $j$ gives $H\preceq_\Omega h$.
Applying (2) to two constructed upper bounds proves (3).

Finally, $u_j-b$ and $H-b$ vanish outside $L$. Choose a smooth
compact cutoff equal to one near $L$. Apply weighted increasing
BT convergence separately to the two weights $u_j$ and $b$ in
each summand of $\mathcal E_b(u_j)$. The cutoff makes this a
finite collection of original compact smooth tests. Their limits
are precisely the summands of $\mathcal E_b(H)$, proving the
last identity.
\end{proof}

The exceptional set in (1) is allowed to depend on $v$.
For an uncountable family, (1) does not assert that the pointwise
supremum of the entire family is bounded by $H$ outside one
common exceptional set.

\subsection{A countable energy construction for the local order}

The preceding theorem uses one fixed capacity domain. The
following construction enforces the estimates in every smaller
domain, without assuming equivalence of capacities on different
domains.

Let $W$ be an allowed ordinary norm-box domain over the fixed
finite $p$-adic field. Define $a\preceq_{S,\mathrm{loc}}d$ when,
for every permitted $H_0\Subset\Omega\subset W$,
\[
 \operatorname{Cap}_S^*(H_0\cap\{a>d\};\Omega)=0.
\]

\begin{theorem}[Local relative capacitary envelope]
\label{thm:r2-local-capacitary-envelope}
Let $\mathcal F\subset\Cb(W)$ be nonempty and uniformly bounded
above on each compact subset of $W$. There is an increasing sequence
$u_j\in\mathcal F^\vee$ whose actual Borel supremum $H$ belongs
to $\mathcal A_{S,b}(W)$ and is the least upper bound of
$\mathcal F$ for $\preceq_{S,\mathrm{loc}}$:
\[
 v\preceq_{S,\mathrm{loc}}H\quad(v\in\mathcal F),
 \qquad
 \bigl[\ v\preceq_{S,\mathrm{loc}}h\ \text{for all }v\in\mathcal F\ \bigr]
 \Longrightarrow H\preceq_{S,\mathrm{loc}}h.
\]
Its local quasi-everywhere class is independent of all choices.
\end{theorem}

\begin{proof}
Fix $v_0\in\mathcal F$ and first replace the family by
$\{\max(v,v_0):v\in\mathcal F\}$. This family is uniformly
bounded on every compact subset, since $v_0$ supplies a common
lower bound. It has exactly the same local quasi-everywhere upper
bounds as the original family, and every finite maximum of its
members is a finite maximum of original members. We may therefore
assume the two-sided bounds in the rest of the construction.

Choose countably many pairs $V_a\Subset\Omega_a\Subset W$
cofinal for compact/domain pairs, as in
Proposition~\ref{prop:r2-completion}. Thus every
$H_0\Subset\Omega\subset W$ is covered by finitely many $V_a$
whose corresponding $\Omega_a$ lie in $\Omega$.
The common bound of $\mathcal F^\vee$ on a slightly larger box
lets us choose, for each $a$, one fixed collar operation
\[
 P_a(v)=\max\{v+c_a,b_a\}\quad\text{on }\Omega_a.
\]
It equals $v+c_a$ on $V_a$ and the bounded smooth psh function
$b_a$ outside one fixed compact subset of $\Omega_a$.
Its construction is uniform over $\mathcal F^\vee$, and
\[
 P_a(\max(v,w))=\max(P_a(v),P_a(w)).
\]
Write $E_a(v)=\mathcal E_{b_a}(P_a(v))$. The relative energy
theorem gives a finite $B_a\geq1$ with $|E_a(v)|\leq B_a$ on
$\mathcal F^\vee$. Define
\[
 \Phi(v)=\sum_{a=1}^{\infty}2^{-a}
             \frac{E_a(v)+B_a}{2B_a+1}.
\]
The series converges uniformly, takes values in $[0,1]$, and is
monotone for pointwise order.

Put $e=\sup_{v\in\mathcal F^\vee}\Phi(v)$. Choose
$v_j\in\mathcal F^\vee$ with $\Phi(v_j)>e-2^{-j}$ and set
$u_j=\max(v_1,\ldots,v_j)$. Then $\Phi(u_j)\uparrow e$.
The functions $u_j$ are locally uniformly bounded, so their
actual limit $H$ belongs to $\mathcal A_{S,b}(W)$ by
Corollary~\ref{cor:r2-local-increasing}.

Fix $v\in\mathcal F$ and set $w_j=\max(u_j,v)$. For each $a$
the difference $E_a(w_j)-E_a(u_j)$ is nonnegative. Consequently
\[
 0\leq E_a(w_j)-E_a(u_j)
 \leq 2^a(2B_a+1)\bigl(\Phi(w_j)-\Phi(u_j)\bigr)
 \leq 2^a(2B_a+1)(e-\Phi(u_j))\longrightarrow0.
\]
For $\delta>0$, the pairwise energy estimate in $\Omega_a$ gives
\[
 \operatorname{Cap}_S^*\bigl(V_a\cap\{v>H+\delta\};\Omega_a\bigr)
 \leq\operatorname{Cap}_S^*\bigl(
       \{P_a(w_j)-P_a(u_j)>\delta\};\Omega_a\bigr)
       \longrightarrow0.
\]
For a requested pair $H_0\Subset\Omega$, take the finite cofinal
cover of $H_0$. Domain monotonicity and finite subadditivity
transfer the zero bounds to $\Omega$. Taking the countable
union over $\delta=1/m$ proves
$v\preceq_{S,\mathrm{loc}}H$.

The least-upper-bound assertion follows from countability,
exactly as in Theorem~\ref{thm:r2-one-domain-envelope}: each
$u_j$ uses finitely many members of $\mathcal F$, and there
are only countably many $j$. The union of the corresponding
local outer-capacity-zero sets remains locally outer-capacity
zero. Uniqueness in the local ordered quotient follows.
\end{proof}

\begin{corollary}[Countable cofinal families]
If $\mathcal F$ has a pointwise cofinal countable subfamily, its
pointwise supremum is Borel and agrees locally quasi-everywhere
with $H$. This includes every countable family.
\end{corollary}

\begin{proof}
Let $v_m\in\mathcal F$ be pointwise cofinal. Then
$\sup\mathcal F=\sup_m v_m$, and $v_m\preceq_{S,\mathrm{loc}}H$
for all $m$. Countable subadditivity gives
$\sup\mathcal F\preceq_{S,\mathrm{loc}}H$. Conversely every
$u_j$ is a finite maximum of members of $\mathcal F$, so
$H\leq\sup\mathcal F$ pointwise.
\end{proof}

\subsection{Obstacle families and the scope of the result}

Let $\mathcal F$ be as in the local theorem, and suppose every
member satisfies the actual pointwise condition $v\leq\psi$
on a specified set $E$. Nonempty finite maxima preserve this
condition. The constructed representative $H=\sup_j u_j$
therefore satisfies $H\leq\psi$ at every point of $E$.
This gives an extremal object in the capacitary closure which
retains a prescribed pointwise obstacle while respecting the
background $S$.

These envelope constructions are recorded as \texttt{NEW\_PROVED}
consequences of the relative energy estimates and the completed
calculus. They do not assert that $H$ is classically psh, that
$H=(\sup\mathcal F)^*$, or that its Monge--Amp\`ere measure is
supported on a contact set. They do not establish compact outer
regularity or weak-closure domination for the original Borel
capacity for an arbitrary background. For the ambient background,
the later negligible-set theorem identifies $H$ with the ordinary
usc envelope quasi-everywhere, and the exhaustion-domain extremal
theorems supply the stated equilibrium and capacity formulas.
These later conclusions use additional arguments.

% END INCLUDED SOURCE: base_r2/sections/relative_capacitary_envelopes


% BEGIN INCLUDED SOURCE: base_r2/sections/domain_comparison
\section{Local comparison with the capacity of a fixed affine domain}
\label{sec:r2-domain-comparison}

Fix a smooth integral affine $K$-variety $Y$ of dimension $n\geq1$
and an ordinary open set
\[
 V\subset Y^{\mathrm{an}},\qquad
                    \overline V\Subset Y^{\mathrm{an}}.
 \tag{D}\label{eq:r2-domain-comparison-scope}
\]
Let $S$ be any symmetric original-positive, separately $d'$- and
$d''$-closed current of bidegree $(p,p)$ on $V$, and put
$q=n-p\geq1$. Restrictions of $S$ to smaller ordinary domains are
understood in the notation $\operatorname{Cap}_S(\,\cdot\,;\Omega)$.
All functions retain their actual point values and the all-extensions
classical definition.

The purpose of the construction is to compare admissible functions
on different domains without assuming that they extend unchanged.
An explicit smooth pluriharmonic correction on the inner core is
sufficient, because that correction disappears from every
differentiated BT slot over an arbitrary closed background.

\subsection{The effect of a change of frame}

\begin{lemma}[Smooth pluriharmonic invariance]
\label{lem:r2-harmonic-frame-invariance}
Let $O\subset V$ be ordinary open. Let $a_i,b_i\in\Cb(O)$ satisfy
$a_i=b_i+\psi_i$, where $\psi_i$ is original smooth and
$D\psi_i=0$. Then, for $0\leq k\leq q$,
\[
 T_{S,k}(a_1,\ldots,a_k)=T_{S,k}(b_1,\ldots,b_k)
 \quad\text{on }O.
\]
In particular this applies when each $\psi_i$ is a real constant
plus a real multiple of the logarithm of the absolute value of an
analytic unit.
\end{lemma}

\begin{proof}
For every separately closed current $Q$ the original smooth product
rule gives
\[
 D(\psi_i Q)=D\psi_i\wedge Q=0.
\]
All cross terms vanish because $d'Q=d''Q=0$. This is an identity
against all original compact smooth complementary tests. The result
now follows inductively from the ordered recursion: first replace
the inner $k-1$ slots by induction and then use the displayed identity
for the outer slot. Only multiplication by an original smooth
function is used in this last step. A unit logarithm is original
smooth and has zero $D$, as it is linear in its unit moment
coordinate. Its positive and negative are classical psh after every
complete extension.
\end{proof}

This \texttt{DIRECT\_COROLLARY} of the ordered calculus is much
narrower than a claim that arbitrary-background products depend only
on the ambient currents $D[a_i]$ for arbitrary potentials. Only the
specified smooth zero-curvature differences are used below.

\subsection{Uniform extension from one inner norm box}

Choose a closed affine embedding $Y\hookrightarrow\mathbf A^N_K$.
Let $Z$ be the normalization of its projective closure, let
\[
 A=\nu^*\mathcal O_{\mathbf P^N}(1),
\]
and write $e_\infty=\nu^*X_0$ for its fixed frame on $Y$. The line
$A$ is ample, and $Y$ identifies with the inverse image of the
standard affine chart because it is already normal. We do not assume
that $Z$ is smooth. No current or BT product on its singular locus
will be used.

The polynomial norm boxes below are those in the constructive part of
\textup{B}, Lemmas D.1--D.2. They use one finite list of polynomial
upper and lower bounds, with strict slack between the successive
thresholds and with upper bounds on the affine coordinates. In
particular the closed boxes are compact. They may always be chosen
locally inside a prescribed ordinary neighborhood.

\begin{lemma}[A uniform extension up to one fixed frame term]
\label{lem:r2-uniform-framed-extension}
Let $\Omega\subset V$ be ordinary open, and fix nested polynomial
norm boxes
\[
 C_1\subset\operatorname{Int}C_{\mathrm{mid}},\qquad
 C_{\mathrm{mid}}\subset\operatorname{Int}C_2,\qquad
 C_2\Subset\Omega,
\]
with the strict-threshold properties just described. Put
$W=\operatorname{Int}C_1$. There are a number $A_0>0$ and one
original smooth pluriharmonic function $\psi$ on a neighborhood of
$C_1$, depending on the boxes and $V$ but not on $v$ or $S$, such
that every $v\in\Cb(\Omega)$ with $0\leq v\leq1$ admits
$\widetilde v\in\Cb(V)$ satisfying
\[
 0\leq\widetilde v\leq1\text{ on }V,
 \qquad
 \widetilde v=A_0^{-1}v+\psi\text{ on a neighborhood of }C_1.
 \tag{E}\label{eq:r2-uniform-framed-extension}
\]
Consequently, for every $k\leq q$,
\[
 T_{S,k}(\widetilde v,\ldots,\widetilde v)
       =A_0^{-k}T_{S,k}(v,\ldots,v)\quad\text{on }W.
 \tag{P}\label{eq:r2-extension-product}
\]
\end{lemma}

\begin{proof}
We give the geometric construction in full to avoid invoking either a
smooth compactification or \textup{B}, Theorem~4.4.

Choose intermediate thresholds between those defining $C_1$ and
$C_{\mathrm{mid}}$. By clearing rational powers, the polynomial
collar can be written
\[
 b=\max\{0,\log|P_a^{N_a}/c_a|,
                       \log|d_l/Q_l^{M_l}|\},
 \qquad c_a,d_l\in K^\times.
\]
Every denominator $Q_l$ is nonzero on a neighborhood of $C_2$.
The chosen strict gaps give
\[
 b=0\text{ near }C_1,
 \qquad b\geq\delta>0\text{ on }C_2\setminus\operatorname{Int}
                                    C_{\mathrm{mid}}.
\]
The function is finite and continuous on a neighborhood of $C_2$.
It is classical psh there: each nonconstant branch is a logarithm of
an analytic function or of an analytic unit; the constant branch
makes branches with a zero inactive in a neighborhood of that zero.
These assertions persist after every complete extension.

Clear the denominators and homogenize in the fixed projective
embedding. After taking a common sufficiently high degree $m_0$,
all displayed branches are ratios $s_a/s_0$ of global sections of
$A^{m_0}$. Here $s_0$ is a product of homogenized denominator
polynomials and a power of $X_0$, pulled back to $Z$, and is nowhere
zero on a neighborhood of $C_2$. Include the branch $s_0/s_0=1$.
Append small nonzero $K$-scalar multiples of a finite global
generating tuple of $A^{m_0}$. The scalars can be chosen so that
all their ratios to $s_0$ have absolute value strictly less than one
on $C_2$, by compactness. They do not change the maximum $b$ there,
and the full tuple generates the line everywhere on $Z$.

Its pure Fubini--Study metric $q_0$ on $A^{m_0}$ has weight $b$
in the frame $s_0$ near $C_2$. This metric is classical on $Z$ in
the literal sense that every local frame weight is classical psh
after every complete extension. Indeed, pullback along any analytic
curve gives the finite maximum of logarithms of analytic functions
with no common zero. This argument uses generation and the sheaf
and maximum properties, and needs no smoothness of $Z$.

Fix $c>0$, put $M=c+1$, and choose an integer $m>0$ with
$m\delta>M+1$. Set
\[
 d=mm_0,\qquad L=A^d,\qquad e=s_0^m,\qquad Q=mq_0.
\]
All these objects are independent of $v$. Define a metric $R_v$ on
$L$ using the two ordinary opens
$\operatorname{Int}C_2$ and $Z^{\mathrm{an}}\setminus C_{\mathrm{mid}}$:
its $e$-weight on the first is
\[
 \max\{v+c,mb\},
\]
and on the second it equals $Q$. The opens cover $Z^{\mathrm{an}}$
because $C_{\mathrm{mid}}\subset\operatorname{Int}C_2$. On their
overlap the inequality $mb\geq m\delta>M\geq v+c$ makes the two
prescriptions identical. They therefore glue as metrics on the same
identified line. Their weights are all-extensions classical by
restriction, maxima, and locality on curves. This is an explicit
gluing construction on the normal variety, not an application of a
global approximation or envelope theorem there.

The metric difference is a globally defined bounded Borel function
and satisfies
\[
 0\leq R_v-Q\leq M.
\]
On $\operatorname{Int}C_2$ this is the identity
$R_v-Q=\max\{v+c-mb,0\}$ and the bound follows from $b\geq0$.
On the other open the difference is zero.

Now use the fixed frame $e_\infty^d$ on all of $Y$. Let
\[
 w_v=\operatorname{wt}_{e_\infty^d}(R_v),\qquad
 Q_\infty=\operatorname{wt}_{e_\infty^d}(Q).
\]
The function $w_v$ is classical psh on $Y^{\mathrm{an}}$, and
$Q_\infty$ is continuous there. In particular they are bounded on
$V$, uniformly in $v$, because
\[
 Q_\infty\leq w_v\leq Q_\infty+M
 \quad\text{and}\quad\overline V\Subset Y^{\mathrm{an}}.
\]
Write
\[
 a=\min_{\overline V}Q_\infty,
 \qquad
 A_0=1+\operatorname{osc}_{\overline V}(Q_\infty)+M,
 \qquad
 \widetilde v=\frac{w_v-a}{A_0}\Big|_V.
\]
This gives $0\leq\widetilde v\leq1$ on the original domain $V$.
The actual bounds and $A_0$ are independent of the competitor.

Near $C_1$ we have $b=0$ and $v+c>0$, so the $e$-weight of $R_v$
is exactly $v+c$. The ratio
\[
 f=\frac{s_0^m}{e_\infty^d}
\]
is an algebraic unit near $C_2$. Since $e=f e_\infty^d$, the metric
transition law gives
\[
 w_v=v+c+\log|f|\quad\text{near }C_1.
\]
Thus \eqref{eq:r2-uniform-framed-extension} holds with the fixed
function
$\psi=A_0^{-1}(c+\log|f|-a)$. It is original smooth
pluriharmonic. Positive homogeneity and
Lemma~\ref{lem:r2-harmonic-frame-invariance} prove
\eqref{eq:r2-extension-product}, with the original background and on
all original tests. The normal projective variety was used only to
construct the metric and its frame; every current in this final step
is on the smooth open $Y$.
\end{proof}

The lemma is recorded as \texttt{NEW\_PROVED} in the stated uniform
framed-affine scope, by the constructive collar mechanism of
\textup{B}, Lemma D.2. Its proof does not use that lemma's smooth
projective ambient hypothesis, a resolution of singularities, or
\textup{B}, Theorem~4.4. The following capacity consequences are
\texttt{DIRECT\_COROLLARY}.

\subsection{Borel and outer capacities on different domains}

For arbitrary subsets of a domain $O$, write
\[
 \operatorname{Cap}_S^*(E;O)
   =\inf_{\substack{G\supset E\\G\text{ ordinary open in }O}}
                 \operatorname{Cap}_S(G;O).
\]
This is the outer capacity by definition; no assertion that it agrees
with the original Borel capacity is used here.

\begin{theorem}[Uniform local comparison with a fixed domain]
\label{thm:r2-domain-capacity-comparison}
Assume \eqref{eq:r2-domain-comparison-scope}. For every ordinary
$\Omega\subset V$ and every ordinary $W$ with
$\overline W\Subset\Omega$, there is a finite constant
$C=C(Y,V,\Omega,W,q)>0$, independent of $S$, such that
\[
 \operatorname{Cap}_S(E;V)
 \leq\operatorname{Cap}_S(E;\Omega)
 \leq C\operatorname{Cap}_S(E;V)
 \quad(E\subset W\text{ Borel}).
 \tag{C}\label{eq:r2-domain-capacity-comparison}
\]
For arbitrary $E\subset W$ one also has
\[
 \operatorname{Cap}_S^*(E;V)
 \leq\operatorname{Cap}_S^*(E;\Omega)
 \leq C\operatorname{Cap}_S^*(E;V).
 \tag{C*}\label{eq:r2-domain-outer-comparison}
\]
For a single inner core in
Lemma~\ref{lem:r2-uniform-framed-extension}, one may take
$C=A_0^q$.
\end{theorem}

\begin{proof}
The first inequality of \eqref{eq:r2-domain-capacity-comparison}
is restriction of competitors: a normalized function on $V$ remains
normalized and classical on $\Omega$, and its product restricts
there by ordinary locality.

First suppose that $W$ is the core of one nested-box construction.
For an arbitrary competitor $0\leq v\leq1$ on $\Omega$, the extension
lemma and equality of actual Radon measures on $W$ give
\[
 \int_E T_{S,q}(v,\ldots,v)
 =A_0^q\int_E T_{S,q}(\widetilde v,\ldots,\widetilde v)
 \leq A_0^q\operatorname{Cap}_S(E;V).
\]
Taking the supremum over $v$ proves the second inequality for that
core. The constant is uniform over all competitors and all allowed
backgrounds.

For general $W$, cover the compact $\overline W$ by finitely many
inner cores $W_1,\ldots,W_N$ of nested polynomial norm boxes whose
outer boxes lie in $\Omega$. Let $A_i$ be the constants from their
extension lemmas. The same competitor $v$ on $\Omega$ can be
extended separately on each core. Positivity and a finite cover give
\[
 \begin{split}
 \int_E T_{S,q}(v,\ldots,v)
 &\leq\sum_{i=1}^N\int_{E\cap W_i}T_{S,q}(v,\ldots,v)\\
 &\leq\sum_{i=1}^NA_i^q\operatorname{Cap}_S(E\cap W_i;V)\\
 &\leq\left(\sum_{i=1}^NA_i^q\right)
                                 \operatorname{Cap}_S(E;V).
 \end{split}
\]
Take the supremum. Thus $C=\sum_iA_i^q$ works. This finite-cover
step avoids assuming that an arbitrary compact core lies in one
strictly smaller box of a previously chosen presentation. In
particular, no strict-inequality description of an ordinary interior
is silently imposed.

For outer capacities, every ordinary open of $\Omega$ is ordinary
open in $V$, so restriction and the first Borel inequality give the
first inequality of \eqref{eq:r2-domain-outer-comparison}. Conversely,
if $G\subset V$ is ordinary open and contains $E$, then $G\cap W$
is an ordinary open of $\Omega$ containing $E$. The already proved
Borel comparison and monotonicity give
\[
 \operatorname{Cap}_S^*(E;\Omega)
 \leq\operatorname{Cap}_S(G\cap W;\Omega)
 \leq C\operatorname{Cap}_S(G\cap W;V)
 \leq C\operatorname{Cap}_S(G;V).
\]
Take the infimum over such $G$. This argument invokes the definition
of outer capacity and no outer-regularity theorem.
\end{proof}

The endpoint $p=n$ can be included by defining the degree-zero
capacity to be $S(E)$; domain restriction then makes both Borel
capacities identical, with constant one.

\subsection{Equivalent local convergence and null sets}

\begin{corollary}[One fixed capacity domain suffices locally]
\label{cor:r2-fixed-domain-local-convergence}
Assume \eqref{eq:r2-domain-comparison-scope}, and let $a_j,a$ be
real Borel functions on $V$. The following are equivalent:
\begin{enumerate}
\item For every compact $L\Subset V$ and every $\delta>0$,
\[
 \operatorname{Cap}_S(L\cap\{|a_j-a|>\delta\};V)
                          \longrightarrow0.
\]
\item The functions converge in the all-domain local $S$-capacity
      sense of Definition~\ref{def:local-capacity-convergence},
      restricted to $V$.
\end{enumerate}
The same equivalence holds with the outer capacities throughout.
No uniform boundedness or classicality of the functions $a_j,a$ is
needed for this equivalence of convergence conditions.
\end{corollary}

\begin{proof}
Suppose (1) holds. Given a permitted relatively compact domain
$\Omega\subset V$ and a compact $L\Subset\Omega$, choose an
ordinary $W$ with $L\subset W$ and $\overline W\Subset\Omega$.
The second inequality of
\eqref{eq:r2-domain-capacity-comparison}, with this fixed geometry,
implies the required convergence using $\Omega$ as admissibility
domain.

Conversely, cover a prescribed compact $L\Subset V$ by finitely many
ordinary $W_i$ with compact closure in permitted small norm-box
domains $\Omega_i\Subset V$. Domain monotonicity and finite
subadditivity give
\[
 \begin{split}
 \operatorname{Cap}_S(L\cap\{|a_j-a|>\delta\};V)
 &\leq\sum_i\operatorname{Cap}_S
       (\overline W_i\cap\{|a_j-a|>\delta\};\Omega_i).
 \end{split}
\]
Each summand tends to zero by (2). The same proof uses
\eqref{eq:r2-domain-outer-comparison} and subadditivity of outer
capacity for the outer version.
\end{proof}

\begin{corollary}[Equivalence of local null classes]
\label{cor:r2-fixed-domain-null-classes}
For a Borel set $E\subset V$, its capacity with domain $V$ is zero
if and only if
\[
 \operatorname{Cap}_S(E\cap L;\Omega)=0
\]
for every permitted small domain $\Omega\subset V$ and compact
$L\Subset\Omega$. For arbitrary subsets $E$, the same assertion
holds with outer capacities.
\end{corollary}

\begin{proof}
One direction is the comparison theorem. For the converse, $V$ is
second countable in the present finite $p$-adic affine setting.
Choose countably many ordinary cores $W_i$ covering $V$, with
$\overline W_i\Subset\Omega_i$ for permitted small norm-box domains
$\Omega_i\Subset V$. Monotonicity under restriction and countable
subadditivity give
\[
 \operatorname{Cap}_S(E;V)
 \leq\sum_i\operatorname{Cap}_S(E\cap\overline W_i;\Omega_i)=0.
\]
The outer-capacity proof is identical. No ambient upper-semicontinuous
representative is substituted in this argument.
\end{proof}

These statements identify the local convergence and local null-set
notions on the qualified fixed domain $V$. They do not assert uniform
constants as the core approaches the boundary of $V$, convergence in
capacity on the entire noncompact domain from local convergence, or
capacity nondegeneracy. They remain valid for $S=0$ and backgrounds
supported on a proper analytic subvariety. They do not apply to an
arbitrary proper ambient space with only constant untwisted global
psh competitors. Every comparison is between ordinary domains; the
sets being measured may be arbitrary Borel sets, including non-open
strict gaps.

% END INCLUDED SOURCE: base_r2/sections/domain_comparison


% BEGIN INCLUDED SOURCE: base_r2/sections/support_envelopes
\section{Analytic supports and a relative-envelope obstruction}
\label{sec:r2-support-envelopes}

\begin{proposition}[BT products and capacity on an analytic retract]
\label{prop:r2-retract}
Let $i:Y\hookrightarrow U$ be a closed analytic embedding of allowed
pure-dimensional quasi-smooth spaces, of dimensions $m$ and $n$.
Let $R$ be a symmetric original-positive separately closed current of
bidegree $(a,a)$ on $Y$, with $0\leq a\leq m$, and put
\[
 S=i_*R,\qquad p=n-m+a,\qquad q=m-a=n-p.
\]
For $0\leq k\leq q$ and $u_1,\ldots,u_k\in\Cb(U)$,
\begin{equation}
\label{eq:r2-retract-product}
 T^U_{S,k}(u_1,\ldots,u_k)
   =i_*T^Y_{R,k}(u_1\circ i,\ldots,u_k\circ i)
\end{equation}
on all original compact smooth complementary tests.
If there is an analytic retraction $r:U\to Y$ with $r\circ i=\mathrm{id}$,
then, for every Borel $E\subset U$,
\begin{equation}
\label{eq:r2-retract-capacity}
 \operatorname{Cap}_{S}(E;U)
     =\operatorname{Cap}_{R}(i^{-1}E;Y).
\end{equation}
The corresponding equality holds for outer capacities of arbitrary
subsets. In degree zero the capacity is simply the measure of the set.
\end{proposition}

\noindent\textbf{Status: \texttt{DIRECT\_COROLLARY}.}
This is a functorial consequence of the fixed-background recursion,
not a restriction of that recursion to integration-current backgrounds.

\begin{proof}
The map $i$ is proper. Its pullback preserves the original positive
test cone and commutes with the two differentials; hence $S$ has the
stated positivity and closedness. For any coefficient-Radon current
$P$ on $Y$ and locally bounded Borel function $f$ on $U$,
\[
 f\,i_*P=i_*((f\circ i)P).
\]
Indeed, the coefficient measure of $i_*P$ against a test $\eta$ is
the pushforward of the coefficient measure of $P$ against $i^*\eta$.
This follows first on smooth scalar tests and then on Borel functions
by Radon uniqueness. Apply this identity and commutation of $D$ with
proper pushforward inductively through the recursion. Analytic
pullback preserves classical psh, so every restricted slot is allowed.
This proves \eqref{eq:r2-retract-product}.

Every normalized potential on $U$ restricts to a normalized potential
on $Y$, giving the inequality $\leq$ in
\eqref{eq:r2-retract-capacity}. Conversely every normalized $v$ on
$Y$ extends to $v\circ r$ on $U$, still classical and between zero
and one. Equation~\eqref{eq:r2-retract-product} gives the opposite
inequality. For outer capacities, restricting any ordinary open
neighborhood of $E$ gives one of $i^{-1}E$. Conversely, if $H\subset Y$
is open and contains $i^{-1}E$, choose an open $G_0\subset U$ with
$i^{-1}G_0=H$. The open set $G_0\cup(U\setminus i(Y))$ contains
$E$ and has capacity $\operatorname{Cap}_{R}(H;Y)$. Taking infima
proves the assertion.
\end{proof}

\begin{proposition}[Ambient negligible sets can have positive relative capacity]
\label{prop:r2-negligible-positive}
For every $m\geq1$, there is an allowed $(m+1)$-dimensional domain
$U$, a nonzero positive closed background $S$ of bidegree $(1,1)$,
and an increasing sequence $a_j\in\Cb(U)\cap C^0(U)$, with
$-1\leq a_j\leq0$, such that, writing $h=\sup_j a_j$,
\[
 h^*=0,\qquad
 \operatorname{Cap}_{S}(\{h<h^*\};U)>0.
\]
There is even a compact singleton $K\subset U$ of positive relative
capacity for which the ambient relative extremal function
\[
 h_K=\sup\{v\in\Cb(U):v\leq0,\ v|_K\leq-1\}
\]
satisfies $h_K^*=0$.
\end{proposition}

\noindent\textbf{Status of the assertions refuted: \texttt{FALSE}.}
The construction refutes both universal relative negligibility of
ambient envelope exceptions and the unqualified formula
$\operatorname{Cap}_S(K;U)=\int_U T_{S,m}((h_K^*)^m)$.

\begin{proof}
Choose $0<r<1<R$ and set
\[
 U=\{|z_0|<1\}\times\prod_{i=1}^m\{r<|z_i|<R\},
 \qquad Y=\{z_0=0\},\qquad S=i_*[1]_Y.
\]
The projection onto the last $m$ coordinates is an analytic
retraction. Put
\[
 a_j=\max\{-1,j^{-1}\log|z_0|\}.
\]
The logarithm is interpreted as $-\infty$ on $Y$. The maximum is
continuous and classical, it equals $-1$ on $Y$, and it increases
to zero at every point outside $Y$. Hence $h=-1$ on $Y$ and $h=0$
outside $Y$.

The complement of $Y$ is dense in $U$. For an explicit verification
at $y\in Y$, use the Gauss seminorm in the $z_0$ variable over the
completed residue field $\mathcal H(y)$, with radius $s>0$. It defines
a point $y_s$ of the original $K$-analytic product with
$|z_0(y_s)|=s$. As $s\downarrow0$, evaluations of polynomials tend
to their values at $y$, and so do local analytic functions by uniform
approximation on a compact product neighborhood. Thus $y_s\to y$.
Consequently $h^*=0$.

Let $\xi_Y$ be the product Gauss point of $Y$ and $K=\{i(\xi_Y)\}$.
For $u_i=\max\{0,\log|z_i|\}$, the standard model metric on
$(\mathbf P^1_{K^\circ})^m$ gives
\[
 T_{S,m}(u_1,\ldots,u_m)=\delta_{i(\xi_Y)}.
\]
This is exactly the mixed model calculation of the first record,
Proposition~\ref{prop:polar-relative-mass}, with the original
absolute value and calibration. Put $M=m\log R$ and
$v=M^{-1}\sum_i u_i$. Then $0\leq v\leq1$ on $U$ and positivity
and multilinearity give
\[
 T_{S,m}(v^m)(K)\geq \frac{m!}{(m\log R)^m}>0.
\]
In particular both $K$ and $Y=\{h<h^*\}$ have positive relative
capacity.

Every $a_j$ is a competitor for $h_K$. Thus $h_K=0$ off $Y$, while
$h_K\leq0$ everywhere. Density again gives $h_K^*=0$. The constant
$-1$ is a competitor and every competitor is at most $-1$ on $K$,
so $h_K=-1$ on $K$ itself. All differentiated slots of $h_K^*=0$
vanish, whereas the displayed capacity of $K$ is positive. This
proves the claimed failure without an appeal to an unspecified
pluripolar exceptional set.
\end{proof}

\paragraph{A correct replacement in this setting.}
For a background supported on an analytic retract, the intrinsic
capacity and extremal problem on $Y$ are the appropriate comparison
objects, by Proposition~\ref{prop:r2-retract}. Ambient upper
regularization can erase values on $Y$ that the background detects.
For a general positive closed current its support need not be an
analytic space; this observation does not by itself define an
intrinsic psh envelope on that support. Such a construction remains a
separate analytic problem.

% END INCLUDED SOURCE: base_r2/sections/support_envelopes


% BEGIN INCLUDED SOURCE: base_r2/sections/gap_detection
\section{Detection of strict gaps and actual point values}
\label{sec:r2-gap-detection}

All fields, currents, and test forms retain the conventions of the first
round. In particular, $K/\mathbf Q_p$ is finite with its prescribed
absolute value, $D=d'd''$, and $\Cb$ denotes the actual pointwise,
locally bounded, all-extensions classical psh class. No representative
is altered on a null set. This section supplies a sufficient geometric
condition for the strict-gap hypothesis in
Corollary~\ref{cor:conditional-pointwise-domination}. The condition is
proved for $S=[1]$ on the domains specified below, and for an explicit
family of positive closed backgrounds. It is not asserted for every
positive closed background or every ordinary domain.

\subsection{Shilov boundaries determine pointwise order}

\begin{definition}[A family of affinoid neighborhoods]
Let $V$ be an ordinary open in a good strictly analytic space. A family
$\mathscr A$ of strict affinoid domains contained in $V$ is called an
affinoid neighborhood basis if, whenever $x\in O\subset V$ and $O$ is
ordinary open, some $A\in\mathscr A$ satisfies
\[
 x\in\operatorname{Int}_V(A)\subset A\subset O.
\]
Write $\Gamma(A)$ for its finite Shilov boundary and put
$\mathscr D=\bigcup_{A\in\mathscr A}\Gamma(A)$.
\end{definition}

The external input used here is \cite[Theorem~6.4, pp.~33--34]{GR25}:
a classical psh function on an affinoid space attains its maximum at a
Shilov point. Its definition of classical psh uses curve pullback after
every complete extension. The theorem applies to the restriction of a
function in $\Cb(V)$ to each $A$; it does not require that the function
be continuous. The following lemma is a \texttt{DIRECT\_COROLLARY} of
that theorem and upper semicontinuity.

\begin{lemma}[A strict gap contains a Shilov point]
\label{lem:r2-shilov-gap}
Let $a$ be a finite classical psh function on $V$, let $b$ be a finite
upper-semicontinuous function on $V$, and let $\delta\geq0$. If
$a(x)>b(x)+\delta$, then every sufficiently small ordinary neighborhood
of $x$ contains a point $z\in\mathscr D$ with
$a(z)>b(z)+\delta$. Consequently,
\[
 a\leq b\text{ on }\mathscr D\quad\Longrightarrow\quad a\leq b
 \text{ on }V.
\]
In particular, two members of $\Cb(V)$ that agree on $\mathscr D$
agree at every point of $V$.
\end{lemma}

\begin{proof}
Choose $\epsilon>0$ with
$b(x)+\delta+2\epsilon<a(x)$. Upper semicontinuity gives an ordinary
neighborhood $O$ of $x$ on which $b<b(x)+\epsilon$. Choose
$A\in\mathscr A$ with $x\in\operatorname{Int}_V(A)$ and $A\subset O$.
The affinoid maximum principle gives $z\in\Gamma(A)$ such that
$a(z)=\max_A a\geq a(x)$. Therefore
\[
 a(z)\geq a(x)>b(x)+\delta+2\epsilon>b(z)+\delta.
\]
The same construction can be made inside any prescribed ordinary
neighborhood of $x$, by intersecting it with $O$. This proves the first
assertion. If pointwise order failed, apply it with $\delta=0$ to
contradict order on $\mathscr D$. Apply that implication in both orders
to obtain equality. The strict-gap set itself has nowhere been treated
as ordinary open.
\end{proof}

For example, for a fixed $x\in V$ the same maximum principle and upper
semicontinuity give the recovery formula
\[
 a(x)=\inf_{\substack{A\in\mathscr A\\
                         x\in\operatorname{Int}_V(A)}}
                  \max_{z\in\Gamma(A)}a(z).
\]
The infimum is over the members of the full neighborhood basis. This
is a pointwise statement, rather than an assertion about equality
classes modulo any measure.

\subsection{An atom at a prescribed divisorial point}

For the fixed-domain result, assume that $Y/K$ is a smooth integral
affine variety of dimension $n\geq1$ and that
\[
 V\subset Y^{\mathrm{an}}
 \quad\text{is ordinary open with}\quad
 \overline V\Subset Y^{\mathrm{an}}.
 \tag{G}\label{eq:r2-framed-domain}
\]
The closure here is the ordinary topological closure. More generally,
the argument applies when $V$ has compact closure inside an algebraic
open of a normal projective variety on which a fixed ample line bundle
has a specified nowhere-vanishing frame, and that open is smooth.
No convexity, exhaustion, or regularity of the boundary of $V$ is
assumed. Small norm boxes in the algebraic neighborhoods supplied by
\textup{B} satisfy \eqref{eq:r2-framed-domain}, after choosing one
affine chart.

We use the following external facts with their original field:
\cite[Proposition~A.3, p.~365; Corollary~A.6, p.~367;
Proposition~A.9, p.~368]{GM19} identify Shilov points of strict
affinoids with divisorial points, show locality of divisoriality, and
realize divisorial points on algebraic $K^\circ$-models of a proper
variety. The ample extension theorem is
\cite[Proposition~4.11, p.~353]{GM19}. The ordinary CLD model-measure
formula is \cite[Theorem~6.9.3 and Proposition~6.9.7]{CLD12}; its
generically smooth and stable-field hypotheses apply here. For smooth
projective varieties the same normalization is explicitly recorded in
\textup{R}, Lemma~3.3, and \textup{B}, equation~(4.1). We only use the
ordinary-current measure, not equality on all $\delta$-test forms.

\begin{lemma}[A bounded model tester on the original domain]
\label{lem:r2-atomic-tester}
Assume \eqref{eq:r2-framed-domain}. For every divisorial point
$z\in V$ there exists a continuous semipositive real PL function
$h_z\in\Cb(V)$ such that
\[
 0\leq h_z\leq1\quad\text{on }V,
 \qquad (Dh_z)^n(\{z\})>0.
\]
The function is defined and bounded on a neighborhood of $\overline V$.
If a larger relatively compact ordinary $V^+$ with
$\overline V\subset V^+\Subset Y^{\mathrm{an}}$ is fixed in advance,
one can arrange $0\leq h_z\leq1$ on this same $V^+$ for every $z$.
\end{lemma}

\begin{proof}
Choose a closed embedding $Y\hookrightarrow\mathbf A^N_K$ and take
the normalization $Z$ of its closure in $\mathbf P^N_K$. Because $Y$
is normal, the inverse image of the standard affine chart in $Z$ is
identified with $Y$. If $\nu:Z\to\overline Y$ is the normalization,
then $L=\nu^*\mathcal O_{\mathbf P^N}(1)$ is ample, and the pullback
of the homogeneous coordinate $X_0$ gives a fixed frame $e$ of $L$
on $Y$. Thus no modification of the open line or its frame has been
made.

The point $z$ is divisorial in $Z^{\mathrm{an}}$. By the cited model
realization theorem there is an algebraic $K^\circ$-model
$\mathscr Z_0$ of $Z$ on which $z$ is associated with an irreducible
component of the special fiber. Normalize this model. Normalization is
finite in the present excellent discretely valued situation, and the
center selected by $z$ is still a special-fiber component. Apply the
ample extension theorem to obtain a dominating model $\mathscr Z$
and an ample line bundle $\mathscr L$ with a specified identification
\[
 \mathscr L|_Z\simeq L^{\otimes m},\qquad m>0.
\]
Normalize once more if necessary and pull back $\mathscr L$; finite
pullback preserves ampleness. A dominating model of a normal model
preserves each codimension-one divisorial valuation represented on the
latter. Indeed, after localizing the target at the generic point of
its component one has a discrete valuation ring, and the proper
birational morphism is an isomorphism at that generic point. Hence $z$
is still associated with a special-fiber component $E$ of
$\mathscr Z$.

For clarity, the ample extension used here does not require resolution
of mixed-characteristic models. One may first choose a projective
model carrying an ample extension of a power of $L$, take a vertical
blowup dominating $\mathscr Z_0$, and tensor a sufficiently high
pullback power with its relatively ample tautological line. The center
of the blowup lies in the special fiber, so the tautological line has
trivial generic fiber. This produces an ample extension of the stated
power of $L$, before normalization. This is an existence of a
dominating model, not a claim about an arbitrary previously fixed
model.

The model metric of $\mathscr L$ gives the weight
\[
 \phi=-\log\|e^{\otimes m}\|_{\mathscr L}
 \quad\text{on }Y^{\mathrm{an}}.
\]
It is continuous and semipositive PL. In particular it is classical
psh after all complete extensions by
\cite[Theorem~4.3, pp.~22--23]{GR25}. Equivalently, replace
$\mathscr L$ by a positive power that is generated by model sections.
Its weight is then a pure Fubini--Study maximum in the same frame;
the inactive-branch rule handles zeros of the sections. This also
shows directly that all functions used are defined over the original
$K$.

On a normal algebraic model the original CLD model formula gives the
atom
\[
 (D\phi)^n(\{z\})
   =b_E\deg_{\widetilde K,E}\bigl((\mathscr L|_E)^n\bigr)>0,
\]
where $b_E$ is the special-fiber multiplicity and the structural degree
retains residue-field degrees. Ampleness makes the displayed degree
strictly positive. This formula is interpreted in the original
calibration; neither the absolute value nor the integration current
is rescaled. The continuous BT product on the smooth open $Y$ is this
ordinary CLD measure by the accepted continuous compatibility of
\textup{R} and \textup{B}. If a power was used to obtain generating
sections, its tensor degree is retained in this formula and in the
frame, or divided out in every slot.

Fix $V^+$ as in the statement, taking one if none was prescribed.
The continuous function $\phi$ is bounded on $\overline {V^+}$. Choose
real numbers $A<B$ with
$A\leq\phi\leq B$ there, and set $h_z=(\phi-A)/(B-A)$. A positive
constant multiple and a constant translate preserve classicality and
semipositive real PL structure. They multiply the displayed atom by
$(B-A)^{-n}>0$. This proves all assertions on the fixed domain $V$.
\end{proof}

The lemma is a \texttt{DIRECT\_COROLLARY} of the cited model geometry,
model-measure formula, and the accepted continuous compatibility. It
does not solve a prescribed Dirac Monge--Amp\`ere equation: other
positive atoms of $(Dh_z)^n$ are allowed.

\begin{theorem}[Strict-gap detection for the ambient capacity]
\label{thm:r2-ambient-gap-detection}
Assume \eqref{eq:r2-framed-domain}. If $a,b\in\Cb(V)$ and
$\{a>b+\delta\}$ is nonempty for some $\delta\geq0$, then
\[
 \operatorname{Cap}_{[1]}(\{a>b+\delta\};V)>0.
\]
Here the supremum defining capacity uses functions on the original
domain $V$. In particular the strict-gap hypothesis of
Corollary~\ref{cor:conditional-pointwise-domination} holds for
$S=[1]$ on $V$.
\end{theorem}

\begin{proof}
Choose a strict affinoid neighborhood basis in $V$. By
Lemma~\ref{lem:r2-shilov-gap}, the gap contains a Shilov point $z$
of one of these affinoids. Such a point is divisorial in
$Y^{\mathrm{an}}$ by the cited locality of divisorial points. Choose
the tester $h_z$ from Lemma~\ref{lem:r2-atomic-tester}. The gap is
Borel, since both potentials are upper semicontinuous. Thus
\[
 \operatorname{Cap}_{[1]}(\{a>b+\delta\};V)
 \geq\int_{\{a>b+\delta\}}(Dh_z)^n
 \geq(Dh_z)^n(\{z\})>0.
\]
Every tester in this inequality is defined on $V$. There is no use of
an implication from positive capacity on a smaller domain to positive
capacity on $V$.
\end{proof}

This is a proved consequence of the two lemmas, independently of any
outer regularity or identification of polar sets. In the project
ledger it is a \texttt{DIRECT\_COROLLARY} of those inputs and the
definition of capacity.

\begin{corollary}[Pointwise ambient domination and uniqueness]
\label{cor:r2-pointwise-ambient}
Assume \eqref{eq:r2-framed-domain} and let $u,v\in\Cb(V)$. If
\[
 \overline{\{u<v-\epsilon\}}\Subset V\quad(\epsilon>0),
 \qquad (Du)^n(\{u<v\})=0,
\]
then $u\geq v$ at every point of $V$. If the compact-gap condition
holds in both orders and $(Du)^n=(Dv)^n$, then $u=v$ at every point.
\end{corollary}

\begin{proof}
Apply Proposition~\ref{prop:capacity-domination} or
Proposition~\ref{prop:capacity-uniqueness}, respectively, and then
Theorem~\ref{thm:r2-ambient-gap-detection} to any prohibited strict
inequality with a positive margin. This uses only Borel restrictions
of the measures.
\end{proof}

\subsection{Faithfulness of original trace currents}

\begin{proposition}[Actual classical representatives are determined by
their original currents]
\label{prop:r2-current-faithfulness}
Let $U$ have the local algebraic neighborhoods supplied by \textup{B},
and let $u,v\in\Cb(U)$. Then
\[
 [u]\leq[v]\text{ as original CLD }(0,0)\text{-currents}
 \quad\Longleftrightarrow\quad u\leq v\text{ pointwise on }U.
\]
On the left, $[v]-[u]$ is nonnegative on every positive original
compactly supported $(n,n)$-form. In particular $[u]=[v]$ implies
actual pointwise equality. This proposition uses the fixed-background
weighted monotone theorem of the first round; it makes no assertion
of stability for arbitrary varying backgrounds.
\end{proposition}

\begin{proof}
Pointwise order plainly implies current order. For the converse it is
enough to work in a domain satisfying
\eqref{eq:r2-framed-domain}, inside an algebraic neighborhood of
each point. Fix a divisorial point $z$ there, and choose its model
tester $h$ by Lemma~\ref{lem:r2-atomic-tester}. On a fixed smaller
ordinary neighborhood $W$ of $z$, the smooth approximation theorem of
\textup{R} gives smooth psh functions converging uniformly to $h$.
The pure-FS realization in the proof of the lemma allows this also
to be obtained directly from \textup{R}, Proposition~A.1. After
taking a subsequence $s_j$ with
$\|s_j-h\|_W\leq4^{-j}$, set
\[
 h_j=s_j+2\cdot4^{-j}.
\]
Then $h_j\downarrow h$ on the same $W$, since
$h+4^{-j}\leq h_j\leq h+3\cdot4^{-j}$.
Each $h_j$ is smooth and belongs to the all-extensions classical
class: locally it is a convex, coordinatewise nondecreasing smooth
function of logarithms of analytic units, with zero-section branches
inactive. Its curvature is an original positive $(1,1)$-form.

For every nonnegative original smooth scalar
$\chi\in A_c^{0,0}(W)$, the form $\chi(Dh_j)^n$ is positive and
compactly supported. The assumed current order therefore gives
\[
 \big\langle (v-u)(Dh_j)^n,\chi\big\rangle\geq0.
\]
Apply Theorem~\ref{thm:relative-monotone} with fixed background
$[1]$, all differentiated slots equal to $h_j$, and the weight equal
to the constant sequence $u$ or $v$. It yields
\[
 (v-u)(Dh_j)^n\longrightarrow(v-u)(Dh)^n
\]
on all original scalar tests. The limit is therefore a positive Radon
measure. Its atom at $z$ is
$(v(z)-u(z))(Dh)^n(\{z\})$, and the second factor is strictly
positive. Hence $u(z)\leq v(z)$.

This holds at every point in the Shilov determining set of the local
domain. Lemma~\ref{lem:r2-shilov-gap} gives $u\leq v$ throughout
that domain, and these domains cover $U$ by \textup{B}. This proves
the converse. Equality follows by interchanging $u$ and $v$.
\end{proof}

Thus equality almost everywhere for all local original trace measures
does determine the actual classical representatives: it implies
equality of their $(0,0)$-currents, to which the proposition applies.
The essential passage to an atomic tester uses weighted convergence
of a decreasing potential with fixed background $[1]$. Weak
convergence of arbitrary background measures would not justify it.
The proposition is a \texttt{DIRECT\_COROLLARY} of that theorem,
the atomic testers, and the Shilov maximum principle.

\subsection{A precise condition on an arbitrary background}

Let $S$ be an original symmetric positive separately $d'$- and
$d''$-closed current of bidegree $(p,p)$ on $V$, with
$q=n-p\geq1$. An explicit sufficient condition is the following:
there is an affinoid neighborhood basis $\mathscr A$ such that, for
each $z\in\mathscr D=\bigcup_{A\in\mathscr A}\Gamma(A)$, there is
$h_z\in\Cb(V)$ satisfying
\[
 0\leq h_z\leq1,
 \qquad T_{S,q}(h_z,\ldots,h_z)(\{z\})>0.
 \tag{A}\label{eq:r2-relative-atomic-condition}
\]
The condition is stated in actual Borel measures on $V$, not only in
one tropical image.

\begin{proposition}[Pointwise conclusions under atomic detection]
\label{prop:r2-relative-gap-detection}
If \eqref{eq:r2-relative-atomic-condition} holds, every nonempty
strict gap of two members of $\Cb(V)$ has positive
$\operatorname{Cap}_S(\,\cdot\,;V)$. Consequently the pointwise
domination and uniqueness conclusions of
Corollary~\ref{cor:r2-pointwise-ambient} hold with
$(Da)^n$ replaced by $T_{S,q}(a,\ldots,a)$.
\end{proposition}

\begin{proof}
Lemma~\ref{lem:r2-shilov-gap} supplies a point $z$ of the gap in
$\mathscr D$. Its tester in \eqref{eq:r2-relative-atomic-condition}
assigns positive measure to the gap. The definition of capacity and
the first-round domination and uniqueness propositions finish the
proof.
\end{proof}

For a general $S$ this is \texttt{CONDITIONAL}: the condition is not
deduced from positivity, nonzero mass, or full support alone. The
next construction shows that it is nonvacuous in every allowed degree.

\begin{proposition}[Backgrounds with detection in every degree]
\label{prop:r2-faithful-backgrounds}
Assume \eqref{eq:r2-framed-domain}. There is a sequence of normalized
continuous semipositive PL functions $h_j$ on a fixed neighborhood
of $\overline V$ such that, for every $0\leq p<n$, the series
\[
 S_p=\sum_{j\geq1}2^{-j}T_{[1],p}(h_j,\ldots,h_j)
\]
defines a symmetric positive separately closed original current on
$V$ satisfying \eqref{eq:r2-relative-atomic-condition}. Here
$T_{[1],0}=[1]$. More generally, any positive closed background $S$
of the same degree with $S-\epsilon S_p$ positive for some
$\epsilon>0$ satisfies that condition.

The measure
\[
 \sigma=\sum_{j\geq1}2^{-j}(Dh_j)^n
\]
is a finite positive Radon measure on $V$ with full support. If two
members of $\Cb(V)$ are ordered $\sigma$-almost everywhere, they
are ordered at every point of $V$. In particular this one fixed
measure determines their actual representatives.
\end{proposition}

\begin{proof}
The space $Y^{\mathrm{an}}$ is second countable in the present setting.
Here is the elementary reason needed for this assertion. A finite
extension of $\mathbf Q_p$ has a countable dense subfield. On each
bounded coordinate region, polynomials with coefficients in that
subfield approximate arbitrary $K$-polynomials uniformly. Their
absolute-value functions therefore determine the analytic topology,
and inverse images of intervals with rational endpoints give a
countable basis. Passing to a closed affine subvariety and an ordinary
open preserves second countability.

For each member $O$ of a countable ordinary basis, cover $O$ by
interiors of strict affinoid domains contained in $O$. Good strict
analyticity supplies the cover, and second countability makes $O$
Lindel\"of, so a countable subcover suffices. The resulting countable
family is an affinoid neighborhood basis. Its Shilov determining set
$\mathscr D=\{z_1,z_2,\ldots\}$ is countable. Fix
$\overline V\subset V^+\Subset Y^{\mathrm{an}}$ and choose the
testers $h_j=h_{z_j}$ of Lemma~\ref{lem:r2-atomic-tester}, normalized
on this one $V^+$.

Relative CLN with background $[1]$ bounds the coefficient variations of
$T_{[1],p}(h_j,\ldots,h_j)$ uniformly in $j$ on every fixed compact
finite original presentation in $V^+$. The bounds use only
$0\leq h_j\leq1$ there. Thus the positive series defining $S_p$
converges absolutely in each such coefficient measure. The resulting
measures satisfy the original relations and refinement compatibilities,
and so define an original current. Positivity and symmetry pass term
by term. Closedness follows by applying $d'$ or $d''$ to a fixed
compact smooth test and summing the absolutely convergent pairings.
The same estimate on the compact set $\overline V\subset V^+$
shows that $\sigma(V)<\infty$.

Fix $j$ and let $R_j=S_p-2^{-j}T_{[1],p}(h_j,\ldots,h_j)$.
The current $R_j$ is positive and separately closed. Finite additivity
in the background follows directly from the defining recursion and
linearity of coefficient Borel multiplication. Therefore, for
$q=n-p$,
\[
 \begin{split}
 T_{S_p,q}(h_j,\ldots,h_j)
 &=2^{-j}T_{T_{[1],p}(h_j,\ldots,h_j),q}(h_j,\ldots,h_j)
     +T_{R_j,q}(h_j,\ldots,h_j)\\
 &\geq 2^{-j}(Dh_j)^n.
 \end{split}
\]
The first term is the indicated $n$-fold product by its literal ordered
recursion. In particular it has a positive atom at $z_j$. This proves
\eqref{eq:r2-relative-atomic-condition} for $S_p$. If
$S-\epsilon S_p$ is positive, the same finite decomposition proves it
for $S$. No assertion of continuity under arbitrary varying
backgrounds is used to pass from the series to these products.

Finally, $\sigma(\{z_j\})>0$ for every $j$. The set $\mathscr D$ is
dense because every nonempty ordinary open contains one of the
affinoids in the basis and hence a Shilov point. Thus $\sigma$ has
full support. Almost-everywhere order forces order at each of its
positive atoms, and Lemma~\ref{lem:r2-shilov-gap} then gives order
everywhere. The conclusion uses these particular determining atoms,
not merely the assertion that $\sigma$ has full support.
\end{proof}

This construction and its proof are recorded as
\texttt{NEW\_PROVED} relative to the adopted inputs, without a claim
of bibliographic priority. They exhibit backgrounds to which the
conditional criterion applies. They do not identify all backgrounds
with the desired property.

\subsection{Limits of the conclusion and excluded arguments}

\begin{enumerate}
\item The domain in \eqref{eq:r2-framed-domain} is essential to the
fixed-domain proof. On a connected proper positive-dimensional
analytic space, the global maximum principle makes every finite
global classical psh function constant, so the untwisted ambient
capacity defined using such global functions is zero. Distinct
constants then give nonempty capacity-zero strict gaps, while their
Monge--Amp\`ere measures both vanish. The first-round statements modulo
capacity zero remain true; a pointwise uniqueness assertion on a proper
space needs a normalization. No domain-independent detection theorem
is claimed.

\item A positive closed background may vanish off a proper analytic
subset. In an ordinary ball disjoint from that subset, every product
with this background vanishes by restriction. Nonempty strict gaps
there have zero relative capacity. In particular ambient positivity
does not furnish condition \eqref{eq:r2-relative-atomic-condition}.
The ambient-polar concentration example of the first round is fully
compatible with the present results.

\item Positive capacity on all nonempty ordinary opens, or full
support of one trace measure, is not substituted for strict-gap
detection. A gap of two upper-semicontinuous functions need not be
ordinary open. Lemma~\ref{lem:r2-shilov-gap} is the step that actually
relates its pointwise content to detectable points.

\item Local positivity of capacity on a smaller domain does not by
itself prove positivity for the capacity of a larger prescribed
domain: admissible competitors must extend to the latter. The ample
frame and bounded closure in Lemma~\ref{lem:r2-atomic-tester} were used
precisely to construct the tester on that larger fixed domain.

\item An unconditional pointwise domination or uniqueness theorem for
arbitrary $S$ is \texttt{FALSE}. For example, take $S=0$ on the open
unit disc, $u=\max\{\log|T|,-1\}$, and $v=0$. Both functions are
continuous bounded classical psh; their relative Monge--Amp\`ere
measures vanish. For $0<\epsilon<1$, the ordinary closure of
$\{u<v-\epsilon\}$ is contained in the compact closed disc of radius
$e^{-\epsilon}<1$;
for $\epsilon\geq1$ the set is empty. The reverse gaps are empty.
Thus the compact-gap hypotheses hold in both orders, but $u<v$
throughout the domain. What remains \texttt{OPEN} here is finding
necessary conditions, or weaker useful sufficient conditions, for
backgrounds not covered by
\eqref{eq:r2-relative-atomic-condition}. The atomic condition may be
stronger than necessary. No implication from an unspecified smooth
strict-positivity condition to detection is asserted.
\end{enumerate}

% END INCLUDED SOURCE: base_r2/sections/gap_detection


% BEGIN INCLUDED SOURCE: base_r2/sections/gap_detection_addendum
\subsection{A connected example with a nonzero background}
\label{subsec:r2-nonzero-background-counterexample}

The failure of unconditional pointwise domination is not limited to
the zero current or disconnected spaces.

\begin{example}[A nonzero connected-domain counterexample]
\label{ex:r2-nonzero-domination-failure}
Let $n\geq2$, let $R>1$, and let
\[
 U=\{\,|z_i|<R:0\leq i<n\,\}\subset(\mathbf A_K^n)^{\mathrm{an}},
 \qquad Y=\{z_0=0\}\cap U,
 \qquad S=[Y].
\]
The open polydisc $U$ is connected and $S$ is a nonzero symmetric
positive separately closed original current of bidegree $(1,1)$.
Choose $0<r<1$ and $0<c<\log(1/r)$, and set
\[
 b=\max\bigl\{0,\log(|z_0-1|/r),
                    \log(|z_1|/r),\ldots,
                    \log(|z_{n-1}|/r)\bigr\},
 \qquad u=b,\quad v=\max\{c,b\}.
\]
Then $u,v$ are continuous bounded classical psh on $U$. They satisfy
the compact-gap condition in both orders and
\[
 T_{S,n-1}(u,\ldots,u)=T_{S,n-1}(v,\ldots,v),
 \qquad
 T_{S,n-1}(u,\ldots,u)(\{u<v\})=0,
\]
but $u<v$ on a nonempty relatively compact ordinary polydisc.
Consequently unconditional pointwise relative domination and
uniqueness are \texttt{FALSE} even for a nonzero background on a
connected domain satisfying \eqref{eq:r2-framed-domain}.
\end{example}

\begin{proof}
Truncation by the constant zero makes the displayed maxima finite and
continuous across the zero sets of the coordinates. Logarithms of
analytic functions are classical psh after all complete extensions,
and finite maxima and positive constant translations preserve that
class. The coordinate bounds on $U$ give actual uniform upper bounds
for $b$; it is bounded below by zero.

Put $\rho=re^c<1$. Then
\[
 \{u<v\}=\{b<c\}
  =\{\,|z_0-1|<\rho,
              |z_i|<\rho\ (1\leq i<n)\,\}.
\]
This polydisc is nonempty. Its closure is contained in the closed
polydisc with the same center and radius. That closed polydisc is
compact and lies in $U$, since $|z_0|=1<R$ on it and
$|z_i|<1<R$ for $i\geq1$.
It is disjoint from $Y$. More precisely, for $0<\epsilon<c$ the gap
$\{u<v-\epsilon\}$ is the same polydisc with radius
$re^{c-\epsilon}<1$, while for $\epsilon\geq c$ it is empty.
The reverse gaps are empty because $v\geq u$. These facts check every
compact-gap quantifier.

On the ordinary neighborhood $\{|z_0|<1\}\cap U$ of $Y$, the
non-Archimedean triangle inequality gives $|z_0-1|=1$. Hence
$b\geq\log(1/r)>c$ there and $u=v$ there. Every recursively defined
product with background $[Y]$ is supported on $Y$: Borel
multiplication does not enlarge the support, and distributional
differentiation does not enlarge it. The products for $u$ and $v$
are equal on this neighborhood by ordinary locality, and vanish
outside $Y$; hence they are equal on $U$. Their common measure assigns
zero mass to the displayed gap because it is disjoint from $Y$.
No base change and no modified calibration enter the example.
\end{proof}

\subsection{Detection by backgrounds constructed from smooth factors}
\label{subsec:r2-smooth-factor-detection}

Atomic detection by $S$-products is a sufficient condition, but the
proof can also use a controlled family of smooth factors approaching
the ambient atomic testers. The following construction makes no claim
about whether its resulting measures have atoms.

Continue to assume \eqref{eq:r2-framed-domain}. Use the determining
set $\mathscr D=\{z_j\}$ and the continuous model testers $h_j$ in
Proposition~\ref{prop:r2-faithful-backgrounds}, normalized on the fixed
$V^+$. For each $j$ choose smooth classical psh functions $g_{j,m}$ on
$Y^{\mathrm{an}}$ such that
\[
 g_{j,m}\downarrow h_j,\qquad
 h_j+4^{-m}\leq g_{j,m}\leq h_j+3\cdot4^{-m}.
 \tag{S}\label{eq:r2-smooth-detectors}
\]
Here $h_j$ denotes its original model weight with the same constant
normalization, also defined on $Y^{\mathrm{an}}$. These approximants
exist on the whole required domain: after taking a generated power,
the model metric is a finite pure-FS maximum of algebraic section
ratios in its specified frame. The smooth regularized maxima of
\textup{R}, Proposition~A.1 have uniformly bounded error everywhere;
zeros of sections are inactive locally. Apply the constant-translation
construction of \textup{R}, Lemma~A.2 to that global uniform error.
The functions remain all-extensions classical by the convex
unit-logarithm description. Set
\[
 f_{j,m}=\tfrac12 g_{j,m}.
\]
Thus $0\leq f_{j,m}\leq1$ on $V^+$, and
$f_{j,m}\downarrow h_j/2$ on the same domain.

\begin{proposition}[A sufficient condition using smooth factors]
\label{prop:r2-smooth-factor-condition}
Let $0\leq p<n$ and $q=n-p$. Suppose that a symmetric positive
separately closed original $(p,p)$-current $S$ on $V$ has the following
property: for every $j,m$ there is a constant $a_{j,m}>0$ such that
\[
 S-a_{j,m}(Df_{j,m})^p\text{ is a positive current on }V.
 \tag{F}\label{eq:r2-smooth-factor-condition}
\]
As usual, the right-hand factor is interpreted as
$(Df_{j,m})^p\wedge[1]$, and equals $[1]$ when $p=0$.
Then every nonempty strict gap of two members of $\Cb(V)$ has
positive $\operatorname{Cap}_S(\,\cdot\,;V)$. The pointwise relative
domination and uniqueness conclusions hold under the first-round
compact-gap hypotheses.
\end{proposition}

\begin{proof}
Finite additivity in the background and positivity give
\[
 T_{S,q}(f_{j,m},\ldots,f_{j,m})
       \geq a_{j,m}(Df_{j,m})^n.
\]
Let $u,v\in\Cb(V)$ and $\delta\geq0$, and suppose that
$E=\{u>v+\delta\}$ has $S$-capacity zero. Since each $f_{j,m}$ is an
admissible capacity tester, $(Df_{j,m})^n(E)=0$ for every $j,m$.
Put $w=\max\{u,v+\delta\}\in\Cb(V)$. The nonnegative Borel
function $w-v-\delta$ is zero off $E$, so
\[
 (w-v-\delta)(Df_{j,m})^n=0.
\]
Fix $j$. Theorem~\ref{thm:relative-monotone}, applied with fixed
background $[1]$ and differentiated slots $f_{j,m}\downarrow h_j/2$,
passes this identity to the limit. The two nonconstant weights $w$
and $v$ are constant sequences in that theorem; the scalar constant
$\delta$ uses its unweighted case. Thus
\[
 (w-v-\delta)(D(h_j/2))^n=0.
\]
All statements are local on compact original tests, so local
boundedness of $u,v,w$ suffices. The limiting measure has a positive
atom at $z_j$, which forces
$w(z_j)-v(z_j)-\delta=0$, or $u(z_j)\leq v(z_j)+\delta$.
Lemma~\ref{lem:r2-shilov-gap} then gives $u\leq v+\delta$ everywhere
on $V$. Hence $E$ was empty. The pointwise comparison consequences
follow from the first-round results modulo capacity zero.
\end{proof}

\begin{corollary}[Explicit backgrounds and a faithful smooth trace measure]
\label{cor:r2-smooth-factor-backgrounds}
For every $0\leq p<n$, the series
\[
 \widehat S_p=\sum_{j,m\geq1}2^{-j-m}(Df_{j,m})^p
\]
defines a positive separately closed original current on $V$ that
satisfies \eqref{eq:r2-smooth-factor-condition}. So does any
background $S\geq\epsilon\widehat S_p$ with $\epsilon>0$.
Moreover,
\[
 \widehat\sigma=\sum_{j,m\geq1}2^{-j-m}(Df_{j,m})^n
\]
is a finite positive Radon measure of full support on $V$, formed from
original smooth top-form trace measures. For $u,v\in\Cb(V)$,
\[
 u\leq v\quad\widehat\sigma\text{-almost everywhere}
 \quad\Longleftrightarrow\quad u\leq v\quad\text{everywhere on }V.
\]
\end{corollary}

\begin{proof}
Uniform coefficient bounds and convergence of the positive series
follow from relative CLN, exactly as in
Proposition~\ref{prop:r2-faithful-backgrounds}, since all $f_{j,m}$
lie in $[0,1]$ on the fixed $V^+$. Removing one summand leaves a
positive separately closed current, which proves
\eqref{eq:r2-smooth-factor-condition} with $a_{j,m}=2^{-j-m}$.
The same argument applies after multiplication by $\epsilon$ and
addition of a positive background.

If $u\leq v$ almost everywhere for $\widehat\sigma$, the positive
coefficients in its defining sum imply
$(Df_{j,m})^n(\{u>v\})=0$ for every $j,m$. Apply the weighted limit
argument in the preceding proof with $\delta=0$. It gives pointwise
order on $\mathscr D$, hence everywhere. The reverse implication is
immediate.

Finally take a nonempty ordinary open $O\subset V$, choose
$z_j\in O\cap\mathscr D$, and choose a nonnegative original smooth
cutoff supported in $O$ and equal to one near $z_j$. Unweighted
monotone convergence gives
\[
 \int\chi(Df_{j,m})^n
     \longrightarrow 2^{-n}\int\chi(Dh_j)^n>0.
\]
Some summand assigns positive mass to $O$, and consequently
$\widehat\sigma(O)>0$. This proves full support without asserting
any atomic or non-atomic description of these smooth-factor products.
\end{proof}

The counterexample above and this construction are
\texttt{NEW\_PROVED} relative to the adopted inputs. The implication
under condition \eqref{eq:r2-smooth-factor-condition} is a
\texttt{DIRECT\_COROLLARY} of fixed-background weighted convergence
and the established Shilov detectors. These are sufficient conditions;
a necessary characterization of pointwise-faithful backgrounds is
not proved here.

% END INCLUDED SOURCE: base_r2/sections/gap_detection_addendum


% BEGIN INCLUDED SOURCE: base_r2/sections/weighted_l1
\section{A relative weighted $L^1$ Chern--Levine--Nirenberg estimate}
\label{sec:r2-weighted-l1}

The background is fixed throughout. Let $X$ be an allowed analytic
space over the prescribed finite extension $K/\mathbf Q_p$, of pure
dimension $n$, and let $S$ be a symmetric, original-positive,
separately $d'$- and $d''$-closed current of bidegree $(p,p)$.
Put $q=n-p$ and $D=d'd''$. All products, coefficient measures and
tests below belong to the original CLD calculus. No extension of $S$
to the additional test forms in CLD 2025 is used.

We prove a relative version of the weighted-form argument in
CLD 2025, Proposition 8.4.3. The passage to bounded classical
potentials uses the fixed-background calculus proved in R1. The
result is recorded as \texttt{NEW\_PROVED} in R2, with this method
and these dependencies stated explicitly; no claim of literature
priority is made.

\subsection{The positive cutoff forms}

Let $O$ be an ordinary open set and let $\rho$ be an original smooth
psh function on $O$, meaning that $A=D\rho$ is a positive
$(1,1)$-form. Suppose that
\[
 V=\{\rho<0\}\Subset O.
\]
Write $t=(-\rho)_+$ and $B=d'\rho\wedge d''\rho$. Fix $a=4$.
For $0\leq k\leq q$, put
\begin{equation}
\label{eq:r2-cutoff-forms}
 \beta_k=t^{k+a}A^{q-k}
       +(k+a)t^{k+a-1}B\wedge A^{q-k-1},
\end{equation}
where the second term is omitted for $k=q$. The forms are extended
by zero off $V$. They have compact support in $\overline V$, are
strongly positive, and have continuous coefficients. In fact their
coefficients are of class $C^2$ on each original moment presentation:
the smallest power occurring is $t^3$.

The coefficient-Radon property of $S$ defines $\beta_k\wedge S$ and
its products with smooth forms. Integration by parts with these
$C^2$ cutoff forms is legitimate. To check this within the original
test category, approximate their one-variable coefficient functions
in $C^2$ by smooth functions of $\rho$, and multiply by one fixed
original smooth cutoff equal to one near $\overline V$.
The approximants have one compact support. The coefficient measures
permit passage in every term involving at most two derivatives.
The approximating forms need not be positive; positivity of the
limiting forms follows directly from \eqref{eq:r2-cutoff-forms}.

On $V$, set $b=k+a$, $r=q-k$, and
\[
 c_k=\begin{cases}2,&k<q,\\1,&k=q.\end{cases}
\]
For $1\leq k\leq q$, direct differentiation gives
\begin{align}
 d''\beta_k
   &=-c_k b\,t^{b-1}d''\rho\wedge A^r,
       \label{eq:r2-dbeta}\\
 D\beta_k
   &=c_k b\{-t^{b-1}A^{r+1}
                 +(b-1)t^{b-2}B\wedge A^r\}.
       \label{eq:r2-ddbeta}
\end{align}
The endpoint $k=q$ is distinguished because $\beta_q=t^{q+a}$
has no second summand. The following proof uses both cases explicitly.

\begin{lemma}[Smooth relative weighted estimate]
\label{lem:r2-smooth-weighted-cln}
Assume $q\geq1$. Let $F\leq0$ and $u_1,\ldots,u_k$ be original
smooth psh functions on a neighborhood of $\overline V$, where
$0\leq k\leq q$. Put $M_i=\sup_V|u_i|$. Then
\begin{equation}
\label{eq:r2-smooth-l1}
 \int(-F)\,\beta_k\wedge Du_1\wedge\cdots\wedge Du_k\wedge S
 \leq C(k,a)\prod_{i=1}^kM_i
                 \int(-F)\,\beta_0\wedge S,
\end{equation}
where $C(0,a)=1$ and
\begin{equation}
\label{eq:r2-constant}
 C(k,a)=\prod_{i=1}^{k}
           \frac{4e(a+i)(a+i+1)}{a+i-1}.
\end{equation}
The constants do not depend on $S$, $F$, or the chosen functions.
\end{lemma}

\begin{proof}
All integrals have compact support. Write
$\Psi_{k-1}=Du_1\wedge\cdots\wedge Du_{k-1}$ and denote the
left-hand side of \eqref{eq:r2-smooth-l1} by $J_k$.
We first estimate the last slot when its function $u$ is strictly
positive on a neighborhood of $\overline V$. Twice integrating by
parts against the closed current $\Psi_{k-1}\wedge S$ gives
\begin{equation}
\label{eq:r2-integration-identity}
 J_k
 =\int\bigl\{u(FD\beta_k-DF\wedge\beta_k)
              +2F\,d'u\wedge d''\beta_k\bigr\}
                       \wedge\Psi_{k-1}\wedge S.
\end{equation}
Here and below cross terms can be replaced by their symmetric parts
because the remaining current is symmetric. For example, the two
cross terms in $D(F\beta_k)$ have equal pairings. Stokes applied to
$Fu\,d''\beta_k\wedge\Psi_{k-1}\wedge S$ yields
\eqref{eq:r2-integration-identity}, with the displayed signs.

Put $w=-F\geq0$ and
\[
 \Gamma(f,g)=\tfrac12(d'f\wedge d''g+d'g\wedge d''f),
 \qquad \Gamma(f)=d'f\wedge d''f.
\]
For every positive scalar $h$, the positive-square identity gives
\[
 2\Gamma(u,\rho)\leq h^{-1}\Gamma(u)+hB.
\]
Apply this pointwise on $V$ with $h=c_kbu/t$. Equations
\eqref{eq:r2-dbeta}--\eqref{eq:r2-integration-identity}, and
positivity of $uDF\wedge\beta_k$, give
\begin{align*}
 J_k
 &\leq\int wu\{-D\beta_k+c_k^2b^2t^{b-2}B\wedge A^r\}
                          \wedge\Psi_{k-1}\wedge S\\
 &\hspace{1cm}
       +\int w\,t^b u^{-1}\Gamma(u)\wedge A^r
                          \wedge\Psi_{k-1}\wedge S.
\end{align*}
The form in braces is
\[
 c_kb\,t^{b-1}A^{r+1}
  +c_kb\bigl((c_k-1)b+1\bigr)t^{b-2}B\wedge A^r.
\]
Since $c_k$ is either $1$ or $2$, it is bounded above by
\[
 C_k\beta_{k-1},\qquad
 C_k=\frac{2b(b+1)}{b-1}.
\]
Also $t^bA^r\leq\beta_k$. Wedging each displayed inequality with
the other strongly positive factors and then with the
original-positive current $S$ preserves it. We obtain
\begin{equation}
\label{eq:r2-positive-slot}
 J_k(F;u_1,\ldots,u_{k-1},u)
 \leq C_k\|u\|_{V}J_{k-1}
       +\int w\,\beta_k\wedge\Psi_{k-1}\wedge
                          u^{-1}\Gamma(u)\wedge S.
\end{equation}

Now let $u_k$ have arbitrary sign and set $M=\sup_V|u_k|$.
If $M=0$, its curvature vanishes on $V$ and there is nothing to
prove. Otherwise put $z=\exp(B_0u_k)$, where $B_0>0$. Then
\[
 Dz=z^{-1}\Gamma(z)+B_0z\,Du_k
       \geq z^{-1}\Gamma(z)+B_0e^{-B_0M}Du_k.
\]
Apply \eqref{eq:r2-positive-slot} to $z$. The energy term cancels
with the identical term in this lower bound for $Dz$, leaving
\[
 B_0e^{-B_0M}J_k\leq C_ke^{B_0M}J_{k-1}.
\]
Choose $B_0=(2M)^{-1}$. This gives
\[
 J_k\leq\frac{4e(a+k)(a+k+1)}{a+k-1}\,M_kJ_{k-1}.
\]
Induction from $J_0$ proves the assertion.
\end{proof}

\subsection{Bounded potentials and negative psh witnesses}

We now work in a common algebraic norm-box neighborhood of the kind
used in R1. The compact set $\overline V$ lies in a prescribed
smaller norm box, and all potentials are defined on a larger one.
Thus the compact-box regularization lemma of R1 gives smooth
decreasing classical psh approximants on a neighborhood of the
entire $\overline V$. This domain condition is only local geometry;
it places no restriction on the positive closed current $S$.

For a negative extended-real classical psh function
$\phi:O\to[-\infty,0]$, put
\[
 \phi_N=\max\{\phi,-N\}\in\Cb(O),
 \qquad
 A_{S,\rho}(\phi)=\int(-\phi)\,\beta_0\wedge S\in[0,\infty].
\]
This is an integral of the actual Borel function against one
specified positive Radon measure. It is not an ambient-volume
integrability condition.

\begin{theorem}[Uniform truncation estimate over an arbitrary background]
\label{thm:r2-weighted-l1}
Assume $q\geq1$ and $A_{S,\rho}(\phi)<\infty$.
Let $K\Subset V$ and let $u_1,\ldots,u_q\in\Cb(O)$ have
pointwise bounds $|u_i|\leq M_i$ on a fixed neighborhood of
$\overline V$. Then, with
$\delta_K=\min_K(-\rho)>0$,
\begin{equation}
\label{eq:r2-uniform-truncation}
 \sup_{N\geq1}\int_K(-\phi_N)
       T_{S,q}(u_1,\ldots,u_q)
 \leq \delta_K^{-(q+a)} C(q,a)
             \prod_{i=1}^qM_i\,A_{S,\rho}(\phi).
\end{equation}
In particular, $\phi$ is integrable on $K$ for every such mixed
measure, and the same bound holds with $-\phi$ in place of
$-\phi_N$.
\end{theorem}

\begin{proof}
First fix $N$ and $\epsilon>0$. The bounded classical psh function
$F=\phi_N-\epsilon$ is at most $-\epsilon$.
Choose a smooth decreasing row $F_j\downarrow F$ on one
neighborhood of $\overline V$. After a prefix is discarded,
$F_j\leq-\epsilon/2$ there. This follows from compactness applied
to the increasing open sets where the continuous majorants are
less than $-\epsilon/2$.

Simultaneously choose smooth decreasing classical psh rows
$u_{i,j}\downarrow u_i$ on the same neighborhood. Given
$\varepsilon>0$, a prefix makes $|u_{i,j}|\leq M_i+\varepsilon$
on $\overline V$. The lower bound follows from $u_{i,j}\geq u_i$;
the upper bound is the same compactness argument. All choices are
made for this fixed $N$. No regularization theorem for the
unbounded function $\phi$ is assumed.

Apply Lemma~\ref{lem:r2-smooth-weighted-cln} with $k=q$.
The R1 weighted decreasing theorem, in top degree against
continuous compact scalar tests, passes the left side to
\[
 \int(-\phi_N+\epsilon)\,\beta_q\,
                          T_{S,q}(u_1,\ldots,u_q).
\]
On the right, bounded convergence is against the \emph{fixed}
measure $\beta_0\wedge S$. Let $\varepsilon\downarrow0$ and
then $\epsilon\downarrow0$. We obtain
\[
 \int(-\phi_N)\,\beta_q\,
                          T_{S,q}(u_1,\ldots,u_q)
 \leq C(q,a)\prod_iM_i\int(-\phi_N)\,\beta_0\wedge S.
\]
Because $\beta_q=t^{q+a}\geq\delta_K^{q+a}$ on $K$ and
$-\phi_N\leq-\phi$, this is
\eqref{eq:r2-uniform-truncation}. Finally $-\phi_N\uparrow-\phi$.
Monotone convergence for each fixed mixed measure proves the
integrability assertion and its bound.
\end{proof}

\begin{corollary}[Uniformity over normalized potentials]
\label{cor:r2-normalized-witness}
Under the same hypotheses,
\[
 \sup_{N\geq1}\ \sup_{\substack{v\in\Cb(O)\\0\leq v\leq1}}
 \int_K(-\phi_N)\,T_{S,q}(v,\ldots,v)
 \leq \delta_K^{-(q+a)}C(q,a)A_{S,\rho}(\phi).
\]
The same assertion holds with different normalized potentials in
the differentiated slots.
\end{corollary}

\begin{corollary}[A precise relative polar conclusion]
\label{cor:r2-relative-polar}
For the capacity with the \emph{fixed admissibility domain} $O$,
\[
 \operatorname{Cap}_{S}(E;O)
   =\sup_{\substack{v\in\Cb(O)\\0\leq v\leq1}}
          \int_E T_{S,q}(v,\ldots,v),
\]
one has, for every $N>0$,
\[
 \operatorname{Cap}_{S}\bigl(K\cap\{\phi\leq-N\};O\bigr)
 \leq \frac{\delta_K^{-(q+a)}C(q,a)}{N}
                   A_{S,\rho}(\phi).
\]
Hence $K\cap\{\phi=-\infty\}$ has zero relative capacity and
zero mass for every bounded mixed measure in
Theorem~\ref{thm:r2-weighted-l1}.
\end{corollary}

\begin{proof}
On the indicated set, $-\phi\geq N$. Apply the integrability
bound in Theorem~\ref{thm:r2-weighted-l1} and take the supremum.
The pole set is contained in every such level set. Its zero mass
also follows directly from the proved integrability of $\phi$.
\end{proof}

\begin{corollary}[Outer capacity and quasicontinuity of an integrable witness]
\label{cor:r2-witness-outer-capacity}
Keep the hypotheses of Theorem~\ref{thm:r2-weighted-l1}, and let
$W$ be ordinary open with $\overline W\subset K\Subset V$.
Define outer capacity by
\[
 \operatorname{Cap}_S^*(E;O)
  =\inf_{\substack{G\subset O\text{ ordinary open}\\E\subset G}}
                  \operatorname{Cap}_S(G;O).
\]
Write $C_K=\delta_K^{-(q+a)}C(q,a)A_{S,\rho}(\phi)$.
Then, for $N>0$,
\begin{equation}
\label{eq:r2-open-sublevel-bound}
 \operatorname{Cap}_S(W\cap\{\phi<-N\};O)\leq C_K/N.
\end{equation}
Consequently $W\cap\{\phi=-\infty\}$ has outer capacity zero.
The truncations $\phi_N$ converge to $\phi$ locally in outer
capacity: for every $\delta>0$,
\[
 \operatorname{Cap}_S^*
       (W\cap\{\phi_N-\phi>\delta\};O)\leq C_K/N
                  \longrightarrow0,
\]
where the difference is $+\infty$ at a pole of $\phi$.
Finally, for every $\epsilon>0$ there is an ordinary open
$G\subset O$ with $\operatorname{Cap}_S(G;O)<\epsilon$ such
that $\phi|_{W\setminus G}$ is finite and continuous.
These are \texttt{DIRECT\_COROLLARY} assertions; no outer-regularity
theorem for arbitrary Borel sets is assumed.
\end{corollary}

\begin{proof}
The function $\phi$ is upper semicontinuous, so
$H_N=W\cap\{\phi<-N\}$ is ordinary open. It is contained in
$K\cap\{\phi\leq-N\}$, so
Corollary~\ref{cor:r2-relative-polar} proves
\eqref{eq:r2-open-sublevel-bound}. Each $H_N$ contains the pole
set in $W$, giving its zero outer capacity. Since $\phi_N=\phi$
off $\{\phi<-N\}$, the same ordinary open set covers every
displayed truncation exceptional set.

For the last assertion choose $N$ with $C_K/N<\epsilon/2$.
The bounded classical function $\phi_N$ is quasicontinuous by R1.
Precisely, cover $\overline W$ by finitely many smaller norm-box
cores in $O$ and apply the bounded local quasicontinuity theorem
on each larger box with a total error less than $\epsilon/2$.
Restriction of admissible functions gives
$\operatorname{Cap}_S(E;O)\leq\operatorname{Cap}_S(E;\Omega)$
for sets in any of these boxes $\Omega\subset O$. Thus the union
$G_N\subset O$ of their ordinary open exceptional sets has
$\operatorname{Cap}_S(G_N;O)<\epsilon/2$, and $\phi_N$ is
continuous on $W\setminus G_N$. Put $G=G_N\cup H_N$.
Countable subadditivity bounds its capacity by $\epsilon$.
On $W\setminus G$ one has $\phi=\phi_N\geq-N$, proving
the assertion in the subspace topology.
\end{proof}

\subsection{Intrinsic local hypotheses and scope}

One sufficient hypothesis, independent of a particular barrier, is
local coefficient integrability of $\phi$ with respect to $S$:
\begin{equation}
\label{eq:r2-coefficient-integrability}
 \int|\phi|\,d\mu_{S,\Theta}<\infty
 \quad\hbox{for every positive original compact smooth
 $(q,q)$-test $\Theta$.}
\end{equation}
The finite positive controls of R imply that this is equivalent
to integrability against all local coefficient variations.
They also dominate the continuous-coefficient compact form
$\beta_0$ by finitely many positive smooth controls. Therefore
\eqref{eq:r2-coefficient-integrability} implies
$A_{S,\rho}(\phi)<\infty$ for every barrier used above.

The R1 collar construction supplies the required barriers locally:
take a smooth psh function $b$ which is zero on the inner norm box
and large on an outer collar, and put $\rho=b-c$ with $c>0$
between the two levels. Then $\{\rho<0\}\Subset O$ and contains
the chosen core. No strict positivity of $D\rho$ is required.
A finite covering gives the stated integrability and mass bounds
on arbitrary compact subsets for which
\eqref{eq:r2-coefficient-integrability} holds.

Lower-degree coefficient estimates follow by inserting bounded
smooth quadratic moment potentials into the remaining slots and
using R's finite-test domination. Thus the same hypothesis makes
$\phi T_{S,k}(u_1,\ldots,u_k)$ a coefficient-Radon current for
every bounded tuple and every $k\leq q$.

For $q=0$, the asserted uniform truncation bound is simply
integrability of $\phi$ against the fixed top-degree measure $S$;
no differentiated slot is present. Finally, the capacity conclusion
keeps its admissibility domain $O$. To obtain the corresponding
statement on a smaller domain, choose its collar and verify its
integrability condition there; monotonicity under domain restriction
does not transfer the displayed upper bound automatically.

The hypothesis on $\phi$ and $S$ is essential to the scope of this
result. If $S=[Y]$ and $\phi$ is identically $-\infty$ on $Y$,
it is generally not satisfied. Nothing here asserts that arbitrary
relative products avoid every ambient pluripolar set. Nor does the
uniform bound alone establish weighted convergence when both an
unbounded witness and its background measure vary.


% END INCLUDED SOURCE: base_r2/sections/weighted_l1


% BEGIN INCLUDED SOURCE: base_r2/sections/polar_restriction
\section{Restriction of a positive closed current to a complete pole set}
\label{sec:r2-polar-restriction}

Let $S$ be a symmetric original-positive, separately $d'$- and
$d''$-closed current of bidegree $(p,p)$ on an allowed ordinary
open $U$, and put $q=n-p$. Let $\phi:U\to[-\infty,\infty)$ be
an extended-real classical psh function, and set
\[
 P=\{\phi=-\infty\}.
\]
All restrictions in this section mean Borel multiplication of the
actual coefficient-Radon measures. In particular, they are not
restrictions to a tropical image. The proof uses the bounded
capacitary closure developed in this record and the original CLD
positive cone of R throughout.

\subsection{A first-order estimate on a fixed test presentation}

For a bounded classical psh function $a$, define the distribution
\begin{equation}
\label{eq:r2-relative-gamma-current}
 \Gamma_S(a)=\tfrac12D(a^2S)-aD(aS).
\end{equation}
It is a positive coefficient-Radon current of bidegree $(p+1,p+1)$.
Indeed, locally add a constant so that $a\geq0$ and use a decreasing
smooth classical psh row. The square of each term is classical psh
and decreases to $a^2$. The first-round weighted theorem passes
\eqref{eq:r2-relative-gamma-current} to the limit of the positive
currents $d'a_j\wedge d''a_j\wedge S$. Adding the constant does
not change \eqref{eq:r2-relative-gamma-current}, since $S$ is
separately closed. Positivity and R's coefficient-Radon theorem
give the assertion. Wedging with a fixed closed smooth trace form
recovers the first-round relative energy measures.

\begin{lemma}[First-order relative Cauchy--Schwarz estimate]
\label{lem:r2-first-order-CS}
Suppose $q\geq1$. Localize an original compact smooth test
$\eta$ of bidegree $(q-1,q)$ to one finite moment presentation
with coordinates $t_1,\ldots,t_r$. Let
\[
 \beta=\sum_{a=1}^r d't_a\wedge d''t_a,
\]
and choose a nonnegative compact original smooth cutoff $\chi$
equal to one near $\operatorname{supp}\eta$. There is a constant
$C_\eta<\infty$, independent of $S$ and of the bounded classical
psh function $a$, such that
\begin{equation}
\label{eq:r2-first-order-CS}
 \bigl|\langle d'(aS),\eta\rangle\bigr|^2
 \leq C_\eta
       \langle\Gamma_S(a),\chi\beta^{q-1}\rangle
       \langle S,\chi\beta^q\rangle.
\end{equation}
The same assertion holds for $d''$ and tests of bidegree $(q,q-1)$.
The constant remains valid when the moment presentation is refined
to include additional coordinates used by $a$ or its approximants.
\end{lemma}

\begin{proof}
First let $a$ be smooth and write $\xi=d'a$.
The finite expansion of $\eta$ has terms
$f_{I,J}d't_I\wedge d''t_J$, where $|I|=q-1$ and $|J|=q$.
Under the standard symmetric-matrix convention for original
$(q,q)$-forms, the symmetric part of $\xi\wedge\eta$ is a
finite sum, with fixed signs, of the symmetric tensors between
\[
 A_I=\xi\wedge d't_I,\qquad B_J=d't_J.
\]
For two arbitrary $q$-covectors $A,B$, the square
$(\sqrt\epsilon A\pm\epsilon^{-1/2}B)^2$ is in the original
positive cone. The associated matrix Young inequality therefore
bounds either sign of their symmetric cross tensor by
$\epsilon A^2+\epsilon^{-1}B^2$, up to a fixed harmless factor.
Apply it to the finitely many terms and use the fixed bounds of
the coefficients $f_{I,J}$ on the test support. Summing the
$A_I^2$ terms gives a fixed multiple of
$d'a\wedge d''a\wedge\beta^{q-1}$; summing the $B_J^2$ terms
is bounded by a fixed multiple of $\beta^q$.
After rescaling $\epsilon$, we obtain, for every $\epsilon>0$,
\[
 \pm\operatorname{Sym}(d'a\wedge\eta)
 \leq \epsilon\chi\,d'a\wedge d''a\wedge\beta^{q-1}
             +\frac{C'\!_\eta}{\epsilon}\chi\beta^q.
\]
Outside the test support the left side is zero. The cutoff makes
the right side a compact positive test. No strict positivity in
additional analytic directions is used: $A_I$ contains the full
covector $d'a$, and only its remaining $q-1$ factors come from
the fixed frame. The same matrix-square argument works on every
common refinement, with the same finite index sets and constants.

The current $S$ annihilates the antisymmetric component, by the
original-positive coefficient argument in R. Evaluation and
separate closedness give
\[
 |\langle d'(aS),\eta\rangle|
 \leq\epsilon E+\epsilon^{-1}C'\!_\eta M,
 \quad
 E=\langle\Gamma_S(a),\chi\beta^{q-1}\rangle,
 \quad M=\langle S,\chi\beta^q\rangle.
\]
Optimizing proves \eqref{eq:r2-first-order-CS}; if $E$ or $M$
is zero, take a limit in $\epsilon$ to obtain the same conclusion.

For bounded classical $a$, choose a smooth decreasing row on
a fixed neighborhood of the compact test stage, as in R1.
Multiplication against the fixed coefficient measures of $S$
gives $a_jS\to aS$. The weighted first-round theorem gives
$\Gamma_S(a_j)\to\Gamma_S(a)$, as explained above.
Pass to the limit in the smooth inequality. The involution $J$,
or the identical calculation with primed factors interchanged,
proves the $d''$ assertion. A finite partition handles every
original compact test, not only those in one initially fixed chart.
\end{proof}

\subsection{The idempotent identity}

\begin{lemma}[Cubic chain identity in the bounded capacitary closure]
\label{lem:r2-cubic-chain}
Let $0\leq a\leq M$ be a member of $\mathcal A_{S,b}$, locally
with a common bound for its defining approximants. Then
\begin{equation}
\label{eq:r2-cubic-chain}
 D(a^3S)=3aD(a^2S)-3a^2D(aS).
\end{equation}
The square and cube belong to the same local capacitary closure.
Moreover \eqref{eq:r2-relative-gamma-current} is a positive
coefficient-Radon current, and it is continuous under bounded
outer-capacity convergence of classical approximants.
\end{lemma}

\begin{proof}
For a smooth function, the usual product rule gives
\[
 D(a^2S)=2aD(aS)+2d'a\wedge d''a\wedge S,
\]
\[
 D(a^3S)=3a^2D(aS)+6a\,d'a\wedge d''a\wedge S,
\]
which imply \eqref{eq:r2-cubic-chain}.
For bounded nonnegative classical $a$, choose decreasing smooth
nonnegative classical approximants. Their squares and cubes are
again classical psh and decrease to the corresponding powers.
Pass the two weighted one-slot terms in
\eqref{eq:r2-cubic-chain} by the first-round weighted theorem;
pass the left side distributionally against the fixed $S$.

For a nonnegative member of the capacitary closure, replace its
classical defining approximants by their maxima with zero. The
approximants remain uniformly bounded and converge in outer
capacity. Squaring and cubing are Lipschitz on that bounded
range, so their powers also converge in outer capacity. They
are classical psh by the preceding paragraph and decreasing
closure, hence define $a^2,a^3$ in the capacitary class.
The completed weighted theorem now passes every term of the
classical identity. The same argument passes
$\tfrac12D(a_j^2S)-a_jD(a_jS)$ to
\eqref{eq:r2-relative-gamma-current}. Positivity passes from the
classical energy currents, and R gives coefficient-Radon type.
\end{proof}

\begin{theorem}[Positive closed restriction to a complete pole set]
\label{thm:r2-polar-restriction}
For the arbitrary background $S$ and the complete pole set $P$
specified above, both currents
\[
 S_P=\mathbf1_P S,\qquad S_{P^c}=\mathbf1_{U\setminus P}S
\]
are symmetric, original-positive, and separately $d'$- and
$d''$-closed. For every $0\leq k\leq q$ and
$u_1,\ldots,u_k\in\Cb(U)$, one has
\begin{equation}
\label{eq:r2-polar-restriction-BT}
 \mathbf1_P T_{S,k}(u_1,\ldots,u_k)
       =T_{S_P,k}(u_1,\ldots,u_k),
\end{equation}
and the analogous identity for $P^c$, on every original compact
smooth complementary test form.
\end{theorem}

\noindent\textbf{Status: \texttt{NEW\_PROVED}, relative to R1 and the
proved R2 capacitary calculus.}
This is a restriction theorem for a current already defined on
$U$. It is not an extension theorem for a current initially
defined only on $U\setminus P$.

\begin{proof}
Positivity, symmetry and coefficient-Radon type of both Borel
restrictions are immediate. In top degree $q=0$, separate
closedness is automatic and there are no differentiated slots.
Suppose therefore $q\geq1$.

Work locally and subtract a constant from $\phi$ so that
$\phi\leq0$ on the neighborhood. Set
\[
 w_N=1+\frac{\max\{\phi,-N\}}{N},\qquad N\geq1.
\]
Then $0\leq w_N\leq1$, the functions are classical psh, and
they increase pointwise to the actual Borel function
$w=\mathbf1_{P^c}$. The local increasing-convergence theorem
puts $w$ in $\mathcal A_{S,b}$ and gives outer-capacity
convergence of $w_N$ to $w$. Its square and cube equal itself.
The cubic identity therefore gives
\[
 D(wS)=3wD(wS)-3wD(wS)=0,
 \qquad
 \Gamma_S(w)=\tfrac12D(wS)-wD(wS)=0.
\]
Completed weighted convergence gives
\[
 w_NS\longrightarrow wS,
 \qquad \Gamma_S(w_N)\longrightarrow0.
\]
Fix an arbitrary first-order original compact test and localize
it to the presentations of Lemma~\ref{lem:r2-first-order-CS}.
The fixed factor $\langle S,\chi\beta^q\rangle$ is finite,
whereas the energy factor tends to zero. The lemma implies
$d'(w_NS)\to0$ on that test. Distributional continuity also
gives $d'(w_NS)\to d'(wS)$, so $d'(wS)=0$.
The same proof gives $d''(wS)=0$. Subtracting from the separately
closed $S$ proves closedness of $S_P$. These local conclusions
glue because the Borel restrictions have already been defined
globally.

For \eqref{eq:r2-polar-restriction-BT}, first take $k=1$ and
a smooth classical psh function $u$. Separate closedness just
proved and smooth Borel multiplication give
\[
 D(uS_P)=Du\wedge S_P
              =\mathbf1_P(Du\wedge S)=\mathbf1_PD(uS).
\]
For general $u\in\Cb$, use a decreasing smooth classical row
$u_j\downarrow u$ on a fixed compact norm-box stage. It has
one common bound on the slightly larger domain. The left side
converges by the fixed-background first-round theorem for $S_P$.
The currents $D(u_jS)$ converge to $D(uS)$ and satisfy uniform
coefficient bounds by the top relative capacity of that same
domain. The multiplier $\mathbf1_P=1-w$ is quasicontinuous for
this capacity because $w\in\mathcal A_{S,b}$. The R2
quasicontinuous multiplier theorem therefore passes the right
side to $\mathbf1_PD(uS)$. This proves the one-slot identity
on all original tests without asserting stability for arbitrary
varying backgrounds.

For several slots, apply the one-slot identity to the arbitrary
positive closed background
$Q=T_{S,k-1}(u_2,\ldots,u_k)$. The restriction theorem already
applies to this $Q$ and the same classical pole set. Induction
gives
\[
 \mathbf1_PQ=T_{S_P,k-1}(u_2,\ldots,u_k),
\]
and hence
\[
 \mathbf1_PD(u_1Q)=D(u_1\mathbf1_PQ)
                         =T_{S_P,k}(u_1,\ldots,u_k).
\]
The complement follows in the same way, or by finite linearity
in the decomposition $S=S_P+S_{P^c}$.
\end{proof}

\begin{corollary}[Exact preservation of pole-set mass]
\label{cor:r2-pole-mass-preservation}
If $\mathbf1_PS=0$, then
$\mathbf1_PT_{S,k}(u_1,\ldots,u_k)=0$ for every bounded
classical tuple. This conclusion needs no integrability estimate
for the original witness. Conversely, the theorem allows nonzero
$S_P$; it does not assert that arbitrary relative products avoid
ambient pole sets.
\end{corollary}

\begin{proof}
Apply \eqref{eq:r2-polar-restriction-BT} with $S_P=0$.
\end{proof}

\paragraph{The ambient vanishing needs a separate argument.}
For $S=[1]$, the theorem shows that $\mathbf1_P[1]$ is a
positive separately closed original current of degree zero.
This calculation alone does not imply its vanishing. Under the
connectedness and nontriviality assumptions on the witness,
Theorem~\ref{thm:r2-ambient-polar-traces} proves that vanishing
by an independent argument using divisorial points and bounded
model approximations. No distributional elliptic regularity is invoked.
CLD 2025, Remark 7.8.3, explicitly distinguishes the equation
$D[u]=0$ from a continuity theorem on its general tropical
spaces, and Proposition 8.1.12 retains local integrability as
an extra condition in its psh class. The proof in the cited
ambient trace theorem supplies the additional analytic input,
rather than silently interpreting a closed degree-zero density
as a continuous pluriharmonic function.

% END INCLUDED SOURCE: base_r2/sections/polar_restriction


% BEGIN INCLUDED SOURCE: base_r2/sections/divisorial_poles
\section{Divisorial points are not poles of a nontrivial classical psh function}
\label{sec:r2-divisorial-poles}

The field, ordinary topology, and all-extensions classical definition
remain fixed. The geometric lemma below uses strict affinoid domains
inside the algebraic neighborhoods supplied by \textup{B}. Its
normalization step is essential: a formal model with only one
irreducible special-fiber component need not have a unique Shilov
point unless the corresponding generic-point fiber is controlled.

\begin{lemma}[A domain with one prescribed Shilov point]
\label{lem:r2-one-shilov-domain}
Let $U$ be an ordinary open in an allowed good, Hausdorff,
boundaryless, quasi-smooth, pure $n$-dimensional $K$-analytic space,
with $n\geq1$. If $z\in U$ is divisorial, there exists a strict
affinoid domain $A\subset U$ such that
\[
 \Gamma(A)=\{z\},\qquad \operatorname{Int}_U(A)\ne\varnothing.
\]
The lemma does not assert that $z$ itself belongs to this ordinary
interior. In fact it belongs to the boundary of $A$.
\end{lemma}

\begin{proof}
Work in a strict algebraic neighborhood of $z$, as permitted by
\textup{B}, and choose a strict affinoid neighborhood $C\subset U$
of $z$. Divisoriality is unchanged on passing to $C$, by
\cite[Corollary~A.6, p.~367]{GM19}. By the definition of a divisorial
point, equivalently \cite[Proposition~A.5, pp.~366--367]{GM19}, there
is an admissible formal $K^\circ$-model $\mathfrak C$ of $C$ such
that the reduction of $z$ is the generic point of a special-fiber
component.

We may and do make this formal model normal. Indeed, $C$ is normal
because it is quasi-smooth. The formal affine rings are topologically
of finite type over the excellent complete discrete valuation ring
$K^\circ$. Their normalizations in their reduced generic fibers are
finite, and glue to a normal admissible formal model. The generic
fiber remains $C$, since it is already normal. The reduction of $z$
on this new model is still generic in the special fiber: it lies
over the former generic point under a finite integral map, and a
prime above a minimal prime over the uniformizer is again minimal
over that uniformizer, by incomparability. Denote the new generic
point by $\eta$ and its component by $E$.

Remove the other irreducible components of the special fiber and
choose a formal affine open neighborhood $\mathfrak A$ of $\eta$
inside the remaining formal open. Such a choice is possible because
the formal model is locally noetherian and $\eta$ lies on no other
component; one first works in an affine neighborhood and then avoids
its finitely many other components. The special fiber of
$\mathfrak A$ has exactly one irreducible component, an open part
of $E$. Set $A=\mathfrak A_\eta$, its generic fiber. Formal open
immersions induce analytic-domain inclusions on generic fibers, so
$A$ is a strict affinoid domain of $C$, hence of $U$, and it contains
$z$.

The model $\mathfrak A$ is normal and its generic fiber is reduced.
The uniqueness of the point over its special-fiber generic point is
the normal-model assertion in
\cite[Proposition~6.9.7, pp.~119--120]{CLD12}. It can also be read
directly at the height-one local ring of the normal model: the
resulting discrete valuation, normalized by the prescribed absolute
value of the uniformizer, determines that point. By
\cite[Proposition~A.3, p.~365]{GM19}, the points reducing to generic
points of this formal special fiber are exactly the Shilov boundary
of $A$. Hence $\Gamma(A)=\{z\}$.

We check the ordinary-interior assertion separately. The special
fiber of $\mathfrak A$ has a closed point $\xi$, and the reduction
map is surjective, so some $y\in A$ reduces to $\xi$. The reduced
closure of $\xi$ is a finite scheme over $\widetilde K$, hence is
proper. The formal reduction criterion
\cite[Lemma~6.5.1, p.~105]{CLD12} therefore gives
$y\in\operatorname{Int}(A/K)$. Equivalently, one can use the
analytic Nullstellensatz and the fact that rigid points of a strict
analytic space over a nontrivially valued field do not belong to its
analytic boundary; see \cite[Remark~4.3.9, p.~67]{CLD12}.

For the analytic-domain inclusion $A\hookrightarrow U$, relative
interior is ordinary topological interior. Since $U$ is boundaryless
and locally separated, the composition formula gives
\[
 \operatorname{Int}_U(A)
   =\operatorname{Int}(A/U)=\operatorname{Int}(A/K);
\]
these are \cite[\S2.11(3)--(4), pp.~7--8]{GR25}. Thus the ordinary
interior is nonempty. Finally $A$ is pure of dimension $n$, and its
Shilov points belong to its analytic boundary by
\cite[Lemma~6.2, p.~31]{GR25}. The same identification of interiors
shows that $z\notin\operatorname{Int}_U(A)$, as asserted.
\end{proof}

The geometric lemma is a \texttt{DIRECT\_COROLLARY} of the cited
formal-model and boundary facts, with the needed normalization
spelled out. It makes no field extension. Rigid or reduction points
with a finite residue-field extension, when used in the proof, are
points of the original $K$-analytic space.

\begin{remark}[Why one component alone is insufficient]
Over the same discrete valuation ring with uniformizer $\pi$, the
formal scheme
\[
 \operatorname{Spf}\bigl(K^\circ\langle T\rangle/(T^2-\pi^2)\bigr)
\]
has only one irreducible special-fiber component, but its generic
fiber consists of two rigid points, $T=\pi$ and $T=-\pi$, and its
Shilov boundary consists of both. The characteristic of $K$ is zero,
so these points are distinct even when the residue characteristic is
two. This formal scheme is not normal. Thus replacing the
normalization paragraph by a claim that ``one component implies one
Shilov point'' would be an invalid argument.
\end{remark}

\begin{proposition}[A nontrivial pole set has empty ordinary interior]
\label{prop:r2-no-interior-poles}
Let $U$ satisfy the allowed analytic hypotheses and be connected.
Let $\phi:U\to\mathbf R\cup\{-\infty\}$ be classical psh and not
identically $-\infty$. Then
\[
 P=\{\phi=-\infty\}
\]
has empty ordinary interior in $U$.
\end{proposition}

\begin{proof}
For integers $N\geq1$ define
\[
 h_N=\frac{\max\{\phi,-N\}}{N}.
\]
These are members of $\Cb(U)$ and are bounded below by $-1$.
Upper semicontinuity of $\phi$, whose values do not include $+\infty$,
gives a finite upper bound $M$ on a neighborhood of each point. On
such a neighborhood,
\[
 h_N\leq \frac{\max\{M,0\}}{N}.
\]
In particular the family is locally uniformly bounded above, although
it need not be increasing. Apply the bounded upper-envelope theorem,
Theorem~\ref{thm:r2-local-upper-envelope}, to each tail and put
\[
 g_j=\left(\sup_{N\geq j}h_N\right)^*\in\Cb(U).
\]
The star is ordinary upper-semicontinuous regularization. The tail
envelopes decrease, and on the same neighborhood the estimates give
\[
 -1\leq g_j\leq\frac{\max\{M,0\}}j.
\]
Consequently $g=\lim_j g_j$ is in $\Cb(U)$ by decreasing closure
of the classical class, and $-1\leq g\leq0$ everywhere. At every
point where $\phi$ is finite, $h_N$ tends to zero, so every tail
supremum, and hence every $g_j$, is at least zero there. Thus $g=0$
at that point. Since $\phi$ is not identically $-\infty$, $g$
attains its maximum zero somewhere on $U$.

The set $\{g=0\}$ is closed by upper semicontinuity and open by
the local maximum principle \cite[Theorem~6.1, pp.~30--31]{GR25},
because all points of $U$ are interior points. It is nonempty, so
connectedness implies $g\equiv0$. If $P$ contained a nonempty
ordinary open $O$, every $h_N$ and every tail supremum would equal
$-1$ on $O$. Locality of ordinary upper regularization would give
$g_j=-1$ on $O$ for all $j$, hence $g=-1$ there, a contradiction.
\end{proof}

\begin{theorem}[No divisorial poles]
\label{thm:r2-no-divisorial-poles}
Under the assumptions of Proposition~\ref{prop:r2-no-interior-poles},
the value $\phi(z)$ is finite at every divisorial point $z\in U$.
\end{theorem}

\begin{proof}
Suppose instead that $\phi(z)=-\infty$ at such a point. Choose
$A\subset U$ from Lemma~\ref{lem:r2-one-shilov-domain}. For every
$N\geq1$, the finite classical truncation
$\phi_N=\max\{\phi,-N\}$ is bounded on the compact domain $A$.
It takes its maximum on $A$ at a Shilov
point, by \cite[Theorem~6.4, pp.~33--34]{GR25}. Since the unique
Shilov point is $z$ and $\phi_N(z)=-N$, it follows that
$\phi_N\equiv-N$ on $A$. Letting $N\to\infty$ gives
$\phi\equiv-\infty$ on $A$. Its nonempty ordinary interior is
therefore contained in $P$, contrary to
Proposition~\ref{prop:r2-no-interior-poles}.
\end{proof}

Both pole statements are \texttt{DIRECT\_COROLLARY} of the bounded
upper-envelope theorem already obtained from the accepted internal
inputs, the external maximum principles, and the geometric lemma.
They do not use any capacity-zero theorem for pole sets, a trace
nonconcentration assertion, or a relative increasing-continuity
statement. Thus they may be used as inputs to those subsequent
questions without a capacity argument being fed back into itself.

The conclusion concerns divisorial points in the ambient dimension.
For example, a proper analytic subvariety can still be a pole set and
can carry an integration-current background. Divisorial points of
that smaller subvariety are not ambient divisorial points. This
theorem therefore does not assert nonconcentration on ambient polar
sets for arbitrary positive closed backgrounds.

% END INCLUDED SOURCE: base_r2/sections/divisorial_poles


% BEGIN INCLUDED SOURCE: base_r2/sections/ambient_polar_traces
\section{Ambient pole sets and all original trace measures}
\label{sec:r2-ambient-polar-traces}

This section combines the absence of divisorial poles from
Theorem~\ref{thm:r2-no-divisorial-poles} with the proved local
increasing-capacity theorem. Density of divisorial points alone
would not justify the conclusion. The passage from model atoms
to a smooth trace uses the quasicontinuity of the actual pole
indicator and a uniform bound on one fixed capacity domain.

\begin{theorem}[Vanishing of ambient smooth traces on a classical pole set]
\label{thm:r2-ambient-polar-traces}
Let $U$ be a connected allowed ordinary open and let
$\phi:U\to[-\infty,\infty)$ be classical psh and not identically
$-\infty$. Put $P=\{\phi=-\infty\}$. For every original compact
smooth $(n,n)$-test $\eta$, one has
\[
 |\mu_{[1],\eta}|(P)=0.
\]
Equivalently, $\mathbf1_P[1]=0$ as an original coefficient-Radon
current. No local $L^1$ assertion about the original function
$\phi$ is assumed in this theorem.
\end{theorem}

\noindent\textbf{Status: \texttt{NEW\_PROVED}, relative to the
stated internal inputs and the preceding R2 theorems.}

\begin{proof}
Use a finite partition of the test and ordinary algebraic
localization from B. Work on one moment presentation, with unit
coordinates $t_1,\ldots,t_r$ defined on a fixed larger compact
norm-box neighborhood. Choose a smaller norm-box domain
$\Omega\Subset U$ containing the localized test support, on
which the coordinates are bounded. Set
\[
 \rho=\tfrac12\sum_{a=1}^r t_a^2,
 \qquad \beta=D\rho=\sum_{a=1}^r d't_a\wedge d''t_a.
\]
The function $\rho$ is smooth and classical psh. Indeed its
pullback to every curve after every complete field extension is
a convex function of the harmonic logarithms of analytic units;
the curve chain rule gives nonnegative Laplacian.
The finite matrix domination of R supplies a cutoff $\chi$ and
a finite constant $C_\eta$ such that
\begin{equation}
\label{eq:r2-ambient-trace-domination}
 |\mu_{[1],\eta}|\leq C_\eta\,\chi\beta^n\wedge[1].
\end{equation}
Any antisymmetric component is annihilated, as in R. It therefore
suffices to show that the positive smooth trace
$\mu=\beta^n\wedge[1]$ vanishes on $P$ near the test support.

Apply the common local projective globalization from B, Lemma
D.2, to $\rho$, and then its projective pure-Fubini--Study
regularization theorem. More explicitly, this is the constructive
globalization part of the first-round
Lemma~\ref{lem:compact-box-regularization}, before its final
smoothing step; it applies on the entire prescribed smaller
norm box. In the fixed line frame, after removing
the common additive constant, this gives
\[
 \rho_j\downarrow\rho
\]
on a neighborhood of the entire prescribed compact box. The
functions $\rho_j$ are restrictions of rational pure-Fubini--Study
metric weights on one projective variety and one identified ample
line. They are bounded uniformly on the fixed $\Omega$: the
first row and the continuous limit bound the whole decreasing
row on its compact closure. Choose $M>0$ with
$|\rho_j|\leq M$ there for every $j$.

The positive measures
\[
 \mu_j=T_{[1],n}(\rho_j,\ldots,\rho_j)
\]
are restrictions of the corresponding calibrated model measures.
Each is a finite sum of atoms at divisorial points. Ordinary
algebraic localization, including the possible finite constant
field in B, Lemma D.5, preserves this assertion; the calibration
may acquire its stated fixed degree factor, which does not affect
the support or zero-mass conclusion. By
Theorem~\ref{thm:r2-no-divisorial-poles}, none of these points
belongs to $P$. Consequently
\[
 \mathbf1_P\mu_j=0.
\]
The fixed-background decreasing theorem gives
$\mu_j\to\mu$ on all original scalar tests, and hence weakly
on compact continuous scalar tests.

These weak limits have a uniform bound on the \emph{same}
admissibility domain. In fact
$v_j=(\rho_j+M)/(2M)$ belongs to $\Cb(\Omega)$ and is between
zero and one, so for every Borel $E\subset\Omega$,
\[
 \mu_j(E)\leq(2M)^n\operatorname{Cap}_{[1]}(E;\Omega).
\]
Subtract a constant from $\phi$ so that it is nonpositive on
a neighborhood of $\overline\Omega$, and use
\[
 w_N=1+N^{-1}\max\{\phi,-N\}\uparrow\mathbf1_{P^c}.
\]
The local increasing theorem gives quasicontinuity of this actual
limit for the same capacity. Thus $\mathbf1_P$ is a bounded
quasicontinuous multiplier. The R2 weak-closure multiplier theorem
applies to $\mu_j\to\mu$ and yields
\[
 \mathbf1_P\mu
       =\lim_j\mathbf1_P\mu_j=0
\]
on the compact core. This passage, rather than density of the
model points, is the essential step. Equation
\eqref{eq:r2-ambient-trace-domination} now proves the assertion
for the localized test. The finite partition proves it for every
original compact smooth $(n,n)$-form.
\end{proof}

\begin{corollary}[Ambient bounded BT products avoid classical pole sets]
\label{cor:r2-ambient-bt-poles}
Under the same hypotheses, for every bounded classical tuple
and every $0\leq k\leq n$,
\[
 \mathbf1_P T_{[1],k}(u_1,\ldots,u_k)=0.
\]
Every coefficient measure of these currents has total variation
zero on $P$. In particular every bounded mixed Monge--Amp\`ere
measure gives $P$ zero mass.
\end{corollary}

\begin{proof}
The theorem gives $\mathbf1_P[1]=0$. Apply the exact restriction
identity of Theorem~\ref{thm:r2-polar-restriction}. Equality of
currents and uniqueness of their coefficient-Radon measures give
the coefficient assertion.
\end{proof}

\begin{corollary}[Local outer-capacity zero of ambient pluripolar sets]
\label{cor:r2-ambient-polar-outer-zero}
For every permitted norm-box domain $\Omega\Subset U$ and
compact $H\Subset\Omega$,
\[
 \operatorname{Cap}_{[1]}^*(H\cap P;\Omega)=0.
\]
More generally, every locally pluripolar subset has local outer
capacity zero, where locally pluripolar means locally contained
in the pole set of a classical psh function not identically
$-\infty$ on its connected witness domain. This statement does
not itself assert the converse. The converse and qualified pole
patching on exhaustion domains are proved later in
Section~\ref{sec:r2-polar-converse-choquet}.
\end{corollary}

\begin{proof}
Every normalized bounded potential on $\Omega$ gives a measure
that vanishes on $P$, by the preceding corollary applied locally.
Thus the Borel capacity of $W\cap P$ is zero for an ordinary
open core $H\subset W\Subset\Omega$.
The pole indicator is quasicontinuous on $W$ by the increasing
argument used above. Hence
\[
 W\cap P=W\cap\{\mathbf1_{P^c}<1/2\}
\]
is quasi-open. The proved outer regularity of quasi-open sets
gives outer capacity zero, and monotonicity gives the assertion
on $H$. For a locally pluripolar set, cover its intersection
with a prescribed compact core by smaller open witness cores
whose larger witness domains lie in $\Omega$. Second
countability in the fixed finite $p$-adic setting gives a
countable subcover. Apply the local conclusion in each and
use countable subadditivity together with the correct direction
of capacity domain monotonicity.
\end{proof}

\begin{corollary}[Quasicontinuity of extended classical psh functions]
\label{cor:r2-extended-psh-qc}
Every extended-real classical psh function that is not identically
$-\infty$ on any connected component is locally finite and
continuous outside ordinary open sets of arbitrarily small
ambient outer BT capacity. Its lower truncations converge
locally in outer capacity. No $L^1$ bound on the unattenuated
witness is asserted by this conclusion.
\end{corollary}

\begin{proof}
Work on a compact core, subtract a local upper bound, and thus
write the function as a negative witness $\phi$. Choose the
collar trace $\nu=\beta_0\wedge[1]$ of the weighted theorem.
The ambient trace theorem and finite smooth domination of this
continuous-coefficient trace give $\nu(P)=0$.
Lemma~\ref{lem:r2-attenuation} supplies
\[
 \psi=-\vartheta(-\phi),\qquad
 \vartheta(t)=\sum_{j\geq1}\min\{t/L_j,1\},
 \qquad \int(-\psi)\,d\nu<\infty.
\]
The function $\vartheta$ is increasing and tends to infinity:
for any integer $m$, all of its first $m$ summands equal one
once $t\geq\max_{j\leq m}L_j$.
For every $N>0$,
\[
 \{\phi<-N\}\subset\{\psi\leq-\vartheta(N)\}.
\]
The weighted theorem applied to $\psi$ therefore bounds the
capacity of the ordinary open sublevel set on the core by
$C/\vartheta(N)$, which tends to zero. The original truncation
$\max\{\phi,-N\}$ agrees with $\phi$ outside that open set.
This proves local outer-capacity convergence. Finally combine
one such truncation's bounded quasicontinuity with the small
ordinary open sublevel set, exactly as in
Corollary~\ref{cor:r2-witness-outer-capacity}. The original
function is finite and continuous on the complement.
\end{proof}

% END INCLUDED SOURCE: base_r2/sections/ambient_polar_traces


% BEGIN INCLUDED SOURCE: base_r2/sections/witness_criteria
\section{Attenuated witnesses and an obstruction to exponential tails}
\label{sec:r2-witness-criteria}

The notation and original test category are those of
Section~\ref{sec:r2-weighted-l1}. The first observation weakens the
hypothesis needed for a pole-set conclusion; it does not assert
integrability of the original witness under that weaker hypothesis.

\begin{lemma}[An integrable witness with the same pole set]
\label{lem:r2-attenuation}
Let $\phi:O\to[-\infty,0]$ be classical psh, and let $\nu$ be a
finite positive Radon measure on $O$. If
\[
 \nu(\{\phi=-\infty\})=0,
\]
there is a negative classical psh function $\psi$ on the same
ordinary open set such that
\[
 \{\psi=-\infty\}=\{\phi=-\infty\},\qquad
                         \int(-\psi)\,d\nu<\infty.
\]
More precisely, one can choose $L_j\geq2^j$, tending increasingly
to infinity, and set
\begin{equation}
\label{eq:r2-attenuated-witness}
 \psi=\sum_{j\geq1}\frac{\max\{\phi,-L_j\}}{L_j},
 \qquad \int(-\psi)\,d\nu\leq1.
\end{equation}
\end{lemma}

\noindent\textbf{Status: \texttt{DIRECT\_COROLLARY}.}
This is an elementary measure-theoretic attenuation combined with
decreasing closure of extended-real classical psh functions
(Gubler--Rabinoff, Proposition 4.2(3)--(4)).

\begin{proof}
Put $t=-\phi$. The function $t$ is finite $\nu$-almost everywhere.
For $L\to\infty$, the functions $\min(t,L)/L$ are bounded by one
and converge to zero $\nu$-almost everywhere. Since $\nu$ is finite,
bounded convergence permits the choice of increasing $L_j\geq2^j$
such that
\[
 \int\frac{\min(t,L_j)}{L_j}\,d\nu\leq2^{-j}.
\]
Every finite partial sum in \eqref{eq:r2-attenuated-witness} is a
negative bounded classical psh function. The partial sums decrease,
so their extended-real limit is classical psh. At a point where
$t<\infty$, the absolute sum is bounded by
\[
 t\sum_j L_j^{-1}<\infty.
\]
At a point where $t=\infty$, each summand is $-1$, so the limit
is $-\infty$. This proves equality of the actual pole sets, including
the complete-extension classicality supplied by decreasing closure.
Finally monotone convergence gives
\[
 \int(-\psi)\,d\nu
   =\sum_j\int\frac{\min(t,L_j)}{L_j}\,d\nu\leq1.
\]
\end{proof}



\begin{corollary}[A pole-set criterion using one background trace]
\label{cor:r2-no-pole-trace}
In the geometry of Theorem~\ref{thm:r2-weighted-l1}, let
$\nu=\beta_0\wedge S$ and let $W$ be ordinary open with
$\overline W\Subset V$. If
\[
 (\beta_0\wedge S)(\{\phi=-\infty\})=0,
\]
then $W\cap\{\phi=-\infty\}$ has zero outer capacity for the
fixed admissibility domain $O$. Every bounded mixed top-degree
measure over $S$ gives this set zero mass. No $L^1$ assumption on
the original $\phi$ is necessary for these conclusions.
\end{corollary}

\noindent\textbf{Status: \texttt{DIRECT\_COROLLARY}.}

\begin{proof}
The compactly supported positive trace $\nu$ is finite. Apply
Lemma~\ref{lem:r2-attenuation}, then apply
Theorem~\ref{thm:r2-weighted-l1} and
Corollary~\ref{cor:r2-witness-outer-capacity} to $\psi$.
The pole sets coincide. All domains remain the originally chosen
ones. For a mixed tuple that is only locally bounded, take its
finite bounds on a neighborhood of the compact core, as in the
weighted theorem.
\end{proof}

\begin{corollary}[Preservation of non-concentration on a specified pole set]
\label{cor:r2-background-pole-preservation}
Let $P=\{\phi=-\infty\}$ for a negative classical psh function
$\phi$ on $U$. Suppose that every coefficient measure of $S$
has total variation zero on $P$. Then every coefficient measure
of $T_{S,k}(u_1,\ldots,u_k)$ has total variation zero on $P$,
for $u_i\in\Cb(U)$ and $0\leq k\leq q$.
In particular, a background whose coefficient measures avoid all
such pole sets retains this property after every bounded BT product.
\end{corollary}

\noindent\textbf{Status: \texttt{DIRECT\_COROLLARY}.}

\begin{proof}
Work locally on nested norm boxes with a collar as in the weighted
theorem. The hypothesis and finite positive controls imply
$(\beta_0\wedge S)(P)=0$, including for the continuous-coefficient
cutoff trace. The previous corollary makes every bounded top-degree
mixed measure over $S$ vanish on $P$ in a smaller core.

For a coefficient of a lower-degree product, localize its original
test to a finite moment presentation, and let
$\beta=D\rho_0$ be the quadratic trace form of that presentation,
where $\rho_0$ is smooth psh and bounded on the larger box.
The finite-test domination from R bounds the variation of that
coefficient on the core by a constant multiple of
\[
 T_{S,k}(u_1,\ldots,u_k)\wedge\beta^{q-k}
   =T_{S,q}(u_1,\ldots,u_k,
                 \underbrace{\rho_0,\ldots,\rho_0}_{q-k}).
\]
The equality is compatibility with smooth slots, already proved in
R1. This mixed measure vanishes on $P$ by the preceding paragraph.
Finite local presentations and ordinary restriction give the
assertion for every original coefficient. The case $k=0$ is the
hypothesis itself.
\end{proof}

\paragraph{The ambient specialization.}
For $S=[1]$, the criterion reduces the general polar-set question to
zero mass of pole sets for the specified original smooth volume
traces. That input is proved separately in
Theorem~\ref{thm:r2-ambient-polar-traces}, using divisorial
non-poles, bounded model approximations and the quasicontinuous
multiplier theorem. Neither the attenuation argument nor that
trace theorem asserts local $L^1$ integrability of the original
unattenuated witness. Complex distributional ellipticity is not
used. The criterion also makes clear why an arbitrary background
concentrated on the pole set is excluded.

\begin{proposition}[No uniform exponential improvement over arbitrary backgrounds]
\label{prop:r2-no-exponential-tails}
Over the original finite extension $K/\mathbf Q_p$, there are a
two-dimensional allowed ordinary open set $U$, a bounded continuous
classical psh function $f$, a background $S=D(f[1])$, a negative
classical psh witness $\phi$, and a compact $K_0\Subset U$ with
the following properties. The background trace integrability
condition of Theorem~\ref{thm:r2-weighted-l1} holds for every
fixed collar used below, but there is a normalized continuous
classical psh $v$ such that
\[
 \int_{K_0}\exp(-a\phi)\,T_{S,1}(v)=\infty
                   \quad\text{for every }a>0.
\]
Moreover the relative capacities of the deep sublevels have a
polynomial lower bound. Consequently the weighted $L^1$ theorem
cannot be strengthened, under its stated hypotheses, to an
exponential tail bound uniform over normalized potentials.
\end{proposition}

\noindent\textbf{Status of the exponential assertion refuted:
\texttt{FALSE}.}

\begin{proof}
Choose a uniformizer $\pi$ and put $\lambda=-\log|\pi|>0$.
Choose $0<r<1<R$ and set
\[
 U=\{|z_0|<1\}\times\{r<|z_1|<R\}.
\]
For $j\geq1$, put
\[
 f_j=\max\{\log|z_0|,-j\lambda\},\qquad
 f=\sum_{j\geq1}j^{-4}f_j,\qquad
 \phi=\log|z_0|,\qquad
 v=\frac{\max\{0,\log|z_1|\}}{\log R}.
\]
Each $f_j$ is continuous classical psh and takes values in
$[-j\lambda,0]$. The series for $f$ converges uniformly on $U$
because $\sum j^{-3}<\infty$. Its partial sums decrease, so $f$
is bounded continuous classical psh. The current
$S=D(f[1])$ is symmetric, original-positive and separately closed.
Writing $S_j=D(f_j[1])$, linearity and convergence give
\[
 S=\sum_{j\geq1}j^{-4}S_j.
\]
The positive series converges locally as a coefficient-Radon
current. Indeed, relative CLN bounds its mass against each fixed
positive cutoff by a constant times
$\sum j^{-4}(1+j\lambda)<\infty$.

The support of $S_j$ lies in $\{|z_0|=|\pi|^j\}$: on one
side $f_j$ is constant and on the other it is the logarithm of
the norm of an analytic unit. Thus $-\phi=j\lambda$ on this
support. For a fixed positive collar trace $\beta_0$, finite-test
domination and CLN give a constant $C_\rho$, independent of $j$,
such that
\[
 \int\beta_0\wedge S_j\leq C_\rho(1+j\lambda).
\]
The same is valid for its continuous coefficients by positive
smooth controls. Equality of the positive coefficient measures of
the series and monotone convergence imply
\[
 A_{S,\rho}(\phi)
  =\sum_{j\geq1}j^{-4}\,j\lambda
                       \int\beta_0\wedge S_j
  \leq C_\rho\sum_{j\geq1}j^{-4}j\lambda(1+j\lambda)<\infty.
\]

Let $x_j$ be the product Gauss point with radii $(|\pi|^j,1)$.
The usual model metric of $\mathcal O(1)$ on the first projective
line, after the coordinate rescaling $z_0\mapsto z_0/\pi^j$,
has affine weight $f_j+j\lambda$. Together with the standard
model metric in $z_1$, the mixed model intersection has degree
one. The original calibrated model-measure identity therefore gives
\[
 T_{S_j,1}(v)=(\log R)^{-1}\delta_{x_j}.
\]
Fixed continuous multiplication by $v$ and distributional $D$
pass the positive partial sums to the limit, yielding
\begin{equation}
\label{eq:r2-polynomial-tail-measure}
 T_{S,1}(v)=(\log R)^{-1}
                            \sum_{j\geq1}j^{-4}\delta_{x_j}.
\end{equation}
Choose $r<r_0<1<R_0<R$ and take
\[
 K_0=\{|z_0|\leq|\pi|\}
                      \times\{r_0\leq|z_1|\leq R_0\}.
\]
This is compact in $U$, contains every $x_j$, and lies in local
norm boxes admitting the required collar. Since
$\phi(x_j)=-j\lambda$, equation
\eqref{eq:r2-polynomial-tail-measure} gives
\[
 \int_{K_0}(-\phi)\,T_{S,1}(v)
      =\frac{\lambda}{\log R}\sum_{j\geq1}j^{-3}<\infty,
 \qquad
 \int_{K_0}e^{-a\phi}\,T_{S,1}(v)
      =\frac1{\log R}\sum_{j\geq1}j^{-4}e^{aj\lambda}=\infty.
\]
The potential $v$ is between zero and one on all of $U$. Hence,
for every $N>0$,
\[
 \operatorname{Cap}_S(K_0\cap\{\phi<-N\};U)
 \geq\frac1{\log R}\sum_{j>N/\lambda}j^{-4}
 \geq\frac1{3\log R}(1+N/\lambda)^{-3}.
\]
This cannot be bounded above by $A e^{-aN}$ with fixed finite
$A$ and $a>0$. All powers of $\pi$, all Gauss points, and all
model metrics used here are defined over the original field with
its prescribed absolute value.
\end{proof}

\begin{proposition}[A uniform first moment does not give uniform integrability]
\label{prop:r2-unbounded-weight-failure}
Even for $S=[1]$ on a smooth analytic curve, a fixed negative psh
witness satisfying the collar integrability condition need not
commute with weak limits of normalized bounded BT measures.
There can be positive measures $\mu_j\to0$ with
\[
 \int_{K_0}(-\phi)\,d\mu_j=1\quad\text{for all }j,
 \qquad (-\phi)\mu_j\longrightarrow\delta_0\ne0
\]
on one fixed compact $K_0$ containing their supports.
\end{proposition}

\noindent\textbf{Status of the asserted unbounded-weight stability:
\texttt{FALSE}.}

\begin{proof}
On $U=\{|z|<1\}\subset\mathbf A^{1,\an}_K$, set
$\phi=\log|z|$, with value $-\infty$ at $0$, and put
\[
 v_j=1+\frac{\max\{\log|z|,-j\lambda\}}{j\lambda},
 \qquad \lambda=-\log|\pi|.
\]
These are continuous classical psh functions with $0\leq v_j\leq1$.
Let $\xi_j$ be the Gauss point of radius $|\pi|^j$. The calibrated
one-variable model calculation gives
\[
 \mu_j=D(v_j[1])=(j\lambda)^{-1}\delta_{\xi_j},
 \qquad -\phi(\xi_j)=j\lambda.
\]
All points $\xi_j$ lie in the fixed compact disc
$K_0=\{|z|\leq|\pi|\}\Subset U$ and tend to $0$.
Thus $\mu_j\to0$ even in total variation, whereas
$(-\phi)\mu_j=\delta_{\xi_j}\to\delta_0$ as Radon measures,
and hence on every original compact smooth scalar test.
The witness is finite $\mu_j$-almost everywhere. Its first moment
is exactly one, but for every $M>0$,
\[
 \sup_j\int_{K_0\cap\{-\phi>M\}}(-\phi)\,d\mu_j=1.
\]
This is the precise failure of uniform integrability.

The background trace hypothesis can be checked directly. Choose
a radial smooth psh collar $\rho$ that is constant on a disc
containing $K_0$, with $\{\rho<0\}\Subset U$ still containing
that disc. Such a collar is a convex smooth regularized maximum
of zero and an affine function of $\log|z|$, minus a constant.
For $q=1$, its trace
$\beta_0=(-\rho)_+^4D\rho+
4(-\rho)_+^3d'\rho\wedge d''\rho$ is supported in a compact
annulus away from $0$.
The function $\phi$ is bounded there, so
$\int(-\phi)\beta_0\wedge[1]<\infty$.

The sequence $v_j$ increases. Its actual pointwise supremum is
one off $0$ and zero at $0$, so it is not upper semicontinuous.
It belongs to the bounded capacitary closure by
Corollary~\ref{cor:r2-local-increasing}. No claim is made that
this supremum is a classical psh limit. The bounded-weight
calculus is consistent with $\mu_j\to0$; the failure concerns
only the attempted extension to the fixed unbounded multiplier.
\end{proof}

% END INCLUDED SOURCE: base_r2/sections/witness_criteria


% BEGIN INCLUDED SOURCE: base_r2/sections/obstacle_maximality
\section{Quasicontinuous obstacles and local maximality}
\label{sec:r2-obstacle-maximality}

This section concerns the ambient background $[1]$. It uses the
continuous-obstacle envelope and orthogonality theorem of
\textup{B}, Theorem C.6, without changing its hypotheses. A separate
limit argument extends orthogonality to a class of bounded
discontinuous obstacles. It then supplies the local equivalence
between maximality and vanishing Monge--Amp\`ere measure. No
relative obstacle orthogonality for an arbitrary background is
asserted.

\begin{lemma}[Continuous majorants of a quasicontinuous function]
\label{lem:r2-continuous-majorants-capacity}
Let $f_j\downarrow f$ be continuous real functions on an ordinary
neighborhood of a compact $H$, with $f$ bounded and
$C_\Omega$-quasicontinuous on a compact neighborhood of $H$.
Then
\[
 C_\Omega^*(H\cap\{f_j-f>\delta\})\longrightarrow0
       \qquad(\delta>0).
\]
Here $\Omega$ is the fixed local capacity domain containing that
neighborhood.
\end{lemma}

\begin{proof}
For $\epsilon>0$, choose an ordinary open $G$ of capacity less
than $\epsilon$ such that $f$ is continuous on the chosen compact
neighborhood minus $G$. On this compact set, the continuous
functions $f_j$ decrease to a continuous function. Dini's theorem
gives uniform convergence there. Thus the indicated level set is
contained in $G$ for every sufficiently large $j$, and its outer
capacity is at most $C_\Omega(G)<\epsilon$.
\end{proof}

\begin{theorem}[Orthogonality for bounded quasicontinuous usc obstacles]
\label{thm:r2-qc-orthogonality}
Let $Y/F$ be smooth projective geometrically integral over a finite
extension $F/\mathbf Q_p$, with its prescribed absolute value,
and let $L$ be ample. Let $f$ be a bounded
upper-semicontinuous metric on $L$, relative to a continuous
reference metric. Assume that its local weights are locally
quasicontinuous for the ambient capacities $\operatorname{Cap}_{[1]}$
on the permitted local norm boxes. Put
\[
 p=P_L^{\mathrm b} f,
\]
the greatest bounded classical minorant from
Proposition~\ref{prop:r2-usc-obstacle}. Then
\begin{equation}
\label{eq:r2-qc-orthogonality}
 \int_{Y^{\an}}(f-p)c_1(L,p)^n=0.
\end{equation}
In particular this applies to $f=r+g$, where $r$ is a bounded
classical metric and $g$ is any continuous real function on
$Y^{\an}$.
\end{theorem}

\noindent\textbf{Status: \texttt{NEW\_PROVED}.}
The new step is the extension to these discontinuous obstacles;
the continuous-obstacle theorem is an \texttt{INTERNAL\_INPUT}.

\begin{proof}
Choose uniformly bounded continuous metrics $f_j\downarrow f$.
Let $p_j=P_L f_j$. The accepted envelope theorem makes each
$p_j$ continuous model-semipositive and gives
\[
 \int(f_j-p_j)c_1(L,p_j)^n=0.
\]
Proposition~\ref{prop:r2-usc-obstacle} gives $p_j\downarrow p$,
with a common bound relative to a continuous semipositive
reference metric. Write $\mu_j=c_1(L,p_j)^n$ and
$\mu=c_1(L,p)^n$. The fixed-background monotone theorem, in
local frames, gives $\mu_j\to\mu$ on all original scalar tests.
The measures have the same finite total mass $\deg_F(L^n)$,
by continuous compatibility and passage to the decreasing limit.

Fix a local compact test stage. In its frame the weights of all
$p_j$ have uniform bounds on a slightly larger box. The R1
coefficient-capacity estimate therefore controls $\mu_j$
uniformly by that fixed box's ambient capacity. By
Lemma~\ref{lem:r2-continuous-majorants-capacity}, the local weights
of $f_j$ converge to those of $f$ in outer capacity. These weights
are quasicontinuous, as is their limit. The multiplier theorem
\ref{thm:r2-weak-closure-multipliers} yields convergence of the
integrals with multiplier $f_j$. The weighted decreasing theorem
yields convergence with multiplier $p_j$. Subtract the two
convergences. A finite original smooth partition of unity on the
compact projective space makes the differences of local weights
equal to the global scalar functions $f_j-p_j$ and $f-p$.
Summing gives \eqref{eq:r2-qc-orthogonality}.

Finally, bounded classical local weights are quasicontinuous by
R1, and addition of a continuous function preserves
quasicontinuity and upper semicontinuity. Thus $r+g$ satisfies all
the stated assumptions. No monotone-convergence theorem for
arbitrary discontinuous weights was used: their convergence was
proved using Dini's theorem off ordinary small-capacity opens.
\end{proof}

\begin{definition}[Local maximality]
\label{def:r2-maximality}
A function $u\in\Cb(U)$ is locally maximal if every point has
an ordinary neighborhood $V$ with compact closure in a smooth
affine algebraic chart such that the following holds. Whenever
$v\in\Cb(V)$ and $v=u$ on $V\setminus K$ for some compact
$K\Subset V$, one has $v\leq u$ on $V$.

Equivalently one may use competitors whose strict positive gap
over $u$ has compact closure in $V$: replace such a competitor
by its maximum with $u$. This definition is local and does not
use the whole of a proper analytic space as a boundaryless
comparison domain.
\end{definition}

\begin{theorem}[Maximality and the homogeneous equation]
\label{thm:r2-maximal-ma-zero}
Let $u\in\Cb(U)$ on any allowed good, Hausdorff, boundaryless,
quasi-smooth pure-dimensional space over the fixed $K$. Then
\[
 u\text{ is locally maximal}
       \quad\Longleftrightarrow\quad (Du)^n=0\text{ on }U.
\]
The equality is that of the ordinary original-CLD Radon measure.
No continuity assumption on $u$ is imposed.
\end{theorem}

\noindent\textbf{Status: \texttt{NEW\_PROVED}, relative to the adopted
inputs and the preceding orthogonality extension.}

\begin{proof}
Assume first that $(Du)^n=0$. On a relatively compact framed
affine neighborhood $V$, a competitor as in
Definition~\ref{def:r2-maximality} satisfies every compact-gap
condition of Corollary~\ref{cor:r2-pointwise-ambient}. Its strict
gap has zero $(Du)^n$-mass. That corollary gives $u\geq v$
pointwise. Such neighborhoods exist by \textup{B}, so $u$ is
locally maximal.

Conversely, fix a point where maximality holds on $V$. The
common-collar construction of \textup{B}, Lemma D.2, realizes
$u+c$, on a smaller ordinary neighborhood $W\Subset V$, as the
weight of a bounded classical metric $r$ on an ample line $L$
over a smooth projective variety. Its frame is fixed on $W$.
Choose any nonnegative original smooth scalar $\chi$ compactly
supported in $W$; it extends by zero to the projective space.
For $\epsilon>0$, let
\[
 f=r+\epsilon\chi,\qquad p=P_L^{\mathrm b}f.
\]
The metric $r$ is a permitted minorant, so
$r\leq p\leq r+\epsilon\chi$. In particular $p=r$ off the
compact support of $\chi$. Its weight on $W$, minus $c$, glues
with $u$ outside that support to give a member of $\Cb(V)$
which equals $u$ off a compact subset of $V$. Maximality forces
this function to be at most $u$. Therefore $p=r$ everywhere,
also outside $W$ by the obstacle inequalities.

The obstacle $f=r+\epsilon\chi$ satisfies
Theorem~\ref{thm:r2-qc-orthogonality}. Its conclusion is
\[
 0=\int(f-p)c_1(L,p)^n
   =\epsilon\int\chi\,c_1(L,r)^n
   =\epsilon\int_W\chi\,(Du)^n.
\]
This holds for every such $\chi$, so positivity and Radon
uniqueness give $(Du)^n=0$ on $W$. Cover $U$ by these
neighborhoods. If the local algebraic chart uses the finite
constant-field convention in \textup{B}, Lemma D.5, its original
calibration changes by a positive constant; the vanishing
assertion is unchanged. Thus it holds for the prescribed $K$
and its original absolute value.
\end{proof}

\begin{remark}[Scope for Dirichlet theory]
The theorem identifies bounded solutions of the homogeneous
equation with a local maximality property. It does not prove
that a Perron envelope attains arbitrary prescribed boundary
values, that every ordinary domain admits boundary barriers,
or that a nonzero prescribed measure admits a solution. Those
are separate Dirichlet existence and boundary-regularity
questions. For an arbitrary background $S$, even pointwise
domination can fail, as the connected nonzero-background
counterexample in this record shows; the ambient theorem above
cannot be transferred to that setting without further hypotheses.
\end{remark}

% END INCLUDED SOURCE: base_r2/sections/obstacle_maximality


% BEGIN INCLUDED SOURCE: base_r2/sections/perron_peak_barriers
\section{A homogeneous Dirichlet theorem with explicit peak barriers}
\label{sec:r2-perron}

The results in this section concern the ambient background $[1]$.
The geometric hypothesis is stated separately: no assertion is
made that an arbitrary Berkovich analytic domain has these
barriers. The solution will be bounded and classical psh, with
the stated ordinary boundary limits. Interior continuity is not
part of the conclusion.

Let $V$ be a connected ordinary open set whose closure $\overline V$
is compact and is contained in a smooth affine algebraic chart.
Write
\[
 \partial V=\overline V\setminus V
\]
for its ordinary topological boundary, and assume that this set
is nonempty. All closures and limits in this section refer to
the ordinary Berkovich topology.

\begin{definition}[Continuous psh peak barriers]
\label{def:r2-peak-barriers}
The domain $V$ has continuous psh peak barriers if for each
$\xi\in\partial V$ there is a function
\[
 \beta_\xi\in \Cb(V)\cap C(\overline V)
\]
such that $\beta_\xi(\xi)=0$ and
$\beta_\xi(x)<0$ for every $x\in\overline V\setminus\{\xi\}$.
\end{definition}

\begin{lemma}[Boundary maximum principle]
\label{lem:r2-boundary-maximum}
Suppose that $h\in\Cb(V)$ is bounded above and
\[
 \limsup_{V\ni x\to\xi}h(x)\leq0
       \qquad(\xi\in\partial V).
\]
Then $h\leq0$ on $V$.
\end{lemma}

\begin{proof}
If $h(x_0)>0$, choose $0<a<h(x_0)$. For each boundary point the
limsup condition gives a neighborhood on which $h<a$ in $V$.
The union of these neighborhoods contains the compact boundary.
Thus $\{h\geq a\}$ is contained in a compact subset of $V$; it
is closed there by upper semicontinuity. The bounded usc
function $h$ attains its global positive maximum at a point
of $V$. The strong maximum principle for classical psh
functions, \cite[Theorem 6.1]{GR25}, makes it constant on the
connected domain $V$. A positive constant contradicts the
boundary condition at the nonempty boundary.
\end{proof}

\begin{theorem}[Homogeneous Dirichlet problem under peak barriers]
\label{thm:r2-peak-dirichlet}
Assume that $V$ has the barriers in
Definition~\ref{def:r2-peak-barriers}. For every
$g\in C(\partial V)$ there is a unique bounded classical psh
function $H_g$ on $V$ such that
\[
 (DH_g)^n=0\quad\hbox{on }V,\qquad
 \lim_{V\ni x\to\xi}H_g(x)=g(\xi)
       \quad(\xi\in\partial V).
\]
The measure equation is tested against all original compactly
supported smooth CLD scalar tests. Moreover
\[
 \min_{\partial V}g\leq H_g\leq\max_{\partial V}g,
 \qquad
 \|H_{g_1}-H_{g_2}\|_{L^\infty(V)}
       \leq\|g_1-g_2\|_{L^\infty(\partial V)} .
\]
In particular the solution is monotone with respect to its
boundary values.
\end{theorem}

\noindent\textbf{Status: \texttt{NEW\_PROVED} under the displayed
barrier hypothesis.} The corresponding assertion for general
ordinary analytic domains remains \texttt{OPEN}; no existence
of barriers in general dimension is supplied here.

\begin{proof}
Put $m=\min_{\partial V}g$ and $M=\max_{\partial V}g$, and let
$\mathcal F_g$ consist of all $v\in\Cb(V)$ satisfying
\[
 m\leq v\leq M,\qquad
 \limsup_{V\ni x\to\eta}v(x)\leq g(\eta)
       \quad(\eta\in\partial V).
\]
This family contains the constant $m$. Let
$H=\sup_{v\in\mathcal F_g}v$, pointwise, and let $H^*$ denote
its ordinary upper-semicontinuous regularization on $V$.
Theorem~\ref{thm:r2-local-upper-envelope} gives
$H^*\in\Cb(V)$ and $m\leq H^*\leq M$.

Fix $\xi\in\partial V$ and $\epsilon>0$. Continuity of $g$
gives a boundary neighborhood of $\xi$ on which
$|g-g(\xi)|<\epsilon$. On the compact complement,
$\beta_\xi$ has a strictly negative maximum. Consequently
one can choose a finite $A>0$ such that, on all of $\partial V$,
\[
 b_-=g(\xi)-\epsilon+A\beta_\xi\leq g,\qquad
 b_+=g(\xi)+\epsilon-A\beta_\xi\geq g .
\]
The function $b_-$ is classical psh on $V$, continuous on the
closure, and at most $M$. Thus $\max(m,b_-)\in\mathcal F_g$,
and $H^*\geq b_-$. For every $v\in\mathcal F_g$, the function
\[
 v-b_+=v+A\beta_\xi-g(\xi)-\epsilon
\]
is bounded classical psh and has boundary limsup at most zero.
Lemma~\ref{lem:r2-boundary-maximum} implies $v\leq b_+$.
Since $b_+$ is continuous, also $H^*\leq b_+$. Taking limits
at $\xi$ in
\[
 g(\xi)-\epsilon+A\beta_\xi
       \leq H^*
       \leq g(\xi)+\epsilon-A\beta_\xi
\]
and then letting $\epsilon$ decrease to zero proves
$H^*(x)\to g(\xi)$ as $x\to\xi$ from $V$. Hence
$H^*\in\mathcal F_g$, so $H=H^*$ pointwise.

It remains to prove the equation. Let $W\Subset V$ be an
ordinary neighborhood, and let $v\in\Cb(W)$ agree with $H$
outside a compact set $K\Subset W$. The function equal to
$\max(v,H)$ on $W$ and to $H$ off $K$ is well defined and
classical psh by ordinary gluing and maximum stability.
Call it $w$. It is globally bounded, since its modification
is confined to $K$, and it has boundary limits $g$.
Lemma~\ref{lem:r2-boundary-maximum}, applied to $w-M$,
gives $w\leq M$; also $w\geq H\geq m$. Therefore
$w\in\mathcal F_g$ and $w\leq H$ by the definition of $H$.
It follows that $v\leq H$. This proves local maximality
of $H$, and Theorem~\ref{thm:r2-maximal-ma-zero} yields
$(DH)^n=0$.

If $H_1,H_2$ are two bounded classical solutions with the
same boundary values, then for each $\epsilon>0$ the set
$\{H_1<H_2-\epsilon\}$ has compact closure in $V$.
Indeed, the ordinary boundary limits give a neighborhood
of each boundary point on which this set is empty, and
compactness of the boundary gives a common compact core.
The analogous assertion holds after interchanging the
functions. Since their Monge--Amp\`ere measures vanish,
pointwise ambient domination
(Corollary~\ref{cor:r2-pointwise-ambient}) gives both
$H_1\geq H_2$ and $H_2\geq H_1$.

More generally set
$\delta=\|g_1-g_2\|_{L^\infty(\partial V)}$.
The pair $H_{g_2}+\delta,H_{g_1}$ satisfies the required
one-sided compact-gap condition and has zero measures, so domination
gives $H_{g_1}\leq H_{g_2}+\delta$. Reversing the indices
proves the norm estimate. If $g_1\leq g_2$, apply this
comparison directly with no shift to obtain
$H_{g_1}\leq H_{g_2}$.
\end{proof}

\begin{remark}[What the theorem does not settle]
A classification or construction of suitable domains in
arbitrary dimension is a separate geometric problem.
The proof uses continuous peak barriers to control the
ordinary boundary and the discontinuous-obstacle theorem
to identify maximality with the homogeneous equation.
It supplies neither interior continuity for $H_g$ nor
existence for an arbitrary prescribed nonzero measure.
For an arbitrary positive closed background $S$, the
pointwise domination step would require an additional
gap-detection hypothesis.
\end{remark}


% END INCLUDED SOURCE: base_r2/sections/perron_peak_barriers


% BEGIN INCLUDED SOURCE: base_r2/sections/capacity_noncompactness
\section{A bounded family which is not precompact in capacity}
\label{sec:r2-capacity-noncompactness}

This example uses the fixed ambient background $S=[1]$. All
potentials, including their pointwise limit, are continuous
classical psh. It rules out replacing the capacity-convergence
hypothesis by bounded pointwise convergence, and rules out a
capacity-precompactness argument for the whole normalized test
family. It is not a counterexample to weak-closure domination
of that family's measures.

\subsection{A moving residue branch over the original field}

Let $\pi$ be a uniformizer of the fixed finite extension
$K/\mathbf Q_p$, and retain its original absolute value. Put
\[
 \lambda=-\log|\pi|>0,\qquad
 V=\{x\in\mathbf A_K^{1,\an}:|T(x)|<|\pi|^{-1}\}.
\]
Choose distinct monic irreducible polynomials
$\overline P_i\in\widetilde K[T]$, of degrees $d_i\geq2$,
and monic lifts $P_i\in K^\circ[T]$. Define
\[
 h_i=\frac1{d_i}\max\{\log|P_i(T)|,-d_i\lambda\},
 \qquad
 \phi_i=\lambda^{-1}h_i,\qquad
 \psi_i=\frac{1+\phi_i}{2}.
\]
These are continuous classical psh functions. In fact $h_i$ is
the local weight of the pure Fubini--Study metric on
$\mathcal O_{\mathbf P^1}(1)$ specified, after tensoring by
$d_i$, by the two sections
\[
 P_i^{\mathrm{hom}}(X,Z),\qquad \pi^{d_i}Z^{d_i}.
\]
They have no common zero. The factors $d_i^{-1}$ and
$\lambda^{-1}$ above are factors in the potential, not a
renormalization of the field or of the integration current.

Let $\zeta$ be the unit Gauss point. Let $\xi_i$ be the unique
Shilov point of the closed residue disc
\[
 B_i=\{|P_i(T)|\leq|\pi|^{d_i}\}.
\]
More explicitly, $P_i$ has $d_i$ simple Hensel roots in the
unramified extension of degree $d_i$. Over that extension,
$B_i$ is the disjoint union of the closed discs of radius
$|\pi|^{d_i}$ about those roots. Their conjugate Gauss points
restrict to the same original $K$-point $\xi_i$.
The residue field of its branch has degree $d_i$.

\begin{lemma}[The branch weights and their measures]
\label{lem:r2-moving-branch-measures}
On $V$ one has $0\leq\psi_i<1$ and
\[
 D h_i=\delta_{\xi_i},\qquad
 D\psi_i=\frac1{2\lambda}\delta_{\xi_i}.
\]
Moreover $\xi_i\to\zeta$, and pointwise on $V$
\[
 \psi_i\longrightarrow
 \psi=\frac12+\frac1{2\lambda}\max\{0,\log|T|\}.
\]
The limit is continuous classical psh, and
\[
 D\psi=\frac1{2\lambda}\delta_\zeta.
\]
At the indicated points,
\[
 \psi_i(\xi_i)=0,\qquad
 \psi_j(\xi_i)=\frac12\ (j\ne i),\qquad
 \psi(\xi_i)=\psi(\zeta)=\frac12.
\]
\end{lemma}

\begin{proof}
Irreducibility of the separable reduction gives the unramified
root extension and the described discs by Hensel's lemma.
The closed $K$-domain $B_i$ is the closed disc over this
unramified residue extension, viewed as a $K$-analytic domain;
it has one Shilov point. On its interior $h_i=-\lambda$.
Off $B_i$, the function is $d_i^{-1}\log|P_i|$, which is
harmonic there because $P_i$ has no zero there. Thus the
positive curvature measure of the displayed model metric is
supported at $\xi_i$. Its total mass is
$\deg_K\mathcal O_{\mathbf P^1}(1)=1$ in the original model
formula, so it is $\delta_{\xi_i}$. Equivalently, the radial
slope $d_i^{-1}$ at that branch is multiplied by its residue
degree $d_i$; these two factors cancel. This is the same
normalization, not a degree-one assertion about the residue
field. Scaling gives the formula for $D\psi_i$.

For $|T|\leq1$, all $|P_i|\leq1$, so $-\lambda\leq h_i\leq0$.
For $|T|>1$, monicity and integral coefficients give
$|P_i(T)|=|T|^{d_i}$, hence $h_i=\log|T|$. The asserted
normalization on $V$ follows.

At a point of the closed unit disc, at most one of the numbers
$|P_i|$ can be strictly less than one. Indeed, two distinct
irreducible reductions are coprime; a lifted Bezout identity
would otherwise have left side of norm less than one and right
side of norm one. Therefore $h_i(x)\to0$ at every such point.
Outside that disc all $h_i$ equal $\log|T|$. This proves the
pointwise convergence. The formula for $D\psi$ is the
one-variable canonical model formula at the unit Gauss point.
The same coprimality argument gives all the stated point values.

Finally, take a nonzero $F\in K[T]$ and divide by a coefficient
of maximum norm so that its reduction is a nonzero polynomial.
Only finitely many of the $\overline P_i$ divide that
reduction. For all other $i$,
\[
 |F(\xi_i)|=\max_\nu|\operatorname{coeff}_\nu(F)|
           =|F(\zeta)|.
\]
Polynomial seminorms define the Berkovich topology, so
$\xi_i\to\zeta$.
\end{proof}

\subsection{The counterexample in every dimension}

Let $\Omega=V^n$, an ordinary open in $\mathbf A_K^{n,\an}$.
For $2\leq l\leq n$ put
\[
 b_l(T_l)=\lambda^{-1}\max\{0,\log|T_l|\},
 \qquad
 u_i=\frac1n\left(\psi_i(T_1)+\sum_{l=2}^n b_l(T_l)\right).
\]
For $n=1$ the sum is empty. Then $0\leq u_i<1$ on
$\Omega$, and $u_i$ converges pointwise to the continuous
classical psh function
\[
 u=\frac1n\left(\psi(T_1)+\sum_{l=2}^n b_l(T_l)\right).
\]
Let $z_i$ be the product Gauss point with first coordinate
$\xi_i$ and the remaining coordinates unit Gauss points, and
let $z$ be the all-unit Gauss point. Explicitly, its seminorm
on polynomials in the remaining coordinates is the maximum
of the coefficient seminorms at $\xi_i$. This description
shows that $z_i\to z$.

\begin{proposition}[Failure of capacity precompactness and pointwise
weighted stability]
\label{prop:r2-capacity-noncompactness}
Set
\[
 c_n=\frac{n!}{2n^n\lambda^n}>0.
\]
For the fixed ambient background $[1]$,
\[
 (Du_i)^n=c_n\delta_{z_i},\qquad
 (Du)^n=c_n\delta_z.
\]
For $i\ne j$,
\[
 \operatorname{Cap}_{[1]}
    \left(\left\{|u_i-u_j|>\frac1{4n}\right\};\Omega\right)
       \geq c_n.
\]
Consequently the normalized continuous classical family
$(u_i)$ has no capacity-Cauchy subsequence. The same lower
bound holds for the error sets with $u_j$ replaced by the
pointwise limit $u$.
This failure also holds for local capacity convergence in the
all-domain sense: one compact subset of one fixed smaller
domain already detects every pairwise error.

Although $(Du_i)^n\to(Du)^n$ weakly, the weighted products fail
to converge to the product of the pointwise limits:
\[
 u_i(Du_i)^n=0\quad\text{for every }i,\qquad
 u(Du)^n=\frac{c_n}{2n}\delta_z\ne0.
\]
\end{proposition}

\begin{proof}
The mixed product with one factor from each coordinate is
\[
 D\psi_i(T_1)\wedge Db_2(T_2)\wedge\cdots\wedge Db_n(T_n)
       =\frac1{2\lambda^n}\delta_{z_i}.
\]
This is the product model formula: the mixed intersection
number of the $n$ distinct pullbacks of
$\mathcal O_{\mathbf P^1}(1)$ is one. The first coordinate
retains the residue-degree cancellation proved in the lemma.
There is no factorial in this mixed product. A term with a
repeated coordinate vanishes; it is the product of two
curvatures pulled back from the same curve. Equivalently,
the relevant positive product on $(\mathbf P^1)^n$ has zero
mixed intersection number and hence zero mass.
Expanding the $n$th power of the sum defining $u_i$ therefore
contributes exactly $n!$ equal mixed terms and the factor
$n^{-n}$. This gives $c_n\delta_{z_i}$; the same calculation
gives the formula for $u$.

At $z_i$ one has
\[
 u_i(z_i)=0,\qquad
 u_j(z_i)=u(z_i)=\frac1{2n}\quad(j\ne i).
\]
The error set thus contains $z_i$. The function $u_i$ itself
is admissible in the capacity supremum, so its Monge--Amp\`ere
mass $c_n$ at $z_i$ proves both lower bounds. They hold for
every pair of distinct indices, excluding every
capacity-Cauchy subsequence.

For the local assertion, put
\[
 K_*=\{z\}\cup\{z_i:i\geq1\},\qquad
 \Omega_0=\operatorname{Int}
   \{\,|T_l|\leq|\pi|^{-1/2}\text{ for }1\leq l\leq n\,\}.
\]
The convergence $z_i\to z$ makes $K_*$ compact, and
$K_*\Subset\Omega_0\Subset\Omega$. All $u_i$ remain normalized
admissible tests after restriction to $\Omega_0$. Therefore
\[
 \operatorname{Cap}_{[1]}
  \left(K_*\cap\left\{|u_i-u_j|>\frac1{4n}\right\};\Omega_0\right)
       \geq c_n\qquad(i\ne j),
\]
and the same estimate holds with $u_j$ replaced by $u$.
Thus passing to local capacities does not restore a
capacity-Cauchy subsequence.

Weak convergence of the unweighted measures follows from
$z_i\to z$. The displayed point values give the weighted
identities, with $u(z)=1/(2n)$. These are different ordinary
Radon measures and are distinguished by an original smooth
compact scalar cutoff which equals one near $z$.
\end{proof}

\paragraph{Status and excluded route.}
The claims of capacity precompactness of the whole normalized
family, and of weighted continuity under bounded pointwise
convergence, are \texttt{FALSE}; this is an explicit
counterexample in every dimension. Each individual ingredient
uses the established continuous/model calculus. The example
does not contradict monotone continuity, convergence in
capacity, the completion theorem, or weak-closure domination:
its unweighted weak limit is itself the Monge--Amp\`ere measure
of an admissible continuous potential.

% END INCLUDED SOURCE: base_r2/sections/capacity_noncompactness


% BEGIN INCLUDED SOURCE: base_r2/sections/divisorial_determination
\section{Divisorial determination of the ambient capacitary order}
\label{sec:r2-divisorial-determination}

The statement here is an additional consequence of the model
approximation and controlled-multiplier arguments already proved.
It gives the pointwise reduction used to solve the
increasing-envelope problem in the following section. Density of divisorial points by itself
would not suffice for any assertion about discontinuous functions.

\begin{proposition}[Divisorial test for local quasi-everywhere order]
\label{prop:r2-divisorial-capacitary-order}
Let $V$ be an ordinary open with compact closure in a smooth
affine algebraic chart. Let $a,b$ be locally bounded real Borel
functions on $V$, locally quasicontinuous for the ambient
capacities on all permitted local norm boxes. Then the following
conditions are equivalent:
\begin{enumerate}[label=(\roman*)]
\item $a\leq b$ locally outside sets of outer
$\operatorname{Cap}_{[1]}$-capacity zero;
\item $a(z)\leq b(z)$ for every ambient divisorial point $z\in V$.
\end{enumerate}
In particular the local ambient capacitary equivalence class of
a bounded quasicontinuous function is determined by its actual
values at ambient divisorial points.
\end{proposition}

\noindent\textbf{Status: \texttt{DIRECT\_COROLLARY}.}
The inputs are the atomic testers, the whole-box model
regularization from B, and the controlled quasicontinuous
multiplier theorem; a new regularization theorem is not asserted.

\begin{proof}
For (i)$\Rightarrow$(ii), a divisorial point has positive
singleton capacity on every sufficiently small framed affine
capacity domain containing it, by Lemma~\ref{lem:r2-atomic-tester}.
It therefore cannot lie in a local outer-capacity-zero
exceptional set.

For the converse fix a permitted capacity domain $\Omega\subset V$
and an ordinary open $W\Subset\Omega$ on a compact neighborhood
of whose closure $a$ and $b$ are bounded and
$\operatorname{Cap}_{[1]}(\,\cdot\,;\Omega)$-quasicontinuous.
Let $H\subset W$ be compact. Set $g=(a-b)_+$. This is a bounded
nonnegative quasicontinuous Borel function near $\overline W$,
and it vanishes at every ambient divisorial point there.

Let $v\in\Cb(\Omega)$ with $0\leq v\leq1$. Cover $H$ by finitely
many smaller cores $H_\alpha$ and choose nonnegative original
smooth scalar cutoffs $\chi_\alpha$, supported in slightly
larger norm boxes and equal to one on $H_\alpha$.
The whole-box regularization from B gives, on each such larger
box, a uniformly bounded decreasing sequence of continuous
semipositive rational pure Fubini--Study weights
$v_{\alpha,j}\downarrow v$. Here rational pure-FS is the
model class specified in B, Theorem 4.4(2); arbitrary real
weights in a Fubini--Study presentation are not needed.
Their Monge--Amp\`ere measures
$\mu_{\alpha,j}=(Dv_{\alpha,j})^n$ are restrictions of model
measures and hence are supported at ambient divisorial points.
If the construction shifts $v$ by a fixed constant, subtract
that constant; curvature and the argument are unchanged.
Consequently
\[
 \int \chi_\alpha g\,d\mu_{\alpha,j}=0 .
\]
The measures converge weakly to $(Dv)^n$ and, on the support
of $\chi_\alpha$, have a common bound by the capacity of one
fixed intermediate norm box. The function $g$ is quasicontinuous
for this same box. Theorem~\ref{thm:r2-weak-closure-multipliers}
therefore gives
\[
 \int \chi_\alpha g\,(Dv)^n=0 .
\]
Only a fixed bounded quasicontinuous multiplier has been passed
through this controlled family of measures. No arbitrary
Borel-multiplier continuity is used.

It follows that $(Dv)^n(H\cap\{g>\delta\})=0$ for every
$\delta>0$. Since this holds for every normalized $v$ on
$\Omega$, the raw capacity of this set is zero. More usefully,
the argument applies to every compact core in $W$. Exhausting
$W$ by compact sets and using continuity from below of the raw
capacity shows that $W\cap\{g>\delta\}$ has raw capacity zero
relative to $\Omega$. This set is quasiopen relative to $\Omega$:
off an arbitrarily small ordinary exceptional open, $g$ is
continuous on a neighborhood of $\overline W$, and $W$ is open.
Proposition~\ref{prop:r2-quasiopen-outer} thus gives zero outer
capacity as well. Taking the countable union over positive
rational $\delta$ proves (i). The all-domain local
formulation follows by repeating the argument in every
requested norm box, or by the framed-domain comparison
Theorem~\ref{thm:r2-domain-capacity-comparison}.
\end{proof}

\begin{corollary}[The exact divisorial form of the no-jump problem]
\label{cor:r2-divisorial-nojump-reduction}
Let $u_j\in\Cb(U)$ increase and be locally uniformly bounded.
Put
\[
 h=\sup_j u_j,\qquad U_h=h^*.
\]
The preceding results give local ambient quasicontinuity of
$h$, and $U_h\in\Cb(U)$. The assertion
\[
 h=U_h\quad\hbox{locally }
       \operatorname{Cap}_{[1]}^*\hbox{-quasi-everywhere}
\]
is equivalent to the assertion
\[
 \sup_j u_j(z)
   =\limsup_{x\to z}\sup_j u_j(x)
       \quad\hbox{for every ambient divisorial }z.
\]
When the sequence has a common compact collar, these conditions
are also equivalent to the energy no-jump condition proved in
Proposition~\ref{prop:r2-envelope-energy-criterion}.
\end{corollary}

\begin{proof}
The actual increasing-limit theorem gives the required
quasicontinuity of $h$; the bounded upper-envelope theorem
gives that of $U_h$. Apply the proposition to the pair, using
$h\leq U_h$. For the energy equivalence, the earlier
common-collar criterion first gives quasi-everywhere equality
for its fixed capacity domain. The atomic testers on that
same domain force equality at each divisorial point. The
present proposition then upgrades this to the all-domain
local quasi-everywhere assertion. Conversely that local
assertion implies the fixed-domain one on the compact support
of the common-collar difference by a finite cover and domain
monotonicity. Thus it implies the earlier energy criterion as well.
\end{proof}

\begin{remark}[Why Shilov control alone did not close the problem]
For a fixed psh function, the Shilov maximum principle controls
its values on an affinoid by its Shilov boundary. The one-Shilov
domain constructed at a divisorial point in this record need
not contain that point in its ordinary interior. Thus that
domain alone does not control the limsup over all ordinary
directions appearing in the displayed equality. A uniform
upper-control statement for bounded psh families at divisorial
points is needed in addition to this reduction. It is supplied
by Theorem~\ref{thm:r2-divisorial-upper-control}; the resulting
no-jump conclusion is Theorem~\ref{thm:r2-increasing-no-jump}.
\end{remark}


% END INCLUDED SOURCE: base_r2/sections/divisorial_determination


% BEGIN INCLUDED SOURCE: base_r2/sections/divisorial_no_jump
\section{Divisorial barriers and the absence of an upper-envelope jump}
\label{sec:r2-divisorial-no-jump}

The argument below supplies the missing ordinary-neighborhood control
at a divisorial point. It uses a local model potential whose
Monge--Amp\`ere measure has only one atom on that neighborhood. The
potential need not vanish only at that point: its zero set is an
affinoid domain. The proof uses the already established pointwise
ambient domination principle, not an assertion that negligible sets
are pluripolar.

All model constructions are over the fixed field $K$, with its original
absolute value and integration calibration. Write $R=K^\circ$.
Throughout this section $n\geq1$.

\begin{lemma}[An algebraic affine model inside a prescribed neighborhood]
\label{lem:r2-algebraic-one-shilov}
Let $Y/K$ be smooth integral affine of dimension $n$, let
$O\subset Y^{\an}$ be ordinary open, and let $z\in O$ be divisorial.
There is a normal, integral, flat, finite-type $R$-algebra $B$ with
the following properties. Its generic fiber
$Y_B=\operatorname{Spec}(B\otimes_RK)$ is a smooth affine algebraic
open in $Y$. The strict affinoid generic fiber
\[
 A=\bigl(\operatorname{Spf}\widehat B\bigr)_\eta
       \subset Y_B^{\an}
\]
of its uniformizer-adic completion is contained in $O$ and contains
$z$. The special fiber of $\operatorname{Spec}B$ has exactly one
irreducible component, and its generic point is the reduction of
$z$. In particular $\Gamma(A)=\{z\}$.
\end{lemma}

\begin{proof}
Choose a compact polynomial norm box $C$ with
\[
 z\in\operatorname{Int}(C),\qquad C\subset O.
\]
After clearing rational powers in its thresholds, write its finitely
many defining inequalities as
\[
 |a_i|\leq1\quad(1\leq i\leq s)
\]
on an affine algebraic open $Y_0\subset Y$. Here the $a_i$ are
regular on $Y_0$: an upper constraint uses a ratio $F^N/c$, and a
lower constraint uses $d/G^M$ after removing $G=0$. Bounds on affine
coordinates are included, so the displayed locus is the compact
box $C$. Such thresholds can be chosen strictly around the values
at $z$, since rational powers of $|K^\times|$ are dense in
$\mathbf R_{>0}$. In particular $|a_i(z)|\leq1$.

Take a normal projective compactification $\overline Y$ of $Y$.
Proposition~A.9 of \cite{GM19} realizes $z$ on a proper algebraic
$R$-model of $\overline Y$. Normalize that model. Normalization is
finite because the model is of finite type over the excellent DVR
$R$; it is an isomorphism on the normal generic fiber. The center
of $z$ is still the generic point $\eta_E$ of a special-fiber
component $E$. Equivalently, one can use the strict transform of
the component on a normal dominating model. Its height-one local
ring is the valuation ring of the same divisorial valuation.

For each $i$, the inequality $|a_i(z)|\leq1$ says exactly that
the rational function $a_i$ belongs to the DVR
$\mathcal O_{\mathscr Y,\eta_E}$. It is therefore regular on a
Zariski neighborhood of $\eta_E$. Intersect these finitely many
neighborhoods. Remove the other special-fiber components. Also
remove the closure of the generic boundary
$\overline Y\setminus Y_0$. This last removal does not remove
$\eta_E$. Indeed, the ideal $I$ of the schematic closure is
uniformizer-saturated. At the DVR
$\mathcal O_{\mathscr Y,\eta_E}$ its localization is nonzero,
since the boundary does not contain the generic point of
$\overline Y$. Any nonzero ideal of that DVR contains a power
of the uniformizer of $R$. Saturation then forces
$I_{\eta_E}=\mathcal O_{\mathscr Y,\eta_E}$, so the closure
does not contain $\eta_E$.

Choose an affine open $\operatorname{Spec}B$ containing $\eta_E$
inside the resulting open. Then $B$ is normal, flat and finite type
over $R$, its special fiber has only the component coming from
$E$, and $B\otimes_RK$ defines a smooth affine open $Y_B$ of
$Y_0$. All $a_i$ belong to $B$. Thus $|a_i|\leq1$ on the
generic fiber of $\operatorname{Spf}\widehat B$, so $A\subset C$.
It contains $z$ because its reduction is $\eta_E\in\operatorname{Spec}B$.

Normality supplies the unique point above this special-fiber
generic point, by \cite[Proposition~6.9.7]{CLD12}; the completion
has the corresponding relatively normal formal model, as $R$ is
excellent. The affine formal-model description of the Shilov
boundary in \cite[Proposition~A.3]{GM19} now gives
$\Gamma(A)=\{z\}$. No generic modification has changed $Y_0$:
the only normalization above is an isomorphism on the already
normal generic compactification. In particular $A\subset O$ is
an actual domain of the originally specified analytic space.
\end{proof}

\begin{lemma}[A local model barrier with a single divisorial atom]
\label{lem:r2-local-dirac-barrier}
Under the hypotheses of Lemma~\ref{lem:r2-algebraic-one-shilov},
there are a smooth affine algebraic open $Y_B\subset Y$, a
continuous semipositive PL function $\rho\in\Cb(Y_B^{\an})$,
and a real number $\epsilon>0$ such that
\[
 \rho\geq0,\qquad \rho(z)=0,\qquad
 \Omega=\{\rho<\epsilon\},\qquad
 \{\rho\leq2\epsilon\}\Subset O\cap Y_B^{\an},
\]
and, on the ordinary open neighborhood $\Omega$ of $z$,
\[
 (D\rho)^n=c_z\delta_z\qquad(c_z>0).
\]
Every sublevel $\{\rho\leq t\}$ is compact in $Y_B^{\an}$.
All statements about the measure are ordinary CLD-current
statements, with the accepted continuous BT compatibility.
\end{lemma}

\begin{proof}
Use $B$ from the preceding lemma and choose finitely many
$R$-algebra generators $g_1,\ldots,g_N$. Their closed embedding
\[
 \operatorname{Spec}B\hookrightarrow\mathbf A_R^N
\]
identifies its formal generic fiber with
\[
 A=\{|g_i|\leq1\text{ for every }i\}\subset Y_B^{\an}.
\]
Define
\[
 \rho=\log\max\{1,|g_1|,\ldots,|g_N|\}.
\]
The function is finite, continuous and classical after all complete
extensions: pullback to a curve is a finite maximum of logarithms
of analytic functions and the constant zero. It is semipositive PL.
Its zero set is exactly $A$. Closed sublevels are intersections of
the closed affine variety $Y_B^{\an}$ with compact closed
polydiscs, hence are compact.

Take the projective closure of $\operatorname{Spec}B$ in
$\mathbf P_R^N$, and normalize it. Denote this normal projective
flat model by $\mathscr Z$. Its normalization is finite, and its
standard affine chart $X_0\ne0$ is still $\operatorname{Spec}B$
because $B$ is normal. Let
\[
 \mathscr L=\nu^*\mathcal O_{\mathbf P_R^N}(1).
\]
It is ample and generated, and its model weight in the frame
$X_0$ on $Y_B$ is exactly $\rho$, with $g_i=X_i/X_0$.
The generic fiber is normal and generically smooth over the
characteristic-zero field $K$. The original CLD model formula
therefore applies, since the discretely valued field $K$ is stable;
see \cite[Theorem~6.9.3, \S6.9.6 and Proposition~6.9.7]{CLD12}.
It gives the finite measure
\[
 c_1(\mathscr L)^n
   =\sum_F b_F\deg_{\widetilde K,F}
                    ((\mathscr L|_F)^n)\,\delta_{z_F}.
\]
Here $F$ runs over special-fiber components, $b_F>0$ is its
multiplicity, and the degree retains its structural residue-field
degree. Ampleness makes every displayed coefficient positive.
The points $z_F$ lie in the analytification of every nonempty
generic algebraic open, so $\rho(z_F)$ is finite.

We determine exactly which of these atoms have height zero.
At the generic point of $F$, choose a local frame $e$ of
$\mathscr L$. The sections $X_i/e$ are regular in the
corresponding DVR, and at least one is a unit because the sections
generate $\mathscr L$. Consequently
\[
 \rho(z_F)
  =\log\max_i|X_i/X_0|(z_F)
  =-\log|X_0/e|(z_F).
\]
This equals zero precisely when $X_0/e$ is a unit, that is,
when the generic point of $F$ belongs to the affine chart
$X_0\ne0$. That chart has exactly one special-fiber component.
Its associated point is $z$. Thus $z$ is the only height-zero
atom, and every other atom has strictly positive height.

Choose $2\epsilon$ smaller than the finitely many positive
heights, if there are any. Shrink $\epsilon$ further so that
$\{\rho\leq2\epsilon\}\subset O$: the compact sublevels
decrease to $A\subset O$, so compactness gives this inclusion.
They are compact subsets of the open $O\cap Y_B^{\an}$, as
required. The restriction of the displayed model measure to
$\Omega=\{\rho<\epsilon\}$ is $c_z\delta_z$, where
\[
 c_z=b_E\deg_{\widetilde K,E}((\mathscr L|_E)^n)>0.
\]
On the smooth open $Y_B$ its model curvature measure is the
ordinary continuous BT product by the accepted compatibility
from R and B. This proves the claimed identity against all
original compact smooth scalar tests. No equation with a
prescribed mass has been solved, and no change of field or
normalization has been made.
\end{proof}

\begin{theorem}[Uniform upper control at a divisorial point]
\label{thm:r2-divisorial-upper-control}
Use $\rho,\epsilon,\Omega$ from
Lemma~\ref{lem:r2-local-dirac-barrier}. Given real numbers
$m\leq M$, every $u\in\Cb(\Omega)$ satisfying
\[
 u\leq M\text{ on }\Omega,\qquad u(z)\geq m
\]
obeys
\[
 u(x)\leq u(z)+A\rho(x)\quad(x\in\Omega),
 \qquad A=\frac{2(M-m+1)}\epsilon.
\]
In particular, a family with these same bounds is uniformly
upper semicontinuous at $z$ in the ordinary topology: for each
$\delta>0$ the ordinary neighborhood
$\Omega\cap\{A\rho<\delta\}$ works simultaneously for
every member of the family.
\end{theorem}

\begin{proof}
Put $v=u(z)+A\rho$. This is continuous classical psh and
\[
 (Dv)^n=A^n c_z\delta_z.
\]
On $\{\rho\geq\epsilon/2\}\cap\Omega$ one has
\[
 v\geq m+A\epsilon/2=M+1>u.
\]
Thus for every $\delta>0$,
\[
 \overline{\{v<u-\delta\}}
     \subset\{\rho\leq\epsilon/2\}\Subset\Omega.
\]
This verifies the compact-gap condition at every part of the
ordinary boundary, without treating the discontinuous strict
gap as an ordinary open set. At $z$ one has $v(z)=u(z)$, so
\[
 (Dv)^n(\{v<u\})=0.
\]
The domain $\Omega$ has compact closure in the smooth affine
chart $Y_B$, and therefore satisfies the framed affine
hypothesis of Corollary~\ref{cor:r2-pointwise-ambient}.
Apply that proved pointwise domination principle to the pair
$(v,u)$. It gives $v\geq u$ everywhere. Since $\rho$ is
continuous and vanishes at $z$, the simultaneous-neighborhood
assertion follows.
\end{proof}

The two geometric lemmas are direct consequences of model
realization, normality and the model-measure formula, with
the neighborhood construction made explicit. The uniform upper
control is recorded as \texttt{NEW\_PROVED} relative to the
previous ambient domination theorem. Its proof uses neither
the upper-envelope identification nor a polar-set theorem.

\begin{corollary}[Upper regularization preserves divisorial values]
\label{cor:r2-divisorial-no-jump}
Let $U$ have the allowed analytic hypotheses. Let
$\mathcal F\subset\Cb(U)$ be nonempty and locally uniformly
bounded above, and put
\[
 h=\sup_{u\in\mathcal F}u,\qquad U_h=h^*.
\]
Then
\[
 U_h(z)=h(z)\qquad\text{at every ambient divisorial }z\in U.
\]
The family can be uncountable; no measurability of $h$ is used
in this assertion.
\end{corollary}

\begin{proof}
Fix $z$ and pass to an ordinary smooth affine algebraic
neighborhood, using the accepted B, Lemma D.4. Choose
$u_0\in\mathcal F$ and replace every $u$ by $\max(u,u_0)$.
The supremum is unchanged. The new family is uniformly bounded
above and below on a compact neighborhood of $z$, since
$u_0$ is locally bounded. Apply the preceding theorem on a
smaller $\Omega$, with common constants $m,M,A$ and the
fixed barrier $\rho$. Taking the supremum gives
\[
 h(x)\leq h(z)+A\rho(x)\quad(x\in\Omega).
\]
Ordinary upper regularization at $z$, where $\rho(z)=0$,
yields $U_h(z)\leq h(z)$. The reverse inequality holds by
definition. All algebraic charts and barriers here are over
the original $K$; no geometric-integrality field change is
needed for this step.
\end{proof}

\begin{theorem}[Increasing envelopes have no ambient capacity jump]
\label{thm:r2-increasing-no-jump}
Let $u_j\in\Cb(U)$ increase and be locally uniformly bounded
above. Set $h=\sup_j u_j$ and $U_h=h^*$. Then
\[
 U_h\in\Cb(U),\qquad
 h=U_h\quad\text{locally outside sets of outer }
           \operatorname{Cap}_{[1]}\text{-capacity zero}.
\]
Moreover $u_j\to U_h$ locally in outer ambient capacity on
every permitted norm-box domain. If the sequence has one
common compact collar, its ambient energy satisfies
\[
 \mathcal E_b(U_h)=\sup_j\mathcal E_b(u_j).
\]
\end{theorem}

\begin{proof}
The established bounded upper-envelope theorem gives
$U_h\in\Cb(U)$. Actual increasing convergence gives local
ambient quasicontinuity of the Borel function $h$, and
quasicontinuity of $U_h$ is already known. The preceding
corollary identifies their values at every ambient divisorial
point. Proposition~\ref{prop:r2-divisorial-capacitary-order}
therefore gives the asserted local outer-capacity-zero equality.
Combine it with actual increasing convergence to $h$ and
subadditivity of outer capacity to obtain convergence to $U_h$.
The common-collar energy assertion is then
Proposition~\ref{prop:r2-envelope-energy-criterion}, for the
ambient background $[1]$.
\end{proof}

\begin{theorem}[Negligible upper-envelope sets have ambient outer capacity zero]
\label{thm:r2-negligible-outer-zero}
Let $\mathcal F\subset\Cb(U)$ be nonempty and locally uniformly
bounded above. Write
\[
 h=\sup_{u\in\mathcal F}u,\qquad U_h=h^*.
\]
Then the set $\{h<U_h\}$ has locally zero outer
$\operatorname{Cap}_{[1]}$-capacity, whether or not the raw
supremum $h$ is Borel. Locally, a countable family of finite
maxima of members of $\mathcal F$ has an increasing actual
supremum $H$ with
\[
 H\leq h\leq U_h,
 \qquad H=U_h\quad\text{locally ambient-quasi-everywhere}.
\]
\end{theorem}

\begin{proof}
Work on one permitted local norm-box domain $W$. The local
capacitary-envelope theorem, with background $[1]$, supplies
an increasing sequence of finite maxima of members of
$\mathcal F|_W$ whose actual Borel supremum $H$ belongs to
$\mathcal A_{[1],b}(W)$. It satisfies $H\leq h$ pointwise,
and every $u\in\mathcal F$ satisfies $u\leq H$ locally
ambient-quasi-everywhere.

Fix an ambient divisorial $z\in W$. Its singleton has
positive capacity on each sufficiently small framed affine
domain containing it. Thus the exceptional set for each
individual inequality $u\leq H$ cannot contain $z$.
Consequently $u(z)\leq H(z)$ for every $u\in\mathcal F$,
even though their exceptional sets may depend on $u$.
Taking the supremum and using $H\leq h$ gives
\[
 H(z)=h(z)=U_h(z),
\]
where the last equality is
Corollary~\ref{cor:r2-divisorial-no-jump}.

Both $H$ and $U_h$ are locally bounded, Borel and locally
ambient-quasicontinuous. Proposition~\ref{prop:r2-divisorial-capacitary-order}
gives $H=U_h$ locally outside a set of outer capacity zero.
The pointwise sandwich implies
\[
 \{h<U_h\}\subset\{H<U_h\}.
\]
Outer capacity is defined for arbitrary sets and is monotone,
so no Borel measurability of the left side is needed. This
proves the theorem locally on every requested domain.
\end{proof}

\paragraph{Scope and dependencies.}
The upper-envelope conclusions in this section are
\texttt{NEW\_PROVED} consequences of the model barrier and
the previously proved ambient calculus. They concern the
ambient capacity $\operatorname{Cap}_{[1]}$; the assertion
for unrestricted positive closed backgrounds remains false,
as the subvariety-supported example already shows. Zero
ambient outer capacity has been proved here. A converse
identification of all zero-capacity sets with pluripolar sets
is not used or asserted.

The proof order is model realization and the original model
formula, then the established pointwise ambient domination,
then uniform divisorial upper control. The quasicontinuous
divisorial criterion and actual increasing convergence are
used only after that control is proved. Thus no assertion
about negligible sets is inserted into its own proof.

% END INCLUDED SOURCE: base_r2/sections/divisorial_no_jump


% BEGIN INCLUDED SOURCE: base_r2/sections/ambient_completion_representative
\section{Classical representatives of the ambient capacitary closure}
\label{sec:r2-ambient-completion-representative}

The completion retains actual Borel functions. The result here
identifies their classes modulo ambient outer-capacity-zero sets;
it does not identify all of their actual point values. It is
specific to the ambient background $[1]$.

\begin{corollary}[The unique classical representative]
\label{cor:r2-ambient-classical-representative}
Let $U$ have the allowed analytic hypotheses. Every
$a\in\mathcal A_{[1],b}(U)$ has a unique function
$r_a\in\Cb(U)$ such that
\[
 a=r_a\quad\text{locally outside sets of outer }
          \operatorname{Cap}_{[1]}\text{-capacity zero}.
\]
At each ambient divisorial point $z$ one has $r_a(z)=a(z)$.
If $a\in\mathcal A^M_{[1]}(W)$ on a norm-box domain $W$,
then $|r_a|\leq M$ there. For any defining approximants
$a_j\in\Cb(W)$ with $|a_j|\leq M$, one formula is
\[
 r_a=\lim_{m\to\infty}
           \left(\sup_{j\geq m}a_j\right)^*.
\]
The right side is a decreasing limit of bounded classical psh
functions, and the final actual classical function is independent
of the approximants.
\end{corollary}

\noindent\textbf{Status: \texttt{DIRECT\_COROLLARY}.}
The inputs are divisorial no-jump, the divisorial-capacitary order
criterion, and the already proved existence of capacitary limits.

\begin{proof}
First work on a defining domain $W$ with
$a\in\mathcal A^M_{[1]}(W)$ and approximants $a_j$ as stated.
Set
\[
 v_m=\left(\sup_{j\geq m}a_j\right)^*.
\]
The bounded upper-envelope theorem gives $v_m\in\Cb(W)$,
and $-M\leq v_m\leq M$. The tails decrease, so $v_m$ decreases
to a bounded classical psh function $r$ with the same bound.

Fix an ambient divisorial point $z\in W$ and choose a small
framed affine norm-box capacity domain containing $z$. Its
singleton has positive capacity by the atomic tester lemma.
Local outer-capacity convergence $a_j\to a$, applied to the
compact singleton $\{z\}$, therefore implies
$a_j(z)\to a(z)$. Otherwise one fixed positive error threshold
would put that singleton into infinitely many error sets,
contradicting their capacities tending to zero.
Corollary~\ref{cor:r2-divisorial-no-jump}, applied to each tail,
gives
\[
 v_m(z)=\sup_{j\geq m}a_j(z).
\]
It follows that $r(z)=a(z)$ at every ambient divisorial point.
The function $a$ is locally quasicontinuous by the completion
theorem; $r$ is locally quasicontinuous as a bounded classical
potential. Proposition~\ref{prop:r2-divisorial-capacitary-order}
therefore gives $r=a$ locally ambient-quasi-everywhere.

For uniqueness, suppose $r,s\in\Cb(W)$ have this property.
Positive singleton capacity forces $r(z)=s(z)$ at every
ambient divisorial point. Shilov boundaries of strict affinoid
neighborhoods consist of such points. The actual pointwise
order lemma, Lemma~\ref{lem:r2-shilov-gap}, applied in both
orders, gives $r=s$ at every point of $W$. This also proves
independence of the defining approximants.

For a general local member $a\in\mathcal A_{[1],b}(U)$,
perform the construction on each defining neighborhood.
On an ordinary overlap the two resulting classical functions
are both locally quasi-everywhere equal to $a$, hence are
actually equal by the uniqueness just proved. They glue to
$r_a\in\Cb(U)$ by the classical sheaf property, which retains
the all-extensions curve tests. Its divisorial values and
uniqueness are local and hence hold on $U$.
\end{proof}

\begin{remark}[The representative need not be the usc regularization of the Borel function]
Take a connected ordinary polydisc $W\subset\mathbf A_K^{n,\an}$,
$n\geq1$, containing the $K$-rational origin $o$, and put
$a=\mathbf1_{\{o\}}$. The singleton is locally pluripolar:
the classical function $\log|T_1|$ is not identically $-\infty$
and has a pole there. The already proved ambient polar theorem
therefore makes $\{o\}$ locally outer-capacity zero. Consequently
the constant zero sequence converges to $a$ in local outer
capacity, so $a\in\mathcal A_{[1],b}(W)$ and $r_a=0$.

The point $o$ is closed, so its indicator is upper
semicontinuous and $a^*=a$. This function is not classical psh:
it attains the strict value one at an interior point and is zero
elsewhere, contradicting the classical strong maximum principle
on the connected domain. Thus $r_a$ is not in general $a^*$.
For the actual supremum of an increasing classical sequence,
in contrast, the proved no-jump theorem does identify its
classical representative with that supremum's usc regularization.
This example uses only the ambient polar-zero theorem; it does
not depend on a converse from capacity zero to pluripolarity.
\end{remark}

\begin{corollary}[The completed ambient calculus uses the classical representative]
\label{cor:r2-ambient-completed-calculus}
For $0\leq k\leq n$ and
$a_0,\ldots,a_k\in\mathcal A_{[1],b}(U)$,
\[
 \mathcal T_{[1],k}(a_1,\ldots,a_k)
       =T_{[1],k}(r_{a_1},\ldots,r_{a_k}),
\]
and
\[
 a_0\mathcal T_{[1],k}(a_1,\ldots,a_k)
       =r_{a_0}T_{[1],k}(r_{a_1},\ldots,r_{a_k}).
\]
These are equalities of coefficient-Radon currents on every
original complementary smooth test form. A uniformly bounded
classical sequence that is locally Cauchy in ambient capacity
thus converges locally in outer capacity to a unique actual
classical function, with the established joint weighted BT
continuity whenever all the slots have the stated bounds.
\end{corollary}

\begin{proof}
The completed calculus is independent of changes on local
outer-capacity-zero sets and agrees with the classical calculus
on $\Cb$, by Theorem~\ref{thm:r2-completed-calculus}.
This gives the unweighted identity. The same theorem's
coefficient bounds by outer capacity make every coefficient
measure of this product vanish on the local disagreement set
of $a_0$ and $r_{a_0}$, proving the weighted identity as a
Borel multiplication identity. A local capacity-Cauchy sequence
has a limit in the completed class; the first corollary
identifies its unique classical representative, and zero-capacity
changes of the limit preserve convergence. The weighted
convergence assertion is then the already established completed
calculus theorem.
\end{proof}

No corresponding conclusion for an arbitrary background $S$
is claimed. For example, a background supported on an ambient
pluripolar subvariety can distinguish the actual increasing
limit from its ambient upper regularization. The relative
capacitary closure remains necessary in that setting.

% END INCLUDED SOURCE: base_r2/sections/ambient_completion_representative


% BEGIN INCLUDED SOURCE: base_r2/sections/dirac_dirichlet
\section{A Dirac Dirichlet problem on the constructed divisorial domains}
\label{sec:r2-dirac-dirichlet}

\begin{corollary}[Constant boundary data and one divisorial atom]
\label{cor:r2-dirac-dirichlet}
Let $z,\rho,\epsilon,\Omega$ and $c_z$ be supplied by
Lemma~\ref{lem:r2-local-dirac-barrier}. For every real
$m>0$ and $b\in\mathbf R$, the function
\[
 u_{m,b}=b+\left(\frac m{c_z}\right)^{1/n}
                         (\rho-\epsilon)
\]
is continuous on a neighborhood of $\overline\Omega$,
belongs to $\Cb(\Omega)$, and satisfies
\[
 (Du_{m,b})^n=m\delta_z\quad\text{on }\Omega,
 \qquad u_{m,b}=b\quad\text{on }\partial\Omega.
\]
Here the closure and boundary are ordinary ones in the smooth
affine chart $Y_B^{\an}$. It is the unique bounded classical
solution with this measure and the boundary condition
\[
 \lim_{\substack{x\to y\\x\in\Omega}}w(x)=b
                   \quad(y\in\partial\Omega).
\]
In particular uniqueness holds among classical solutions
continuous on $\overline\Omega$ with boundary value $b$.
\end{corollary}

\noindent\textbf{Status: \texttt{DIRECT\_COROLLARY}.}
The domain and its barrier have been constructed together.
This corollary concerns that family of domains and a single
divisorial atom. It makes no prescribed-domain assertion for
general atomic or diffuse measures.

\begin{proof}
Positive scaling and translation preserve classicality,
and the proved homogeneity of the BT product gives
\[
 (Du_{m,b})^n=\frac m{c_z}(D\rho)^n=m\delta_z.
\]
Continuity is inherited from $\rho$. Since
$\Omega=\{\rho<\epsilon\}$ and $\rho$ is continuous,
each point of its ordinary boundary satisfies
$\rho=\epsilon$. This proves the boundary value.

Suppose $w$ is another solution with the stated boundary
limit. Both $w$ and $u_{m,b}$ have the same finite boundary
limit at every point. For each $\delta>0$, every boundary
point therefore has an ordinary neighborhood on which their
difference has absolute value less than $\delta$, at points
of $\Omega$. The boundary is compact, being a closed subset
of the compact $\overline\Omega$. A finite union of these
neighborhoods covers it. Consequently both strict-gap sets
with margin $\delta$ have closure compactly contained in
$\Omega$. The equality of the Monge--Amp\`ere measures
and Corollary~\ref{cor:r2-pointwise-ambient} now give
$w=u_{m,b}$ at every point.
\end{proof}

% END INCLUDED SOURCE: base_r2/sections/dirac_dirichlet


% BEGIN INCLUDED SOURCE: base_r2/sections/ambient_extremals
\section{Ambient extremals on domains with a bounded negative exhaustion}
\label{sec:r2-ambient-extremals}

This section concerns $S=[1]$. The domain hypothesis is explicit
and weaker than a strictly psh exhaustion. Its conclusions use
the proved divisorial no-jump theorem; none is claimed here for
an arbitrary positive closed background.

\begin{definition}[The exhaustion domains used in this section]
\label{def:r2-exhaustion-domain}
Let $Y/K$ be smooth affine of dimension $n\geq1$. An exhaustion
domain is a connected ordinary open $V\Subset Y^{\an}$ with
nonempty ordinary boundary and a function
\[
 \rho\in\Cb(V)\cap C(\overline V),\qquad
 \rho<0\text{ on }V,\qquad \rho=0\text{ on }\partial V.
\]
In particular $\rho$ is bounded and
$\{\rho\leq-\epsilon\}\Subset V$ for every $\epsilon>0$.
No positivity stronger than classical psh is required of $\rho$,
and no equation for $(D\rho)^n$ is assumed.
\end{definition}

\begin{lemma}[Local supply of exhaustion domains]
\label{lem:r2-local-exhaustion-domains}
Every point in an allowed ordinary algebraic neighborhood has
arbitrarily small exhaustion domains of the preceding type.
This does not assert that every ordinary analytic domain is an
exhaustion domain.
\end{lemma}

\begin{proof}
Inside a prescribed ordinary neighborhood $O$, choose a compact
rational polynomial norm box $C$ containing the point in its
ordinary interior. After passing to an affine algebraic open
$Y_0$ and clearing powers in the rational thresholds, write
\[
 C=\{|a_i|\leq1\text{ for }1\leq i\leq N\}.
\]
The regular functions $a_i$ include the required bounds on
affine coordinates and on inverses used for lower constraints.
Consequently each locus $\{|a_i|\leq e^t\text{ for all }i\}$,
with $t\geq0$, is compact in $Y_0^{\an}$. Equivalently, one
can include scaled powers of a set of affine algebra generators.
Put
\[
 \rho_0=\log\max(1,|a_1|,\ldots,|a_N|).
\]
This is continuous classical psh by the finite maximum rule
after complete extensions and curve pullbacks. Its compact
sublevels decrease to $C$ as $t\downarrow0$. Choose a small
positive rational valuation threshold $\epsilon$ such that
$\{\rho_0\leq2\epsilon\}\subset O$.

Analytic spaces are locally connected, as recalled in
\cite{GR25} from Berkovich's local connectedness theorem.
Let $V$ be the connected component containing the specified
point of $\{\rho_0<\epsilon\}$. It is ordinary open, and its
compact closure is contained in $\{\rho_0\leq\epsilon\}$.
Since a connected component is relatively closed in that
ordinary open, its boundary is contained in
$\{\rho_0=\epsilon\}$. Thus $\rho=\rho_0-\epsilon$ is the
required bounded continuous negative exhaustion.

For completeness the boundary is nonempty. Otherwise $V$
would be a compact open-and-closed subdomain of the strict
affinoid $\{\rho_0\leq2\epsilon\}$, hence an affinoid
component. It is pure of positive dimension and therefore has
a Shilov point on its analytic boundary, by
\cite[Lemma~6.2]{GR25}. This contradicts its being an ordinary
open in the boundaryless space $Y_0^{\an}$ and the
interior-composition rule \cite[\S2.11(3)--(4)]{GR25}.
\end{proof}

The lemma is a \texttt{DIRECT\_COROLLARY} of the ordinary
norm-box geometry and classical maximum calculus already used
in the record. It does not supply strict positivity or peak
barriers at every boundary point.

Fix an exhaustion domain $V$ for the rest of this section and
keep the capacity domain fixed:
\[
 C(E)=\operatorname{Cap}_{[1]}(E;V),\qquad C^*(E)=C_V^*(E).
\]
For any subset $E\subset V$, define
\[
 \mathcal F_E=\{v\in\Cb(V):-1\leq v\leq0,\ v=-1\text{ on }E\},
 \qquad h_E=\sup_{v\in\mathcal F_E}v,\qquad U_E=h_E^*.
\]
The constraints in $\mathcal F_E$ are actual pointwise
constraints, including when $E$ is not Borel. The family
contains the constant $-1$. When $E=\varnothing$, $U_E=0$.

\begin{lemma}[Basic extremal properties and a countable selection]
\label{lem:r2-extremal-basics}
One has $U_E\in\Cb(V)$ and $-1\leq U_E\leq0$.
If $E\subset F$, then $U_F\leq U_E$ pointwise.
The exceptional subset
\[
 E\cap\{U_E>-1\}
\]
has locally zero outer ambient capacity. If $E\Subset V$,
then $U_E$ has ordinary boundary limit zero, and the
exceptional subset has $C^*$-capacity zero in the fixed domain.
There is an increasing sequence $u_j\in\mathcal F_E$ with
\[
 \left(\sup_j u_j\right)^*=U_E,\qquad
 u_j\longrightarrow U_E
   \quad\text{locally in outer ambient capacity}.
\]
\end{lemma}

\begin{proof}
The bounded upper-envelope theorem gives classicality.
Monotonicity follows from inclusion of the admissible families.
Since $h_E=-1$ on $E$, the exceptional subset is contained
in $\{h_E<U_E\}$. The arbitrary-family negligible-set theorem
gives its local outer-capacity zero, without requiring that
$h_E$ or $E$ be Borel. On a relatively compact $E$, a finite
local cover and domain monotonicity give zero for the fixed
domain $V$.

If $E\Subset V$, choose $c>0$ such that $c\rho\leq-1$ on
$\overline E$. The function $\max(-1,c\rho)$ belongs to
$\mathcal F_E$. It tends to zero at every ordinary boundary
point, and
\[
 \max(-1,c\rho)\leq U_E\leq0.
\]
This proves the boundary assertion. The same lower barrier
works simultaneously for all $E$ contained in one fixed
compact subset of $V$.

Here is an elementary countable selection, independent of
Choquet's capacitability theorem. The finite $p$-adic field
makes $V$ second countable. For a countable ordinary basis
$(B_l)$, let
\[
 M_l=\sup\{v(x):v\in\mathcal F_E,\ x\in B_l\}.
\]
For each $l,m$ choose $v_{l,m}\in\mathcal F_E$ and
$x_{l,m}\in B_l$ with $v_{l,m}(x_{l,m})>M_l-1/m$.
Enumerate these functions and take their successive finite
maxima $u_j$. They remain in $\mathcal F_E$.
For $a=\sup_j u_j$, one has
$\sup_{B_l}a=M_l=\sup_{B_l}h_E$ for each $l$.
The neighborhood formula for upper regularization therefore
gives $a^*=h_E^*=U_E$. The proved increasing no-jump theorem
then gives the claimed local outer-capacity convergence.
\end{proof}

\begin{lemma}[Order of total Monge--Amp\`ere masses]
\label{lem:r2-total-mass-order}
Let $u,v\in\Cb(V)$ be bounded, with
$u\leq v\leq0$ and ordinary boundary limits zero. Then
\[
 \int_V(Dv)^n\leq\int_V(Du)^n.
\]
The assertion allows infinite masses and contains no
subtraction of infinities.
\end{lemma}

\begin{proof}
If $v\equiv0$ the assertion is immediate. Otherwise the
strong maximum principle on the connected domain gives
$v<0$ everywhere. For $0<t<1$, put $w=(1-t)v$.
Then $w>v\geq u$ on all of $V$. The boundary limits imply
\[
 \overline{\{u<w-\epsilon\}}\Subset V\qquad(\epsilon>0).
\]
Indeed each such gap is absent on a neighborhood of the
compact ordinary boundary. The comparison principle gives
\[
 (1-t)^n\int_V(Dv)^n
  =\int_{\{u<w\}}(Dw)^n
  \leq\int_{\{u<w\}}(Du)^n
  =\int_V(Du)^n.
\]
Let $t\downarrow0$.
\end{proof}

\begin{theorem}[The compact extremal and its equilibrium mass]
\label{thm:r2-compact-extremal}
For every compact $K\Subset V$, the function $U_K$ has
ordinary boundary limit zero and is locally maximal on
$V\setminus K$. Its measure
\[
 \mu_K=(DU_K)^n
\]
is supported in $K$, is finite, and satisfies
\begin{equation}
\label{eq:r2-compact-capacity-mass}
 C(K)=\mu_K(K)=\mu_K(V).
\end{equation}
The equality $U_K=-1$ on $K$ is asserted outside a
$C^*$-zero subset; it is not asserted at every point of $K$.
\end{theorem}

\begin{proof}
The boundary statement and the zero-capacity exceptional set
are in Lemma~\ref{lem:r2-extremal-basics}. Choose there an
increasing sequence $u_j\in\mathcal F_K$ whose upper-regularized
supremum is $U_K$.

We first prove that the measure vanishes off $K$.
Work in a small ordinary $W\Subset V\setminus K$ on which
the constructive common projective globalization applies,
and fix a nonnegative original smooth scalar $\chi$ compactly
supported in $W$. Globalize the whole bounded family $u_j$
on one identified ample line, obtaining bounded classical
metrics $R_j$ whose weights on a neighborhood of
$\operatorname{supp}\chi$ are $u_j+c$ with one fixed $c$.
The construction is a maximum with one common collar, so
the metrics increase and have one common bound.
Extend $\chi$ by zero to the projective space. For a fixed
$\epsilon>0$, put
\[
 p_j=P_L^{\mathrm b}(R_j+\epsilon\chi).
\]
The exact obstacle inequalities give
$R_j\leq p_j\leq R_j+\epsilon\chi$, hence $p_j=R_j$
off $\operatorname{supp}\chi$. Envelope monotonicity makes
$p_j$ increasing.

Its local weight minus $c$ glues with $u_j$ off that support
to a function $w_j\in\Cb(V)$. This modification is compact
and uniformly bounded. It satisfies $w_j\geq u_j\geq-1$,
and has boundary limsup at most zero. The boundary maximum
principle gives $w_j\leq0$. Since the modification is away
from $K$, the constraint $w_j=-1$ on $K$ holds at every
point. Therefore $w_j\in\mathcal F_K$, and
\[
 u_j\leq w_j\leq h_K,\qquad
 \left(\sup_jw_j\right)^*=U_K.
\]
Both increasing rows converge locally in outer capacity
to $U_K$ by the no-jump theorem.

The proved quasicontinuous-obstacle orthogonality gives,
in local weights,
\[
 0=\int_W(u_j+\epsilon\chi-w_j)(Dw_j)^n .
\]
The integrand vanishes off the fixed compact support of
$\chi$. Insert a fixed smooth cutoff equal to one near
that support. Weighted increasing BT convergence sends
the terms with weights $u_j$ and $w_j$ to the same limit,
while weak convergence handles the fixed continuous
weight $\chi$. Consequently
\[
 \epsilon\int_W\chi(DU_K)^n=0.
\]
If the globalization uses B's finite constant-field
convention, its local calibration multiplies this zero
identity by a positive constant, exactly as in the
maximality proof. The conclusion is for the original
$K$-measure. Such $W$ cover $V\setminus K$, and all original
nonnegative compact scalar tests are allowed there.
Thus $\mu_K$ is zero on $V\setminus K$.
It is finite by local CLN on $K$, and local maximality
follows from the proved homogeneous-equation criterion.

This argument deliberately lifts the ACTUAL admissible
functions $u_j$. Directly modifying $U_K$ and treating a
quasi-everywhere obstacle inequality as a pointwise
admissibility condition would not be justified.

Since $U_K+1$ is a normalized test on the original $V$,
$\mu_K(K)\leq C(K)$. For the reverse inequality take
$v\in\Cb(V)$ with $0\leq v\leq1$ and $0<t<1$, and set
\[
 w=(1-t)v-1+t/2.
\]
It lies between $-1+t/2$ and $-t/2$. On $K$ outside the
exceptional zero-capacity set, $U_K=-1<w$.
That exceptional Borel set has zero $(Dv)^n$-mass, because
$(Dv)^n(E)\leq C(E)\leq C^*(E)$.
The boundary limit $U_K=0$ gives the compact-gap condition
for $(U_K,w)$. Comparison and the established support give
\[
 (1-t)^n\int_K(Dv)^n
 \leq\int_{\{U_K<w\}}(Dw)^n
 \leq\int_{\{U_K<w\}}\mu_K
 \leq\mu_K(K).
\]
Take the supremum over normalized $v$ and let $t\downarrow0$.
This proves \eqref{eq:r2-compact-capacity-mass}.
\end{proof}

\begin{theorem}[Compact continuity and compact outer regularity]
\label{thm:r2-compact-capacity-continuity}
If $K_j\Subset V$ are decreasing compact sets with
$K=\bigcap_jK_j$, then
\[
 C(K_j)\downarrow C(K),\qquad
 (\,\sup_j U_{K_j}\,)^*=U_K.
\]
Moreover $U_{K_j}\to U_K$ locally in outer ambient capacity
and their Monge--Amp\`ere measures converge weakly on all
original compact scalar tests. Every compact $K\Subset V$
satisfies
\[
 C^*(K)=C(K).
\]
All these capacities use the same original domain $V$.
\end{theorem}

\begin{proof}
The case of an empty intersection is eventually empty and
is immediate by compactness. Otherwise put
$U=(\sup_jU_{K_j})^*$. Monotonicity gives $U\leq U_K$.
For $v\in\mathcal F_K$ and $\epsilon>0$, upper
semicontinuity makes
$O_\epsilon=\{v<-1+\epsilon\}$ an ordinary open
neighborhood of $K$. Nested compactness gives
$K_j\subset O_\epsilon$ for all sufficiently large $j$.
The function
\[
 v_\epsilon=\max(-1,v-\epsilon)
\]
then belongs to $\mathcal F_{K_j}$. Hence
$v_\epsilon\leq U_{K_j}\leq U$. Letting
$\epsilon\downarrow0$ gives $v\leq U$.
Taking the supremum and upper regularization gives
$U_K\leq U$, proving equality.

The increasing no-jump and BT theorems now give the stated
capacity and measure convergence. Every measure involved
is supported in the one compact set $K_1$. Testing with
a compact smooth cutoff equal to one near $K_1$ therefore
gives convergence of their total masses. Formula
\eqref{eq:r2-compact-capacity-mass} proves compact
continuity of the raw capacity.

Choose decreasing compact neighborhoods $L_j\Subset V$
of $K$ with intersection $K$, using local compactness
and metrizability. Since $K\subset\operatorname{Int}L_j$,
\[
 C(K)\leq C^*(K)
 \leq C(\operatorname{Int}L_j)
 \leq C(L_j)\longrightarrow C(K).
\]
This proves compact outer regularity without interchanging
an infimum of neighborhoods with a supremum of individual
Radon measures.
\end{proof}

\begin{theorem}[Open-set extremals and their total mass]
\label{thm:r2-open-extremal}
Let $G\Subset V$ be ordinary open. Then $U_G=-1$ at every
point of $G$, $U_G$ has ordinary boundary limit zero on
$\partial V$, and
\[
 \operatorname{supp}(DU_G)^n\subset\partial_VG,
 \qquad
 \int_V(DU_G)^n=C(G).
\]
Here $\partial_VG=\overline G\setminus G$ is the ordinary
relative boundary. In particular the mass formula is a
TOTAL-mass formula; $(DU_G)^n(G)=0$.
If $K_m\Subset G$ increase with
$G=\bigcup_m\operatorname{Int}K_m$, then
$U_{K_m}\downarrow U_G$ pointwise.
\end{theorem}

\begin{proof}
The family constraint makes $h_G=-1$ on the ordinary open
$G$, so its upper regularization is also exactly $-1$
there. The exhaustion barrier from
Lemma~\ref{lem:r2-extremal-basics}, using $\overline G$,
gives the boundary limit zero.

Choose the indicated compact exhaustion and put
$U=\lim_m U_{K_m}$. This is a bounded decreasing limit
of classical psh functions and is therefore classical.
For each $x\in G$, one has $x\in\operatorname{Int}K_m$
for some $m$. On that interior $U_{K_m}=-1$ exactly,
by the same open-set regularization observation.
Thus $U=-1$ on $G$, and $U\in\mathcal F_G$.
Conversely $\mathcal F_G\subset\mathcal F_{K_m}$ for
every $m$, so $U_G\leq U_{K_m}$ and hence $U_G\leq U$.
The admissibility of $U$ gives the reverse inequality.

The decreasing BT theorem gives convergence of the measures.
They are all supported in the fixed compact $\overline G$.
A cutoff equal to one near it gives convergence of total
masses. Since the raw capacity is continuous from below,
\[
 \int_V(DU_G)^n
 =\lim_m\int_V(DU_{K_m})^n
 =\lim_m C(K_m)=C(G).
\]
The limit measure is supported in $\overline G$. Ordinary
locality and the actual equality $U_G=-1$ on $G$ make
the measure zero on $G$, proving the boundary support.
\end{proof}

\paragraph{Status and scope.}
The compact and open extremal mass formulas and compact
continuity are \texttt{NEW\_PROVED} on the displayed
exhaustion domains, relative to the reviewed ambient no-jump
and obstacle-orthogonality theorems. The results do not
identify every arbitrary ordinary domain as such a domain,
do not assert a strictly psh exhaustion, and do not transfer
these mass formulas to an arbitrary background $S$.
Compact outer regularity is proved here for the raw ambient
capacity. Full Borel outer regularity and Choquet
capacitability require the additional steps stated in the
subsequent section.

% END INCLUDED SOURCE: base_r2/sections/ambient_extremals


% BEGIN INCLUDED SOURCE: base_r2/sections/polar_converse_choquet
\section{Polar witnesses and capacitability from ambient extremals}
\label{sec:r2-polar-converse-choquet}

\noindent\textbf{Status: \texttt{NEW\_PROVED}, relative to the adopted
inputs and the preceding reviewed R2 theorems.}
The witness construction, arbitrary-set extremal mass formula, and
ordinary-open proof of increasing continuity are supplied below.
The Choquet and patching consequences are
\texttt{DIRECT\_COROLLARY}. The two theorem groups
\textup{(N)} and \textup{(E)} are separated to display dependencies;
both have been proved in the preceding sections. Neither a general
Josefson theorem nor Borel outer regularity is an input.

\subsection{Domain and the two theorem inputs}

Let $V$ be a connected ordinary open with compact closure in a smooth
affine algebraic chart over the fixed finite extension $K/\mathbf Q_p$.
Assume that its ordinary boundary is nonempty and that there is
\[
 \rho\in\Cb(V)\cap C(\overline V),\qquad
 \rho<0\text{ on }V,\qquad \rho=0\text{ on }\partial V.
\]
Thus each negative sublevel of $\rho$ has compact closure in $V$.
All capacities in this section use the ambient background $[1]$ and
the \emph{fixed admissibility domain $V$}. Write
\[
 C(E)=\operatorname{Cap}_{[1]}(E;V),\qquad
 C^*(E)=\inf_{\substack{G\supset E\\G\text{ ordinary open in }V}}C(G).
\]
The first notation is initially for Borel sets; the second is defined
for arbitrary subsets. All Monge--Amp\`ere measures and their weak
limits use the original CLD scalar test functions.

For any subset $A\subset V$, define
\[
 \mathcal F_A=\{v\in\Cb(V):-1\leq v\leq0,\ v=-1\text{ on }A\},
 \qquad h_A=\sup_{v\in\mathcal F_A}v,\qquad U_A=h_A^*.
\]
The ordinary upper regularization $U_A$ is bounded classical by
Theorem~\ref{thm:r2-local-upper-envelope}. The conditions $v=-1$
and $v\leq-1$ are identical for this normalized family. The raw
supremum is never substituted for its regularization without a
stated argument.

The two theorem groups are recorded separately to make the dependency
order explicit:
\begin{enumerate}[label=\textup{(\Alph*)}]
\item[\textup{(N)}] For every locally uniformly upper-bounded family
of bounded classical psh functions, the difference between its actual
supremum and its ordinary usc regularization is contained locally in
a Borel set of ambient outer capacity zero. In particular a locally
uniformly bounded increasing sequence has its actual supremum equal
to its regularization outside such a set, and their values agree at
every ambient divisorial point. Local and fixed-domain null classes
are identified only through
Corollary~\ref{cor:r2-fixed-domain-null-classes}.
These are Theorem~\ref{thm:r2-negligible-outer-zero},
Theorem~\ref{thm:r2-increasing-no-jump}, and
Corollary~\ref{cor:r2-divisorial-no-jump}.

\item[\textup{(E)}] If $G\Subset V$ is ordinary open, then $U_G=-1$
at every point of $G$, $U_G$ has ordinary boundary limit zero,
$\operatorname{supp}(DU_G)^n\subset\overline G$, and
\[
 \int_V(DU_G)^n=C(G).
\]
For a compact $K\Subset V$, the corresponding compact extremal
identity is
\[
 \int_V(DU_K)^n=C(K).
\]
The open-set identity concerns \emph{total mass}. It does not assert
$\int_G(DU_G)^n=C(G)$: the function is constant on $G$.
These are Theorems~\ref{thm:r2-open-extremal} and
\ref{thm:r2-compact-extremal}, with the domain and boundary
properties of Lemma~\ref{lem:r2-extremal-basics}.
\end{enumerate}

For every fixed $A\subset W\Subset V$, choose $a>0$ sufficiently
large that $a\rho\leq-1$ on $\overline W$. Then
\[
 b=\max(-1,a\rho)\in\mathcal F_A,
 \qquad b\leq U_A\leq0.
\]
This gives the ordinary boundary limit zero of $U_A$, uniformly for
all such $A$. In particular $\{U_A<-1+\eta\}\Subset V$ for every
$0<\eta<1$, with a common compact bound for fixed $W$ and $\eta$.
The capacity of $W$ is finite by relative CLN on a slightly larger
compact stage.

We also use the total-mass order of
Lemma~\ref{lem:r2-total-mass-order}: if $u\leq v\leq0$ are bounded
classical functions with ordinary boundary limits zero, then
$\int_V(Dv)^n\leq\int_V(Du)^n$, including infinite masses.

\subsection{A nontrivial pole witness for an outer-capacity-zero set}

\begin{theorem}[A witness on the specified exhaustion domain]
\label{thm:r2-capacity-zero-polar-witness}
Let $E\subset W\Subset V$ be an arbitrary subset with $C^*(E)=0$.
For every fixed ambient divisorial point $z\in V$ there is a nonpositive
extended classical psh function $\phi$ on $V$ such that
\[
 E\subset\{\phi=-\infty\},\qquad -1\leq\phi(z)\leq0.
\]
In particular the witness is not identically $-\infty$.
\end{theorem}

\begin{proof}
Choose an ordinary $W'$ with
$\overline W\subset W'\Subset V$. From $C^*(E)=0$ choose ordinary
opens $O_j\supset E$ with $C(O_j)<2^{-j}$ and set
\[
 G_j=W'\cap\bigcap_{i=1}^jO_i,\qquad u_j=U_{G_j}.
\]
The $G_j$ decrease, while $-1\leq u_j\uparrow h\leq0$.
They share a lower barrier $b=\max(-1,a\rho)$ as above.
Their measures are supported in the one compact set
$\overline W'$, and their total masses tend to zero by \textup{(E)}.

Put $U=h^*\in\Cb(V)$. By \textup{(N)}, $h=U$ locally outside
sets of outer capacity zero. The actual increasing-limit theorem
\ref{cor:r2-local-increasing}, the completed calculus
\ref{thm:r2-completed-calculus}, and independence under such changes
of representatives give
\[
 (Du_j)^n\longrightarrow(DU)^n
\]
on all original compact scalar tests. Their vanishing total masses
imply $(DU)^n=0$. Also $b\leq U\leq0$, so $U$ has ordinary
boundary limit zero. Pointwise ambient domination
\ref{cor:r2-pointwise-ambient}, applied to $U$ and $0$, now gives
$U=0$ on $V$. Indeed all sets $\{U<-\epsilon\}$ are relatively
compact by the lower barrier, and $(DU)^n$ vanishes on the whole
strict gap.

By the divisorial part of \textup{(N)}, $u_j(z)\uparrow0$.
Choose a strictly increasing sequence $j_l$ such that
$u_{j_l}(z)>-2^{-l}$. The finite sums
\[
 \phi_m=\sum_{l=1}^m u_{j_l}
\]
are nonpositive classical psh and decrease with $m$. Their pointwise
limit $\phi$ is extended classical psh by the classical decreasing
closure, and $\phi(z)\geq-\sum_l2^{-l}=-1$.
For $x\in E$, every $G_j$ contains $x$ and the open-set extremal
satisfies $u_j(x)=-1$ exactly. Thus $\phi_m(x)=-m$ and
$\phi(x)=-\infty$. No integrability of the original set or a
global pole-patching theorem has been used.
\end{proof}

\begin{corollary}[The local zero-set equivalence on this domain]
\label{cor:r2-local-polar-equivalence}
On the same exhaustion domain, for every arbitrary $E\subset W\Subset V$,
\[
 C^*(E)=0
 \quad\Longleftrightarrow\quad
 E\subset\{\phi=-\infty\}
 \text{ for some nonpositive classical }\phi\text{ on }V,
 \ \phi\not\equiv-\infty.
\]
The implication from right to left is the already proved ambient
polar theorem, not a new nonconcentration hypothesis.
\end{corollary}

\begin{proof}
The forward implication is the theorem. For the reverse implication,
Corollary~\ref{cor:r2-ambient-polar-outer-zero} gives local outer
capacity zero on every compact/domain pair. The fixed-domain
comparison \ref{cor:r2-fixed-domain-null-classes} gives $C^*(E)=0$.
\end{proof}

\begin{definition}[Local negligibility in the ordinary topology]
\label{def:r2-local-negligibility}
Let $E$ be an arbitrary subset of an allowed ordinary open $U$.
It is locally negligible if every $x\in U$ has an ordinary
neighborhood $O\subset U$ and a nonempty locally uniformly
upper-bounded family $\mathcal F\subset\Cb(O)$ such that
\[
 E\cap O\subset
 \left\{\sup_{v\in\mathcal F}v<
                    \left(\sup_{v\in\mathcal F}v\right)^*\right\}.
\]
It is locally increasing-negligible if the family can be chosen as
one increasing locally uniformly bounded sequence in $\Cb(O)$.
Local pluripolarity means that every $x\in U$ has a connected ordinary
neighborhood $O\subset U$ and an extended classical psh function
$\phi$ on $O$, not identically $-\infty$ there, with
$E\cap O\subset\{\phi=-\infty\}$.
Local outer-capacity zero means that
\[
 \operatorname{Cap}_{[1]}^*(E\cap H;\Omega)=0
\]
for every permitted ordinary norm-box domain $\Omega\subset U$
and every compact $H\Subset\Omega$.
\end{definition}

\begin{corollary}[Local polar, negligible, and outer-zero sets]
\label{cor:r2-polar-negligible-equivalence}
On every allowed ordinary open $U$, the following properties of an arbitrary subset
$E\subset U$ are equivalent:
\begin{enumerate}[label=\textup{(\roman*)}]
\item $E$ is locally pluripolar;
\item $E$ has local ambient outer-capacity zero;
\item $E$ is locally increasing-negligible;
\item $E$ is locally negligible.
\end{enumerate}
These are local witness assertions. A single global psh witness on
$U$ is not asserted.
\end{corollary}

\begin{proof}
The implication (i)$\Rightarrow$(ii) is
Corollary~\ref{cor:r2-ambient-polar-outer-zero}.
For (ii)$\Rightarrow$(i), fix $x\in U$ and use
Lemma~\ref{lem:r2-local-exhaustion-domains} in an ordinary algebraic
neighborhood to choose an exhaustion domain $V$ containing $x$.
Choose a connected ordinary $W$ with $x\in W\Subset V$ and an
ambient divisorial point $z\in W$. The local null-class comparison
gives $\operatorname{Cap}_{[1]}^*(E\cap W;V)=0$.
Theorem~\ref{thm:r2-capacity-zero-polar-witness} supplies a
nonpositive classical witness $\phi$ on $V$, finite at $z$, whose
pole set contains $E\cap W$. Its restriction to $W$ is a
nontrivial classical witness, as required.

For (i)$\Rightarrow$(iii), shrink a connected witness neighborhood
if necessary and subtract a finite upper bound so that $\phi\leq0$.
The pole set $P=\{\phi=-\infty\}$ has empty ordinary interior by
Proposition~\ref{prop:r2-no-interior-poles}; hence the witness remains
nontrivial on every smaller connected ordinary neighborhood.
The actual bounded classical functions
\[
 w_N=1+\frac{\max(\phi,-N)}{N}
\]
increase to $\mathbf1_{P^c}$. Because $P$ has empty ordinary
interior, the ordinary usc regularization of this actual limit is
the constant one. Its discrepancy set is exactly $P$, which
contains the specified part of $E$. This proves increasing
negligibility. The implication (iii)$\Rightarrow$(iv) is immediate.

Finally (iv)$\Rightarrow$(ii) follows from
Theorem~\ref{thm:r2-negligible-outer-zero}. To express it on a
prescribed compact/domain pair $(H,\Omega)$, restrict the local
witness families to smaller permitted domains contained in $\Omega$.
Choose countably many ordinary cores covering $E\cap H$; such a
choice exists in the locally second-countable finite $p$-adic
setting. Each core intersection has outer capacity zero on its
own domain and therefore on $\Omega$ by domain monotonicity.
Countable subadditivity gives the required zero capacity on $H$.
No uncountable union of exceptional sets is taken.
\end{proof}

\begin{corollary}[Pole patching on one exhaustion domain]
\label{cor:r2-exhaustion-domain-pole-patching}
On the same connected exhaustion domain $V$, let $E\subset V$
have local ambient outer-capacity zero. In particular this applies
to a locally pluripolar $E$. Then there is one nonpositive extended
classical psh function $\phi$ on $V$, not identically $-\infty$,
such that $E\subset\{\phi=-\infty\}$. No relative compactness
of $E$ itself is required for this conclusion.
\end{corollary}

\noindent\textbf{Provenance: \texttt{DIRECT\_COROLLARY}.}
The domain $V$ and its exhaustion hypothesis are retained.

\begin{proof}
Choose ordinary $W_m\Subset V$ increasing to $V$ and fix one
ambient divisorial point $z\in V$. Put $E_m=E\cap W_m$.
The fixed-domain null-class comparison gives $C^*(E_m)=0$.
The witness theorem gives nonpositive classical $\phi_m$ on $V$
with $E_m\subset\{\phi_m=-\infty\}$ and $\phi_m(z)\geq-1$.
The series
\[
 \phi=\sum_{m\geq1}2^{-m}\phi_m
\]
is a decreasing limit of nonpositive classical partial sums and
is therefore extended classical. It satisfies $\phi(z)\geq-1$,
so it is nontrivial. Every point of $E$ belongs to some $E_m$,
and the corresponding positive summand forces $\phi=-\infty$
there.
\end{proof}

\begin{remark}[Why the domain qualification is essential]
Untwisted pole globalization on every analytic space is
\texttt{FALSE}. For example, the rational point $0$ of
$\mathbf P^{1,\an}_K$ is locally pluripolar, with local witness
$\log|t|$. A nontrivial global extended classical psh function
on this proper connected space attains a finite maximum and is
constant by the strong maximum principle
\cite[Theorem 6.1]{GR25}. It therefore has no
pole. The preceding corollary is globalization on the one
specified exhaustion domain; it is not that false unrestricted
assertion, nor a theorem about globally twisted psh functions.
\end{remark}

\subsection{An outer-capacity formula for arbitrary relatively compact sets}

\begin{theorem}[The extremal total-mass formula]
\label{thm:r2-arbitrary-extremal-mass}
For every arbitrary $A\subset W\Subset V$,
\[
 C^*(A)=\int_V(DU_A)^n<\infty.
\]
This assertion neither assumes Borel measurability of $A$ nor
identifies the raw capacity with its outer regularization beforehand.
\end{theorem}

\begin{proof}
If $G\Subset V$ is ordinary open and contains $A$, then
$\mathcal F_G\subset\mathcal F_A$, hence $U_G\leq U_A\leq0$.
Both functions have boundary limit zero. The mass-order lemma and
\textup{(E)} imply
\[
 \int_V(DU_A)^n\leq\int_V(DU_G)^n=C(G).
\]
Taking the infimum over such $G$ gives
$\int_V(DU_A)^n\leq C^*(A)$. Relatively compact opens suffice here:
intersect any open cover with a fixed relatively compact open
neighborhood of $\overline W$. This also proves finiteness.

Since $h_A=-1$ on $A$, input \textup{(N)} shows that
\[
 A\cap\{U_A>-1\}\subset N
\]
for some set $N\subset W$ with $C^*(N)=0$. We may use a Borel
outer-zero cover of this defect if $A$ itself is not Borel.
Fix $0<\eta<1$ and put
\[
 G_\eta=\{U_A<-1+\eta\},\qquad
 v_\eta=\max\left(-1,\frac{U_A}{1-\eta}\right).
\]
The set $G_\eta$ is ordinary open and relatively compact, because
$U_A$ is usc and bounded below by the common exhaustion barrier.
The function $v_\eta$ is classical, has boundary value zero,
belongs to $[-1,0]$, and equals $-1$ on $G_\eta$. Consequently
\[
 \frac{U_A}{1-\eta}\leq v_\eta\leq U_{G_\eta}.
\]
Twice applying the mass-order lemma gives
\[
 C(G_\eta)=\int_V(DU_{G_\eta})^n
 \leq\int_V(Dv_\eta)^n
 \leq(1-\eta)^{-n}\int_V(DU_A)^n.
\]
Finally $A\subset G_\eta\cup N$. Monotonicity and subadditivity
of outer capacity, and its equality with raw capacity on ordinary
open sets, imply
\[
 C^*(A)\leq C(G_\eta)
 \leq(1-\eta)^{-n}\int_V(DU_A)^n.
\]
Let $\eta$ decrease to zero. The two inequalities prove the formula.
\end{proof}

\subsection{The missing continuity axiom and the Choquet conclusion}

\begin{theorem}[Continuity from below on arbitrary subsets of a core]
\label{thm:r2-outer-capacity-increasing-sets}
If $A_j\subset A_{j+1}\subset W\Subset V$ are arbitrary subsets
and $A=\bigcup_jA_j$, then
\[
 C^*(A)=\lim_j C^*(A_j).
\]
\end{theorem}

\begin{proof}
Write $U_j=U_{A_j}$; these bounded classical functions decrease
with $j$. By \textup{(N)}, each defect
$A_j\cap\{U_j>-1\}$ is contained in an outer-zero set $N_j$.
The countable union of the $N_j$ has outer capacity zero.
For $\epsilon>0$, choose one ordinary open set $G$ containing
that union with $C(G)<\epsilon$.

Fix $0<\eta<1$ and put
\[
 O_j=G\cup\{U_j<-1+\eta\}.
\]
These are \emph{increasing ordinary open sets}, since the $U_j$
decrease. Each $O_j$ contains $A_j$. The truncation estimate in
the preceding proof and its mass formula give
\[
 C(O_j)\leq\epsilon+
      (1-\eta)^{-n}C^*(A_j).
\]
The raw capacity is continuous from below on Borel sets, directly
from its definition as a supremum of Radon measures. Applying
this only to the increasing \emph{open} sets $O_j$ yields
\[
 C^*(A)\leq C\Bigl(\bigcup_jO_j\Bigr)
       =\lim_j C(O_j)
       \leq\epsilon+(1-\eta)^{-n}\lim_j C^*(A_j).
\]
Let $\epsilon$ and then $\eta$ decrease to zero. Monotonicity
supplies the reverse inequality.

This is the ordinary-open argument of the classical capacitability
proof, \cite[\S8, Theorem 8.2]{BT82}, with its analytic hypotheses
proved separately here. An infinite intersection of open covers is
not declared open, and no infimum/supremum interchange is made.
\end{proof}

\begin{corollary}[Local Choquet capacitability and Borel outer regularity]
\label{cor:r2-local-choquet}
Fix an ordinary core $W\Subset V$. The set function
\[
 A\longmapsto C^*(A),\qquad A\subset W,
\]
is a finite monotone countably subadditive Choquet capacity: it is
continuous from below on arbitrary increasing sets and from above
on decreasing compact sets. Every Suslin subset $A\subset W$
(analytic in the descriptive-set-theoretic sense) is capacitable:
\[
 C^*(A)=\sup_{\substack{K\subset A\\K\text{ compact}}}C(K).
\]
In particular every Borel $A\subset W$ satisfies
\[
 C(A)=C^*(A).
\]
\end{corollary}

\begin{proof}
Monotonicity and countable subadditivity are already proved for
outer capacity. Compact continuity from above is
Lemma~\ref{lem:r2-outer-facts}; increasing continuity is the preceding
theorem. Finiteness follows by covering $\overline W$ with one
relatively compact ordinary open and applying CLN to its closure.
The core $W$ is a locally compact second-countable metrizable space
in the fixed finite $p$-adic algebraic setting. The Choquet
capacitability theorem, in the form recalled in
\cite[\S8, following Theorem 8.2]{BT82}, therefore gives
\[
 C^*(A)=\sup_{K\subset A\text{ compact}}C^*(K)
\]
for such Suslin $A$. This terminology concerns analytic sets of the
Polish space $W$, not analytic subspaces over the base field $K$.
For compact $K$, the mass formula of the
preceding theorem and the compact part of \textup{(E)} give
$C^*(K)=C(K)$.

For Borel $A$, raw inner regularity is elementary:
\[
 C(A)=\sup_{v}\sup_{K\subset A\text{ compact}}
       \int_K(Dv)^n
     =\sup_{K\subset A\text{ compact}}C(K).
\]
Only two suprema are interchanged in this equality. Combining it
with capacitability gives $C(A)=C^*(A)$.
\end{proof}

\paragraph{Scope and dependency order.}
The capacity-zero witness uses only the open extremal identity,
ambient no-jump, the proved increasing calculus, and pointwise
ambient domination. The arbitrary-set mass formula and Choquet
continuity do not use that witness theorem: their only null-set
input is \textup{(N)}. Thus the argument does not prove negligible
sets polar by first assuming that conclusion. The capacitability
statement is on each fixed relatively compact core of this explicitly
qualified domain. It does not assert a global Josefson theorem on an
arbitrary analytic space, or discard the negative-exhaustion-domain
hypothesis.

\paragraph{A distinction in abstract capacity theory.}
For a Borel capacity $c$, its Borel-envelope extension
$\inf\{c(B): A\subset B,\ B\text{ Borel}\}$ is different from
the ordinary-open extension used here. The familiar increasing-set
argument using $\bigcap_{k\geq j}G_k$ yields Borel sets, not open
ones. Raw Borel inner regularity and raw Borel continuity from below
therefore do not, by that argument, supply the missing continuity
axiom for the ordinary-open extension. Theorem
\ref{thm:r2-outer-capacity-increasing-sets} supplies the required
axiom by actual ordinary sublevel sets of extremal functions.

% END INCLUDED SOURCE: base_r2/sections/polar_converse_choquet


% BEGIN INCLUDED SOURCE: base_r2/sections/round2_status
\section{Second-round theorem and status increment}
\label{sec:r2-status}

This section updates, rather than replaces, the source-by-source comparison
in Section~\ref{sec:comparison}. The first-round statements remain frozen
with their original hypotheses. An earlier \texttt{OPEN} entry is closed
only by the particular second-round theorem cited below. No claim of
literature priority is encoded by \texttt{NEW\_PROVED}: that label records
a proof supplied in this project, relative to its accepted inputs.

\subsection{Two independent status axes and the fixed category}

The labels \texttt{EXTERNAL\_KNOWN}, \texttt{INTERNAL\_INPUT},
\texttt{DIRECT\_COROLLARY}, and \texttt{NEW\_PROVED} record provenance.
The labels \texttt{PROVED}, \texttt{CONDITIONAL}, \texttt{OPEN}, and
\texttt{FALSE} assess a precisely qualified proposed analogy. A theorem
proved under an explicit additional geometric hypothesis is
\texttt{PROVED} with that hypothesis; the corresponding general-domain
analogy is \texttt{CONDITIONAL} or \texttt{OPEN}. A counterexample to an
unrestricted assertion does not change the status of its qualified
replacement.

All local assertions retain the finite extension $K/\mathbf Q_p$, its
prescribed absolute value, and the allowed good, Hausdorff, boundaryless,
quasi-smooth pure-dimensional analytic spaces. Unless a narrower theorem
states otherwise, $S$ is an arbitrary \emph{fixed} symmetric current of
bidegree $(p,p)$, positive in the original CLD cone and separately
$d'$- and $d''$-closed. Products and current equalities are tested on
\emph{all original compactly supported smooth complementary CLD forms}.
Neither a change of positivity cone nor an enlargement of the test sheaf
is implicit. Put $q=n-p$; the relative top-capacity statements below use
$q\geq1$.

R and B remain \texttt{INTERNAL\_INPUT}. In particular, the bounded
classical class and its $[1]$-background calculus are not second-round
discoveries. The completed geometric adelic nefness theorem A2 also
remains an \texttt{INTERNAL\_INPUT} under its own number-field,
integrability, and boundary hypotheses. The original classical, CLD,
GR, and GK theorems recorded in Section~\ref{sec:comparison} remain
\texttt{EXTERNAL\_KNOWN} in their stated categories.

\subsection{Increment for the eleven Bedford--Taylor directions}

\begin{longtable}{@{}>{\raggedright\arraybackslash}p{.045\textwidth}>{\raggedright\arraybackslash}p{.19\textwidth}>{\raggedright\arraybackslash}p{.40\textwidth}>{\raggedright\arraybackslash}p{.28\textwidth}@{}}
\hline
 & Target & Qualified second-round increment & Remaining qualification or obstruction\\
\hline
\endfirsthead
\hline
 & Target & Qualified second-round increment & Remaining qualification or obstruction\\
\hline
\endhead
1 & Locally bounded psh functions &
\texttt{PROVED / DIRECT\_COROLLARY}: upper regularizations of locally
uniformly upper-bounded classical families are classical, by
Theorem~\ref{thm:r2-local-upper-envelope}; the projective bounded usc
obstacle envelope is Proposition~\ref{prop:r2-usc-obstacle}.
\texttt{PROVED / DIRECT\_COROLLARY}: original trace-current order determines
the actual bounded classical representative,
Proposition~\ref{prop:r2-current-faithfulness}. &
The new capacitary class is relative to $S$ and retains Borel
representatives. For $S=[1]$, every such class now has a unique
actual classical representative modulo local outer-capacity zero,
Corollary~\ref{cor:r2-ambient-classical-representative}
(\texttt{DIRECT\_COROLLARY}). Actual Borel values on the exceptional
set are not identified. No corresponding assertion for arbitrary
$S$ is made.\\[5pt]

2 & Positive closed currents and mixed products &
The arbitrary-fixed-$S$ bounded classical calculus is already
\texttt{PROVED} in R1. Theorem~\ref{thm:r2-completed-calculus} extends
it to $\mathcal A_{S,b}$, with well-defined Borel recursion,
quasi-everywhere independence, symmetry, and restriction/gluing:
\texttt{PROVED / NEW\_PROVED}. &
This does not define the wedge of two arbitrary positive closed
currents. Nor does it settle dependence only on the currents $D[u_i]$:
the general discontinuous potential-quotient assertion remains
\texttt{OPEN}.\\[5pt]

3 & Chern--Levine--Nirenberg estimates &
\texttt{PROVED / NEW\_PROVED}: the weighted $L^1$ estimate and uniform
truncation estimate of Theorem~\ref{thm:r2-weighted-l1}, under the
specified collar-trace integrability condition. The deep-sublevel
outer-capacity estimate and integrable-witness consequences are
\texttt{DIRECT\_COROLLARY},
Corollaries~\ref{cor:r2-normalized-witness} and
\ref{cor:r2-witness-outer-capacity}. &
The integrability hypothesis concerns an explicit trace of the fixed
$S$; it is not automatic for arbitrary $S$. Uniform exponential tails
and unbounded-weight stability from a uniform first moment alone are
\texttt{FALSE}, Propositions~\ref{prop:r2-no-exponential-tails} and
\ref{prop:r2-unbounded-weight-failure}.\\[5pt]

4 & Symmetry and weighted monotone convergence &
R1 already proves symmetry and weighted joint decreasing continuity
for actual bounded classical limits and fixed $S$.
\texttt{PROVED / NEW\_PROVED}: bounded increasing sequences converge
in outer relative capacity to their \emph{actual} Borel supremum,
Theorem~\ref{thm:r2-increasing-cauchy} and
Corollary~\ref{cor:r2-local-increasing}; weighted increasing continuity
in the completed class is Corollary~\ref{cor:r2-increasing-weighted}.
The ambient no-jump theorem
\ref{thm:r2-increasing-no-jump} identifies this supremum with its usc
regularization locally quasi-everywhere:
\texttt{PROVED / NEW\_PROVED}.
Theorem~\ref{thm:r2-weak-closure-multipliers} proves the qualified
quasicontinuous multiplier passage through weak closure. &
Ambient no-jump and its common-collar energy equality are closed
questions in this record. Replacing the actual supremum by its
ambient usc regularization remains \texttt{FALSE} without
qualification for arbitrary $S$. Unrestricted varying-$S$ weighted
stability remains \texttt{FALSE}.\\[5pt]

5 & Locality and maximum principles &
R1 Borel maximum locality and ordinary restriction/locality remain
\texttt{PROVED}. Theorem~\ref{thm:r2-polar-restriction} proves that
$\mathbf1_P S$ and $\mathbf1_{P^c}S$ are separately closed for a
classical pole set $P$, and that BT products commute with this
restriction: \texttt{PROVED / NEW\_PROVED}.
Theorem~\ref{thm:r2-maximal-ma-zero} proves local classical maximality
equivalent to $(Du)^n=0$ for the ambient background. &
Strict inequality sets of discontinuous potentials are Borel and
locally quasi-open, not automatically ordinary open. The polar
restriction theorem starts with a current defined on the entire
domain; it is not a general extension theorem for a current defined
only off the pole set.\\[5pt]

6 & Bedford--Taylor capacities &
\texttt{PROVED / DIRECT\_COROLLARY}: outer regularity on quasi-open
sets, Proposition~\ref{prop:r2-quasiopen-outer}, and the elementary
outer-capacity properties in Lemma~\ref{lem:r2-outer-facts}.
\texttt{PROVED / NEW\_PROVED}: two-sided local comparison with one
fixed relatively compact affine capacity domain,
Theorem~\ref{thm:r2-domain-capacity-comparison}. For the ambient
capacity on an exhaustion domain, compact raw outer regularity and
compact continuity are Theorem~\ref{thm:r2-compact-capacity-continuity};
outer-capacity continuity for arbitrary increasing subsets of a
fixed relatively compact core is
Theorem~\ref{thm:r2-outer-capacity-increasing-sets}.
Choquet capacitability and Borel outer regularity on that core are
Corollary~\ref{cor:r2-local-choquet}. &
Raw compact outer regularity and the corresponding Choquet theory
for \emph{arbitrary relative backgrounds $S$} remain \texttt{OPEN}.
The ambient result retains the exhaustion-domain and fixed-core
scope. General-domain nondegeneracy remains \texttt{FALSE}; the R1
proper-space example still applies.\\[5pt]

7 & Quasicontinuity and capacity convergence &
R1 bounded quasicontinuity is retained. Theorems
\ref{thm:r2-increasing-cauchy} and \ref{thm:r2-completed-calculus}
give the increasing and capacitary-completion extensions.
Corollary~\ref{cor:r2-extended-psh-qc} gives local finite
quasicontinuity and lower-truncation outer-capacity convergence for
every nontrivial extended classical psh function:
\texttt{PROVED / DIRECT\_COROLLARY} of the new polar-trace and
weighted estimates. &
Capacity precompactness of the full normalized continuous psh family,
and weighted BT stability under bounded pointwise convergence alone,
are \texttt{FALSE},
Proposition~\ref{prop:r2-capacity-noncompactness}. This counterexample
does not refute compact-set domination by the weak closure of the
admissible measure family.\\[5pt]

8 & Pluripolar and negligible sets &
\texttt{PROVED / NEW\_PROVED}: every original smooth ambient trace
vanishes on a nontrivial classical pole set,
Theorem~\ref{thm:r2-ambient-polar-traces}.
\texttt{DIRECT\_COROLLARY}: all ambient bounded BT coefficients avoid
such sets, and every locally pluripolar subset has local outer
capacity zero, Corollaries~\ref{cor:r2-ambient-bt-poles} and
\ref{cor:r2-ambient-polar-outer-zero}.
Theorem~\ref{thm:r2-capacity-zero-polar-witness} constructs a
nontrivial witness for an outer-zero set. The resulting local
equivalence of pluripolar, outer-zero, and negligible sets is
Corollary~\ref{cor:r2-polar-negligible-equivalence}.
Pole patching on one exhaustion domain is
Corollary~\ref{cor:r2-exhaustion-domain-pole-patching}. &
These local ambient equivalences are \texttt{PROVED}.
Unqualified untwisted pole globalization on every analytic space is
\texttt{FALSE}, already on $\mathbf P^{1,\an}_K$.
Arbitrary relative backgrounds may charge ambient pluripolar sets;
their mass is governed by
Theorem~\ref{thm:r2-polar-restriction}, not forced to vanish.\\[5pt]

9 & Comparison and domination &
The R1 compact-gap comparison theorem for arbitrary fixed $S$
remains \texttt{PROVED}.
\texttt{PROVED / DIRECT\_COROLLARY}: ambient strict-gap detection and
pointwise domination/uniqueness on the stated relatively compact
framed affine domains, Theorem~\ref{thm:r2-ambient-gap-detection}
and Corollary~\ref{cor:r2-pointwise-ambient}.
Relative atomic and smooth-factor detection criteria and explicit
faithful backgrounds are proved in
Propositions~\ref{prop:r2-relative-gap-detection} and
\ref{prop:r2-smooth-factor-condition}. &
Pointwise domination for arbitrary $S$ remains
\texttt{CONDITIONAL} on detection; the unrestricted version is
\texttt{FALSE} even for nonzero $S$ on a connected domain,
Example~\ref{ex:r2-nonzero-domination-failure}.
Capacity-modulo uniqueness and actual pointwise uniqueness are
distinct assertions.\\[5pt]

10 & Relative extremal functions &
\texttt{PROVED / NEW\_PROVED}: least upper bounds in the local
$S$-quasi-everywhere order, represented in $\mathcal A_{S,b}$,
Theorem~\ref{thm:r2-local-capacitary-envelope}.
Ambient bounded usc obstacle envelopes follow directly from B;
orthogonality for bounded quasicontinuous usc projective obstacles
is Theorem~\ref{thm:r2-qc-orthogonality},
\texttt{PROVED / NEW\_PROVED}.
For the ambient background on an exhaustion domain, compact/open
extremal equilibrium formulas are Theorems
\ref{thm:r2-compact-extremal} and \ref{thm:r2-open-extremal};
arbitrary relatively compact sets satisfy the outer-capacity
total-mass formula \ref{thm:r2-arbitrary-extremal-mass}. &
The ambient and relative categories remain distinct. General
\emph{relative} equilibrium-measure/contact-set formulas and their
capacity identities remain \texttt{OPEN}; unrestricted replacement
of an $S$-capacitary envelope by the ambient envelope is \texttt{FALSE},
Proposition~\ref{prop:r2-negligible-positive}.\\[5pt]

11 & Monge--Amp\`ere and Dirichlet theory &
\texttt{PROVED / NEW\_PROVED}: the ambient local equivalence between
maximality and the homogeneous equation,
Theorem~\ref{thm:r2-maximal-ma-zero}.
The explicit peak-barrier theorem
\ref{thm:r2-peak-dirichlet} gives a bounded classical solution of
$(DH_g)^n=0$, ordinary boundary values, uniqueness, and the boundary
sup-norm stability estimate under
Definition~\ref{def:r2-peak-barriers}.
Corollary~\ref{cor:r2-dirac-dirichlet} gives an explicit continuous
solution and uniqueness for a single positive divisorial atom and
constant boundary data on the domains constructed with its barrier:
\texttt{PROVED / DIRECT\_COROLLARY}. &
As a general-domain analogy the Dirichlet conclusion is
\texttt{CONDITIONAL} on those explicit barriers. Their general
geometric construction, interior continuity for the Perron solution,
and existence for a general prescribed nonzero measure remain
\texttt{OPEN}. Bounded negative exhaustion domains are explicitly
defined and locally supplied; they are not silently assumed to have
peak barriers at every boundary point.\\[3pt]
\hline
\end{longtable}

\subsection{Scope distinctions that survive the second round}

\paragraph{Fixed backgrounds and controlled varying backgrounds.}
The primary calculus keeps $S$ fixed. The order-dominated replacement
in Proposition~\ref{prop:dominated-background-stability} allows
$0\leq S_j\leq C R_0$, $S_j\to S$, with one fixed positive closed
$R_0$ and a fixed differentiated tuple; its multiplier hypotheses
must also be retained. The second-round weak-closure multiplier
theorem instead requires uniform coefficient bounds by a fixed local
capacity and quasicontinuous weights. Neither statement removes the
counterexample in Corollary~\ref{cor:varying-background-false}.

\paragraph{Actual representatives and ambient upper regularization.}
An actual increasing supremum is now an element of the bounded
capacitary closure, and its BT products are established there. For the
ambient background, Theorem~\ref{thm:r2-increasing-no-jump} identifies
it with its ordinary usc regularization locally outside outer-zero
sets and proves the common-collar energy no-jump equality. More
generally, Corollary~\ref{cor:r2-ambient-classical-representative}
assigns each $a\in\mathcal A_{[1],b}$ a unique classical representative
$r_a$ modulo such sets. For an arbitrary Borel representative $a$,
$r_a$ need not equal $a^*$; arbitrary modifications on polar sets are
allowed in that capacitary class. This distinction does not affect
the actual increasing-supremum theorem.
Theorem~\ref{thm:r2-local-capacitary-envelope} obtains a least upper
bound without taking an uncountable union of exceptional sets: each
competitor has its own quasi-everywhere comparison, and a countable
energy construction determines the resulting class. For an arbitrary
relative background $S$, the capacitary representative cannot in
general be replaced by the ambient usc envelope.

\paragraph{Raw capacity and outer capacity.}
For arbitrary $S$, the identity $C^*(E)=C(E)$ is proved for quasi-open
$E$; full raw compact outer regularity remains open. For $S=[1]$ on
an exhaustion domain, it is now proved for compact sets and for
every Borel subset of each fixed relatively compact core. The
ordinary-open proof of Theorem~\ref{thm:r2-outer-capacity-increasing-sets}
supplies the additional continuity axiom needed for the local
Choquet theorem; compact continuity of $C^*$ alone would not supply
that proof. The raw-capacity closure criterion in
Proposition~\ref{prop:capacity-weak-closure} is therefore resolved
for the stated ambient compact stages but remains an exact target
for arbitrary relative $S$. The explicit potential-noncompactness
example has an admissible unweighted weak limit and does not refute
compact-set domination by the admissible measure family's weak closure.

\paragraph{Polar sets and arbitrary support.}
The general ambient pole-trace problem is now \texttt{PROVED}; it is
no longer a missing hypothesis in the ambient nonconcentration result.
The proof uses the absence of divisorial poles, whole-box model
approximation, and the capacity-controlled multiplier passage, as
recorded in Theorem~\ref{thm:r2-ambient-polar-traces}. For arbitrary
$S$, the exact statement is
$\mathbf1_P T_{S,k}=T_{\mathbf1_PS,k}$ from
Theorem~\ref{thm:r2-polar-restriction}. In particular, a background
already carried by $P$ is not made nonpolar by applying bounded BT
operators. Ambient outer-zero sets now admit local nontrivial pole
witnesses, and the local pluripolar/negligible/outer-zero classes agree
by Corollary~\ref{cor:r2-polar-negligible-equivalence}. There is also
pole patching on a single connected exhaustion domain. A global
untwisted version for all analytic spaces is false on proper connected
spaces; these local and qualified global statements are not conflated.

\paragraph{Potential quotients.}
Proposition~\ref{prop:r2-current-faithfulness} concerns $[u]$ and
$[v]$, not merely $D[u]$ and $D[v]$. General independence under
replacing a differentiated potential by another actual bounded
classical potential with the same $D$-current remains \texttt{OPEN}.
The continuous harmonic-difference result
\ref{lem:continuous-harmonic-multipliers}, the ambient-mixed and smooth
background cases \ref{lem:potential-independence-ambient-mixed} and
\ref{lem:potential-independence-smooth}, and the dominated-closure
case \ref{prop:potential-independence-dominated-closure} remain proved
partial answers. None is silently promoted to the general assertion.

\paragraph{Comparison with other differential calculi.}
Agreement on ordinary CLD tests with compatible model and
Chambert--Loir measures is retained where the accepted R/B comparison
applies. Equality on every $\delta$-test form is a different statement.
The report does not extend arbitrary original currents to an enlarged
CLD25 test sheaf, or claim the nonpositive, nonclosed, or nondiagonal
parts of the GK calculus. The completed A2 adelic theorem is an input
to subsequent geometric applications, not an unresolved item in this
local analytic program.

\subsection{Three minimal targets for the next round}

The former ambient no-jump and local polar-witness targets are now
proved. The following remaining targets are not assumed anywhere in
the second-round arguments.

\begin{enumerate}
\item \textbf{Compact domination for the raw capacity's measure
closure for arbitrary relative backgrounds.} For an arbitrary allowed
fixed $S$, on one fixed permitted capacity domain and compact test
stage, prove or refute
\[
 \nu(K)\leq\operatorname{Cap}_S(K;\Omega)
 \quad\text{for every compact }K
\]
for every weak limit $\nu$ of normalized top BT measures (with the
localization conventions of
Proposition~\ref{prop:capacity-weak-closure}). It is enough to resolve
this precise closure statement; a capacity-convergent subsequence of
all normalized potentials is not available. A counterexample must
violate the displayed inequality, rather than only potential
precompactness or weighted pointwise continuity. The ambient result
on exhaustion domains is already closed and is not assigned again
as new work.

\item \textbf{The general potential-quotient lemma.} Given actual
$u,v\in\Cb(U)$ with $D[u]=D[v]$ as original currents, establish
or refute
\[
 D\bigl((u-v)S\bigr)=0
\]
for every allowed fixed positive closed background $S$.
The continuous-difference and specified-background cases are already
proved. Equality of the original currents $[u]$ and $[v]$, and the
unique ambient capacitary representative, do not by themselves
settle this statement about equality of their $D$-currents.

\item \textbf{Nonzero data on one prescribed domain.} Starting with
an explicitly qualified domain and boundary barriers, construct and
control solutions for a sum of at least two prescribed positive
divisorial atoms on that \emph{same} domain. The single-atom result
\ref{cor:r2-dirac-dirichlet} constructs its domain together with the
atom and barrier. It does not provide this next fixed-domain step;
adding its potentials produces mixed terms. A passage from finite
atomic data to more general measures would then require uniform
potential bounds, boundary control, and a proved capacity stability
estimate for that family of solutions.
\end{enumerate}

The homogeneous peak-barrier theorem yields bounded classical
solutions with ordinary boundary limits; interior continuity has not
been proved there. General prescribed nonzero measures and general
ordinary domains still require additional geometric and analytic
hypotheses. The relative equilibrium/contact problem and comparison
on enlarged $\delta$-test categories remain separate subsequent tasks.

\subsection{Excluded routes and the conclusions they actually exclude}

\begin{enumerate}
\item \texttt{FALSE}: arbitrary varying-background weighted stability,
even at order zero; see
Corollary~\ref{cor:varying-background-false}. Fixed-background
monotonicity and the qualified order-dominated replacement survive.

\item \texttt{FALSE}: relative nonconcentration on every ambient
pluripolar set, or universal relative negligibility of the discrepancy
between an actual supremum and its ambient usc envelope; see
Proposition~\ref{prop:polar-relative-mass} and
Proposition~\ref{prop:r2-negligible-positive}. The valid replacement
is support restriction and the $S$-capacitary envelope.

\item \texttt{FALSE}: pointwise domination solely from relative
measure inequalities for an arbitrary nonzero $S$ on a connected
domain; see Example~\ref{ex:r2-nonzero-domination-failure}.
Relative gap detection is a real additional condition.

\item \texttt{FALSE}: nondegeneracy of the untwisted capacity on every
analytic domain, including proper ones; see
Proposition~\ref{prop:compact-capacity-zero}. Local framed domains
and projective metric capacities must be stated separately.

\item \texttt{FALSE}: compactness in capacity of the whole normalized
continuous psh family, or weighted BT continuity under bounded
pointwise convergence; see
Proposition~\ref{prop:r2-capacity-noncompactness}.
\texttt{OPEN}: the compact raw-capacity weak-closure inequality for
arbitrary relative $S$. The latter is not contradicted by this example;
the qualified ambient compact inequality is now proved.

\item \texttt{FALSE}: a uniform first-moment estimate implies uniform
integrability of a fixed unbounded multiplier, or universal exponential
tail control for arbitrary $S$; see
Propositions~\ref{prop:r2-unbounded-weight-failure} and
\ref{prop:r2-no-exponential-tails}. The proved weighted estimate is
used with its actual trace-integrability and truncation conclusions.

\item \texttt{FAILED\_ATTEMPT}: deriving raw outer regularity by
interchanging an infimum over ordinary neighborhoods with a supremum
over a measure family. Individual Radon outer regularity does not
justify that interchange. For arbitrary $S$, the quasi-open argument
and weak-closure characterization remain the available replacements.
For the ambient exhaustion-domain capacity, the extremal mass formula
and increasing ordinary sublevels supply the additional proof.

\item \texttt{FALSE}: untwisted pole globalization on every analytic
space. The rational point of $\mathbf P^{1,\an}_K$ in
Section~\ref{sec:r2-polar-converse-choquet} is locally pluripolar,
while nontrivial global classical psh functions there are constant.
Corollary~\ref{cor:r2-exhaustion-domain-pole-patching} proves the
qualified replacement on one connected exhaustion domain.

\item \texttt{FAILED\_ATTEMPT}: inferring smooth-trace nullity merely
from density of divisorial points, or passing an indicator through
weak convergence without coefficient control. The successful polar
proof additionally uses quasicontinuity of the actual indicator and
uniform same-domain capacity estimates.

\item \texttt{FAILED\_ATTEMPT}: swapping two discontinuous moving
limits by formal integration by parts, using the desired joint
continuity to prove its own energy estimate, or treating an
uncountable union of quasi-everywhere exceptional sets as negligible.
The established R1 trace argument and R2 countable energy-envelope
construction avoid these inferences.
\end{enumerate}

The second-round record therefore closes several analytic gaps while
retaining explicit frontier statements. It does not label the full
non-Archimedean Bedford--Taylor program complete.

% END INCLUDED SOURCE: base_r2/sections/round2_status

\clearpage
\part*{Part II. Third-round proofs and counterexamples}
\addcontentsline{toc}{part}{Part II. Third-round proofs and counterexamples}

% BEGIN INCLUDED SOURCE: sections/energy_convex_blocks
\section{Mixed energy controls capacity in every remaining degree}
\label{sec:r3-all-degree-energy}

The fixed background in this section is arbitrary: $S$ is symmetric,
original-positive, of bidegree $(p,p)$, and separately closed, and
$q=n-p\geq1$. The original field, ordinary topology, and all original
CLD tests are retained. The estimates resolve the varying-tester
energy obstruction mentioned in
Section~\ref{sec:r3-relative-compact-capacity}. They do not identify
nonlinear weak limits of the original unaveraged measures.

Expressions such as $(Da)^i(Db)^jS$ abbreviate the already constructed
mixed products over the fixed $S$. An energy expression always means
$\Gamma_Q(f,g)$ for the specified mixed background $Q$ of bidegree
$(n-1,n-1)$. It never means a wedge of an arbitrary energy current
with another current.

\subsection{A common-collar energy distance}

Let $\Omega$ be a sufficiently small ordinary norm-box domain and
$L\Subset\Omega$ a compact set in a strictly smaller norm box. Fix
a smooth classical psh function $h$ with semipositive original
curvature form. Consider $u,v\in\Cb(\Omega)$ with
\[
 u=v=h\quad\hbox{on }\Omega\setminus L.
\]
There is one fixed bound for $u,v,h$ near $L$. Put $f=u-v$ and
\begin{equation}
 \label{eq:r3-I-energy}
 I_S(u,v)=\int f\bigl((Dv)^qS-(Du)^qS\bigr).
\end{equation}
This integral is finite and has compact support.

\begin{lemma}[The positive energy identity]
\label{lem:r3-I-positive}
For the actual point values of the potentials one has
\[
 I_S(u,v)=\sum_{a=0}^{q-1}
    \int\Gamma_{(Du)^a(Dv)^{q-1-a}S}(f)\geq0.
\]
\end{lemma}

\begin{proof}
Symmetry and multilinearity of the proved mixed calculus give
\[
 (Du)^qS-(Dv)^qS
   =D\left(f\sum_{a=0}^{q-1}(Du)^a(Dv)^{q-1-a}S\right).
\]
Apply $-\int fD(fQ)=\int\Gamma_Q(f)$ to each positive separately
closed summand. No weakly varying arbitrary background is used.
\end{proof}

\begin{theorem}[An all-degree mixed-energy capacity estimate]
\label{thm:r3-I-capacity}
Under these common-collar assumptions, for every $\delta>0$,
\begin{equation}
 \label{eq:r3-all-degree-I-capacity}
 \operatorname{Cap}_S^*(\{|u-v|>\delta\};\Omega)
       \leq C\delta^{-2} I_S(u,v)^{\,2^{-q}}.
\end{equation}
The finite constant depends on the fixed compact collar, the common
bound for the potentials, $q$, and finitely many positive coefficient
evaluations of $S$. It is independent of the normalized capacity
tester and all regularization indices. Restriction to a smaller
compact core gives the corresponding cutoff version.
\end{theorem}

\begin{proof}
Fix $w\in\Cb(\Omega)$ with $0\leq w\leq1$. For $0\leq r\leq q-1$ put
\[
 E_r(w)=\int\Gamma_{(Dw)^r(D(u+v))^{q-1-r}S}(f).
\]
These are nonnegative integrals on the same compact support, since
$f=0$ off $L$. Expansion and the positive energy identity give
\begin{equation}
 \label{eq:r3-E-zero}
 E_0(w)\leq2^{q-1}I_S(u,v).
\end{equation}
For $r\geq1$ write
$Q=(Dw)^{r-1}(D(u+v))^{q-1-r}S$, of bidegree $(n-2,n-2)$.
First suppose that the potentials are smooth. Stokes for the
compactly supported first-order expression gives
\begin{equation}
 \label{eq:r3-E-transfer}
 E_r(w)=\int d'f\wedge d''w\wedge(Du-Dv)\wedge Q.
\end{equation}
For the sign, apply $d'$ to
$d'f\wedge d''f\wedge d''w\wedge Q$; the result is
$-d'f\wedge Df\wedge d''w\wedge Q+
d'f\wedge d''f\wedge Dw\wedge Q$.

Choose a fixed nonnegative compact smooth cutoff $\kappa$ equal
to one near $L$. Symmetry and Cauchy--Schwarz, separately for
the positive backgrounds $Du\,Q$ and $Dv\,Q$, imply
\begin{align*}
 E_r(w)
 &\leq\left(\int\Gamma_{Du\,Q}(f)\right)^{1/2}
       \left(\int\kappa\,\Gamma_{Du\,Q}(w)\right)^{1/2}\\
 &\quad+\left(\int\Gamma_{Dv\,Q}(f)\right)^{1/2}
       \left(\int\kappa\,\Gamma_{Dv\,Q}(w)\right)^{1/2}.
\end{align*}
The first factors squared sum to $E_{r-1}(w)$. The second factors
squared have a uniform finite sum: for either $R=Du\,Q$ or
$R=Dv\,Q$ the cutoff identity is
\[
 \int\kappa\,\Gamma_R(w)
  =\tfrac12\langle R,w^2D\kappa\rangle-\int\kappa wD(wR).
\]
The subtracted term is nonnegative. The relative finite-test CLN
estimate bounds the first term for the fixed form $D\kappa$.
All slots of $R$ are bounded copies of $u,v,w$ or their sum.
Thus, for a fixed $A_r$ independent of $w$,
\begin{equation}
 \label{eq:r3-E-recurrence}
 E_r(w)\leq A_r E_{r-1}(w)^{1/2}.
\end{equation}
For $q=1$ this recurrence is empty.

The compact Green identity for $f^2$ now gives
\[
 \int f^2(Dw)^qS
       =-2\int f\,\Gamma_{(Dw)^{q-1}S}(f,w).
\]
Cauchy--Schwarz with the same compact cutoff and another uniform
CLN bound yield
\begin{equation}
 \label{eq:r3-E-square-final}
 \int f^2(Dw)^qS\leq A E_{q-1}(w)^{1/2}.
\end{equation}
The common bound for $f$ is absorbed in $A$. Iteration proves
\begin{equation}
 \label{eq:r3-all-degree-square-control}
 \int f^2(Dw)^qS\leq C I_S(u,v)^{\,2^{-q}}.
\end{equation}

For bounded potentials, use the protected prescribed-box
construction for simultaneous smooth decreasing rows. If necessary,
first modify $u,v$ outside a neighborhood of $L$ to a common
auxiliary psh collar, preserving them up to the same constant
near $L$ and preserving their difference $f$ everywhere. Their
rows then have an exact common smooth collar outside a fixed
larger compact set $L'\Subset\Omega$. This does not assert that
arbitrary smoothing preserves the originally prescribed $h$.
All limiting weighted and energy expressions in the argument
are supported in $L$, so the auxiliary modification leaves
them unchanged. Choose all cutoffs for the larger fixed stage.
Take a simultaneous smooth decreasing row for $w$ there.
After discarding a prefix the latter lies in $[0,2]$,
so all constants remain uniform. Every background in
the preceding formulas is an explicit mixed product over the
same $S$. Expand each energy as
\[
 \Gamma_R(a,b)=\tfrac12\{D(abR)-aD(bR)-bD(aR)\}.
\]
Its scalar factors are polynomials in finitely many decreasing
psh rows. The functional-weighted joint monotone theorem passes
every term to the asserted bounded-potential expression on a
fixed cutoff. The differences retain fixed compact support.
This proves the identities and inequalities without using
capacity continuity to justify the energy calculation.

Chebyshev and the supremum over $w$ give the asserted raw
capacity bound. The strict gap set is relatively quasi-open
by quasicontinuity of $u,v$. Proposition~\ref{prop:r2-quasiopen-outer}
therefore identifies its raw and outer capacities. No such
set was assumed ordinarily open.
\end{proof}

\subsection{Concavity selects simultaneous convex blocks}

For the fixed collar define
\[
 \mathcal E_{S,h}(u)
   =\frac1{q+1}\sum_{j=0}^q
       \int(u-h)(Du)^j(Dh)^{q-j}S.
\]
All nonzero weighted integrands lie on the fixed compact core.
No nondegeneracy assumption on $Dh$ is imposed.

\begin{lemma}[Quantitative midpoint concavity]
\label{lem:r3-energy-midpoint}
The energy is concave and
\begin{equation}
 \label{eq:r3-midpoint-I-bound}
 \mathcal E_{S,h}\!\left(\frac{u+v}{2}\right)
  -\frac{\mathcal E_{S,h}(u)+\mathcal E_{S,h}(v)}2
      \geq c_q I_S(u,v),\qquad c_q=\frac q{4^{q+1}}.
\end{equation}
For $u_t=(1-t)u+tv$,
\[
 \frac{d}{dt}\mathcal E_{S,h}(u_t)=\int(v-u)(Du_t)^qS,
 \qquad
 \frac{d^2}{dt^2}\mathcal E_{S,h}(u_t)
       =-q\int\Gamma_{(Du_t)^{q-1}S}(v-u).
\]
\end{lemma}

\begin{proof}
The energy is a polynomial in $t$. In its derivative, compact
Green exchange replaces
\[
 \int(u_t-h)D(v-u)(Du_t)^{j-1}(Dh)^{q-j}S
\]
by
$\int(v-u)(Du_t-Dh)(Du_t)^{j-1}(Dh)^{q-j}S$.
The finite sum telescopes to the first derivative. Differentiating
again and using compact Green gives the second derivative and
concavity. These are identities in the established bounded
mixed calculus; smooth common-collar rows give the same proof.

The one-variable midpoint Green formula expresses the gap as
\[
 \frac q2\int_0^1\min(t,1-t)
       \left(\int\Gamma_{(Du_t)^{q-1}S}(u-v)\right)dt.
\]
On $[1/4,3/4]$, expansion into positive mixed energy measures gives
\[
 \int\Gamma_{(Du_t)^{q-1}S}(u-v)
 \geq4^{-(q-1)}\int\Gamma_{(D(u+v))^{q-1}S}(u-v)
 \geq4^{-(q-1)}I_S(u,v).
\]
The weight is at least $1/4$ on this interval of length $1/2$.
This yields the displayed constant.
\end{proof}

\begin{corollary}[An energy criterion for capacity convergence]
\label{cor:r3-energy-capacity-topology}
For common-collar $u,v$ put
\[
 J_{S,h}(u,v)=\mathcal E_{S,h}(u)-\mathcal E_{S,h}(v)
                 +\int(v-u)(Du)^qS.
\]
Then
\[
 \frac{I_S(u,v)}{q+1}
       \leq J_{S,h}(u,v)
       \leq\frac{qI_S(u,v)}{q+1}.
\]
For a uniformly bounded common-collar sequence $u_j$ and $u$,
the following conditions are equivalent:
\begin{enumerate}
\item $I_S(u_j,u)\to0$;
\item $u_j\to u$ locally in outer relative capacity on all
      permitted compact/domain pairs;
\item $\mathcal E_{S,h}(u_j)\to\mathcal E_{S,h}(u)$ and
      $\int(u_j-u)(Du)^qS\to0$.
\end{enumerate}
In particular, for an actual pointwise convergent sequence with
these bounds and collar, convergence of its energies suffices
for capacity convergence.
\end{corollary}

\begin{proof}
Taylor's formula and the preceding second derivative give
\[
 J_{S,h}(u,v)=
 \sum_{a=0}^{q-1}\frac{q-a}{q+1}
       \int\Gamma_{(Dv)^a(Du)^{q-1-a}S}(v-u).
\]
Indeed the coefficient is
$q\binom{q-1}{a}\int_0^1t^a(1-t)^{q-a}dt=(q-a)/(q+1)$.
Each measure is positive, proving the two bounds.
Condition (3) therefore implies (1). The all-degree
estimate gives convergence on the fixed capacity domain
$\Omega$ from (1), and
Corollary~\ref{cor:r2-fixed-domain-local-convergence} gives
the all-domain local assertion (2).

Conversely, the completed weighted capacity-continuity
theorem applies to every finite term in the energy, with
a compact cutoff equal to one near $L$. It gives both
limits in (3), and also the defining limit for $I$.
The common support and local bounds justify all these
compact pairings. Under actual pointwise convergence,
the integral against the fixed measure $(Du)^qS$ tends
to zero by bounded convergence on $L$. This proves the
last assertion. This corollary is
\texttt{DIRECT\_COROLLARY} of the preceding estimates
and the established completed calculus.
\end{proof}

\begin{theorem}[Convex-block compactness in every positive remaining degree]
\label{thm:r3-all-degree-convex-blocks}
Let $\Omega$ be an ordinary open in one allowed smooth affine
algebraic neighborhood, and let $S$ be the fixed background above.
Every sequence $u_j\in\Cb(\Omega)$ uniformly bounded on each
compact neighborhood admits finite tail convex combinations
\[
 v_N\in\operatorname{conv}\{u_j:j\geq N\}
\]
and an actual Borel $a\in\mathcal A_{S,b}(\Omega)$ such that
$v_N\to a$ locally in outer relative capacity on every permitted
compact/domain pair.

For an ordinary $\Omega$ as in this theorem, use the same definition
of $\mathcal A_S^{[0,1]}(\Omega)$ as in
Definition~\ref{def:r3-normalized-completion}: its normalized
approximating functions must be defined on the entire $\Omega$.
If $0\leq u_j\leq1$ on that entire original $\Omega$, then
$0\leq v_N\leq1$ there and one may take
$a\in\mathcal A_S^{[0,1]}(\Omega)$. Its normalized approximants
are defined on that same domain. The completed weighted mixed
calculus applies to these convex blocks and their limit.
\end{theorem}

\begin{proof}
Choose a countable cofinal family of norm-box stages
$V_l\Subset\Omega_l\Subset\Omega$ as in the R2 completion theorem.
Every permitted compact/domain pair is covered by finitely many
$V_l$ whose $\Omega_l$ lie inside its prescribed capacity domain.

For each $l$, the common bounds on the original sequence and
the smooth collar construction give $c_l$, a smooth psh collar
$h_l$, and a fixed compact $L_l\Subset\Omega_l$, such that
\[
 u_{j,l}=\max(u_j+c_l,h_l)
\]
equals $u_j+c_l$ near $\overline V_l$ and equals $h_l$ off $L_l$.
These functions have a common bound near $L_l$. The construction
uses the uniform bound only; no monotonicity of $u_j$ is required.

Put $U_j=(u_{j,l})_{l\geq1}$. A finite convex vector
$W=\sum_jb_jU_j$ has components $W_l=\sum_jb_ju_{j,l}$ with
the same collar $h_l$. On $V_l$ it agrees with the global
convex combination $\sum_jb_ju_j+c_l$. These componentwise
convex combinations are retained; the nonlinear patching
operation is not applied a second time to the global sum.

Write $\mathcal E_l=\mathcal E_{S,h_l}$. Relative CLN gives
$B_l<\infty$ such that $|\mathcal E_l(W_l)|\leq B_l$ for
all these convex vectors. The functional
\[
 F(W)=\sum_{l\geq1}\frac{2^{-l}}{1+B_l}\mathcal E_l(W_l)
\]
is bounded and concave. The series converges uniformly and
absolutely on the convex hull. Let
\[
 m_N=\sup\{F(W):W\in\operatorname{conv}(U_j:j\geq N)\}.
\]
Then $m_N\downarrow m_\infty$ with finite limit. By the
definition of supremum, choose a finite vector $W_N$ in
the indicated tail with $F(W_N)\geq m_N-\epsilon_N$,
where $\epsilon_N=2^{-N}$. No compactness of the convex
hull or existence of a maximizing vector is assumed.

For $M\geq N$ the midpoint of $W_N,W_M$ belongs to the
$N$-tail. Its concavity gap is nonnegative and at most
\[
 \frac{m_N-m_M+\epsilon_N+\epsilon_M}{2}\longrightarrow0
       \quad(N,M\to\infty).
\]
Every summand of the gap is nonnegative. The midpoint
estimate consequently gives, for each fixed $l$,
\[
 I_S(W_{N,l},W_{M,l})\longrightarrow0.
\]
The all-degree estimate makes the component sequences
Cauchy in capacity on their fixed stages.

Use the same finite weights defining $W_N$ to put
$v_N=\sum_jb_{N,j}u_j$ on all of $\Omega$. It is in
the required tail convex hull and equals $W_{N,l}-c_l$
on $V_l$. Thus, for every $l$ and every $\delta>0$,
\[
 \operatorname{Cap}_S^*
   (V_l\cap\{|v_N-v_M|>\delta\};\Omega_l)\longrightarrow0.
\]
Finite subadditivity and domain monotonicity give this
Cauchy property on every prescribed compact/domain pair.
Run the summable-subsequence construction of the R2 completion
theorem on this one global countable cofinal list: at its
$m$th step make the first $m$ capacity errors summable.
The pointwise limsup of that single subsequence is one
actual Borel function on all of $\Omega$. It is the local
capacity limit on every stage, and preserves the common
local bounds and the interval $[0,1]$ when present.
The approximants remain the global $v_N$. This proves
all assertions. Weighted mixed convergence is then an
application of the already proved completed calculus.
\end{proof}

\begin{corollary}[Actual bounded pointwise limits have a completed calculus]
\label{cor:r3-pointwise-completed}
In the preceding setting, suppose $u_j\in\Cb(\Omega)$ are locally
uniformly bounded and converge at every point to $u$. Then the
actual Borel function $u$ belongs to $\mathcal A_{S,b}(\Omega)$.
It has local relative quasicontinuity and the completed weighted
BT calculus. There are finite tail convex combinations of the
given sequence converging locally in outer $S$-capacity to these
same point values modulo relative outer-null sets.

This does not assert weighted or higher-degree unweighted
BT convergence for the original unaveraged $u_j$.
\end{corollary}

\begin{proof}
Choose $v_N$ and its capacity limit $a$ from the theorem.
At every fixed point, every finite convex combination of
the $N$th tail tends to $u$ as $N\to\infty$, so $v_N\to u$
pointwise. The summable-subsequence construction for the
known capacity convergence gives a subsequence converging
pointwise to $a$ off locally relative-outer-null sets on
the same global cofinal list. Thus $u=a$ outside these
sets. Subadditivity shows that $v_N\to u$ in capacity
on each exact compact/domain pair. This is membership
in the completed class with the given Borel values.
The remaining assertions are the completed-calculus
and quasicontinuity theorems, and are
\texttt{DIRECT\_COROLLARY}.
\end{proof}

\paragraph{The nonlinear distinction.}
For $q=1$, linearity identifies a weak limit of $D(u_jS)$
with $D(aS)$, as in Theorem~\ref{thm:r3-linear-convex-blocks}.
For $q>1$ the measure of a convex combination contains mixed
terms. The new theorem does not identify a weak limit of
$(Du_j)^qS$ with $(Da)^qS$. It provides capacity compactness
after convex averaging, not capacity precompactness of the
original sequence; the latter remains false even for the
ambient background and continuous potentials.

\paragraph{Status and dependence.}
The all-degree energy-capacity estimate and the simultaneous
concave-tail selection argument are \texttt{NEW\_PROVED}.
The positive energy identity and the Aubin--Mabuchi variation
formulas are \texttt{DIRECT\_COROLLARY} of the bounded mixed
and common-collar Green calculus. No compact outer regularity,
strictly psh exhaustion, arbitrary-current wedge, or arbitrary
varying-background stability is used.


% END INCLUDED SOURCE: sections/energy_convex_blocks


% BEGIN INCLUDED SOURCE: sections/strong_relative_topology
% Reconstructed after the scratch reset from the complete proof record.
% Previous independently reviewed source SHA256:
% b5b8bb8d7156c7574dc8297ee05c5759aab617fdacda65f5b87d5df7db2e4dfe.
\section{Strong psh-dual topology for arbitrary fixed backgrounds}
\label{sec:r3-strong-relative-topology}

Keep the fixed finite extension $K/\mathbf Q_p$, its original
absolute value, and the allowed analytic spaces. Let $S$ be a
fixed symmetric original-positive, separately $d'$- and
$d''$-closed current of bidegree $(p,p)$, and put $q=n-p\geq1$.
All identities below use the full original compact smooth test
sheaf. Write $\mu_u=(Du)^qS$ for $u\in\mathcal C_b(U)$.
For a compactly supported finite signed measure $\sigma$ set
\[
 \|\sigma\|_{\mathrm{psh},U}
 =\sup_{\substack{a\in\mathcal C_b(U)\\0\leq a\leq1}}
          \left|\int a\,d\sigma\right|,
 \qquad
 d_{\mathrm{psh},U}(\nu,\nu')
       =\|\nu-\nu'\|_{\mathrm{psh},U}.
\]
This is a seminorm; no faithfulness assertion is implicit.

\subsection{The common-collar estimate}

Let $\Omega$ be a prescribed ordinary norm-box domain. Let
$L\Subset\Omega$ lie in a strictly smaller norm box, and let $h$
be smooth classical psh with semipositive original curvature form.
Suppose
\[
 u,v\in\mathcal C_b(\Omega),\qquad |u|,|v|\leq B,\qquad
 u=v=h\quad\hbox{on }\Omega\setminus L.
\]
Choose original smooth $0\leq\chi\leq1$ of compact support,
equal to one on a neighborhood of $L$. Locality shows that
$\mu_u-\mu_v$ is supported in $L$. Thus
\[
 \mathfrak d_S(u,v)
   :=\|\mu_u-\mu_v\|_{\mathrm{psh},\Omega}
    =d_{\mathrm{psh},\Omega}(\chi\mu_u,\chi\mu_v)
\]
does not depend on the cutoff. Put
$I_S(u,v)=\int(u-v)(\mu_v-\mu_u)$.

\begin{theorem}[Energy and strong psh-dual distance]
\label{thm:r3-relative-energy-dual}
Under these hypotheses there is a finite constant $C$, depending
only on the fixed compact stages, the common bound, $q$, and
finitely many positive coefficient evaluations of $S$, such that
\begin{equation}
 \mathfrak d_S(u,v)\leq C I_S(u,v)^{1/2},
 \qquad I_S(u,v)\leq4B\,\mathfrak d_S(u,v).
 \label{eq:r3-relative-energy-dual}
\end{equation}
The constant is independent of the normalized tester and of
regularization indices.
\end{theorem}

\noindent\textbf{Status: \texttt{NEW\_PROVED}} for the uniform
forward estimate; the reverse estimate is a
\texttt{DIRECT\_COROLLARY} of normalization and positivity of $I_S$.

\begin{proof}
Set $f=u-v$ and
\[
 Q=\sum_{b=0}^{q-1}(Du)^b(Dv)^{q-1-b}S .
\]
Each summand is an already defined positive separately closed
current of bidegree $(n-1,n-1)$. Symmetry, telescoping, and the
compact Green identity give
\[
 D(fQ)=\mu_u-\mu_v,\qquad
 \int\Gamma_Q(f)=I_S(u,v)\geq0 .
\]
For a normalized actual classical tester $a$, the protected
common-collar approximation and the bounded Green identity yield
\[
 \int a(\mu_u-\mu_v)=-\int\chi\,\Gamma_Q(a,f).
\]
There is no cutoff error: $\Gamma_Q(\chi,f)=0$, since $\chi$ is
constant near $L$ and $f=0$ off $L$. This identity is legitimate
for bounded, possibly discontinuous $a,u,v$ by the fixed-background
energy calculus; one first uses smooth rows on the prescribed box
and then its weighted convergence. No varying arbitrary background
is passed to a limit.

For $0\leq a\leq1$, the chain identity and positivity give
\[
 0\leq\int\chi\Gamma_Q(a)
 =\tfrac12\langle a^2Q,D\chi\rangle
                  -\int\chi a\,D(aQ)
 \leq\tfrac12|\mu_{Q,D\chi}|(\operatorname{supp}\chi).
\]
Relative CLN bounds the last expression uniformly in $u,v,a$.
Cauchy--Schwarz therefore bounds the preceding pairing by
$C I_S(u,v)^{1/2}$. Taking the supremum proves the first inequality.
If $B>0$, the functions $(u+B)/(2B)$ and $(v+B)/(2B)$ are
normalized testers. Their difference is $f/(2B)$, so
$I_S(u,v)\leq4B\mathfrak d_S(u,v)$. The case $B=0$ is trivial.
\end{proof}

\begin{corollary}[Equivalent relative topologies]
\label{cor:r3-strong-relative-topology}
For uniformly bounded common-collar potentials, convergence of
$I_S(u_j,u)$ to zero, convergence of $\mathfrak d_S(u_j,u)$ to
zero, and local outer $S$-capacity convergence on all permitted
compact/domain pairs are equivalent. The corresponding Cauchy
conditions are equivalent as well. These statements concern
relative capacity classes, not equality of actual values at
capacity-zero points.
\end{corollary}

\begin{proof}
Combine Theorem~\ref{thm:r3-relative-energy-dual} with
Theorem~\ref{thm:r3-I-capacity} and
Corollary~\ref{cor:r3-energy-capacity-topology}. For completeness,
the capacity-Cauchy implication for pairs follows directly from
uniform coefficient-capacity control:
\[
 I_S(u_j,u_k)
 \leq\int_L|u_j-u_k|(\mu_{u_j}+\mu_{u_k})
 \leq C\delta+
 C'\operatorname{Cap}_S^*(L\cap\{|u_j-u_k|>\delta\};\Omega).
\]
First let $j,k$ tend to infinity, then $\delta\downarrow0$.
Reverse framed domain comparison supplies the all-domain local
form of fixed-domain capacity convergence.
\end{proof}

\subsection{Localized strong convergence without a collar}

\begin{theorem}[Capacity convergence implies strong measure convergence]
\label{thm:r3-relative-capacity-strong-dual}
Let $u_j,u\in\mathcal C_b(U)$ have one common actual bound on every
compact subset. Suppose $u_j\to u$ locally in outer $S$-capacity,
on all permitted compact/domain pairs. For every original smooth
compact scalar $0\leq\chi\leq1$,
\[
 d_{\mathrm{psh},U}(\chi\mu_{u_j},\chi\mu_u)\longrightarrow0.
\]
Neither continuity nor monotonicity of the potentials is required.
\end{theorem}

\noindent\textbf{Status: \texttt{NEW\_PROVED}.}

\begin{proof}
A finite nonnegative smooth partition reduces the cutoff to
one prescribed regularization box. On this box choose original
compact cutoffs $0\leq\eta\leq\kappa\leq1$, with $\eta=1$ near
$\operatorname{supp}\chi$ and $\kappa=1$ near
$\operatorname{supp}\eta$. Put
\[
 f_j=u_j-u,\qquad
 Q_j=\sum_{b=0}^{q-1}(Du_j)^b(Du)^{q-1-b}S,\qquad
 e_j=\int\eta\,\Gamma_{Q_j}(f_j).
\]
All these expressions use the established bounded energy calculus.

First let $a$ be a smooth classical psh function with $0\leq a\leq2$
on a neighborhood of the cutoff stages. Stokes gives
\begin{equation}
 \int a\chi\,D(f_jQ_j)
 =-\int\chi\,\Gamma_{Q_j}(a,f_j)
          -\int a\,\Gamma_{Q_j}(\chi,f_j).
 \label{eq:r3-local-strong-green}
\end{equation}
Here is a precise limiting justification. Fix $j$ and the background
$Q_j$ first, and smooth only $u_j,u$ downward on the prescribed
compact box. The smooth identity passes by distributional and
energy convergence against the fixed smooth tests $\chi$ and
$\kappa a$. A smooth cutoff is locally a difference of two smooth
classical functions after adding a sufficiently large convex
quadratic in the unit moment coordinates. The established
bilinear energy identity therefore applies. This passage does not
pair a changing nonsmooth multiplier with a changing coefficient
measure.

The chain identity, $a\geq0$, and relative CLN show
\[
 \int\eta\,\Gamma_{Q_j}(a)
 \leq\int\kappa\,\Gamma_{Q_j}(a)
 \leq2|\mu_{Q_j,D\kappa}|(\operatorname{supp}\kappa)\leq C.
\]
The energy of the fixed smooth $\chi$ is uniformly bounded by the
same coefficient estimates. Apply the weighted Cauchy--Schwarz
inequality to each term in
\eqref{eq:r3-local-strong-green}; $a\leq2$ controls the second.
One obtains
\begin{equation}
 \left|\int a\chi\,d(\mu_{u_j}-\mu_u)\right|
                  \leq C_\chi e_j^{1/2}.
 \label{eq:r3-local-strong-bound}
\end{equation}
The constant is independent of this tester.

Now let $a\in\mathcal C_b(U)$ be any actual normalized tester.
Its restriction has a smooth decreasing classical row on the
fixed compact stages; late terms lie between zero and two there.
For fixed $j$, bounded convergence against the fixed finite signed
measure $\chi(\mu_{u_j}-\mu_u)$ passes
\eqref{eq:r3-local-strong-bound} to $a$. No energy limit involving
this tester is needed. The same bound holds uniformly for every
normalized actual tester on $U$.

It remains to prove $e_j\to0$. The energy expansion is
\[
 e_j=\tfrac12\langle f_j^2Q_j,D\eta\rangle
                    -\int\eta f_j\,d(\mu_{u_j}-\mu_u).
\]
All coefficient measures in this formula belong to explicit mixed
products over the one fixed $S$. Formula
\eqref{eq:r2-approx-coefficient-cap}, with the common local bounds,
gives uniform mass and capacity control on the fixed stages.
Splitting into $|f_j|\leq\delta$ and $|f_j|>\delta$ bounds the first
term in absolute value by
$C\delta^2+C'\operatorname{Cap}_S^*(\{|f_j|>\delta\}\cap H;\Omega)$
and the second by
$C\delta+C'\operatorname{Cap}_S^*(\{|f_j|>\delta\}\cap H;\Omega)$,
for a fixed compact stage $H$. Let $j\to\infty$ and then
$\delta\downarrow0$. Thus $e_j\to0$.
Taking the supremum in \eqref{eq:r3-local-strong-bound} and summing
the finitely many cutoff pieces proves the theorem.
\end{proof}

\begin{corollary}[Bounded monotone and mixed strong convergence]
\label{cor:r3-bounded-mixed-strong-dual}
The localized psh-dual monotone convergence theorem
\ref{thm:r3-psh-dual-monotone} holds for locally uniformly bounded
actual classical decreasing rows, without requiring that their
terms be continuous. More generally, mixed top-degree products
whose finitely many slots converge locally in outer $S$-capacity
with common bounds converge in the localized psh-dual seminorm.
\end{corollary}

\begin{proof}
The proved decreasing convergence in relative capacity and
Theorem~\ref{thm:r3-relative-capacity-strong-dual} give the first
assertion. For the second use finite polarization,
\[
 q!T_{S,q}(a_1,\ldots,a_q)
 =\sum_{J\subset\{1,\ldots,q\}}(-1)^{q-|J|}
       \left(D\sum_{i\in J}a_i\right)^qS.
\]
Every sum has common local bounds and the required capacity
convergence. The triangle inequality finishes the proof.
\end{proof}

\paragraph{Scope.}
The background may charge ambient pluripolar sets. The theorem
does not replace relative quasi-everywhere equality by pointwise
equality, solve the general potential-quotient problem, or allow
an arbitrary changing background $S_j$. For $q=0$ the differentiated
measure is the fixed background itself, so the corresponding
convergence assertion is immediate.

% END INCLUDED SOURCE: sections/strong_relative_topology


% BEGIN INCLUDED SOURCE: sections/completed_strong_topology
\section{Strong convergence to relative completed potentials}
\label{sec:r3-completed-strong-topology}

Keep the arbitrary fixed background $S$ and the original psh-dual
seminorm of Section~\ref{sec:r3-strong-relative-topology}. For
$a\in\mathcal A_{S,b}(U)$, write
\[
 \mathcal\mu_a=\mathcal T_{S,q}(a,\ldots,a),\qquad q=n-p\geq1.
\]
This section does not assume that $a$ has an ambient classical
representative.

\begin{corollary}[Strong convergence to the completed measure]
\label{cor:r3-completed-strong-dual}
Let $u_j\in\mathcal C_b(U)$ have common actual bounds on compact
subsets and converge locally in outer $S$-capacity, on all
permitted compact/domain pairs, to a locally bounded Borel function
$a$. Then $a\in\mathcal A_{S,b}(U)$ and, for every original
smooth compact scalar $0\leq\chi\leq1$,
\begin{equation}
 d_{\mathrm{psh},U}
       \bigl(\chi(Du_j)^qS,\chi\mathcal\mu_a\bigr)
                     \longrightarrow0.
 \label{eq:r3-completed-strong-dual}
\end{equation}
The limit uses the actual completed relative potential, not its
ambient usc regularization. In particular, whenever the proved
bounded relative-capacity completion supplies such a limit to a
capacity-Cauchy sequence, its localized measures are Cauchy and
converge in this strong seminorm as well.
\end{corollary}

\noindent\textbf{Status: \texttt{DIRECT\_COROLLARY}} of the
pair estimate in
Theorem~\ref{thm:r3-relative-capacity-strong-dual} and the
completed weighted BT calculus.

\begin{proof}
Capacity convergence implies the corresponding capacity Cauchy
condition. Apply the proof of
Theorem~\ref{thm:r3-relative-capacity-strong-dual} to a pair
$u_j,u_k$. On each of its finitely many cutoff stages put
\[
 f_{j,k}=u_j-u_k,\qquad
 Q_{j,k}=\sum_{b=0}^{q-1}(Du_j)^b(Du_k)^{q-1-b}S.
\]
The constants are uniform for both indices, and the proof gives
\[
 d_{\mathrm{psh},U}\bigl(\chi(Du_j)^qS,\chi(Du_k)^qS\bigr)
       \leq C\left(\int\eta\,\Gamma_{Q_{j,k}}(f_{j,k})
                                                   \right)^{1/2}
\]
on a single stage; for a general cutoff one sums finitely many
such bounds. The energy expansion and uniform coefficient-capacity
estimates in that proof show that every displayed energy tends
to zero as $j,k\to\infty$. This step involves only classical
potentials and explicit mixed backgrounds over the one fixed $S$.
No energy measure of an unproved Borel representative is introduced.

The completed weighted calculus already supplies
$\mathcal\mu_a$ as a positive Radon measure. For each fixed
normalized tester $b\in\mathcal C_b(U)$ it gives
\[
 \int b\chi(Du_k)^qS\longrightarrow\int b\chi\,d\mathcal\mu_a.
\]
Indeed $b$ is simply a constant sequence in the weighted slot.
Given $\epsilon>0$, choose $N$ so that the preceding pair
seminorm is at most $\epsilon$ for all $j,k\geq N$. For
fixed $j\geq N$ and fixed $b$, let $k\to\infty$ to obtain
\[
 \left|\int b\chi\,d\bigl((Du_j)^qS-\mathcal\mu_a\bigr)\right|
                         \leq\epsilon.
\]
The bound is independent of $b$. Taking its supremum only now
proves \eqref{eq:r3-completed-strong-dual}. Thus the argument
does not interchange a weak limit with an uncontrolled supremum.
\end{proof}

\begin{corollary}[Stability within the bounded completed class]
\label{cor:r3-strong-within-completion}
Suppose $a_j,a\in\mathcal A_{S,b}(U)$ converge locally in outer
$S$-capacity and, locally, their defining classical approximants
have one common actual bound. Then
\[
 d_{\mathrm{psh},U}(\chi\mathcal\mu_{a_j},
                         \chi\mathcal\mu_a)\longrightarrow0
\]
for every cutoff as above. The corresponding statement holds for
mixed top-degree products with finitely many convergent slots and
common defining bounds.
\end{corollary}

\begin{proof}
Work first on one of the fixed local cutoff stages. For every
$j$, choose a sufficiently late classical defining approximant
$v_j$ of $a_j$ so that its capacity error tends to zero on the
finitely many larger compact stages used by the local estimate,
and its localized measure differs from that of $a_j$ in
psh-dual seminorm by at most $1/j$. Such a choice is possible
by Corollary~\ref{cor:r3-completed-strong-dual}, applied to the
defining sequence for this fixed $j$. The common defining bound
keeps all constants uniform. A countable diagonal, if necessary,
also makes $v_j\to a$ in the local all-domain capacity sense.
The same corollary applied to $v_j$ gives its strong convergence
to $\mathcal\mu_a$. The triangle inequality proves the assertion.
Local tester estimates apply to restrictions of normalized testers
on $U$, and a finite partition of the compact cutoff gives the
stated global seminorm bound. Finally use the finite polarization
identity of Corollary~\ref{cor:r3-bounded-mixed-strong-dual}
for the mixed statement. No classical representative of the
completed limits is required.
\end{proof}

\paragraph{Scope.}
This strengthens stability of the already defined relative
completed operator. It does not make the ambient classical class
closed under these limits, and it does not add an arbitrary
varying-background convergence theorem. The explicit completed
solution in Proposition~\ref{prop:r3-relative-completed-dirichlet-example}
is one application where an ambient classical zero-boundary
solution is impossible.

% END INCLUDED SOURCE: sections/completed_strong_topology


% BEGIN INCLUDED SOURCE: sections/pluripolar_decomposition
\section{A canonical pluripolar decomposition of positive closed currents}
\label{sec:r3-pluripolar-decomposition}

The background in this section is arbitrary. The purpose of the
decomposition is to locate the part on which replacing an actual
increasing limit by its ambient upper regularization can lose
information. The construction uses actual coefficient measures,
not an unspecified ambient volume measure.

We work over the prescribed finite extension $K/\mathbf Q_p$.
A current is positive in the original CLD cone, is symmetric,
and is separately $d'$- and $d''$-closed. All tests are original
compactly supported smooth CLD forms. A current $T$ is said to be
\emph{carried by} a Borel set $E$ if $\mathbf1_E T=T$.
This is a coefficient-measure assertion; it does not assert that
the topological support of $T$ is contained in a nonclosed $E$.

\subsection{A finite measure detecting all coefficient-null sets}

\begin{lemma}[A countable family of coefficient controls]
\label{lem:r3-countable-trace-controls}
Let $U$ be an ordinary open in the analytification of a smooth
integral affine $K$-variety of dimension $n$. Fix $0\leq p\leq n$.
There is a countable family of positive original forms
$\Theta_l\in A_c^{n-p,n-p}(U)$ with this property: for every
original compact smooth complementary test $\eta$, there are
finitely many indices $l_1,\ldots,l_N$ and a finite constant
$C_\eta$ such that every original-positive $(p,p)$-current $T$
satisfies the inequality of Radon measures
\[
 |\mu_{T,\eta}|
       \leq C_\eta\sum_{j=1}^N\mu_{T,\Theta_{l_j}}.
 \tag{1}\label{eq:r3-countable-trace-control}
\]
The controls depend on $U,p$, but not on $T$.
\end{lemma}

\begin{proof}
Choose a countable dense subfield $K_0\subset K$ containing the
finitely many coefficients needed to define the affine variety
and its fixed affine embedding. Such a field exists since $K$ is
finite over $\mathbf Q_p$ and is separable as a topological field.
There are countably many rational functions with coefficients
in $K_0$ and countably many finite lists of these functions.

Every finite analytic unit presentation can locally be replaced
by one of these lists without changing any logarithmic absolute
value. Indeed B, Lemma D.1, first replaces the analytic units
on a compact affinoid neighborhood by rational functions over
$K$. Write their numerators and denominators in the fixed
affine coordinates. The numerators and denominators are units
on that compact neighborhood. Approximate their finitely many
coefficients by elements of $K_0$ closely enough that the
uniform errors are smaller than the positive minima of the
original absolute values. The ultrametric inequality preserves
each numerator and denominator norm, and hence each rational
function norm, at every point of the neighborhood. Their
original smooth logarithmic differentials also agree on its
ordinary interior. This is an exact replacement of the moment
coordinates, not a limiting identification of currents.

For each finite rational list, let $O$ be its algebraic unit
locus intersected with $U$, and write its logarithmic
coordinates as $t_1,\ldots,t_s$. The space $O$ is locally
compact and second countable, so choose a countable compact
exhaustion with interiors covering $O$, and nonnegative
original smooth compact cutoffs $\chi_j$ equal to one near
these compact sets. Put $r=n-p$ and include the forms
\[
 \chi_j\frac1{r!}
       \left(\sum_{i=1}^s d't_i\wedge d''t_i\right)^r
\]
in the countable family, extending them by zero off $O$.
For $r=0$ the parenthesized factor is one. Original smooth
cutoffs and compact exhaustions are available by the same
ordinary topology and partition-of-unity inputs as in R1.

Now divide a given compact test $\eta$ into finitely many
pieces on neighborhoods where the exact rational replacements
apply. Each piece has one of the listed moment presentations.
On its fixed compact support, its symmetric coefficient
matrix is bounded by a multiple of the positive definite
matrix of the displayed power. Choose $\chi_j=1$ near that
support. R's finite-test cone domination, recalled in
Lemma~\ref{lem:finite-positive-controls}, then gives
\eqref{eq:r3-countable-trace-control}. More explicitly, its
cone inequalities remain valid after multiplication by any
original smooth scalar on a common refinement; Riesz
uniqueness therefore makes the domination an inequality of
coefficient measures. Positive currents annihilate the
antisymmetric null part. Summing the finitely many pieces
proves the assertion.
\end{proof}

\begin{corollary}[A faithful finite control measure]
\label{cor:r3-faithful-control-measure}
For every such positive current $S$, there is a finite positive
Radon measure $\lambda_S$ on $U$ such that, for every Borel
$E\subset U$,
\[
 \lambda_S(E)=0
 \quad\Longleftrightarrow\quad
 \mathbf1_E S=0.
 \tag{2}\label{eq:r3-faithful-null-ideal}
\]
The measure need not be canonical; its null ideal is canonical.
\end{corollary}

\begin{proof}
Set $\mu_l=\mu_{S,\Theta_l}$. Each is a finite positive
Radon measure, since $\Theta_l$ has compact support. Define
\[
 \lambda_S=\sum_{l\geq1}
       \frac{2^{-l}}{1+\mu_l(U)}\,\mu_l.
\]
Its total mass is at most one. It is Radon: finite initial
sums are Radon, and the tails have uniformly small total
mass, so compact inner approximation passes to the sum.
If $\lambda_S(E)=0$, then $\mu_l(E)=0$ for every $l$;
the preceding lemma makes every coefficient restriction
$\mathbf1_E\mu_{S,\eta}$ vanish. Conversely,
$\mathbf1_E S=0$ makes all $\mu_l(E)$ vanish. This proves
\eqref{eq:r3-faithful-null-ideal}.
\end{proof}

These two control statements are \texttt{DIRECT\_COROLLARY}
of the original coefficient positivity in R, exact moment
replacement in B, and separability of the prescribed field.
No globally strictly positive form is used. The countable
family may contain arbitrarily many different moment
presentations.

\subsection{The maximal pluripolar part}

Call $S$ \emph{non-pluripolar} if
$\mathbf1_E S=0$ for every Borel locally pluripolar set $E$.
This term refers to all its coefficient measures, not only
to one selected top-degree trace.

\begin{theorem}[Pluripolar decomposition on an exhaustion domain]
\label{thm:r3-pluripolar-decomposition}
Let $V$ be a connected exhaustion domain of
Definition~\ref{def:r2-exhaustion-domain}, of positive dimension.
Every allowed positive closed current $S$ on $V$ has a unique
decomposition
\[
 S=S_{\rm pp}+S_{\rm np}
 \tag{3}\label{eq:r3-polar-decomposition}
\]
into positive symmetric separately closed currents such that
$S_{\rm pp}$ is carried by a complete pole set in $V$ and
$S_{\rm np}$ is non-pluripolar.
There is a nonpositive, nontrivial extended classical psh
function $\phi$ on $V$, with $Q=\{\phi=-\infty\}$, for which
\[
 S_{\rm pp}=\mathbf1_QS,
 \qquad S_{\rm np}=\mathbf1_{V\setminus Q}S.
 \tag{4}\label{eq:r3-polar-split-by-witness}
\]
The two currents are independent of the control measure and
the chosen witness. They commute with ordinary restriction,
addition of positive closed currents, and nonnegative scaling.
\end{theorem}

\begin{proof}
Use the finite control measure $\lambda_S$ above. Let
$\mathscr P$ consist of the complete pole sets of nonpositive,
nontrivial extended classical psh functions on $V$, and put
\[
 m=\sup_{P\in\mathscr P}\lambda_S(P)\leq\lambda_S(V).
\]
The empty set belongs to $\mathscr P$. Choose $P_j\in\mathscr P$
with $\lambda_S(P_j)\to m$, and choose nonpositive witnesses
$\phi_j$. Fix one ambient divisorial point $z\in V$.
By the R2 no-divisorial-poles theorem, every $\phi_j(z)$ is
finite. Choose positive $b_j$ such that
$b_j|\phi_j(z)|\leq2^{-j}$. The decreasing sum
\[
 \phi=\sum_{j\geq1}b_j\phi_j
\]
is extended classical psh and is nontrivial, since
$\phi(z)\geq-1$. This is the negative-series construction
already used in the R2 pole-witness theorem. Its pole set
$Q$ contains every $P_j$, so
$m\leq\lambda_S(Q)\leq m$. In particular the supremum is
attained by this one complete pole set.

Let $E\subset V$ be Borel and locally pluripolar.
The R2 same-domain pole-patching theorem provides a member
of $\mathscr P$ containing $E$. Applying the same negative
sum construction to that witness and $\phi$ gives
$R\in\mathscr P$ containing $E\cup Q$. Maximality gives
\[
 \lambda_S(E\setminus Q)
       \leq\lambda_S(R)-\lambda_S(Q)=0.
\]
By \eqref{eq:r3-faithful-null-ideal},
$\mathbf1_{E\setminus Q}S=0$. Thus the second current in
\eqref{eq:r3-polar-split-by-witness} is non-pluripolar.
The R2 pole-restriction theorem makes both displayed
restrictions positive and separately closed. This proves
existence, including the actual coefficient interpretation.

For uniqueness, suppose $S=A+B$ is another such
decomposition, with $A$ carried by a locally pluripolar
Borel set $P$ and $B$ non-pluripolar. Positivity gives
$0\leq A\leq S$. Since
$\mathbf1_{P\setminus Q}S=0$, positivity of coefficient
controls gives $\mathbf1_{P\setminus Q}A=0$. Hence
$\mathbf1_Q A=A$. Also $\mathbf1_QB=0$. Restricting
$S=A+B$ to $Q$ now gives $A=S_{\rm pp}$ and
$B=S_{\rm np}$.

These properties are local. In particular restriction of
a non-pluripolar current to an ordinary open $W\subset V$
is still non-pluripolar. To check this precisely, exhaust
$W$ by compact subsets: a locally pluripolar set intersected
with one of these compacts is locally pluripolar in $V$,
and coefficient measures pass to the countable union.
Uniqueness therefore identifies the restricted decomposition
with the one constructed on a smaller exhaustion domain.
For sums, the sum of the non-pluripolar parts is
non-pluripolar, and the sum of the pluripolar parts is
carried by the union of their two pole sets, contained in
one complete pole set by the negative-sum construction.
Uniqueness proves additivity. Scaling is immediate.
\end{proof}

\begin{corollary}[Local decomposition on the allowed analytic spaces]
\label{cor:r3-local-polar-decomposition}
On any allowed analytic space, the preceding local
decompositions glue to a canonical decomposition
$S=S_{\rm pp}+S_{\rm np}$. Locally, $S_{\rm pp}$ is carried
by a complete pole set; $S_{\rm np}$ gives zero mass to
every locally pluripolar Borel set in all coefficient measures.
No global countability assumption on the analytic space is
required for this local assertion.
\end{corollary}

\begin{proof}
Use B's ordinary algebraic neighborhoods and the R2 local
basis of exhaustion domains. The uniqueness just proved
identifies the decompositions on every smaller such domain
in an overlap. Currents and their positivity and separate
closedness are local for ordinary opens, so the pieces glue.
In dimension zero the only locally pluripolar set is empty;
the assertion there is $S_{\rm pp}=0$, $S_{\rm np}=S$.
\end{proof}

\subsection{Compatibility with the bounded Bedford--Taylor calculus}

\begin{theorem}[Decomposition commutes with bounded mixed products]
\label{thm:r3-polar-decomposition-bt}
For every $0\leq k\leq n-p$ and $u_1,\ldots,u_k\in\Cb(U)$,
\[
 \bigl(T_{S,k}(u_1,\ldots,u_k)\bigr)_{\rm pp}
       =T_{S_{\rm pp},k}(u_1,\ldots,u_k),
\]
\[
 \bigl(T_{S,k}(u_1,\ldots,u_k)\bigr)_{\rm np}
       =T_{S_{\rm np},k}(u_1,\ldots,u_k).
\]
These are current equalities against all original tests.
In particular all mixed products over a non-pluripolar
background have non-pluripolar coefficient measures.
\end{theorem}

\begin{proof}
Work on an exhaustion neighborhood where
$S_{\rm pp}=\mathbf1_QS$ as above. Linearity of the
background in the recursive construction gives
\[
 T_{S,k}=T_{S_{\rm pp},k}+T_{S_{\rm np},k}.
\]
The R2 pole-restriction identity gives
$\mathbf1_QT_{S_{\rm pp},k}=T_{S_{\rm pp},k}$.
For any locally pluripolar Borel set $E$, choose a complete
pole set $P$ containing it on this neighborhood. Since
$\mathbf1_PS_{\rm np}=0$, the same identity gives
\[
 \mathbf1_PT_{S_{\rm np},k}
      =T_{\mathbf1_PS_{\rm np},k}=0.
\]
Thus the two terms have exactly the required polar and
non-pluripolar properties. Uniqueness of the decomposition,
applied to the resulting positive closed current, identifies
them with its canonical parts. Locality proves the assertion
on $U$.
\end{proof}

\paragraph{Status and limits.}
The decomposition and its BT compatibility are recorded as
\texttt{NEW\_PROVED} relative to the proved R2 pole
restriction and pole-patching theorems. Original smooth
positive backgrounds and bounded ambient BT backgrounds are
non-pluripolar by R2's ambient trace theorem. Subvariety
integration backgrounds can instead be entirely pluripolar.
This is not an analytic-stratum or Siu decomposition: no
analyticity of a pole set is asserted. Nor is the capacity
of a sum asserted to be the sum of the capacities, since
the same test potential enters both summands.


% END INCLUDED SOURCE: sections/pluripolar_decomposition


% BEGIN INCLUDED SOURCE: sections/pluripolar_zero_extension
\section{Zero extension across a closed complete pluripolar set}
\label{sec:r3-zero-extension}

The restriction theorem of R2 starts with a current on the whole
ordinary open set. Here the current is initially defined only on the
complement of the exceptional set. Its closedness across that set is
part of the conclusion. The proof adapts the classical El Mir
argument, but the finite-mass hypothesis and the first-order estimate
are formulated for all original CLD coefficient tests.

The base field, analytic category, and original positivity cone are
unchanged. A closed set $P\subset U$ is \emph{locally complete
pluripolar} here if, near each of its points, it is exactly the pole
set of a nontrivial extended classical psh function. On a smaller
neighborhood such a witness can be taken nonpositive. No assertion
below replaces this equality by containment in an unspecified
larger pole set.

\subsection{The mass condition and a first-order estimate}

Let $S$ be a symmetric original-positive $(p,p)$-current on
$U\setminus P$, and put $q=n-p$. For an original compact smooth
$(q,q)$-form $\eta$ on $U$, denote by $\mu_{S,\eta}$ the coefficient
measure on $U\setminus P$ obtained by restricting the form there.
This notation does not require that its restricted support be
compact in $U\setminus P$.

We say that $S$ has \emph{locally finite coefficient mass across
$P$} if, locally on $U$,
\begin{equation}
 \label{eq:r3-finite-mass-across-pole}
 |\mu_{S,\eta}|(U\setminus P)<\infty
 \quad\hbox{for every }\eta\in A_c^{q,q}(U).
\end{equation}
Equivalently, it suffices to require this for the finite positive
controls of each original test stage. In particular the condition
includes controls for the coordinates of compact cutoffs. It is
not a finite-mass assumption for just one initially selected
tropical trace.

\begin{lemma}[Coefficientwise extension and weighted first-order control]
\label{lem:r3-zero-extension-controls}
Under \eqref{eq:r3-finite-mass-across-pole}, the coefficientwise
zero extension
\[
 \langle\widetilde S,\eta\rangle
       =\int_{U\setminus P}d\mu_{S,\eta}
\]
is an original-positive symmetric coefficient-Radon current on $U$,
with $\mathbf1_P\widetilde S=0$.

Suppose $q\geq1$. Fix a moment presentation of an original
$(q-1,q)$-test $\alpha$ near its compact support, let $t_1,\ldots,t_s$
be its moment coordinates, and put
\[
 \beta=\sum_i d't_i\wedge d''t_i.
\]
Let $\chi\geq0$ be an original compact cutoff equal to one near
$\operatorname{supp}\alpha$, with support in that moment
presentation's ordinary domain. For every original smooth real $f$
and real Borel $b$ for which the displayed integrals are finite,
\begin{equation}
 \label{eq:r3-weighted-first-order-control}
 \left|\int_{U\setminus P}b\,d'f\wedge\alpha\wedge S\right|^2
 \leq C_\alpha
 \left(\int_{U\setminus P}\chi\,d'f\wedge d''f
                           \wedge\beta^{q-1}\wedge S\right)
 \left(\int_{U\setminus P}\chi b^2\beta^q\wedge S\right).
\end{equation}
There is the analogous inequality for $(q,q-1)$-tests and $d''f$.
The constant depends on the fixed test presentation and its bounded
coefficients, and not on additional moment coordinates used by $f$
or by bounded approximations to $b$. For a test with several
presentations one uses a finite partition and the corresponding
finite sum of these controls.
\end{lemma}

\begin{proof}
Extend each coefficient measure by zero on the closed set $P$.
The local finite-mass assumption makes these extensions Radon.
Their linear relations and common-refinement identities hold
off $P$ and both sides have zero restriction to $P$, so the same
identities hold on $U$. On each fixed test stage, the positive
controls in Lemma~\ref{lem:finite-positive-controls} give
an order-zero bound with finite right hand side. Hence the
extended functionals are continuous in the original test
topology and define a current. Positivity and symmetry pass
through zero extension. This proves the first assertion without
using either closedness or a Bedford--Taylor product.

For the second assertion, expand $\alpha$ as a finite sum of
bounded smooth coefficients times
$d't_I\wedge d''t_J$, where $|I|=q-1$ and $|J|=q$.
Use the standard Lagerberg signs for squares. On a common
refinement, set
\[
 A_I=d'f\wedge d't_I,\qquad B_J=b\,d't_J.
\]
The coefficient matrix of $S$ defines a positive semidefinite
bilinear measure on these $q$-vectors: the positivity needed
here is duality with the original positive form cone. The
two-by-two square inequality, followed by the scalar
Cauchy--Schwarz inequality for its densities, bounds each
mixed term by the square root of
\[
 \int\chi\, [A_I]^2\wedge S
       \quad\hbox{times}\quad
 \int\chi\, [B_J]^2\wedge S.
\]
Here $[A_I]^2$ denotes the positive square with that standard
sign. Bounded coefficients of $\alpha$ contribute a fixed
factor. Summing the first squares is bounded by a fixed
multiple of $d'f\wedge d''f\wedge\beta^{q-1}$, and summing
the second squares is bounded by a fixed multiple of
$b^2\beta^q$. The finite-sum Cauchy--Schwarz inequality
gives \eqref{eq:r3-weighted-first-order-control}.

One may first take $b$ smooth and then use coefficient-Radon
measure identities and bounded approximation, or apply the
same two-by-two inequality to Radon densities directly.
Pullback to additional moment coordinates preserves the
original positive cone. Thus the bounds use the fixed
coordinates of $\alpha$ but impose no common fixed
presentation on the varying $f$. The $d''$ assertion follows
by the CLD involution. A finite original partition proves
the last assertion.
\end{proof}

\subsection{Smooth cutoffs adapted to the pole set}

\begin{lemma}[Bounded psh cutoffs with a zero neighborhood]
\label{lem:r3-psh-zero-cutoffs}
Work on a sufficiently small ordinary algebraic neighborhood,
and let $H$ be a compact norm-box core in a larger such
neighborhood. Suppose $\phi\leq0$ is an extended classical
psh function there and $P=\{\phi=-\infty\}$ is closed.
On one fixed ordinary neighborhood of $H$ there are original
smooth classical psh functions $w_N$, with semipositive
original $(1,1)$-forms $Dw_N$, such that
\[
 0\leq w_N\leq1+N^{-2}\quad\hbox{on }H,\qquad
 w_N=0\quad\hbox{near }P\cap H,
\]
and $w_N(x)\to1$ for every $x\in H\setminus P$.
No monotonicity in $N$ is asserted or needed.
\end{lemma}

\begin{proof}
The truncation $\phi_N=\max(\phi,-2N)$ is bounded classical
psh on the larger neighborhood. By the prescribed-box
regularization in Lemma~\ref{lem:compact-box-regularization},
it has smooth classical psh approximants with semipositive
original curvature forms decreasing to
$\phi_N$ on a fixed smaller neighborhood containing $H$.
The smaller geometric box can be chosen independently of
$N$; the projective line and regularization index need
not be independent of $N$.

Choose one approximant $\psi_N\geq\phi_N$ such that
\[
 \psi_N<N^{-1}\quad\hbox{on }H,\qquad
 \psi_N<-3N/2\quad\hbox{on }P\cap H.
\]
Both choices follow from compactness and the decreasing
continuous rows: their limits are at most zero on $H$
and equal $-2N$ on $P\cap H$. The second inequality
holds on an ordinary neighborhood of $P\cap H$.

Let $F$ be the convolution of $t\mapsto\max(t,0)$ with
an even nonnegative smooth probability kernel supported
in $(-1/4,1/4)$. Then $F$ is nonnegative, smooth, convex
and nondecreasing, equals zero for $t\leq-1/4$, and
equals $t$ for $t\geq1/4$. Put
\[
 w_N=F(1+\psi_N/N).
\]
Convex nondecreasing composition preserves the classical
psh property, by the all-extensions curve definition;
the result is also an original smooth function. Moreover,
$Dw_N=F'(1+\psi_N/N)D\psi_N/N+
F''(1+\psi_N/N)d'\psi_N\wedge d''\psi_N/N^2$
is a semipositive original form.
The stated upper bound follows by monotonicity of $F$,
and the negative inequality on $\psi_N$ gives the zero
neighborhood. If $\phi(x)$ is finite, then, for all large
$N$, $\psi_N(x)\geq\phi(x)$ and
$1+\phi(x)/N>1/4$. Hence
\[
 1+\phi(x)/N\leq w_N(x)\leq1+N^{-2}.
\]
This proves the pointwise limit. No diagonal approximation
at all points is involved.
\end{proof}

\begin{theorem}[A Skoda--El Mir zero-extension theorem]
\label{thm:r3-zero-extension}
Let $P$ be a closed locally complete pluripolar subset of
an allowed ordinary open $U$. Let $S$ be a symmetric
original-positive $(p,p)$-current on $U\setminus P$,
separately $d'$- and $d''$-closed there. Assume
\eqref{eq:r3-finite-mass-across-pole}. Then its
coefficientwise zero extension $\widetilde S$ to $U$
is separately $d'$- and $d''$-closed.
All assertions are on all original compact smooth tests.
\end{theorem}

\begin{proof}
Lemma~\ref{lem:r3-zero-extension-controls} constructs $\widetilde S$ and proves
its positivity, symmetry and coefficient-Radon property.
If $p=n$, separate closedness follows from degree. Suppose
$q=n-p\geq1$. Closedness is local for ordinary opens.
Near a point of $P$, choose a nonpositive complete pole
witness and nested norm boxes as in
Lemma~\ref{lem:r3-psh-zero-cutoffs}. A finite partition
reduces each compact test to this situation and, within
it, to one fixed moment presentation.

Fix an original $(q-1,q)$-test $\alpha$. Choose a
nonnegative compact original smooth $\chi$, equal to
one near its support, with support in a compact core
$H$ on which the cutoffs $w_N$ are defined and within
the ordinary domain of that fixed moment presentation. Put
$\beta=\sum_i d't_i\wedge d''t_i$ for the coordinates
of $\alpha$. The positive current
$S\wedge\beta^{q-1}$ is separately closed on the
complement of $P$. The function $w_N$ is zero on
an ordinary neighborhood of $P\cap H$, so every term
below can be integrated by parts on that complement.
One can insert a further compact cutoff away from $P$
equal to one on the relevant supports; its differentiated
terms multiply zero functions.

The smooth identity
$D(w_N^2)=2w_NDw_N+2d'w_N\wedge d''w_N$ gives
\begin{align*}
 E_N
 &:=\int_{U\setminus P}
        \chi\,d'w_N\wedge d''w_N\wedge\beta^{q-1}\wedge S\\
 &=\frac12\int_{U\setminus P}
        w_N^2D\chi\wedge\beta^{q-1}\wedge S
   -\int_{U\setminus P}
        \chi w_NDw_N\wedge\beta^{q-1}\wedge S\\
 &\leq 2\,|\mu_{S,D\chi\wedge\beta^{q-1}}|
                         (\operatorname{supp}\chi)
 =:C<\infty .
\end{align*}
Here $0\leq w_N\leq2$ on the fixed compact set, and
the last term being subtracted is nonnegative.
The finite constant concerns the one fixed smooth
form $D\chi\wedge\beta^{q-1}$, including the moment
coordinates of $\chi$. This step has used closedness
only where $S$ was originally given.

Set
\[
 a_N=2w_N-w_N^2=1-(1-w_N)^2.
\]
On $H$ one has $0\leq a_N\leq1$; it is zero near
$P\cap H$, and tends to one at every point of
$H\setminus P$. Consequently $a_N\alpha$ is a
valid compact test on $U\setminus P$. Separate
closedness there gives, with the harmless sign
determined by total degree,
\[
 \left|\int_{U\setminus P}a_Nd'\alpha\wedge S\right|
 =2\left|\int_{U\setminus P}
              (1-w_N)d'w_N\wedge\alpha\wedge S\right|.
\]
By \eqref{eq:r3-weighted-first-order-control}, its
square is at most
\[
 4C_\alpha E_N
       \int_{U\setminus P}\chi(1-w_N)^2\beta^q\wedge S .
\]
The latter integral tends to zero by dominated
convergence for the fixed finite coefficient measure
$\chi\beta^q\wedge\widetilde S$. Its restriction to
$P$ is zero by construction. Uniform boundedness of
$E_N$ therefore makes the preceding error tend to zero.
Dominated convergence for the fixed signed coefficient
measure of $d'\alpha$ identifies the left hand side
in the limit with $|\langle\widetilde S,d'\alpha\rangle|$.
This proves $d'\widetilde S=0$. The $d''$ estimate and
the same argument prove $d''\widetilde S=0$.

Each localization and inequality uses original
forms and preserves common-refinement identities.
No convergence statement for a varying background
current or an unbounded multiplier was invoked.
\end{proof}

\subsection{Minimality and extension of bounded products}

\begin{corollary}[The least positive extension]
\label{cor:r3-minimal-zero-extension}
If $R$ is any symmetric original-positive separately closed
extension of $S$ to $U$, then
\[
 R=\widetilde S+A,
\]
where $A$ is positive, symmetric, separately closed and
carried by $P$. Conversely, adding any such $A$ gives
an extension. Thus $\widetilde S$ is the least positive
extension. It is the unique extension whose coefficient
measures give zero mass to $P$.
\end{corollary}

\begin{proof}
Coefficientwise, the difference equals $\mathbf1_PR$,
which is positive and symmetric. The theorem proves
that $\widetilde S$ is closed, so the difference is
separately closed by subtraction. It is carried by
$P$. The converse and uniqueness follow immediately
from the coefficient restrictions. No assertion is
made that a positive current carried by $P$ must
have an analytic-support description.
\end{proof}

\begin{corollary}[Bounded mixed products extend without extra mass]
\label{cor:r3-zero-extension-bt}
Let $u_1,\ldots,u_k\in\Cb(U)$, where $0\leq k\leq n-p$.
The current $T_{S,k}(u_1|_{U\setminus P},\ldots,
u_k|_{U\setminus P})$ has locally finite coefficient
mass across $P$, and its coefficientwise zero extension
is
\[
 T_{\widetilde S,k}(u_1,\ldots,u_k).
\]
This current gives zero coefficient mass to $P$.
The same coefficient equality holds after multiplication
by any locally bounded Borel weight on $U$.
\end{corollary}

\begin{proof}
The theorem makes $\widetilde S$ an admissible fixed
background on $U$. Its bounded mixed products exist
and have locally finite coefficients by the established
relative calculus and relative CLN estimate. Ordinary
restriction identifies them with the asserted products
on $U\setminus P$. The R2 pluripolar restriction identity
gives
\[
 \mathbf1_P T_{\widetilde S,k}(u_1,\ldots,u_k)
 =T_{\mathbf1_P\widetilde S,k}(u_1,\ldots,u_k)=0.
\]
These two coefficient assertions identify the product
with the zero extension and prove the mass claim.
Actual Borel multiplication preserves the equality.
The potentials in this corollary are required on $U$;
it does not assert a removable-singularity theorem
for arbitrary potentials given only off $P$.
\end{proof}

\paragraph{Status and dependencies.}
The zero-extension theorem is \texttt{NEW\_PROVED} in
the prescribed non-Archimedean setting. The proof uses
R's original coefficient cone and B's prescribed-box
regularization. Its cutoff mechanism is the classical
Skoda--El Mir mechanism; no novelty is claimed for
that classical argument. The mass and first-order
lemma makes the required transfer to all original
tests explicit. The last two corollaries are
\texttt{DIRECT\_COROLLARY}, the last also using the
R2 pole-restriction identity. Mere support avoidance
in CLD's zero-extension convention does not prove
the finite-mass extension theorem stated here.

% END INCLUDED SOURCE: sections/pluripolar_zero_extension


% BEGIN INCLUDED SOURCE: sections/zero_extension_counterexamples
\section{Two sharp scope tests for the zero-extension theorem}
\label{sec:r3-zero-extension-counterexamples}

These examples test two hypotheses of
Theorem~\ref{thm:r3-zero-extension}. They do not concern the
previously proved restriction theorem for a current already
defined on the whole space. The field and its absolute value
remain the prescribed finite $p$-adic ones; every form below
is an original CLD form.

\begin{proposition}[Closed pluripolar containment is insufficient]
\label{prop:r3-zero-extension-noncomplete-false}
There are a closed pluripolar set $E$ in a smooth two-dimensional
ordinary analytic domain $U$ and an original-positive symmetric
$(1,1)$-current $S$ on $U\setminus E$ such that:
\begin{enumerate}
\item $S$ is separately $d'$- and $d''$-closed;
\item all its original coefficient masses are locally finite
      across $E$;
\item its coefficientwise zero extension to $U$ is not closed.
\end{enumerate}
The exceptional set is not locally complete pluripolar.
Thus replacing \emph{closed locally complete pluripolar} by
\emph{closed pluripolar} in Theorem~\ref{thm:r3-zero-extension}
is \texttt{FALSE}.
\end{proposition}

\begin{proof}
Let $\pi$ be a uniformizer and $\lambda=-\log|\pi|>0$. Use
coordinates $(t,z)$ on $\mathbf G_{m,K}\times\mathbf A^1_K$ and put
\[
 s=\log|t|,\qquad
 U=\{-\lambda<s<\lambda,\ |z|<1\},\qquad
 Y=U\cap\{z=0\},\qquad E=Y\cap\{s=0\}.
\]
The set $U$ is an ordinary open in the smooth affine surface,
and $Y$ is its closed smooth analytic curve. No normalization
of the prescribed absolute value has been changed.
The set $E$ is closed in $U$. It is pluripolar because
\[
 E\subset Y=\{\log|z|=-\infty\},
\]
and $\log|z|$ is a nontrivial extended classical psh function
on $U$.

On $U\setminus E$, define the current of integration over
the component
\[
 Z=Y\cap\{s>0\}.
\]
Equivalently, this current is $[Y]$ locally where $s>0$
and is zero locally where $s<0$ or $z\ne0$. These neighborhoods
cover $U\setminus E$, so it is a well-defined symmetric
original-positive separately closed current $S$.
There is no boundary point of $Z$ inside $U\setminus E$.
Its zero extension is the actual Borel multiplication
\[
 T=\mathbf1_{\{s>0\}}[Y]\quad\hbox{on }U.
\]
For every original compact smooth $(1,1)$-test $\eta$,
\[
 |\mu_{S,\eta}|(U\setminus E)
       \leq|\mu_{[Y],\eta}|(U)<\infty.
\]
Thus the full coefficient-mass condition holds, including all
test presentations, not just the radial one used below.

Choose $f\in C_c^\infty((-\lambda,\lambda))$ with $f(0)=1$.
Choose an original smooth scalar cutoff $\kappa(z)$ on
$\{|z|<1\}$, compactly supported there and equal to one on a
neighborhood of $z=0$. One explicit construction takes a
smooth function of $\log|z|$ on $z\ne0$, equal to one for
$\log|z|\leq-\lambda$ and zero for
$\log|z|\geq-\lambda/2$, and gives it value one at zero.
It is original smooth because it is constant on an ordinary
neighborhood of zero. Now put
\[
 \alpha=\kappa(z)f(s)d''s\in A_c^{0,1}(U).
\]
Its support is compact: it is contained in the product of a
closed smaller $z$-disc and a closed smaller $t$-annulus.
Pullback to $Y$ gives
\[
 d'\alpha|_Y=f'(s)d's\wedge d''s.
\]
The original calibration of the identity tropicalization of
$\mathbf G_m$ has multiplicity one. Hence the defining
CLD integral, equivalently its one-dimensional Lebesgue
formula, gives
\begin{equation}
\label{eq:r3-annular-boundary-current}
 \langle T,d'\alpha\rangle
   =\int_{\{s>0\}\subset Y}f'(s)d's\wedge d''s
   =\int_0^\lambda f'(s)\,ds=-1.
\end{equation}
Here $f$ vanishes near $\lambda$.
Thus $d'T\ne0$. Symmetry, or the same computation with a
$(1,0)$-test, also gives $d''T\ne0$. The nonzero boundary
functional is an original test functional; no additional
$\delta$-test is being introduced.

For completeness, the local-completeness failure can be
verified independently of the extension theorem. Let
$x\in Y$ be the unit Gauss point of the $t$-annulus.
Suppose on an ordinary neighborhood $W$ of $x$ that
$E\cap W=\{\phi=-\infty\}$ for an extended classical psh
function $\phi$. Choose a connected ordinary neighborhood
$C$ of $x$ in $Y\cap W$. Such neighborhoods exist by local
connectedness of analytic curves; the same fact is used in
\cite[proof of Proposition~3.5(3)]{GR25}.
The restriction $\phi|_C$ is extended classical psh.

The map sending $r$ to the radius-$e^r$ Gauss point is
continuous. Since $C$ is a neighborhood of $x$, it contains
points of this map with sufficiently small $r>0$, which are
outside $E$. Thus $\phi|_C$ is not identically $-\infty$.
But $x$ is a divisorial point of this smooth curve and
$\phi(x)=-\infty$, contradicting the proved
Theorem~\ref{thm:r2-no-divisorial-poles} applied on $C$.
This proves that $E$ is not locally complete pluripolar.
\end{proof}

\paragraph{Interpretation.}
The proposition is an explicit \texttt{NEW\_PROVED}
counterexample within the allowed original analytic category.
Its purpose is the necessity of the complete-pole hypothesis
for this extension statement. It is not a counterexample to
the extension theorem with all its hypotheses. The current
has already finite mass; strengthening only the mass bound
cannot repair the false version.

\begin{example}[The local mass hypothesis cannot be omitted]
\label{ex:r3-zero-extension-infinite-mass}
Let $U=\{|t|<1\}\subset\mathbf A_K^{1,\mathrm{an}}$ and
$P=\{0\}=\{\log|t|=-\infty\}$. For $j\geq1$ let
$x_j=\zeta_{0,|\pi|^j}$ be the corresponding Gauss point and set
\[
 \mu=\sum_{j=1}^{\infty}\delta_{x_j}
       \quad\hbox{on }U\setminus P.
\]
The points $x_j$ tend to zero. Every compact subset of
$U\setminus P$ has a positive minimum for $|t|$, hence meets
only finitely many of the $x_j$. Thus $\mu$ is a positive
Radon measure there, and defines an original top-degree
$(1,1)$-current; separate closedness is automatic in that
degree. The closed set $P$ is a complete pole set.

Any positive Radon extension to $U$ would give infinite mass
to the compact closed disc $\{|t|\leq|\pi|\}$, since that disc
contains every $x_j$. This is impossible. Therefore no
positive coefficient-Radon extension exists. This elementary
example has remaining degree $q=0$; it tests deletion of the
mass hypothesis from the all-bidegrees theorem and is not
presented as a new remaining-degree-one obstruction.
\end{example}


% END INCLUDED SOURCE: sections/zero_extension_counterexamples


% BEGIN INCLUDED SOURCE: sections/potential_quotient
\section{Potential quotients: normal traces and the remaining regularity gate}
\label{sec:r3-potential-quotient}

Fix the original absolute value on a finite extension $K/\mathbf Q_p$.
All currents, positive cones, and compactly supported tests in this
section are original CLD objects. The functions in $\Cb$ retain their
actual values and the all-complete-extensions classical definition.
For the normal-coordinate statements, the spaces are ordinary opens
in smooth algebraic analytifications over this same $K$.

The question is whether
\begin{equation}
 D[u]=D[v],\quad u,v\in\Cb(U)
 \quad\Longrightarrow\quad
 D\bigl((u-v)S\bigr)=0
 \label{eq:r3-general-potential-quotient}
\end{equation}
for every symmetric original-positive, separately $d'$- and
$d''$-closed current $S$ of bidegree $(p,p)$, $p<n$.
The full implication remains \texttt{OPEN}. The results below give
additional backgrounds and a precise sufficient regularity statement.
They do not replace the quantifier over arbitrary $S$ by a restricted
one in the main research question.

\subsection{Bounded classical weights through a normal Gauss kernel}

Let $B$ be an ordinary open in a smooth $d$-dimensional
$K$-variety, let $\mathbb D_s=\{|t|<s\}$, and write
\[
 W=B\times_K\mathbb D_s,\qquad
 \pi:W\longrightarrow B,\qquad i:B\longrightarrow W
\]
for the projection and zero section. For $a\in B$ and $0<r<s$,
let $\sigma_r(a)$ be the image in $W$ of the radius-$r$ Gauss
point of the $\mathscr H(a)$-disc over $a$.

Choose a smooth convex function $\psi:\mathbf R\to\mathbf R$
which equals $0$ on $(-\infty,-\epsilon]$ and equals the
identity on $[\epsilon,\infty)$, where $\epsilon>0$ is fixed.
Put
\[
 b_N(t)=-N+\psi(\log|t|+N),\qquad \beta_N=Db_N.
\]
The function is extended by the constant $-N$ at $t=0$.
For sufficiently large $N$, $\beta_N$ is a smooth strongly positive
closed form supported in
$e^{-N-\epsilon}\leq|t|\leq e^{-N+\epsilon}<s$.

\begin{lemma}[The normal trace for a bounded classical weight]
\label{lem:r3-bounded-normal-trace}
For $u\in\Cb(W)$ and every original
$\eta\in A_c^{d,d}(B)$,
\begin{equation}
 \lim_{N\to\infty}\int_W u\,\beta_N\wedge\pi^*\eta
       =\int_B (u\circ i)\eta.
 \label{eq:r3-normal-trace}
\end{equation}
The same conclusion holds for any finite real linear combination
of such $u$. All integrals use the prescribed Borel representatives.
\end{lemma}

\begin{proof}
Let $\mu_\eta$ be the signed Radon measure of the smooth top
form $\eta$ on $B$. The product calibration and Borel coarea give
\begin{equation}
 \int_W u\,\beta_N\wedge\pi^*\eta
 =\int_B\left(\int_{\mathbf R}
      u\bigl(\sigma_{e^x}(a)\bigr)\psi''(x+N)\,dx\right)
      d\mu_\eta(a).
 \label{eq:r3-normal-coarea}
\end{equation}
Here the inner integral is restricted to the support of
$\psi''(x+N)$, which lies below $\log s$.
For clarity, this is coarea for smooth forms, with a Borel scalar
coefficient, rather than an external product of arbitrary currents.
On a finite moment presentation of the base top form, keep its
calibrated skeleton cells separate. Adjoining the coordinate $t$
gives their products with the radial interval; integration on a
product cell is the base calibrated measure times $dx$. The normal
coefficient is $\psi''(x+N)$. Borel Fubini proves
\eqref{eq:r3-normal-coarea} for bounded Borel coefficients, exactly
as in the cellwise argument of \textup{B}, Lemma A.3. A finite
partition and signed linearity recover every $\eta$. Any structural
degree of the base is already in $\mu_\eta$; the normal factor has
mass one, since $\int\psi''(x+N)\,dx=1$.
The Gauss lift is continuous in $(a,r)$ on an affinoid base patch:
for a polynomial in $t$ its value is
$|\sum_i f_it^i|(\sigma_r(a))=\max_i|f_i(a)|r^i$,
and uniform approximation gives the assertion for analytic
functions on a compact radial range. Thus the Borel integrand
in \eqref{eq:r3-normal-coarea} is evaluated on actual analytic
points and is measurable.

Fix $a\in B$. Pull $u$ back to the $\mathscr H(a)$-disc over
$a$. This is a locally bounded subharmonic function by the
all-extensions classical hypothesis. Its maximum on a closed
radius-$r$ disc occurs at the Gauss point, so
\[
 u(i(a))\leq u(\sigma_r(a)).
\]
As $r\downarrow0$, $\sigma_r(a)\to i(a)$ in the ordinary
topology. Upper semicontinuity gives the opposite limiting upper
bound. Therefore
\begin{equation}
 u(\sigma_r(a))\longrightarrow u(i(a)).
 \label{eq:r3-radial-actual-limit}
\end{equation}
This argument concerns only the pullback of the given function to
a curve. It does not assume base extension of a current equality.

Choose $s_0<s$ and a compact base neighborhood of
$\operatorname{supp}\eta$. Its product with the closed
radius-$s_0$ disc is contained in a finite union of compact
analytic domains in $W$. Local boundedness of $u$ gives one
bound there. For large $N$, every normal radius in
\eqref{eq:r3-normal-coarea} is at most $s_0$.
The inner measures are probability measures supported at radii
tending uniformly to zero. Equation
\eqref{eq:r3-radial-actual-limit} and dominated convergence against
the fixed finite measure $|\mu_\eta|$ prove
\eqref{eq:r3-normal-trace}. Finite real linearity proves the
last assertion.
\end{proof}

This is a short \texttt{NEW\_PROVED} extension of the continuous
normal-kernel argument in \textup{R}, Lemma 5.1. The additional
input is the actual radial limit supplied by classical
subharmonicity; weak convergence of $\beta_N$ alone would not
justify the weighted limit.

\begin{proposition}[Restriction of a bounded difference with positive curvature]
\label{prop:r3-normal-restriction}
Let $f=u-v$, with $u,v\in\Cb(W)$. If $D[f]$ is an
original-positive current on $W$, then $D[f\circ i]$ is
an original-positive current on $B$. If $D[u]=D[v]$, then
\[
 D[u\circ i]=D[v\circ i]
\]
on every original compact smooth complementary test on $B$.

Consequently, if $j:Z\hookrightarrow X$ is a smooth closed
immersion of smooth $K$-varieties and $u,v\in\Cb(U)$ on an
ordinary $U\subset X^{\an}$ satisfy $D[u]=D[v]$, then
\begin{equation}
 D[u|_{Z^{\an}\cap U}]=D[v|_{Z^{\an}\cap U}].
 \label{eq:r3-subvariety-restriction}
\end{equation}
\end{proposition}

\begin{proof}
The assertion has no curvature content when $d=0$.
For $d\geq1$, take a positive original
$\alpha\in A_c^{d-1,d-1}(B)$. The form
$\beta_N\wedge\pi^*\alpha$ is positive, smooth, and compactly
supported in $W$. Separate closedness of $\beta_N$ and
distributional integration by parts give
\[
 0\leq\langle D[f],\beta_N\wedge\pi^*\alpha\rangle
   =\int_W f\,\beta_N\wedge\pi^*(D\alpha).
\]
The normal trace lemma, applied to $u$ and $v$ separately, gives
the limit $\langle D[f\circ i],\alpha\rangle$.
This proves positivity. In the equality case the same computation
starts with zero and applies to every original test $\alpha$,
without a positivity restriction; it proves the full current equality.

By \textup{R}, Lemma 5.2, each point of the smooth immersion
has one ordinary analytic product neighborhood with a normal
polydisc over an ordinary neighborhood in $Z^{\an}$.
Apply the zero-section result successively to its normal
coordinates. At every step the restricted functions remain actual
locally bounded classical psh functions by their definition.
Locality gives \eqref{eq:r3-subvariety-restriction}.
No equality on tests after an arbitrary extension of $K$ has been
used or asserted.
\end{proof}

\begin{corollary}[Backgrounds carried by a smooth subvariety]
\label{cor:r3-subvariety-potential-quotient}
In the last setting, put $h=u-v$. Then
\[
 D(h[Z^{\an}\cap U])=0.
\]
More generally the same conclusion holds for $S=j_*Q$ when,
locally on $Z^{\an}\cap U$, $Q$ is a finite linear combination
of bounded-potential ambient BT products on $Z$, or is a smooth
separately closed form current. The pushforward is taken only
in the degrees where it is defined. In particular this includes
every positive closed background of either displayed kind.
\end{corollary}

\begin{proof}
Borel multiplication and the proper local pushforward satisfy
$h\,j_*Q=j_*((h|_Z)Q)$ on all original tests, with the
original calibration. The restricted curvature equality is
\eqref{eq:r3-subvariety-restriction}. For ambient bounded-potential
products on $Z$, the symmetry argument of
Lemma~\ref{lem:potential-independence-ambient-mixed} puts that
curvature in the innermost differentiated slot; for a smooth
closed form, use the distributional Leibniz rule. Thus
$D((h|_Z)Q)=0$. Pushforward commutes with $D$, proving the claim.
The first assertion is the case $Q=[1]$.
\end{proof}

The corollary is a \texttt{DIRECT\_COROLLARY} of the normal trace
and the previously proved restricted quotient statements. The
relative background $[Z]$ is allowed to charge the ambient
pluripolar set $Z$; ambient nonconcentration has not been imposed.

\subsection{A further closure class supplied by R2}

\begin{corollary}[Bounded ambient BT weak closures]
\label{cor:r3-bounded-bt-quotient-closure}
Let $u,v\in\Cb(U)$ satisfy $D[u]=D[v]$. Suppose that locally
on each fixed norm-box domain one has
\[
 S_j=T_{[1],p}(a_{1,j},\ldots,a_{p,j})\longrightarrow S,
 \qquad 0\leq a_{i,j}\leq1,
 \qquad 0\leq p<n,
\]
on all original compact smooth tests. Then
$D((u-v)S)=0$. No monotonicity or convergence of the
generating potentials is required.
\end{corollary}

\begin{proof}
For every $j$, ambient BT symmetry gives $D(uS_j)=D(vS_j)$.
Corollary~\ref{cor:r2-fixed-slots-closure}, applied separately
to the fixed weights $u$ and $v$, gives
$uS_j\to uS$ and $vS_j\to vS$. Distributional differentiation
passes the equality to the limit. The local argument glues.
\end{proof}

This is a \texttt{DIRECT\_COROLLARY} of R2, not a regularization
theorem for arbitrary currents. The uniform bounded-potential
family provides the capacity control in that result. Neither an
arbitrary positive closed current nor a current supported on an
ambient polar set has been shown to admit such an approximation.

\subsection{Continuous representatives: a precise sufficient gate}

\begin{lemma}[A continuous current representative is the actual difference]
\label{lem:r3-continuous-representative-actual}
Let $u,v\in\Cb(U)$ and let $H$ be a finite continuous function
on $U$. If
\[
 [u]-[v]=[H]
\]
on every original compact smooth top-degree test, then
$u-v=H$ at every point of $U$.
No classical-psh hypothesis on $H$ is required.
\end{lemma}

\begin{proof}
Work in a framed affine ordinary neighborhood as in
Section~\ref{sec:r2-gap-detection}. At a prescribed divisorial
point $z$, choose its continuous model tester $a$ with
$(Da)^n(\{z\})>0$. On one smaller neighborhood, its pure-FS
construction gives smooth classical psh $a_j\downarrow a$.
For every original compact smooth scalar $\chi$,
\[
 \int \chi(u-v-H)(Da_j)^n=0.
\]
The fixed-background weighted decreasing theorem passes the
$u$ and $v$ terms to their actual Borel limits. The $H$ term
passes by continuity of the fixed multiplier and weak convergence
of $(Da_j)^n$. Hence
$(u-v-H)(Da)^n=0$ as a signed Radon measure. Its atom at $z$
gives $u(z)-v(z)=H(z)$.

This holds at every divisorial point in the Shilov determining
family of the neighborhood. The function $v+H$ is finite and
upper semicontinuous, although it has not been declared psh.
Lemma~\ref{lem:r2-shilov-gap}, applied to the classical function
$u$ and the upper-semicontinuous function $v+H$, gives
$u\leq v+H$ everywhere. Apply the same lemma to $v$ and
$u-H$ for the reverse inequality. Locality finishes the proof.
\end{proof}

The lemma is a \texttt{DIRECT\_COROLLARY} of weighted convergence,
the divisorial testers, and Shilov determination. It is a recovery
statement \emph{after} a continuous representative has been supplied.

\begin{proposition}[A sufficient local regularity statement]
\label{prop:r3-quotient-regularity-gate}
Assume the following statement on the permitted analytic spaces:
\begin{quote}
For $u,v\in\Cb(U)$ with $D[u]=D[v]$, every point has an
ordinary neighborhood $W$ and a finite continuous $H_W$ such that
$[u]-[v]=[H_W]$ on $W$.
\end{quote}
Then \eqref{eq:r3-general-potential-quotient} holds for every
allowed positive separately closed background $S$.
\end{proposition}

\begin{proof}
Lemma~\ref{lem:r3-continuous-representative-actual} gives
$u-v=H_W$ at the actual points of $W$. Distributional
differentiation gives $D[H_W]=0$. The continuous-harmonic
multiplier lemma from R1, based on \textup{R}, Theorem 2.9,
then gives $D(H_WS)=0$ on $W$. Indeed $H_W$ and $-H_W$
are locally uniform limits of smooth functions with positive
curvature forms; multiplying those forms by $S$, passing the
uniform limit through its coefficient-Radon measures, and using
the pointed positive cone gives both signs of the required zero
current. These local equalities glue.
\end{proof}

The proposition is \texttt{CONDITIONAL}; the displayed existence
of a continuous representative remains \texttt{OPEN}. This is
strictly a sufficient reduction, not a claimed equivalence.
The known faithfulness of $u\mapsto[u]$ does not construct
$H_W$. Likewise $u-v$ is not a psh slot merely because both
$u$ and $v$ are psh.

\begin{remark}[The external harmonicity statements have different scopes]
In CLD's 2025 revision, Section 7.8 begins with a boundaryless
\emph{tropical space} of level $n$. Definition 7.8.2 defines its
harmonic sheaf using locally real linear combinations of
paralinear harmonic functions. Remark 7.8.3 discusses both the
possible gap between that sheaf and continuous functions with
$D[u]=0$, and the unresolved continuity, up to current-equivalence,
of locally integrable functions with $D[u]=0$ in that tropical
setting. It is not an established theorem about the present
bounded differences of classical psh functions over $K$, and it
is not used as a proof that our specific gate is equivalent to
that broader problem.

Similarly, GR, Definition 7.1 and Proposition 7.2, make a function
pluriharmonic when it and its negative are classically psh; its
continuity then follows from the two upper-semicontinuity
conditions. The current equality $D[u]=D[v]$ has not been shown
to put either sign of $u-v$ in that class. These are source
comparisons, not substitute regularity proofs.
\end{remark}

\subsection{A false extension to arbitrary ambient-capacitary representatives}

\begin{example}[Ambient current zero does not imply a relative multiplier zero]
\label{ex:r3-polar-representative-quotient}
Let $R>1$ and
\[
 U=\{|z|<R,\ |t|<R\}\subset(\mathbf A^2_K)^{\an},
 \qquad Y=\{z=0\}\cap U,
 \qquad q(t)=\max(0,\log|t|).
\]
Define the actual Borel function
\[
 f(x)=\begin{cases}q(t(x)),&x\in Y,\\0,&x\notin Y.\end{cases}
\]
It is bounded and upper semicontinuous: $Y$ is closed, $q$ is
continuous and nonnegative, and the outside value is zero.
It is locally ambient-capacity-quasicontinuous, since it vanishes
off the analytic polar set $Y$, which has local outer ambient
capacity zero. It represents the zero class in the ambient
capacitary completion.

Every original smooth top-degree trace measure on $U$ gives
zero mass to $Y$. Hence $[f]=0$ and $D[f]=0$. Nevertheless,
with the original positive separately closed background $S=[Y]$,
\[
 D(fS)=i_*D[q]=\delta_{(0,\zeta_{0,1})}\ne0,
\]
where $\zeta_{0,1}$ is the unit Gauss point of the $t$-disc.
The atom has mass one in the original $K$ calibration, by the
one-variable model formula for $\max(0,\log|t|)$.

This does not refute \eqref{eq:r3-general-potential-quotient}:
$f$ cannot be a difference $a-b$ of two members of $\Cb(U)$.
Otherwise $[a]=[b]$ and R2 current-faithfulness would force
$a=b$ at every point, whereas $f$ is nonzero on part of $Y$.
The example instead proves \texttt{FALSE} the shortcut which
replaces the actual difference hypothesis by boundedness,
upper semicontinuity, and ambient quasicontinuity. Its ingredients
are direct consequences of the already proved polar and model
measure facts; it is not counted as a new BT construction.
\end{example}

\subsection{What the normal trace does not supply}

The normal-kernel proof gives restriction along a smooth immersion
over the original field. It does not establish that an equality
of original currents pulls back to every complete field extension,
or to an arbitrary analytic curve map after such an extension.
Nor does it approximate an arbitrary positive closed current by
smooth-subvariety integration currents with multiplier control.
The bounded-difference restriction statement and the continuous
representative gate are therefore kept separate. The unresolved
step in \eqref{eq:r3-general-potential-quotient} is not hidden in
an unproved change of test category or in an exchange of two
discontinuous limits.

% END INCLUDED SOURCE: sections/potential_quotient


% BEGIN INCLUDED SOURCE: sections/potential_quotient_projection
\section{The projection route for potential quotients: a precise test gate}
\label{sec:r3-potential-quotient-projection}

This section supplements Section~\ref{sec:r3-potential-quotient}.
It does not resolve its unrestricted potential-independence question.
All spaces, analytic products, and absolute values are over the original
finite extension $K/\mathbf Q_p$. All tests are original CLD smooth tests.
Let $U$ be an allowed ordinary smooth $n$-dimensional algebraic
analytic neighborhood, $n\geq1$, let $\mathbb D_s$ be an ordinary
disc, and put $W=U\times_K\mathbb D_s$, with projection $\pi$.
Suppose $u,v\in\Cb(U)$ and $D[u]=D[v]$, and set $f=u-v$.
The functions $u\circ\pi,v\circ\pi$ are classical by the
all-extensions definition. Pulling back the equality of currents is
a separate question.

\begin{proposition}[A sufficient projected-test criterion]
\label{prop:r3-projection-continuous-coefficients}
For $\alpha\in A_c^{n,n}(W)$ define the compactly supported
ordinary current $P_\alpha=\pi_*[\alpha]$ of bidegree $(n-1,n-1)$.
The pushforward is defined because the projection is proper on the
support of $\alpha$. The current
\[
 \nu_\alpha=DP_\alpha=\pi_*[D\alpha]
\]
is a finite signed Radon measure, and
\begin{equation}\label{eq:r3-projection-exact-reduction}
 \langle D[f\circ\pi],\alpha\rangle
       =\int_U f\,d\nu_\alpha.
\end{equation}
Suppose, as a compact current on $U$, that
\[
 P_\alpha=\sum_{\ell=1}^N c_\ell[\eta_\ell],\qquad
 \eta_\ell\in A_c^{n-1,n-1}(U),
\]
where the $c_\ell$ are continuous and bounded on the relevant
compact supports. Then $\langle D[f\circ\pi],\alpha\rangle=0$.
\end{proposition}

\noindent\textbf{Status: \texttt{DIRECT\_COROLLARY}.}
The conclusion for every $\alpha$ is \texttt{CONDITIONAL} on the
displayed representation, or an equally strong pairing-continuity
statement, for every projected test.

\begin{proof}
The first identity is the definition of $D$ on the locally bounded
Borel function $f\circ\pi$, followed by pushforward of the fixed
signed measure $D\alpha$. It uses no pullback of nonsmooth currents.
Use a finite original smooth partition on the compact support to
reduce to a neighborhood on which the accepted regularization supplies
smooth psh $u_j\downarrow u$ and $v_j\downarrow v$, with uniform
local bounds. Write $f_j=u_j-v_j$.
Since $Du_j$ and $Dv_j$ converge weakly to the same positive current,
the fixed-continuous-multiplier theorem of R gives
\[
 \sum_\ell\int_U c_\ell Df_j\wedge\eta_\ell\longrightarrow0.
\]
For the smooth functions $f_j$, ordinary smooth pullback and Stokes
identify this sum with
\[
 \langle P_\alpha,Df_j\rangle
      =\int_U f_j\,d\nu_\alpha.
\]
Bounded convergence against the one fixed finite signed measure
$\nu_\alpha$ gives the asserted zero limit. If a representation is
initially stated only near the support, multiply it by a compact
cutoff equal to one there before using it as an identity of compact
currents. The finite partition reduction preserves continuity of the
coefficients and the degrees of the test forms.
\end{proof}

\begin{lemma}[Separated radial tests]
\label{lem:r3-radial-separated-kernel}
Let $I\Subset(-\infty,\log s)$, $x=\log|t|$, and let
$\gamma$ be a compactly supported Lagerberg form of bidegree
$(r,s')$ on $I$, where $r,s'\in\{0,1\}$. If
$\eta\in A_c^{n-r,n-s'}(U)$, then the form
\[
 \alpha=\pi^*\eta\wedge\gamma(x),
\]
extended by zero off the annulus, satisfies
$\langle D[f\circ\pi],\alpha\rangle=0$.
The same holds for finite sums, limits in the original compact test
topology, and analytic changes of centre $t'=t-a(x)$ for which
$a$ is analytic and $|a|<s$ on a neighborhood of the base support.
\end{lemma}

\noindent\textbf{Status: \texttt{DIRECT\_COROLLARY}.}

\begin{proof}
Use separate moment coordinates on the base and the radial annulus.
On each calibrated product cell the integration formula is iterated
integration in the base coordinates and the radial coordinate.
Consequently it is valid also for the bounded Borel multiplier pulled
back from $f$: first approximate a bounded Borel function by simple
functions in the resulting fixed coefficient measures.

If $r=s'=1$, the only contribution is the base pairing
$\langle D[f],\eta\rangle$ times the radial integral of $\gamma$,
and is zero. If $r=s'=0$, the only possible contribution contains
$D\gamma$ and its integral is zero by one-dimensional Stokes. In
the remaining two bidegrees the possible cross term has radial
factor $d'\gamma$ or $d''\gamma$, again with zero integral.
Degree considerations eliminate the other terms. This proves the
claim without differentiating $f$. Finite sums and convergence in
the test topology preserve the equality by the definition of a
current. Under the stated bound on $a$, translation is an analytic
automorphism of the disc over the relevant base neighborhood and
preserves $\pi$; applying the same computation in $t'$ proves the
last assertion. No assertion that these tests generate all original
tests on the analytic product is made.
\end{proof}

\begin{example}[The analytic product is not the topological product]
\label{ex:r3-analytic-product-not-topological-product}
The assertion that products of functions of the two factors are
dense in all original smooth functions on the Berkovich analytic
product is \texttt{FALSE}.
Choose a uniformizer $\varpi$, and work in an ordinary bidisc of
radius $R>1$ with coordinates $(x,t)$. Let $z_0$ be the two-variable
unit Gauss point. Let $z_1$ be the image of the two-variable unit
Gauss point in $(X,S)$ under
\[
 x=X,\qquad t=X+\varpi S.
\]
Both points project to the unit Gauss point in each factor. Indeed,
for a polynomial $P$, the Gauss norm of $P(X+\varpi S)$ is at most
the maximum coefficient norm of $P$, and its coefficient of $S^0$
is $P(X)$, which gives the opposite inequality. However,
\[
 |t-x|(z_0)=1,\qquad |t-x|(z_1)=|\varpi|.
\]
Thus the original smooth function $\log|t-x|$ on the ordinary
nonvanishing locus of $t-x$ separates the two points. Multiplying by an
original compact cutoff equal to one near both gives a compact
smooth separating function. Every finite sum of products of
functions from the two factors, and every uniform limit of such
sums on the compact set $\{z_0,z_1\}$, takes equal values at the
two points. Equivalently, the nonzero top-degree current
$\delta_{z_0}-\delta_{z_1}$ annihilates every such tensor test.
\end{example}

The example disproves this density shortcut, not the desired
projection-kernel implication. It also does not contradict the
translated radial tests, which can involve $t-x$.

\paragraph{What the published projection results supply.}
CLD's 2025 manuscript, Section 9.1, constructs the product of its
topological tropical spaces. Its factor-generated refinements in
Section 9.2 (in particular Corollary 9.2.4), and the tensor-current
construction and differential identity in 9.4.4 and 9.4.8, have
that product as their domain. The preceding example prevents their
factor-density conclusion from being substituted for the missing
assertion on the analytic fiber product.

Mihatsch, \emph{$\delta$-Forms on Lubin--Tate spaces},
arXiv:2112.10018v4, Theorem 3.23, printed page 58, provides flat
refined skeletons for locally dominant analytic morphisms.
Definition 2.18, printed page 22, is formulated with piecewise
smooth fiber integrals. Example 2.45(3), printed pages 41--42,
also illustrates that a balanced pushforward need not be a
$\delta$-form. It is an abstract sheaf-level warning, not a
counterexample for the present disc projection. These results do
not state the continuous-coefficient original-test representation
in Proposition~\ref{prop:r3-projection-continuous-coefficients}.
Gubler--K\"unnemann, Section 6.10, concerns proper pushforward of
$\delta$-currents and likewise does not prove that representation.

\noindent\textbf{Status: \texttt{FAILED\_ATTEMPT}.}
The inference ``flat after refinement, hence an admissible projected
original smooth test'' is not justified. The missing statement is
pairing continuity for the projected test through cell changes and
refinements, in exactly the original sheaf.

\paragraph{The conditional graph consequence.}
Suppose the projection-kernel implication has been established for
iterated disc products. Let $C$ be an ordinary open in a smooth
algebraic $K$-curve, with a closed immersion
$\iota:C\hookrightarrow\mathbb D^m$, and let $g:C\to U$ be an
analytic map. Assume explicitly that the graph $(g,\iota)$ is
locally in the smooth algebraic closed-immersion scope of the
already proved normal-trace restriction theorem; for example,
both maps may be induced by algebraic maps with an algebraic
closed embedding for $\iota$. This hypothesis is stronger than
merely requiring an arbitrary analytic closed immersion.
Its graph is then a permitted smooth closed immersion into
$U\times_K\mathbb D^m$.
Applying that normal-trace theorem to the pulled-back kernel
identity would give $D[f\circ g]=0$. This is a
\texttt{CONDITIONAL} same-$K$ consequence. It neither supplies
the kernel identity nor changes the ground field; it therefore
does not by itself verify the all-complete-extensions curve test
in Gubler--Rabinoff, Definition 4.1.

\begin{proposition}[The positive square of a kernel difference]
\label{prop:r3-kernel-positive-square}
Under $D[u]=D[v]$, the established bounded energy calculus gives
\begin{equation}\label{eq:r3-kernel-positive-square}
 \Gamma_{[1]}(f)
       =\tfrac12D[f^2]-fD[f]
       =\tfrac12D[f^2]\geq0.
\end{equation}
For every fixed smooth positive closed $(1,1)$-form $\beta$ and
every nonnegative compact smooth $\chi$ on its domain,
\[
 0\leq\langle\Gamma_{[1]}(f),\chi\beta^{n-1}\rangle
  =\tfrac12\int f^2D\chi\wedge\beta^{n-1}
  \leq\tfrac12\|f\|_{\infty,\operatorname{supp}D\chi}^2
                 \int|D\chi\wedge\beta^{n-1}|.
\]
\end{proposition}

\noindent\textbf{Status: \texttt{DIRECT\_COROLLARY}.}

\begin{proof}
On a regularization neighborhood use smooth decreasing psh rows
$u_j,v_j$ and $f_j=u_j-v_j$. The smooth currents
$d'f_j\wedge d''f_j$ are positive. Expand $f_jDf_j$ into its
four weighted psh terms; the established fixed-background weighted
monotone theorem gives the limit $fD[f]$. Bounded convergence gives
$D[f_j^2]\to D[f^2]$. The smooth square identity thus gives the
first formula and its positivity, which glue by locality. Pairing
with $\chi\beta^{n-1}$ and using that $\beta$ is separately closed
gives the second formula and the displayed elementary bound.
\end{proof}

This estimate is not zero energy. On a relatively compact annulus,
$f=\log|t|$, $u=f+C$ and $v=C$ are bounded classical potentials
and $D[f]=0$, whereas
$\Gamma_{[1]}(f)=d'\log|t|\wedge d''\log|t|$ has positive radial
integral. Nor is the Borel function $f^2$ shown continuous by the
positive-square identity. Consequently R's continuous-potential
input cannot be applied to it without a separate regularity proof.
CLD 2025, Remark 7.8.3, is in its enlarged tropical test category;
Gubler--Rabinoff 2025, Definition 7.1 and Proposition 7.2, concern
functions whose two signs satisfy the classical psh conditions.
Neither statement turns the original current equation $D[f]=0$
alone into that hypothesis.

\paragraph{Remaining gate.}
The unrestricted implication
$D[u]=D[v]\Rightarrow D((u-v)S)=0$ for an arbitrary allowed
positive closed background remains \texttt{OPEN}. The proved
projected-test criterion, the radial cases, and the positive-square
identity are retained; neither the tensor-density shortcut nor a
formal passage to another test sheaf closes the gap.

% END INCLUDED SOURCE: sections/potential_quotient_projection


% BEGIN INCLUDED SOURCE: sections/relative_compact_capacity
\section{Relative compact regularity: a linear-degree reduction and a contact obstruction}
\label{sec:r3-relative-compact-capacity}

The field, actual classical class $\Cb$, original CLD tests, and
original-positive cone remain those of the first two records. The
background is fixed throughout. No equality on additional
$\delta$-tests is asserted. The target of this section is the sharp
inequality
\[
 \nu(K)\leq\operatorname{Cap}_S(K;\Omega)
 \tag{RC}\label{eq:r3-sharp-relative-closure}
\]
for weak limits of normalized top relative BT measures on one fixed
domain. We do not prove \eqref{eq:r3-sharp-relative-closure} for an
arbitrary $S$. We prove a representation and an exact reduction when
the remaining degree is one, and give counterexamples to a proposed
contact-measure proof in every intermediate degree.

\paragraph{External comparison.}
For complex positive closed currents, Dabbek--Elkhadhra define the
analogous capacity in \cite[Definition~2.1, p.~471]{DE06}. Their list
of open problems on p.~472 explicitly asks for continuity on
decreasing compact sets, equality with a contact-constrained
capacity, and attainment of the capacity of a compact set by one
potential. This is an \texttt{EXTERNAL\_KNOWN} statement about that
paper; it is not a claim that all these questions remain open in the
present literature. In particular, that paper cannot be cited as
already proving the compact outer regularity sought here.

\subsection{An energy estimate for one remaining degree}

For this subsection let $S$ have bidegree $(n-1,n-1)$. It is
arbitrary within the fixed original-positive, symmetric, separately
closed class. Thus
\[
 C_{\Omega'}(E)=\operatorname{Cap}_S(E;\Omega')
   =\sup_{\substack{w\in\Cb(\Omega')\\0\leq w\leq1}}
        D(wS)(E).
\]
The relative energy measure $\Gamma_S(f,g)$ for differences and
real linear combinations of bounded classical potentials is the one
in Lemma~\ref{lem:relative-energy-measures}.
For an original smooth real scalar $\chi$ we also use
\[
 \Gamma_S(\chi)=d'\chi\wedge d''\chi\wedge S.
\]
This is a positive measure: the $(1,1)$-factor is strongly positive,
so pairing it with the specified original-positive $S$ is legitimate.
The cutoff is not being assumed psh.

Choose permitted local stages
\[
 V\Subset\Omega'\Subset\Omega,
\]
with $V$ ordinary open, and original smooth cutoffs
$0\leq\chi,\kappa\leq1$ compactly supported in $\Omega'$ such
that $\chi=1$ near $\overline V$ and $\kappa=1$ near
$\operatorname{supp}\chi$. The cutoffs may be chosen with support
inside a smaller compact norm box, so prescribed-domain
regularization applies near their entire supports. Put
\begin{equation}
\label{eq:r3-cutoff-energy-norm}
 N_\chi(f)^2
    =\int\chi\,d\Gamma_S(f)
       +\int f^2\,d\Gamma_S(\chi).
\end{equation}
This is the square of a positive semidefinite quadratic seminorm
on the real span of bounded classical potentials. The associated
bilinear form is obtained by polarization.

\begin{lemma}[A fixed-energy bound for relative capacity]
\label{lem:r3-linear-energy-capacity}
For the preceding stage there is a finite constant $A_{\chi,\kappa,S}$
such that, whenever $a,b\in\Cb(\Omega)$ satisfy $0\leq a,b\leq1$,
$f=a-b$, and $\delta>0$, one has
\begin{align}
 \int\chi f^2\,D(wS)
   &\leq A_{\chi,\kappa,S} N_\chi(f)
       &&(w\in\Cb(\Omega'),\ 0\leq w\leq1),
       \label{eq:r3-energy-square-bound}\\
 C_{\Omega'}^*(V\cap\{|f|>\delta\})
   &\leq A_{\chi,\kappa,S}\delta^{-2}N_\chi(f).
       \label{eq:r3-energy-capacity-bound}
\end{align}
The same bound holds with any compact subset of $V$ in place of
$V$. The constant is independent of $a,b,w$ and of all
regularization indices.
\end{lemma}

\begin{proof}
First take smooth $a,b,w$. For compactly supported expressions,
Stokes and symmetry of $S$ give
\begin{equation}
\label{eq:r3-square-green}
 \int\chi f^2D(wS)
    =-2\int\chi f\,d\Gamma_S(f,w)
       -\int f^2\,d\Gamma_S(\chi,w).
\end{equation}
There is no boundary term. The positive $2\times2$ matrices of
energy measures give their measure-valued Cauchy--Schwarz
inequality, also after bounded Borel multiplication. Since
$\chi\leq\kappa$ and $\kappa=1$ near the support of
$d\chi$, it follows that
\begin{align*}
 \left|\int\chi f\,d\Gamma_S(f,w)\right|
   &\leq \|f\|_\infty
       \left(\int\chi\,d\Gamma_S(f)\right)^{1/2}
       \left(\int\kappa\,d\Gamma_S(w)\right)^{1/2},\\
 \left|\int f^2\,d\Gamma_S(\chi,w)\right|
   &\leq \|f\|_\infty
       \left(\int f^2\,d\Gamma_S(\chi)\right)^{1/2}
       \left(\int\kappa\,d\Gamma_S(w)\right)^{1/2}.
\end{align*}
All norms in this proof can be taken on the fixed compact
neighborhood of the cutoff supports.

The energy of a normalized nonnegative tester is uniformly bounded.
Indeed, the defining identity for $\Gamma$ gives
\[
 \int\kappa\,d\Gamma_S(w)
  =\frac12\langle S,w^2D\kappa\rangle
       -\int\kappa w\,D(wS).
\]
The last integral is nonnegative. The first is bounded in absolute
value by the variation of the fixed coefficient measure
$\mu_{S,D\kappa}$ times $\|w\|_\infty^2/2$.
This measure is finite on the compact stage by the accepted
coefficient-Radon property. In particular one finite bound
$M_{\kappa,S}$ works for all $0\leq w\leq2$ near the stage.
Combining the preceding inequalities yields, in the smooth case,
\[
 \int\chi f^2D(wS)
 \leq \sqrt{5M_{\kappa,S}}\,\|f\|_\infty N_\chi(f).
\]

For general bounded classical $a,b,w$, use the simultaneous smooth
decreasing rows of Lemma~\ref{lem:compact-box-regularization} on
one neighborhood of the whole compact stage. After discarding a
prefix, each row lies in $[0,2]$ there. To see the upper bound,
the continuous approximants decrease to a function at most one;
their open sublevel sets below two form an increasing cover of
the compact stage. A finite subcover gives the uniform index.
The smooth differences therefore have absolute value at most two.

Proposition~\ref{prop:functional-weights}, applied to
$F(x,y)=(x-y)^2$, passes the left side to
$\int\chi f^2D(wS)$. The first summand of
$N_\chi(f)^2$ converges by the distributional formula for
$\Gamma$ and the weighted fixed-background theorem. The second
converges by bounded convergence against the fixed measure
$\Gamma_S(\chi)$. Passing the smooth inequality proves
\eqref{eq:r3-energy-square-bound}, for example with
$A_{\chi,\kappa,S}=2\sqrt{5M_{\kappa,S}}$.
This argument uses no continuity theorem with a varying background.

Since $\chi=1$ on $V$, Chebyshev's inequality and the supremum
over normalized $w$ give the estimate for the raw capacity of
$V\cap\{|f|>\delta\}$. Both $a$ and $b$ are relatively
quasicontinuous. That strict level set on the ordinary core $V$
is therefore quasi-open, and
Proposition~\ref{prop:r2-quasiopen-outer} identifies its raw and
outer capacities. This proves
\eqref{eq:r3-energy-capacity-bound}. Monotonicity gives the compact
version. In particular no strict level set has been silently
treated as ordinarily open.
\end{proof}

\paragraph{Status.}
Lemma~\ref{lem:r3-linear-energy-capacity} is
\texttt{NEW\_PROVED}. Its inputs are the established original-cone
energy measures, weighted monotone calculus, coefficient-Radon
control, and quasicontinuity. It needs no global strictly psh
exhaustion and no compact outer regularity.

\subsection{Convex blocks give actual capacitary limits}

The following class preserves the normalization on the original
domain; it is narrower than merely asking that an arbitrary
member of $\mathcal A_S^1(\Omega)$ take values in $[0,1]$.

\begin{definition}
\label{def:r3-normalized-completion}
Let $\Omega$ be a permitted ordinary algebraic norm-box domain.
Write $\mathcal A_S^{[0,1]}(\Omega)$ for the Borel functions
$a:\Omega\to[0,1]$ admitting a sequence
\[
 a_j\in\Cb(\Omega),\qquad 0\leq a_j\leq1\text{ on }\Omega,
\]
converging to $a$ locally in outer $S$-capacity in the all-domain
sense of Definition~\ref{def:r2-outer-convergence}. Actual Borel
representatives are retained. Its BT calculus is the restriction
of Theorem~\ref{thm:r2-completed-calculus}.
\end{definition}

\begin{theorem}[Linear-degree convex-block compactness]
\label{thm:r3-linear-convex-blocks}
Let $S$ have bidegree $(n-1,n-1)$ on the fixed norm-box domain
$\Omega$. Given any sequence $u_j\in\Cb(\Omega)$ with
$0\leq u_j\leq1$, there exist a subsequence $(u_{j_l})$, finite
convex combinations
\[
 v_m\in\operatorname{conv}\{u_{j_l}:l\geq m\},
\]
and $a\in\mathcal A_S^{[0,1]}(\Omega)$ such that
$v_m\to a$ locally in outer $S$-capacity on every permitted
compact/domain pair in $\Omega$.

If in addition $D(u_jS)\to\nu$ locally weakly on all original
compact smooth scalar tests, then
\begin{equation}
\label{eq:r3-linear-weak-representation}
 \nu=D(aS),\qquad
 D(v_mS)\longrightarrow D(aS),\qquad
 v_mD(v_mS)\longrightarrow aD(aS).
\end{equation}
The last convergence is for the convex blocks; it is not a
weighted convergence assertion about the original sequence.
\end{theorem}

\begin{proof}
Choose the countable cofinal pairs
$V_l\Subset\Omega_l\Subset\Omega$ used in
Proposition~\ref{prop:r2-completion}, and cutoffs
$\chi_l,\kappa_l$ as in the preceding lemma. Let $\mathscr V$
be the real span of all normalized classical potentials on
$\Omega$. For each $l$, write $N_l=N_{\chi_l}$. Positivity of
$\Gamma$ and the elementary $L^2$ term in
\eqref{eq:r3-cutoff-energy-norm} show that $N_l$ is a Hilbert
seminorm on $\mathscr V$.

There is a finite $B_l$ such that
\[
 N_l(u)^2\leq B_l
 \quad(u\in\Cb(\Omega),\ 0\leq u\leq1).
\]
For the first summand use the cutoff energy identity from the
preceding proof with $\chi_l$ in place of $\kappa$; for the
second use $u^2\leq1$ and the finite mass of
$\Gamma_S(\chi_l)$. Define
\begin{equation}
\label{eq:r3-total-hilbert-norm}
 \|f\|_{\mathcal H}^2
   =\sum_{l=1}^{\infty}
          \frac{2^{-l}}{1+B_l}N_l(f)^2.
\end{equation}
This is finite on $\mathscr V$: for
$f=\sum_{i=1}^r c_i u_i$, the seminorm triangle inequality
gives $N_l(f)^2\leq(\sum_i|c_i|)^2 B_l$.
Quotient by its null space and complete to obtain a real Hilbert
space $\mathcal H$. Each normalized $u_j$ has norm at most one.
This construction makes no claim that a Hilbert vector is
already a pointwise function or that a separately chosen
Dirichlet form is closable in a trace $L^2$ space.

The closed span of the sequence in $\mathcal H$ is separable.
Weak sequential compactness of its unit ball gives a subsequence
$u_{j_l}$ converging weakly to $h\in\mathcal H$. For every
$m$, $h$ belongs to the norm closure of the convex hull of the
tail $\{u_{j_l}:l\geq m\}$. Here is the needed elementary
Hilbert argument. If $h$ were outside that closed convex set and
$p$ its nearest point, then
\[
 \langle h-p,v-p\rangle\leq0
 \quad\text{for every }v\text{ in the set}.
\]
Passing to the weak limit along its tail would give
$\|h-p\|^2\leq0$, a contradiction.
Consequently choose finite convex combinations $v_m$ as stated
with $\|v_m-h\|_{\mathcal H}<2^{-m}$.

The $v_m$ are actual classical potentials on the whole original
$\Omega$, still in $[0,1]$. For each fixed $l$,
\[
 N_l(v_m-v_r)\longrightarrow0\quad(m,r\to\infty).
\]
Lemma~\ref{lem:r3-linear-energy-capacity} makes the sequence
Cauchy in outer capacity on $V_l$ with capacity domain
$\Omega_l$. For any permitted $H\Subset\Omega'\subset\Omega$,
finitely many such $V_l$ cover $H$, with $\Omega_l\subset\Omega'$.
Finite subadditivity and domain monotonicity therefore give the
Cauchy property on this exact pair $(H,\Omega')$. Thus it holds
on all permitted pairs, not only for one larger capacity domain.

Proposition~\ref{prop:r2-completion} now supplies an actual Borel
limit locally in outer capacity. Its construction takes a
pointwise limsup of a subsequence of the $v_m$, so the resulting
representative has values in $[0,1]$ at every point. It is in
$\mathcal A_S^{[0,1]}(\Omega)$ by the very same normalized
approximants. This step supplies the pointwise representative
which the abstract Hilbert argument alone would not supply.

Suppose finally that $D(u_jS)\to\nu$. Linearity in one
differentiated slot makes $D(v_mS)$ the corresponding finite
convex combination of tail measures. For each fixed original
test its scalar pairings converge to the pairing with $\nu$.
The completed calculus instead identifies their limit as
$D(aS)$, hence $\nu=D(aS)$. Its weighted capacity-continuity
theorem also gives the last limit in
\eqref{eq:r3-linear-weak-representation}. All conclusions use
the fixed $S$ and all original tests.
\end{proof}

\paragraph{Status.}
Theorem~\ref{thm:r3-linear-convex-blocks} is
\texttt{NEW\_PROVED}. It holds in every ambient dimension, but
its remaining degree is specifically one. Convex combinations
are essential: the theorem does not contradict the normalized
family without a capacity-Cauchy subsequence in
Proposition~\ref{prop:r2-capacity-noncompactness}.

\subsection{What the representation does and does not settle}

Define, for Borel $E\subset\Omega$,
\[
 C_{\mathcal A}(E;\Omega)
   =\sup_{a\in\mathcal A_S^{[0,1]}(\Omega)}D(aS)(E).
\]
There is no enlargement of the capacity domain in this definition.

\begin{corollary}[The outer capacity on compact sets is the completed capacity]
\label{cor:r3-linear-completed-capacity}
In the degree-one setting of
Theorem~\ref{thm:r3-linear-convex-blocks},
\begin{align}
 C_\Omega(E)&\leq C_{\mathcal A}(E;\Omega)
                   \leq C_\Omega^*(E)
                  &&(E\text{ Borel}),\label{eq:r3-completed-cap-sandwich}\\
 C_{\mathcal A}(K;\Omega)&=C_\Omega^*(K)
                  &&(K\Subset\Omega\text{ compact}).
                  \label{eq:r3-completed-compact-cap}
\end{align}
Consequently raw compact outer regularity for the original
classical capacity is equivalent here to the no-gain assertion
\begin{equation}
\label{eq:r3-completion-no-gain}
 \sup_{a\in\mathcal A_S^{[0,1]}(\Omega)}D(aS)(K)
       \leq
 \sup_{\substack{u\in\Cb(\Omega)\\0\leq u\leq1}}D(uS)(K)
 \quad(K\Subset\Omega).
\end{equation}
This last assertion is not proved for an arbitrary $S$.
\end{corollary}

\begin{proof}
The first inequality of \eqref{eq:r3-completed-cap-sandwich}
follows from constant approximating sequences. If $a$ is in
the normalized completion, choose its normalized approximants
$a_j$ on $\Omega$. The completed calculus and local mass bounds
give local weak convergence of $D(a_jS)$ to $D(aS)$, including
against compact continuous scalar tests. For every ordinary
open $G\subset\Omega$, the open-set inequality implies
\[
 D(aS)(G)\leq\liminf_jD(a_jS)(G)\leq C_\Omega(G).
\]
Taking open neighborhoods of a Borel $E$ and then taking the
supremum over $a$ proves the second inequality. This argument
does not use the unavailable compact-set inequality for weak
limits.

Fix compact $K\Subset\Omega$. Choose an ordinary neighborhood
$W$ of $K$ and a compactly supported original smooth cutoff
$0\leq\zeta\leq1$ equal to one near $\overline W$. Put
\[
 \mathcal M_\zeta
    =\{\zeta D(uS):u\in\Cb(\Omega),\ 0\leq u\leq1\}.
\]
On a compact metrizable neighborhood of the cutoff support this
is a uniformly mass-bounded family. Its capacity on subsets of
$W$ is exactly $C_\Omega$; open neighborhoods of $K$ may be
restricted to $W$. Proposition~\ref{prop:capacity-weak-closure}
therefore gives
\[
 C_\Omega^*(K)
       =\sup_{\widetilde\nu\in
                    \overline{\mathcal M_\zeta}^{\,w}}
                    \widetilde\nu(K).
\]

Take a sequence $\zeta D(u_jS)$ converging to one such
$\widetilde\nu$. Relative CLN bounds $D(u_jS)$ on every compact
subset of $\Omega$. By a countable compact exhaustion and
diagonal weak compactness, pass to a subsequence for which
$D(u_jS)$ converges locally weakly on all of $\Omega$ to a
positive Radon measure $\nu$. More explicitly, use compact
cutoffs equal to one on successive compact neighborhoods;
the finite weak limits agree on their overlapping interiors
and hence define $\nu$. The already convergent cutoff sequence
then satisfies $\widetilde\nu=\zeta\nu$.

Theorem~\ref{thm:r3-linear-convex-blocks} represents
$\nu=D(aS)$ for some
$a\in\mathcal A_S^{[0,1]}(\Omega)$. Since $\zeta=1$ on $K$,
\[
 \widetilde\nu(K)=D(aS)(K)\leq C_{\mathcal A}(K;\Omega).
\]
Take the supremum over the weak closure and combine with
\eqref{eq:r3-completed-cap-sandwich} to obtain
\eqref{eq:r3-completed-compact-cap}. Finally, the equivalence
with \eqref{eq:r3-completion-no-gain} is immediate from that
equality and the inclusion of classical potentials.
\end{proof}

\paragraph{Scope and minimal remaining gap.}
The corollary is a \texttt{DIRECT\_COROLLARY} of the new convex-block
theorem and the proved measure-family reduction. It replaces a
weak-limit question by the exact completed-class inequality
\eqref{eq:r3-completion-no-gain}. It does not replace that
inequality by a constant-factor estimate or a smaller-domain
capacity. For an arbitrary $S$, an actual completed representative
has not been shown to have a classical representative on the same
domain. Ambient upper regularization can change values on a set
charged by $S$, as the support examples of
Section~\ref{sec:r2-support-envelopes} show. The ambient
$S=[1]$ representative theorem is therefore not an input here.

\subsection{The non-pluripolar-background case closes in degree one}

The next two assertions allow every remaining degree $q=n-p\geq1$.
We say that a fixed positive closed background $S$ is
\emph{locally non-pluripolar} if the total variations of all its
coefficient-Radon measures give zero mass to every locally
pluripolar Borel set. Equivalently,
on each ordinary connected witness domain one has
$\mathbf1_P S=0$ whenever
$P=\{\phi=-\infty\}$ and $\phi$ is extended classical psh and
not identically $-\infty$. The equivalence follows by the
definition of local pluripolarity, countable local covers in
the present finite $p$-adic setting, and coefficient-Radon
multiplication. This is an explicit extra hypothesis, not a
property of all positive closed currents.

\begin{lemma}[Ambient polar sets are relatively outer-null under nonconcentration]
\label{lem:r3-nonpolar-relative-null}
Let $S$ be locally non-pluripolar. Every locally pluripolar set
has locally zero outer $\operatorname{Cap}_S$-capacity, on every
permitted compact/domain pair.
\end{lemma}

\begin{proof}
First work where $\phi\leq0$ is a nontrivial classical witness
and put $P=\{\phi=-\infty\}$. By the hypothesis and
Theorem~\ref{thm:r2-polar-restriction},
\[
 \mathbf1_P T_{S,q}(v,\ldots,v)
     =T_{\mathbf1_PS,q}(v,\ldots,v)=0
\]
for every normalized classical tester on this same domain.
Thus the raw capacity of $P$ is zero there.

It remains essential to obtain an outer-zero assertion.
The bounded increasing classical sequence
\[
 w_N=1+\frac{\max(\phi,-N)}{N}
\]
has actual limit $w=\mathbf1_{P^c}$.
Corollary~\ref{cor:r2-local-increasing} and the completion
theorem put $w$ in $\mathcal A_{S,b}$ and make it locally
relatively quasicontinuous. On an ordinary relatively compact
core $V$,
\[
 P\cap V=V\cap\{w<1/2\}
\]
is quasi-open. Proposition~\ref{prop:r2-quasiopen-outer}
therefore turns its raw zero capacity into outer zero
capacity. This argument does not assume compact outer
regularity.

Finally, fix an arbitrary permitted pair $H\Subset\Omega$.
Cover the part of the locally pluripolar set in $H$ by
countably many ordinary cores with witness domains contained
in $\Omega$. On each core the preceding argument applies.
Domain monotonicity and countable subadditivity of the outer
capacity give zero capacity on the original pair.
Subsets need not be Borel, since outer capacity is monotone.
\end{proof}

\begin{proposition}[Classical representatives for a non-pluripolar background]
\label{prop:r3-nonpolar-classical-representative}
Let $S$ be locally non-pluripolar on the fixed norm-box domain
$\Omega$, in any remaining degree $q\geq1$. For every
$a\in\mathcal A_S^{[0,1]}(\Omega)$ there is
$r\in\Cb(\Omega)$, $0\leq r\leq1$ on $\Omega$, such that
$a=r$ locally outside sets of outer $S$-capacity zero.
The representative is not asserted to be unique.

For any finite tuple of such completed potentials, any choices
of these representatives give the same weighted and unweighted
BT currents on all original compact smooth complementary tests.
In particular the normalized completed class does not enlarge
the capacity of any Borel set in this case.
\end{proposition}

\begin{proof}
Choose normalized defining approximants $a_j$ on the entire
$\Omega$. The summable-subsequence construction in
Proposition~\ref{prop:r2-completion}, now applied to their known
capacity limit $a$, gives a subsequence $a_{j_l}$ converging
pointwise to $a$ outside a locally relative-outer-null set $N$.
For clarity, on the first $l$ members of a countable cofinal
list of compact/domain pairs require the capacity of
$\{|a_{j_l}-a|>2^{-l}\}$ to be at most $2^{-l}$.
Open covers and the countable subadditivity bound on their
limsup sets give the asserted $N$ on every pair.

Set
\[
 h_m=\sup_{l\geq m}a_{j_l},\qquad r_m=h_m^*.
\]
The bounded upper-envelope theorem,
Theorem~\ref{thm:r2-local-upper-envelope}, gives
$r_m\in\Cb(\Omega)$ with $0\leq r_m\leq1$.
The sequence decreases, so
$r=\lim_m r_m$ is classical on the same domain with the same
normalization. For each $m$, the discrepancy set
\[
 E_m=\{h_m<r_m\}
\]
is locally pluripolar by
Theorem~\ref{thm:r2-negligible-outer-zero} and
Corollary~\ref{cor:r2-polar-negligible-equivalence}.
Lemma~\ref{lem:r3-nonpolar-relative-null} makes $E_m$ locally
relative-outer-null. Outside $N\cup\bigcup_mE_m$ one has
\[
 r=\lim_m h_m=\limsup_l a_{j_l}=a.
\]
This proves the asserted equality modulo local relative outer
capacity zero, without replacing any of the given Borel point
values in the original completed class.

The completed calculus is independent of precisely these
changes of representatives, including the weight slot and
the degree-zero background multiplication, by
Theorem~\ref{thm:r2-completed-calculus}. It therefore identifies
each completed product with the product of the chosen
classical representatives. In particular, if $a$ is normalized,
its measure $T_{S,q}(a,\ldots,a)$ is the measure of a normalized
classical $r$ on the original $\Omega$. Taking suprema over
Borel sets gives equality of the two capacity families.
Uniqueness is intentionally absent: for $S=0$ all functions
are relatively quasi-everywhere equal.
\end{proof}

\noindent\textbf{Status.}
The preceding lemma and representative proposition are
\texttt{DIRECT\_COROLLARY} consequences of the proved R2
polar-restriction, envelope, negligible-set, and completion
theorems. They retain the stated nonconcentration hypothesis.

\begin{theorem}[Sharp compact regularity for a non-pluripolar degree-one background]
\label{thm:r3-nonpolar-linear-compact-regularity}
Let $S$ be an arbitrary fixed locally non-pluripolar,
symmetric original-positive separately closed current of
bidegree $(n-1,n-1)$ on the permitted fixed norm-box domain
$\Omega$. Every locally weak limit of
\[
 D(u_jS),\qquad u_j\in\Cb(\Omega),\quad0\leq u_j\leq1,
\]
is $D(rS)$ for some $r\in\Cb(\Omega)$ with $0\leq r\leq1$.
Consequently, for every compact $K\Subset\Omega$,
\[
 \operatorname{Cap}_S^*(K;\Omega)
       =\operatorname{Cap}_S(K;\Omega).
\]
The raw capacity is continuous on decreasing compact subsets
of $\Omega$. Equivalently, every localized weak-closure
measure satisfies the sharp bound
\eqref{eq:r3-sharp-relative-closure} on the core where its
cutoff is one.
\end{theorem}

\begin{proof}
Theorem~\ref{thm:r3-linear-convex-blocks} represents any such
weak limit by a normalized completed potential on $\Omega$.
Proposition~\ref{prop:r3-nonpolar-classical-representative}
replaces its class by a normalized classical potential on
this same domain, preserving the current.
The compact assertion follows either from
Corollary~\ref{cor:r3-linear-completed-capacity} or directly
from Proposition~\ref{prop:capacity-weak-closure}.
That proposition also gives continuity from above on compact
sets. Its localized version is exactly the final assertion.
\end{proof}

\noindent\textbf{Status: \texttt{NEW\_PROVED}.}
The conclusion holds in all ambient dimensions under the
explicit non-pluripolarity hypothesis and with remaining
degree one. The representative proposition alone does not
give weak-closure representation for $q>1$.

\subsection{Why an ambient contact measure is not a relative one}

We first give a compact-obstacle example in remaining degree one.
It applies to all classical competitors, not merely toric ones.

\begin{proposition}[An ambient compact extremal with relative mass off the obstacle]
\label{prop:r3-compact-relative-contact-false}
Let $n\geq2$, choose a uniformizer $\pi$ of the original $K$,
put $\lambda=-\log|\pi|>0$, and work on the analytic torus
$X=(\mathbf G_{m,K}^n)^{\mathrm{an}}$. Set
\[
 t_i=\log|T_i|,\qquad
 \theta_i=d't_i\wedge d''t_i,\qquad
 \beta=\sum_{i=1}^n\theta_i,
 \qquad S=[\beta^{n-1}].
\]
For $0<a<\lambda$ define the ordinary domain and its compact core
\[
 \Omega=\{\max_i|t_i|<\lambda\},\qquad
 K_a=\{\max_i|t_i|\leq a\},
\]
and the continuous classical psh function
\[
 H=\max\left\{-1,
          \frac{\max_i|t_i|-\lambda}{\lambda-a}\right\}.
\]
Then $H$ is the actual pointwise supremum of all
$v\in\Cb(\Omega)$ satisfying $v\leq0$ on $\Omega$ and
$v\leq-1$ on $K_a$. Nevertheless
\[
 D(HS)(\Omega\setminus K_a)>0.
\]
Thus the ambient compact extremal does not in general have a
relative Monge--Amp\`ere measure supported on its obstacle.
\end{proposition}

\begin{proof}
The domain is an ordinary open with compact closure in the
smooth affine torus, and $K_a$ is compact in it. If a strict
norm-box core is desired, take, for example, $a=\lambda/2$
and express its bounds using $T_i^2$ and the original value
$|\pi|$; no valuation is rescaled. Logarithms of analytic units
are harmonic after every complete extension and curve pullback.
Finite maxima of their affine combinations and constants are
classical psh. This proves that $H\in\Cb(\Omega)$, that
$-1\leq H<0$, that $H=-1$ on $K_a$, and that $H$ has ordinary
boundary value zero.

We claim $(DH)^n=0$ off $K_a$. Near any point there the constant
branch $-1$ is inactive. Among the branches $t_i$ and $-t_i$
of the maximum, at most one sign for each $i$ is active, since
the maximum is strictly positive. Delete inactive branches
on a neighborhood. If the remaining affine functions are
$\ell_1,\ldots,\ell_m$, then $m\leq n$, and
\[
 \max(\ell_1,\ldots,\ell_m)
    =\ell_m+\max(\ell_1-\ell_m,\ldots,
                         \ell_{m-1}-\ell_m,0).
\]
The differences are logarithms of analytic units. Smooth the
finite maximum by translation-invariant convex regularized
maxima. Their curvature factors through at most $m-1$ such
moment coordinates, so its $n$th exterior power vanishes
identically as an original form. The uniform continuous
fixed-background calculus passes this equality to $H$.
This proves the claim without a regularity assertion for
arbitrary psh functions on the skeleton.

Let $v$ be an admissible competitor. For every $\epsilon>0$,
the set $\{H<v-\epsilon\}$ is contained in
$\{H<-\epsilon\}$, whose closure is compact in $\Omega$.
Also $\{H<v\}$ is disjoint from $K_a$, because $H=-1$ and
$v\leq-1$ there. Thus $(DH)^n(\{H<v\})=0$.
Corollary~\ref{cor:r2-pointwise-ambient}, on this relatively
compact affine domain, gives $v\leq H$ at every point.
Since $H$ itself is admissible, it is the stated actual
supremum. This comparison includes all non-toric competitors.

To compute relative mass outside $K_a$, choose a tropical open
rectangle meeting
\[
 t_1=t_2=r,\qquad a<r<\lambda,
 \qquad |t_i|<r\ (i\geq3),
\]
small enough that the only potentially maximal branches are
$t_1,t_2$, and that their maximum is greater than $a$.
There
\[
 H=\frac{t_2-\lambda+\max(0,t_1-t_2)}{\lambda-a}.
\]
The one-variable distributional second derivative of
$\max(0,s)$ is $\delta_0$. One may equivalently obtain the
following formula by smooth convex regularization and the
original CLD tropical integration formula. For a nonnegative
original smooth tester $\psi\circ\operatorname{trop}$ supported
in that rectangle,
\[
 \int(\psi\circ\operatorname{trop})\,D(HS)
   =\frac{2(n-1)!}{\lambda-a}
     \int\psi(t_2,t_2,t_3,\ldots,t_n)
                    \,dt_2\,dt_3\cdots dt_n.
\]
Indeed, the square of the gradient of $s=t_1-t_2$ in the
diagonal form $\beta$ is two, and
$d's\wedge d''s\wedge\beta^{n-1}
 =2(n-1)!\theta_1\wedge\cdots\wedge\theta_n$.
Choose $\psi$ positive at a point of the indicated diagonal.
The displayed integral is strictly positive and is supported
outside $K_a$, proving the assertion.
\end{proof}

There is also a smooth common-collar obstruction in every
intermediate remaining degree. It directly tests the relative
energy envelope rather than only the ambient compact extremal.

\begin{proposition}[A relative energy maximum need not satisfy contact orthogonality]
\label{prop:r3-relative-energy-contact-false}
On the same torus and domain, let $1\leq q<n$ and put
\[
 S=[\beta^{n-q}],\qquad b=\sum_{i=1}^q t_i^2.
\]
Choose nonzero $\varphi\in C_c^\infty((-\lambda,\lambda)^n)$,
$\varphi\geq0$, and let $\chi=\varphi\circ\operatorname{trop}$,
$L=\operatorname{supp}\chi$, and $f=b+\chi$. Consider
\[
 \mathcal F=\{v\in\Cb(\Omega):
       b-1\leq v\leq f,\quad v=b\text{ on }\Omega\setminus L\}.
\]
The greatest actual member and the relative capacitary energy
envelope of this uniformly bounded common-collar family are $b$
(the latter modulo local relative capacity zero). However
\begin{equation}
\label{eq:r3-failed-relative-orthogonality}
 \int_\Omega(f-b)T_{S,q}(b,\ldots,b)
    =2^q q!(n-q)!\int_{\mathbf R^n}\varphi(t)\,dt>0.
\end{equation}
In particular, relative energy maximization over ambient
classical psh competitors does not imply the expected contact
orthogonality, even for a smooth strongly positive background.
\end{proposition}

\begin{proof}
The smooth function $b$ is classical psh: it is a convex smooth
function of harmonic logarithms of units, after every complete
extension and curve pullback. It is bounded on $\Omega$ and
\[
 Db=2\sum_{i=1}^q\theta_i,
 \qquad (Db)^n=0,
 \qquad T_{S,q}(b,\ldots,b)
       =2^q q!(n-q)!\theta_1\wedge\cdots\wedge\theta_n.
\]
For any $v\in\mathcal F$, the closures of all positive-margin
sets $\{b<v-\epsilon\}$ lie in the fixed compact $L$.
Pointwise ambient domination again gives $v\leq b$ on all of
$\Omega$. Since $b\in\mathcal F$, it is the greatest member.
Relative Aubin--Mabuchi energy is monotone on this common-collar
class, so $\mathcal E_b(v)\leq\mathcal E_b(b)=0$. The constant
sequence $b$ is a maximizing sequence in
Theorem~\ref{thm:r2-one-domain-envelope}, identifying its
capacitary envelope with $b$.

The original CLD integration formula for the identity tropical
chart of the torus has multiplicity one and integrates
$\theta_1\wedge\cdots\wedge\theta_n$ as Lebesgue measure
with its positive orientation. Substitution gives
\eqref{eq:r3-failed-relative-orthogonality}. Its strict positivity
uses only the chosen nonzero nonnegative smooth $\varphi$.
All products in this example are ordinary original currents.
\end{proof}

\paragraph{Status of the counterexamples.}
The two counterexamples are \texttt{NEW\_PROVED}; the contact or
orthogonality implications they disprove have status
\texttt{FALSE}. Neither is a counterexample to
\eqref{eq:r3-sharp-relative-closure}. The second concerns a
continuous obstacle and a common-collar family; it does not
identify every possible relative compact-capacity extremal.
The smooth backgrounds in these examples are locally
non-pluripolar, by the ambient smooth-trace polar-zero theorem.
Thus the degree-one contact failure coexists with the compact
outer regularity proved above; it is a failure of the proposed
proof mechanism, rather than a failure of that regularity.

\subsection{The remaining nonlinear weak-limit obstruction}

For $q>1$, applying the same square integration-by-parts argument
to
\[
 \int\chi f^2(Dw)^q\wedge S
\]
uses the energy of $f$ relative to
\[
 Q_w=(Dw)^{q-1}\wedge S.
\]
This is a positive closed background of bidegree $(n-1,n-1)$,
but it changes with the capacity tester $w$. The direct Hilbert
argument above therefore does not by itself give a fixed
seminorm controlling every tester: convergence in countably
many selected seminorms would not justify that supremum.
The independent mixed-energy recurrence in
Theorem~\ref{thm:r3-I-capacity} resolves this uniform-control
problem. Together with quantified concavity, it gives the
all-degree convex-block theorem
\ref{thm:r3-all-degree-convex-blocks}. Thus tester-dependent
energy is no longer an open obstruction in this record.

The separate nonlinear weak-limit problem remains. For $q>1$,
the MA measure of a convex combination contains mixed terms;
it is not the corresponding convex combination of the original
MA measures. The new capacity-convergent convex blocks do not
therefore identify an arbitrary weak limit of those original
measures with the MA measure of their completed limit.

The current status of arbitrary-background compact outer
regularity remains \texttt{OPEN}. The precise reusable advances
are the fixed-energy capacity estimate, normalized convex-block
representation in degree one, its exact completed-capacity
reduction, and sharp compact regularity under the additional
non-pluripolarity hypothesis in that degree. The stronger
convex-block compactness result for every remaining degree is
proved separately in Section~\ref{sec:r3-all-degree-energy}.
The contact
counterexamples rule out the ambient contact and relative-energy
orthogonality routes above.



% END INCLUDED SOURCE: sections/relative_compact_capacity


% BEGIN INCLUDED SOURCE: sections/relative_contact_inequality
% Reconstructed after the scratch reset. Original reviewed SHA256:
% e004c695696c6709a129f0d7f3bd93cffe7a96ca6f9b70813743372f21eeec43.
\section{A contact inequality for every fixed positive background}
\label{sec:r3-relative-contact-inequality}

Let $S$ be any allowed symmetric original-positive separately
closed $(p,p)$-current on an ordinary open set $U$, and put
$q=n-p\geq1$. Write $\mu_u=T_{S,q}(u,\ldots,u)$ for actual
$u\in\mathcal C_b(U)$. Restrictions below are Borel restrictions
of the coefficient-Radon top measures.

\begin{proposition}[Relative maximum and contact inequalities]
\label{prop:r3-relative-contact-inequality}
For $u,v\in\mathcal C_b(U)$,
\begin{equation}
 \mu_{\max(u,v)}
       \geq{\bf1}_{\{u\geq v\}}\mu_u
                  +{\bf1}_{\{u<v\}}\mu_v.
 \label{eq:r3-relative-max-inequality}
\end{equation}
In particular, if $u\leq h$,
\begin{equation}
 {\bf1}_{\{u=h\}}\mu_u
             \leq{\bf1}_{\{u=h\}}\mu_h.
 \label{eq:r3-relative-contact-inequality}
\end{equation}
\end{proposition}

\noindent\textbf{Status: \texttt{DIRECT\_COROLLARY}} of the
proved Borel maximum locality and decreasing BT continuity.

\begin{proof}
Put $w_j=\max(u+1/j,v)$. This sequence decreases to
$w=\max(u,v)$. On the Borel set $E=\{u\geq v\}$ the strict
inequality $u+1/j>v$ holds, so Borel maximum locality in each
differentiated slot and invariance under constants give
\begin{equation}
 {\bf1}_E\mu_{w_j}={\bf1}_E\mu_u.
 \label{eq:r3-contact-before-limit}
\end{equation}
The top measures converge weakly against compact continuous
functions, by the original-current decreasing theorem and the
coefficient-Radon convergence result. For a compact $K\subset E$,
local Portmanteau therefore gives
\[
 \mu_w(K)\geq\limsup_j\mu_{w_j}(K)=\mu_u(K).
\]
For an arbitrary Borel subset of $E$, inner regularity of the
fixed Radon measure $\mu_u$ extends this inequality from its
compact subsets. Work first in a relatively compact stage and
then exhaust if necessary. This step does not require $E$ to
be closed or ordinary open. On the complementary Borel set
$\{u<v\}$, strict Borel maximum locality gives
$\mu_w=\mu_v$. The two disjoint restrictions prove
\eqref{eq:r3-relative-max-inequality}.
If $u\leq h$, take $v=h$, so that $w=h$ and
$\{u\geq h\}=\{u=h\}$, giving
\eqref{eq:r3-relative-contact-inequality}.
\end{proof}

\begin{corollary}[Zero mass on contact with a maximal potential]
\label{cor:r3-relative-zero-contact-mass}
If $u\leq h$ are actual bounded classical potentials and
$T_{S,q}(h,\ldots,h)=0$, then
$ {\bf1}_{\{u=h\}}T_{S,q}(u,\ldots,u)=0$.
\end{corollary}

The sets in this statement need only be Borel. No arbitrary
varying-background theorem, capacity regularity assertion, or
exchange of discontinuous limits is used.

% END INCLUDED SOURCE: sections/relative_contact_inequality


% BEGIN INCLUDED SOURCE: sections/nonzero_dirichlet
\section{Nonzero atomic Dirichlet problems on one fixed domain}
\label{sec:r3-nonzero-dirichlet}

Fix a finite extension $K/\mathbf Q_p$ with its original absolute
value. Let $Y/K$ be smooth integral affine of pure dimension $n\geq1$.
Throughout this section the domain is fixed in advance:
\[
 V\Subset Y^{\an}
\]
is connected and ordinary open, has nonempty ordinary boundary, and
admits
\[
 \rho\in\Cb(V)\cap C(\overline V),\qquad
 \rho<0\text{ on }V,\qquad \rho=0\text{ on }\partial V.
 \tag{H}\label{eq:r3-dirichlet-domain}
\]
Closures and boundary limits are ordinary ones in $Y^{\an}$.
The points and masses below are chosen after $V$ has been fixed.
We use $D=d'd''$ and the ambient background $[1]$, with the original
CLD calibration. All weak equalities and limits are tested against
the full original compactly supported smooth scalar class.

The inputs from the preceding records are the bounded mixed calculus,
ordinary upper-envelope closure, divisorial no-jump, local
maximality equivalence $(Du)^n=0$, the boundary maximum principle,
and pointwise ambient comparison, domination and uniqueness.
These inputs retain the all-extensions classical definition. The
projective input B, Theorems C.6 and 4.4, has already been used in
proving those local results; no new projective or residue-characteristic
hypothesis is inserted here.

\subsection{Point obstacles and their actual envelopes}

Fix pairwise distinct ambient divisorial points
$Z=\{z_1,\ldots,z_r\}\subset V$, where $r\geq1$. For
$a=(a_1,\ldots,a_r)\in[0,\infty)^r$ put $A(a)=\max_i a_i$ and
\[
 \mathcal F_a=\{v\in\Cb(V):-A(a)\leq v\leq0,
                         \ v(z_i)\leq-a_i\text{ for every }i\},
 \qquad U_a=\left(\sup_{v\in\mathcal F_a}v\right)^*.
\]
The constant $-A(a)$ makes this family nonempty.

\begin{proposition}[Actual weighted point envelopes]
\label{prop:r3-weighted-point-envelope}
For every $a$ the function $U_a$ belongs to $\mathcal F_a$ and
\[
 U_a=\sup_{v\in\mathcal F_a}v
\]
at every point. It is also the greatest member of $\Cb(V)$ with
$v\leq0$ and $v(z_i)\leq-a_i$, even without imposing the lower
bound in the definition of $\mathcal F_a$. It has ordinary
boundary limit zero.

For all $a,b\in[0,\infty)^r$ and $t\geq0$,
\[
 \|U_a-U_b\|_{L^\infty(V)}\leq\|a-b\|_\infty,
 \qquad U_{ta}=tU_a.
 \tag{1}\label{eq:r3-envelope-lipschitz}
\]
For every fixed $M<\infty$, all parameters $a\in[0,M]^r$ have
one common continuous lower barrier with zero boundary limit.
\end{proposition}

\begin{proof}
Upper-envelope closure gives $U_a\in\Cb(V)$ and
$-A(a)\leq U_a\leq0$. Divisorial no-jump gives the crucial
actual constraint
\[
 U_a(z_i)=\sup_{v\in\mathcal F_a}v(z_i)\leq-a_i.
\]
Thus $U_a$ itself is in $\mathcal F_a$. The pointwise supremum
is consequently at least $U_a$, and is at most $U_a$ by the
definition of upper regularization. This proves the equality
and the greatest-member assertion for the displayed family.
If $v\in\Cb(V)$ is any other nonpositive admissible function,
then $\max(v,-A(a))\in\mathcal F_a$ and dominates $v$.
This proves the assertion without the artificial lower bound.

Set
\[
 d_Z=\min_i(-\rho(z_i))>0.
\]
For $M>0$ the continuous classical function
\[
 b_M=\max\left(-M,\frac{M}{d_Z}\rho\right)
\]
equals $-M$ at every $z_i$ and has boundary limit zero.
For $a\in[0,M]^r$, the greatest-member characterization gives
$b_M\leq U_a\leq0$. This proves the common barrier and the
boundary limit. When $a=0$, the family consists of zero and
$U_0=0$.

Let $\delta=\|a-b\|_\infty$. The function
\[
 w=\max(-A(b),U_a-\delta)
\]
lies between $-A(b)$ and zero. At each $z_i$, both branches
are at most $-b_i$, because $b_i\leq a_i+\delta$.
Hence $w\in\mathcal F_b$ and $U_b\geq U_a-\delta$.
Interchanging the parameters proves the first inequality in
\eqref{eq:r3-envelope-lipschitz}. For $t>0$, multiplication by
$t$ is a bijection between the two admissible families, and
commutes with ordinary upper regularization. The case $t=0$
has already been noted.
\end{proof}

\begin{proposition}[A continuous homogeneous map of obstacle depths to masses]
\label{prop:r3-atomic-mass-map}
There are nonnegative real functions $c_i(a)$ such that
\[
 (DU_a)^n=\sum_{i=1}^r c_i(a)\delta_{z_i}.
 \tag{2}\label{eq:r3-atomic-mass-map}
\]
Each $c_i$ is locally Lipschitz on $[0,\infty)^r$, satisfies
$c_i(ta)=t^n c_i(a)$, and obeys
\[
 a_i=0\ \Longrightarrow\ c_i(a)=0.
 \tag{3}\label{eq:r3-mass-face}
\]
For every $a\ne0$ the total mass
$m(a)=\sum_i c_i(a)$ is strictly positive.
\end{proposition}

\begin{proof}
Write $Z_+(a)=\{z_i:a_i>0\}$. We first prove that $U_a$
is locally maximal off $Z_+(a)$. Take an ordinary
$W\Subset V\setminus Z_+(a)$ and a competitor
$v\in\Cb(W)$ which equals $U_a$ outside a compact
$L\Subset W$. Glue $\max(v,U_a)$ on $W$ to $U_a$ off
$L$. The resulting function $w\in\Cb(V)$ is bounded above:
only the compact set $L$ has changed. It has boundary limit
zero, so the boundary maximum principle gives $w\leq0$.
Moreover $w\geq U_a\geq-A(a)$. At each positive-depth
constraint it equals $U_a$; each zero-depth constraint is
automatic from $w\leq0$. Thus $w\in\mathcal F_a$.
The actual greatest-member property forces $w\leq U_a$,
and hence $v\leq U_a$. This is local maximality.

The proved homogeneous-equation criterion now gives
$(DU_a)^n=0$ on $V\setminus Z_+(a)$. This positive Radon
measure is supported on a finite compact set, so it is finite
and has the form \eqref{eq:r3-atomic-mass-map}. Its support
also proves \eqref{eq:r3-mass-face}. Homogeneity follows from
\eqref{eq:r3-envelope-lipschitz} and the proved BT homogeneity.

Here is a direct continuity estimate, with no varying-background
assertion. Choose original smooth compact scalar cutoffs
$\chi_i$ which equal one near $z_i$ and zero near every other
point of $Z$. Such cutoffs exist in the original CLD smooth
sheaf. For $a,b\in[0,M]^r$, the bounded mixed calculus gives
\[
 (DU_a)^n-(DU_b)^n
   =D\left((U_a-U_b)
          \sum_{j=0}^{n-1}(DU_a)^j\wedge(DU_b)^{n-1-j}\right).
\]
Pair with $\chi_i$. The coefficient-Radon estimates for the
fixed smooth test $D\chi_i$ and relative CLN, on a compact
neighborhood of its support, give a constant $C_{i,M}$ with
\[
 \begin{split}
 |c_i(a)-c_i(b)|
 &\leq C_{i,M}\|U_a-U_b\|_{L^\infty(V)}\\
 &\leq C_{i,M}\|a-b\|_\infty.
 \end{split}
\]
The constant depends on the fixed test stage, $M,n$ and the
domain geometry, not on $a,b$. Borel multiplication by the
bounded difference is legitimate for these coefficient-Radon
currents. When $n=1$, the sum contains the single background
$[1]$ and the same argument applies. Alternatively, uniform
convergence in \eqref{eq:r3-envelope-lipschitz}, the proved
BT stability, and the fixed finite support give continuity
of every coefficient directly. No mass can escape to the
boundary because the support is always contained in $Z$.

Finally suppose $a\ne0$ and $m(a)=0$. Then $(DU_a)^n=0$.
The common lower barrier makes every set
$\{U_a<-\delta\}$ relatively compact in $V$. Pointwise
ambient domination applied to $(U_a,0)$ gives $U_a\geq0$.
This contradicts $U_a(z_i)\leq-a_i<0$ for an index with
$a_i>0$. Therefore $m(a)>0$.
\end{proof}

\subsection{A finite-dimensional surjectivity argument}

\begin{lemma}[Surjectivity of a face-preserving simplex map]
\label{lem:r3-simplex-surjectivity}
Let
\[
 \Delta^{r-1}=\{a\in[0,\infty)^r:\textstyle\sum_i a_i=1\}.
\]
If a continuous $F:\Delta^{r-1}\to\Delta^{r-1}$ satisfies
$a_i=0\Rightarrow F_i(a)=0$, then $F$ is surjective.
\end{lemma}

\begin{proof}
First let $y$ be in the relative interior of the simplex.
Put
\[
 d_i(a)=(y_i-F_i(a))_+,
 \qquad \sigma(a)=\sum_i d_i(a),
 \qquad
 T_i(a)=\frac{a_i+d_i(a)}{1+\sigma(a)}.
\]
This is a continuous self-map of the simplex. Brouwer's
fixed-point theorem, used here as an external topological
input \cite{Brouwer11}, gives a fixed point $a$. Its equations are
\[
 d_i(a)=\sigma(a)a_i.
\]
If $a_i=0$, the face condition gives $F_i(a)=0$, whereas
$y_i>0$ gives $d_i(a)>0$, a contradiction. Thus all
$a_i>0$. If $\sigma(a)>0$, every $d_i(a)>0$, so
$y_i>F_i(a)$ for every $i$, contradicting that both
vectors have sum one. Hence $\sigma(a)=0$. Then
$y_i\leq F_i(a)$ for every $i$, and equality of their sums
gives $F(a)=y$.

For a boundary target, restrict to its support face. The
face condition makes that face invariant, and the preceding
argument applies to its relative interior. A one-point
simplex is immediate. This proves surjectivity.
\end{proof}

For $r=2$ only the intermediate value theorem is needed:
the normalized first mass along $a=(t,1-t)$ is continuous,
equals zero at $t=0$, and equals one at $t=1$.

\begin{theorem}[Finite atomic Dirichlet problem on the fixed domain]
\label{thm:r3-finite-atomic-dirichlet}
Assume \eqref{eq:r3-dirichlet-domain}. Let $z_1,\ldots,z_r$
be distinct ambient divisorial points of $V$, let $m_i\geq0$,
and let $b\in\mathbf R$. There exists a unique bounded
classical psh function $u\in\Cb(V)$ such that
\[
 (Du)^n=\sum_{i=1}^r m_i\delta_{z_i},\qquad
 \lim_{\substack{x\to\xi\\x\in V}}u(x)=b
                         \quad(\xi\in\partial V).
 \tag{4}\label{eq:r3-atomic-dirichlet}
\]
The equality of measures has the original calibration.
For a fixed finite set $Z$, there is a finite constant
$C_Z$, depending on $V$ and $Z$ but not on the masses or
$b$, such that
\[
 \|u-b\|_{L^\infty(V)}
       \leq C_Z\left(\sum_i m_i\right)^{1/n}.
 \tag{5}\label{eq:r3-finite-atomic-bound}
\]
With $V,Z$ fixed, the solution depends continuously in the
uniform norm on the nonnegative mass vector and on $b$.
No continuity of $u$ at arbitrary interior points is asserted
by this theorem alone.
\end{theorem}

\begin{proof}
It is enough first to take $b=0$. If all masses are zero,
take $u=0$. Otherwise let $M_\mu=\sum_i m_i>0$ and
$y_i=m_i/M_\mu$. On $\Delta^{r-1}$ define
\[
 F_i(a)=\frac{c_i(a)}{m(a)}.
\]
The denominator is positive; the map is continuous and
preserves faces by Proposition~\ref{prop:r3-atomic-mass-map}.
Choose $a\in\Delta^{r-1}$ with $F(a)=y$, using the
simplex lemma. Put
\[
 t=\left(\frac{M_\mu}{m(a)}\right)^{1/n},\qquad u=tU_a.
\]
Homogeneity gives every prescribed coefficient $m_i$, and
the boundary limit is zero. Adding $b$ proves existence in
\eqref{eq:r3-atomic-dirichlet}. The same previously fixed
domain has been used throughout; no domains adapted to
individual atoms have been glued together.

For uniqueness, two solutions have equal ordinary boundary
limits. For every positive margin, both strict-gap closures
are therefore compact in $V$: cover the compact ordinary
boundary by finitely many neighborhoods on which the
difference has smaller absolute value. Pointwise ambient
uniqueness gives equality at every point.

Since $m(a)$ is continuous and positive on the compact
simplex, $m_Z=\min_{\Delta^{r-1}}m(a)>0$. Also
$\|U_a\|_\infty\leq A(a)\leq1$ there. Thus
$C_Z=m_Z^{-1/n}$ proves \eqref{eq:r3-finite-atomic-bound}.

To prove stability, let the nonnegative mass vectors tend
to a vector $\mu$. If their total masses tend to zero,
\eqref{eq:r3-finite-atomic-bound} gives uniform convergence
to zero. Otherwise represent their zero-boundary solutions
as $t_jU_{a_j}$ with $a_j\in\Delta^{r-1}$ as above.
The parameters $t_j$ are bounded. From any subsequence one
can extract a further subsequence with $a_j\to a$ and
$t_j\to t$. The envelope Lipschitz estimate gives uniform
convergence to $tU_a$. The atomic mass map and the common
boundary barrier show that this limit solves the limiting
problem. Uniqueness identifies it with the specified
solution. All subsequential limits therefore agree, which
proves uniform convergence of the entire sequence.
Translations handle the varying boundary constant.
\end{proof}

The analytic construction and the simplex argument are recorded
as \texttt{NEW\_PROVED}, relative to the preceding R1/R2
theorems. Brouwer's theorem is \texttt{EXTERNAL\_KNOWN}.
The construction does not sum single-atom solutions. Its
coefficients $c_i(a)$ come from the full nonlinear BT measure,
so every mixed term is included before prescribing the masses.

\subsection{Order and a summable countable atomic class}

\begin{lemma}[Monge--Amp\`ere order reverses zero-boundary solution order]
\label{lem:r3-measure-order}
Let $u_1,u_2\in\Cb(V)$ be bounded with boundary limits zero.
If $(Du_1)^n\geq(Du_2)^n$ as Radon measures, then
$u_1\leq u_2$ at every point. The measures need not have
finite total mass on $V$.
\end{lemma}

\begin{proof}
Write $\mu_i=(Du_i)^n$. For $t,\delta>0$ put
\[
 w=(1+t)u_1-\delta,\qquad E=\{u_2<w\}.
\]
The boundary limit of $u_2-w$ is $\delta$, so $E$ itself
has compact closure in $V$. The compact-gap comparison
principle gives
\[
 (1+t)^n\mu_1(E)\leq\mu_2(E)\leq\mu_1(E).
\]
The masses being compared are finite on this compact set.
Hence $\mu_1(E)=\mu_2(E)=0$. Pointwise ambient domination
for $(u_2,w)$ gives $u_2\geq w$ everywhere. Let
$\delta\downarrow0$ and then $t\downarrow0$. The extra
constant shift is what permits the cancellation without
assuming a finite total mass on the whole domain.
\end{proof}

Write $G(\mu)$ for a zero-boundary solution when one has
been constructed. Positive homogeneity and uniqueness give
$G(s\mu)=s^{1/n}G(\mu)$ for $s\geq0$. If the relevant
solutions exist, positivity of all mixed BT terms gives
\[
 G(\mu)+G(\nu)
       \leq G(\mu+\nu)
       \leq\min\{G(\mu),G(\nu)\}.
 \tag{6}\label{eq:r3-sum-sandwich}
\]
Indeed the full mixed expansion makes the Monge--Amp\`ere
measure of the sum at least $\mu+\nu$, so the lemma gives
the lower inequality. The two measure-order comparisons
give the upper inequalities. This explains precisely why
adding single-atom solutions is a useful subsolution but
does not compute the solution in dimension greater than one.

\begin{corollary}[A countable atomic Dirichlet class]
\label{cor:r3-summable-atomic-dirichlet}
Let $z_i\in V$ be pairwise distinct ambient divisorial points
and $m_i>0$, for $i\geq1$. Let
$g_i=G(\delta_{z_i})$ be their single-atom zero-boundary
solutions on this same fixed domain. Suppose
\[
 \sum_{i=1}^{\infty}m_i^{1/n}\|g_i\|_{L^\infty(V)}<\infty.
 \tag{7}\label{eq:r3-atomic-summability}
\]
Then $\mu=\sum_i m_i\delta_{z_i}$ is a locally finite
Radon measure on $V$. It has a unique bounded classical
zero-boundary solution. The solutions for its finite
partial sums decrease uniformly to that solution, with the
explicit error bound
\[
 0\leq u_N-u\leq
 \sum_{i>N}m_i^{1/n}\|g_i\|_{L^\infty(V)}.
 \tag{8}\label{eq:r3-atomic-tail-bound}
\]
Their weighted BT measures converge on every original compact
smooth test.
Adding a constant gives the same conclusion for any constant
boundary value.
\end{corollary}

\begin{proof}
Set $\mu_N=\sum_{i=1}^N m_i\delta_{z_i}$ and
$u_N=G(\mu_N)$, using the finite atomic theorem. The measure
order lemma makes $u_N$ decrease. The mixed sandwich gives
\[
 \sum_{i=1}^N m_i^{1/n}g_i\leq u_N\leq0.
\]
Condition \eqref{eq:r3-atomic-summability} makes the series
$s=\sum_i m_i^{1/n}g_i$ uniformly convergent and bounded
on $V$. Its finite sums decrease and are classical, so $s$
is classical. Its boundary limit is zero: first make the
uniform tail small, then use the boundary limits of the
finite initial sum. Thus
\[
 s\leq u_N\downarrow u\leq0
\]
gives a bounded classical limit with zero boundary limit.
In fact, repeated use of the mixed sandwich gives, for every
$N,M\geq1$,
\[
 u_N+\sum_{i=N+1}^{N+M}m_i^{1/n}g_i
       \leq u_{N+M}\leq u_N.
\]
Let $M\to\infty$. This proves
\eqref{eq:r3-atomic-tail-bound}, and hence uniform convergence.

The fixed-background decreasing BT theorem gives
$(Du_N)^n\to(Du)^n$. For a nonnegative original smooth
compact scalar test $\chi$ it follows that
\[
 \sum_i m_i\chi(z_i)
    =\lim_N\int\chi\,d\mu_N
    =\int\chi(Du)^n<\infty.
\]
Choose such a cutoff at least one near any prescribed
compact subset of $V$. This proves local finiteness of the
atomic sum; on this locally compact metrizable space it is
the corresponding Radon measure. The same identity gives
$(Du)^n=\mu$ on all original tests. The weighted decreasing
theorem also gives $u_N(Du_N)^n\to u(Du)^n$ locally.
Uniqueness follows from the two compact-gap conditions and
pointwise ambient uniqueness, without a global total-mass
assumption.
\end{proof}

For the singleton, the R2 compact-extremal formula identifies
\[
 g_z=\operatorname{Cap}_{[1]}(\{z\};V)^{-1/n}U_{\{z\}},
 \qquad
 \|g_z\|_\infty
       =\operatorname{Cap}_{[1]}(\{z\};V)^{-1/n}.
\]
The point constraint is exact at the divisorial point, so
$U_{\{z\}}(z)=-1$. Thus \eqref{eq:r3-atomic-summability}
is equivalently the explicit condition
\[
 \sum_i
 \left(\frac{m_i}
 {\operatorname{Cap}_{[1]}(\{z_i\};V)}\right)^{1/n}<\infty.
\]
The order, sandwich and countable-atomic conclusions are
\texttt{DIRECT\_COROLLARY} of the finite atomic theorem and
the proved positive mixed calculus, with their proofs given
above. No arbitrary diffuse-measure existence or stability
under moving supports is asserted here.

\begin{corollary}[A quantitative bound for atomic changes of mass]
\label{cor:r3-atomic-quantitative-stability}
Let $\mu=\sum_i m_i\delta_{z_i}$ and
$\nu=\sum_i n_i\delta_{z_i}$ be locally finite nonnegative atomic
measures on one finite or countable set of distinct divisorial
points in $V$. Suppose bounded classical zero-boundary solutions
$u_\mu,u_\nu$ exist. For
\[
 \kappa_i=\operatorname{Cap}_{[1]}(\{z_i\};V)>0,
 \qquad
 d_+=\sum_{m_i>n_i}\left(\frac{m_i-n_i}{\kappa_i}\right)^{1/n},
 \quad
 d_-=\sum_{n_i>m_i}\left(\frac{n_i-m_i}{\kappa_i}\right)^{1/n},
\]
assume $d_++d_-<\infty$. Then
\[
 -d_+\leq u_\mu-u_\nu\leq d_-,
 \qquad
 \|u_\mu-u_\nu\|_{L^\infty(V)}
       \leq\max(d_+,d_-)
       \leq\sum_i\left(\frac{|m_i-n_i|}{\kappa_i}\right)^{1/n}.
 \tag{9}\label{eq:r3-atomic-stability-bound}
\]
In particular this applies to the finite atomic solutions and to
any two countable solutions constructed under
\eqref{eq:r3-atomic-summability} on their combined support.
\end{corollary}

\begin{proof}
Put $\alpha=(\mu-\nu)_+$ and $\beta=(\nu-\mu)_+$, with
the positive and negative parts taken coefficient by coefficient.
The summability hypothesis and the preceding corollary produce
zero-boundary solutions $G(\alpha),G(\beta)$ with
\[
 -d_+\leq G(\alpha)\leq0,
 \qquad -d_-\leq G(\beta)\leq0.
\]
Positivity of every mixed BT term gives
\[
 (D(u_\nu+G(\alpha)))^n\geq\nu+\alpha\geq\mu.
\]
The measure-order lemma therefore yields
$u_\nu+G(\alpha)\leq u_\mu$. Interchanging $\mu,\nu$
gives $u_\mu+G(\beta)\leq u_\nu$. The displayed bounds
on $G(\alpha),G(\beta)$ prove the conclusion. This proof does
not need the existence of a separate solution for $\mu+\nu$.
\end{proof}

On the extendible-exhaustion domains of
Theorem~\ref{thm:r3-local-atomic-continuity}, every finite partial
solution extends continuously to $\overline V$. The uniform
tail bound consequently makes every countable solution satisfying
\eqref{eq:r3-atomic-summability} continuous on $\overline V$
as well. This last statement uses that additional domain hypothesis;
the uniform estimate itself holds on every domain fixed at the
start of this section.

% END INCLUDED SOURCE: sections/nonzero_dirichlet


% BEGIN INCLUDED SOURCE: sections/atomic_contact
\section{Contact and unique depth parameters for positive atomic data}
\label{sec:r3-atomic-contact}

Keep the fixed domain $V$, the distinct divisorial points
$z_1,\ldots,z_r$, and the actual envelopes $U_a$ from
Section~\ref{sec:r3-nonzero-dirichlet}. The statements below are
\texttt{DIRECT\_COROLLARY} of that construction and local maximality.
They do not require interior continuity of $U_a$.

\begin{proposition}[Complementarity at a finite point obstacle]
\label{prop:r3-atomic-complementarity}
For every $a\in[0,\infty)^r$ and every index $i$,
\[
 c_i(a)>0\quad\Longrightarrow\quad U_a(z_i)=-a_i.
\]
Equivalently, the atomic measure of $U_a$ is supported on those
point constraints at which equality holds.
\end{proposition}

\begin{proof}
Suppose that $U_a(z_i)<-a_i$, and write
$\delta=-a_i-U_a(z_i)>0$. Choose an ordinary
$W\Subset V$ containing $z_i$ and no other $z_j$. Let
$v\in\Cb(W)$ agree with $U_a$ off a compact subset
$L\Subset W$. The function $\max(v,U_a)-U_a$ vanishes off
$L$ and is bounded there, so it has a finite upper bound $M\geq0$.
Choose $0<t\leq1$ with $tM<\delta$. On $W$ put
\[
 w=(1-t)U_a+t\max(v,U_a),
\]
and set $w=U_a$ off $L$. The formulas agree on the ordinary
overlap. Positive linear combinations, the maximum rule, and
ordinary gluing show that $w$ is classical psh. It is bounded,
has ordinary boundary limit zero, and satisfies $w\geq U_a$.
The boundary maximum principle gives $w\leq0$.
At $z_i$ one has
\[
 w(z_i)\leq U_a(z_i)+tM<-a_i;
\]
at every other constraint point $w=U_a$. Thus
$w\in\mathcal F_a$, using $w\geq U_a\geq-A(a)$ for the
lower bound. The actual greatest-member property implies
$w\leq U_a$. Since $t>0$, it follows that $v\leq U_a$.

This proves local maximality of $U_a$ on $W$. The proved
maximality criterion gives $(DU_a)^n|_W=0$, so $c_i(a)=0$.
The contrapositive is the assertion. Only actual constraints
were used; no quasi-everywhere constraint was substituted.
\end{proof}

\begin{corollary}[Unique depth coordinates for strictly positive masses]
\label{cor:r3-active-mass-homeomorphism}
Let
\[
 \mathcal A=\{a\in[0,\infty)^r:c_i(a)>0\text{ for all }i\}.
\]
Then $\mathcal A$ is an open subset of $(0,\infty)^r$, and
the full mass map
\[
 c:\mathcal A\longrightarrow(0,\infty)^r,
 \qquad a\longmapsto(c_1(a),\ldots,c_r(a)),
\]
is a homeomorphism. Its inverse is the actual point-depth map
\[
 (m_1,\ldots,m_r)\longmapsto
 \bigl(-G(\textstyle\sum_i m_i\delta_{z_i})(z_j)\bigr)_{j=1}^r.
\]
Consequently the normalized mass map $F$ restricts to a
homeomorphism from $\mathcal A\cap\Delta^{r-1}$ onto the
relative interior of $\Delta^{r-1}$. For a target
$y\in\operatorname{relint}\Delta^{r-1}$ its unique normalized
depth vector is
\[
 a_i(y)=
 \frac{-G(\sum_jy_j\delta_{z_j})(z_i)}
 {\sum_{\ell=1}^r-G(\sum_jy_j\delta_{z_j})(z_\ell)}.
\]
\end{corollary}

\begin{proof}
Continuity of the coefficients makes $\mathcal A$ open, and
face preservation puts it inside $(0,\infty)^r$. Given a
strictly positive mass vector, the finite atomic theorem
constructs its unique solution as $tU_a=U_{ta}$ for some
$t>0$ and $a\in\Delta^{r-1}$. Its mass vector is exactly the
target, so $ta\in\mathcal A$. Complementarity gives
$-U_{ta}(z_i)=ta_i$ for every $i$. This proves surjectivity
and the displayed inverse formula.

If $c(a)=c(b)$ for $a,b\in\mathcal A$, actual Dirichlet
uniqueness gives $U_a=U_b$. Complementarity then gives
$a_i=-U_a(z_i)=-U_b(z_i)=b_i$ for every $i$. Thus the
map is injective. The inverse is continuous because the
fixed-support solution map is continuous in the uniform
norm, and evaluation at a fixed point is continuous for
that norm. The mass map itself is continuous by its proved
local Lipschitz estimate.

For normalized masses, scaling the unique solution scales
all its depths by the same positive factor. Dividing those
depths by their sum gives the displayed point of the simplex.
The sum is positive because the full depth vector is in
$\mathcal A$. The formula is continuous, and homogeneity
shows it is the inverse of the restricted normalized map.
\end{proof}

This does not assert that the mass map is injective on the
whole nonnegative cone. A constraint with strict slack has
zero mass, and its prescribed depth need not be recoverable
from the solution. The inverse theorem is restricted to the
region on which every prescribed atom has positive mass.

% END INCLUDED SOURCE: sections/atomic_contact


% BEGIN INCLUDED SOURCE: sections/divisorial_obstacle_continuity
\section{Continuous envelopes with divisorial point constraints}
\label{sec:r3-divisorial-obstacle-continuity}

This section uses the continuous ample-envelope theorem of B and the
divisorial no-jump theorem proved in R2. The first result is a
\texttt{DIRECT\_COROLLARY} of those inputs. The local application also
uses the explicit order-preserving common-collar construction. No
continuity theorem for general homogeneous Dirichlet solutions is
assumed or obtained.

\subsection{A continuous insertion argument}

\begin{lemma}[Continuous insertion with a uniform error]
\label{lem:r3-continuous-insertion}
Let $Z$ be a compact Hausdorff space. Let $a$ be a bounded upper
semicontinuous real function, and let $f$ be a bounded lower
semicontinuous real function with $a\leq f$. For every $\epsilon>0$
there is a continuous function $g$ such that
\[
 a-\epsilon\leq g\leq f.
\]
\end{lemma}

\begin{proof}
For each $x\in Z$, choose a neighborhood $O_x$ on which
$a(y)<a(x)+\epsilon/4$ and $f(y)>f(x)-\epsilon/4$.
The constant $c_x=a(x)-\epsilon/2$ then satisfies
$a(y)-\epsilon<c_x<f(y)$ on $O_x$. Take a finite subcover
and a continuous partition of unity subordinate to it. The same
partition-weighted sum of the constants $c_x$ is the required $g$.
\end{proof}

Let $Y/K$ be smooth, projective and geometrically integral, and let
$L$ be ample. All metric inequalities below are inequalities of local
weights; their differences are globally defined real functions.
The fixed field and its absolute value are unchanged.

\begin{theorem}[A lower semicontinuous obstacle with an attained envelope]
\label{thm:r3-lsc-obstacle-continuity}
Let $f$ be a lower semicontinuous metric on $L$ whose difference from
a continuous reference metric is bounded. Put
\[
 h_f=\sup\{r:r\text{ is a bounded classical metric on }L, r\leq f\},
 \qquad P_f=h_f^*.
\]
If $P_f\leq f$ at every point, then $h_f=P_f$ and $P_f$ is continuous.
More precisely, for every $\epsilon>0$ there is a continuous obstacle
$g_\epsilon\leq f$ such that
\[
 P_f-\epsilon\leq P_Lg_\epsilon\leq P_f.
\]
\end{theorem}

\begin{proof}
Choose a continuous semipositive reference metric $r_0$.
A sufficiently negative translate of $r_0$ is a competitor, and
every competitor is bounded above by $r_0+C$ for one constant $C$.
Replacing competitors by their maxima with one such negative
translate gives a uniformly bounded family without changing the
supremum. Corollary~\ref{cor:r2-projective-upper-envelope} therefore
makes $P_f$ a bounded classical metric. Under the stated hypothesis
it is itself a competitor. Thus $h_f=P_f$ pointwise.

Apply Lemma~\ref{lem:r3-continuous-insertion} to $P_f-r_0$ and
$f-r_0$, and add $r_0$ back. This gives a continuous metric
$g_\epsilon$ with $P_f-\epsilon\leq g_\epsilon\leq f$.
By B, Theorem C.6, $P_Lg_\epsilon$ is continuous and semipositive.
The actual-competitor comparison
Lemma~\ref{lem:r2-envelope-competitor} gives
$P_f-\epsilon\leq P_Lg_\epsilon$.
On the other hand, $P_Lg_\epsilon\leq g_\epsilon\leq f$, so it is
a competitor for $h_f$ and is at most $P_f$. These uniform
inequalities prove continuity of $P_f$.
\end{proof}

\begin{corollary}[Discontinuities confined to divisorial points]
\label{cor:r3-divisorial-lsc-obstacle}
In the preceding theorem, suppose that every discontinuity point
of $f$ is ambient divisorial. Then $P_f\leq f$ pointwise, and the
conclusion of the theorem holds.
\end{corollary}

\begin{proof}
At a continuity point of $f$, upper regularization of $h_f\leq f$
gives $P_f\leq f$. At a divisorial point $z$, the local
divisorial no-jump result, Corollary~\ref{cor:r2-divisorial-no-jump},
gives $P_f(z)=h_f(z)\leq f(z)$. These two cases cover the space.
\end{proof}

\begin{corollary}[Finite divisorial obstacles]
\label{cor:r3-finite-divisorial-obstacle}
Let $F$ be a continuous semipositive metric on $L$, let
$z_1,\ldots,z_N$ be distinct divisorial points, and let $a_i\geq0$.
The greatest bounded classical metric subject to
\[
 r\leq F,\qquad r(z_i)\leq F(z_i)-a_i\quad(1\leq i\leq N)
\]
exists and is continuous. Denote it by $P_a$. Then
\[
 F-\max_i a_i\leq P_a\leq F,
 \qquad
 \|P_a-P_b\|_\infty\leq\max_i|a_i-b_i|.
\]
\end{corollary}

\begin{proof}
The obstacle $f_a=F-\sum_i a_i\mathbf1_{\{z_i\}}$ is lower
semicontinuous and is continuous off the finite divisorial set.
The preceding corollary applies. The displayed lower bound is an
admissible continuous metric. If $\delta=\max_i|a_i-b_i|$, then
$P_a-\delta$ is admissible for the $b$-constraints, whence
$P_a-\delta\leq P_b$. Interchange $a$ and $b$.
\end{proof}

The lower semicontinuity hypothesis is deliberate. The point
obstacle has downward jumps and is generally not upper
semicontinuous. It must not be inserted into the upper
semicontinuous-obstacle theorem without the preceding argument.

\subsection{An application on domains with an extendible exhaustion}

\begin{lemma}[A framed smooth compactification of an affine variety]
\label{lem:r3-framed-compactification}
A smooth geometrically integral affine variety $Y/K$ has a smooth
projective compactification $\overline Y$, an ample line $A$ on
$\overline Y$, and a section $s$ with
$\overline Y\setminus\operatorname{div}(s)=Y$.
After taking a power, $A$ is very ample and $s$ is one coordinate
section of its projective embedding.
\end{lemma}

\begin{proof}
Embed $Y$ in affine space, normalize its projective closure, and
pull back the hyperplane line and its section at infinity.
Normalization is finite and is an isomorphism on the smooth open
$Y$. The pulled-back line is ample, and the nonvanishing locus of
its section is exactly $Y$.

Use the characteristic-zero projective resolution preserving $Y$
already invoked in B, Proposition 4.5. The ample-boundary property
can be retained at each blowup in this resolution. Indeed, suppose
$H$ is ample and $s_H$ has nonvanishing locus $Y$, and let
$\pi:Z'\to Z$ be a blowup of an ideal $I$ whose cosupport is in
the boundary. Write $I\mathcal O_{Z'}=\mathcal O_{Z'}(-E)$.
For sufficiently large $N$, the line
\[
 H'=\pi^*H^N\otimes\mathcal O_{Z'}(-E)
\]
is ample. Increasing $N$ if necessary, the section $s_H^N$
belongs locally to $I H^N$: this follows from the containment
of the cosupport in $\operatorname{div}(s_H)$ and the
Noetherian radical criterion, using finitely many affine charts.
The induced section of $H'$ has divisor
$N\pi^*\operatorname{div}(s_H)-E$. For still larger $N$ this
divisor is effective and has exactly the full boundary as its
support. Thus the property persists. The same finite-pullback
argument handles a normalization if one is inserted. The
resulting smooth projective compactification has the required
ample line and section. A sufficiently high very ample power,
with the powered section included in a basis, gives the last
assertion.
\end{proof}

Let $Y/K$ now be smooth, geometrically integral and affine.
Suppose that $\rho$ is a continuous classical psh function on
$Y^{\an}$ and that $V$ is a connected component of
$\{\rho<0\}$ with
\[
 V\Subset Y^{\an},\qquad
 \varnothing\ne\partial V\subset\{\rho=0\}.
\]
This includes the norm-box exhaustion domains constructed from
$\log\max(1,|g_1|,\ldots,|g_N|)-\epsilon$ in R2. It is stronger
than merely having an exhaustion defined on $V$.

For distinct divisorial points $z_i\in V$ and $a_i\geq0$, let
\[
 U_a=\sup\{u\in\Cb(V):u\text{ is bounded},\ u\leq0,
                         \ u(z_i)\leq-a_i\text{ for every }i\}.
\]
The divisorial no-jump theorem makes this actual supremum
classical and admissible, as in the finite-atomic construction.

\begin{theorem}[Continuity of the finite-point Dirichlet envelopes]
\label{thm:r3-local-atomic-continuity}
Under the displayed extendible-exhaustion hypotheses, $U_a$ is
continuous on $\overline V$, with value zero on $\partial V$.
Consequently the finite divisorial-atomic Dirichlet solutions with
constant boundary values constructed in this round are continuous
on $\overline V$ for these domains.
\end{theorem}

\begin{proof}
Choose $C>0$ with $C\rho(z_i)\leq-a_i$ for every $i$.
The bounded function $C\rho|_V$ is a competitor. Hence
\[
 C\rho\leq U_a\leq0\quad\hbox{on }V.
\]
These inequalities already give the claimed boundary values.

For $\delta>0$, define a function on $Y^{\an}$ by
\[
 w_\delta=
 \begin{cases}
 \max(U_a,C\rho+\delta),&\text{on }V,\\
 C\rho+\delta,&\text{off }V.
 \end{cases}
\]
Near every point of $\partial V$, one has $C\rho+\delta>0\geq U_a$
on the part lying in $V$. Thus the two prescriptions agree on an
ordinary neighborhood of the boundary. Ordinary gluing and the
maximum rule show that $w_\delta$ is classical psh.
As $\delta\downarrow0$, these functions decrease locally boundedly
to the classical psh function
\[
 w=
 \begin{cases}U_a,&\text{on }V,\\ C\rho,&\text{off }V.
 \end{cases}
\]
Moreover $w\leq g:=\max(0,C\rho)$, and $g$ is continuous
classical psh on $Y^{\an}$.

Use the framed compactification of
Lemma~\ref{lem:r3-framed-compactification}. Its one affine
coordinate chart contains the entire $Y$. Large upper bounds
on its finitely many affine coordinates give nested compact
norm boxes whose innermost ordinary interior contains
$\overline V$; all outer boxes remain in $Y^{\an}$.
The functions $w$ and $g$ are defined on a neighborhood of the
outer box and have common bounds there. Apply the explicit
common-collar construction from B, Lemma D.2, in the
order-preserving form used in
Theorem~\ref{thm:r2-local-upper-envelope}.
For clarity, in one fixed frame that construction is
\[
 \mathcal G(f)=\max(f+c,mb)
\]
on the outer box and the same fixed metric off its middle box.
Here $b=0$ on the inner box, the common translation $c$ makes
$f+c>0$ there, and $m$ forces equality with the fixed metric on
the outer collar. It gives bounded classical metrics
\[
 R=\mathcal G(w)\leq F=\mathcal G(g)
\]
on one ample line. The metric $F$ is continuous and classical.
B, Theorem 4.4, gives decreasing continuous model-semipositive
metrics converging to $F$; compactness and Dini's theorem make
this convergence uniform. Thus $F$ is a continuous semipositive
metric, as required by the projective finite-point corollary.
Their weights
on $V$ are respectively $U_a+c$ and $c$, since $g=0$ on $V$.

Let $P_a$ be the continuous projective finite-point envelope of
Corollary~\ref{cor:r3-finite-divisorial-obstacle}, with upper
obstacle $F$ and depths $a_i$. The metric $R$ is a competitor,
so $R\leq P_a$. Conversely, the local weight $P_a-c$ on $V$
is bounded classical, is at most zero, and satisfies the actual
point constraints. It is therefore at most $U_a$. On $V$ this
gives
\[
 U_a+c=R\leq P_a\leq U_a+c.
\]
Thus $U_a=P_a-c$ on $V$, proving interior continuity. Combined
with the previously established ordinary boundary limits, this
proves continuity on $\overline V$.
\end{proof}

\paragraph{Scope.}
The projective envelope statement is a direct consequence of B
and R2. The local theorem is a \texttt{DIRECT\_COROLLARY} after
the extension and common-collar construction above. Its extra
exhaustion hypothesis is explicit. This argument neither proves
interior continuity for the general peak-barrier Perron solution
of R2 nor replaces arbitrary boundary values by constant ones.

% END INCLUDED SOURCE: sections/divisorial_obstacle_continuity


% BEGIN INCLUDED SOURCE: sections/skeletal_dirichlet
\section{Non-atomic Dirichlet data on fixed Gauss skeleton boxes}
\label{sec:r3-skeletal-dirichlet}

Fix the field, the smooth integral affine $Y/K$, and the connected
ordinary exhaustion domain $V\Subset Y^{\an}$ of
\eqref{eq:r3-dirichlet-domain}. In particular $n=\dim Y\geq1$,
$\rho<0$ on $V$, and $\rho$ is continuous on $\overline V$ with
ordinary boundary value zero. The domain remains fixed throughout.
All Monge--Amp\`ere measures use $[1]$, the original absolute value,
and the original CLD tests and calibration.

The new step is a compactness mechanism for a specified class of
supports. It proves weak-measure stability in that class; it does not
assume such stability. The main existence theorem is recorded as
\texttt{NEW\_PROVED}, relative to the finite atomic construction and
the preceding R1/R2 results.

\subsection{The support geometry and its function estimates}

For a bounded open rectangular box $I\subset\mathbf R^n$, put
\[
 \mathbb A(I)=
 \{x\in(\mathbf G_{m,K}^n)^{\an}:
              (\log|T_1(x)|,\ldots,\log|T_n(x)|)\in I\}.
\]
Its Gauss section $\sigma:I\to\mathbb A(I)$ is characterized by
\[
 \left|\sum_{m\in\mathbf Z^n}c_mT^m\right|(\sigma(x))
          =\max_m |c_m|e^{m\cdot x}
\]
for Laurent polynomials. Suppose that, for $1\leq\ell\leq L$,
there are $K$-analytic isomorphisms onto ordinary opens
\[
 \Phi_\ell:\mathbb A(I_\ell)\xrightarrow{\sim}\Omega_\ell
          \subset V
\]
and compact rectangular boxes $J_\ell\Subset I_\ell$, each with
nonempty Euclidean interior. Set
\[
 K_\ell=\Phi_\ell(\sigma(J_\ell)),\qquad
 K=\bigcup_{\ell=1}^L K_\ell\Subset V.
 \tag{S}\label{eq:r3-gauss-box-support}
\]
The larger full polyannulus, and the positive margin between its box
and $J_\ell$, are part of the hypotheses. A compact subset merely
called a skeleton is not being substituted for these data.
The measure may be supported on any closed subset of this fixed $K$;
the approximating divisorial points below are allowed to lie anywhere
in $K$. Thus an irrational-radius face need not itself contain a
dense set of divisorial points.

\begin{lemma}[Uniform continuity on an inner Gauss box]
\label{lem:r3-gauss-box-equicontinuity}
If $u\in\Cb(\mathbb A(I))$ and $-B\leq u\leq0$, then
$F_u(x)=u(\sigma(x))$ is convex in each coordinate separately.
For $J\Subset I$ as above there is a constant $C_{I,J}$,
independent of $u$ and $B$, such that
\[
 |F_u(x)-F_u(y)|\leq C_{I,J}B\|x-y\|_1
                         \quad(x,y\in J).
\]
In particular a uniformly bounded family of classical functions
restricts to an equicontinuous family on $K$ in
\eqref{eq:r3-gauss-box-support}.
\end{lemma}

\begin{proof}
Fix a coordinate $i$ and fix all the other real coordinates $x_j$.
Let $F/K$ be the completed residue field of the weighted Gauss
point with these $n-1$ coordinates. For $n=1$, take $F=K$.
The corresponding coordinate elements $c_j\in F^\times$ have
$|c_j|=e^{x_j}$ and the Gauss independence property. The map from
the $F$-annulus with $\log|t|\in I_i$ given by
\[
 T_i=t,\qquad T_j=c_j\quad(j\ne i)
\]
lands in $\mathbb A(I)_F$. At its radius-$e^s$ Gauss point,
the image in the original $K$-space is
$\sigma(x_1,\ldots,s,\ldots,x_n)$. This follows directly by
grouping a Laurent polynomial according to its exponent of $T_i$
and applying the Gauss norm first to the remaining coefficients.

The all-complete-extensions classical hypothesis makes the
pullback of $u$ classical subharmonic on that annulus. Let
$a<b$ be any two logarithmic radii strictly inside $I_i$.
Subtract the affine function of $\log|t|$ interpolating its
values at the two Gauss endpoints. This affine function is
pluriharmonic, so the difference is still classical psh.
Its values at the two Shilov points of the closed subannulus
are zero. The affinoid maximum principle
\cite[Theorem~6.4]{GR25} makes the difference nonpositive.
This is exactly the convexity inequality for the $i$th
coordinate of $F_u$.

A real convex function with range in $[-B,0]$ on an interval
has uniformly bounded secant slopes on every inner interval:
comparison with secants to the outer endpoints bounds their
absolute values by the oscillation divided by the margin.
If $\delta>0$ is the smallest coordinate margin from $J$ to
the endpoints of $I$, the bound $2B/\delta$ is sufficient.
Change one coordinate at a time inside $J$ and add these
estimates. This proves the displayed inequality.

Each Gauss section is continuous and is a homeomorphism onto
its image, since its coordinates recover $x$. Pullback by
$\Phi_\ell$ preserves the classical definition. Apply the
estimate on each of the finitely many boxes. All auxiliary
complete extensions above were used only for the curve tests
of functions, not for any asserted base change of currents.
\end{proof}

\begin{lemma}[A uniform singleton capacity on the original domain]
\label{lem:r3-gauss-box-point-capacity}
There is $c_K>0$, depending on the data in
\eqref{eq:r3-gauss-box-support} and on $V$, such that
\[
 \operatorname{Cap}_{[1]}(\{z\};V)\geq c_K
                                      \quad(z\in K).
\]
Moreover $K$ has a dense set of ambient divisorial points over
the original field $K$.
\end{lemma}

\begin{proof}
First work on $\mathbb A(I_\ell)$. Choose $A_\ell>0$ with
\[
 A_\ell>\sum_{i=1}^n
       (\sup (I_\ell)_i-\inf(J_\ell)_i).
\]
For $a\in J_\ell$ the continuous classical function
\[
 h_a=\frac1{A_\ell}\sum_{i=1}^n
                       \max(0,\log|T_i|-a_i)
\]
lies in $[0,1]$ on the same $\mathbb A(I_\ell)$, independently
of $a$. Its ordinary Monge--Amp\`ere measure is
\[
 (Dh_a)^n=\frac{n!}{A_\ell^n}\delta_{\sigma(a)}.
 \tag{10}\label{eq:r3-gauss-box-test-atom}
\]
Here is a check of that normalization on all original tests.
Replace each maximum by a smooth convex regularization
$\gamma_\epsilon(\log|T_i|-a_i)$, with second derivative
a nonnegative approximate identity of mass one. The ordinary
product calibration on the split torus gives, for every
original compact smooth scalar test $\chi$,
\[
 \int_{\mathbb A(I_\ell)}\chi
       \bigwedge_{i=1}^nD\gamma_\epsilon(\log|T_i|-a_i)
  =\int_{I_\ell}\chi(\sigma(x))
       \prod_{i=1}^n\gamma_\epsilon''(x_i-a_i)\,dx.
\]
To check the formula when $\chi$ requires additional moments,
adjoin those moments to the coordinate map and subdivide its
tropicalization. On cells whose projection to the coordinate
box has rank less than $n$, the displayed top form vanishes.
On a full-rank cell the canonical calibration is the degree
over the coordinate Gauss skeleton. The coordinate map here
is the identity inclusion into the split torus, so its degree
is one and its point over $x$ is precisely $\sigma(x)$.
The construction of the canonical calibration and its
compatibility with projections
\cite[\S3.5.1, Lemma~3.5.2, and \S\S3.7--3.8]{CLD12}
therefore give the formula with the actual value
$\chi(\sigma(x))$, also after every required refinement.
For small
$\epsilon$ its support lies in $I_\ell$, and the right side
tends to $\chi(\sigma(a))$. Locally uniform BT convergence
gives the corresponding mixed product of the maxima.
Repeated coordinate factors vanish already for the smooth
rank-one forms. Expansion of the $n$th power leaves precisely
$n!$ distinct-coordinate terms, proving
\eqref{eq:r3-gauss-box-test-atom}. No finite extension or
change in the integration current was made.

Choose an intermediate open box $J_\ell^+$ with
$J_\ell\Subset J_\ell^+\Subset I_\ell$ and put
$W_\ell=\Phi_\ell(\mathbb A(J_\ell^+))$.
Its closure is compact in $\Omega_\ell$. The reverse
direction of Theorem~\ref{thm:r2-domain-capacity-comparison}
provides a fixed $C_\ell<\infty$ such that
\[
 \operatorname{Cap}_{[1]}(E;\Omega_\ell)
       \leq C_\ell\operatorname{Cap}_{[1]}(E;V)
                                  \quad(E\subset W_\ell).
\]
The isomorphism $\Phi_\ell$ has degree one and preserves
the original products and tests. Apply the last inequality
to a singleton and use \eqref{eq:r3-gauss-box-test-atom}.
One may take
\[
 c_K=\min_\ell\frac{n!}{C_\ell A_\ell^n}>0.
\]
This lower bound uses the proved reverse collar comparison;
ordinary domain monotonicity alone has the opposite direction.

Write $\lambda=-\log|\pi|>0$ for a uniformizer of the fixed
field. The points $a\in J_\ell\cap(\lambda\mathbf Q)^n$
are dense in $J_\ell$. Their Gauss points are the unique
Shilov points of strict fixed-radius polyannuli: rational
valuation thresholds are expressed by finitely many integral
power inequalities over $K$. They are therefore divisorial,
and remain so on ordinary neighborhoods and under the analytic
isomorphism, by \cite[Proposition~A.3 and Corollary~A.6]{GM19}.
This gives a dense divisorial subset of $K$, without extending
the ground field used for the data or currents.
\end{proof}

\subsection{From skeletal compactness to uniform compactness of solutions}

\begin{lemma}[A common support determines the uniform distance]
\label{lem:r3-support-uniform-isometry}
Let $u,v\in\Cb(V)$ be bounded, with ordinary boundary limit
zero, and suppose both $(Du)^n$ and $(Dv)^n$ are supported
on a compact subset $H\Subset V$. Then
\[
 \|u-v\|_{L^\infty(V)}=\sup_H|u-v|.
\]
\end{lemma}

\begin{proof}
Put $b=\sup_H|u-v|$. For every $\epsilon>b$, the Borel gap
$E=\{u<v-\epsilon\}$ is disjoint from $H$, so
$(Du)^n(E)=0$. The boundary values make its closure compact
in $V$, and give every compact-gap condition required for
pointwise ambient domination. That principle gives
$u\geq v-\epsilon$ on $V$. Interchange $u,v$ and let
$\epsilon\downarrow b$. The reverse norm inequality follows
from $H\subset V$. No strict-gap set was treated as an
ordinary open set.
\end{proof}

\begin{lemma}[Uniform bounds for finite atomic data in $K$]
\label{lem:r3-skeletal-finite-bounds}
Let $\nu$ be a finite nonnegative sum of divisorial atoms in
$K$, and let $u_\nu$ be its zero-boundary solution from
Theorem~\ref{thm:r3-finite-atomic-dirichlet}. Put
$M=\nu(K)$ and
\[
 d_K=\min_K(-\rho)>0,
 \qquad b_K=\max(-1,\rho/d_K).
\]
Then
\[
 \|u_\nu\|_\infty\leq(M/c_K)^{1/n},
 \qquad (M/c_K)^{1/n}b_K\leq u_\nu\leq0.
 \tag{11}\label{eq:r3-skeletal-uniform-bounds}
\]
In particular all such solutions with bounded total mass have
one common ordinary boundary barrier.
\end{lemma}

\begin{proof}
The assertion is immediate for $M=0$. Otherwise discard
zero-mass points. The finite atomic construction writes
$u_\nu=tU_a$, with $t>0$ and $a\in\Delta^{r-1}$.
Put $A=\max_i a_i>0$ and choose $i$ with $a_i=A$.
The defining lower bound and the actual constraint give
$U_a(z_i)=-A$. Consequently $B=\|u_\nu\|_\infty=tA$.
The function $u_\nu/B$ lies in $[-1,0]$ and equals $-1$
at $z_i$. It is a competitor for the singleton extremal,
so $u_\nu/B\leq U_{\{z_i\}}\leq0$. Both functions have
ordinary boundary limit zero. The exact total-mass order
Lemma~\ref{lem:r2-total-mass-order} and the compact-extremal
formula give
\[
 c_K\leq\operatorname{Cap}_{[1]}(\{z_i\};V)
      =\int_V(DU_{\{z_i\}})^n
      \leq\int_V(D(u_\nu/B))^n=\frac{M}{B^n}.
\]
All masses here are finite; the cited order theorem in fact
also permits infinite masses without cancellation.
This proves the norm bound.

At every constraint point $z_j\in K$, the function $b_K$
equals $-1$, which is at most $-a_j/A$. The greatest-member
description of $U_{a/A}$ therefore gives
$b_K\leq U_{a/A}=u_\nu/B$. Since $b_K\leq0$ and
$B\leq(M/c_K)^{1/n}$, the second bound follows. The same
$b_K$ works for every choice of the finite support in $K$.
\end{proof}

\subsection{Arbitrary finite measures on the prescribed skeleton}

\begin{theorem}[A non-atomic fixed-domain Dirichlet class]
\label{thm:r3-skeletal-dirichlet}
Assume \eqref{eq:r3-gauss-box-support}. For every finite
nonnegative Radon measure $\mu$ supported on $K$ and every
$b\in\mathbf R$, there is a unique bounded actual classical
psh function $u\in\Cb(V)$ with
\[
 (Du)^n=\mu,
 \qquad\lim_{\substack{x\to\xi\\x\in V}}u(x)=b
                                      \quad(\xi\in\partial V).
\]
The bounds \eqref{eq:r3-skeletal-uniform-bounds} hold for
$u-b$ with $M=\mu(K)$.

If $\mu_j\to\mu$ weakly as finite measures on this same
compact $K$, and $b_j\to b$, then their solutions converge
uniformly on $V$. Their weighted measures also satisfy
\[
 u_j(Du_j)^n\longrightarrow u(Du)^n
\]
against every compact continuous scalar test, hence against
every original compact smooth scalar test.

If $V$ additionally satisfies the extendible-exhaustion
hypotheses of Theorem~\ref{thm:r3-local-atomic-continuity},
all these solutions are continuous on $\overline V$.
\end{theorem}

\begin{proof}
First take $b=0$. The compact space $K$ is metrizable; it is
a finite union of compact Gauss boxes. Its dense divisorial
subset permits finite atomic measures $\mu_j$ supported in
$K$, with $\mu_j(K)=\mu(K)$, that converge weakly to $\mu$.
For example take finite Borel partitions of successively
smaller diameter and replace each part's mass by a nearby
divisorial point in $K$. Uniform continuity of functions in
$C(K)$ verifies the asserted weak convergence.

Let $u_j$ be the finite atomic solutions. Their uniform
bound and common boundary barrier are
\eqref{eq:r3-skeletal-uniform-bounds}. Lemma
\ref{lem:r3-gauss-box-equicontinuity} makes their restrictions
equicontinuous on each $K_\ell$. Arzel\`a--Ascoli and a
finite diagonal extraction give a subsequence uniformly
Cauchy on all of $K$. Lemma
\ref{lem:r3-support-uniform-isometry} then makes this
subsequence uniformly Cauchy on the entire original domain
$V$. Denote its uniform limit by $u$.

Uniform limits here are classical. Explicitly, take a further
subsequence with uniform errors at most $4^{-j}$; adding
$2\cdot4^{-j}$ makes it decrease to $u$. The decreasing
closure of the classical class applies. The common finite
bound makes $u$ locally bounded, and the common $b_K$
barrier gives its ordinary zero boundary limit.

Uniform convergence implies local capacity convergence on
every prescribed small norm box. The proved fixed-background
BT stability, Theorem~\ref{thm:capacity-stability}, therefore
gives $(Du_j)^n\to(Du)^n$ on all original compact tests.
But $(Du_j)^n=\mu_j\to\mu$ on those tests. The associated
Radon measures are determined by the original smooth scalar
tests, so $(Du)^n=\mu$. This passage uses uniform convergence
of the solutions obtained above, not an assumed stability
under weak convergence of their right-hand sides.
The bounds pass to the uniform limit. Pointwise ambient
uniqueness, with the two boundary compact-gap conditions,
gives uniqueness among all bounded classical solutions.

Now let arbitrary $\mu_j\to\mu$ on $K$. Their total masses
are bounded. The solutions just constructed satisfy the same
uniform bounds, boundary barrier, and skeletal equicontinuity.
From any subsequence, the preceding compactness and support
argument extracts a further subsequence converging uniformly
on $V$. Uniform BT stability identifies its measure with
$\mu$, and uniqueness identifies its limit with $u$.
Thus the entire sequence converges uniformly.

Every bounded classical function is continuous on $K$, by
Lemma~\ref{lem:r3-gauss-box-equicontinuity}. For a fixed
continuous compact scalar $\chi$, write
\[
 \int\chi u_j\,d\mu_j-\int\chi u\,d\mu
  =\int_K\chi(u_j-u)\,d\mu_j
       +\int_K\chi u\,d(\mu_j-\mu).
\]
The first term tends to zero by the uniform convergence and
mass bound; the second tends to zero because $\chi u$ is
continuous on $K$. This proves the stated weighted limit
with the actual representatives. Adding the constants $b_j$
proves all assertions for varying constant boundary values.

Under the additional extendible-exhaustion hypotheses, every
finite atomic approximant extends continuously to $\overline V$,
with value zero on the boundary. Uniform convergence on $V$
is also uniform on $\overline V$ after those extensions.
Their limit is therefore continuous on $\overline V$.
\end{proof}

For fixed $M$, the solutions with $\mu\geq0$, support in $K$,
and $\mu(K)\leq M$ form a compact set for the uniform norm:
the finite positive measures of bounded mass on compact $K$
are weakly compact, and the theorem gives a continuous
solution map. This is compactness of this prescribed-measure
class, not compactness of every uniformly bounded psh family.

\begin{example}[Non-atomic data in every dimension]
\label{ex:r3-skeletal-nonatomic}
Let $R>1$ and fix
\[
 V=\{\max_i|T_i|<R\}\subset(\mathbf A_K^n)^{\an},
 \qquad
 \rho=\log\max(1,|T_1|,\ldots,|T_n|)-\log R.
\]
This is an extendible-exhaustion domain. Fix finite
$a<b<\log R$ and let $K=\sigma([a,b]^n)$.
Every finite positive Radon measure on this $K$ has a
unique bounded classical Dirichlet solution with any
constant boundary value, continuous on $\overline V$.
For instance
\[
 \mu=\sigma_*(f(x)\,dx),\qquad
       0\leq f\in L^1([a,b]^n),
\]
is non-atomic and nonzero when $f\not\equiv0$ in $L^1$.
No equation on a new domain adapted to its mass distribution
has been used. The local testers in
\eqref{eq:r3-gauss-box-test-atom} are already bounded on
this whole $V$, so with $A=n(\log R-a)$ one can use
$c_K=n!/A^n$ in the uniform norm bound.
\end{example}

\paragraph{The remaining diffuse-data scope.}
The argument has not proved existence for every compactly
supported measure, or for every measure dominated by
$(Dw)^n$ for a continuous bounded classical $w$.
Outside the stated Gauss-box geometry, the uniform skeletal
continuity needed for Arzel\`a--Ascoli has not been supplied.
The earlier moving-branch counterexample rules out replacing
that step by unrestricted compactness of bounded psh families.
Nor may an atomization of a non-atomic measure be assumed to
remain dominated by that measure: a positive point mass
already contradicts such domination on its singleton.
The present proof obtains its testing continuity and uniform
compactness from the specified support geometry instead of
assuming a general variational compactness theorem.

% END INCLUDED SOURCE: sections/skeletal_dirichlet


% BEGIN INCLUDED SOURCE: sections/moving_atomic_dirichlet
% Reconstructed after the scratch reset. Original reviewed SHA256:
% cc54ba0d930f0684881340e6964306db5eea9e0d9e1b9846b3d47a3bd9c4e55e.
\section{Moving atoms and failure of weak-data Dirichlet stability}
\label{sec:r3-moving-atomic-dirichlet}

Let $\pi$ be a uniformizer, $r=|\pi|$, $\lambda=-\log r$,
$R=r^{-1}$, and
$\Omega=\{|T_l|<R:1\leq l\leq n\}$.
Choose pairwise distinct monic irreducibles $\overline P_i$ of
degrees $d_i\geq2$ over the residue field, and monic integral
lifts $P_i$. Let $\xi_i$ be the Shilov point of
$\{|P_i(T_1)|\leq r^{d_i}\}$ in the original analytic affine line.
It is a single original $K$-analytic point of residual degree
$d_i$. After the unramified splitting extension, it consists of
the conjugate Gauss points of discs centered at the roots of $P_i$
and of radius $r^{d_i}$.

Let $z_i$ be the product Gauss point with first factor $\xi_i$
and remaining coordinate radii $r$, and let $z$ be the product
Gauss point with first radius one and remaining radii $r$.
Let $y$ have all coordinate radii one. The distinct residue branches
give $\xi_i\to\zeta_{0,1}$, hence $z_i\to z$.
For completeness, a fixed polynomial has only finitely many
irreducible residue factors; its norm at $\xi_i$ eventually equals
its unit Gauss norm. Uniform polynomial approximation on compact
polydiscs gives the asserted convergence for the product points.

Define
\begin{equation}
 g_i=\max\left\{d_i^{-1}\log|P_i(T_1)|-\lambda,\,
       \log|T_2|-\lambda,\ldots,\log|T_n|-\lambda,\,-2\lambda\right\}.
 \label{eq:r3-moving-atomic-potentials}
\end{equation}
The list of remaining coordinates is empty when $n=1$.

\begin{theorem}[Explicit zero-boundary weak instability]
\label{thm:r3-moving-atomic-dirichlet}
The functions $g_i$ are continuous classical psh on $\overline\Omega$,
take values in $[-2\lambda,0]$, have ordinary boundary value zero,
and satisfy
\[
 (Dg_i)^n=\delta_{z_i}.
\]
They converge pointwise, indeed eventually at each point, to
\[
 g=\max\{0,\log|T_1|,\ldots,\log|T_n|\}-\lambda,\qquad
 (Dg)^n=\delta_y.
\]
Put $B=2^{(n-1)/n}\lambda$ and
\[
 H=\max\left\{
     \frac B\lambda(\log|T_1|-\lambda),\,
     \frac B{2\lambda}(\log|T_l|-\lambda)\ (2\leq l\leq n),\,
     -B\right\}.
\]
Then $H$ is the continuous zero-boundary solution of
$(DH)^n=\delta_z$, but
\[
 \|g_i-H\|_\infty\geq2\lambda-B>0.
\]
For $n\geq2$, $z\ne y$ and $g(z)=-\lambda\ne-B=H(z)$.
Thus weak convergence of the data on a fixed compact support
does not imply uniform or capacity convergence of the solutions;
in dimension at least two it does not even imply pointwise
convergence to the solution of the limit datum.
\end{theorem}

\noindent\textbf{Status: \texttt{NEW\_PROVED}} counterexample;
the corresponding unrestricted weak-data stability assertion is
\texttt{FALSE}.

\begin{proof}
Finite maxima of logarithms with a finite floor are continuous
classical after every complete extension. At an ordinary boundary
point, some $|T_l|=R$. If $l=1$, monicity and integrality give
$|P_i(T_1)|=|T_1|^{d_i}$; otherwise the displayed $l$th term
is zero. All other terms are nonpositive on the closed polydisc.
This proves the boundary and boundedness assertions.

We compute the atom over the original field. The separable
irreducible reduction splits in the unramified extension of
degree $d_i$. Its roots are units, and distinct roots have
distance one. Near a conjugate root $\alpha$,
$\log|P_i(T_1)|=\log|T_1-\alpha|$ on the corresponding open
residue disc, since every other factor is a unit of norm one.
The tie of all terms in \eqref{eq:r3-moving-atomic-potentials}
occurs at radii $r^{d_i},r,\ldots,r$. The slope simplex has
volume $1/(d_i n!)$. Original torus integration contributes
$n!$, and passage to the original degree-$d_i$ point contributes
the factor $d_i$. Its original mass is exactly one.

There is a useful independent check that no other mass was
omitted. On $\mathbf P_K^n$, homogenize the sections
\[
 \pi^{d_i}P_i(T_1),\quad
 \pi^{d_i}T_l^{d_i}\ (2\leq l\leq n),\quad \pi^{2d_i}.
\]
They are base-point free sections of $\mathcal O(d_i)$:
on the affine chart the last is nonzero; at infinity their
leading terms are the coordinate $d_i$th powers. Their pure-FS
weight divided by $d_i$ is the displayed $g_i$ in this frame.
Its positive ordinary top measure has total mass
$\deg(\mathcal O(1)^n)=1$, by the established ordinary-current
model compatibility. The original atom already has mass one.
This proves $(Dg_i)^n=\delta_{z_i}$ on all original tests,
without assuming a pullback formula for an arbitrary delta form.

If $|T_1|>1$, monicity gives $g_i=g$. If $|T_1|\leq1$, at most
one of the distinct residue polynomials can have norm below one.
For every other index the first logarithmic term is $-\lambda$,
again giving $g_i=g$. Thus the equality is eventual at every
point, and always
\[
 0\leq g-g_i\leq\lambda.
\]
The slope simplex for $g$ at $y$ is the standard unit simplex;
elsewhere the active slope differences have rank below $n$.
Original tropical calibration gives $(Dg)^n=\delta_y$.
For $H$, the floor ties precisely at $z$. The slopes are
$B/\lambda$ in the first coordinate and $B/(2\lambda)$ in
the others, so its top mass is
\[
 \frac{B^n}{2^{n-1}\lambda^n}=1.
\]
The same rank calculation excludes further mass. Real slopes
are justified by smooth convex regularization and the proved
bounded continuity, so no rationality of $B/\lambda$ is needed.
The ordinary boundary values are zero. Actual zero-boundary
uniqueness identifies $g_i$ and $H$ as the respective solutions.
At $z_i$, one has $g_i=-2\lambda$ and $H=-B$, proving the
uniform lower bound.

Finally set
$K_*=\{z,z_1,z_2,\ldots\}$ and
$\Omega_0=\{|T_l|<r^{-1/2}\text{ for all }l\}\Subset\Omega$.
This is a fixed compact set in $\Omega_0$. The tester
$1+g_i/(2\lambda)$, restricted to $\Omega_0$, is normalized and
has mass $(2\lambda)^{-n}$ at $z_i$. Therefore any strict error
set on $K_*$ containing $z_i$ has outer capacity at least this
constant on $\Omega_0$. Taking an error threshold less than
$2\lambda-B$ rules out local capacity convergence to $H$.
\end{proof}

\begin{corollary}[Convex blocks do not preserve the original weak limit]
\label{cor:r3-moving-convex-block-mismatch}
For $n\geq2$ the example has tail convex blocks converging
uniformly to $g$, whose top measures converge to $\delta_y$,
whereas the unaveraged top measures converge weakly to
$\delta_z\ne\delta_y$.
\end{corollary}

\begin{proof}
The open residue branches
$O_i=\{|P_i(T_1)|<1\}\cap\Omega$ are pairwise disjoint.
The difference $g-g_i$ vanishes off $O_i$ and is at most
$\lambda$. Thus for any $N$ distinct indices,
\[
 0\leq g-\frac1N\sum g_i\leq\frac{\lambda}{N}
\]
uniformly. Uniform BT continuity gives the measure limit
$\delta_y$. The original measures have weak limit $\delta_z$.
This explicitly excludes deducing a nonlinear unaveraged
weak-limit representation solely from capacity compactness
of convex blocks.
\end{proof}

\paragraph{Additional scope checks.}
The weighted measures $g_i(Dg_i)^n$ tend to $-2\lambda\delta_z$.
For distinct $i,j$, $g_j(z_i)=-\lambda$, so the same fixed
compact capacity testers rule out a capacity-Cauchy subsequence.
The compact $K_*$ is not contained in finitely many fixed full
Gauss boxes on which all these bounded potentials have a common
interior Lipschitz estimate: such equicontinuity would contradict
$z_i\to z$ and $g_i(z_i)-g_i(z)=-\lambda$.
There is no contradiction with sharp weak-closure capacity
domination, since $\delta_z$ is itself the measure of $H$.
Nor is there a common locally bounded Monge--Ampere subsolution
dominating all $\delta_{z_i}$: its locally finite measure would
have mass at least one at infinitely many points of $K_*$.

% END INCLUDED SOURCE: sections/moving_atomic_dirichlet


% BEGIN INCLUDED SOURCE: sections/bounded_subsolution
\section{A bounded-subsolution theorem on extendible domains}
\label{sec:r3-bounded-subsolution}

The final Dirichlet theorem in this section has ambient background
$[1]$.  Its preliminary uniform weighted estimate holds over every
fixed allowed positive closed background.  The topology of the data
is stronger than ordinary weak convergence, which is insufficient
even for the explicit moving-atom examples.

We keep the original finite field extension $K/\mathbf Q_p$, its
absolute value, and all original CLD compact smooth tests.  Positive
backgrounds are symmetric, original-positive, and separately
$d'$- and $d''$-closed.  Actual bounded classical psh functions are
used throughout.  Products and energies mean the already established
fixed-background BT products, not wedges of arbitrary currents.

\subsection{Uniformity over bounded psh weights}

For finite signed Radon measures carried by a fixed compact subset
of an ordinary domain $O$, define
\[
 \|\sigma\|_{\mathrm{psh},O}
   =\sup_{\substack{a\in\mathcal C_b(O)\\0\leq a\leq1}}
                 \left|\int a\,d\sigma\right|,
 \qquad
 d_{\mathrm{psh},O}(\mu,\nu)=\|\mu-\nu\|_{\mathrm{psh},O}.
\]
The bound $0\leq a\leq1$ is on the same full domain $O$.

\begin{theorem}[Uniform psh-weight convergence for decreasing rows]
\label{thm:r3-psh-dual-monotone}
Let $S$ be a fixed positive background of bidegree $(p,p)$ on an
allowed ordinary domain $O$, and put $q=n-p\geq1$.  Suppose
$u_j\in\mathcal C_b(O)\cap C^0(O)$ decrease to
$u\in\mathcal C_b(O)$.  If $f$ is original smooth with compact
support in $O$ and $0\leq f\leq1$, then
\[
 \sup_{\substack{a\in\mathcal C_b(O)\\0\leq a\leq1}}
 \left|\int af\bigl(T_{S,q}(u_j^q)-T_{S,q}(u^q)\bigr)\right|
       \longrightarrow0.
\]
Here $u_j^q$ denotes $q$ repeated slots, not a scalar power.
All constants used in the proof depend only on fixed compact
stages, local bounds, and finitely many positive evaluations of $S$.
\end{theorem}

\noindent\textbf{Status: \texttt{NEW\_PROVED}.}

\begin{proof}
Split $f$ by a finite original smooth nonnegative partition subordinate
to smaller norm boxes covering its support.  It suffices to prove
the assertion for one summand, still denoted $f$, supported strictly
inside one such box.  Choose an original nonnegative compact cutoff
$\chi$ equal to one near its support.  All potentials are uniformly
bounded on a fixed larger compact stage.  For $j\leq k$ put
\[
 g_{jk}=u_j-u_k,\quad
 Q_{jk}=\sum_{r=0}^{q-1}(Du_j)^r(Du_k)^{q-1-r}S,\quad
 \mu_j=(Du_j)^qS.
\]
Every summand of $Q_{jk}$ is a positive separately closed current
of bidegree $(n-1,n-1)$, and the proved mixed calculus gives
\[
 D(g_{jk}Q_{jk})=\mu_j-\mu_k.
\]
Let
\[
 e_{jk}=\int\chi\,\Gamma_{Q_{jk}}(g_{jk})\geq0.
\]
For an original smooth psh tester $a$ with $0\leq a\leq2$ on
the fixed stage, compact Green and the product rule give
\[
 \int af(\mu_j-\mu_k)
   =-\int f\,\Gamma_{Q_{jk}}(a,g_{jk})
     -\int a\,\Gamma_{Q_{jk}}(f,g_{jk}).
\]
The pairing involving the original smooth $f$ has its usual smooth
first factor; no differentiated arbitrary Borel function is used.
Cauchy--Schwarz gives
\begin{equation}
 \left|\int af(\mu_j-\mu_k)\right|\leq A e_{jk}^{1/2},
 \label{eq:r3-uniform-weight-energy}
\end{equation}
with $A$ independent of $a,j,k$.  Indeed the required factors other
than $e_{jk}$ are bounded by
\[
 \int\chi\,\Gamma_{Q_{jk}}(a),\qquad
 \int\chi\,\Gamma_{Q_{jk}}(f).
\]
For the first, the cutoff identity is
\[
 \int\chi\,\Gamma_{Q_{jk}}(a)
 =\tfrac12\langle Q_{jk},a^2D\chi\rangle
       -\int\chi aD(aQ_{jk})
 \leq\tfrac12\left|\langle Q_{jk},a^2D\chi\rangle\right|.
\]
The subtracted measure is nonnegative.  The finite-test CLN estimate
bounds the last expression uniformly, using $a^2\leq4$ and the
fixed local bounds for the slots of $Q_{jk}$.  The second factor
is bounded by the same CLN theorem against the fixed original form
$\chi\,d'f\wedge d''f$.  The product rule and Cauchy--Schwarz can
first be applied to smooth approximants of $u_j,u_k$ and then passed
to their limits by functional weighted monotone convergence.  Every
background occurring in this passage is an explicit mixed product
over the same $S$.

For an actual $a\in\mathcal C_b(O)$ with $0\leq a\leq1$, take a
smooth decreasing majorizing row on the prescribed box.  After a
prefix it is between zero and two on the compact stage.  For fixed
$j,k$, dominated convergence against the fixed finite measures
$f\mu_j,f\mu_k$ passes \eqref{eq:r3-uniform-weight-energy} to $a$.
The bound $A$ has not changed.  This avoids introducing a first-order
product with a simultaneously varying discontinuous tester.

The energy has the exact scalar expression
\begin{equation}
 e_{jk}
   =\tfrac12\langle Q_{jk},g_{jk}^2D\chi\rangle
       -\int\chi g_{jk}(\mu_j-\mu_k).
 \label{eq:r3-dual-energy-expansion}
\end{equation}
Fix $j$ and let $k\to\infty$.  Functional weighted joint monotone
convergence applies to every term after expanding the scalar
polynomials in $u_j,u_k$.  It gives a finite nonnegative limit
$e_j$.  Letting $j\to\infty$ in that resulting expression gives
$e_j\to0$, because its scalar difference is now $u_j-u\to0$ and
the same functional weighted theorem applies to the mixed tuples
$(u_j,u)$.  The constant row $u$ is an allowed decreasing row.

For each fixed bounded psh tester $a$, weighted monotone continuity
gives $a\mu_k\to a\mu$, where $\mu=(Du)^qS$, against the test $f$.
Pass $k\to\infty$ in \eqref{eq:r3-uniform-weight-energy}; then
\[
 \left|\int af(\mu_j-\mu)\right|\leq A e_j^{1/2}
 \quad\text{for every normalized }a.
\]
Take the supremum only after this fixed-tester limit.  The right
side tends to zero.  Summing the finitely many initial localized
estimates proves the assertion on $O$.  No capacity convergence
was used to prove the uniform weighted estimate.
\end{proof}

\begin{lemma}[The psh dual distance detects weak convergence]
\label{lem:r3-dual-determines-weak}
Let $V$ be relatively compact in a smooth affine algebraic analytification
and fix a compact $L\subset V$.  On finite positive Radon measures
carried by $L$, $d_{\mathrm{psh},V}$ is a metric.  Convergence in this
metric implies ordinary weak convergence of these measures.
\end{lemma}

\begin{proof}
Constants are testers, so convergence controls total masses.  The
differences of bounded continuous classical psh functions on $V$
form a real vector lattice: for $a,b,c,d$ in that cone,
\[
 \max(a-b,c-d)=\max(a+d,c+b)-(b+d).
\]
This lattice contains constants and separates the points of $L$.
For the separation assertion, an affine polynomial $P$ whose
absolute values differ at two Berkovich points yields, for a suitable
finite $c$, the bounded continuous classical function
$\max(\log|P|,c)$, still separating them.  It is bounded above on
$\overline V$.  Affine polynomial seminorms separate the points of
the affine analytification.  The real lattice Stone--Weierstrass
theorem therefore gives uniform density of this lattice in $C(L)$.
Every bounded psh function in its definition can be shifted and
positively rescaled into $[0,1]$ on $V$.  The metric consequently
controls each fixed lattice test.  Uniform approximation and the
bounded masses give convergence against every continuous test on
$L$.  The same argument with distance zero proves equality of the
measures.  Symmetry and the triangle inequality are immediate.
\end{proof}

\subsection{The domain and the common exterior}

For the remainder of the section, $Y$ is a smooth geometrically
integral affine $K$-variety and $\rho$ is continuous classical psh
on all of $Y^{\an}$.  Let $V$ be a connected ordinary component
of $\{\rho<0\}$ with compact closure in $Y^{\an}$, nonempty ordinary
boundary, and $\rho=0$ on that boundary.  This is the extendible
exhaustion-domain hypothesis.  No strict psh exhaustion is assumed.
Write $\operatorname{MA}(v)=(Dv)^n$ on $V$.

\begin{lemma}[Weak boundary order for a subsolution]
\label{lem:r3-weak-boundary-subsolution-order}
Let $a,b\in\mathcal C_b(V)$ be bounded.  Suppose $b$ has ordinary
boundary limit zero, $\limsup_{x\to\partial V}a(x)\leq0$, and
$\operatorname{MA}(a)\geq\operatorname{MA}(b)$.  Then $a\leq b$
at every actual point of $V$.
\end{lemma}

\begin{proof}
The boundary maximum principle gives $a,b\leq0$.  Fix $t,\delta>0$
and put $w=(1+t)a-\delta$.  The strict gap set $E=\{b<w\}$ is
relatively compact: the boundary limsup of $w-b$ is at most
$-\delta$, and the boundary is compact.  It need not be ordinary
open.  The proved comparison theorem gives
\[
 (1+t)^n\operatorname{MA}(b)(E)
 \leq\operatorname{MA}(w)(E)
 \leq\operatorname{MA}(b)(E).
\]
These masses are finite, so $\operatorname{MA}(b)(E)=0$.
Pointwise ambient domination with the strict boundary gap yields
$w\leq b$.  Let $\delta\downarrow0$ and then $t\downarrow0$.
\end{proof}

\begin{lemma}[A common lower barrier and exterior]
\label{lem:r3-common-exterior-extension}
Fix a compact $L\Subset V$ and $B>0$.  Every bounded classical
zero-boundary solution $v$ with
\[
 -B\leq v\leq0,\qquad\operatorname{MA}(v)\text{ carried by }L
\]
satisfies $v\geq C\rho$ on $V$, where $C$ depends only on
$B,L,\rho$.  It has a classical extension $\widetilde v$ to
$Y^{\an}$, equal to $v$ on $V$ and to $C\rho$ on $Y^{\an}\setminus V$.
For the whole family one can make a common prescribed-box collar
extension $W_v$ such that, on a fixed neighborhood of $\overline V$,
$W_v=\widetilde v+c$ for one constant $c$, and all $W_v$ equal one
fixed smooth classical collar outside one larger compact set.
In particular $W_v-W_{v'}$ vanishes at every point outside $V$,
including every point of $\partial V$.
\end{lemma}

\begin{proof}
If $L$ is empty, domination gives $v=0$ and the assertion is
immediate.  Otherwise choose $C$ with
$C\min_L(-\rho)\geq B$.  The function
$b=\max(-B,C\rho)$ has zero boundary value and satisfies $b\leq v$
on $L$.  The measure $\operatorname{MA}(v)$ therefore vanishes on
$\{v<b\}$.  Pointwise domination gives $b\leq v$, and in particular
$C\rho\leq v$.

For $\varepsilon>0$, use $\max(v,C\rho+\varepsilon)$ inside $V$
and $C\rho+\varepsilon$ outside it.  Near every boundary point the
second branch is active, by the ordinary boundary limit of $v$.
The ordinary gluing theorem makes this a global classical function.
Let $\varepsilon\downarrow0$.  Decreasing closure gives precisely
the claimed actual extension, including its value zero on the
boundary.  The family has uniform bounds on every fixed compact
neighborhood of $\overline V$: inside $V$ use $B$, and outside it
all functions equal $C\rho$.

Apply the constructive common-collar part of the prescribed-box
regularization lemma (the B extension step) to this one uniformly
bounded family.  Choose the inner compact stage to contain
$\overline V$ in its interior and the outer barrier above the
common bound on the outer collar.  The same shift and maximum
construction preserve every $\widetilde v+c$ near $\overline V$
and give the same smooth outer collar.  Outside $V$ all inputs
were already identical, so the same pointwise construction keeps
them identical there.  This proves the final assertion.  The
algebraic compactification used for this collar does not change
the field or the original current test category.
\end{proof}

\begin{theorem}[Compact-support stability in the psh dual metric]
\label{thm:r3-psh-dual-dirichlet-stability}
On the specified extendible $V$, let $v_j$ be bounded actual
classical zero-boundary solutions, with one common bound
$-B\leq v_j\leq0$.  Suppose all $\mu_j=\operatorname{MA}(v_j)$
are carried by one compact $L\Subset V$ and form a Cauchy sequence
for $d_{\mathrm{psh},V}$.  Then there is a unique bounded actual
classical zero-boundary $v$ such that
\[
 v_j\longrightarrow v\quad\text{in all-domain local outer ambient
 capacity},\qquad
 d_{\mathrm{psh},V}(\mu_j,\operatorname{MA}(v))\longrightarrow0.
\]
Its measure is carried by $L$.  In particular, if
$d_{\mathrm{psh},V}(\mu_j,\mu)\to0$ for a finite positive measure
$\mu$ carried by $L$, then $\operatorname{MA}(v)=\mu$.
\end{theorem}

\noindent\textbf{Status: \texttt{NEW\_PROVED}.}

\begin{proof}
The case $B=0$ is immediate.  Use the common exterior and collar
from Lemma~\ref{lem:r3-common-exterior-extension}.  All $W_j=W_{v_j}$
are uniformly bounded on one fixed collar domain $\Omega$.  Their
differences vanish outside $V$, including its boundary.  Hence the
positive mixed energy has the exact expression
\[
 \begin{split}
 I_{[1]}(W_j,W_k)
 &=\int_V(v_j-v_k)(\mu_k-\mu_j)\\
 &\leq2B\,d_{\mathrm{psh},V}(\mu_j,\mu_k).
 \end{split}
\]
For the bound, write $v_j-v_k=B(a_j-a_k)$ with
$a_j=1+v_j/B$, so both $a_j,a_k$ are normalized psh testers on $V$.
No equality of the extension measures on $\partial V$ is asserted:
their actual multiplier $W_j-W_k$ is zero there, which is the
reason no boundary mass occurs in the identity.

The all-degree estimate of Theorem~\ref{thm:r3-I-capacity} makes
$W_j$ Cauchy in outer ambient capacity on $\Omega$.  Domain
comparison gives all-domain local convergence and the completed
calculus supplies a limit in $\mathcal A_{[1],b}$.  The proved
ambient representative theorem supplies its unique actual classical
representative.  More explicitly, take a summably capacity-convergent
subsequence and use its decreasing usc-regularized tail suprema.
Near $\overline V$ the continuous barriers
\[
 C\rho+c\leq W_j\leq\max(0,C\rho)+c
\]
pass through this construction at actual points.  Subtracting $c$
on $V$ gives a bounded classical representative $v$ satisfying
$C\rho\leq v\leq0$.  Thus it has ordinary boundary limit zero.
The Cauchy property identifies the same limit for the full sequence.

The completed weighted calculus gives weak convergence of
$\operatorname{MA}(v_j)$ to $\operatorname{MA}(v)$ on every original
compact test in $V$.  These measures are supported in $L$, so the
limit is carried by $L$ and is finite.  For every fixed normalized
actual psh tester $a$ on $V$, insert an original smooth cutoff equal
to one near $L$.  Weighted capacity continuity then gives
\[
 \int a\,d\mu_k\longrightarrow
              \int a\,d\operatorname{MA}(v).
\]
For fixed $j$, pass $k\to\infty$ in the defining dual inequality,
and only afterwards take the supremum in $a$.  This yields
\[
 d_{\mathrm{psh},V}(\mu_j,\operatorname{MA}(v))
       \leq\limsup_{k\to\infty}d_{\mathrm{psh},V}(\mu_j,\mu_k)
       \longrightarrow0.
\]
This proves strong convergence without assuming in advance that
the measure metric is complete.  It identifies any prescribed
metric limit by Lemma~\ref{lem:r3-dual-determines-weak}.
Uniqueness of the actual zero-boundary solution is the proved
ambient comparison and domination theorem.
\end{proof}

\subsection{Strong atomic approximation from a subsolution}

\begin{theorem}[The bounded-subsolution criterion]
\label{thm:r3-bounded-subsolution-dirichlet}
On the specified extendible domain $V$, let $\mu$ be a finite
positive Radon measure with compact support in $V$.  The following
are equivalent:
\begin{enumerate}
\item There is a bounded actual classical $u$ on $V$ with ordinary
boundary limit zero and $(Du)^n=\mu$.
\item There is a bounded actual classical $w$ on $V$ with
$(Dw)^n\geq\mu$.
\end{enumerate}
The solution in the first assertion is unique.  No boundary
condition is required of the subsolution in the second assertion.
This theorem asserts a bounded classical solution; it does not
assert interior continuity from the bounded-subsolution hypothesis
alone.
\end{theorem}

\noindent\textbf{Status: \texttt{NEW\_PROVED}.}

\begin{proof}
Necessity is immediate.  For sufficiency the zero measure has
solution zero, so suppose $\mu\ne0$, and let $K=\operatorname{supp}\mu$.
Subtract a constant from $w$ so that $-M\leq w\leq0$.  Choose
$C$ so large that $C\rho<-M-1$ on an ordinary neighborhood of $K$.
Then
\[
 b=\max(w,C\rho)
\]
is bounded classical, equals $w$ on that neighborhood, and has
ordinary boundary value zero.  Ordinary locality and positivity
give $(Db)^n\geq\mu$ on all of $V$: on the neighborhood of $K$
this follows from the subsolution, and outside $K$ the datum is
zero.  This uses no presumed ordinary openness of a comparison
set involving $w$; the displayed strict constant bound supplies
the required open neighborhood.

The same epsilon-maximum construction as in the preceding lemma
extends $b$ to a global classical function $\widetilde b$, equal
to $C\rho$ off $V$.  We now use the accepted constructive input
from B, with the entire compact $\overline V$ protected.  Extend
the local potential to one semipositive metric on a projective
algebraic model in a fixed frame, and take its decreasing pure
Fubini--Study rows.  In that frame, after the common constant or
unit-log correction, their local potentials $b_j$ are continuous
classical on a fixed neighborhood of $\overline V$ and satisfy
\[
 b_j\downarrow\widetilde b.
\]
Their ordinary top-degree measures are finite positive sums of
Dirac masses at divisorial points of the original $K$-analytic
space.  This last fact is the accepted semipositive PL/model
measure compatibility applied to these particular algebraic
pure-Fubini--Study rows.  It is not a claim that every locally
defined PL function has globally finite atomic curvature.

On the compact boundary $\partial V$, the continuous functions
$b_j$ decrease to the continuous zero function.  Dini's theorem
on that compact set gives
\[
 \varepsilon_j=\max_{\partial V}b_j\downarrow0.
\]
The maximum principle gives $b_j\leq\varepsilon_j$ on $V$.
Also $b_j\geq b$ there.  Only the boundary restriction is used in
this Dini argument; continuity of $b$ in the interior is not assumed.

Put $\nu=(Db)^n$.  Since $\mu\leq\nu$, Radon--Nikodym gives
$\mu=f\nu$ for a Borel $0\leq f\leq1$, chosen zero outside $K$.
Choose a compact $L\Subset V$ with $K\subset L^{\circ}$.
There exist original smooth $f_l$ with common compact support in
$L$, $0\leq f_l\leq1$, and
\[
 \int|f_l-f|\,d\nu\longrightarrow0.
\]
Here only the finite measure on a compact collar is involved.
Radon regularity first approximates a bounded Borel function by
continuous compactly supported functions; original smooth
partitions and smooth approximation on finitely many moment
presentations, followed by a smooth truncation into $[0,1]$ and
a fixed support cutoff, give the asserted original smooth functions.
Thus neither a changed test sheaf nor an infinite-volume density
assertion is needed.

For each fixed $l$, Theorem~\ref{thm:r3-psh-dual-monotone}, with
$S=[1]$, gives
\[
 d_{\mathrm{psh},V}\bigl(f_l(Db_j)^n,f_l\nu\bigr)
                  \longrightarrow0\quad(j\to\infty).
\]
Choose strictly increasing $j(l)$ so that this distance is at most
$2^{-l}$, and set
\[
 \mu_l=f_l(Db_{j(l)})^n.
\]
These are finite positive divisorial atomic measures carried by
the same compact $L$, and
\[
 d_{\mathrm{psh},V}(\mu_l,\mu)
    \leq2^{-l}+\int|f_l-f|\,d\nu\longrightarrow0.
\]
This is the needed strong approximation, not merely weak
approximation by atomic measures.  Also
$\mu_l\leq(Db_{j(l)})^n$ on all of $V$.

The finite-divisorial-atom Dirichlet theorem on this same
exhaustion domain gives bounded actual zero-boundary classical
solutions $v_l$ with $\operatorname{MA}(v_l)=\mu_l$.
Apply Lemma~\ref{lem:r3-weak-boundary-subsolution-order} with
$a=b_{j(l)}-\varepsilon_{j(l)}$ and $b=v_l$.  Its hypotheses hold
by the preceding measure domination and boundary maximum.  Hence
\[
 0\geq v_l\geq b_{j(l)}-\varepsilon_{j(l)}
              \geq b-\varepsilon_1.
\]
This is one common global bound, independent of the atoms and
their masses.  Theorem~\ref{thm:r3-psh-dual-dirichlet-stability}
now gives a bounded actual classical zero-boundary $u$ with
$\operatorname{MA}(u)=\mu$.  It also supplies capacity convergence
of the constructed solutions.  Actual pointwise uniqueness follows
from the ambient domination theorem.
\end{proof}

\paragraph{Scope of stability and of the algebraic input.}
The compact support, common uniform bound, and extendible exhaustion
are used explicitly in the common-exterior energy argument.  The
construction retains the original field and invokes B only for its
accepted algebraic collar and pure-Fubini--Study approximation, with
the already known original-current/model-measure comparison.  It
does not derive arbitrary-background Dirichlet existence, ordinary
weak-data stability, or continuity under merely bounded subsolutions.
The preliminary uniform weighted theorem, in contrast, is valid for
arbitrary fixed $S$ and all its permitted remaining degrees.

For the moving-atom measures $\delta_{z_j}\rightharpoonup\delta_z$
of the explicit counterexample, the normalized tester
$1+g_j/(2\lambda)$ has values zero at $z_j$ and $1/2$ at $z$.
Consequently their psh dual distances are at least $1/2$.
The strong topology used here therefore excludes that actual weak
convergence defect.

% END INCLUDED SOURCE: sections/bounded_subsolution


% BEGIN INCLUDED SOURCE: sections/capacity_dirichlet
\section{The fixed-domain Dirichlet problem under power capacity growth}
\label{sec:r3-capacity-dirichlet}

Throughout this section $K/\mathbf Q_p$ is the prescribed finite
extension with its original absolute value. Let $Y/K$ be smooth,
geometrically integral and affine, of dimension $n\geq1$. Let
$\rho$ be a continuous classical psh function on all of $Y^{\an}$.
Fix a connected component $V$ of $\{\rho<0\}$ such that
\[
 V\Subset Y^{\an},\qquad
 \varnothing\ne\partial V\subset\{\rho=0\}.
\]
All boundaries are ordinary boundaries. Write
$C(E)=\operatorname{Cap}_{[1]}(E;V)$. The background in this
Dirichlet section is $[1]$, not an arbitrary relative current.
The currents and their convergence use all original CLD tests;
the classical functions retain actual values and all extension tests.

\begin{lemma}[An affine projective collar for a fixed lower obstacle]
\label{lem:r3-capdir-affine-collar}
Fix $B>0$ and put
\[
 \mathcal F_B=\{u\in\Cb(V):B\rho\leq u\leq0\}.
\]
There is a smooth projective compactification $Z$ of $Y$, an ample
line $L$ on $Z$, a fixed frame on an ordinary neighborhood of
$\overline V$, and an affine order-preserving map
$u\mapsto U$ from $\mathcal F_B$ to bounded classical metrics
on $L$ with the following properties. In that frame the weight
of $U$ is $u+c$ on $V$, with one constant $c$ independent of $u$.
All the metrics are equal on $Z^{\an}\setminus V$. They have
common bounds relative to a continuous reference metric, and the
metrics $F$ and $L_B$ associated with $0$ and $B\rho$ are
continuous. One has $L_B\leq U\leq F$.

Conversely, any bounded classical metric $P$ with
$L_B\leq P\leq F$ and $P=F$ outside $V$ corresponds in this
frame to a member of $\mathcal F_B$. The affine family has energy
\[
 \mathcal E(U)=\frac1{n+1}\sum_{j=0}^n
  \int_{Z^{\an}}(U-F)c_1(L,U)^j c_1(L,F)^{n-j}.
\]
It satisfies the mixed energy identities, midpoint inequality and
local energy-to-capacity estimates already proved.
\end{lemma}

\begin{proof}
For $u\in\mathcal F_B$ and $\epsilon>0$, glue
$\max(u,B\rho+\epsilon)$ on $V$ to $B\rho+\epsilon$ outside
$V$. Near every ordinary boundary point the latter function is
positive while $u\leq0$, so these prescriptions agree on an
ordinary neighborhood. The glued functions are classical and
decrease locally boundedly as $\epsilon\downarrow0$ to
\[
 \widetilde u=u\text{ on }V,\qquad
 \widetilde u=B\rho\text{ off }V.
\]
The lower bound guarantees the actual first equality. This gives
a classical function on $Y^{\an}$. The map $u\mapsto\widetilde u$
is affine and order preserving. Its upper member is the continuous
function equal to zero on $V$ and to $B\rho$ off $V$; its lower
member is the continuous function $B\rho$.

Use the framed smooth compactification
Lemma~\ref{lem:r3-framed-compactification}, with one frame whose
nonvanishing locus contains all of $Y$. Choose nested compact
norm boxes in that frame, the innermost containing $\overline V$
in its ordinary interior. On the outer box the family
$\widetilde u$ has one common bound. The constructive collar of
B, Lemma D.2, in the form used in
Theorem~\ref{thm:r2-local-upper-envelope}, gives on one ample
power the weights
\[
 U=\max(\widetilde u+c,mb)
\]
on the outer box, and a fixed metric outside the middle box.
Here $b=0$ on the inner box, $c$ makes every
$\widetilde u+c$ positive there, and the integer $m$ makes the
fixed branch strictly dominate on the outer collar. Since the
inputs all agree outside $V$, the outputs all agree there.
On $V$ the maximum is inactive and its weight is $u+c$.
Thus this particular family remains affine, even though the
collar operation on arbitrary inputs is a maximum. The upper
and lower members are continuous, and ordinary restriction
proves the converse assertion.

Global mixed Green exchange gives the stated energy identities;
differences of the metrics vanish outside the fixed compact
$\overline V$. Equivalently one can compute in this frame on a
larger norm box. When applying the smooth-collar formulation of
Theorem~\ref{thm:r3-I-capacity}, use its prescribed-box auxiliary
collar reduction: preserve the weights near $\overline V$, up
to a common constant, and replace only their already common
exterior by a protected smooth collar. Their difference and all
energy integrands supported by that difference are unchanged.
The constants are uniform for the present bounded family. This
does not assert that the original $B\rho$ is smooth.
\end{proof}

\begin{lemma}[A constrained maximizer]
\label{lem:r3-capdir-maximizer}
Let $H\Subset V$ be compact, and let $\mu$ be a finite positive
Radon measure carried by $H$, with
\[
 \mu(E)\leq A C(E)^{1+\alpha}\quad
 (E\subset H\text{ Borel}),\qquad A,\alpha>0.
\]
For every $B>0$, the functional
\[
 \mathcal F(U)=\mathcal E(U)-\int_H u\,d\mu
\]
on the preceding affine family attains its supremum.
\end{lemma}

\begin{proof}
The common bounds, the fixed projective intersection mass and
finiteness of $\mu$ bound the functional above and below. Choose
a maximizing sequence $U_j$, with local weights $u_j+c$ on $V$.
The family is convex. The integral term is affine, so its midpoint
gap cancels, and Lemma~\ref{lem:r3-energy-midpoint} gives
\[
 0\leq c_n I(U_j,U_k)
 \leq\sup\mathcal F-
              \frac{\mathcal F(U_j)+\mathcal F(U_k)}2
 \longrightarrow0\quad(j,k\to\infty).
\]
The energy-to-capacity estimate on the fixed surrounding norm box
therefore makes the weights capacity Cauchy on every smaller
compact stage. The reverse framed domain comparison,
Corollary~\ref{cor:r2-fixed-domain-local-convergence}, upgrades
this to all-domain local capacity Cauchy before taking any
completion. The ambient completion and unique classical
representative, Corollary~\ref{cor:r2-ambient-classical-representative},
produce a bounded classical limit metric $U$.

Its actual bounds $L_B\leq U\leq F$ and its equality to $F$
outside $V$ follow, for example, from the decreasing tail-usc
formula for that representative: the two bounding metrics are
continuous, so both bounds pass through the tail regularization
and its decreasing limit. Thus the limit is an admissible member
of the same affine family. The completed weighted calculus on
a finite cover of $\overline V$ gives
$\mathcal E(U_j)\to\mathcal E(U)$; outside $V$ all weighted
differences in this formula vanish.

Let $M=\mu(V)$ and let $L$ be a bound for $|u_j|,|u|$. For
$\delta>0$, set $E_j=H\cap\{|u_j-u|>\delta\}$. Local capacity
convergence and a finite cover of $H$ give $C^*(E_j)\to0$ on
the fixed domain $V$. Hence
\[
 \int_H|u_j-u|\,d\mu
 \leq\delta M+2LA\,C^*(E_j)^{1+\alpha}.
\]
First let $j\to\infty$, then $\delta\downarrow0$. This proves
continuity of the measure term, so the admissible limit attains
the supremum. No weak convergence theorem for arbitrary varying
Monge--Amp\`ere measures was used.
\end{proof}

\begin{lemma}[The upper variation]
\label{lem:r3-capdir-upper-variation}
If $\mu\ne0$ and $U$ is the constrained maximizer, with local
function $u$, then $u<0$ everywhere on $V$ and
\[
 (Du)^n\leq\mu.
\]
In particular its Monge--Amp\`ere measure is carried by $H$.
\end{lemma}

\begin{proof}
The zero function is not a maximizer. Indeed the competitor
$tB\rho$, $0<t\leq1$, is admissible. The reference metric $F$
has constant weight on $V$, and all differences are zero off $V$.
Consequently its energy is
\[
 \mathcal E(U_{tB\rho})=
 \frac{t^{n+1}}{n+1}\int_V B\rho\,(BD\rho)^n.
\]
This is a finite constant times $t^{n+1}$. The change in the
measure term is $-t\int_H B\rho\,d\mu>0$, of first order in
$t$. Thus $u\ne0$. The classical strong maximum principle on
the connected domain makes the nonpositive $u$ strictly negative.

Let $\phi\geq0$ be an original compact smooth scalar test in
$V$, extended by zero to $Z$. Since $u$ is usc and negative,
its maximum on $\operatorname{supp}\phi$ is strictly negative.
For small $t>0$, $u+t\phi\leq0$ there. Let
\[
 P_t=P_L^{\mathrm b}(U+t\phi).
\]
The bounded-usc obstacle theorem supplies this actual greatest
classical minorant, and $U\leq P_t\leq U+t\phi$. It is equal
to $U$ off the support of $\phi$, and is an admissible member
of the same constrained family. Put $\Delta_t=P_t-U$.
Quasicontinuous-obstacle orthogonality,
Theorem~\ref{thm:r2-qc-orthogonality}, gives
\[
 \int\Delta_t c_1(L,P_t)^n
       =t\int\phi c_1(L,P_t)^n.
\]
Energy concavity and $0\leq\Delta_t\leq t\phi$ yield
\[
 \int\phi c_1(L,P_t)^n
 \leq\frac{\mathcal E(P_t)-\mathcal E(U)}t
 \leq\int\phi c_1(L,U)^n.
\]
Uniform convergence $P_t\to U$ and the proved bounded capacity
continuity identify both limits with $\int_V\phi(Du)^n$.
Maximality of $\mathcal F$ gives
\[
 \mathcal E(P_t)-\mathcal E(U)
       \leq\int_H\Delta_t\,d\mu
       \leq t\int_H\phi\,d\mu.
\]
Thus $(Du)^n\leq\mu$ on all nonnegative original scalar tests.
The coefficient-Radon identification turns this into the claimed
measure inequality. All perturbations have support inside $V$;
no boundary mass is discarded in these calculations.
\end{proof}

\begin{lemma}[An estimate independent of the lower obstacle]
\label{lem:r3-capdir-apriori}
Let $v\in\Cb(V)$ be bounded, nonpositive, with zero ordinary
boundary limit, and suppose $(Dv)^n\leq\mu$. For $M=\mu(V)>0$
put
\[
 S=\left(1+\frac2{1-2^{-\alpha}}\right)
                  (AM^\alpha)^{1/[n(1+\alpha)]}.
\]
Then $v\geq-S$. If its measure is carried by $H$, then with
$d_H=\min_H(-\rho)>0$ one has
\[
 v\geq\max\{-S,(S/d_H)\rho\}.
\]
\end{lemma}

\begin{proof}
For $s>0$ put $f(s)=C(\{v<-s\})^{1/n}$. These sets have compact
closure because of the ordinary boundary limit. Given a tester
$a\in\Cb(V)$, $0\leq a\leq1$, compare $v$ with
$w=-s-t+ta$. One has
\[
 \{v<-s-t\}\subset\{v<w\}\subset\{v<-s\}.
\]
The middle gap has compact closure. Comparison gives
\[
 t^n\int_{\{v<-s-t\}}(Da)^n
 \leq\int_{\{v<w\}}(Dw)^n
 \leq(Dv)^n(\{v<w\})\leq\mu(\{v<-s\}).
\]
Taking the supremum and using the measure bound on
$H\cap\{v<-s\}$ yields
\[
 t f(s+t)\leq A^{1/n} f(s)^{1+\alpha}.
\]
The same estimate without the capacity bound, followed by
$s\downarrow0$ and increasing-set continuity of raw capacity,
gives $t f(t)\leq M^{1/n}$.

Set $T=(AM^\alpha)^{1/[n(1+\alpha)]}$ and start with $s_0=T$.
As long as $f(s_j)>0$, take
$t_j=2A^{1/n}f(s_j)^\alpha$ and $s_{j+1}=s_j+t_j$.
Then $f(s_{j+1})\leq f(s_j)/2$, and the initial estimate gives
$t_j\leq2T2^{-j\alpha}$. Their sum is at most
$2T/(1-2^{-\alpha})$. If the iteration stops, all deeper
sublevels already have zero capacity; otherwise $f(S)=0$ by
monotonicity and the geometric decay. The strict sublevel
$\{v<-S\}$ is ordinary open because $v$ is usc. Ambient
capacity nondegeneracy on the fixed framed domain makes it empty.

The bounded classical barrier
$b=\max(-S,(S/d_H)\rho)$ equals $-S$ on $H$ and has zero
boundary limit. Thus $(Dv)^n(\{v<b\})=0$ when the measure is
carried by $H$. Pointwise domination, applied first with the
positive margin $b-\delta$ and then $\delta\downarrow0$,
gives $v\geq b$. No strict gap between two discontinuous
potentials has been called ordinary open.
\end{proof}

\begin{theorem}[Power-capacity Dirichlet existence]
\label{thm:r3-capacity-dirichlet}
Let $\mu$ be a finite positive Radon measure carried by a compact
$H\Subset V$, and assume the power bound in
Lemma~\ref{lem:r3-capdir-maximizer}. There is a unique bounded
$u\in\Cb(V)$ with zero ordinary boundary limit and
\[
 (Du)^n=\mu.
\]
For $M>0$ it satisfies
$\max(-S,(S/d_H)\rho)\leq u\leq0$, with $S$ as above.
For $M=0$ the solution is zero. All equalities hold on the full
original test sheaf. Continuity on $\overline V$ is proved later
by the one-sided Dini estimate.
\end{theorem}

\noindent\textbf{Status: \texttt{NEW\_PROVED}.}

\begin{proof}
Assume $M>0$ and choose $B>S/d_H$. The constrained maximizer
exists, and its upper variation gives $(Du)^n\leq\mu$.
The preceding lemma gives $u\geq(S/d_H)\rho$. Therefore
\[
 u-B\rho\geq(B-S/d_H)(-\rho)>0
 \quad\text{at every point of }V.
\]
This is a uniform positive margin on every compact perturbation
support, and releases the artificial lower obstacle there.

Let $\phi\geq0$ be original compact smooth in $V$. For sufficiently
small $t>0$, the global lower metric satisfies
$L_B\leq U-t\phi$. Put $P_t=P_L^{\mathrm b}(U-t\phi)$.
Since $L_B$ is a competitor and $U-t\|\phi\|_\infty$ is also
a global classical minorant,
\[
 L_B\leq P_t\leq U-t\phi\leq U,
 \qquad U-t\|\phi\|_\infty\leq P_t.
\]
Outside $V$ the lower metric equals $U$, so $P_t=U$ there.
Thus $P_t$ is admissible. Write $\Delta_t=U-P_t$.
The obstacle orthogonality identity is now
\[
 \int\Delta_t c_1(L,P_t)^n
       =t\int\phi c_1(L,P_t)^n.
\]
Since $\Delta_t\geq t\phi$, concavity gives
\[
 \int\phi c_1(L,U)^n
 \leq\frac{\mathcal E(U)-\mathcal E(P_t)}t
 \leq\int\phi c_1(L,P_t)^n.
\]
Uniform convergence again shows that the middle expression tends
to $\int_V\phi(Du)^n$. Maximality gives the opposite measure
inequality:
\[
 \mathcal E(U)-\mathcal E(P_t)
       \geq\int_H\Delta_t\,d\mu
       \geq t\int_H\phi\,d\mu.
\]
Hence $(Du)^n\geq\mu$, proving equality. The original sandwich
$B\rho\leq u\leq0$ gives the ordinary zero boundary limit,
and the a priori lemma gives the stated sharper bounds.

If $M=0$, take $u=0$. For uniqueness, if $u,v$ have the same
measure and zero boundary values, compare $v$ with
$(1+\epsilon)u-\delta$ on their compact strict gap. The
comparison inequality and the common measure force that gap to
have zero measure; ambient pointwise domination makes it empty.
Let $\delta\downarrow0$ and then $\epsilon\downarrow0$ and
repeat in the other order. All cancellations occur on compact
gaps of finite mass. This proves uniqueness without a weak-data
stability assumption.
\end{proof}

% END INCLUDED SOURCE: sections/capacity_dirichlet


% BEGIN INCLUDED SOURCE: sections/dini_capacity_dirichlet
\section{Dini capacity moduli for the fixed-domain Dirichlet problem}
\label{sec:r3-dini-capacity-dirichlet}

Retain the extendible domain $(Y,\rho,V)$ and the original ambient
background $[1]$ of Section~\ref{sec:r3-capacity-dirichlet}. Write
$C(E)=\operatorname{Cap}_{[1]}(E;V)$. All functions retain their
actual classical values over the prescribed field.

\begin{definition}[Dini capacity modulus]
\label{def:r3-dini-capacity-modulus}
A finite nondecreasing function $\omega:[0,\infty)\to[0,\infty)$
with $\omega(0)=0$ is a Dini capacity modulus if
\begin{equation}\label{eq:r3-dini-integrability}
 \int_0^a\omega(t)\,\frac{dt}{t}<\infty
\end{equation}
for some $a>0$. No continuity or strict positivity is imposed.
\end{definition}

Monotonicity implies measurability. The condition forces
$\omega(t)\to0$ at zero, and then holds for every finite positive
upper endpoint, since the integrand is bounded on compact intervals
away from zero.

\begin{theorem}[Dirichlet solvability under a Dini modulus]
\label{thm:r3-dini-capacity-dirichlet}
Let $H\Subset V$ be compact, and let $\mu$ be a finite positive
Radon measure carried by $H$. Suppose
\begin{equation}\label{eq:r3-dini-capacity-hypothesis}
 \mu(E)\leq C(E)\omega(C(E)^{1/n})^n
       \qquad(E\subset H\text{ Borel})
\end{equation}
for a Dini capacity modulus. There is a unique bounded
$u\in\Cb(V)$ with ordinary boundary limit zero and $(Du)^n=\mu$.
For $M=\mu(V)>0$, put
\begin{equation}\label{eq:r3-dini-S}
 S_\omega(s)=s+2\sum_{j\geq0}\omega(2^{-j}M^{1/n}/s),
 \qquad S_* =\inf_{s>0}S_\omega(s).
\end{equation}
Each series is finite. With $d_H=\min_H(-\rho)$, one has
\begin{equation}\label{eq:r3-dini-solution-bound}
 \max\{-S_*,(S_*/d_H)\rho\}\leq u\leq0.
\end{equation}
For $M=0$, $u=0$. The infimum defining $S_*$ need not be attained.
\end{theorem}

\noindent\textbf{Status: \texttt{DIRECT\_COROLLARY}} of the
preceding variational construction, with the scalar iteration below.
Continuity is established separately in
Theorem~\ref{thm:r3-dini-solution-continuity}.

\begin{lemma}[Dini sublevel iteration]
\label{lem:r3-dini-sublevel-iteration}
If $v\in\Cb(V)$ is bounded, nonpositive, has zero boundary limit,
and $(Dv)^n\leq\mu$, then $\|v\|_\infty\leq S_\omega(s)$ for
every $s>0$. If its measure is carried by $H$, it also satisfies
$v\geq\max(-S_\omega(s),(S_\omega(s)/d_H)\rho)$.
\end{lemma}

\begin{proof}
Set $f(s)=C(\{v<-s\})^{1/n}$. The compact-gap comparison used in
Lemma~\ref{lem:r3-capdir-apriori}, with normalized testers on this
same domain, gives
\begin{equation}\label{eq:r3-dini-recurrence}
 t f(s+t)\leq f(s)\omega(f(s)),\qquad
 t f(t)\leq M^{1/n}\quad(s,t>0).
\end{equation}
For the first inequality apply the measure bound on
$H\cap\{v<-s\}$ and use monotonicity of
$c\mapsto c\omega(c^{1/n})^n$. The second follows by increasing
continuity of raw capacity as the shift $s$ decreases to zero.

Starting at $s_0>0$, while $f(s_j)>0$, take
$t_j=2\omega(f(s_j))$ and $s_{j+1}=s_j+t_j$. If
$\omega(f(s_j))=0$, the first inequality makes $f(s_j+t)=0$
for every $t>0$. Increasing-set continuity then gives $f(s_j)=0$,
so the iteration has already stopped. Otherwise
\[
 f(s_{j+1})\leq\tfrac12 f(s_j),\qquad
 t_j\leq2\omega(2^{-j}M^{1/n}/s_0).
\]
The sum is finite because
\begin{equation}\label{eq:r3-dini-dyadic-bound}
 \sum_{j\geq0}\omega(2^{-j}a)
 \leq\omega(a)+\frac1{\log2}\int_0^a\omega(t)\,\frac{dt}{t}.
\end{equation}
Indeed the integral over $[2^{-j}a,2^{-j+1}a]$ is at least
$(\log2)\omega(2^{-j}a)$ for $j\geq1$. Monotonicity of $f$
now gives $f(S_\omega(s_0))=0$, whether the iteration terminates
or continues indefinitely. The strict sublevel
$\{v<-S_\omega(s_0)\}$ is an ordinary open set, since $v$ is usc.
Ambient capacity nondegeneracy makes it empty. The barrier
$b=\max(-S_\omega(s_0),(S_\omega(s_0)/d_H)\rho)$ equals
$-S_\omega(s_0)$ on $H$. Thus $v\geq b$ on the carrying set
of its measure, and compact-gap pointwise domination gives
$v\geq b$ throughout $V$.
\end{proof}

\begin{proof}[Proof of the theorem]
For $M>0$, fix $s_0>0$ and choose $B>S_\omega(s_0)/d_H$.
Use the class $\mathcal F_B$, its common affine collar and energy
from Lemma~\ref{lem:r3-capdir-affine-collar}. The power bound in
the constrained-maximizer argument was used for continuity of
integration against $\mu$ under bounded capacity convergence on
$H$. Here, if $|v_j|,|v|\leq L$ and
$c_j=C(H\cap\{|v_j-v|>\delta\})$, then
\[
 \int_H|v_j-v|\,d\mu
 \leq\delta M+2L c_j\omega(c_j^{1/n})^n.
\]
The last term tends to zero when $c_j\to0$, with value zero at
$c_j=0$. The same midpoint argument and ambient representative
therefore give a maximizer of the energy minus $\int v\,d\mu$.
The test $tB\rho$ again proves that it is not the zero function,
and the classical strong maximum principle makes it strictly negative.

The positive obstacle variation of
Lemma~\ref{lem:r3-capdir-upper-variation} yields $(Du_B)^n\leq\mu$.
The preceding lemma makes $u_B\geq(S_\omega(s_0)/d_H)\rho$.
The strict choice of $B$ leaves a positive lower-obstacle margin
on each compact perturbation support. The negative obstacle
variation in Theorem~\ref{thm:r3-capacity-dirichlet} therefore
applies without any change of domain or envelope input and gives
$(Du_B)^n\geq\mu$. These variations use the already proved ambient
quasicontinuous-obstacle orthogonality, not an arbitrary-relative
orthogonality assertion. Ambient pointwise comparison gives uniqueness.

Apply the a priori lemma to this one solution for every $s>0$.
Infimizing gives $u\geq-S_*$. Along any sequence of parameters
with $S_\omega(s)\to S_*$, the barrier inequalities give the
other bound in \eqref{eq:r3-dini-solution-bound}. No minimum is needed.
\end{proof}

\begin{remark}
For $\omega(t)=\mathbf1_{[1,\infty)}(t)$ and $M=1$,
$S_\omega(s)=s$ when $s>1$, whereas $S_\omega(s)\geq s+2$
when $0<s\leq1$. Thus its infimum is $1$ and is not attained.
\end{remark}

\begin{corollary}[Power specialization]
\label{cor:r3-dini-power-specialization}
For $\omega(t)=A^{1/n}t^\alpha$, $A,\alpha>0$, the measure
condition is $\mu(E)\leq A C(E)^{1+\alpha}$. For $M>0$,
\[
 \|u\|_\infty\leq
 (1+\alpha)\alpha^{-\alpha/(1+\alpha)}
 \left(\frac2{1-2^{-\alpha}}\right)^{1/(1+\alpha)}
 (AM^\alpha)^{1/[n(1+\alpha)]}.
\]
\end{corollary}
\begin{proof}
Sum the geometric series to obtain $S_\omega(s)=s+Qs^{-\alpha}$,
where $Q=2A^{1/n}M^{\alpha/n}/(1-2^{-\alpha})$, and minimize at
$s=(\alpha Q)^{1/(1+\alpha)}$.
\end{proof}

\begin{corollary}[Logarithmic modulus]
\label{cor:r3-dini-logarithmic}
For $L>0$ and $\beta>1$, the sufficient condition is
\begin{equation}\label{eq:r3-dini-logarithmic-hypothesis}
 \mu(E)\leq
 \frac{L^n C(E)}{[\log(e+C(E)^{-1/n})]^{n\beta}}.
\end{equation}
For $M>0$, $s\geq eM^{1/n}$ and $\ell=\log(s/M^{1/n})$,
\begin{equation}\label{eq:r3-dini-log-bound}
 \|u\|_\infty\leq
 s+2L\left(\ell^{-\beta}
       +\frac{\ell^{1-\beta}}{(\beta-1)\log2}\right).
\end{equation}
\end{corollary}
\begin{proof}
Use $\omega(t)=L[\log(e+1/t)]^{-\beta}$ for $t>0$ and
$\omega(0)=0$. Its Dini integral converges precisely for $\beta>1$.
For $a\leq e^{-1}$, bound $\omega(a)$ by
$L(\log(1/a))^{-\beta}$ and the Dini integral by
$L(\log(1/a))^{1-\beta}/(\beta-1)$. Insert these bounds into
\eqref{eq:r3-dini-dyadic-bound} with $a=M^{1/n}/s$.
\end{proof}

It is enough to prescribe the modulus and the measure condition
at capacities with $C(E)^{1/n}<a$. Extend it for $t\geq a$ by
$\max\{\omega(a),M^{1/n}/a\}$. The remaining sets have mass at
most $M$, so the extended condition holds. This does not assert
that every individual datum under a non-Dini bound is unsolvable;
the separate radial construction concerns a uniform guarantee.

% END INCLUDED SOURCE: sections/dini_capacity_dirichlet


% BEGIN INCLUDED SOURCE: sections/dini_stability_continuity
\section{One-sided capacity stability and continuity of Dini solutions}
\label{sec:r3-dini-stability-continuity}

Here the background is $[1]$. All functions are actual bounded
classical psh functions over the fixed original field, and all
currents use the original CLD tests. On a relatively compact ordinary
open $V$ in a permitted smooth affine neighborhood, write
$C(E)=\operatorname{Cap}_{[1]}(E;V)$. For a Dini modulus put
\[
 \Phi_\omega(b)=2\sum_{j\geq0}\omega(2^{-j}b),\qquad
 \Phi_\omega(0)=0.
\]
The dyadic integral comparison proves
\begin{equation}\label{eq:r3-stability-dyadic-tail}
 \Phi_\omega(b)\leq2\omega(b)+\frac2{\log2}
          \int_0^b\omega(r)\,\frac{dr}{r},\qquad
 \Phi_\omega(b)\longrightarrow0\quad(b\downarrow0).
\end{equation}

\begin{theorem}[One-sided support-capacity estimate]
\label{thm:r3-one-sided-dini-stability}
Let $u,v\in\Cb(V)$ be bounded and suppose
\[
 \overline{\{u<v-\delta\}}\Subset V\qquad(\delta>0).
\]
Suppose $\mu=(Du)^n$ is carried by a compact $H\Subset V$ and
\begin{equation}\label{eq:r3-stability-dini-measure}
 \mu(E)\leq C(E)\omega(C(E)^{1/n})^n
             \quad(E\subset H\text{ Borel}).
\end{equation}
Then for every $s>0$,
\begin{equation}\label{eq:r3-one-sided-dini-bound}
 \sup_V(v-u)\leq s+\Phi_\omega\bigl(
                  C(H\cap\{u<v-s\})^{1/n}\bigr).
\end{equation}
In particular, if $v_j$ satisfies the same gap condition and
converges to $u$ in outer capacity on $H$, with capacity domain
$V$, then $\limsup_j\sup_V(v_j-u)\leq0$.
\end{theorem}

\noindent\textbf{Status: \texttt{NEW\_PROVED}.}

\begin{proof}
Put $E_s=\{u<v-s\}$ and $f(s)=C(H\cap E_s)^{1/n}$.
It is finite by CLN on $H$. For any normalized classical tester
$0\leq b\leq1$, set $w=v-s-t+tb$. Then
\[
 E_{s+t}\subset\{u<w\}\subset E_s.
\]
The middle set has compact closure. Positivity of the mixed
expansion and compact-gap comparison give
\[
 t^n\int_{E_{s+t}}(Db)^n
 \leq\int_{\{u<w\}}(Dw)^n
 \leq\mu(\{u<w\})\leq\mu(H\cap E_s).
\]
Taking the supremum over $b$ gives the stronger full-gap bound
\begin{equation}\label{eq:r3-stability-recurrence}
 t^n C(E_{s+t})\leq f(s)^n\omega(f(s))^n,
 \qquad t f(s+t)\leq f(s)\omega(f(s)).
\end{equation}
Starting at $s$, repeat the halving iteration of
Lemma~\ref{lem:r3-dini-sublevel-iteration}, with increments
$2\omega(f(s_j))$. The total increment is at most
$\Phi_\omega(f(s))$. A zero increment with $f(s_j)>0$ is
impossible: the recurrence and increasing-set continuity give
$f(s_j)=0$ in that case. Hence, for
$L=s+\Phi_\omega(f(s))$, one has $f(L)=0$.

The measure condition gives $\mu(E_L)=\mu(H\cap E_L)=0$.
The strict gap $E_L$ is only Borel; it is not called ordinary
open. Apply the proved ambient pointwise domination theorem to
$u$ and $v-L$. Its gap has compact closure, and the measure of
$u$ gives it no mass, so $u\geq v-L$ at every point. This proves
\eqref{eq:r3-one-sided-dini-bound}. The convergence statement
follows first for any fixed $s>0$ from
\eqref{eq:r3-stability-dyadic-tail}, and then by $s\downarrow0$.
\end{proof}

\begin{theorem}[Continuity on the ordinary closure]
\label{thm:r3-dini-solution-continuity}
On the extendible exhaustion domains of
Theorem~\ref{thm:r3-dini-capacity-dirichlet}, its solution extends
to a continuous function on $\overline V$ with value zero on
$\partial V$. In particular this holds for power-capacity data.
\end{theorem}

\begin{proof}
The proved compact-support barrier gives $u\geq C\rho$ for a
finite $C$. The common-exterior extension lemma produces a
classical function $U$ on a neighborhood of $\overline V$,
equal to $u$ on $V$ and to $C\rho$ on the exterior. The function
is zero on the ordinary boundary. This is a bounded classical
extension, not yet an assertion of interior continuity.

B's regularization provides continuous classical $w_j\downarrow U$
on one neighborhood of $\overline V$, with a common bound there.
Since $U=0$ on the compact boundary, Dini's lemma applied only to
the continuous restrictions on $\partial V$ gives
\[
 \epsilon_j=\max_{\partial V}w_j\downarrow0.
\]
Put $v_j=w_j-\epsilon_j$ on $V$. The boundary upper limit of
$v_j-u$ is at most zero. Thus each positive-margin gap has compact
closure. Decreasing relative capacity convergence of the $w_j$,
and the fixed-domain comparison on the compact $H$, give
$v_j\to u$ in outer ambient capacity there. The preceding theorem
implies $\limsup_j\sup_V(v_j-u)\leq0$.
Since $w_j\geq u$, it follows that $w_j\to u$ uniformly on $V$.
On the boundary the errors are at most $\epsilon_j$. Hence the
continuous $w_j|_{\overline V}$ converge uniformly to the extension
of $u$ by zero, proving continuity on the whole ordinary closure.
\end{proof}

\begin{theorem}[Strong-data uniform stability for power growth]
\label{thm:r3-power-dual-uniform-stability}
Fix an extendible $(Y,\rho,V)$, a compact $H\Subset V$, $B>0$,
and $A,\alpha>0$. Suppose $u,v\in\Cb(V)$ have ordinary boundary
limit zero, $-B\leq u,v\leq0$, and their measures $\mu,\nu$
are carried by $H$ and obey
\[
 \mu(E),\nu(E)\leq A C(E)^{1+\alpha}
             \qquad(E\subset H\text{ Borel}).
\]
With
\[
 d=d_{\mathrm{psh},V}(\mu,\nu)
   =\sup_{\substack{b\in\Cb(V)\\0\leq b\leq1}}
           \left|\int b\,d(\mu-\nu)\right|,
\]
there is a finite constant depending only on the fixed data and
the specified collar and finite CLN controls such that
\begin{equation}\label{eq:r3-power-dual-uniform-bound}
 \|u-v\|_\infty\leq C_0
             d^{\alpha/[2^n(n+2\alpha)]}.
\end{equation}
No optimality of this exponent is asserted.
\end{theorem}

\begin{proof}
The common compact support and bound give, by
Lemma~\ref{lem:r3-common-exterior-extension}, one larger domain
$\Omega$ and classical extensions $U,V'$ agreeing with a fixed
collar function outside $V$. Their mixed energy is
\[
 I=\int_H(u-v)\,d(\nu-\mu)\leq2B d.
\]
Indeed $(u+B)/B$ and $(v+B)/B$ are normalized testers on $V$;
the two constant terms cancel in their difference. The extensions
have one bound on the chosen larger compact stage. Apply
Theorem~\ref{thm:r3-I-capacity}, then reverse domain comparison
only on $H\Subset V$, to obtain
\begin{equation}\label{eq:r3-dual-gap-capacity}
 C(H\cap\{|u-v|>s\})\leq C_1s^{-2}d^{2^{-n}}.
\end{equation}
The reverse comparison has been used on a compact core, not at
the boundary of $V$.

For the power modulus, $\Phi_\omega(b)=
2A^{1/n}b^\alpha/(1-2^{-\alpha})$. Apply the one-sided estimate
in both orders and use \eqref{eq:r3-dual-gap-capacity}:
\[
 \|u-v\|_\infty\leq
 s+C_2s^{-2\alpha/n}d^{\alpha/(n2^n)}.
\]
For $0<d\leq1$ take $s=d^{\alpha/[2^n(n+2\alpha)]}$.
For $d=0$ let $s\downarrow0$ in the gap estimate and domination,
or use that the dual distance determines compact measures and
ambient uniqueness. For $d>1$, the uniform bound $2B$ is absorbed
in the constant. This proves the claimed exponent.
\end{proof}

\begin{corollary}[Strong-data stability under a common Dini modulus]
\label{cor:r3-dini-dual-uniform-stability}
Under the same domain, compact support and uniform potential bound,
replace the two power bounds by one common Dini bound. If
$d_{\mathrm{psh},V}(\mu_j,\mu)\to0$, the corresponding bounded
zero-boundary solutions converge uniformly on $\overline V$.
\end{corollary}
\begin{proof}
The energy argument still gives
\eqref{eq:r3-dual-gap-capacity}. For a fixed $s>0$, the capacity
on the right tends to zero. Apply the two one-sided Dini estimates
and then \eqref{eq:r3-stability-dyadic-tail}; finally let $s$ decrease
to zero. The boundary values agree. This is a
\texttt{DIRECT\_COROLLARY}, not a weak-measure stability theorem.
\end{proof}

% END INCLUDED SOURCE: sections/dini_stability_continuity


% BEGIN INCLUDED SOURCE: sections/maximal_boundary_dirichlet
\section{Nonconstant boundary data with a continuous maximal lift}
\label{sec:r3-maximal-boundary-dirichlet}

Keep the original field and the extendible domain $(Y,\rho,V)$
of Section~\ref{sec:r3-capacity-dirichlet}. Let $h$ be a continuous
classical psh function on all of $Y^{\an}$ and suppose
\[
 (Dh)^n=0\quad\text{on }V.
\]
Write $g=h|_{\partial V}$. This is a hypothesis on a specified
classical lift of the boundary function; it is not an assertion
that every continuous boundary function admits such a lift.

\begin{theorem}[Dini data with a continuous maximal boundary lift]
\label{thm:r3-maximal-boundary-dirichlet}
Let $K_0\Subset V$ be compact and let $\mu$ be a finite positive
Radon measure carried by $K_0$. Suppose, for a Dini modulus
$\omega$ as in Definition~\ref{def:r3-dini-capacity-modulus}, that
\[
 \mu(E)\leq C(E)\omega(C(E)^{1/n})^n,
 \qquad E\subset K_0\text{ Borel},\qquad
 C(E)=\operatorname{Cap}_{[1]}(E;V).
\]
There is a unique bounded classical function $u$ on $V$,
continuous on $\overline V$, with boundary value $g$ and
$(Du)^n=\mu$. If $M=\mu(V)>0$, set
\[
 d=\min_{K_0}(-\rho),\qquad
 S_\omega(s)=s+2\sum_{j\geq0}\omega(2^{-j}M^{1/n}/s),
 \qquad S_* =\inf_{s>0}S_\omega(s).
\]
Then the actual pointwise bounds are
\[
 h+\max\{-S_*,(S_*/d)\rho\}\leq u\leq h.
\]
If $M=0$, the solution is $u=h$. The uniqueness assertion holds
among all bounded classical solutions with the stated ordinary
boundary limits, without an assumed interior continuity.
\end{theorem}

\noindent\textbf{Status: \texttt{NEW\_PROVED}} for this explicit
nonconstant-boundary extension. The mixed calculus, contact
inequality, projective obstacle theorems and scalar Dini iteration
are previously established inputs.

\begin{proof}
Assume $M>0$. Fix $s_0>0$, put $S=S_\omega(s_0)$, and choose
$B>S/d$. The constrained family is
\[
 \mathcal F_B^h=\{v\in\Cb(V):h+B\rho\leq v\leq h\}.
\]
Its members already have the ordinary boundary limit $g$.

\emph{The common affine extension.}
For $\epsilon>0$ glue $\max(v,h+B\rho+\epsilon)$ on $V$
to $h+B\rho+\epsilon$ off $V$. Near each boundary point the
second branch is larger than $h$, hence larger than $v$.
The prescriptions agree on an ordinary neighborhood, and their
decreasing limit is classical, equal to $v$ on $V$ and to
$h+B\rho$ outside. All members agree outside $V$ and have one
bound on every chosen surrounding compact box. The proof of
Lemma~\ref{lem:r3-capdir-affine-collar} therefore gives a single
ample line on a smooth projective compactification, and metrics
$U_v$ with weight $v+c$ on a neighborhood of $\overline V$.
The map remains affine, since its maximum collar is inactive
on $V$ and all inputs are identical outside $V$. Its continuous
lower and upper metrics $L_B,F$ have weights $h+B\rho+c$ and
$h+c$ on $V$, and agree off $V$. Use the global mixed energy
$\mathcal E_F$ relative to $F$.

\emph{Existence of a constrained maximizer.}
Maximize
\[
 \mathcal F(U_v)=\mathcal E_F(U_v)-\int_{K_0}v\,d\mu.
\]
The common bounds and projective intersection mass bound this
functional. The midpoint argument of
Lemma~\ref{lem:r3-capdir-maximizer} applies without a change:
the integral is affine, the midpoint energy gap controls $I$,
and the prescribed-box estimate followed by reverse framed domain
comparison makes a maximizing sequence locally capacity Cauchy
on every permitted local domain. The ambient completion has a
unique actual classical representative. The continuous upper and
lower metrics pass through its tail-usc formula, so this limit
is admissible and equals the common exterior metric off $V$.
Weighted capacity continuity gives continuity of the energy.
For the measure term, if $|v_j|,|v|\leq L$ on $K_0$ and
$c_j=C^*(K_0\cap\{|v_j-v|>\delta\})$, then
\[
 \int_{K_0}|v_j-v|\,d\mu
 \leq\delta M+2L c_j\omega(c_j^{1/n})^n\longrightarrow\delta M.
\]
Letting $\delta\downarrow0$ proves the required continuity.
Thus a constrained maximizer $U$ exists, with local function $u$.

\emph{The positive variation and the contact set.}
The set $O=\{u<h\}$ is ordinary open because $u-h$ is usc and
$h$ is continuous. For a nonnegative original smooth $\phi$
with compact support in $O$, choose $t>0$ small enough that
$u+t\phi\leq h$ there. Extend $\phi$ by zero to the projective
space and set $P_t=P_L^{\mathrm b}(U+t\phi)$. The actual
usc-obstacle envelope satisfies $U\leq P_t\leq U+t\phi$, is
unchanged off the perturbation support, and is admissible.
Orthogonality and energy concavity give, exactly with these
global metrics,
\[
 \int\phi c_1(L,P_t)^n
 \leq\frac{\mathcal E_F(P_t)-\mathcal E_F(U)}t
 \leq\int\phi c_1(L,U)^n.
\]
Uniform convergence of $P_t$ to $U$ gives the common limit
$\int\phi(Du)^n$. Maximality bounds the middle expression by
$\int\phi\,d\mu$. Thus $(Du)^n\leq\mu$ on $O$.

On the Borel contact set $A=\{u=h\}$ the proved contact
inequality, Corollary~\ref{cor:r3-relative-zero-contact-mass},
gives
\[
 \mathbf1_A(Du)^n\leq\mathbf1_A(Dh)^n=0.
\]
Combining the two disjoint pieces yields $(Du)^n\leq\mu$
on all of $V$. In particular this measure is carried by $K_0$.
There is no use of a strong maximum principle for $u-h$, which
is not asserted to be psh.

\emph{The estimate for the difference without differentiating it.}
Put $f(s)=C(\{u<h-s\})^{1/n}$. For a normalized classical
tester $0\leq a\leq1$ on $V$, compare $u$ with the classical
function
\[
 w=h-s-t+ta.
\]
The mixed expansion is positive and gives
$(Dw)^n\geq t^n(Da)^n$. Moreover
\[
 \{u<h-s-t\}\subset\{u<w\}\subset\{u<h-s\}.
\]
All positive-margin gaps have compact closure by the common
boundary values. Comparison and $(Du)^n\leq\mu$ therefore
give
\[
 t f(s+t)\leq f(s)\omega(f(s)),\qquad
 t f(t)\leq M^{1/n}.
\]
The second estimate follows by decreasing the auxiliary positive
shift to zero and using increasing-set continuity of raw capacity.
Apply the scalar halving argument from
Lemma~\ref{lem:r3-dini-sublevel-iteration}: starting at $s_0$,
the increments $2\omega(f(s_j))$ make $f$ decrease by factors
of two, and their sum is bounded by
$2\sum_{j\geq0}\omega(2^{-j}M^{1/n}/s_0)$. If the modulus
vanishes at a positive value encountered, the recurrence makes
all deeper capacities zero and the iteration stops. In either
case $f(S)=0$. The set $\{u<h-S\}$ is ordinary open, so
ambient capacity nondegeneracy makes it empty. Thus $u\geq h-S$.
This argument has used the psh functions $u,h,w,a$ themselves;
it has not treated $u-h$ as a psh potential.

The classical barrier
\[
 b=h+\max\{-S,(S/d)\rho\}
\]
equals $h-S$ on $K_0$, so $b\leq u$ there. Since $(Du)^n$
is carried by $K_0$, the strict gap $\{u<b\}$ has zero
$(Du)^n$-mass. Pointwise domination applied first to $b-\delta$
and then with $\delta\downarrow0$ gives $b\leq u$ everywhere.
In particular
\[
 u-(h+B\rho)\geq(B-S/d)(-\rho)>0\quad\text{on }V.
\]

\emph{Release of the lower obstacle.}
For an arbitrary nonnegative original compact smooth $\phi$ in
$V$, this strict margin permits $L_B\leq U-t\phi$ for all
sufficiently small $t>0$. Let $P_t=P_L^{\mathrm b}(U-t\phi)$.
The lower metric and the global constant shift give
\[
 L_B\leq P_t\leq U-t\phi\leq U,
 \qquad U-t\|\phi\|_\infty\leq P_t.
\]
Outside $V$ the lower metric equals $U$, so $P_t=U$ there and
$P_t$ is admissible. Orthogonality and concavity, now in the
negative direction, give
\[
 \int\phi c_1(L,U)^n
 \leq\frac{\mathcal E_F(U)-\mathcal E_F(P_t)}t
 \leq\int\phi c_1(L,P_t)^n.
\]
For clarity, orthogonality reads
$\int(U-P_t)c_1(L,P_t)^n=t\int\phi c_1(L,P_t)^n$;
the lower inequality uses $U-P_t\geq t\phi$.
The common limit is $\int\phi(Du)^n$. Maximality gives
\[
 \mathcal E_F(U)-\mathcal E_F(P_t)
       \geq\int_{K_0}(U-P_t)\,d\mu
       \geq t\int_{K_0}\phi\,d\mu.
\]
Thus $(Du)^n\geq\mu$, and equality follows. All these energy
differences are global expressions on one line; no boundary mass
is inferred to vanish merely by restriction.

The construction gives one solution. Repeat only the a priori
estimate for that solution for each $s>0$ and take a sequence
with $S_\omega(s)\to S_*$. This gives the two displayed bounds
with $S_*$; no minimizing parameter is required.

\emph{Continuity and uniqueness.}
The constructed global metric admits bounded continuous classical
decreasing regularizations on one neighborhood of $\overline V$
by B. In the fixed frame write their local functions as
$w_j\downarrow u$ on $V$; on the boundary their limit is $h$.
Compactness of the boundary and Dini's lemma give
\[
 \epsilon_j=\max_{\partial V}(w_j-h)\downarrow0.
\]
The functions $v_j=w_j-\epsilon_j$ have boundary upper limit
at most $g$. Their positive-margin gaps above $u$ have compact
closure. Decreasing capacity convergence and domain comparison
give $v_j\to u$ in outer capacity on $K_0$, with domain $V$.
Theorem~\ref{thm:r3-one-sided-dini-stability} applies; its gap
hypothesis does not require constant boundary values. It gives
$\limsup_j\sup_V(v_j-u)\leq0$. Since $w_j\geq u$ and
$\epsilon_j\to0$, the $w_j$ converge uniformly on $V$, and
also on its boundary, to the extension of $u$ by $g$. This is
the required continuity on $\overline V$.

If two bounded classical solutions have the same ordinary
boundary limits $g$, the gaps of one above the other with any
positive margin have compact closure. Compact-gap comparison
and ambient pointwise domination give equality. One can verify
this with the scaling argument after subtracting one sufficiently
large constant from both functions: their boundary values then
are strictly negative, so the comparison of $b$ with
$(1+\epsilon)a-\delta$ has compact gap. Equality of their
measures and their finite mass on that compact gap force the
gap to have zero measure, and domination removes it. Let
$\delta,\epsilon\downarrow0$ and interchange $a,b$.
For $M=0$, $h$ itself solves the problem and the same uniqueness
argument applies. This completes the proof.
\end{proof}

\begin{corollary}[An upper lift with curvature dominated by the datum]
\label{cor:r3-dominated-curvature-boundary-lift}
The preceding theorem remains valid if the assumption $(Dh)^n=0$
is replaced by
\[
 (Dh)^n\leq\mu\quad\text{on }V.
\]
All other hypotheses, the ordinary boundary values, continuity,
uniqueness and estimates are unchanged.
\end{corollary}

\noindent\textbf{Status: \texttt{DIRECT\_COROLLARY}.}

\begin{proof}
The contact inequality now gives
$\mathbf1_A(Du)^n\leq\mathbf1_A(Dh)^n\leq\mathbf1_A\mu$.
Combined with the positive variation on $\{u<h\}$ it still
gives $(Du)^n\leq\mu$ everywhere. Every subsequent step uses
that inequality, the positivity of mixed terms, and the fixed
boundary lift, and is unchanged. If $\mu=0$, the new hypothesis
also forces $(Dh)^n=0$. The direction of the curvature inequality
is substantive: this is an upper lift of smaller curvature than
the prescribed datum, not an arbitrary continuous extension.
\end{proof}

\begin{example}[A nonconstant maximal lift in every dimension at least two]
On $Y=\mathbf A_K^n$, $n\geq2$, take
$h=\max(0,\log|T_1|)$ and
\[
 \rho=\max\{-1,\log|T_1|-\log R,\ldots,
                       \log|T_n|-\log R\},\qquad R>1.
\]
These are continuous classical psh functions, and
$V=\{|T_i|<R\text{ for every }i\}$ is an extendible domain
of the required type. The curvature $Dh$ is nonzero, but its
top power is zero because $h$ depends on one coordinate only
and $n\geq2$. Its boundary restriction is nonconstant: it is
zero at a boundary point with $T_1=0$ and $|T_2|=R$, and is
$\log R$ where $|T_1|=R$. Thus the theorem includes explicit
nonconstant boundary data with a nonconstant classical lift.
\end{example}

% END INCLUDED SOURCE: sections/maximal_boundary_dirichlet


% BEGIN INCLUDED SOURCE: sections/subsolution_exhaustion
\section{Exhaustion of subsolution-dominated data}
\label{sec:r3-subsolution-exhaustion}

Keep the extendible domain $(Y,\rho,V)$ of
Section~\ref{sec:r3-capacity-dirichlet}, the original field and
absolute value, and the ambient background $[1]$. Boundary limits
in this section are ordinary boundary limits at every point of
$\partial V$. The measures need only be locally finite.

\begin{theorem}[The bounded subsolution criterion without compact support]
\label{thm:r3-exhausted-subsolution-dirichlet}
Let $\mu$ be a locally finite positive Radon measure on $V$.
There exists a bounded $u\in\Cb(V)$ with zero ordinary boundary
limit and $(Du)^n=\mu$ if and only if there exists a bounded
$w\in\Cb(V)$ with zero ordinary boundary limit and
\[
 (Dw)^n\geq\mu.
\]
The solution is unique and satisfies $w\leq u\leq0$ for each
such subsolution $w$. The measure $\mu$ is not required to have
finite total mass or compact support. A constant boundary value
is handled by adding that constant to the solution and subsolution.
\end{theorem}

\noindent\textbf{Status: \texttt{DIRECT\_COROLLARY}.}
The inputs are the proved compact-data bounded-subsolution theorem,
the fixed-background decreasing BT theorem, and ambient pointwise
comparison. The matching boundary condition on $w$ is essential
to this exhaustion argument.

\begin{proof}
The reverse implication follows by taking $w=u$.
For the forward construction choose nested compact sets
$H_j\Subset V$ with $H_j\subset H_{j+1}^{\circ}$ and
$\bigcup_jH_j^{\circ}=V$, and put $\mu_j=\mathbf1_{H_j}\mu$.
These are finite compactly carried measures and $\mu_j\uparrow\mu$.
The same bounded $w$ is a subsolution for each $\mu_j$.
Theorem~\ref{thm:r3-bounded-subsolution-dirichlet} supplies bounded
classical $u_j$ with zero boundary limit and $(Du_j)^n=\mu_j$.
The weak-boundary order lemma
\ref{lem:r3-weak-boundary-subsolution-order} gives
\[
 w\leq u_{j+1}\leq u_j\leq0.
\]
Consequently $u_j$ decreases to an actual bounded classical
function $u$, with $w\leq u\leq0$. The boundary sandwich implies
that $u$ has zero limit at every ordinary boundary point. The
decreasing BT theorem gives $(Du_j)^n\to(Du)^n$ on every original
compact smooth scalar test. If such a test has support $H$, then
$H\subset H_j$ for every sufficiently large $j$, so its pairing
with $\mu_j$ equals its pairing with $\mu$. Thus $(Du)^n=\mu$.
No passage of total mass through the boundary is used.

For completeness, the comparison argument for uniqueness uses only
finite mass on compact gaps, even if $\mu(V)=\infty$.
Let $a,b$ be two bounded zero-boundary classical functions with
$(Da)^n\geq(Db)^n$. For $\epsilon,\delta>0$ the strict gap
\[
 E=\{b<(1+\epsilon)a-\delta\}
\]
has compact closure in $V$ by their boundary limits. The compact-gap
comparison theorem gives
\[
 (1+\epsilon)^n(Da)^n(E)\leq(Db)^n(E)\leq(Da)^n(E).
\]
All three quantities are finite on this compact gap, so both
measures vanish on $E$. Ambient pointwise domination, with lower
function $b$, implies $b\geq(1+\epsilon)a-\delta$ on $V$.
Let $\delta\downarrow0$ and then $\epsilon\downarrow0$ to obtain
$a\leq b$. Equality of the measures gives equality of the
functions by applying the argument in both orders. The same
argument verifies the order used in the exhaustion, and uses
strict gaps as Borel sets, without asserting that they are ordinary
open. Translation by a constant proves the final assertion.
\end{proof}

\begin{corollary}[Mixed data dominated by bounded potentials]
\label{cor:r3-mixed-datum-dirichlet}
Let $u_1,\ldots,u_n\in\Cb(V)$ be globally bounded and let
$f$ be a Borel function with $0\leq f\leq C<\infty$. Set
\[
 \mu=f\,Du_1\wedge\cdots\wedge Du_n.
\]
If $f$ is compactly supported, this measure has a unique bounded
zero-boundary classical solution. No boundary condition on the
$u_i$ is needed in this compactly supported case. If instead every
$u_i$ has zero ordinary boundary limit, the same conclusion holds
without compact support of $f$ and without a finite-total-mass
assumption.
\end{corollary}

\noindent\textbf{Status: \texttt{DIRECT\_COROLLARY}.}

\begin{proof}
If $C=0$, take the zero solution. Otherwise put
\[
 w=\left(\frac{C}{n!}\right)^{1/n}\sum_{i=1}^n u_i.
\]
This is a bounded classical function. Positivity of all mixed
terms in the expansion and the coefficient $n!$ of the product
with all $n$ different factors give
\[
 (Dw)^n
  =\frac{C}{n!}\left(\sum_i Du_i\right)^n
  \geq C\,Du_1\wedge\cdots\wedge Du_n\geq\mu.
\]
The mixed measure is locally finite by the local CLN theorem, so
bounded Borel multiplication defines a locally finite Radon measure.
For compactly supported $f$, apply
Theorem~\ref{thm:r3-bounded-subsolution-dirichlet}, which requires
only a bounded subsolution in this case. If all the $u_i$ have
zero boundary limit, then so does $w$, and the preceding theorem
applies to the possibly noncompact data. This is an ambient
diagonal Dirichlet conclusion; it makes no arbitrary-background
relative Dirichlet existence assertion.
\end{proof}

% END INCLUDED SOURCE: sections/subsolution_exhaustion


% BEGIN INCLUDED SOURCE: sections/capacity_growth_instability
% Reconstructed after the scratch reset. Original reviewed SHA256:
% 988aa4a454cc77b8fefb752f2fcdc210c46e80b5eb1cb542d69f07af2a470f6a.
\section{Uniform power capacity growth does not imply weak-data stability}
\label{sec:r3-capacity-growth-instability}

Keep $\Omega,\lambda,g_i,H,z_i,z,K_*,\Omega_0$ from
Section~\ref{sec:r3-moving-atomic-dirichlet}. Write
$\mu_i=\delta_{z_i}$ and $\mu=\delta_z$.

\begin{corollary}[Weak instability under every fixed power bound]
\label{cor:r3-power-capacity-weak-instability}
For each fixed $\alpha>0$ there is one constant
$A_\alpha=(2\lambda)^{n(1+\alpha)}$ such that all Borel
$E\subset\Omega$ satisfy
\begin{equation}
 \mu_i(E),\ \mu(E)
       \leq A_\alpha\operatorname{Cap}_{[1]}(E;\Omega)^{1+\alpha}
       \quad\hbox{for every }i.
 \label{eq:r3-common-power-capacity-bound}
\end{equation}
Nevertheless the corresponding continuous bounded zero-boundary
solutions $g_i$ do not converge uniformly or locally in outer
capacity to $H$, and have no local capacity-Cauchy subsequence.
For $n\geq2$ they do not converge pointwise to $H$.
Moreover
\[
 d_{\mathrm{psh},\Omega}(\mu_i,\mu)\geq\tfrac12
 \quad\hbox{for every }i .
\]
\end{corollary}

\noindent\textbf{Status: \texttt{DIRECT\_COROLLARY}} of the
explicit moving-atom construction. The assertion that weak data
convergence, common compact support, common zero boundary values,
and a uniform power capacity bound force capacity or uniform
solution convergence is \texttt{FALSE}.

\begin{proof}
The functions $b_i=1+g_i/(2\lambda)$ and
$b=1+H/(2\lambda)$ are normalized classical psh on $\Omega$,
because $B<2\lambda$. Their top measures are
$(2\lambda)^{-n}\delta_{z_i}$ and
$(2\lambda)^{-n}\delta_z$. Hence every set containing the
respective atom has capacity at least $(2\lambda)^{-n}$.
This proves \eqref{eq:r3-common-power-capacity-bound}; sets
without that atom have zero measure. The convention at infinite
capacity is harmless.

Let $0<\epsilon<2\lambda-B$. The set
$K_*\cap\{|g_i-H|>\epsilon\}$ contains $z_i$.
The same normalized tester restricted to the smaller ordinary
domain $\Omega_0$ shows that its raw, and thus its outer,
capacity on that fixed domain is at least $(2\lambda)^{-n}$.
This is failure of local all-domain capacity convergence, not
merely an assertion on the original large domain.
For $i\ne j$, $|g_i-g_j|(z_i)=\lambda$; a fixed threshold less
than $\lambda$ gives the same lower bound for every pair and
rules out a capacity-Cauchy subsequence. Uniform and, for
$n\geq2$, pointwise failures were explicitly computed in the
preceding theorem. Finally $b_i(z_i)=0$ and $b_i(z)=1/2$,
so this one normalized tester proves the psh-dual lower bound.
\end{proof}

This is consistent with the proved capacity-growth existence
theorems and their a priori bounds. It excludes ordinary weak
data convergence as a substitute for the strong psh-dual
topology in the corresponding stability assertion.

% END INCLUDED SOURCE: sections/capacity_growth_instability


% BEGIN INCLUDED SOURCE: sections/linear_capacity_obstruction
% Reconstructed after the scratch reset. Original reviewed SHA256:
% 21e71bd66b0528578d68192e12f48767d20d020f720e5d8bcd8ea5c555cf2248.
\section{Linear capacity domination does not ensure a bounded solution}
\label{sec:r3-linear-capacity-obstruction}

Fix $r=|\pi|$, $\lambda=-\log r>0$, $R=r^{-1}$, and the
ordinary polydisc $\Omega=\{|T_l|<R\}_{l=1}^n$. Set
\[
 t_j=2^j\lambda,\qquad r_j=R e^{-t_j}=r^{\,2^j-1},\qquad
 s=\max_l\log|T_l|-\log R .
\]
Let $z_j$ be the diagonal radius-$r_j$ Gauss point, let
$B_j=\{|T_l|\leq r_j\}_{l=1}^n$, and let $o$ be the rigid origin.
Thus $z_j\to o$. Write $K_j=\{o,z_j,z_{j+1},\ldots\}$.

\begin{lemma}[Exact radial compact capacities]
\label{lem:r3-radial-compact-capacity}
For every $j$,
\[
 U_j=\max(-1,s/t_j),\qquad
 (DU_j)^n=t_j^{-n}\delta_{z_j},
\]
and, on the same fixed domain $\Omega$,
\[
 \operatorname{Cap}_{[1]}(\{z_j\};\Omega)
 =\operatorname{Cap}_{[1]}(K_j;\Omega)
 =\operatorname{Cap}_{[1]}(B_j;\Omega)=t_j^{-n}.
\]
\end{lemma}

\begin{proof}
The ordinary torus slope simplex has volume $t_j^{-n}/n!$;
the original factor $n!$ gives the measure identity, with no
mass on lower-rank strata or coordinate zero loci. The function
$1+U_j$ is a normalized tester, giving the singleton lower bound.
For an arbitrary normalized $a\in\mathcal C_b(\Omega)$ and
$\epsilon>0$, compare $v=(1+\epsilon)U_j$ with $a-1$.
The compact set $B_j$ lies in $\{v<a-1\}$. For every
$\delta>0$ the closure of $\{v<a-1-\delta\}$ is compactly
contained in $\Omega$, since it lies in $\{v<-\delta\}$.
Thus the proved comparison principle gives
\[
 \int_{B_j}(Da)^n
 \leq\int_{\{v<a-1\}}(Da)^n
 \leq\int_{\{v<a-1\}}(Dv)^n
 \leq(1+\epsilon)^n t_j^{-n}.
\]
Let $\epsilon\downarrow0$, then take the supremum over $a$.
The inclusions $\{z_j\}\subset K_j\subset B_j$ give all three
equalities. No outer regularity interchange is used.
\end{proof}

\begin{theorem}[A finite nonpolar linearly dominated datum]
\label{thm:r3-linear-capacity-no-solution}
The measure
\begin{equation}
 \mu=\sum_{j=1}^\infty m_j\delta_{z_j},\qquad
 m_j=t_j^{-n}-t_{j+1}^{-n}>0
 \label{eq:r3-linear-capacity-datum}
\end{equation}
is finite, supported in $K_1\Subset\Omega$, and gives no mass
to locally pluripolar sets. It satisfies, for every Borel $E$,
\begin{equation}
 \mu(E)\leq\operatorname{Cap}_{[1]}(E;\Omega).
 \label{eq:r3-linear-capacity-bound}
\end{equation}
There is no locally bounded actual classical psh solution with
ordinary boundary limit zero and $(Du)^n=\mu$.
For every $\alpha>0$ there is no finite constant $A$ for which
$\mu(E)\leq A\operatorname{Cap}_{[1]}(E;\Omega)^{1+\alpha}$
holds for all Borel $E$.
\end{theorem}

\noindent\textbf{Status: \texttt{NEW\_PROVED}} counterexample.
The assertion that finite nonpolar data and linear capacity
domination suffice for bounded zero-boundary solvability is
\texttt{FALSE} in every dimension.

\begin{proof}
The total mass is $t_1^{-n}$. Every $z_j$ is divisorial and
has positive singleton capacity, so an ambient locally polar
set contains none of these points. The origin has zero mass.
If $E$ contains any atoms, let $j$ be its least atom index.
Then
\[
 \mu(E)\leq\sum_{l\geq j}m_l=t_j^{-n}
       =\operatorname{Cap}_{[1]}(\{z_j\};\Omega)
       \leq\operatorname{Cap}_{[1]}(E;\Omega).
\]
Otherwise $\mu(E)=0$. Moreover
$\mu(K_j)=\operatorname{Cap}_{[1]}(K_j;\Omega)=t_j^{-n}$,
which disproves every uniform higher power bound.

We compute the finite partial solutions, rather than assume
weak stability for their data. Put
\[
 A_j^{(N)}=\left(\sum_{k=j}^N m_k\right)^{1/n}
 \quad(1\leq j\leq N),\qquad A_{N+1}^{(N)}=0
\]
and
\begin{equation}
 u_N=\sum_{j=1}^N
       (A_j^{(N)}-A_{j+1}^{(N)})\max(s,-t_j).
 \label{eq:r3-radial-partial-solution}
\end{equation}
It is continuous classical, has ordinary zero boundary values,
and is bounded for each fixed $N$. Its scalar radial slope
between successive breaks is $A_j^{(N)}$.
At $z_j$ the centered local tropical function, up to an additive
constant, is
\[
 \max\{A x_1,\ldots,A x_n,B x_1,\ldots,B x_n\},
 \qquad A=A_j^{(N)},\quad B=A_{j+1}^{(N)}.
\]
The convex hull of the slopes is
$\{\xi_l\geq0,\ B\leq\sum_l\xi_l\leq A\}$, with volume
$(A^n-B^n)/n!$. Its original top mass is therefore
$A^n-B^n=m_j$.

Here the calibration can be verified directly by smoothing the
finite maximum with log-sum-exp. Its Hessian is positive definite
when the slope hull has interior. For any interior slope target
$\xi$, the smooth function minus $\xi\cdot x$ is coercive and
strictly convex, so its gradient is a diffeomorphism onto that
interior. Change of variables gives the displayed volume.
The original coordinate tropical integration has multiplicity
one and factor $n!$. Scaling the smoothing parameter concentrates
this mass at the original Gauss point. At other strata the active
slope differences have rank below $n$; coordinate zero branches
are locally inactive. Smooth original integration under
refinements and bounded BT convergence therefore prove, on all
original compact tests,
\begin{equation}
 (Du_N)^n=\sum_{j=1}^N m_j\delta_{z_j}.
 \label{eq:r3-radial-partial-measure}
\end{equation}
This computation applies to any positive finite weights at these
radial points and will be used below with other weights.

With $t_0=0$, summation by parts gives the actual origin value
\[
 -u_N(o)=\sum_{j=1}^N A_j^{(N)}(t_j-t_{j-1}).
\]
Since $A_j^{(N)}\geq m_j^{1/n}
 =(1-2^{-n})^{1/n}/t_j$, it follows that
\[
 -u_N(o)\geq(1-2^{-n})^{1/n}
                  \left(1+\frac{N-1}{2}\right)\longrightarrow\infty.
\]
If $u$ were a locally bounded zero-boundary solution for $\mu$,
measure order and the ordinary boundary limits would imply
$u\leq u_N$ for every $N$, because
$(Du)^n\geq(Du_N)^n$. This contradicts the finite actual value
$u(o)$. The comparison is valid on compact gaps, so it does
not require subtracting infinite total masses near the boundary.
\end{proof}

% END INCLUDED SOURCE: sections/linear_capacity_obstruction


% BEGIN INCLUDED SOURCE: sections/general_modulus_obstruction
% Reconstructed after the scratch reset. Original reviewed SHA256:
% 82c69b3cb16af56d2894d025da513a65b7cbe1cb5948b142aa0a235754ae8f05.
\section{A non-Dini modulus obstruction in every dimension}
\label{sec:r3-nondini-capacity-obstruction}

Use the same fixed ordinary polydisc, radial points, and scales
$t_j=2^j\lambda$ as in Section~\ref{sec:r3-linear-capacity-obstruction}.
Write $q_j=t_j^{-1}$ and $C_j=t_j^{-n}$. For a finite nondecreasing
function $\omega:(0,\infty)\to(0,\infty)$ set
\[
 F(x)=x\,\omega(x^{1/n})^n\quad(0<x<\infty),\qquad
 F(0)=0,\qquad F(\infty)=\infty.
\]

\begin{theorem}[Failure of the universal guarantee outside the Dini range]
\label{thm:r3-nondini-modulus-obstruction}
Suppose
\begin{equation}
 \int_0^a\omega(t)\,\frac{dt}{t}=\infty
 \quad\hbox{for some }a>0.
 \label{eq:r3-nondini-hypothesis}
\end{equation}
There is a finite nonpolar positive measure $\mu$ supported in
the fixed compact $K_1=\{o,z_1,z_2,\ldots\}\Subset\Omega$ such that
every Borel $E\subset\Omega$ satisfies
\begin{equation}
 \mu(E)\leq F\bigl(\operatorname{Cap}_{[1]}(E;\Omega)\bigr),
 \label{eq:r3-nondini-growth}
\end{equation}
but $\mu$ has no locally bounded actual classical zero-boundary
Monge--Ampere solution.
\end{theorem}

\noindent\textbf{Status: \texttt{NEW\_PROVED}} counterexample,
using the already calibrated radial finite solutions. The uniform
existence guarantee for this prescribed non-Dini modulus is
\texttt{FALSE}; this is not an if-and-only-if criterion for an
individual measure.

\begin{proof}
Set
\[
 W_j=C_j\omega(q_j)^n,\qquad m_j=W_j-W_{j+1},\qquad
 \mu=\sum_{j=1}^{\infty}m_j\delta_{z_j}.
\]
Monotonicity of $\omega$ and $q_{j+1}=q_j/2$ imply
\begin{equation}
 W_{j+1}\leq2^{-n}W_j,\qquad
 m_j\geq(1-2^{-n})W_j>0.
 \label{eq:r3-modulus-tail-ratio}
\end{equation}
Since $\omega(q_j)\leq\omega(q_1)<\infty$, $W_j\to0$.
The total mass is $W_1$, and its tail from index $j$ is $W_j$.
As before it is nonpolar: all its atoms are divisorial, and the
origin has no mass.

If a Borel set contains atoms, take its least atom index $j$.
Lemma~\ref{lem:r3-radial-compact-capacity} and monotonicity of
$F$ give
\[
 \mu(E)\leq W_j=F(C_j)
 =F(\operatorname{Cap}_{[1]}(\{z_j\};\Omega))
 \leq F(\operatorname{Cap}_{[1]}(E;\Omega)).
\]
Sets containing no atom have zero mass. This proves the growth
bound, including the cases of zero or infinite capacity.

For the first $N$ weights use the explicit finite solution
\eqref{eq:r3-radial-partial-solution}, now with
$A_j^{(N)}=(\sum_{k=j}^Nm_k)^{1/n}$.
The calibrated identity \eqref{eq:r3-radial-partial-measure}
holds for these weights as well. Its origin depth obeys
\[
 \begin{split}
 -u_N(o)
 &=\sum_{j=1}^N A_j^{(N)}(t_j-t_{j-1})\\
 &\geq(1-2^{-n})^{1/n}
       \left(\omega(q_1)+\frac12\sum_{j=2}^N\omega(q_j)\right).
 \end{split}
\]
For every $j$, monotonicity gives
\[
 (\log2)\omega(q_{j+1})
 \leq\int_{q_{j+1}}^{q_j}\omega(t)\,\frac{dt}{t}
 \leq(\log2)\omega(q_j).
\]
The integral over any compact interval bounded away from zero
is finite. Hypothesis \eqref{eq:r3-nondini-hypothesis} therefore
forces $\sum_j\omega(q_j)=\infty$. The depths tend to infinity.
Any bounded zero-boundary solution for $\mu$ would lie below
every finite partial solution by measure order, contradicting
its finite actual value at the origin.
\end{proof}

\begin{corollary}[The logarithmic threshold]
\label{cor:r3-logarithmic-modulus-obstruction}
For $0<\beta\leq1$, there is such a measure with
\[
 \mu(E)\leq
 \frac{\operatorname{Cap}_{[1]}(E;\Omega)}
 {\bigl[\log(e+\operatorname{Cap}_{[1]}(E;\Omega)^{-1/n})\bigr]^{n\beta}}
\]
at finite positive capacity, yet with no bounded zero-boundary
solution. At zero and infinite capacity use the preceding
conventions.
\end{corollary}

\begin{proof}
Take $\omega(t)=[\log(e+1/t)]^{-\beta}$. It is finite,
positive, and nondecreasing, and its Dini integral diverges
exactly for $0<\beta\leq1$ in this range.
\end{proof}

Allowing a finite nonnegative nondecreasing modulus does not
enlarge the non-Dini case: if it vanished at any positive
argument, it would vanish on the whole smaller interval and
its Dini integral there would be finite. Combined with the
proved Dini sufficiency theorem, the construction gives the
sharp universal-modulus threshold in the stated fixed-domain,
compact-support class. It makes no assertion that every
individual datum outside the Dini range fails.

% END INCLUDED SOURCE: sections/general_modulus_obstruction


% BEGIN INCLUDED SOURCE: sections/universal_modulus_criterion
\section{The sharp universal capacity-modulus criterion on a polydisc}
\label{sec:r3-universal-modulus-criterion}

Fix $n\geq1$, the original finite extension $K/\mathbf Q_p$, a
uniformizer $\pi$, and the same polydisc as in the radial obstruction:
\[
 R=|\pi|^{-1},\qquad
 \Omega=\{|T_l|<R:1\leq l\leq n\},\qquad
 C(E)=\operatorname{Cap}_{[1]}(E;\Omega).
\]
Let $\omega:[0,\infty)\to[0,\infty)$ be finite, nondecreasing,
and zero at zero. Put $F(c)=c\omega(c^{1/n})^n$ for finite
$c>0$, with $F(0)=0$ and $F(+\infty)=+\infty$.

\begin{corollary}[A universal rather than individual equivalence]
\label{cor:r3-universal-modulus-criterion}
The following assertions are equivalent.
\begin{enumerate}
\item $\int_0^a\omega(t)\,dt/t<\infty$ for some $a>0$.
\item Every finite positive Radon measure $\mu$ carried by a compact
subset of $\Omega$ and satisfying
$\mu(E)\leq F(C(E))$ for every Borel $E\subset\Omega$ has
a bounded actual classical zero-boundary solution of $(Du)^n=\mu$.
\item Every measure in (2) has such a solution continuous on
$\overline\Omega$.
\end{enumerate}
Such bounded zero-boundary solutions are unique. If (1) fails,
a counterexample can be chosen to be a sum of positive atoms at
original divisorial points tending to the rigid origin, with no
mass at the origin or on any ambient locally pluripolar set.
\end{corollary}

\noindent\textbf{Status: \texttt{DIRECT\_COROLLARY}} of the
proved sufficient Dini theorem, its continuity improvement, and
the non-Dini radial construction. The quantifier is over all
measures with the specified bound, not over one individual measure.

\begin{proof}
The function
\[
 \rho=\max\{-1,\log|T_1|-\log R,\ldots,
                          \log|T_n|-\log R\}
\]
is continuous classical on the whole affine space and makes
$\Omega=\{\rho<0\}$ an extendible exhaustion domain. Therefore
(1) implies (3) by Theorems~\ref{thm:r3-dini-capacity-dirichlet}
and \ref{thm:r3-dini-solution-continuity}. Compactness of the
ordinary closure makes (3) imply (2).

If (1) fails, $\omega(t)>0$ for every $t>0$. Otherwise
$\omega(t_0)=0$ would make it vanish on $(0,t_0]$, and finite
monotonicity would make its Dini integral finite on every compact
interval away from zero. Thus the restriction of $\omega$ to
$(0,\infty)$ satisfies the hypotheses of
Theorem~\ref{thm:r3-nondini-modulus-obstruction}. Its explicit
measure on this same polydisc contradicts (2), and has the stated
nonpolar atomic form. Uniqueness is ambient pointwise comparison.
\end{proof}

The zero measure, for example, satisfies every such growth bound
and has the zero solution. Thus failure of the Dini condition
does not forbid solutions for all individual data. The assertion
also does not extend to arbitrary relative backgrounds with ambient
classical solutions, as
Theorem~\ref{thm:r3-relative-dirichlet-no-classical} shows.

% END INCLUDED SOURCE: sections/universal_modulus_criterion


% BEGIN INCLUDED SOURCE: sections/relative_representative_obstruction
\section{A nonzero polar background with no classical representative}
\label{sec:r3-relative-representative-obstruction}

Let
\[
 V=\{|x|<1,\ |y|<1\}\subset(\mathbf A_K^2)^{\an},
 \qquad Y_1=\{x=0\}\cap V,\quad Y_2=\{y=0\}\cap V,
\]
and $S=[Y_1]+[Y_2]$. This is a nonzero symmetric original-positive
separately closed $(1,1)$-current on a connected smooth space.
The remaining degree is one. All capacities retain normalized
ambient actual classical testers.

\begin{proposition}[Incompatible branch values in a completed class]
\label{prop:r3-no-relative-classical-representative}
There is $a\in\mathcal A_S^{[0,1]}(V)$ which is an actual
increasing limit of continuous normalized classical functions,
but no actual classical psh function is equal to $a$ locally
outside an outer $S$-capacity-zero set. One can take
\[
 a_j=\max\{0,1+j^{-1}\log|x|\}\uparrow a,
 \qquad a=\mathbf1_{\{x\ne0\}}.
\]
Its completed relative curvature is $D(aS)=0$.
\end{proposition}

\noindent\textbf{Status: \texttt{NEW\_PROVED}} counterexample;
the unrestricted classical-representative assertion is \texttt{FALSE}.

\begin{proof}
Each $a_j$ is continuous classical and lies in $[0,1]$. The finite
floor removes the logarithmic singularity. Because $\log|x|<0$
on $V$, the rows increase to the displayed actual Borel limit.
The proved increasing-capacity theorem places it in the normalized
relative completion on every permitted local compact/domain pair.

Let $\xi_{i,r}$ be the centered Gauss point of radius $0<r<1$
on $Y_i$. Each singleton has strictly positive relative capacity.
For example, with $\ell_r=-\log r$, the normalized tester
\[
 b_r(y)=\max\{0,1+\ell_r^{-1}\log|y|\}
\]
has $D(b_rS)=\ell_r^{-1}\delta_{\xi_{1,r}}$; its restriction
to $Y_2$ is zero. Exchanging $x,y$ gives the other singleton.
These are ordinary original-field curve calibrations. The same
positive atomic mass can be tested on a smaller local domain
containing the given point, so a local outer-null exceptional set
cannot contain any $\xi_{i,r}$.

If a classical $w$ represented the relative class, it would satisfy
$w(\xi_{1,r})=0$ and $w(\xi_{2,r})=1$ for every $r$. Restrict
$w$ to the first analytic axis. The classical maximum theorem on
its closed radius-$r$ disc gives $w(o)\leq w(\xi_{1,r})=0$.
On the other axis, $\xi_{2,r}\to o$ as $r\downarrow0$.
Upper semicontinuity gives $w(o)\geq1$. This contradiction proves
nonrepresentability without assuming that $w$ is continuous.

On $Y_1$ every row is zero; on $Y_2$ its Laplacian is the atom
$j^{-1}\delta_{\xi_{2,e^{-j}}}$. Hence $D(a_jS)\to0$ on all
original compact tests. The completed increasing calculus identifies
the limit with $D(aS)$, proving the last assertion.
\end{proof}

This is not a counterexample to compact capacity regularity: its
completed curvature is zero. Nor is $a$ a difference of two locally
bounded classical functions. For any such function $h$, the maximum
theorem and upper semicontinuity imply
$h(\xi_{i,r})\to h(o)$ on each axis. A difference must therefore
have the same limiting value on both axes, whereas $a$ has limits
zero and one. Thus this example does not refute the still open
potential quotient for differences of classical potentials.

% END INCLUDED SOURCE: sections/relative_representative_obstruction


% BEGIN INCLUDED SOURCE: sections/arbitrary_relative_dirichlet_obstruction
% Reconstructed after the scratch reset. Original reviewed SHA256:
% d67c630596611b08c8baaf9ec86d267d55b867bd1472763e1fcf1f2c47ed4ab6.
\section{A relative Dirichlet obstruction in every positive codimension}
\label{sec:r3-arbitrary-relative-dirichlet-obstruction}

Fix integers $p,q\geq1$, put $n=p+q$, and let
\[
 V=\{|T_1|<1,\ldots,|T_n|<1\}\subset(\mathbf A_K^n)^{\rm an}.
\]
Choose a uniformizer $\pi$, and write $r=|\pi|$ and
$\lambda=-\log r>0$. The actual function
\[
 G=\max\{\log|T_1|,\ldots,\log|T_n|,\log r\}
\]
is bounded continuous classical psh on $V$, lies between
$-\lambda$ and zero, and has ordinary boundary limit zero.
Let
\[
 \begin{split}
 Z_1&=\{T_1=T_{q+2}=\cdots=T_n=0\}\cap V,\\
 Z_2&=\{T_2=T_{q+2}=\cdots=T_n=0\}\cap V ,
 \end{split}
\]
with the additional equations absent when $p=1$. These are smooth
$q$-dimensional coordinate polydiscs. Write $i_i:Z_i\hookrightarrow V$
and $S=[Z_1]+[Z_2]$. This is a nonzero symmetric original-positive
separately closed $(p,p)$-current with connected support.

Let $z_i$ be the Gauss point of the radius-$r$ closed polydisc in
the $q$ free coordinates of $Z_i$. Both are points of the original
$K$-analytic space. They lie outside $Z_1\cap Z_2$ and are
divisorial intrinsically on $Z_i$. Put
\[
 \mu=\delta_{z_1}+2^q\delta_{z_2}.
\]
All statements below concern all original compact smooth tests.

\begin{theorem}[A bounded matching subsolution does not suffice]
\label{thm:r3-relative-dirichlet-no-classical}
There is no $u\in\mathcal C_b(V)$ with ordinary boundary limit zero
and
\[
 T_{S,q}(u,\ldots,u)=\mu.
\]
Nevertheless the bounded continuous classical function $w=2G$ has
the same boundary values and satisfies
\begin{equation}
 T_{S,q}(w,\ldots,w)
       =2^q\delta_{z_1}+2^q\delta_{z_2}\geq\mu.
 \label{eq:r3-relative-matching-subsolution}
\end{equation}
For every $\alpha>0$, all Borel $E\subset V$ satisfy
\begin{equation}
 \mu(E)\leq A_\alpha\operatorname{Cap}_S(E;V)^{1+\alpha},
 \qquad A_\alpha=(1+2^q)\lambda^{q(1+\alpha)}.
 \label{eq:r3-relative-power-datum}
\end{equation}
\end{theorem}

\noindent\textbf{Status: \texttt{NEW\_PROVED}} counterexample.
Thus the arbitrary-background classical bounded-subsolution
Dirichlet assertion is \texttt{FALSE}, even with this power
capacity bound in every remaining degree $q\geq1$.

\begin{proof}
On $Z_i$, let $g_i=G|_{Z_i}$. The original tropical calibration of
the radial maximum gives
\begin{equation}
 (Dg_i)^q=\delta_{z_i}.
 \label{eq:r3-relative-branch-measure}
\end{equation}
Indeed near the radius-$r$ Gauss point its slope polytope is
$\operatorname{conv}(0,e_1,\ldots,e_q)$ of volume $1/q!$;
the original CLD top-product normalization is $q!$ and the
coordinate tropical multiplicity is one. Elsewhere the slope
differences have rank below $q$. One can smooth the maximum and
pass by the proved bounded BT continuity, so the equality is an
ordinary-current equality on all original tests.

For a hypothetical $u$, analytic restriction preserves the actual
classical class. Integration-current compatibility and background
linearity give
\begin{equation}
 T_{S,q}(u,\ldots,u)
 =i_{1*}(D(u|_{Z_1}))^q+i_{2*}(D(u|_{Z_2}))^q .
 \label{eq:r3-relative-axis-splitting}
\end{equation}
Both summands are positive and carried by their respective
coordinate subspaces. Since $z_1\notin Z_2$ and $z_2\notin Z_1$,
equality with $\mu$ forces the intrinsic measures to be
$\delta_{z_1}$ and $2^q\delta_{z_2}$, respectively. In particular
there is no mass on the intersection: a positive summand cannot
cancel mass there.

The ordinary boundary of each coordinate polydisc lies in the
ordinary boundary of $V$, so both restrictions have boundary
limit zero. The proved intrinsic ambient uniqueness theorem,
which determines actual bounded classical potentials, and
\eqref{eq:r3-relative-branch-measure} imply
\[
 u|_{Z_1}=g_1,\qquad u|_{Z_2}=2g_2 .
\]
At the common rigid origin these values are $-\lambda$ and
$-2\lambda$, a contradiction.

Formula \eqref{eq:r3-relative-matching-subsolution} follows from
the same integration identity and positive homogeneity. Finally
$b=1+G/\lambda$ is an actual normalized tester on $V$ and
\[
 T_{S,q}(b,\ldots,b)
           =\lambda^{-q}(\delta_{z_1}+\delta_{z_2}).
\]
If $E$ contains an atom of $\mu$, its capacity is at least
$\lambda^{-q}$, whereas $\mu(E)\leq1+2^q$. This gives
\eqref{eq:r3-relative-power-datum}; if it contains no atom the
inequality is immediate.
\end{proof}

\begin{proposition}[An actual completed solution]
\label{prop:r3-relative-completed-dirichlet-example}
Define, with the actual value $\log0=-\infty$,
\[
 a_j=\max\{2G,\ G+j^{-1}\log|T_2|\}.
\]
These are bounded continuous classical psh functions, increase
pointwise, and have boundary limit zero. Their actual limit is
\[
 a=
 \begin{cases}
 G,&T_2\ne0,\\
 2G,&T_2=0.
 \end{cases}
\]
It belongs to the bounded relative capacitary completion
$\mathcal A_{S,b}(V)$, has boundary limit zero, and satisfies
\begin{equation}
 \mathcal T_{S,q}(a,\ldots,a)=\mu.
 \label{eq:r3-relative-completed-solution}
\end{equation}
The sequence converges locally in outer $S$-capacity on all
permitted compact/domain pairs. In contrast, $a^*=G$ and
$T_{S,q}(a^*,\ldots,a^*)=\delta_{z_1}+\delta_{z_2}\ne\mu$.
\end{proposition}

\noindent\textbf{Status: \texttt{DIRECT\_COROLLARY}} of the actual
increasing-limit calculus, integration compatibility, and intrinsic
ambient polar zero, applied to this explicit construction.

\begin{proof}
The maximum with the bounded floor handles the logarithmic pole.
It is classical by truncating the logarithm and decreasing to the
bounded maximum. It is continuous at the zero locus because
the logarithmic branch is inactive on a neighborhood of each such
point. Since $\log|T_2|\leq0$ on $V$, the sequence increases and
\[
 2G\leq a_j\leq G\leq0.
\]
These bounds give its boundary limit and the stated actual
pointwise limit. The relative increasing-limit theorem gives
membership in $\mathcal A_{S,b}$ and local all-domain capacity
convergence.

To identify the measure, restrict only the classical approximants.
On $Z_2$, every restriction is $2g_2$. On $Z_1$ their actual
increasing limit equals $g_1$ off the proper intrinsic analytic
divisor $Z_1\cap Z_2$, and equals $2g_1$ on that divisor.
The divisor has intrinsic ambient outer capacity zero. Hence the
intrinsic completed class and its measure are those of $g_1$.
Intrinsic increasing continuity gives
\[
 (D(a_j|_{Z_1}))^q\longrightarrow\delta_{z_1},\qquad
 (D(a_j|_{Z_2}))^q=2^q\delta_{z_2}.
\]
Apply \eqref{eq:r3-relative-axis-splitting} to each classical $a_j$.
Pushforward against all original tests, followed by the relative
completed limit, proves
\eqref{eq:r3-relative-completed-solution}. No pullback theorem for
an arbitrary completed Borel representative is assumed.

The complement of $\{T_2=0\}$ is dense: normal Gauss points of
positive radius approach every point of that coordinate
hyperplane. Continuity of $G$ and $a\leq G$ therefore give
$a^*=G$. The last measure formula follows from
\eqref{eq:r3-relative-branch-measure}.
\end{proof}

\paragraph{Scope.}
This example works for every positive codimension $p$ and every
remaining degree $q$. It does not refute the ambient $S=[1]$
Dirichlet theorem. It supplies a particular completed solution,
not a general arbitrary-background completed Dirichlet theorem,
and it does not define a wedge of two arbitrary currents.

% END INCLUDED SOURCE: sections/arbitrary_relative_dirichlet_obstruction


% BEGIN INCLUDED SOURCE: sections/relative_dirichlet_tv_limit
% Reconstructed after the scratch reset. Original reviewed SHA256:
% 19209703b0dd390cd7d0bd8249c5e0d6827340143d494ebd5182ca58e517a86a.
\section{Total-variation limits may require a completed solution}
\label{sec:r3-relative-dirichlet-tv-limit}

Use $p,q,V,Z_i,S,G,a_j,\mu$ from
Section~\ref{sec:r3-arbitrary-relative-dirichlet-obstruction}.
Let $\zeta_j\in Z_1$ be the Gauss point whose free coordinate
$T_2$ has radius $r^j$ and whose other $q-1$ free coordinates
have radius $r$. Thus $\zeta_1=z_1$.

\begin{proposition}[The exact measures of the approximants]
\label{prop:r3-relative-dirichlet-tv-formula}
For every integer $j\geq1$ one has the ordinary measure identity
\begin{equation}
 T_{S,q}(a_j,\ldots,a_j)
      =\mu+\frac{2^q-1}{j}\delta_{\zeta_j}.
 \label{eq:r3-relative-dirichlet-tv-formula}
\end{equation}
Consequently
\begin{equation}
 \|T_{S,q}(a_j^q)-\mu\|_{\rm TV}
 =d_{\mathrm{psh},V}(T_{S,q}(a_j^q),\mu)
 =\frac{2^q-1}{j}.
 \label{eq:r3-relative-tv-and-dual}
\end{equation}
All these measures have one common compact support in $V$.
\end{proposition}

\noindent\textbf{Status: \texttt{NEW\_PROVED}} for the exact
formula. The convergence and closure obstruction below are
\texttt{DIRECT\_COROLLARY}.

\begin{proof}
On $Z_2$, $a_j=2g_2$ and its intrinsic top measure is
$2^q\delta_{z_2}$. On $Z_1$ write $t_1=T_2$ and
$t_l=T_{l+1}$ for $2\leq l\leq q$. On its torus let
$s_l=\log|t_l|<0$ and set
\[
 h=\max(-\lambda,s_2,\ldots,s_q),\quad
 g=\max(s_1,h),\quad k_j=\max(s_1/j,h).
\]
For $q=1$, $h=-\lambda$. Since $s_1\leq s_1/j$, a direct
three-case comparison with $h$ gives
\begin{equation}
 a_j|_{Z_1}=\max(2g,g+s_1/j)=g+k_j.
 \label{eq:r3-relative-radial-two-maxima}
\end{equation}
The equality extends over the coordinate zero loci by continuity.
For $j=1$ this is $2g$, giving $2^q\delta_{z_1}$, exactly the
asserted formula when $\zeta_1=z_1$.

Suppose $j>1$. At $z_1$, all $s_l=-\lambda$, and $k_j=s_1/j$
is affine near this tropical point. Hence the mass there is the
mass of $g$, namely one. At $\zeta_j$ the tropical coordinates
are $s_1=-j\lambda$ and $s_l=-\lambda$ for $l\geq2$. In centered
coordinates $x_1=s_1+j\lambda$, $x_l=s_l+\lambda$ one has
\[
 g+k_j=-2\lambda+
       \max(0,x_2,\ldots,x_q)
       +\max(0,x_1/j,x_2,\ldots,x_q).
\]
The slope polytope at this vertex is the Minkowski sum
\begin{equation}
 P_j=\operatorname{conv}(0,e_2,\ldots,e_q)
           +\operatorname{conv}(0,e_1/j,e_2,\ldots,e_q).
 \label{eq:r3-relative-shell-polytope}
\end{equation}
For $q=1$ it is the interval $[0,1/j]$. In general its slice at
$\xi_1=t$, $0\leq t\leq1/j$, is
\[
 \xi_2,\ldots,\xi_q\geq0,\qquad
             \xi_2+\cdots+\xi_q\leq2-jt.
\]
Indeed a point of the second simplex with first coordinate $t$
has remaining coordinate sum at most $1-jt$, and the first
simplex contributes at most one; conversely any nonnegative
vector with the indicated total can be split between these
two bounds. Thus
\[
 \operatorname{vol}_q(P_j)
 =\int_0^{1/j}\frac{(2-jt)^{q-1}}{(q-1)!}\,dt
 =\frac{2^q-1}{j q!}.
\]

Here the use of the original torus normalization can be checked
without an unproved delta pullback. Smooth each maximum by its
log-sum-exp regularization. The sum of the two regularizations is
the log-sum-exp for all sums of their slope vectors. Its Hessian
is positive definite, since the slopes affinely span $\mathbf R^q$.
Its gradient maps $\mathbf R^q$ diffeomorphically onto
$\operatorname{int}P_j$: for an interior target $\xi$, the
regularization minus $\xi\cdot x$ is strictly convex and coercive,
and hence has a unique critical point. The change of variables
formula integrates its Hessian determinant to
$\operatorname{vol}_q(P_j)$. The original CLD top wedge on the
coordinate tropical chart has the factor $q!$ and multiplicity
one. Rescaling the regularization parameter concentrates the
mass at the displayed Gauss point. The smooth local computation,
the original tropical integration formula under refinements,
and bounded BT continuity identify the limit on every original
compact smooth test. Its mass is therefore $(2^q-1)/j$.

For completeness there is no other top-degree mass. Away from
the ties among $s_1$, $s_1/j$, and $h$, both terms are affine,
or one is affine and the other depends on fewer than $q$
directions, or both depend only on $s_2,\ldots,s_q$. At
$s_1=h<s_1/j$, the second term is affine, so the only possible
top vertex is the radial vertex $z_1$ already counted.
At $s_1<h=s_1/j$ with $h>-\lambda$, the constant floor is
inactive. Every active total gradient satisfies
$j\xi_1+\xi_2+\cdots+\xi_q=2$; their differences have rank
at most $q-1$. If $h=-\lambda$ but fewer than $q-1$ of the
remaining coordinates equal the floor, fewer than $q$ directions
are active. The only remaining case is $\zeta_j$.
These rank conclusions are also read from the smooth
regularizations and the same bounded limit. Near a zero of
some $t_l$, its logarithmic branch is inactive; near $t_1=0$
both maxima equal $h$. Thus there is no additional mass on
coordinate zero loci.

The intrinsic result on $Z_1$ is
$\delta_{z_1}+(2^q-1)j^{-1}\delta_{\zeta_j}$.
Integration-current compatibility and linearity in $S$ prove
\eqref{eq:r3-relative-dirichlet-tv-formula} on $V$.
The extra measure is positive; its total variation equals its
total mass, and the normalized constant tester one attains the
same psh-dual value. Finally $\zeta_j$ tends to the point of
$Z_1\cap Z_2$ with $t_1=0$ and the remaining coordinates of
radius $r$. This proves common compact support, including when
$q=1$, where the limit is the rigid origin.
\end{proof}

\begin{corollary}[Failure of classical closure in total variation]
\label{cor:r3-relative-dirichlet-not-tv-closed}
For this fixed background, top measures of uniformly bounded
continuous classical potentials with common ordinary zero
boundary values need not form a closed family, even under
total-variation convergence on a fixed compact support.
The same failure holds for strong psh-dual convergence. The
limit measure has a continuous matching subsolution and satisfies
every fixed power capacity growth condition in
\eqref{eq:r3-relative-power-datum}, but has no ambient classical
zero-boundary solution. Its actual relative completed solution
is $a$ from Proposition~\ref{prop:r3-relative-completed-dirichlet-example}.
\end{corollary}

\begin{proof}
Use \eqref{eq:r3-relative-tv-and-dual} and
Theorem~\ref{thm:r3-relative-dirichlet-no-classical}.
The common bounds and boundary values of the $a_j$, their
capacity convergence to $a$, and the completed equation were
proved in the cited proposition.
\end{proof}

% END INCLUDED SOURCE: sections/relative_dirichlet_tv_limit


% BEGIN INCLUDED SOURCE: sections/finite_carrier_completed_gluing
\section{Constructive gluing on finitely many analytic carriers}
\label{sec:r3-finite-carrier-completed-gluing}

This construction gives completed relative potentials with
independently prescribed branch operators. The hypotheses explicitly
use local \emph{outer} capacity zero; raw capacity zero is not
silently substituted for it.

\begin{proposition}[Logarithmic separation of finitely many carriers]
\label{prop:r3-finite-carrier-completed-gluing}
Let $U$ be an allowed ordinary open set and let
$S_1,\ldots,S_N$ be symmetric original-positive separately closed
currents of one bidegree $(p,p)$. Put
$S=\sum_iS_i$ and $q=n-p\geq1$. Suppose that $S_i$ is
carried by a closed set $A_i$, and there are analytic functions
$f_i$ on $U$, with $|f_i|\leq1$, such that
\[
 A_i\subset F_i:=\{f_i=0\}.
\]
Assume each $\log|f_i|$ is an allowed extended classical psh
function, and, for $i\ne l$, assume that $F_l$ has locally
zero outer $S_i$-capacity on all permitted compact/domain pairs.

Let $h_1,\ldots,h_N,w\in\mathcal C_b(U)$ satisfy
$w\leq h_i$ on $A_i$. Define
\[
 \psi_i=\sum_{l\ne i}\log|f_l|,
 \qquad
 a_j=\max\{w,h_1+j^{-1}\psi_1,\ldots,
                         h_N+j^{-1}\psi_N\}.
\]
For $N=1$ the empty sum is zero. Then $a_j$ is classical,
locally uniformly bounded, and increases to an actual Borel
function $a\in\mathcal A_{S,b}(U)$. For every $0\leq k\leq q$,
\begin{equation}
 \mathcal T_{S,k}(a,\ldots,a)
       =\sum_{i=1}^N T_{S_i,k}(h_i,\ldots,h_i)
 \label{eq:r3-finite-carrier-glued-operator}
\end{equation}
on all original compact smooth complementary tests. If the
$h_i$ and $w$ are continuous, each $a_j$ is continuous as well.
\end{proposition}

\noindent\textbf{Status: \texttt{DIRECT\_COROLLARY}} of the
proved increasing-limit and completed-calculus theorems, with
an explicit gluing construction. It is not a wedge product of
two arbitrary currents.

\begin{proof}
The functions $\psi_i$ are extended classical psh and
nonpositive. Taking their sums with $h_i$, and then the finite
maximum with the locally bounded floor $w$, gives an actual
locally bounded classical function. Moreover
\[
 w\leq a_j\leq\max(w,h_1,\ldots,h_N),
\]
so the bounds are uniform on each compact set, and $a_j$
increases with $j$. The relative increasing-limit theorem gives
all-domain local outer-capacity convergence to the actual limit
$a$, both for $S$ and separately for each fixed $S_i$.
Thus $a$ lies in every completed class required below.

At a point of $A_i\setminus\bigcup_{l\ne i}F_l$, the
penalty $\psi_i$ is finite. Every other branch penalty contains
$\log|f_i|=-\infty$ and is inactive. Hence
\[
 a=\max(w,h_i)=h_i
        \quad\hbox{on }A_i\setminus\bigcup_{l\ne i}F_l.
\]
The ordinary open set $U\setminus A_i$ has zero local outer
$S_i$-capacity: the background and every one of its BT products
are zero there by ordinary restriction, and the set is itself
an open zero-capacity neighborhood. The assumed outer-null
sets $F_l$, $l\ne i$, therefore account for all possible
discrepancies between $a$ and $h_i$ in the $S_i$-class.
Finite subadditivity gives
$a=h_i$ locally outside an outer $S_i$-capacity-zero set.
The completed calculus is invariant under exactly these changes,
so
\[
 \mathcal T_{S_i,k}(a,\ldots,a)
                 =T_{S_i,k}(h_i,\ldots,h_i).
\]

For every classical approximant $a_j$, the original recursion
is linear in its background:
\[
 T_{S,k}(a_j,\ldots,a_j)
          =\sum_{i=1}^N T_{S_i,k}(a_j,\ldots,a_j).
\]
Pass to the completed limits along this common increasing
sequence and use the preceding identities. This proves
\eqref{eq:r3-finite-carrier-glued-operator}; for $k=0$ it is
the identity $S=\sum_iS_i$.

If $w,h_i$ are continuous, the branches are continuous off
their logarithmic zero sets. Near a point where a penalty is
$-\infty$, that branch lies below the continuous finite floor
on a sufficiently small neighborhood, because the other
summands are locally bounded above. It can therefore be deleted
locally. The remaining finite maximum is continuous, proving
the last assertion.
\end{proof}

\begin{corollary}[Completed Dirichlet gluing on coordinate branches]
\label{cor:r3-coordinate-completed-dirichlet}
Let $V$ be the open unit polydisc in
$(\mathbf A_K^n)^{\an}$, and let $Z_1,\ldots,Z_N$ be
distinct coordinate subspaces of one dimension $q\geq1$ and
positive codimension $p=n-q$. Put
$S=\sum_i[Z_i]$, with all $Z_i$ restricted to $V$.
For each $i$, suppose a finite positive Radon measure $\nu_i$
on $Z_i$ has a bounded continuous classical solution $h_i$
with intrinsic ordinary boundary value zero, and suppose
\[
 (Dh_i)^q=\nu_i,\qquad
 C_iG_i\leq h_i\leq0
\]
for some finite $C_i\geq0$, where, for one fixed $0<r<1$,
\[
 G=\max(\log|T_1|,\ldots,\log|T_n|,\log r),
 \qquad G_i=G|_{Z_i}.
\]
Then the relative completed class contains an actual bounded
solution with ambient ordinary boundary limit zero of
\[
 \mathcal T_{S,q}(a,\ldots,a)
                  =\sum_i(i_i)_*\nu_i.
\]
It is an increasing limit of uniformly bounded continuous
ambient classical zero-boundary potentials. No matching of
the intrinsic values of $h_i$ at intersections is required.

In particular, this conclusion applies when each $\nu_i$ has
compact support in $Z_i$ and satisfies an intrinsic Dini
capacity bound on that coordinate polydisc.
\end{corollary}

\begin{proof}
Let $\pi_i:V\to Z_i$ be the coordinate retraction and put
$H_i=h_i\circ\pi_i$. These are bounded continuous classical
psh functions on $V$, and $H_i\leq0$. Choose
$C\geq\max(1,C_1,\ldots,C_N)$ and $w=CG$. Then
$w|_{Z_i}\leq h_i=H_i|_{Z_i}$.

For every $i$ one can choose a linear function $f_i$ which
vanishes on $Z_i$ and is not identically zero on any $Z_l$
with $l\ne i$. Indeed, the linear forms vanishing on $Z_i$
form a finite-dimensional $K$-space. For each distinct
same-dimensional coordinate space $Z_l$, the forms also
vanishing there form a proper subspace. An infinite field
is not a finite union of such proper subspaces. Scale the
chosen form so that its coefficients are integral; then
$|f_i|\leq1$ on $V$.

For $i\ne l$, the locus $\{f_l=0\}\cap Z_i$ is a proper
analytic subset of the smooth $Z_i$. It has zero local outer
ambient capacity there. For each compact subset of $V$, its
intersection with this locus consequently has outer capacity zero
with capacity domain $Z_i$, by finite localization and domain
monotonicity. Retraction compatibility
Proposition~\ref{prop:r2-retract} gives outer
$[Z_i]$-capacity zero on that compact set with domain $V$.
The fixed-domain null-class comparison
Corollary~\ref{cor:r2-fixed-domain-null-classes} gives the
same local all-domain assertion. Thus the preceding proposition
applies with carrier $A_i=Z_i$, branch function $H_i$, and
floor $CG$. It constructs continuous ambient classical $a_j$
and an actual completed limit with the required operator, by
integration-current compatibility.

Since all penalties are nonpositive,
\[
 CG\leq a_j\leq0,\qquad CG\leq a\leq0.
\]
The continuous floor $CG$ tends to zero at every ordinary
boundary point of $V$, proving the ambient boundary limits.
It also gives a uniform global bound. This argument does not
require the projected branch functions $H_i$ to vanish on
every ambient boundary face.

Finally the intrinsic unit polydiscs have the continuous
global classical exhaustions $G_i$. The proved Dini Dirichlet
theorem and its continuity improvement supply $h_i$, and their
compact-support barrier estimate supplies the finite constants
$C_i$. This verifies the last assertion using only the
already established ambient theory on each intrinsic branch.
\end{proof}

\paragraph{Scope.}
The construction is finite and its carrier intersections satisfy
an explicit relative outer-null condition. It does not identify
ambient pluripolar sets with relative null sets for an arbitrary
background. It also does not supply an ambient classical
representative when branch values are incompatible. The examples
of Sections~\ref{sec:r3-arbitrary-relative-dirichlet-obstruction}
and \ref{sec:r3-relative-dirichlet-tv-limit} show why the
completed target class is necessary even under favorable
intrinsic data hypotheses.

% END INCLUDED SOURCE: sections/finite_carrier_completed_gluing


% BEGIN INCLUDED SOURCE: sections/analytic_carrier_capacity_additivity
\section{Capacity additivity on finitely many analytic carriers}
\label{sec:r3-analytic-carrier-capacity}

This section gives an all-degree class of backgrounds for which raw
relative capacity is outer regular on compact sets.  The argument
uses separate ambient testers on the individual carriers and glues
them into one normalized ambient classical potential.  It does not
replace an actual relative representative by an ambient usc envelope.

Work over the fixed finite extension $K/\mathbf Q_p$, with its original
absolute value.  Let $\Omega$ be an ordinary open subset of an allowed
good, Hausdorff, boundaryless, quasi-smooth, pure $n$-dimensional
$K$-analytic space.  Fix $0\leq p<n$ and $q=n-p$.  All backgrounds
below are symmetric, original CLD-positive, and separately
$d'$- and $d''$-closed currents of bidegree $(p,p)$.  All current
identities are against the full original compact smooth test sheaf.
For such a background $R$, write
\[
 \mu_{R,h}=T_{R,q}(h,\ldots,h),\qquad
 C_R(E)=\sup_{\substack{h\in\mathcal C_b(\Omega)\\0\leq h\leq1}}
             \mu_{R,h}(E)
\]
for Borel $E\subset\Omega$, and
\[
 C_R^*(A)=\inf_{\substack{G\supset A\\G\text{ ordinary open in }\Omega}}
                  C_R(G)
\]
for arbitrary $A\subset\Omega$.  These are capacities on the one
fixed domain $\Omega$.  Infinite values are permitted unless a
relatively compact core is specified.

\begin{theorem}[Finite analytic-carrier additivity]
\label{thm:r3-analytic-carrier-additivity}
Let $S_1,\ldots,S_m$ be backgrounds as above.  Suppose there are
analytic functions $f_1,\ldots,f_m$ on $\Omega$ such that
\[
 |f_i|\leq1,\qquad A_i=\{f_i=0\},
\]
and every coefficient measure of $S_i$ is carried by $A_i$.
Assume, for all distinct $i,l$, that
\begin{equation}
 C_{S_i}(A_l)=0.
 \label{eq:r3-carrier-cross-null}
\end{equation}
Put $S=\sum_{i=1}^m S_i$.  Then every Borel $E\subset\Omega$
satisfies
\begin{equation}
 C_S(E)=\sum_{i=1}^m C_{S_i}(E).
 \label{eq:r3-carrier-raw-additivity}
\end{equation}
For every arbitrary subset $A\subset\Omega$ one also has
\begin{equation}
 C_S^*(A)=\sum_{i=1}^m C_{S_i}^*(A).
 \label{eq:r3-carrier-outer-additivity}
\end{equation}
\end{theorem}

\noindent\textbf{Status: \texttt{NEW\_PROVED}.}
This label applies to the simultaneous testing assertion
\eqref{eq:r3-carrier-raw-additivity}.

The outer identity \eqref{eq:r3-carrier-outer-additivity} is a
\texttt{DIRECT\_COROLLARY}.  The carrier condition is a condition
on actual coefficient measures, not an identification of an abstract
current support with an intrinsic analytic space.

\begin{proof}
Linearity of the fixed-background recursion gives, for every
normalized $h$,
\[
 \mu_{S,h}=\sum_{i=1}^m\mu_{S_i,h}.
\]
Hence $C_S(E)\leq\sum_iC_{S_i}(E)$.  We prove the reverse inequality
by constructing one tester for any chosen finite tuple of testers.

Fix $h_i\in\mathcal C_b(\Omega)$ with $0\leq h_i\leq1$, and fix
$0<\delta<1/2$.  Give analytic logarithms their actual value
$-\infty$ on their zero sets.  Put
\[
 \psi_i=\sum_{l\ne i}\log|f_l|\leq0,
 \qquad
 b_{\varepsilon,\delta}
   =\max\bigl\{0,\,
         \delta+(1-2\delta)h_i+\varepsilon\psi_i:
         1\leq i\leq m\bigr\},\qquad\varepsilon>0.
\]
An empty sum is zero.  This function belongs to
$\mathcal C_b(\Omega)$ and satisfies
\begin{equation}
 0\leq b_{\varepsilon,\delta}\leq1-\delta<1.
 \label{eq:r3-carrier-pasted-normalized}
\end{equation}
Here membership can be checked entirely with bounded classical
functions: replace each $\log|f_l|$ by
$\max\{\log|f_l|,-N\}$, take the same finite maximum, and let
$N\to\infty$.  The resulting bounded classical functions decrease
to $b_{\varepsilon,\delta}$ and have the common lower bound zero.
The accepted decreasing closure of the actual classical class gives
the assertion.  Truncated analytic logarithms, positive linear
combinations, and finite maxima are classical after every complete
extension, so this argument uses the required classical class.

For $R>0$ set
\[
 A_{i,R}=A_i\cap\bigcap_{l\ne i}\{f_l\ne0\}
                    \cap\{\psi_i>-R\}.
\]
Choose $\varepsilon R<\delta/2$ and consider the ordinary open set
\begin{equation}
 N_i=
 \bigcap_{l\ne i}\{f_l\ne0\}
       \cap\{\psi_i>-R\}
       \cap\{|f_i|<\exp(-1/\varepsilon)\}.
 \label{eq:r3-carrier-active-open}
\end{equation}
The finite sum $\psi_i$ is continuous on the first intersection,
so this is indeed ordinary open, and it contains $A_{i,R}$.
On $N_i$ the $i$th branch of the maximum is greater than
$\delta-\varepsilon R>\delta/2$.  If $l\ne i$, all logarithms
are nonpositive and $\psi_l\leq\log|f_i|$.  The $l$th branch is
therefore less than
\[
 \delta+(1-2\delta)\cdot1-1=-\delta.
\]
Consequently, as an equality of actual functions on this ordinary
open neighborhood,
\begin{equation}
 b_{\varepsilon,\delta}
      =\delta+(1-2\delta)h_i+\varepsilon\psi_i
      \quad\text{on }N_i.
 \label{eq:r3-carrier-active-branch}
\end{equation}
The bounds $0\leq h_i\leq1$ are what make this an ordinary-open
argument even when the $h_i$ are discontinuous.  No set defined by
comparing two discontinuous psh functions is declared open.

Every $f_l$, $l\ne i$, is an analytic unit on $N_i$.  Thus $\psi_i$
is original smooth and $D\psi_i=0$ there.  Ordinary locality,
positive homogeneity, and the smooth pluriharmonic invariance of
Lemma~\ref{lem:r2-harmonic-frame-invariance} imply
\begin{equation}
 \mu_{S_i,b_{\varepsilon,\delta}}|_{N_i}
       =(1-2\delta)^q\mu_{S_i,h_i}|_{N_i}.
 \label{eq:r3-carrier-local-measure}
\end{equation}
This uses only unit logarithms and the original smooth product rule
over separately closed tail currents.  It does not assume potential
independence for an arbitrary difference whose ambient current has
zero $D$.

The measure $\mu_{S_i,h_i}$ is carried by $A_i$: the background
vanishes on the ordinary open complement of $A_i$, and restriction
and recursion give vanishing of its products there.  Moreover
\eqref{eq:r3-carrier-cross-null} gives
\[
 \mu_{S_i,h_i}(A_l)=0\quad(l\ne i).
\]
The increasing union of $A_{i,R}$, as $R\to\infty$, is
$A_i\setminus\bigcup_{l\ne i}A_l$, which therefore carries the
whole measure.  Positivity, background linearity,
\eqref{eq:r3-carrier-pasted-normalized}, and
\eqref{eq:r3-carrier-local-measure} yield, for every Borel $E$,
\[
 \begin{split}
 C_S(E)
 &\geq\mu_{S,b_{\varepsilon,\delta}}(E)\\
 &\geq(1-2\delta)^q
            \sum_{i=1}^m\mu_{S_i,h_i}(E\cap A_{i,R}).
 \end{split}
\]
The current identities on the ordinary neighborhoods imply the
displayed identities on their Borel subsets by uniqueness of the
coefficient-Radon measures.  For each $R$ choose a permitted
$\varepsilon$; no convergence theorem for these changing testers
is needed.  Monotone convergence on the right as $R\to\infty$,
followed by $\delta\downarrow0$, gives
\[
 C_S(E)\geq\sum_{i=1}^m\mu_{S_i,h_i}(E).
\]
Taking independent suprema over the finite tuple $(h_1,\ldots,h_m)$
proves \eqref{eq:r3-carrier-raw-additivity}, including infinite
values.

Apply the raw identity to ordinary open sets.  Taking the infimum
over one common open neighborhood of $A$ first gives
$C_S^*(A)\geq\sum_iC_{S_i}^*(A)$.  If the right side is finite,
choose, for any $\eta>0$, ordinary opens $G_i\supset A$ with
\[
 C_{S_i}(G_i)\leq C_{S_i}^*(A)+\eta/m.
\]
The \emph{finite} intersection $G=\bigcap_iG_i$ is ordinary open,
contains $A$, and satisfies
\[
 C_S(G)=\sum_iC_{S_i}(G)
       \leq\sum_iC_{S_i}^*(A)+\eta.
\]
Let $\eta\downarrow0$.  If the right side is infinite the earlier
lower bound already proves equality.  This proves
\eqref{eq:r3-carrier-outer-additivity}.
\end{proof}

\begin{corollary}[Transfer of regularity through finitely many carriers]
\label{cor:r3-carrier-regularity-transfer}
Under Theorem~\ref{thm:r3-analytic-carrier-additivity}, fix an ordinary
core $W\Subset\Omega$.  If each $C_{S_i}$ is Borel outer regular
on $W$, the same is true of $C_S$.  If each $C_{S_i}^*$ on subsets
of $W$ is a finite Choquet capacity, the same is true of $C_S^*$.
If, in addition, each such capacity has compact inner capacitability
on Suslin subsets of $W$, so does $C_S^*$.
\end{corollary}

\noindent\textbf{Status: \texttt{DIRECT\_COROLLARY}.}

\begin{proof}
Borel outer regularity follows from the two additive formulas.
Finiteness, monotonicity, countable subadditivity, increasing-set
continuity, and decreasing-compact continuity are preserved by a
finite sum.  For compact inner capacitability, approximate each of
the finitely many summands from below by a compact $K_i\subset A$.
The finite union of the $K_i$ is a compact subset of $A$ and
simultaneously gives the required lower approximation for the sum.
\end{proof}

\begin{corollary}[The crossing coordinate supports have regular capacity]
\label{cor:r3-crossing-support-capacity}
In the notation of
Theorem~\ref{thm:r3-relative-dirichlet-no-classical}, let
$S=[Z_1]+[Z_2]$ on the ordinary unit polydisc $V$.  For every Borel
$E\subset V$ one has the exact same-domain identity
\begin{equation}
 \operatorname{Cap}_S(E;V)
   =\operatorname{Cap}_{[1]_{Z_1}}(i_1^{-1}E;Z_1)
     +\operatorname{Cap}_{[1]_{Z_2}}(i_2^{-1}E;Z_2).
 \label{eq:r3-crossing-capacity-sum}
\end{equation}
The corresponding equality holds for outer capacities of arbitrary
subsets.  On every fixed ordinary core $W\Subset V$, the relative
capacity is Borel outer regular, and its outer capacity is a finite
Choquet capacity with compact inner capacitability on Suslin subsets
of $W$.
\end{corollary}

\noindent\textbf{Status: \texttt{DIRECT\_COROLLARY}.}

\begin{proof}
Take $f_1=T_1$ and $f_2=T_2$.  The additional equations defining
$Z_i$ when $p>1$ cause no problem: the theorem only asks that
$[Z_i]$ be carried by the corresponding zero set, not that this
zero set equal $Z_i$.  The restrictions of $f_2$ to $Z_1$ and of
$f_1$ to $Z_2$ are nonzero coordinate functions.  Their zero sets
are proper intrinsic analytic divisors, whose intrinsic ambient
capacity is zero.  This follows from the proved ambient polar
theorem, Corollary~\ref{cor:r2-ambient-polar-outer-zero}, on a
countable relatively compact exhaustion of these polydiscs; every
normalized tester on the full polydisc restricts to each local
domain.  Thus the full raw intrinsic capacities of the divisors
are zero.

Each coordinate projection $V\to Z_i$ is an analytic retraction.
Proposition~\ref{prop:r2-retract} identifies raw and outer relative
capacities for $[Z_i]$ with the respective intrinsic capacities.
It verifies \eqref{eq:r3-carrier-cross-null} and, together with
Theorem~\ref{thm:r3-analytic-carrier-additivity}, proves
\eqref{eq:r3-crossing-capacity-sum} and its outer version.

For fixed $W\Subset V$, its inverse image $i_i^{-1}W$ is an ordinary
relatively compact core in the intrinsic smooth polydisc $Z_i$.
Apply the ambient exhaustion-domain result
Corollary~\ref{cor:r2-local-choquet} on each such intrinsic core,
and then Corollary~\ref{cor:r3-carrier-regularity-transfer}.
Pullback of a decreasing compact family to the closed subspace
$Z_i$ is again a decreasing compact family, and inverse image
preserves increasing unions.  For Suslin sets, intersection with
the closed coordinate support remains Suslin; compact inner
approximants on the two supports have compact images in $W$.
These observations verify every stated core property for the
ambient relative capacity.
\end{proof}

In particular, for this background and any compact $H\subset V$,
\[
 \operatorname{Cap}_S(H;V)=\operatorname{Cap}_S^*(H;V).
\]
If $0\leq\chi\leq1$ is original smooth with compact support and
$\chi=1$ near $H$, every ordinary weak limit $\nu$ of the finite
measures $\chi\,\mu_{S,h}$, $0\leq h\leq1$, satisfies
\[
 \nu(H)\leq\operatorname{Cap}_S(H;V).
\]
Indeed, for an ordinary open neighborhood $G$ of $H$ contained in
$\{\chi=1\}^{\circ}$, the weak open-set inequality gives
$\nu(H)\leq\nu(G)\leq\operatorname{Cap}_S(G;V)$; take the
infimum over such $G$ and use compact outer regularity.  The
local CLN estimate gives the required uniform finite mass for this
compactly localized family.

\paragraph{Scope.}
This is an all-$q$ positive result for an explicit class of
backgrounds carried by ambient polar analytic sets.  It coexists
with the failure of an ambient classical relative Dirichlet solution
in Theorem~\ref{thm:r3-relative-dirichlet-no-classical} and with the
change of measure under ambient usc regularization in
Proposition~\ref{prop:r3-relative-completed-dirichlet-example}.
Thus good raw capacity regularity does not by itself supply an
ambient classical representative for every relative completed
potential.  The proof supplies neither compact outer regularity
for an arbitrary fixed $S$ nor a countable-carrier additivity
theorem.  Its finite pasting and pairwise relative-capacity-null
hypothesis are substantive.

% END INCLUDED SOURCE: sections/analytic_carrier_capacity_additivity


% BEGIN INCLUDED SOURCE: sections/capacity_order_approximation
\section{Order approximation of backgrounds and countable carrier sums}
\label{sec:r3-capacity-order-approximation}

All capacities below use the same normalized actual classical testers
on the fixed ordinary domain $\Omega$ and the original CLD test sheaf.
Fix $q=n-p\geq1$.  A positive background means a symmetric,
original-positive, separately $d'$- and $d''$-closed current of
bidegree $(p,p)$.

\begin{proposition}[Uniform capacity control under order approximation]
\label{prop:r3-capacity-order-approximation}
Let $S$ and $S_N$ be positive backgrounds on $\Omega$ such that
\[
 0\leq S_N\leq S,\qquad S_N\longrightarrow S
\]
as currents against every original compact smooth complementary
test.  The sequence need not be increasing.  Fix a compact set $K_0$
and ordinary opens satisfying
\[
 K_0\subset W,\qquad \overline W\subset W_1\Subset\Omega.
\]
There are a constant $C$ and finitely many original positive compact
smooth complementary test forms $\theta_\ell$, supported in $W_1$,
which depend only on these compact stages, the degree, and the
chosen finite CLN presentations, such that
\[
 e_N=C\sum_\ell\langle S-S_N,\theta_\ell\rangle\longrightarrow0
\]
and
\begin{align}
 0\leq C_S(E)-C_{S_N}(E)&\leq e_N
       &&\text{for every Borel }E\subset K_0,\label{eq:r3-order-raw}\\
 0\leq C_S^*(A)-C_{S_N}^*(A)&\leq e_N
       &&\text{for every subset }A\subset K_0.\label{eq:r3-order-outer}
\end{align}
All capacities in these formulas are finite.  Consequently Borel
outer regularity on $K_0$ passes from all the $C_{S_N}$ to $C_S$.
The properties of being a finite Choquet capacity and of compact
inner capacitability on Suslin subsets of $K_0$ likewise pass from
the corresponding outer capacities to $C_S^*$.
\end{proposition}

\noindent\textbf{Status: \texttt{DIRECT\_COROLLARY}.}

\begin{proof}
Put $R_N=S-S_N$.  This is a positive background.  The finite-test
relative CLN estimate, with the normalized bound $0\leq h\leq1$,
gives
\begin{equation}
 C_{R_N}(W)
    \leq C\sum_\ell\langle R_N,\theta_\ell\rangle=e_N.
 \label{eq:r3-order-tail-cln}
\end{equation}
The forms and the constant are fixed before choosing $N$ or the
background: the background occurs in this estimate only through
its displayed positive pairings.  One obtains the estimate on $W$
by a finite covering of $\overline W$ and the CLN cutoffs supported
in $W_1$.  Weak convergence on the original test sheaf proves
$e_N\to0$.  The same estimate for $S$ proves finiteness on $W$.

For every normalized tester $h$, background linearity gives the
identity of positive Radon measures
\[
 \mu_{S,h}=\mu_{S_N,h}+\mu_{R_N,h}.
\]
Consequently every Borel $F\subset W$ satisfies
\[
 C_{S_N}(F)\leq C_S(F)
       \leq C_{S_N}(F)+C_{R_N}(W)
       \leq C_{S_N}(F)+e_N.
\]
This proves \eqref{eq:r3-order-raw}.

For $A\subset K_0$, the ordinary opens containing $A$ and contained
in $W$ are cofinal for each outer-capacity infimum: replace an
arbitrary open neighborhood $G$ by $G\cap W$.  Applying the last
inequalities to these ordinary opens and taking infima proves
\eqref{eq:r3-order-outer}.  There is no exchange of an infimum with
an infinite sum and no infinite intersection of open sets.

Suppose every $C_{S_N}$ is Borel outer regular on $K_0$.  For Borel
$E\subset K_0$ the two estimates give
\[
 C_S(E)\leq C_S^*(E)
       \leq C_{S_N}^*(E)+e_N
       =C_{S_N}(E)+e_N\leq C_S(E)+e_N.
\]
Let $N\to\infty$.  This proves Borel outer regularity.

Outer capacities are monotone and countably subadditive.  Uniform
convergence in \eqref{eq:r3-order-outer} preserves increasing-set
continuity and decreasing-compact continuity.  Explicitly, if
$A_j\uparrow A\subset K_0$, then for every $N$,
\[
 C_S^*(A)\leq C_{S_N}^*(A)+e_N
       =\lim_j C_{S_N}^*(A_j)+e_N
       \leq\lim_j C_S^*(A_j)+e_N;
\]
letting $N\to\infty$ proves continuity from below.  For decreasing
compacts $L_j\downarrow L\subset K_0$, the analogous inequality
\[
 \lim_j C_S^*(L_j)
       \leq\lim_j C_{S_N}^*(L_j)+e_N
       =C_{S_N}^*(L)+e_N\leq C_S^*(L)+e_N
\]
proves continuity from above.  Finiteness was already proved.

Finally, if $A\subset K_0$ is Suslin and each approximating outer
capacity has compact inner capacitability, then
\[
 C_S^*(A)\leq C_{S_N}^*(A)+e_N
       =\sup_{L\subset A\text{ compact}}C_{S_N}^*(L)+e_N
       \leq\sup_{L\subset A\text{ compact}}C_S^*(L)+e_N.
\]
Let $N\to\infty$ and use monotonicity for the reverse inequality.
If compact outer regularity also holds, the final compact terms can
be written as raw capacities.  This proves all the assertions.
\end{proof}

\begin{corollary}[Countably many separated analytic carriers]
\label{cor:r3-countable-carrier-capacity}
Let $(S_i)_{i\geq1}$ be positive backgrounds and suppose the partial
sums $P_N=\sum_{i=1}^N S_i$ converge weakly, on every original
compact test, to a locally coefficient-Radon positive background
$S$.  Assume that there are analytic functions $f_i$ on $\Omega$
with $|f_i|\leq1$ such that $S_i$ is carried by $\{f_i=0\}$ and
\[
 C_{S_i}(\{f_l=0\};\Omega)=0\qquad(i\ne l).
\]
On every compact $K_0\subset\Omega$ one then has
\begin{align*}
 C_S(E;\Omega)&=\sum_{i\geq1}C_{S_i}(E;\Omega)
       &&\text{for Borel }E\subset K_0,\\
 C_S^*(A;\Omega)&=\sum_{i\geq1}C_{S_i}^*(A;\Omega)
       &&\text{for arbitrary }A\subset K_0.
\end{align*}
The series are finite and their tails tend to zero uniformly over
the stated subsets of $K_0$.  If each component capacity has Borel
outer regularity and the stated Choquet properties on such compact
cores, the same is true of $S$.
\end{corollary}

\noindent\textbf{Status: \texttt{DIRECT\_COROLLARY}.}

\begin{proof}
Positivity of the series implies $0\leq P_N\leq S$: for fixed $N$,
$S-P_N$ is the weak limit of its positive tails.  It is separately
closed.  Apply Proposition~\ref{prop:r3-capacity-order-approximation}
to $P_N$.  For every finite $N$, Theorem
\ref{thm:r3-analytic-carrier-additivity} identifies the raw and outer
capacities of $P_N$ with the corresponding finite sums.  The uniform
errors of Proposition~\ref{prop:r3-capacity-order-approximation}
then give exactly the two asserted series identities, together with
uniform convergence of their tails.  Finite-component regularity
passes first to each $P_N$ and then to $S$ by the proposition.
\end{proof}

\paragraph{Scope.}
The additional step for a countable family is the uniform CLN tail
estimate, not an infinite logarithmic maximum or an infinite
intersection of ordinary open neighborhoods.  The order-approximation
result does not assert that an arbitrary positive closed background
admits approximants with already regular capacity.  In particular
it does not close the remaining unrestricted arbitrary-background
compact outer-regularity problem.

% END INCLUDED SOURCE: sections/capacity_order_approximation


% BEGIN INCLUDED SOURCE: sections/weak_closure_capacity_attempt
\section{What weak compactness does and does not give for outer capacity}
\label{sec:r3-weak-closure-capacity-attempt}

The following abstract reduction isolates a remaining relative-capacity
gate.  Its counterexample concerns a family of ordinary Radon measures;
it is not a realization by non-Archimedean BT measures over one fixed
background.  In particular it does not refute arbitrary-background
compact outer regularity.

Let $X$ be a compact metrizable space and let $\mathcal M$ be a
nonempty family of finite positive Radon measures with
$\sup_{\mu\in\mathcal M}\mu(X)<\infty$.  Put
$\mathcal N=\overline{\mathcal M}^{\,w}$, a weakly compact family.
For Borel $E$ and arbitrary $A\subset X$, respectively, define
\[
 c(E)=\sup_{\mu\in\mathcal M}\mu(E),\qquad
 c^{\mathrm o}(A)=\inf_{\substack{G\supset A\\G\text{ open}}}c(G),
 \qquad
 \widehat c(A)=\sup_{\nu\in\mathcal N}\nu^*(A).
\]
Here $\nu^*$ is the Radon outer measure, equivalently its open-cover
outer regularization.  The open outer capacity $c^{\mathrm o}$ must
not be confused with a Borel-envelope extension.

\begin{proposition}[The weakly closed measure-family capacity]
\label{prop:r3-weak-closure-capacity-facts}
The function $\widehat c$ is finite, monotone, countably subadditive,
continuous from below on arbitrary increasing subsets, and continuous
from above on decreasing compact sets.  Moreover
\begin{align}
 \widehat c(G)&=c(G)&&\text{for open }G,\label{eq:r3-closure-open}\\
 \widehat c(K)&=c^{\mathrm o}(K)&&\text{for compact }K.\label{eq:r3-closure-compact}
\end{align}
For Borel $E$ one has
\[
 c(E)\leq\widehat c(E)\leq c^{\mathrm o}(E),\qquad
 \widehat c(E)=\sup_{K\subset E\text{ compact}}c^{\mathrm o}(K).
\]
The latter compact inner-capacitability identity also holds for Suslin
subsets, with $\widehat c$ defined above on arbitrary subsets.
\end{proposition}

\noindent\textbf{Status: \texttt{DIRECT\_COROLLARY}} of elementary
Radon regularity, weak compactness, and the Choquet capacitability
theorem.  This is not a new analytic estimate for BT products.

\begin{proof}
Finiteness, monotonicity, and countable subadditivity follow directly
from those of the individual finite Radon outer measures.  We spell
out continuity from below on arbitrary, possibly nonmeasurable sets.
For a fixed finite Radon measure $\nu$ and $A_j\uparrow A$, choose
Borel hulls $B_j\supset A_j$ with $\nu(B_j)=\nu^*(A_j)$, and put
$H_j=\bigcap_{k\geq j}B_k$.  Then $A_j\subset H_j$, the Borel sets
$H_j$ increase, and $\nu(H_j)=\nu^*(A_j)$.  Thus
\[
 \nu^*(A)\leq\nu\!\left(\bigcup_jH_j\right)
             =\lim_j\nu^*(A_j)\leq\nu^*(A).
\]
Taking suprema commutes with this increasing supremum, so
$\widehat c(A_j)\uparrow\widehat c(A)$.

Let $K_j\downarrow K$ be compact.  Choose
$\nu_j\in\mathcal N$ with
$\nu_j(K_j)\geq\widehat c(K_j)-1/j$ and pass to a weakly convergent
subsequence with limit $\nu\in\mathcal N$.  For every fixed $l$,
eventually $K_j\subset K_l$, and the closed-set weak inequality gives
\[
 \lim_j\widehat c(K_j)\leq\limsup_j\nu_j(K_l)\leq\nu(K_l).
\]
Let $l\to\infty$.  The result is at most $\nu(K)\leq\widehat c(K)$;
monotonicity supplies the reverse inequality.

If $\nu\in\mathcal N$, take $\mu_j\in\mathcal M$ converging weakly
to it.  For open $G$, the open-set weak inequality gives
$\nu(G)\leq\liminf_j\mu_j(G)\leq c(G)$.  The inclusion
$\mathcal M\subset\mathcal N$ gives \eqref{eq:r3-closure-open}.
It follows at once that $\widehat c(K)\leq c^{\mathrm o}(K)$.
For the reverse inequality let
$G_j=\{x:\operatorname{dist}(x,K)<1/j\}$ and choose
$\mu_j\in\mathcal M$ with $\mu_j(G_j)\geq c(G_j)-1/j$.
After passing to a weak limit $\nu\in\mathcal N$, compare $G_j$
with each fixed compact neighborhood
$F_l=\{x:\operatorname{dist}(x,K)\leq1/l\}$.  The same closed-set
inequality yields
\[
 c^{\mathrm o}(K)=\lim_jc(G_j)\leq\nu(F_l).
\]
Let $l\to\infty$ to get $c^{\mathrm o}(K)\leq\nu(K)$, proving
\eqref{eq:r3-closure-compact}.  Empty sets are immediate.

The two inequalities for Borel $E$ follow from the inclusion of
families and from open covers.  For Borel inner regularity, only
two suprema are interchanged:
\[
 \widehat c(E)
 =\sup_{\nu\in\mathcal N}\sup_{K\subset E\text{ compact}}\nu(K)
 =\sup_{K\subset E\text{ compact}}\widehat c(K).
\]
For Suslin sets the established finite Choquet axioms imply the
same inner-capacitability formula by the Choquet theorem, in the
classical form recalled in \cite[\S8]{BT82}.  Suslin means analytic
in descriptive set theory, not an analytic subspace over $K$.
\end{proof}

\begin{proposition}[The exact missing outer-regularity implication]
\label{prop:r3-outer-closure-equivalence}
The ordinary-open extension $c^{\mathrm o}$ is finite, monotone,
countably subadditive, outer regular by definition, and continuous
from above on decreasing compact sets.  The following are equivalent:
\begin{enumerate}
\item $c^{\mathrm o}$ is continuous from below on all increasing subsets;
\item $c^{\mathrm o}(E)=\widehat c(E)$ for every Borel $E$.
\end{enumerate}
Under either condition the equality also holds for every Suslin set.
In contrast, the sharp compact weak-closure bound
\[
 \nu(K)\leq c(K)\quad(\nu\in\mathcal N,\ K\text{ compact})
\]
implies $c(E)=\widehat c(E)$ for Borel $E$, but by itself does not
imply $c(E)=c^{\mathrm o}(E)$ for Borel $E$.
\end{proposition}

\begin{proof}
The elementary outer-capacity properties follow by taking open covers;
for countable subadditivity choose errors summing to a prescribed
$\varepsilon$.  If $K_j\downarrow K$ are compact, every open
neighborhood of $K$ contains some $K_j$, which proves compact
continuity from above.

Under the first condition, $c^{\mathrm o}$ is itself a finite Choquet
capacity.  Its compact inner-capacitability on Borel sets, together
with \eqref{eq:r3-closure-compact}, gives
\[
 c^{\mathrm o}(E)=\sup_{K\subset E\text{ compact}}c^{\mathrm o}(K)
                 =\widehat c(E).
\]
Conversely suppose the second condition.  For arbitrary
$A_j\uparrow A$, choose Borel sets $B_j\supset A_j$ with
$c^{\mathrm o}(B_j)=c^{\mathrm o}(A_j)$.  Such hulls are obtained
by taking countable intersections of nearly minimizing open covers;
they are Borel and are not asserted open.  Put
$H_j=\bigcap_{k\geq j}B_k$ and $H=\bigcup_jH_j$.
Then $A_j\subset H_j$, $c^{\mathrm o}(H_j)=c^{\mathrm o}(A_j)$,
and $A\subset H$.  All $H_j,H$ are Borel.  Therefore
\[
 c^{\mathrm o}(A)\leq c^{\mathrm o}(H)=\widehat c(H)
 =\lim_j\widehat c(H_j)=\lim_jc^{\mathrm o}(A_j).
\]
The reverse inequality is monotonicity.  Suslin equality follows
from the same Choquet inner-capacitability formulas.

Finally the sharp compact bound and
\eqref{eq:r3-closure-compact} give $c(K)=\widehat c(K)$ on compacts.
Both raw Borel capacities have compact inner regularity, so they
agree on Borel sets.  The last assertion is proved false by the
following explicit example.
\end{proof}

\begin{example}[A weakly compact family with a Borel outer-capacity jump]
\label{ex:r3-vertical-fiber-capacity}
Take $X=[0,1]^2$ and
\[
 \mathcal M=\{\mu_x:x\in[0,1]\},\qquad
 \mu_x=\delta_x\otimes\mathcal L^1|_{[0,1]}.
\]
This family is weakly compact, and its raw capacity is
$c(E)=\sup_x\mathcal L^1(E_x)$ for Borel $E$.  Define
\[
 E=\bigl(\{0\}\times[0,\tfrac12]\bigr)
    \ \cup\ \bigcup_{j\geq1}
               \bigl(\{1/j\}\times(\tfrac12,1]\bigr).
\]
Then
\[
 c(E)=\widehat c(E)=\tfrac12,
 \qquad c^{\mathrm o}(E)=1.
\]
The set $E$ is both $F_\sigma$ and $G_\delta$.  There are compact
$K_l\uparrow E$ with $c^{\mathrm o}(K_l)=1/2$ for every $l$.
The same example remains valid after replacing $\mathcal M$ by its
weakly compact convex hull.
\end{example}

\noindent\textbf{Status: \texttt{FALSE}} for the proposed abstract
implication ``weakly compact measure family implies ordinary-open
Borel outer regularity.''  Its use to close the BT capacity problem
is a \texttt{FAILED\_ATTEMPT}.

\begin{proof}
For every continuous function on the compact square, its fiber
integral depends continuously on $x$.  Thus $x\mapsto\mu_x$ is
weakly continuous and its image is compact.  All nonempty fibers
of $E$ have length $1/2$, proving the raw and weakly closed values.

Let $G$ be any relatively open subset of the square containing $E$.
It contains the compact set $\{0\}\times[0,1/2]$, so for some
$\eta>0$ it contains
$[0,\eta)\times[0,1/2+\eta)$.  For sufficiently large $j$,
the remainder $\{1/j\}\times(1/2,1]$ is already contained in $E$.
Hence the entire fiber at $1/j$ is contained in $G$, and $c(G)=1$.
This proves $c^{\mathrm o}(E)=1$.

Put
\[
 K_l=\bigl(\{0\}\times[0,\tfrac12]\bigr)
       \cup\bigcup_{1\leq j\leq l}
          \bigl(\{1/j\}\times[\tfrac12+1/(l+2),1]\bigr).
\]
These are increasing compacts with union $E$.  Their raw capacities
are $1/2$, and \eqref{eq:r3-closure-compact} gives
$c^{\mathrm o}(K_l)=1/2$.  This is the claimed failure of continuity
from below, even along an increasing sequence of compact sets.

The displayed union makes $E$ an $F_\sigma$.  Let
$D=\{0\}\cup\{1/j:j\geq1\}$, a closed subset of $[0,1]$.
The complement of $E$ is the union of
\[
 ([0,1]\setminus D)\times[0,1],\qquad
 \{0\}\times(\tfrac12,1],\qquad
 \bigcup_{j\geq1}\{1/j\}\times[0,\tfrac12].
\]
Each is $F_\sigma$, so $E$ is also $G_\delta$.

The weakly compact convex hull is
$\{\tau\otimes\mathcal L^1:\tau\text{ a probability on }[0,1]\}$.
For Borel $E$, integration of the fiber lengths against $\tau$
has supremum exactly their pointwise supremum, since Dirac
probabilities are allowed.  Thus the raw capacity on all Borel
sets, and hence its ordinary-open extension, is unchanged.
\end{proof}

The family need not even give a strongly subadditive capacity.
For example, let $F_1=\{0,1\}\times[0,1/2]$ and
$F_2=(\{0\}\times[1/2,1])\cup(\{1\}\times[0,1/2])$.
Then $c(F_1)=c(F_2)=1/2$, whereas
$c(F_1\cup F_2)=1$ and $c(F_1\cap F_2)=1/2$.

\paragraph{Application and limitation for the present project.}
For a compactly localized family of normalized BT measures, the
relative CLN estimate supplies the finite mass hypothesis.  Weak
closure then gives precisely the identities on opens and compacts
proved above.  It does not identify the resulting inner Choquet
capacity with the original ordinary-open outer capacity on Borel
sets.  The abstract example is not asserted to satisfy any BT
maximum or locality structure, and is not a counterexample to an
arbitrary-$S$ theorem.  A BT-specific approximation, contact, or
separation argument is still required for the unrestricted problem.
The analytic-carrier results in this round provide such arguments
for a concrete class; they do not rely on the disproved shortcut.

% END INCLUDED SOURCE: sections/weak_closure_capacity_attempt


% BEGIN INCLUDED SOURCE: sections/round3_status
\section{Current theorem correspondence and the advance beyond R2}
\label{sec:r3-status}

Section~\ref{sec:comparison} gives the original theorem-by-theorem
comparison with BT76/82, Demailly, CLD12/25, GK, GR, and Wu.
Section~\ref{sec:r2-status} records the frozen second-round state.
This section states the present assessment. An old open entry is
closed only with the hypotheses of the specific theorem cited below.

The original finite extension $K/\mathbf Q_p$, its absolute value,
and all original compact smooth complementary CLD tests are retained.
An arbitrary background means a fixed symmetric current in the
original positive cone, separately $d'$- and $d''$-closed.
Put $q=n-p$. The all-degree relative energy and top-capacity
statements have $q\geq1$. They do not construct the wedge of two
arbitrary positive closed currents.

\subsection{What was already complete before this continuation}

R and B are \texttt{INTERNAL\_INPUT}; their classical regularization,
original coefficient-Radon control and background-$[1]$ mixed calculus
are not discoveries of this round. A2's geometric adelic nefness
theorem is also complete under its stated global hypotheses.

R1 already established the seven requested fixed-background bounded
BT properties, including symmetry and weighted joint decreasing
continuity. Its finite-test relative CLN estimate is a
\texttt{DIRECT\_COROLLARY} of R/B. R2 supplied the completed class
$\mathcal A_{S,b}$, actual increasing limits, polar restriction,
fixed-domain capacity comparison, ambient classical representatives,
the qualified ambient polar/negligible/capacity-zero equivalence,
ambient extremal mass formulas and local Choquet theory. Those
results remain frozen dependencies.

\subsection{The eleven directions after the third round}

\begin{longtable}{@{}>{\raggedright\arraybackslash}p{.04\textwidth}>{\raggedright\arraybackslash}p{.19\textwidth}>{\raggedright\arraybackslash}p{.40\textwidth}>{\raggedright\arraybackslash}p{.28\textwidth}@{}}
\hline
 & Direction & Current qualified result & Remaining limitation\\
\hline\endfirsthead
\hline
 & Direction & Current qualified result & Remaining limitation\\
\hline\endhead
1 & Locally bounded psh &
\texttt{PROVED}: B's actual classical class and R2's ambient
representative theorem. Proposition
\ref{prop:r3-no-relative-classical-representative} gives a nonzero
polar background whose completion has no classical representative. &
The unrestricted arbitrary-$S$ representative assertion is
\texttt{FALSE}, even on a connected smooth bidisc.\\[4pt]
2 & Positive closed currents and products &
\texttt{PROVED}: the arbitrary-fixed-$S$ bounded calculus and its
completion. New locally finite coefficient-mass zero extension across a closed locally
complete pole set is Theorem~\ref{thm:r3-zero-extension}. &
The arbitrary-background potential quotient remains \texttt{OPEN}.
Normal traces settle smooth algebraic-immersion and specified weak-closure
backgrounds, not every positive closed current.\\[4pt]
3 & CLN estimates &
The finite-test mass estimate remains a \texttt{DIRECT\_COROLLARY}.
Theorem~\ref{thm:r3-I-capacity} adds a genuinely relative
all-degree energy-to-capacity bound. &
Constants retain the finite local controls, common bound and collar
hypotheses. No global strictly psh form is assumed.\\[4pt]
4 & Symmetry and weighted continuity &
The requested decreasing theorem was already proved. New results
give simultaneous tail-convex-block compactness and strong local
measure convergence in the normalized-psh dual distance, including
completed limits. &
Nonlinear weak limits need not survive convexification. Unrestricted
varying-background weighted stability remains \texttt{FALSE}.\\[4pt]
5 & Locality and maxima &
Ordinary locality, Borel maximum locality and ambient maximality
are retained. Proposition~\ref{prop:r3-relative-contact-inequality}
adds the contact measure inequality as a direct corollary. &
Strict gap sets of discontinuous potentials are not automatically
ordinary open. Closed polar sets need not satisfy the locally
complete hypothesis needed for zero extension.\\[4pt]
6 & BT capacities &
Ambient core Choquet theory is retained. New results treat nonpolar
remaining degree one, separated analytic carriers in every degree,
and order limits with a uniform CLN tail. See
Theorems~\ref{thm:r3-nonpolar-linear-compact-regularity} and
\ref{thm:r3-analytic-carrier-additivity}, and
Proposition~\ref{prop:r3-capacity-order-approximation}. &
For entirely arbitrary $S$, sharp compact regularity and the further
Borel outer-regularity gate remain \texttt{OPEN}. Abstract weak
compactness alone does not supply them.\\[4pt]
7 & Quasicontinuity and capacity convergence &
Bounded relative quasicontinuity and monotone capacity convergence
are retained. The common-collar energy, psh-dual distance and
relative capacity determine the same bounded topology. &
The full normalized psh family is not capacity precompact.
The limit class may be $\mathcal A_{S,b}$ without an ambient
classical representative.\\[4pt]
8 & Pluripolar and negligible sets &
The qualified ambient equivalences are retained. Section
\ref{sec:r3-pluripolar-decomposition} gives a canonical polar/nonpolar
decomposition with its countable coefficient-control hypothesis. &
Relative measures may charge ambient polar analytic sets. A current
carried by a pole set need not have topological support contained
in that nonclosed set.\\[4pt]
9 & Comparison and domination &
Fixed-$S$ compact-gap comparison and q.e. domination are retained.
Theorem~\ref{thm:r3-one-sided-dini-stability} gives one-sided
ambient stability; power bounds yield a quantitative uniform
estimate for strong convergence of data. &
Ordinary weak-data convergence is insufficient, even under one
uniform power-capacity bound. General relative pointwise
domination still requires detection of its capacity ambiguity.\\[4pt]
10 & Extremal functions &
Ambient obstacle orthogonality and capacity/mass formulas are
retained. On the prescribed exhaustion domains, weighted divisorial obstacles
give a continuous mass map, a finite-data existence proof, and
active-set complementarity. &
The ambient extremal-contact formula is false for general relative
capacity, including specified smooth backgrounds. No arbitrary-$S$
orthogonality theorem is used in the existence proofs.\\[4pt]
11 & Dirichlet problems &
\texttt{PROVED}: ambient compactly supported data with a bounded
subsolution, and compactly supported Dini-capacity data on extendible
domains; the latter solutions are continuous. Specified maximal
boundary lifts give nonconstant boundary versions. A matching bounded
zero-boundary subsolution permits noncompact data. &
An arbitrary-relative classical existence theorem from a continuous
matching subsolution is \texttt{FALSE} in every $p,q\geq1$.
Completed branch gluing supplies a qualified replacement;
general domains, boundary data and completed equations remain open.\\
\hline
\end{longtable}

\subsection{Principal new proofs and their dependencies}

\paragraph{Relative energy and strong topology.}
For bounded potentials with one common smooth classical collar value,
or with the proved protected reduction to such a collar,
\[
 I_S(u,v)=\int (u-v)
          \bigl(\operatorname{MA}_S(v)-\operatorname{MA}_S(u)\bigr)
       \geq0.
\]
The mixed carré-du-champ identity and repeated Cauchy--Schwarz
transfer give
\[
 \operatorname{Cap}_S^*(\{|u-v|>\delta\};\Omega)
       \leq C\delta^{-2} I_S(u,v)^{2^{-q}}.
\]
The midpoint energy gap yields Cauchy tail convex blocks. The
normalized-psh dual distance satisfies
$d_S(u,v)\leq C I_S(u,v)^{1/2}$ and $I_S(u,v)\leq4B d_S(u,v)$.
The localized strong-convergence proof retains the cutoff derivative
term. For each pair of potentials it first fixes the already
constructed mixed background, and only then regularizes. This
prevents a circular use of an unproved varying Borel energy multiplier.

\paragraph{Relative capacity on analytic carriers.}
For $S_i$ carried by $\{f_i=0\}$, $|f_i|\leq1$, and pairwise
relative-null cross carriers, normalized testers $h_i$ are combined by
\[
 b_{\varepsilon,\delta}
  =\max\left(0,
       \delta+(1-2\delta)h_i+\varepsilon
                        \sum_{l\ne i}\log|f_l|:1\leq i\leq m\right).
\]
Explicit uniform margins make one branch active on an ordinary open
neighborhood of the relevant part of its carrier, even when $h_i$
is discontinuous. There the logarithms are smooth units with zero
curvature. This proves exact raw capacity additivity on all Borel
sets and outer additivity on arbitrary subsets. For countable sums,
the extra ingredient is the uniform CLN estimate for the positive
tail background, not an infinite maximum or an infinite open
intersection. These are actual normalized testers on the same domain.

\paragraph{Ambient nonzero Dirichlet existence.}
An extendible exhaustion gives a common exterior for the energy.
The constrained energy maximum is attained using the midpoint
capacity estimate and the ambient representative theorem. Positive
obstacle variation gives one measure inequality. The sublevel
capacity iteration gives a bound independent of the initially imposed
floor. The floor is then inactive on every compact perturbation,
and negative variation yields the other inequality. The obstacle
orthogonality is the already proved ambient theorem based on B;
no arbitrary-relative orthogonality is assumed.

For $\mu(E)\leq C(E)\omega(C(E)^{1/n})^n$, the sufficient scalar
condition is $\int_0^a\omega(t)\,dt/t<\infty$. Its dyadic estimate
gives an explicit bound and one-sided stability. Uniform convergence
of B's regularizing rows is then proved; it is not assumed from the
start. Dini's elementary uniform-convergence lemma is used only on
the already continuous ordinary boundary. The radial measure in
Theorem~\ref{thm:r3-nondini-modulus-obstruction} proves the converse
for a universal guarantee on a fixed polydisc. It does not classify
each individual measure satisfying a weaker bound.

\paragraph{A necessary change in the relative solution class.}
For two intersecting coordinate $q$-planes and $S=[Z_1]+[Z_2]$,
the datum $\delta_{z_1}+2^q\delta_{z_2}$ forces the two intrinsic
zero-boundary solutions to take incompatible values at their common
origin. Yet a continuous matching subsolution and every fixed
positive power-capacity bound are available. This proves failure of
the proposed ambient classical solution class. The explicit increasing
completed solution has continuous approximants with
\[
 \operatorname{MA}_S(a_j)
   =\mu+\frac{2^q-1}{j}\delta_{\zeta_j}.
\]
Thus the classical relative Dirichlet image is not even closed under
total variation with one common compact support and potential bound.
The completed gluing theorem handles independent intrinsic branch
solutions, without requiring incompatible intersection values to agree.
The same crossing background has regular capacity by analytic-carrier
additivity. Capacity regularity and classical representability are
therefore separate issues.

\subsection{Additional original-source pins}

El Mir's extension theorem is Theorem II.1, printed p.~15,
with proof on pp.~15--20; the restriction corollary II.5 is on
pp.~21--22. Demailly III.2.3, pp.~139--140, is the corresponding
complex template. The CLD12 support/dimension lemma 3.2.5 and
CLD25 7.1.12--13 have their particular hypotheses; they are not
silently substituted for the arbitrary-original-current extension
proved here.\cite[II.1, II.5]{ElMir84}\cite[III.2.3]{Dem12}

Ko\l odziej 1996, Theorem 1 and Proposition 2, pp.~13--15, give
the capacity-sublevel template. Its Theorem 2, p.~16, has additional
local-capacity and modulus hypotheses for existence. In the 1998
paper, Lemma 2.3.1 is on pp.~81--83; Theorem B, p.~96, supplies
bounds and continuity for an existing solution. Theorem C, p.~97,
with proof in \S3.4, pp.~103--111, is the bounded-subsolution
existence theorem, with the boundary qualification stated on p.~111.
The present non-Archimedean existence argument is supplied in full,
rather than inferred from the name of either classical theorem.
\cite[Theorems 1--2]{Kol96}\cite[Theorems B--C]{Kol98}

Dabbek--Elkhadhra 2006, Definition 2.1 and Proposition 2.2,
p.~471, with questions on p.~472, provide a precise relative-capacity
template. The follow-up source register is a scoped search, not a
claim that those classical questions are globally open today.
\cite[\S2]{DE06}

The projection discussion checks the actual topological-product
category in CLD25 \S9.1 and its tensor-density results. Mihatsch
Theorem 3.23, printed p.~58, gives refined flat skeleton maps;
Definition 2.18 gives piecewise smooth fibre integrals, and
Example 2.45(3), pp.~41--42, distinguishes balanced pushforwards
from $\delta$-forms. These do not establish the original projected-test
continuity required by the potential-quotient argument.
\cite[\S9.1--9.4]{CLD25}\cite[2.18, 2.45(3), 3.23]{Mih24}

\subsection{The smallest unresolved gates}

\paragraph{Q: potential quotients for arbitrary original backgrounds.}
For $u,v\in\Cb(U)$ with $D[u]=D[v]$ on all original tests,
prove $D((u-v)S)=0$ for every allowed $S$. A continuous representative
of the bounded classical difference would suffice. Normal traces
settle smooth algebraic immersions over the original field. The projection
route has an explicit continuous-coefficient sufficient criterion,
but no theorem covering every original projected test. An analytic
fibre product is not the topological product of its factor spaces,
and a same-field graph argument is not the all-extensions curve test.

\paragraph{C1: sharp compact control for a completely arbitrary $S$.}
For the fixed localized family of normalized relative BT measures,
show that every weak-closure measure $\nu$ satisfies
$\nu(K)\leq\operatorname{Cap}_S(K;\Omega)$ on compact cores.
The all-degree convex-block result does not preserve the original
nonlinear weak limit. The analytic-carrier and order-approximation
theorems do not assert that every positive closed current admits
approximants of those types.

\paragraph{C2: Borel outer regularity beyond compact control.}
Establish increasing-set continuity of open outer relative capacity,
or the equivalent Borel identification with the weak-closure capacity
in Proposition~\ref{prop:r3-outer-closure-equivalence}. The explicit
vertical-fibre example has a convex weakly compact measure family
and sharp compact control, but fails this outer-regularity step.
It is an abstract counterexample to that shortcut, not an actual
non-Archimedean BT counterexample.

\paragraph{D: the appropriate relative completed Dirichlet problem.}
The unrestricted ambient classical assertion is false. The replacement
must specify a completed solution class, a boundary condition which
survives its limits, and a justified uniqueness statement. The finite
analytic-branch construction gives a concrete model, not the general
existence theorem. For the ambient problem, broader domains and
arbitrary continuous boundary data still require proved barriers or
an appropriate homogeneous solution.

\paragraph{Comparison with the other calculi.}
An equality on all GK $\delta$-tests needs its own comparison theorem.
The full nonpositive, nonclosed and off-diagonal GK calculus is not a
consequence of the positive BT construction. A2 remains a completed
geometric input. Strong Green forms and general heights remain later
applications and do not replace any of Q, C1, C2 or D.

The next continuation should work on these specific gates and test
each proposed mechanism against the moving-atom, radial-modulus,
intersecting-support, vertical-fibre and analytic-product examples.
Reproving the accepted bounded calculus or auditing R/B/A2 afresh
would not address the remaining obstacles.

% END INCLUDED SOURCE: sections/round3_status

\clearpage

% BEGIN INCLUDED SOURCE: references_all
\begin{thebibliography}{BGJK21}
\bibitem[BT76]{BT76}
E.~Bedford and B.~A.~Taylor,
\emph{The Dirichlet problem for a complex Monge--Amp\`ere equation},
Invent. Math. \textbf{37} (1976), 1--44.
\url{https://www.math.stonybrook.edu/~ebedford/PapersForM655/BT1.pdf}.

\bibitem[BT82]{BT82}
E.~Bedford and B.~A.~Taylor,
\emph{A new capacity for plurisubharmonic functions},
Acta Math. \textbf{149} (1982), 1--40.
\url{https://www.math.stonybrook.edu/~ebedford/PapersForM655/BTNewCapacity.pdf}.

\bibitem[Dem12]{Dem12}
J.-P.~Demailly,
\emph{Complex Analytic and Differential Geometry},
online book, version June 21, 2012.
\url{https://www-fourier.univ-grenoble-alpes.fr/~demailly/manuscripts/agbook.pdf}.

\bibitem[CLD12]{CLD12}
A.~Chambert-Loir and A.~Ducros,
\emph{Formes diff\'erentielles r\'eelles et courants sur les espaces de Berkovich},
arXiv:1204.6277v1 (2012).
\url{https://arxiv.org/pdf/1204.6277v1}.

\bibitem[CLD25]{CLD25}
A.~Chambert-Loir and A.~Ducros,
\emph{Formes diff\'erentielles r\'eelles et courants sur les espaces de Berkovich},
arXiv:1204.6277v2 (2025), 608 pp.
\url{https://arxiv.org/pdf/1204.6277v2}.

\bibitem[GK17]{GK17}
W.~Gubler and K.~K\"unnemann,
\emph{A tropical approach to non-Archimedean Arakelov geometry},
Algebra Number Theory \textbf{11} (2017), 77--180.
The theorem numbering used here is checked in arXiv:1406.7637v3 (2016).
\url{https://arxiv.org/pdf/1406.7637v3}.

\bibitem[GR25]{GR25}
W.~Gubler and J.~Rabinoff,
\emph{Classically psh and pluriharmonic functions on Berkovich spaces},
arXiv:2506.13548v2 (September 17, 2025).
\url{https://gubler.app.uni-regensburg.de/Publikationen/2506.13548.pdf}.

\bibitem[Wu26]{Wu26}
L.~Wu,
\emph{On the local theory of non-Archimedean psh functions},
arXiv:2609.27610v1 (September 23, 2026).
\url{https://arxiv.org/pdf/2609.27610v1}.

\bibitem[BGJK21]{BGJK21}
J.~I.~Burgos Gil, W.~Gubler, P.~Jell, and K.~K\"unnemann,
\emph{Pluripotential theory for tropical toric varieties and non-Archimedean
Monge--Amp\`ere equations}, arXiv:2102.07392v1 (February 15, 2021).
The cited theorem numbers refer to this author-hosted version.
\url{https://gubler.app.uni-regensburg.de/Publikationen/2102.07392.pdf}.
\bibitem[GM19]{GM19}
W.~Gubler and F.~Martin,
\emph{On Zhang's semipositive metrics},
Doc. Math. \textbf{24} (2019), 331--372.
\url{https://ems.press/content/serial-article-files/26466}.


\bibitem[Brouwer11]{Brouwer11}
L.~E.~J.~Brouwer, \emph{\"Uber Abbildung von Mannigfaltigkeiten},
Math. Ann. \textbf{71} (1911), 97--115.
\url{https://doi.org/10.1007/BF01456931}.
\bibitem[DE06]{DE06}
K.~Dabbek and F.~Elkhadhra,
\emph{Capacit\'e associ\'ee \`a un courant positif ferm\'e},
Doc. Math. \textbf{11} (2006), 469--486.
\url{https://ems.press/content/serial-article-files/26001?nt=1}.
\bibitem[ElMir84]{ElMir84}
H.~El Mir, \emph{Sur le prolongement des courants positifs ferm\'es},
Acta Math. \textbf{153} (1984), 1--45.
\url{https://doi.org/10.1007/BF02392374}.
\bibitem[Kol96]{Kol96}
S.~Ko\l odziej, \emph{Some sufficient conditions for solvability
of the Dirichlet problem for the complex Monge--Amp\`ere operator},
Ann. Polon. Math. \textbf{65} (1996), 11--21.
\url{https://matwbn.icm.edu.pl/ksiazki/apm/apm65/apm6512.pdf}.
\bibitem[Kol98]{Kol98}
S.~Ko\l odziej, \emph{The complex Monge--Amp\`ere equation},
Acta Math. \textbf{180} (1998), 69--117.
\url{https://archive.ymsc.tsinghua.edu.cn/pacm_download/117/6558-11511_2006_Article_BF02392879.pdf}.
\bibitem[Mih24]{Mih24}
A.~Mihatsch, \emph{$\delta$-Forms on Lubin--Tate spaces},
arXiv:2112.10018v4 (7 September 2024).
\url{https://arxiv.org/pdf/2112.10018v4}.
\bibitem[FST91]{FST91}
M.~Fukushima, K.-I.~Sato, and S.~Taniguchi,
\emph{On the closable parts of pre-Dirichlet forms and the fine
supports of underlying measures}, Osaka J. Math. \textbf{28}
(1991), 517--535.
\url{https://ir.library.osaka-u.ac.jp/repo/ouka/all/6785/ojm28_03_03.pdf}.

\end{thebibliography}

% END INCLUDED SOURCE: references_all

\end{document}
